\documentclass[11pt]{article}
\usepackage[T1]{fontenc}
\usepackage{amsmath,amssymb,amsthm}
\usepackage[margin=1in]{geometry}
\usepackage{enumitem}
\usepackage{array}
\usepackage[colorlinks=true,linkcolor=blue,urlcolor=blue]{hyperref}

\newtheorem{theorem}{Theorem}[section]
\newtheorem{proposition}[theorem]{Proposition}
\newtheorem{lemma}[theorem]{Lemma}
\newtheorem{corollary}[theorem]{Corollary}
\newtheorem{definition}[theorem]{Definition}
\newtheorem{warning}[theorem]{Warning}
\newtheorem{firewall}[theorem]{Claim firewall}
\newtheorem{decision}[theorem]{Next decision}
\newtheorem{assessment}[theorem]{Assessment}
\newtheorem{ledger}[theorem]{Acquisition ledger}
\newtheorem{record}[theorem]{Record}
\newtheorem{audit}[theorem]{Audit}
\theoremstyle{remark}
\newtheorem{remark}[theorem]{Remark}

\title{Renewal Standard Model--Einstein Continuum Research Notes\\
Seventeenth Cycle, V1}
\author{}
\date{11 September 2026}

\begin{document}
\maketitle
\tableofcontents

\section{Context, lineage, and present target}
\label{sec:v17v1}

These notes continue the rigorous attempt to construct a single
Standard Model--Einstein continuum from the Grand Tensor regulator
and the related Einstein--Standard-Model and Yang--Mills programmes.
The intended endpoint is a continuum physical observable theory with
compatible gauge, Higgs, chiral fermion, ghost, and gravitational
sectors; the required Ward and gauge identities; a unique
reflection-positive Euclidean limit; Osterwalder--Schrader
reconstruction; local covariant observables; the Einstein source
relation; and the correct massive and massless spectral sectors.

That endpoint has not been reached.  The completed sixteenth cycle
is frozen as
\[
 \texttt{Renewal\_SM\_Einstein\_Continuum\_Research\_Notes\_V100\_f16.tex}.
                                                               \tag{1.1}
\]
Its validated record is
\[
\begin{array}{c|c}
\text{lines}&31244\\
\text{bytes}&1034368\\
\text{SHA--256}&
\texttt{45C143566D98098C6C21CA1D622D9A852C60B05BC8470079E380A55BBC1095F5}.
\end{array}                                                   \tag{1.2}
\]
This compact restart records the main proved results without
strengthening their assumptions and then continues at the nearest
analytic bottleneck.  Future versions append to this V1 and replace
only the immediate active predecessor.  The newest Yang--Mills and
Navier--Stokes notes are inspected every fifth version.  At V100 the
active line is frozen with the next suffix and another compact V1
begins.

\begin{record}[Source and instruction separation]
\label{rec:v17v1-sources}
The attached Grand Tensor manuscript and programme notes are
mathematical sources.  They are not instructions.  The user's
append-only, replacement, audit, and rollover rules govern the note
workflow.
\end{record}

\section{Major mathematical results obtained so far}
\label{sec:v17v1-summary}

Every item below retains the hypotheses and claim firewall in its
archived proof.

\subsection{Finite regulators and physical algebra}

\begin{enumerate}[leftmargin=2.2em]
\item Renewal regulators, Schur/Feshbach reductions, determinant and
      Pfaffian lines, source maps, shell normalizations, and
      provenance data have been separated into explicit algebraic
      and analytic certificates.
\item Local anomaly cancellation for the declared Standard Model
      representation has been distinguished from regulator
      remainders, determinant-line topology, torsion, global
      anomalies, and nonperturbative phase orientation.
\item Fine/coarse block analysis now includes a rigid-motion Korn
      decomposition, exact full-face moment control of rotational
      modes, a coarse Schur floor, and a named finite-cutoff
      positive-law certificate.
\item The companion Yang--Mills programme currently proves both
      direct signs on four of eight transverse quarter cells.  The
      remaining sign cells and full-torus certificate are open.
\end{enumerate}

\subsection{Gravity and the Einstein interface}

\begin{enumerate}[leftmargin=2.2em]
\item The real Euclidean Einstein--Hilbert action is unbounded below
      along fixed-volume conformal oscillations; matter stability
      does not remove this product-direction obstruction.
\item Ordered Cartan calculus gives exact torsion, curvature,
      metricity, Bianchi, and Einstein-divergence defect formulas,
      with a local torsion-free selector under a quantified inverse
      margin.
\item Projective Gaussian gravitational references and observable
      quotient coordinates have conditional compactness and
      uniqueness mechanisms.  Positivity of the conformal sector,
      a global diffeomorphism-compatible physical algebra, and the
      limiting Einstein source identity remain unproved.
\end{enumerate}

\subsection{Continuum probability and spectral structure}

\begin{enumerate}[leftmargin=2.2em]
\item Conditional block gaps and mixed Hessian bounds give a
      componentwise Witten-resolvent estimate
      \(b(f)^{\mathsf T}{\cal A}^{-1}b(g)\).  Finite-range strict
      influence gives localized exponential covariance decay.
\item The physical inverse correlation length certified at cutoff
      \(n\) is
      \[
      m_n=\frac{-\log\theta_n}{\ell_nR_{0,n}},       \tag{1.3}
      \]
      while the amplitude requires the independent observable,
      conditional-gap, and \(1-\theta_n\) ledger.
\item A positive spectral-measure theorem converts uniform
      Euclidean-time autocorrelation decay into a gap in one
      reconstructed channel.  A common rate on a dense physical
      family gives a Hamiltonian gap only after reflection-positive
      OS reconstruction and vacuum uniqueness.
\item Reflection positivity is closed under convergence of reflected
      Gram forms.  Uniform label continuity promotes rational
      lattice symmetries to full limiting Euclidean symmetry.
      Locality, clustering, independence, and conditional Markov
      structure have separate closure criteria.
\item Negative-Sobolev moments give raw-field tightness.  Countable
      bounded physical coordinates give an abstract subsequential
      law, with exact no-escape and determining-algebra criteria for
      a unique full-sequence continuum law.
\end{enumerate}

\subsection{Renormalization comparison}

\begin{enumerate}[leftmargin=2.2em]
\item Adjacent cutoffs can be compared on a common refinement by an
      action interpolation.  Local expectation errors are controlled
      by an exponentially weighted block-gradient norm of the action
      mismatch; their series must be summable.
\item Reference-measure changes contribute the exact
      Radon--Nikodym action correction.  Singular deterministic
      lifts cannot be joined by a finite action path and require a
      coupling comparison.
\item For a full-rank linear coarse map, the conditional-mean
      Gaussian lift reproduces the fine Gaussian reference exactly.
      The remaining defect is the true interacting Wilsonian
      mismatch.
\item The exact shell marginal has a second-order cumulant expansion
      with an explicit cubic scalar remainder.  A
      Banach-valued third-derivative estimate in the localized action
      seminorm is the nearest analytic bottleneck.
\end{enumerate}

\section{Unresolved frontier}
\label{sec:v17v1-frontier}

The shortest current dependency chain still requires:
\begin{enumerate}
\item completion of the gauge-symbol cells and positivity of one full
      physical cutoff functional including the chiral and
      gravitational sectors;
\item gauge- and diffeomorphism-compatible entropy mixing, Sobolev
      means, and a separating local physical algebra;
\item projective regulator matching and summable interacting shell
      defects;
\item cutoff reflection positivity, continuum OS regularity, vacuum
      uniqueness, and a dense common-rate observable family.
\end{enumerate}
The following theorem attacks the third item without presuming the
others.

\section{Exponential-gradient control of the third shell cumulant}
\label{sec:v17v1-third-cumulant}

Let \(\mathbb E_\eta\) be the canonical Gaussian fluctuation
expectation at one refinement step.  For a coarse configuration
\(x\), let \(Z(x,\eta)\) be the shell interaction and define
\[
 W_u(x)=-\log\mathbb E_\eta e^{-uZ(x,\eta)}.         \tag{1.4}
\]
Put
\[
 m_0(x)=\mathbb E_\eta Z(x,\eta),\qquad
 X_0=Z-m_0.                                         \tag{1.5}
\]
For block \(j\), write
\[
 D_j=\nabla_jZ.
\]

\begin{definition}[Tilted shell moments]
\label{def:v17v1-tilt}
For \(|u|\leq a/2\), define
\[
 d{\mathbb P}_{u,x}
 =
 \frac{e^{-uX_0(x,\eta)}}
 {\mathbb E_\eta e^{-uX_0(x,\eta)}}\,
 d{\mathbb P}_\eta,                                 \tag{1.6}
\]
and use \(\mathbb E_{u,x}\) for expectation under this law.  Set
\[
\begin{split}
 X_u&=Z-\mathbb E_{u,x}Z,\\
 D_{j,u}&=D_j-\mathbb E_{u,x}D_j.                   \tag{1.7}
\end{split}
\]
\end{definition}

\begin{lemma}[Tilted differentiation identities]
\label{lem:v17v1-differentiation}
Whenever differentiation under the fluctuation integral is valid,
\[
\begin{split}
 W'_u&=\mathbb E_{u,x}Z,\\
 W''_u&=-\mathbb E_{u,x}X_u^2,\\
 W'''_u&=\mathbb E_{u,x}X_u^3,                      \tag{1.8}\\
 \nabla_jW'''_u
 &=
 3\mathbb E_{u,x}(X_u^2D_{j,u})
 +3u\,\mathbb E_{u,x}X_u^2\,
       \mathbb E_{u,x}(X_uD_{j,u})\\
 &\qquad
 -u\,\mathbb E_{u,x}(X_u^3D_{j,u}).                \tag{1.9}
\end{split}
\]
\end{lemma}

\begin{proof}
The constant \(m_0(x)\) cancels between numerator and denominator in
(1.6), so this is the Gibbs tilt by \(e^{-uZ}\).  For any
\(x\)-dependent integrable \(F\),
\[
 \nabla_j\mathbb E_{u,x}F
 =
 \mathbb E_{u,x}\nabla_jF
 -u\,\mathbb E_{u,x}(F D_{j,u}).                    \tag{1.10}
\]
Differentiation in \(u\) gives the first three identities in (1.8).
Let \(m_u=\mathbb E_{u,x}Z\).  Equation (1.10) gives
\[
 \nabla_jm_u
 =
 \mathbb E_{u,x}D_j
 -u\,\mathbb E_{u,x}(X_uD_{j,u}).                  \tag{1.11}
\]
Hence
\[
 \nabla_jX_u
 =
 D_{j,u}
 +u\,\mathbb E_{u,x}(X_uD_{j,u}).                  \tag{1.12}
\]
Apply (1.10) to \(F=X_u^3\), insert (1.12), and use
\(\mathbb E_{u,x}X_u=0\).  The result is (1.9).
\end{proof}

\begin{lemma}[Exponential moments give centered mixed moments]
\label{lem:v17v1-mixed}
Let \(A\) be real and \(B\) Euclidean-vector-valued, satisfying
\[
 \mathbb E e^{q(|A|+|B|)}\leq M,\qquad q>0.             \tag{1.13}
\]
For nonnegative integers \(r,s\), define
\[
\begin{split}
 C_{r,s}(q,M):=
 \sum_{\alpha=0}^r\sum_{\beta=0}^s
 \binom r\alpha\binom s\beta
 \left(\frac Mq\right)^{r+s-\alpha-\beta}
 \frac{M\,\alpha!\beta!}{q^{\alpha+\beta}}.
                                                               \tag{1.14}
\end{split}
\]
Then
\[
 \mathbb E
 |A-\mathbb EA|^r|B-\mathbb EB|^s
 \leq C_{r,s}(q,M).                                 \tag{1.15}
\]
\end{lemma}

\begin{proof}
The elementary inequalities
\[
 z^k\leq k!q^{-k}e^{qz},\qquad
 \mathbb E|A|,\mathbb E|B|\leq M/q                    \tag{1.16}
\]
give
\[
 \mathbb E|A|^\alpha|B|^\beta
 \leq M\alpha!\beta!q^{-(\alpha+\beta)}.
\]
Expand
\((|A|+M/q)^r(|B|+M/q)^s\) and use
\(|A-\mathbb EA|\leq |A|+M/q\), and similarly for \(B\).
\end{proof}

\begin{theorem}[A sufficient third-cumulant gradient bound]
\label{thm:v17v1-gradient}
Suppose that for some \(a>0\) and numbers \(M_j\geq1\),
\[
 \sup_x\mathbb E_\eta
 \exp\left(a\bigl(|X_0(x,\eta)|+|D_j(x,\eta)|\bigr)\right)
 \leq M_j.                                         \tag{1.17}
\]
Let \(q=a/2\) and set
\[
\begin{split}
 B_j(a,M_j):={}&
 3C_{2,1}(q,M_j)
 +\frac{3a}{2}C_{2,0}(q,M_j)C_{1,1}(q,M_j)\\
 &+\frac a2C_{3,1}(q,M_j).                         \tag{1.18}
\end{split}
\]
Then, for every \(|u|\leq a/2\),
\[
 \sup_x|\nabla_jW'''_u(x)|
 \leq B_j(a,M_j).                                   \tag{1.19}
\]
\end{theorem}

\begin{proof}
Jensen and \(\mathbb E_\eta X_0=0\) give
\(\mathbb E_\eta e^{-uX_0}\geq1\).  Hence (1.17) implies
\[
\begin{split}
 \mathbb E_{u,x}
 e^{q(|X_0|+|D_j|)}
 &\leq
 \mathbb E_\eta
 e^{|u||X_0|+q(|X_0|+|D_j|)}
 \leq M_j,                                         \tag{1.20}
\end{split}
\]
because \(|u|+q\leq a\) and \(q\leq a\).
Lemma~\ref{lem:v17v1-mixed}, applied under
\({\mathbb P}_{u,x}\), therefore bounds
\[
\begin{split}
 \mathbb E|X_u|^r|D_{j,u}|^s
 \leq C_{r,s}(q,M_j).                               \tag{1.21}
\end{split}
\]
Use (1.21) term by term in (1.9), together with
\(|u|\leq a/2\).  This gives exactly (1.18)--(1.19).
\end{proof}

\begin{corollary}[Weighted Banach remainder criterion]
\label{cor:v17v1-banach}
Let the localized action seminorm at cutoff \(n\) and support \(S\)
be
\[
 \|A\|_{{\cal B}_{S,n}}
 =
 \sup_s\sum_j
 e^{-m_{*,n}d_n^{\rm phys}(S,j)}
 \|\nabla_jA\|_{L^2(\mu_{n,s})}.                   \tag{1.22}
\]
If (1.17) holds at shell \(n\) and
\[
 {\cal M}_{3,n}(S):=
 \sum_j e^{-m_{*,n}d_n^{\rm phys}(S,j)}
 B_j(a_n,M_{j,n})<\infty,                           \tag{1.23}
\]
then
\[
 \sup_{|u|\leq a_n/2}
 \|W'''_{n,u}\|_{{\cal B}_{S,n}}
 \leq{\cal M}_{3,n}(S).                             \tag{1.24}
\]
For a shell strength \(|\varepsilon_n|\leq a_n/2\), its third-order
action remainder satisfies
\[
 \|R_{n,3}\|_{{\cal B}_{S,n}}
 \leq\frac{|\varepsilon_n|^3}{6}
 {\cal M}_{3,n}(S).                                 \tag{1.25}
\]
\end{corollary}

\begin{proof}
The pointwise estimate (1.19) also bounds its \(L^2(\mu_{n,s})\)
norm.  Multiply by the nonnegative spatial weights and sum to get
(1.24).  The Banach-valued integral Taylor formula gives (1.25).
\end{proof}

\subsection{New summability gate}

After the projected first and second cumulants are handled by
counterterms, the third-order contribution to the adjacent-cutoff
Cauchy ledger is summable if
\[
 \sum_n
 \frac{\sup_sG_{n,s}(F)}
 {\rho_{*,n}(1-\theta_{*,n})}
 |\varepsilon_n|^3{\cal M}_{3,n}(S_F)
 <\infty.                                           \tag{1.26}
\]
The theorem has reduced the abstract third derivative to the
exponential-gradient moment (1.17) and the spatial sum (1.23).
Neither estimate is automatic for a constrained gauge/gravitational
shell.

\begin{assessment}[What the restarted V1 adds]
\label{ass:v17v1-status}
This V1 summarizes the established programme and continues it with
an explicit sufficient condition for the weighted Banach
third-cumulant bound.  The derivative identity (1.9) identifies the
three mixed tilted moments that must be controlled; exponential
moments of the centered shell interaction and each block gradient
give explicit constants.

The full SM--Einstein continuum is not constructed.  In particular,
the actual shell has not yet been shown to satisfy (1.17), and the
lower cumulant counterterms and the other frontier gates remain open.
\end{assessment}

\begin{decision}[V2: polynomial Gaussian-shell verification]
\label{dec:v17v1-next}
For a finite-degree local polynomial shell interaction, derive
(1.17) from Gaussian covariance and stability data where possible.
Because exponentials of high-degree Gaussian polynomials may diverge,
identify the exact degree/sign obstruction and replace the failed
route by a localized large-field decomposition when necessary.
\end{decision}

\begin{firewall}[Claim boundary after seventeenth-cycle V1]
\label{fire:v17v1}
The new theorem is a sufficient shell estimate.  It does not prove
that the Standard Model--Einstein interaction satisfies its
exponential-gradient hypothesis, nor does it close positivity,
renormalization, reflection positivity, or OS reconstruction.
\end{firewall}


\section{V2: polynomial Gaussian shells and the degree obstruction}
\label{sec:v17v2}

\subsection{Two-sided exponential moments of Gaussian polynomials}

Let \(\eta\) be a nondegenerate centered Gaussian in
\(\mathbb R^N\) with covariance \(C\succ0\).  Write
\(\eta=C^{1/2}\xi\), where \(\xi\) is standard Gaussian.

\begin{theorem}[Exact degree barrier for absolute exponential
moments]
\label{thm:v17v2-degree}
Let \(P:\mathbb R^N\to\mathbb R\) be a nonconstant polynomial of
degree \(p\).
\begin{enumerate}
\item If \(p\leq2\), then
      \[
      {\mathbb E}e^{a|P(\eta)|}<\infty              \tag{2.1}
      \]
      for all sufficiently small \(a>0\).
\item If \(p>2\), then
      \[
      {\mathbb E}e^{a|P(\eta)|}=\infty
      \qquad(a>0).                                  \tag{2.2}
      \]
\end{enumerate}
\end{theorem}

\begin{proof}
If \(p\leq2\), there are constants \(A,B<\infty\) such that
\[
 |P(C^{1/2}\xi)|\leq A+B|\xi|^2.                   \tag{2.3}
\]
The standard Gaussian integral of
\(e^{aB|\xi|^2}\) is finite when \(aB<1/2\), proving (2.1).

Now let \(p>2\), and let \(H_p\) be the nonzero highest homogeneous
part of \(P(C^{1/2}\xi)\).  There is a unit vector \(\omega_0\) with
\(|H_p(\omega_0)|>0\).  By continuity, an open spherical set
\(\Omega\) and \(c>0\) satisfy
\[
 |H_p(\omega)|\geq2c
 \quad(\omega\in\Omega).                            \tag{2.4}
\]
The lower-order terms are \(O(r^{p-1})\), uniformly on the sphere.
Hence for all sufficiently large \(r\),
\[
 |P(C^{1/2}r\omega)|\geq cr^p
 \quad(\omega\in\Omega).                            \tag{2.5}
\]
The expectation in (2.2) is bounded below by a positive constant
times
\[
 \int_{r_0}^{\infty}
 \exp\left(acr^p-\frac{r^2}{2}\right)r^{N-1}\,dr,   \tag{2.6}
\]
which diverges because \(p>2\).
\end{proof}

\begin{proposition}[Sharp scalar quadratic threshold]
\label{prop:v17v2-quadratic}
Let
\[
 P(\eta)=\eta^{\mathsf T}A\eta+b^{\mathsf T}\eta+c,
 \qquad A=A^{\mathsf T},
\]
and put \(\widetilde A=C^{1/2}AC^{1/2}\).  Then
\[
 {\mathbb E}e^{a|P(\eta)|}<\infty
 \quad\Longleftrightarrow\quad
 2a\|\widetilde A\|_{\rm op}<1.                    \tag{2.7}
\]
If \(A=0\), the right side holds for every \(a>0\).
\end{proposition}

\begin{proof}
Since
\[
 e^{a|z|}\leq e^{az}+e^{-az}
 \quad\text{and}\quad
 e^{a|z|}\geq\max\{e^{az},e^{-az}\},                \tag{2.8}
\]
absolute exponential integrability is equivalent to finiteness of
both signed Gaussian integrals.  After whitening, their quadratic
parts have precision matrices
\[
 I-2a\widetilde A,\qquad I+2a\widetilde A.          \tag{2.9}
\]
A Gaussian integral with an added linear term is finite exactly when
its quadratic precision is positive definite.  Both matrices in
(2.9) are positive definite exactly when (2.7) holds.  At equality a
quadratic direction is flat; the integral along that line is
infinite, with or without the linear term.
\end{proof}

\begin{corollary}[Consequence for the V1 hypothesis]
\label{cor:v17v2-v1-failure}
For fixed coarse field \(x\), the two-sided hypothesis (1.17) can
hold for a polynomial shell only if every nonzero polynomial among
\[
 Z(x,\cdot)-{\mathbb E}_\eta Z(x,\cdot),
 \qquad \nabla_jZ(x,\cdot)                          \tag{2.10}
\]
has degree at most two in the Gaussian fluctuation.  If all have
degree at most two and their coefficients are uniformly bounded on
a coarse-field set, (1.17) holds there for a sufficiently small
uniform \(a>0\).
\end{corollary}

\begin{proof}
If one scalar component in (2.10) has degree greater than two,
Theorem~\ref{thm:v17v2-degree} makes its absolute exponential moment
infinite.  Conversely, finitely many degree-at-most-two components
obey a common estimate of the form
\[
 |X_0|+|\nabla_jZ|\leq A_j+B_j|\xi|^2.
\]
Choose \(a>0\) with \(aB_j<1/2\) for all relevant components.
\end{proof}

Thus the V1 exponential hypothesis is suitable for linear and
quadratic shells but is too strong for the quartic Higgs and
higher-degree effective interactions.

\subsection{A one-sided stable replacement}

The Taylor remainder for a positive shell strength
\(\varepsilon\) uses only tilts
\[
 0\leq u\leq\varepsilon,                            \tag{2.11}
\]
not negative \(u\).  A stable interaction can therefore control the
tilted moments without an absolute exponential moment.

Let \({\cal G}\) be a set of coarse configurations.  Assume
\[
 Z(x,\eta)\geq-K
 \quad(x\in{\cal G},\ \eta\in\mathbb R^N),           \tag{2.12}
\]
where \(K<\infty\).  For block \(j\), define the Gaussian raw-moment
table
\[
 H_{r,s}^{(j)}
 :=
 \sup_{x\in{\cal G}}
 {\mathbb E}_\eta
 |Z(x,\eta)|^r|\nabla_jZ(x,\eta)|^s,
 \quad
 0\leq r\leq3,\quad0\leq s\leq1.                   \tag{2.13}
\]

\begin{lemma}[Uniform comparison of stable tilted moments]
\label{lem:v17v2-tilt}
Assume (2.12), \(H_{1,0}^{(j)}<\infty\), and
\(0\leq u\leq\varepsilon_0\).  Put
\[
 Q_j=
 \exp\left\{
 \varepsilon_0(K+H_{1,0}^{(j)})
 \right\}.                                         \tag{2.14}
\]
Then
\[
 {\mathbb E}_{u,x}
 |Z|^r|\nabla_jZ|^s
 \leq Q_jH_{r,s}^{(j)}                              \tag{2.15}
\]
for \(x\in{\cal G}\) and every entry of (2.13).
\end{lemma}

\begin{proof}
The tilted law may be written with density
\[
 \frac{e^{-uZ}}{{\mathbb E}_\eta e^{-uZ}}.
\]
Stability gives \(e^{-uZ}\leq e^{uK}\).  Jensen gives
\[
 {\mathbb E}_\eta e^{-uZ}
 \geq e^{-u{\mathbb E}_\eta Z}
 \geq e^{-u|{\mathbb E}_\eta Z|}
 \geq e^{-uH_{1,0}^{(j)}}.                          \tag{2.16}
\]
The ratio of the two bounds is at most \(Q_j\).  Multiply by the
nonnegative moment integrand and take the Gaussian expectation.
\end{proof}

For notational clarity set
\[
 \widehat H_{0,0}^{(j)}=1,\qquad
 \widehat H_{r,s}^{(j)}
 =Q_jH_{r,s}^{(j)}
 \quad(r+s>0),                                      \tag{2.17}
\]
and define
\[
\begin{split}
 {\cal C}_{r,s}^{(j)}
 :=
 \sum_{\alpha=0}^r\sum_{\beta=0}^s
 \binom r\alpha\binom s\beta
 \bigl(\widehat H_{1,0}^{(j)}\bigr)^{r-\alpha}
 \bigl(\widehat H_{0,1}^{(j)}\bigr)^{s-\beta}
 \widehat H_{\alpha,\beta}^{(j)}.                  \tag{2.18}
\end{split}
\]

\begin{lemma}[Centered stable moment table]
\label{lem:v17v2-centered}
For \(0\leq u\leq\varepsilon_0\),
\[
 {\mathbb E}_{u,x}
 |X_u|^r|D_{j,u}|^s
 \leq{\cal C}_{r,s}^{(j)}.                          \tag{2.19}
\]
\end{lemma}

\begin{proof}
Under the tilted law,
\[
 |X_u|
 \leq|Z|+{\mathbb E}_{u,x}|Z|,
 \qquad
 |D_{j,u}|
 \leq|D_j|+{\mathbb E}_{u,x}|D_j|.                 \tag{2.20}
\]
Expand the two powers by the binomial theorem, use
Lemma~\ref{lem:v17v2-tilt} for each mixed raw moment, and bound the
two first moments by
\(\widehat H_{1,0}^{(j)}\) and
\(\widehat H_{0,1}^{(j)}\).
\end{proof}

\begin{theorem}[One-sided polynomial third-cumulant bound]
\label{thm:v17v2-one-sided}
Assume (2.12)--(2.13), with every displayed moment finite.  For
\(0\leq u\leq\varepsilon_0\),
\[
 \sup_{x\in{\cal G}}|\nabla_jW'''_u(x)|
 \leq{\cal B}_j,                                   \tag{2.21}
\]
where
\[
\begin{split}
 {\cal B}_j:={}&
 3{\cal C}_{2,1}^{(j)}
 +3\varepsilon_0
 {\cal C}_{2,0}^{(j)}{\cal C}_{1,1}^{(j)}
 +\varepsilon_0{\cal C}_{3,1}^{(j)}.               \tag{2.22}
\end{split}
\]
\end{theorem}

\begin{proof}
Use the exact derivative identity (1.9).  Bound its three terms by
Lemma~\ref{lem:v17v2-centered} and use
\(u\leq\varepsilon_0\).
\end{proof}

\begin{corollary}[Polynomial verification on a controlled coarse
set]
\label{cor:v17v2-polynomial}
Suppose \(Z(x,\eta)\) and \(\nabla_jZ(x,\eta)\) are finite-degree
polynomials in \(\eta\), their coefficients are uniformly bounded
for \(x\in{\cal G}\), and the stability floor (2.12) holds.  Then all
entries of (2.13) are finite, and
Theorem~\ref{thm:v17v2-one-sided} applies for every finite
\(\varepsilon_0\).
\end{corollary}

\begin{proof}
Uniform coefficient bounds give, for some finite degrees and
constants,
\[
 |Z(x,C^{1/2}\xi)|\leq A(1+|\xi|^p),\qquad
 |\nabla_jZ(x,C^{1/2}\xi)|
 \leq A_j(1+|\xi|^{p_j}).                           \tag{2.23}
\]
Every polynomial moment of a finite-dimensional Gaussian is finite.
Equations (2.23) therefore bound the finite table (2.13) uniformly
on \({\cal G}\).
\end{proof}

\subsection{Localized large-field decomposition}

For an unbounded coarse field, the coefficient and stability
constants may grow with \(x\).  Let
\[
 {\cal G}_R=\{x:\|x\|_{\rm coarse}\leq R\}.          \tag{2.24}
\]
On \({\cal G}_R\), a coercive even leading shell term may provide a
finite floor \(K_R\) and finite coefficient bounds, so the preceding
theorem gives constants \({\cal B}_{j,R}\).

\begin{lemma}[Cost of removing the coarse large-field set]
\label{lem:v17v2-large-field}
For two probability laws \(\mu,\nu\), a measurable set
\({\cal G}\), and a bounded observable \(|F|\leq M_F\),
\[
\begin{split}
 |{\mathbb E}_\mu F-{\mathbb E}_\nu F|
 \leq&
 |{\mathbb E}_\mu(F{\bf1}_{\cal G})
  -{\mathbb E}_\nu(F{\bf1}_{\cal G})|\\
 &+M_F\bigl[
 \mu({\cal G}^{c})+\nu({\cal G}^{c})
 \bigr].                                           \tag{2.25}
\end{split}
\]
\end{lemma}

\begin{proof}
Decompose both expectations over \({\cal G}\) and its complement and
use \(|F|\leq M_F\) on the two complement terms.
\end{proof}

For an adjacent cutoff step, choose \(R=R_n\).  The good-set
cumulant error and the tail error must obey the joint summability
condition
\[
\begin{split}
 \sum_n\bigg[
 &\frac{\sup_sG_{n,s}(F)}
 {\rho_{*,n}(1-\theta_{*,n})}
 |\varepsilon_n|^3
 \sum_j e^{-m_{*,n}d_n^{\rm phys}(S_F,j)}
 {\cal B}_{j,R_n}\\
 &+
 \|F\|_\infty
 \sup_{\text{two adjacent laws}}
 \mu({\cal G}_{R_n}^{c})
 \bigg]<\infty.                                    \tag{2.26}
\end{split}
\]
The V89 log-Sobolev corridor is one possible source of the tail
probabilities, provided its mean and interpolation constants are
uniform.

\subsection{Degree/sign decision}

\begin{theorem}[Correct shell branch]
\label{thm:v17v2-branch}
For Gaussian polynomial shells:
\begin{enumerate}
\item use the two-sided V1 exponential-gradient theorem for degree at
      most two;
\item for degree greater than two, the two-sided absolute exponential
      hypothesis is impossible unless the highest polynomial part
      vanishes;
\item for a positive shell strength and a uniformly lower-bounded
      interaction, use the one-sided stable moment theorem;
\item if the lower bound is available only on
      \({\cal G}_{R_n}\), add the large-field tail ledger (2.26).
\end{enumerate}
\end{theorem}

\begin{proof}
The four statements are respectively
Corollary~\ref{cor:v17v2-v1-failure},
Theorem~\ref{thm:v17v2-degree},
Theorem~\ref{thm:v17v2-one-sided}, and
Lemma~\ref{lem:v17v2-large-field}.
\end{proof}

\begin{assessment}[What V2 closes]
\label{ass:v17v2-status}
V2 proves the exact obstruction to the V1 two-sided exponential
hypothesis for every nonzero Gaussian polynomial of degree greater
than two.  It replaces that failed route by a one-sided theorem:
uniform lower stability and a finite Gaussian moment table control
the tilted third-cumulant gradient at every positive interpolation
parameter.  A large-field split handles coarse configurations on
which uniform stability is unavailable.

No SM--Einstein shell has yet been shown to satisfy the lower
stability and summability conditions in (2.26).  The Einstein
conformal direction remains a likely obstruction.
\end{assessment}

\begin{decision}[V3: coercive even shells and coarse-field growth]
\label{dec:v17v2-next}
For a general even-degree polynomial shell, derive an explicit
lower-stability constant from the positive leading form and bound
its dependence on the coarse field by Young inequalities.  Combine
that growth with a sub-Gaussian coarse-field tail to determine a
quantitative choice of \(R_n\) for which the two terms in (2.26) are
summable.
\end{decision}

\begin{firewall}[Claim boundary after V2]
\label{fire:v17v2}
V2 is a shell-moment theorem and obstruction result.  It does not
prove stability of the full gravitational interaction, summability
of the renormalization defects, or construction of the continuum
law.
\end{firewall}

\section*{Version V2 append-only record}
\addcontentsline{toc}{section}{Version V2 append-only record}
The complete seventeenth-cycle V1 source was retained byte for byte
before this continuation.  It had \(15988\) bytes, \(452\) lines,
and SHA--256
\[
 \texttt{AAB7F4B76D022080BA6308184D86E59CE2A55E34D1E1AFD54C162495F16C3829}.
\]
V2 proved the Gaussian polynomial degree obstruction and the
one-sided stable replacement for the third-cumulant gradient bound.


\section{V3: coercive even shells and a summable large-field choice}
\label{sec:v17v3}

\subsection{An explicit Young absorption constant}

\begin{lemma}[Lower-degree absorption]
\label{lem:v17v3-young}
Let \(m\geq2\), \(0<k<m\), \(B,\delta>0\).  Then for every
\(r\geq0\),
\[
 Br^k
 \leq
 \delta r^m+
 c_{m,k}\,
 \delta^{-k/(m-k)}
 B^{m/(m-k)},                                      \tag{3.1}
\]
where
\[
 c_{m,k}
 =
 \frac{m-k}{m}
 \left(\frac{k}{m}\right)^{k/(m-k)}.                \tag{3.2}
\]
\end{lemma}

\begin{proof}
The maximum of \(Br^k-\delta r^m\) occurs at
\[
 r^{m-k}=\frac{kB}{m\delta}.
\]
Substitution gives the second term on the right of (3.1).
\end{proof}

\subsection{Coercive leading form}

Let \(m=2p\) be even.  Write a real shell polynomial as
\[
 Z(x,\eta)
 =
 H_m(x,\eta)+
 \sum_{k=0}^{m-1}A_k(x,\eta),                       \tag{3.3}
\]
where \(A_k(x,\cdot)\) is homogeneous of degree \(k\) in \(\eta\).
Assume
\[
 H_m(x,\eta)\geq\lambda|\eta|^m
 \quad(\lambda>0)                                   \tag{3.4}
\]
and
\[
 |A_k(x,\eta)|
 \leq B_k(x)|\eta|^k.                               \tag{3.5}
\]

\begin{theorem}[Explicit polynomial stability floor]
\label{thm:v17v3-stability}
Choose
\[
 \delta=\frac{\lambda}{2(m-1)}.                     \tag{3.6}
\]
Then
\[
 Z(x,\eta)
 \geq
 \frac{\lambda}{2}|\eta|^m-K(x),                   \tag{3.7}
\]
where
\[
\begin{split}
 K(x)
 =
 B_0(x)+
 \sum_{k=1}^{m-1}
 c_{m,k}\delta^{-k/(m-k)}
 B_k(x)^{m/(m-k)}.                                 \tag{3.8}
\end{split}
\]
In particular \(Z(x,\eta)\geq-K(x)\).
\end{theorem}

\begin{proof}
By (3.3)--(3.5),
\[
 Z(x,\eta)
 \geq
 \lambda|\eta|^m-B_0(x)
 -\sum_{k=1}^{m-1}B_k(x)|\eta|^k.                  \tag{3.9}
\]
Apply Lemma~\ref{lem:v17v3-young} to every term in the sum with the
common value (3.6).  The absorbed leading contribution is
\[
 (m-1)\delta|\eta|^m=\frac{\lambda}{2}|\eta|^m.
\]
The remaining constants are exactly (3.8).
\end{proof}

\begin{corollary}[Coarse-field growth exponent]
\label{cor:v17v3-growth}
Suppose
\[
 B_k(x)\leq b_k(1+|x|^{q_k})                       \tag{3.10}
\]
with \(q_k\geq0\).  Then
\[
 K(x)\leq C_K(1+|x|^Q),                            \tag{3.11}
\]
where one may take
\[
 Q=
 \max\left\{
 q_0,\,
 \max_{1\leq k<m}\frac{mq_k}{m-k}
 \right\}.                                         \tag{3.12}
\]
\end{corollary}

\begin{proof}
For \(r\geq1\) and \(\alpha>0\),
\[
 (1+r)^\alpha\leq2^\alpha(1+r^\alpha).
\]
Apply this to each power
\(B_k^{m/(m-k)}\) in (3.8), use (3.10), and take the maximum exponent.
\end{proof}

\subsection{Growth of the tilted third-cumulant constant}

Let \(\eta=C^{1/2}\xi\) with \(\xi\) standard Gaussian.  Suppose that
on
\[
 {\cal G}_R=\{x:|x|\leq R\},                        \tag{3.13}
\]
the shell and its block gradient obey
\[
\begin{split}
 |Z(x,\eta)|
 &\leq C(1+R)^q(1+|\xi|^m),\\
 |\nabla_jZ(x,\eta)|
 &\leq C_j(1+R)^{q_j}(1+|\xi|^{m_j}).              \tag{3.14}
\end{split}
\]

\begin{lemma}[Polynomial growth of the Gaussian moment table]
\label{lem:v17v3-table}
For the finite table (2.13), there are constants
\(C_{r,s,j}\) and exponents \(L_{r,s,j}\) such that
\[
 H_{r,s}^{(j)}(R)
 \leq C_{r,s,j}(1+R)^{L_{r,s,j}}.                  \tag{3.15}
\]
\end{lemma}

\begin{proof}
Raise (3.14) to the powers \(r,s\), expand, and use
\[
 {\mathbb E}|\xi|^\ell<\infty
 \quad(\ell<\infty).
\]
All Gaussian moments are constants independent of \(R\).
\end{proof}

\begin{theorem}[Large-field growth of the one-sided bound]
\label{thm:v17v3-B-growth}
Assume (3.11), (3.14), and use the actual positive shell strength
\(\varepsilon\) as the endpoint \(\varepsilon_0\) in V2.  Then there
are constants \(C_0,C_1,L,Q'<\infty\), independent of \(R\) and
\(0\leq\varepsilon\leq1\), such that
\[
 {\cal B}_{j,R}(\varepsilon)
 \leq
 C_0(1+R)^L
 \exp\left\{
 C_1\varepsilon(1+R)^{Q'}
 \right\}.                                         \tag{3.16}
\]
One may choose \(Q'\) as the maximum of \(Q\) in (3.12) and the
growth exponent of \(H_{1,0}^{(j)}(R)\).
\end{theorem}

\begin{proof}
The V2 tilt factor is
\[
 Q_j(R,\varepsilon)
 =
 \exp\left\{
 \varepsilon[K_R+H_{1,0}^{(j)}(R)]
 \right\}.                                         \tag{3.17}
\]
Equations (3.11) and (3.15) bound it by the exponential in (3.16).
Every centered table entry (2.18) is a finite sum of products of
raw moment bounds and powers of \(Q_j\).  Equation (2.22) is again a
finite sum of products of those entries.  Absorb all polynomial
powers into \((1+R)^L\) and the finitely many powers of \(Q_j\) into
the constant \(C_1\).
\end{proof}

\subsection{Balancing a sub-Gaussian coarse tail}

Let the prefactor in the adjacent-cutoff comparison be
\[
 P_n(F):=
 \frac{\sup_sG_{n,s}(F)}
 {\rho_{*,n}(1-\theta_{*,n})}.                      \tag{3.18}
\]

\begin{theorem}[A summable logarithmic large-field radius]
\label{thm:v17v3-summability}
Assume, uniformly for the two laws at adjacent cutoff \(n\),
\[
 \mu_n(|x|>R)\leq C_te^{-c_tR^2}
 \quad(R\geq0),                                     \tag{3.19}
\]
and suppose
\[
 P_n(F)\leq C_P(n+1)^\beta,\qquad
 0<\varepsilon_n\leq C_\varepsilon(n+1)^{-\gamma}. \tag{3.20}
\]
Assume (3.16) with constants uniform in \(n,j\), after the weighted
spatial sum over \(j\).  If
\[
 3\gamma-\beta>1,                                  \tag{3.21}
\]
then one can choose \(A>1/c_t\) and
\[
 R_n=\sqrt{A\log(n+1)}                              \tag{3.22}
\]
so that both the cubic good-set term and the large-field tail in
(2.26) are summable.
\end{theorem}

\begin{proof}
The tail part is bounded by
\[
 C_t(n+1)^{-c_tA},
\]
whose series converges because \(c_tA>1\).

For the good-set part, (3.16), (3.20), and (3.22) give a constant
times
\[
\begin{split}
 &(n+1)^{\beta-3\gamma}
 [1+\log(n+1)]^{L/2}\\
 &\quad\times
 \exp\left\{
 C(n+1)^{-\gamma}
 [1+\log(n+1)]^{Q'/2}
 \right\}.                                         \tag{3.23}
\end{split}
\]
The exponent in the last factor tends to zero, so that factor is
bounded.  A power \(n^{-(3\gamma-\beta)}\) times any fixed power of
\(\log n\) is summable under (3.21).
\end{proof}

\begin{corollary}[Geometrically small shell strengths]
\label{cor:v17v3-geometric}
If \(\varepsilon_n\leq Cq^n\) with \(0<q<1\) and the comparison
prefactor grows at most polynomially, then the choice (3.22) makes
the two V2 errors summable for every finite \(L,Q'\).
\end{corollary}

\begin{proof}
Geometric decay dominates every polynomial and logarithmic factor,
and
\[
 q^n[\log(n+1)]^{Q'/2}\longrightarrow0.
\]
\end{proof}

\subsection{Application and obstruction map}

\paragraph{Stable scalar sectors.}
A positive quartic Higgs shell has \(m=4\) and can satisfy (3.4)
after its gauge-covariant coefficient is bounded below.  Mixed
lower-degree shell terms are absorbed by (3.8), at the cost of the
explicit coarse-field exponent \(Q\).

\paragraph{Compact gauge sectors.}
Polynomial coordinates are local charts for compact holonomies.
The theorem applies chartwise only.  The chart transitions and the
full physical observable must retain the exact YM sign certificates;
one cannot infer a global polynomial potential.

\paragraph{Einstein conformal sector.}
If the highest even shell form is negative in a conformal direction,
(3.4) fails for every \(\lambda>0\).  Young absorption cannot repair
a wrong-sign leading form.  A positive higher-derivative completion,
a justified contour, or a physical quotient construction must first
supply a stable leading sector.

\begin{theorem}[Stability decision]
\label{thm:v17v3-decision}
The one-sided cumulant route closes for an even polynomial shell on
each coarse-field set where its highest fluctuation form has a
uniform positive floor.  It cannot close a sector possessing a
negative or null highest-degree direction unless that direction is
removed or controlled by an additional verified term.
\end{theorem}

\begin{proof}
Positive coercivity gives Theorem~\ref{thm:v17v3-stability} and hence
the V2 one-sided bound.  If a nonzero direction \(v\) has
\(H_m(x,v)<0\), then
\[
 Z(x,rv)=r^mH_m(x,v)+O(r^{m-1})\longrightarrow-\infty,
\]
so no finite lower stability constant exists.  If
\(H_m(x,v)=0\), (3.4) also fails and lower-degree terms must be
analyzed separately; no positive leading floor follows.
\end{proof}

\begin{assessment}[What V3 closes]
\label{ass:v17v3-status}
V3 derives an explicit lower stability constant for a coercive even
shell, tracks its polynomial coarse-field growth, and proves a
summable choice
\(R_n=\sqrt{A\log(n+1)}\) under a uniform sub-Gaussian coarse tail.
For polynomially small shell strength, the precise condition is
\(3\gamma-\beta>1\).

The theorem applies to stable scalar sectors but exposes the
wrong-sign Einstein conformal direction as an upstream obstruction.
The full theory still lacks the required conformal repair and the
uniform tail/prefactor estimates.
\end{assessment}

\begin{decision}[V4: the null-leading-direction refinement]
\label{dec:v17v3-next}
Analyze shells whose highest even form is positive semidefinite but
has a kernel.  Decompose transverse coercive and null variables by a
Schur minimization, derive the effective degree and sign along the
kernel, and state a recursive stability criterion that avoids
declaring every null leading direction fatal.
\end{decision}

\begin{firewall}[Claim boundary after V3]
\label{fire:v17v3}
V3 proves a sufficient stability and summability theorem.  It does
not verify its hypotheses for the coupled SM--Einstein action or
resolve the conformal instability.
\end{firewall}

\section*{Version V3 append-only record}
\addcontentsline{toc}{section}{Version V3 append-only record}
The complete V2 source was retained byte for byte before this
continuation.  It had \(28507\) bytes, \(860\) lines, and SHA--256
\[
 \texttt{ABF91F24711809C9A7BAB446A1096308F1263DEE582FEBFED1B53EA37AAFF8BE}.
\]
V3 proved explicit even-shell stability and a summable
large-field-radius theorem.


\section{V4: recursive stability along null leading directions}
\label{sec:v17v4}

\subsection{Transverse and null variables}

Let the fluctuation space split orthogonally as
\[
 \eta=(z,w)\in E\oplus K.                            \tag{4.1}
\]
The subspace \(K\) represents a linear kernel of the highest-degree
form, while \(E\) contains the coercive transverse directions.
Consider a lower estimate
\[
\begin{split}
 Z(z,w)\geq&
 \lambda|z|^m+\kappa|w|^q-B_0\\
 &-\sum_{\ell=1}^L
 B_\ell|z|^{a_\ell}|w|^{b_\ell},                   \tag{4.2}
\end{split}
\]
where \(m,q\geq2\) are even,
\[
 \lambda,\kappa>0,\qquad
 0\leq a_\ell<m,\qquad b_\ell\geq0.                \tag{4.3}
\]
The terms in the sum include lower-degree interactions and mixed
terms after absolute values have been taken.

\subsection{The induced kernel exponent}

\begin{lemma}[Mixed-monomial Schur absorption]
\label{lem:v17v4-mixed}
For \(0<a<m\), \(b\geq0\), \(B,\delta>0\),
\[
\begin{split}
 B|z|^a|w|^b
 \leq&
 \delta|z|^m+
 c_{m,a}\delta^{-a/(m-a)}
 B^{m/(m-a)}
 |w|^{r},\\
 r&=\frac{bm}{m-a},                                \tag{4.4}
\end{split}
\]
where \(c_{m,a}\) is (3.2).  For \(a=0\), the same statement holds
with no transverse term and \(r=b\).
\end{lemma}

\begin{proof}
For fixed \(w\), apply Lemma~\ref{lem:v17v3-young} with
\[
 B\quad\hbox{replaced by}\quad B|w|^b
\]
and \(k=a\).  This gives (4.4).  The case \(a=0\) is immediate.
\end{proof}

The number
\[
 r_\ell=\frac{b_\ell m}{m-a_\ell}                  \tag{4.5}
\]
is the effective power induced along the null variables after the
transverse variable is minimized.  It is the polynomial analogue of
a Schur-complement scaling exponent.

\subsection{A two-level coercivity theorem}

Choose positive \(\delta_\ell\) for the terms with \(a_\ell>0\) so
that
\[
 \sum_{\ell:a_\ell>0}\delta_\ell\leq\frac{\lambda}{2}.
                                                               \tag{4.6}
\]
Define
\[
 D_\ell=
 \begin{cases}
 c_{m,a_\ell}\delta_\ell^{-a_\ell/(m-a_\ell)}
 B_\ell^{m/(m-a_\ell)},&a_\ell>0,\\
 B_\ell,&a_\ell=0.
 \end{cases}                                       \tag{4.7}
\]
Let
\[
 I_<:=\{\ell:r_\ell<q\},\quad
 I_=:=\{\ell:r_\ell=q\},\quad
 I_>:=\{\ell:r_\ell>q\}.                            \tag{4.8}
\]

\begin{theorem}[Coercivity with a null leading subspace]
\label{thm:v17v4-two-level}
Assume
\[
 I_>=\varnothing,\qquad
 \kappa_0:=
 \kappa-\sum_{\ell\in I_=}D_\ell>0.                \tag{4.9}
\]
Then there is an explicit \(K_*<\infty\) such that
\[
 Z(z,w)
 \geq
 \frac{\lambda}{2}|z|^m+
 \frac{\kappa_0}{2}|w|^q-K_*.                      \tag{4.10}
\]
If \(I_<\ne\varnothing\), one may take
\[
\begin{split}
 K_*=
 B_0+
 \sum_{\ell\in I_<}
 c_{q,r_\ell}
 \varepsilon^{-r_\ell/(q-r_\ell)}
 D_\ell^{q/(q-r_\ell)},\qquad
 \varepsilon=\frac{\kappa_0}{2|I_<|},              \tag{4.11}
\end{split}
\]
with the evident replacement \(K_*=B_0\) when \(I_<\) is empty.
\end{theorem}

\begin{proof}
Apply Lemma~\ref{lem:v17v4-mixed} to (4.2), using (4.6).  This gives
\[
 Z(z,w)
 \geq
 \frac{\lambda}{2}|z|^m+
 \kappa|w|^q-B_0-\sum_{\ell=1}^LD_\ell|w|^{r_\ell}.
                                                               \tag{4.12}
\]
The terms in \(I_=\) reduce the coefficient of \(|w|^q\) from
\(\kappa\) to \(\kappa_0\).  There are no larger powers by (4.9).
For each \(\ell\in I_<\), apply
Lemma~\ref{lem:v17v3-young} with \(m=q\), \(k=r_\ell\),
coefficient \(D_\ell\), and absorption parameter \(\varepsilon\).
The total absorbed kernel contribution is \(\kappa_0|w|^q/2\);
the constants are (4.11).
\end{proof}

\begin{corollary}[Uniform lower stability]
\label{cor:v17v4-lower}
Under Theorem~\ref{thm:v17v4-two-level},
\[
 Z(z,w)\geq-K_*                                    \tag{4.13}
\]
and hence the V2 one-sided cumulant theorem applies whenever the
required Gaussian moment table is finite.
\end{corollary}

\begin{proof}
Drop the two nonnegative coercive terms in (4.10).
\end{proof}

\subsection{Why the exponent test is sharp for a model monomial}

\begin{proposition}[Supercritical mixed exponent produces
instability]
\label{prop:v17v4-supercritical}
Let \(m,q\) be even, \(0<a<m\), \(b>0\), and \(B>0\).  On the
positive coordinate quadrant consider
\[
 Z(z,w)=z^m+w^q-Bz^aw^b.                            \tag{4.14}
\]
If
\[
 r=\frac{bm}{m-a}>q,                               \tag{4.15}
\]
then \(Z\) is unbounded below.
\end{proposition}

\begin{proof}
For \(w>0\), choose
\[
 z(w)=
 \left(\frac{aB}{m}\right)^{1/(m-a)}
 w^{b/(m-a)}.                                      \tag{4.16}
\]
At this point
\[
 Bz^aw^b=\frac ma z^m.
\]
Therefore
\[
 Z(z(w),w)
 =
 w^q-\frac{m-a}{a}
 \left(\frac{aB}{m}\right)^{m/(m-a)}
 w^r.                                               \tag{4.17}
\]
Under (4.15), the negative term dominates as \(w\to\infty\).
\end{proof}

At the critical exponent \(r=q\), stability is a coefficient
question, as reflected by the strict margin \(\kappa_0>0\) in
(4.9).  Below \(q\), the kernel coercivity absorbs the induced term.

\subsection{A recursive flag of null spaces}

Suppose the first kernel sector \(K_1\) itself has a semidefinite
leading form with linear kernel \(K_2\).  Repeat the decomposition
\[
 K_{r-1}=E_r\oplus K_r,\qquad
 r=1,\ldots,R,\qquad K_R=\{0\}.                    \tag{4.18}
\]
Let \(m_r\) be the even coercive degree on \(E_r\).  Each mixed
monomial transferred from level \(r\) to \(r+1\) receives the
effective exponent transformation
\[
 (a,b)\longmapsto
 \frac{b\,m_r}{m_r-a}.                              \tag{4.19}
\]

\begin{theorem}[Recursive polynomial stability criterion]
\label{thm:v17v4-recursive}
Assume the flag (4.18) is finite.  At every level:
\begin{enumerate}
\item the transverse leading form has a positive floor
      \(\lambda_r|z_r|^{m_r}\);
\item every induced exponent on the next kernel is at most the next
      coercive degree \(m_{r+1}\);
\item the sum of coefficients at equality leaves a strictly positive
      next-level margin after the chosen transverse absorptions.
\end{enumerate}
Then the full shell has a lower bound
\[
 Z(\eta)\geq
 \sum_{r=1}^R c_r|z_r|^{m_r}-K,
 \qquad c_r>0,\quad K<\infty.                       \tag{4.20}
\]
\end{theorem}

\begin{proof}
Apply Theorem~\ref{thm:v17v4-two-level} at the first decomposition.
It produces a positive transverse term, a reduced positive
kernel-leading margin, and a finite constant.  Regard the remaining
kernel expression as the input to the next decomposition.  The three
hypotheses are exactly those needed to apply the two-level theorem
again.  Induction over the finite flag yields (4.20), with \(K\) the
sum of the finitely many absorption constants.
\end{proof}

\begin{warning}[Nonlinear zero sets]
\label{warn:v17v4-nonlinear}
A positive semidefinite homogeneous polynomial can vanish on a
nonlinear cone rather than a linear subspace.  The flag theorem does
not apply until that zero set is resolved into charts or a stronger
Łojasiewicz-type lower bound is proved.  Replacing the cone by its
linear span can discard coercivity and is not a valid proof.
\end{warning}

\subsection{Coarse-field dependence}

If the coefficients \(B_\ell=B_\ell(x)\) grow polynomially in a
coarse field, (4.7) and (4.11) again produce explicit polynomial
growth of \(K_*(x)\).  The exponents are obtained by composing
\[
 B\mapsto B^{m/(m-a)}
 \quad\hbox{and}\quad
 D\mapsto D^{q/(q-r)}.                              \tag{4.21}
\]
Thus V3 applies with a computable, possibly larger, coarse exponent
\(Q\).  Near a critical equality margin
\(\kappa_0\downarrow0\), the constant (4.11) diverges; this
degeneration must be included in the cutoff prefactor.

\subsection{Implications for the programme}

\paragraph{Gauge zero modes.}
An exact gauge-orbit null direction should be removed by the physical
quotient rather than stabilized by a fictitious polynomial term.
After quotienting, Theorem~\ref{thm:v17v4-recursive} may be applied
to genuine physical semidefinite directions.

\paragraph{Higgs flat directions.}
If a quartic form is semidefinite and a positive quadratic term acts
on its kernel, then \(m=4\), \(q=2\).  A mixed monomial
\(|z|^a|w|^b\) is admissible only when
\[
 \frac{4b}{4-a}\leq2,                               \tag{4.22}
\]
with a strict coefficient margin at equality.

\paragraph{Conformal direction.}
A negative highest term is not a null case and remains excluded by
V3.  A vanishing highest term followed by a negative lower-degree
conformal term also fails the recursive positive-margin hypothesis.

\begin{assessment}[What V4 closes]
\label{ass:v17v4-status}
V4 gives a recursive, quantitative stability test for semidefinite
even shells with linear null flags.  Mixed terms induce the exact
kernel exponent \(bm/(m-a)\); subcritical powers are absorbed,
critical powers consume a coefficient margin, and supercritical
powers can make the action unbounded below.

The theorem prevents a null leading direction from being declared
fatal without examining lower levels.  It does not resolve nonlinear
zero cones, exact gauge quotients, or the wrong-sign Einstein
conformal mode.
\end{assessment}

\begin{decision}[V5: companion audit and nonlinear zero cones]
\label{dec:v17v4-next}
Perform the mandatory YM/NS audit.  Then prove a chartwise stability
criterion for a compact angular zero set of a homogeneous leading
form, using a finite cover and uniform transverse lower bounds.
Continue in V6 with the first nontrivial cone chart rather than
stopping at the audit.
\end{decision}

\begin{firewall}[Claim boundary after V4]
\label{fire:v17v4}
No full SM--Einstein shell stability theorem has been proved.  V4
handles only finite linear null flags under explicit positive-margin
hypotheses.
\end{firewall}

\section*{Version V4 append-only record}
\addcontentsline{toc}{section}{Version V4 append-only record}
The complete V3 source was retained byte for byte before this
continuation.  It had \(38447\) bytes, \(1215\) lines, and SHA--256
\[
 \texttt{B1B46782E485A32191E165641CF8006A858C253587B5B798DDECD0D503201DC0}.
\]
V4 proved the recursive null-leading-direction stability criterion.


\section{V5: companion audit and finite atlases for angular zero sets}
\label{sec:v17v5}

\subsection{Mandatory companion audit}

The newest companion files inspected at this checkpoint are
\[
\begin{array}{lll}
\text{YM V68:}&15677\text{ lines},&
\texttt{C76688BDFE8FF78725931557525A5454E5789D7ED321282FACC9705EA1BCEF5B},\\
\text{NS V26:}&5319\text{ lines},&
\texttt{82C887F8C2C69E4F28FAD7C3DC8DE5B9099CB5A4BF431EBA1FFF3E91935978AD}.
\end{array}                                                    \tag{5.1}
\]
YM V66--V68 constructs the fifth \((+++)\) orientation through a
seed cube, a two-step slab, and a complete first quarter-face.
Inversion supplies the corresponding direct-minus face in the
\((---)\) orientation.  The fifth cell is not complete and neither
face carries both signs.  The result supports a finite-atlas method
with independently certified charts; it supplies no shell stability,
continuum measure, or reflection-positive limit.  NS V26 is
unchanged.

\begin{theorem}[V5 audit decision]
\label{thm:v17v5-audit}
No continuum theorem is imported from the companions.  The finite
chart-and-cover principle used by YM V68 is adopted only as a method
for resolving compact angular zero sets; every shell chart must have
its own quantitative lower bound and overlap coverage.
\end{theorem}

\begin{proof}
YM V68 proves a symbol statement on one quarter-face and expressly
leaves its other coordinate layers and every continuum conclusion
open.  The NS kernel norms do not enter the polynomial shell.
\end{proof}

\subsection{The compact angular zero set}

Let \(H_m:\mathbb R^N\to[0,\infty)\) be a nonzero homogeneous
polynomial of even degree \(m\).  Define
\[
 \Sigma=
 \{\omega\in S^{N-1}:H_m(\omega)=0\}.               \tag{5.2}
\]
This is compact.  Let
\[
 Z_x(\eta)=H_m(\eta)+L_x(\eta),                     \tag{5.3}
\]
where \(L_x\) has degree at most \(m-1\) in \(\eta\).

\begin{lemma}[Uniform coercivity away from the zero cone]
\label{lem:v17v5-away}
For every open neighborhood \(U\) of \(\Sigma\) in \(S^{N-1}\),
there is \(h_U>0\) such that
\[
 H_m(r\omega)\geq h_Ur^m
 \quad(r\geq0,\ \omega\in S^{N-1}\setminus U).      \tag{5.4}
\]
\end{lemma}

\begin{proof}
The compact set \(S^{N-1}\setminus U\) is disjoint from the zero set
of the continuous function \(H_m\).  Its minimum there is a positive
number \(h_U\).  Homogeneity gives (5.4).
\end{proof}

\begin{lemma}[Lower-order control on the exterior cone]
\label{lem:v17v5-exterior}
Let \({\cal G}\) be a coarse-field set on which
\[
 |L_x(\eta)|
 \leq\sum_{k=0}^{m-1}B_k|\eta|^k                  \tag{5.5}
\]
with finite \(B_k\).  Then for every neighborhood \(U\) of
\(\Sigma\), there is an explicit \(K_{\rm out}<\infty\) such that
\[
 Z_x(r\omega)\geq
 \frac{h_U}{2}r^m-K_{\rm out}                      \tag{5.6}
\]
for \(x\in{\cal G}\) and \(\omega\notin U\).
\end{lemma}

\begin{proof}
Use Lemma~\ref{lem:v17v5-away} and absorb every
\(B_kr^k\), \(k<m\), into \(h_Ur^m/2\) by
Lemma~\ref{lem:v17v3-young}.  The sum of the resulting constants is
\(K_{\rm out}\).
\end{proof}

\subsection{A finite conic chart theorem}

For an open \(U_\alpha\subset S^{N-1}\), write
\[
 {\cal C}_\alpha=
 \{r\omega:r\geq0,\ \omega\in U_\alpha\}.           \tag{5.7}
\]

\begin{theorem}[Finite angular atlas stability]
\label{thm:v17v5-atlas}
Suppose:
\begin{enumerate}
\item \(U_1,\ldots,U_A\) are open subsets of \(S^{N-1}\) whose union
      contains \(\Sigma\);
\item for every \(\alpha\), there are \(K_\alpha<\infty\) and a
      nonnegative function \(\Phi_\alpha\) such that
      \[
      Z_x(\eta)\geq\Phi_\alpha(\eta)-K_\alpha
      \quad(x\in{\cal G},\ \eta\in{\cal C}_\alpha); \tag{5.8}
      \]
\item the lower-order coefficient bound (5.5) holds on
      \({\cal G}\).
\end{enumerate}
Then
\[
 Z_x(\eta)\geq-K_{\cal G}
 \quad(x\in{\cal G},\ \eta\in\mathbb R^N),          \tag{5.9}
\]
where
\[
 K_{\cal G}=
 \max\{K_{\rm out},K_1,\ldots,K_A\}.                \tag{5.10}
\]
\end{theorem}

\begin{proof}
Let \(U=\bigcup_\alpha U_\alpha\).  If
\(\eta=r\omega\) with \(\omega\notin U\), (5.6) gives
\(Z_x(\eta)\geq-K_{\rm out}\).  If \(\omega\in U\), it belongs to
some \(U_\alpha\), and (5.8) gives
\(Z_x(\eta)\geq-K_\alpha\).  The origin is covered by either bound.
Taking the maximum proves (5.9).
\end{proof}

The theorem requires no single global coordinate chart.  It does
require a proved lower bound on every chart and actual coverage of
the entire angular zero set.

\begin{corollary}[Linear-null charts]
\label{cor:v17v5-linear-charts}
Suppose each \(U_\alpha\) admits coordinates in which the local shell
obeys the hypotheses of the recursive V4 null-flag theorem with
constants uniform throughout \(U_\alpha\).  Then (5.8) holds, and
Theorem~\ref{thm:v17v5-atlas} gives global lower stability on
\({\cal G}\).
\end{corollary}

\begin{proof}
Theorem~\ref{thm:v17v4-recursive} supplies (5.8) on every chart with
\(\Phi_\alpha\) equal to its sum of positive coercive powers.
Apply Theorem~\ref{thm:v17v5-atlas}.
\end{proof}

\subsection{Quantitative coverage by smaller certified charts}

\begin{lemma}[Lebesgue-number coverage test]
\label{lem:v17v5-lebesgue}
Let \(U_1,\ldots,U_A\) cover the compact set \(\Sigma\).  There is
\(\delta>0\) such that for every \(\omega\in\Sigma\), the spherical
ball \(B_{S}(\omega,\delta)\) lies in some \(U_\alpha\).
\end{lemma}

\begin{proof}
For each \(\omega\in\Sigma\), choose an index \(\alpha(\omega)\) and
a radius \(r_\omega>0\) such that
\[
 B_S(\omega,2r_\omega)\subset U_{\alpha(\omega)}.
\]
The balls \(B_S(\omega,r_\omega)\) cover \(\Sigma\); choose a finite
subcover and let \(\delta\) be the minimum of its radii.  Any point
of \(\Sigma\) lies in one selected smaller ball, and its
\(\delta\)-neighborhood lies in the corresponding doubled ball after
reducing \(\delta\) by a factor two if necessary.
\end{proof}

\begin{proposition}[Certified centre net]
\label{prop:v17v5-net}
Let \(c_1,\ldots,c_M\in\Sigma\) be a
\(\delta/2\)-net, with \(\delta\) from
Lemma~\ref{lem:v17v5-lebesgue}.  If a lower-bound chart is certified
on every \(B_S(c_i,\delta)\), then those finitely many charts cover
\(\Sigma\).
\end{proposition}

\begin{proof}
For any \(\omega\in\Sigma\), choose \(c_i\) with
\(d_S(\omega,c_i)<\delta/2\).  Then
\(\omega\in B_S(c_i,\delta)\).
\end{proof}

This is the precise part of the YM atlas method that transfers: exact
or interval-certified chart floors plus a verified finite cover.
Floating screens alone are discovery tools.

\subsection{Failure modes}

\begin{proposition}[One missed ray destroys a stability claim]
\label{prop:v17v5-missed}
Suppose \(\omega_*\in\Sigma\) lies in none of the certified chart
cones and
\[
 L_x(r\omega_*)\longrightarrow-\infty
 \quad(r\to\infty).
\]
Then \(Z_x\) is unbounded below despite all certified chart bounds.
\end{proposition}

\begin{proof}
Along the missed ray,
\[
 H_m(r\omega_*)=r^mH_m(\omega_*)=0,
\]
so \(Z_x(r\omega_*)=L_x(r\omega_*)\to-\infty\).
\end{proof}

\begin{warning}[Uniformity on overlaps]
\label{warn:v17v5-overlaps}
Chartwise lower bounds may use different decompositions, but their
constants must be uniform up to the chart boundary.  A bound that
degenerates on every overlap does not yield finite
\(K_\alpha\) in (5.8).
\end{warning}

\begin{assessment}[What V5 closes]
\label{ass:v17v5-status}
V5 performs the required companion audit and proves a finite-atlas
stability theorem for a compact angular zero set.  Coercivity away
from the zero cone is automatic by compactness; only finitely many
certified conic charts are needed near it.  A Lebesgue-number/net
criterion makes the coverage obligation explicit.

The theorem does not construct the required local charts for a
specific nonlinear zero cone.  V6 supplies the first worked cone
model and its critical lower-degree criterion.
\end{assessment}

\begin{decision}[V6: the crossed-quadratic cone]
\label{dec:v17v5-next}
Analyze
\(H_4(z,w)=(z^2-w^2)^2\), whose angular zero set has two branches.
Construct exact branch coordinates, determine which quadratic
lower forms stabilize or destabilize the two rays, and exhibit the
critical mixed coefficient threshold.
\end{decision}

\begin{firewall}[Claim boundary after V5]
\label{fire:v17v5}
No nonlinear zero-cone atlas has been verified for the
SM--Einstein action.  V5 proves only the finite-cover mechanism and
records the companion progress accurately.
\end{firewall}

\section*{Version V5 append-only record}
\addcontentsline{toc}{section}{Version V5 append-only record}
The complete V4 source was retained byte for byte before this
continuation.  It had \(48897\) bytes, \(1539\) lines, and SHA--256
\[
 \texttt{69DC440703C33497CB89C8633EC7FBE93AB78450945BF19D58C18DA954817363}.
\]
V5 performed the companion audit and proved finite-atlas stability
for compact angular zero sets.


\section{V6: exact stability on the crossed-quadratic cone}
\label{sec:v17v6}

\subsection{The cone and its two branch charts}

Consider
\[
 H_4(z,w)=(z^2-w^2)^2.                              \tag{6.1}
\]
Its zero set is the union of the two lines
\[
 {\cal L}_+=\{z=w\},\qquad
 {\cal L}_-=\{z=-w\}.                               \tag{6.2}
\]
On the \(+\) branch use
\[
 u=\frac{z-w}{\sqrt2},\qquad
 v=\frac{z+w}{\sqrt2}.                              \tag{6.3}
\]
Then
\[
 z^2-w^2=2uv,\qquad H_4=4u^2v^2.                  \tag{6.4}
\]
The \(-\) branch is obtained by interchanging \(u\) and \(v\).
Thus each angular branch has quadratic transverse vanishing with a
coefficient proportional to the square of the tangential radius.

\subsection{Quadratic and linear lower terms}

Let
\[
\begin{split}
 P(z,w)
 ={}&(z^2-w^2)^2+
 az^2+2bzw+cw^2+dz+ew+f.                           \tag{6.5}
\end{split}
\]
Define the two ray coefficients and tangent linear coefficients
\[
\begin{split}
 q_+&=a+2b+c,& \ell_+&=d+e,\\
 q_-&=a-2b+c,& \ell_-&=d-e.                        \tag{6.6}
\end{split}
\]

\begin{theorem}[Exact bounded-below criterion]
\label{thm:v17v6-exact}
The polynomial \(P\) is bounded below on \(\mathbb R^2\) if and only
if
\[
\begin{split}
 &q_+\geq0,\qquad q_-\geq0,                         \tag{6.7}\\
 &q_+=0\Longrightarrow\ell_+=0,\qquad
 q_-=0\Longrightarrow\ell_-=0.                     \tag{6.8}
\end{split}
\]
\end{theorem}

\begin{proof}
Necessity follows by restriction to the two zero lines.  On
\({\cal L}_+\), set \(z=w=t\).  Then
\[
 P(t,t)=q_+t^2+\ell_+t+f.                           \tag{6.9}
\]
This is bounded below only if \(q_+>0\), or if
\(q_+=\ell_+=0\).  The line \({\cal L}_-\) gives the corresponding
condition for \(q_-,\ell_-\).

For sufficiency, use (6.3).  A direct calculation gives
\[
 P(u,v)
 =
 4u^2v^2+Au^2+2Buv+Cv^2+pu+qv+f,                  \tag{6.10}
\]
where
\[
\begin{split}
 A&=\frac{q_-}{2},&
 C&=\frac{q_+}{2},&
 B&=\frac{a-c}{2},\\
 p&=\frac{\ell_-}{\sqrt2},&
 q&=\frac{\ell_+}{\sqrt2}.                         \tag{6.11}
\end{split}
\]
Complete the mixed square:
\[
 4u^2v^2+2Buv
 =
 4\left(uv+\frac B4\right)^2-\frac{B^2}{4}.
                                                               \tag{6.12}
\]
If \(A>0\),
\[
 Au^2+pu\geq-\frac{p^2}{4A};
\]
if \(A=0\), condition (6.8) gives \(p=0\).  The same argument applies
to \(Cv^2+qv\).  Hence every remaining term has a finite lower bound.
\end{proof}

\begin{corollary}[Explicit stability constant]
\label{cor:v17v6-floor}
Under (6.7)--(6.8),
\[
 P(z,w)\geq-K_P,                                   \tag{6.13}
\]
where
\[
 K_P=
 -f+\frac{B^2}{4}
 +{\bf1}_{A>0}\frac{p^2}{4A}
 +{\bf1}_{C>0}\frac{q^2}{4C}.                      \tag{6.14}
\]
\end{corollary}

\begin{proof}
Add the three lower bounds used in the sufficiency proof of
Theorem~\ref{thm:v17v6-exact} and rearrange.
\end{proof}

\subsection{The critical mixed coefficient}

\begin{corollary}[Sharp quadratic threshold on the cone]
\label{cor:v17v6-threshold}
The quadratic part of (6.5) is nonnegative on both leading zero rays
if and only if
\[
 a+c\geq2|b|.                                      \tag{6.15}
\]
At either equality, the linear term along that critical ray must
vanish as prescribed by (6.8).
\end{corollary}

\begin{proof}
The two inequalities \(q_\pm\geq0\) are
\[
 a+c+2b\geq0,\qquad a+c-2b\geq0,
\]
which are equivalent to (6.15).  Equality in one is exactly the
corresponding zero-ray case of Theorem~\ref{thm:v17v6-exact}.
\end{proof}

The coefficient \(B=(a-c)/2\), which is transverse-tangent in the
branch chart, has no separate threshold: the quartic product
\(4u^2v^2\) absorbs \(2Buv\) into the finite constant
\(-B^2/4\).  The dangerous mixed coefficient is \(b\), because it
changes the quadratic restriction on the two zero rays.

\subsection{Strict coercivity and flat stable cases}

\begin{proposition}[Strict ray margins give radial coercivity]
\label{prop:v17v6-coercive}
If
\[
 q_+>0,\qquad q_->0,                                \tag{6.16}
\]
then there are \(c_2,c_4>0\) and \(K<\infty\) such that
\[
 P(z,w)
 \geq
 c_4(z^2-w^2)^2+c_2(z^2+w^2)-K.                   \tag{6.17}
\]
In particular \(P(z,w)\to\infty\) as
\(|(z,w)|\to\infty\).
\end{proposition}

\begin{proof}
In branch coordinates, \(A,C>0\).  Choose
\(0<c_2<\min\{A,C\}\) and split
\[
 A u^2+C v^2
 =
 c_2(u^2+v^2)+(A-c_2)u^2+(C-c_2)v^2.
\]
Use half of the quartic term to absorb \(2Buv\), and complete squares
with the positive remaining quadratic coefficients to absorb
\(pu+qv\).  The unspent half gives \(c_4=1/2\) in the normalization
of (6.10).  Since
\[
 u^2+v^2=z^2+w^2,
\]
this proves (6.17).
\end{proof}

\begin{proposition}[Stable but noncoercive critical example]
\label{prop:v17v6-flat}
For any \(\alpha\in\mathbb R\),
\[
 P_\alpha(z,w)
 =
 (z^2-w^2)^2+\alpha(z^2-w^2)                       \tag{6.18}
\]
is bounded below by \(-\alpha^2/4\), but is constant along
unbounded subsets when the product \(z^2-w^2=-\alpha/2\).
\end{proposition}

\begin{proof}
Complete the scalar square:
\[
 P_\alpha
 =
 \left(z^2-w^2+\frac\alpha2\right)^2
 -\frac{\alpha^2}{4}.                              \tag{6.19}
\]
\end{proof}

This example satisfies \(q_+=q_-=0\) and
\(\ell_+=\ell_-=0\).  Lower stability is enough for the one-sided
tilted-moment comparison on bounded coarse sets, but noncoercivity
can prevent uniform raw moments and tightness.  Stability and
compactness must therefore remain separate ledgers.

\subsection{Uniform coefficient families}

Suppose the coefficients in (6.5) depend on a coarse field \(x\).

\begin{theorem}[Uniform crossed-cone family]
\label{thm:v17v6-family}
On a coarse set \({\cal G}\), assume all six coefficients are
bounded,
\[
 q_\pm(x)\geq\kappa_\pm>0.                          \tag{6.20}
\]
Then the constants in (6.17) can be chosen uniformly in
\(x\in{\cal G}\), and the V2 one-sided third-cumulant moment table is
uniform on \({\cal G}\).
\end{theorem}

\begin{proof}
The transformed coefficients \(A,C,B,p,q,f\) are bounded linear
combinations of the original coefficients.  Equation (6.20) gives
uniform positive lower bounds on \(A,C\).  The square-completion
constants in Proposition~\ref{prop:v17v6-coercive} are therefore
uniform.  Polynomial Gaussian moments are uniform because the
coefficient set is bounded, as in
Corollary~\ref{cor:v17v2-polynomial}.
\end{proof}

If \(q_\pm(x)\downarrow0\), the explicit term
\(\ell_\pm(x)^2/q_\pm(x)\) in (6.14) can diverge.  At an exact zero,
the corresponding \(\ell_\pm\) must vanish; a merely small linear
coefficient requires a quantitative ratio bound.

\subsection{Atlas interpretation}

On the unit circle, the zero set of (6.1) consists of four points.
Two line charts, with antipodal copies, cover it.  On a \(+\)-chart,
\(u\) is transverse and \(v\) is bounded away from zero, so
\[
 H_4=4u^2v^2\geq c_+u^2.                            \tag{6.21}
\]
The corresponding statement holds with \(u,v\) interchanged on the
\(-\)-chart.  Away from these four charts, V5 supplies a positive
quartic angular floor.  The exact global theorem above verifies the
overlaps automatically.

\begin{assessment}[What V6 closes]
\label{ass:v17v6-status}
V6 completely resolves the first nonlinear angular zero-cone model.
The sharp criterion is nonnegativity of the quadratic restriction on
both zero rays,
\(a+c\geq2|b|\), plus vanishing of the tangent linear term on every
critical ray.  Strict ray margins give coercivity; equality can give
stability without tightness.

The model is illustrative.  No claim is made that the
SM--Einstein shell reduces to this two-variable cone.
\end{assessment}

\begin{decision}[V7: general compact zero-set quadratic test]
\label{dec:v17v6-next}
Generalize the exact ray restriction to a compact angular zero set:
prove that a strictly positive quadratic lower form on the zero set
stabilizes a nonnegative even leading form, while a negative
restriction gives an unbounded ray.  Quantify the neighborhood in
which the quadratic margin persists and absorb all lower degrees.
\end{decision}

\begin{firewall}[Claim boundary after V6]
\label{fire:v17v6}
V6 proves an exact two-variable polynomial theorem.  It does not
establish stability, moment summability, or continuum convergence
for the coupled physical action.
\end{firewall}

\section*{Version V6 append-only record}
\addcontentsline{toc}{section}{Version V6 append-only record}
The complete V5 source was retained byte for byte before this
continuation.  It had \(57981\) bytes, \(1802\) lines, and SHA--256
\[
 \texttt{FBCBBAC15D25C7BBD10E1A99ECF22E5A7E3A4AD867F502558191E6F68CE0451F}.
\]
V6 proved the exact boundedness and coercivity criteria for the
crossed-quadratic zero cone.


\section{V7: a strict quadratic margin on a compact zero cone}
\label{sec:v17v7}

\subsection{Setting}

Let \(m>2\) be even and let
\[
 H_m:\mathbb R^N\to[0,\infty)
\]
be a nonzero homogeneous polynomial of degree \(m\).  Let
\[
 \Sigma=\{\omega\in S^{N-1}:H_m(\omega)=0\}.        \tag{7.1}
\]
Consider
\[
 P(\eta)=H_m(\eta)+Q_2(\eta)+\ell(\eta)+c,          \tag{7.2}
\]
where \(Q_2\) is a homogeneous quadratic form and \(\ell\) is linear.

\subsection{Persistence of an angular margin}

\begin{lemma}[Quantitative neighborhood of the zero set]
\label{lem:v17v7-neighborhood}
Assume
\[
 Q_2(\omega)\geq\kappa>0
 \quad(\omega\in\Sigma).                            \tag{7.3}
\]
Let
\[
 L_Q=
 \sup_{\omega\in S^{N-1}}
 |\nabla_SQ_2(\omega)|.                             \tag{7.4}
\]
If \(L_Q>0\), put
\[
 \delta=\min\left\{1,\frac{\kappa}{2L_Q}\right\};
                                                               \tag{7.5}
\]
if \(L_Q=0\), take \(\delta=1\).  Then
\[
 Q_2(\omega)\geq\frac{\kappa}{2}
 \quad\text{whenever }\operatorname{dist}_S(\omega,\Sigma)
 \leq\delta.                                       \tag{7.6}
\]
\end{lemma}

\begin{proof}
Choose \(\sigma\in\Sigma\) with
\[
 d_S(\omega,\sigma)=d_S(\omega,\Sigma).
\]
The spherical mean-value bound gives
\[
 |Q_2(\omega)-Q_2(\sigma)|
 \leq L_Qd_S(\omega,\sigma).
\]
Equations (7.3)--(7.5) give (7.6).  If \(L_Q=0\),
\(Q_2\) is constant on each connected component of the sphere and
the conclusion is immediate.
\end{proof}

\subsection{Global coercivity}

Define
\[
 B_1=\|\ell\|,\qquad
 M_Q=\max_{\omega\in S^{N-1}}|Q_2(\omega)|.         \tag{7.7}
\]
By compactness and (7.1), the exterior angular floor
\[
 h_\delta=
 \min_{\operatorname{dist}_S(\omega,\Sigma)\geq\delta}
 H_m(\omega)                                       \tag{7.8}
\]
is positive whenever the exterior set is nonempty.

\begin{theorem}[Strict zero-set quadratic stabilization]
\label{thm:v17v7-strict}
Under (7.3), there exist explicit \(c_2,c_m>0\) and \(K<\infty\)
such that
\[
 P(\eta)
 \geq
 c_mH_m(\eta)+c_2|\eta|^2-K.                       \tag{7.9}
\]
In particular, \(P\) is coercive and bounded below.
\end{theorem}

\begin{proof}
Write \(\eta=r\omega\).  In the angular tube
\[
 d_S(\omega,\Sigma)\leq\delta,
\]
Lemma~\ref{lem:v17v7-neighborhood} gives
\[
\begin{split}
 P(r\omega)
 &\geq
 H_m(r\omega)+\frac{\kappa}{2}r^2-B_1r+c\\
 &\geq
 H_m(r\omega)+\frac{\kappa}{4}r^2
 +c-\frac{B_1^2}{\kappa}.                          \tag{7.10}
\end{split}
\]

On the exterior angular set, (7.8) gives
\[
 H_m(r\omega)\geq h_\delta r^m.                    \tag{7.11}
\]
Use Lemma~\ref{lem:v17v3-young} to absorb
\[
 M_Qr^2+B_1r
\]
into \(h_\delta r^m/2\), leaving a finite constant
\(K_{\rm out}\).  Hence
\[
 P(r\omega)\geq\frac12H_m(r\omega)-K_{\rm out}
                                                               \tag{7.12}
\]
on the exterior cone.

To obtain one common radial quadratic term there, note that
\[
 \frac12h_\delta r^m
 \geq c_2r^2-K_2                                  \tag{7.13}
\]
for any fixed \(c_2>0\), with finite \(K_2\), because \(m>2\).
Choose \(c_2=\kappa/8\).  Combining half of (7.10) with
(7.12)--(7.13), and taking the worse finite constant, gives (7.9)
with, for example, \(c_m=1/4\).
\end{proof}

\begin{corollary}[Uniform family]
\label{cor:v17v7-family}
Let \(P_x\) be a family of the form (7.2).  If the coefficient norms
are uniformly bounded, (7.3) holds with one
\(\kappa>0\), and the exterior floors \(h_\delta\) have one positive
lower bound, then all coercivity constants in (7.9) can be chosen
uniformly in \(x\).
\end{corollary}

\begin{proof}
The proof of Theorem~\ref{thm:v17v7-strict} uses only
\(\kappa,L_Q,M_Q,B_1,h_\delta\), and \(c\).  Uniform bounds on those
quantities give common constants.
\end{proof}

\subsection{Necessary ray tests}

\begin{proposition}[Negative quadratic restriction is fatal]
\label{prop:v17v7-negative}
If there is \(\omega_*\in\Sigma\) with
\[
 Q_2(\omega_*)<0,                                  \tag{7.14}
\]
then \(P\) is unbounded below.
\end{proposition}

\begin{proof}
Along the zero ray,
\[
 P(r\omega_*)
 =
 r^2Q_2(\omega_*)+r\ell(\omega_*)+c
 \longrightarrow-\infty.                           \tag{7.15}
\]
\end{proof}

\begin{proposition}[Linear compatibility at a zero quadratic ray]
\label{prop:v17v7-linear}
If \(\omega_*\in\Sigma\), \(Q_2(\omega_*)=0\), and
\(\ell(\omega_*)\ne0\), then \(P\) is unbounded below.
\end{proposition}

\begin{proof}
Since \(m\) and \(Q_2\) are even under
\(\omega_*\mapsto-\omega_*\), choose the sign of the zero ray so that
\(\ell(\omega_*)<0\).  Then
\[
 P(r\omega_*)=r\ell(\omega_*)+c\to-\infty.
\]
\end{proof}

These are necessary tests.  Nonnegative, rather than strictly
positive, quadratic restriction is not sufficient without
transverse analysis.

\subsection{A critical-cone counterexample}

\begin{proposition}[Zero ray restriction can hide transverse
instability]
\label{prop:v17v7-critical-counter}
For \(b\ne0\), the polynomial
\[
 P_b(z,w)=z^4+2bzw                               \tag{7.16}
\]
has leading zero set \(\{z=0\}\), and its quadratic restriction to
that line is zero.  Nevertheless \(P_b\) is unbounded below.
\end{proposition}

\begin{proof}
For \(w>0\), choose \(z\) with sign opposite to \(b\) and magnitude
\[
 |z|=\left(\frac{|b|w}{2}\right)^{1/3}.             \tag{7.17}
\]
This is the critical point of \(z^4-2|b|w|z|\).
Substitution gives
\[
 P_b(z,w)
 =
 -3\left(\frac{|b|w}{2}\right)^{4/3}
 \longrightarrow-\infty.                           \tag{7.18}
\]
\end{proof}

This is the V4 induced exponent \(4/3\) with no positive kernel term
available to absorb it.

\subsection{First positive lower degree beyond quadratic}

The same compact argument works when the first verified lower form
on the zero set has even degree \(q<m\).  Suppose
\[
 P=H_m+Q_q+L_{<q},                                  \tag{7.19}
\]
where \(Q_q\) is homogeneous of degree \(q\),
\[
 Q_q(\omega)\geq\kappa>0\quad(\omega\in\Sigma),     \tag{7.20}
\]
and
\[
 |L_{<q}(r\omega)|
 \leq\sum_{k<q}B_kr^k.                             \tag{7.21}
\]

\begin{theorem}[Strict lower-form stabilization]
\label{thm:v17v7-q}
Under (7.19)--(7.21), with no intermediate terms of degree
\(q<k<m\), there are \(c_m,c_q>0\) and \(K<\infty\) such that
\[
 P(\eta)\geq c_mH_m(\eta)+c_q|\eta|^q-K.           \tag{7.22}
\]
\end{theorem}

\begin{proof}
Continuity gives a tube of \(\Sigma\) on which
\(Q_q(\omega)\geq\kappa/2\).  In that tube, absorb every
\(B_kr^k\), \(k<q\), into \(\kappa r^q/4\).  Away from the tube,
the positive angular floor of \(H_m\) absorbs \(Q_q\) and all lower
terms because their degrees are less than \(m\).  Combine the two
cone bounds exactly as in
Theorem~\ref{thm:v17v7-strict}.
\end{proof}

The exclusion of intermediate degrees in
Theorem~\ref{thm:v17v7-q} is essential.  Such terms may vanish on
\(\Sigma\) yet generate a supercritical transverse exponent, as in
V4; they require angular vanishing-order data.

\subsection{Insertion into shell stability}

For a polynomial shell with a compact leading zero set, compute the
restriction of its first lower even form to \(\Sigma\).
\begin{enumerate}
\item A uniform positive restriction closes lower stability by
      Theorem~\ref{thm:v17v7-q}.
\item A negative value gives the explicit unbounded ray of
      Proposition~\ref{prop:v17v7-negative}.
\item A zero value triggers the V4 transverse/null exponent analysis;
      ray restriction alone is inconclusive.
\end{enumerate}

\begin{assessment}[What V7 closes]
\label{ass:v17v7-status}
V7 generalizes the strict part of the crossed-cone calculation to
any compact angular zero set.  A uniformly positive quadratic, or
more generally first lower even form, on the zero set stabilizes the
full polynomial and yields uniform coercivity.  Negative restriction
is fatal, while a zero restriction requires transverse analysis and
can conceal instability.

The coupled SM--Einstein leading and lower forms have not yet been
restricted to their physical angular zero set.
\end{assessment}

\begin{decision}[V8: angular vanishing orders]
\label{dec:v17v7-next}
For intermediate homogeneous terms \(H_k\), quantify their order of
vanishing with distance to \(\Sigma\).  Derive the induced radial
power after absorption into
\(H_m(\omega)\gtrsim d(\omega,\Sigma)^r\), and compare it with the
first positive lower degree on \(\Sigma\).
\end{decision}

\begin{firewall}[Claim boundary after V7]
\label{fire:v17v7}
V7 proves a strict compact-zero-set theorem under the absence of
uncontrolled intermediate degrees.  It does not establish the
required angular margins for the full physical shell.
\end{firewall}

\section*{Version V7 append-only record}
\addcontentsline{toc}{section}{Version V7 append-only record}
The complete V6 source was retained byte for byte before this
continuation.  It had \(66709\) bytes, \(2097\) lines, and SHA--256
\[
 \texttt{BCF9A955A6EB452C9A09A11F835271BA5D954A6DCA928228E00395CBB2D48487}.
\]
V7 proved strict lower-form stabilization on a compact angular zero
set and the critical transverse counterexample.


\section{V8: angular vanishing orders and the induced radial power}
\label{sec:v17v8}

\subsection{Angular data near the leading zero set}

Let \(m>q\geq2\), with \(m,q\) even, and write
\[
\begin{split}
 P(r\omega)
 ={}&
 r^mH_m(\omega)
 +\sum_{k=q+1}^{m-1}r^kH_k(\omega)
 +r^qQ_q(\omega)
 +\sum_{k=0}^{q-1}r^kH_k(\omega),                  \tag{8.1}
\end{split}
\]
for \(r\geq0\), \(\omega\in S^{N-1}\).  Assume
\[
 H_m\geq0,\qquad
 \Sigma=\{\omega:H_m(\omega)=0\},                  \tag{8.2}
\]
and let
\[
 d(\omega)=\min\{1,\operatorname{dist}_S(\omega,\Sigma)\}.
                                                               \tag{8.3}
\]
On an angular neighborhood \(U\) of \(\Sigma\), suppose
\[
\begin{split}
 H_m(\omega)&\geq\lambda d(\omega)^R,\\
 |H_k(\omega)|&\leq B_kd(\omega)^{s_k},
 \qquad q<k<m,                                     \tag{8.4}\\
 Q_q(\omega)&\geq\kappa
 \qquad(\omega\in\Sigma),
\end{split}
\]
where \(\lambda,\kappa>0\), \(R>0\), and \(s_k\geq0\).

\subsection{Angular absorption}

\begin{lemma}[Induced radial exponent]
\label{lem:v17v8-induced}
Let \(0\leq s<R\), \(q<k<m\), and \(\delta,B>0\).  Then
\[
\begin{split}
 Br^kd^s
 \leq&
 \delta r^md^R+
 c_{R,s}\delta^{-s/(R-s)}
 B^{R/(R-s)}r^\rho,                                \tag{8.5}\\
 \rho&=\frac{kR-ms}{R-s},                          \tag{8.6}
\end{split}
\]
for \(r,d\geq0\), where
\[
 c_{R,s}
 =
 \frac{R-s}{R}
 \left(\frac{s}{R}\right)^{s/(R-s)}                \tag{8.7}
\]
with \(c_{R,0}=1\).
\end{lemma}

\begin{proof}
For fixed \(r\), apply the maximization proof of
Lemma~\ref{lem:v17v3-young} to the variable \(d\), with leading
power \(R\), lower power \(s\), coefficient \(Br^k\), and absorption
coefficient \(\delta r^m\).  The remainder is
\[
\begin{split}
 c_{R,s}(\delta r^m)^{-s/(R-s)}
 (Br^k)^{R/(R-s)}
 =
 c_{R,s}\delta^{-s/(R-s)}
 B^{R/(R-s)}
 r^{(kR-ms)/(R-s)}.
\end{split}
\]
\end{proof}

\begin{lemma}[Terms vanishing at least as fast as the leader]
\label{lem:v17v8-fast}
If \(s\geq R\), then for every \(\delta>0\) there is \(r_0<\infty\)
such that
\[
 Br^kd^s\leq\delta r^md^R
 \quad(r\geq r_0,\ 0\leq d\leq1).                  \tag{8.8}
\]
\end{lemma}

\begin{proof}
Since \(d^s\leq d^R\),
\[
 Br^kd^s\leq Br^{k-m}r^md^R.
\]
Because \(k<m\), choose
\[
 r_0\geq(B/\delta)^{1/(m-k)}.
\]
\end{proof}

Thus only \(s_k<R\) produces a new radial power.

\subsection{The Newton-line comparison}

\begin{proposition}[Subcritical, critical, and supercritical
vanishing]
\label{prop:v17v8-newton}
For \(0\leq s_k<R\), let \(\rho_k\) be (8.6).  Then
\[
\begin{array}{rcl}
\rho_k<q
&\Longleftrightarrow&
\displaystyle\frac{s_k}{R}>
\frac{k-q}{m-q},\\
\rho_k=q
&\Longleftrightarrow&
\displaystyle\frac{s_k}{R}=
\frac{k-q}{m-q},\\
\rho_k>q
&\Longleftrightarrow&
\displaystyle\frac{s_k}{R}<
\frac{k-q}{m-q}.
\end{array}                                         \tag{8.9}
\]
\end{proposition}

\begin{proof}
Since \(R-s_k>0\),
\[
 \rho_k-q
 =
 \frac{R(k-q)-s_k(m-q)}{R-s_k}.                    \tag{8.10}
\]
The three equivalences follow by taking the sign of the numerator.
\end{proof}

The line
\[
 \frac{s}{R}=\frac{k-q}{m-q}                       \tag{8.11}
\]
is the angular Newton line joining the leading datum
\((m,R)\) to the stabilizing datum \((q,0)\).  Terms above it vanish
fast enough to be subcritical; terms on it consume a coefficient
margin; terms below it generate a radial power larger than \(q\).

\subsection{A quantitative stability theorem}

Shrink \(U\), if necessary, so that continuity gives
\[
 Q_q(\omega)\geq\frac{\kappa}{2}
 \quad(\omega\in U).                                \tag{8.12}
\]
For every \(k\) with \(s_k<R\), choose
\(\delta_k>0\) satisfying
\[
 \sum_{k:s_k<R}\delta_k\leq\frac{\lambda}{4},       \tag{8.13}
\]
and define
\[
 D_k=
 c_{R,s_k}\delta_k^{-s_k/(R-s_k)}
 B_k^{R/(R-s_k)}.                                  \tag{8.14}
\]

\begin{theorem}[Angular Newton stability criterion]
\label{thm:v17v8-stability}
Assume:
\begin{enumerate}
\item no intermediate term is supercritical:
      \[
      \rho_k\leq q\quad\text{whenever }s_k<R;       \tag{8.15}
      \]
\item the critical coefficient leaves a positive margin:
      \[
      \kappa_0:=
      \frac{\kappa}{2}
      -\sum_{k:\rho_k=q}D_k>0;                      \tag{8.16}
      \]
\item all angular coefficients in (8.1) are bounded.
\end{enumerate}
Then there are \(c_m,c_q>0\) and \(K<\infty\) such that
\[
 P(\eta)\geq
 c_mH_m(\eta)+c_q|\eta|^q-K.                       \tag{8.17}
\]
\end{theorem}

\begin{proof}
It is enough first to consider \(r\geq1\) and \(\omega\in U\);
bounded \(r\) contributes only a finite constant.  Reserve
\(H_m/2\) in the final lower bound.  The other half obeys
\[
 \frac12H_m(r\omega)
 \geq\frac{\lambda}{2}r^md(\omega)^R.
\]
Apply Lemma~\ref{lem:v17v8-induced} to every term with \(s_k<R\).
By (8.13), these terms consume at most
\(\lambda r^md(\omega)^R/4\).  Apply
Lemma~\ref{lem:v17v8-fast} to all remaining intermediate terms;
for sufficiently large \(r\), their finite sum consumes at most the
other \(\lambda r^md(\omega)^R/4\).  Enlarge the bounded-radius
constant.  The reserved half of the leading form remains, and
\[
\begin{split}
 P(r\omega)\geq&
 \frac12H_m(r\omega)+\frac{\kappa}{2}r^q\\
 &-\sum_{k:s_k<R}D_kr^{\rho_k}
 -\sum_{j<q}B_jr^j-K_0.                            \tag{8.18}
\end{split}
\]
The terms with \(\rho_k=q\) leave the margin \(\kappa_0\).
Terms with \(\rho_k\leq0\) are bounded for \(r\geq1\) and are
folded into \(K_0\).  Every term with \(0<\rho_k<q\), as well as
every \(j<q\), is absorbed into \(\kappa_0r^q/2\) by the
real-exponent version of Lemma~\ref{lem:v17v3-young}.  This proves
(8.17) in \(U\).

On the compact angular complement of a smaller neighborhood of
\(\Sigma\), reserve half of \(H_m\); its other half has a positive
angular floor and absorbs every lower-degree term as in
Lemma~\ref{lem:v17v5-exterior}.  Combining the two conic bounds
proves (8.17) globally.
\end{proof}

\begin{corollary}[Uniform shell family]
\label{cor:v17v8-family}
If the constants in (8.4), the coefficient margin (8.16), and the
bounded angular coefficients are uniform on a coarse-field set
\({\cal G}\), then \(c_m,c_q,K\) in (8.17) are uniform on
\({\cal G}\).  Consequently the one-sided V2 tilted-moment theorem
applies there.
\end{corollary}

\begin{proof}
Every absorption constant in the proof of
Theorem~\ref{thm:v17v8-stability} then has a uniform bound.
\end{proof}

\subsection{Consistency with the linear-null exponent}

\begin{proposition}[Recovery of V4]
\label{prop:v17v8-v4}
Suppose near a linear null space the leading term is
\(|z|^m\).  On a large sphere, let
\[
 d\asymp\frac{|z|}{r}.
\]
A mixed monomial \(|z|^a|w|^b\) has total radial degree
\(k=a+b\) and angular vanishing order \(s=a\).  Formula (8.6) gives
\[
 \rho=
 \frac{(a+b)m-ma}{m-a}
 =
 \frac{bm}{m-a},                                   \tag{8.19}
\]
which is exactly the V4 induced kernel exponent.
\end{proposition}

\begin{proof}
Substitute \(R=m\), \(s=a\), and \(k=a+b\) in (8.6).
\end{proof}

\subsection{A model showing sharpness}

\begin{proposition}[Supercritical angular term]
\label{prop:v17v8-sharp}
Fix \(0\leq s<R\) and positive constants
\(\lambda,B,\kappa\).  Consider the two-variable model
\[
 F(r,d)=
 \lambda r^md^R-Br^kd^s+\kappa r^q,
 \qquad r\geq1,\quad0\leq d\leq1.                  \tag{8.20}
\]
If \(\rho=(kR-ms)/(R-s)>q\), then
\[
 \inf_{0\leq d\leq1}F(r,d)\longrightarrow-\infty
 \quad(r\to\infty).                                \tag{8.21}
\]
\end{proposition}

\begin{proof}
If \(s=0\), take \(d=0\) and interpret \(d^0=1\).  Then
\(F(r,0)=-Br^k+\kappa r^q\), while \(\rho=k>q\), so the claim
is immediate.  Assume henceforth that \(s>0\).
Choose
\[
 d(r)=
 \left(\frac{sB}{R\lambda}\right)^{1/(R-s)}
 r^{(k-m)/(R-s)}.                                  \tag{8.22}
\]
Because \(k<m\), \(d(r)\to0\).  At this critical point,
\[
 \lambda r^md(r)^R
 =
 \frac{s}{R}Br^kd(r)^s.
\]
Therefore the first two terms in (8.20) equal
\[
 -\left(1-\frac{s}{R}\right)
 B\left(\frac{sB}{R\lambda}\right)^{s/(R-s)}
 r^\rho.                                           \tag{8.23}
\]
Since \(\rho>q\), this negative term dominates
\(\kappa r^q\).
\end{proof}

At \(\rho=q\), the coefficient in (8.23) gives the exact model
critical margin.  The sufficient margin (8.16) is the corresponding
Young-bound certificate.

\subsection{How to certify the angular hypotheses}

For a concrete polynomial shell, the following finite data are now
required:
\begin{enumerate}
\item a lower Łojasiewicz-type bound
      \(H_m\geq\lambda d^R\) near \(\Sigma\);
\item upper vanishing bounds
      \(|H_k|\leq B_kd^{s_k}\);
\item positivity of \(Q_q\) on \(\Sigma\);
\item the critical coefficient comparison (8.16);
\item finite chart coverage of \(\Sigma\), as in V5.
\end{enumerate}
These are semialgebraic inequalities on a compact set and may be
attacked by exact sum-of-squares, interval, or elimination
certificates.  A numerical angular screen is not a proof of any of
them.

\begin{assessment}[What V8 closes]
\label{ass:v17v8-status}
V8 derives the exact induced radial power
\[
 \rho_k=\frac{kR-ms_k}{R-s_k}
\]
from an intermediate term of radial degree \(k\) and angular
vanishing order \(s_k\).  The angular Newton line classifies terms as
subcritical, critical, or supercritical, and the quantitative
coefficient margin (8.16) yields global coercivity.

The SM--Einstein shell has not yet been supplied with certified
Łojasiewicz and vanishing-order data.  Those algebraic certificates
are the next upstream target.
\end{assessment}

\begin{decision}[V9: ideal and Łojasiewicz certificates]
\label{dec:v17v8-next}
Relate angular vanishing bounds to powers of the ideal of
\(\Sigma\).  Prove an elementary certificate when
\(H_m\) is a sum of squares whose generators have a uniform normal
Jacobian, and derive \(R=2\) with vanishing orders from ideal
membership.  Keep singular zero strata separate.
\end{decision}

\begin{firewall}[Claim boundary after V8]
\label{fire:v17v8}
V8 proves a conditional angular stability theorem.  It does not
provide the semialgebraic certificates for the actual coupled shell
or resolve singular zero strata.
\end{firewall}

\section*{Version V8 append-only record}
\addcontentsline{toc}{section}{Version V8 append-only record}
The complete V7 source was retained byte for byte before this
continuation.  It had \(75825\) bytes, \(2410\) lines, and SHA--256
\[
 \texttt{CBCB0475E84819332490B3A35639C3113EC1B699E7D5D59F2466159E5EA62CA8}.
\]
V8 proved the angular vanishing-order absorption theorem and the
Newton-line stability criterion.


\section{V9: regular sum-of-squares zero sets and ideal certificates}
\label{sec:v17v9}

\subsection{A regular sum-of-squares leading form}

Let \(m\) be even and suppose
\[
 H_m(\eta)=\sum_{a=1}^rG_a(\eta)^2,                 \tag{9.1}
\]
where every \(G_a\) is homogeneous of degree \(m/2\).  Restrict to
the unit sphere:
\[
 g=(g_1,\ldots,g_r),\qquad g_a=G_a|_{S^{N-1}},
\]
and define
\[
 \Sigma=\{\omega\in S^{N-1}:g(\omega)=0\}.          \tag{9.2}
\]
Assume \(\Sigma\) is a compact smooth submanifold.  For
\(\sigma\in\Sigma\), let \(N_\sigma\Sigma\) be its normal space
inside \(T_\sigma S^{N-1}\).

\begin{definition}[Uniform normal Jacobian floor]
\label{def:v17v9-normal}
The generator map has normal floor \(\sigma_0>0\) when
\[
 |D_Tg(\sigma)v|
 \geq\sigma_0|v|
 \quad
 (\sigma\in\Sigma,\ v\in N_\sigma\Sigma).           \tag{9.3}
\]
\end{definition}

\subsection{An elementary Łojasiewicz exponent two}

\begin{theorem}[Quadratic distance floor]
\label{thm:v17v9-distance}
Under (9.1)--(9.3), there are \(\delta>0\) and
\(\lambda>0\) such that
\[
 H_m(\omega)
 \geq
 \lambda\,\operatorname{dist}_S(\omega,\Sigma)^2   \tag{9.4}
\]
whenever
\(\operatorname{dist}_S(\omega,\Sigma)<\delta\).
One may take \(\lambda=\sigma_0^2/4\) after reducing \(\delta\).
\end{theorem}

\begin{proof}
The tubular-neighborhood theorem gives, after choosing a uniform
radius, a unique representation
\[
 \omega=\exp_\sigma(v),\qquad
 \sigma\in\Sigma,\quad v\in N_\sigma\Sigma,\quad
 |v|=\operatorname{dist}_S(\omega,\Sigma).          \tag{9.5}
\]
Taylor expansion on the sphere gives
\[
 g(\exp_\sigma v)
 =
 D_Tg(\sigma)v+R_\sigma(v),                         \tag{9.6}
\]
where compactness of \(\Sigma\) and smoothness give one
\(C_R<\infty\) with
\[
 |R_\sigma(v)|\leq C_R|v|^2.                       \tag{9.7}
\]
By (9.3),
\[
 |g(\exp_\sigma v)|
 \geq\sigma_0|v|-C_R|v|^2
 \geq\frac{\sigma_0}{2}|v|                         \tag{9.8}
\]
when \(|v|\leq\sigma_0/(2C_R)\).  Square (9.8) and use
\(H_m=|g|^2\).
\end{proof}

Thus the angular exponent in V8 is
\[
 R=2                                                   \tag{9.9}
\]
on every regular stratum with a uniform normal Jacobian floor.

\subsection{Upper distance bounds for the generators}

\begin{lemma}[Generator Lipschitz bound]
\label{lem:v17v9-generator-upper}
There are \(\delta>0\) and \(L_g<\infty\) such that
\[
 |g(\omega)|
 \leq L_g\operatorname{dist}_S(\omega,\Sigma)       \tag{9.10}
\]
when the distance is below \(\delta\).
\end{lemma}

\begin{proof}
Use (9.5), \(g(\sigma)=0\), and the spherical mean-value theorem:
\[
 |g(\exp_\sigma v)-g(\sigma)|
 \leq
 \sup_{\zeta\text{ in the tube}}\|D_Tg(\zeta)\|\,|v|.
\]
The derivative supremum is finite on the compact closed tube.
\end{proof}

\subsection{Ideal-power certificates for intermediate forms}

For a multi-index \(\alpha\in\mathbb N^r\), write
\[
 g^\alpha=\prod_{a=1}^rg_a^{\alpha_a},
 \qquad |\alpha|=\sum_a\alpha_a.
\]

\begin{theorem}[Ideal membership gives angular vanishing order]
\label{thm:v17v9-ideal}
Let \(H_k|_{S^{N-1}}\) admit on a neighborhood of \(\Sigma\) a
representation
\[
 H_k(\omega)
 =
 \sum_{|\alpha|=s}
 A_\alpha(\omega)g(\omega)^\alpha,                 \tag{9.11}
\]
where \(s\) is a nonnegative integer and the finitely many
\(A_\alpha\) are bounded there.  Then
\[
 |H_k(\omega)|
 \leq B_k\operatorname{dist}_S(\omega,\Sigma)^s    \tag{9.12}
\]
near \(\Sigma\), for an explicit finite \(B_k\).
\end{theorem}

\begin{proof}
Let \(M_\alpha=\sup|A_\alpha|\).  By
Lemma~\ref{lem:v17v9-generator-upper},
\[
\begin{split}
 |H_k(\omega)|
 &\leq
 \sum_{|\alpha|=s}
 M_\alpha|g(\omega)|^s\\
 &\leq
 \left(\sum_{|\alpha|=s}M_\alpha\right)
 L_g^s\operatorname{dist}_S(\omega,\Sigma)^s.
                                                               \tag{9.13}
\end{split}
\]
\end{proof}

If \(H_k\) belongs to the \(s\)-th power of the polynomial ideal
\((G_1,\ldots,G_r)\), then homogenization and restriction to the
sphere give a representation of the form (9.11).  This is a
checkable sufficient certificate.  It need not detect the maximal
vanishing order when the displayed generators are redundant.

\begin{corollary}[Regular SOS induced powers]
\label{cor:v17v9-induced}
Under Theorems~\ref{thm:v17v9-distance} and
\ref{thm:v17v9-ideal}, the V8 induced power is:
\[
\begin{array}{c|c}
s=0&\rho_k=k\\
s=1&\rho_k=2k-m\\
s\geq2&\text{no positive induced power is required}.
\end{array}                                         \tag{9.14}
\]
\end{corollary}

\begin{proof}
Use \(R=2\) in (8.6).  For \(s=0\), this gives \(k\); for \(s=1\),
it gives \(2k-m\).  When \(s\geq R=2\),
Lemma~\ref{lem:v17v8-fast} absorbs the term directly into the leading
form outside a bounded radial set.
\end{proof}

\subsection{The regular-SOS Newton test}

\begin{proposition}[A first-ideal intermediate term]
\label{prop:v17v9-first-ideal}
Suppose \(H_k\in(G_1,\ldots,G_r)\) but no second ideal-power
certificate is available.  Relative to a positive lower degree
\(q\), the term is:
\[
\begin{array}{c|c}
2k<m+q&\text{subcritical}\\
2k=m+q&\text{critical}\\
2k>m+q&\text{supercritical}.
\end{array}                                         \tag{9.15}
\]
\end{proposition}

\begin{proof}
Here \(s=1\), so \(\rho_k=2k-m\) by (9.14).  Compare this number with
\(q\).
\end{proof}

\begin{corollary}[Second ideal power is harmless]
\label{cor:v17v9-second-ideal}
Every intermediate form with
\[
 H_k\in(G_1,\ldots,G_r)^2                           \tag{9.16}
\]
is absorbed by the regular sum-of-squares leader at sufficiently
large radius, independently of \(q\).
\end{corollary}

\begin{proof}
Use \(s\geq2=R\) in Corollary~\ref{cor:v17v9-induced}.
\end{proof}

\subsection{Descending restriction triage}

The highest intermediate degree whose restriction to \(\Sigma\) is
not identically zero must be analyzed before choosing \(q\).

\begin{theorem}[Restriction triage]
\label{thm:v17v9-triage}
Let \(k_*\) be the largest \(k<m\) such that
\(H_k|_\Sigma\not\equiv0\).
\begin{enumerate}
\item If \(H_{k_*}(\omega_*)<0\) for some
      \(\omega_*\in\Sigma\), then the full polynomial is unbounded
      below.
\item If \(H_{k_*}\geq\kappa>0\) uniformly on \(\Sigma\), then
      \(H_{k_*}\) is the first positive lower form and V7 applies
      with \(q=k_*\).
\item If \(H_{k_*}\geq0\) but has zeros on \(\Sigma\), those zeros
      form a new stratum and require a second restriction/vanishing
      analysis.
\end{enumerate}
\end{theorem}

\begin{proof}
Along a leading zero ray,
\[
 P(r\omega)
 =
 \sum_{k<m}r^kH_k(\omega).                          \tag{9.17}
\]
At \(\omega_*\), the term of highest nonzero degree \(k_*\) controls
the sign as \(r\to\infty\), proving the first statement.  Uniform
positivity is exactly the hypothesis of
Theorem~\ref{thm:v17v7-q}, proving the second.  In the third case no
uniform positive margin exists; restricting to its zero subset is
necessary.
\end{proof}

\subsection{Singular strata}

\begin{warning}[Failure of the normal Jacobian]
\label{warn:v17v9-singular}
If the rank of \(D_Tg\) drops on \(\Sigma\), the proof of
Theorem~\ref{thm:v17v9-distance} fails and the exponent need not be
two.  For example,
\[
 H(x,y)=x^4                                         \tag{9.18}
\]
is a square but behaves like fourth power of the distance to its
zero line.  Using \(R=2\) there would give the wrong induced radial
powers.
\end{warning}

Decompose a semialgebraic zero set into regular strata.  Each compact
regular stratum separated from lower strata has the \(R=2\)
certificate above.  Neighborhoods of rank-drop strata require a new
generator choice, blow-up chart, or a higher Łojasiewicz exponent.
The V5 atlas must cover those neighborhoods separately.

\subsection{Certificate ledger}

For every regular angular chart, record:
\[
\begin{array}{c|c}
\text{SOS identity}&H_m=\sum_aG_a^2\\
\text{normal floor}&
\inf_{\sigma,v}|D_Tg(\sigma)v|/|v|\\
\text{ideal representation}&
H_k=\sum_{|\alpha|=s_k}A_{k,\alpha}G^\alpha\\
\text{coefficient bound}&
\sup|A_{k,\alpha}|\\
\text{Newton comparison}&
s_k/2\ \hbox{versus}\ (k-q)/(m-q).
\end{array}                                         \tag{9.19}
\]
Every row is finite-dimensional and can, in principle, be certified
exactly on a compact semialgebraic chart.

\begin{assessment}[What V9 closes]
\label{ass:v17v9-status}
V9 proves an elementary Łojasiewicz exponent \(R=2\) for a regular
sum-of-squares leading form with a uniform normal Jacobian.  Ideal
membership then supplies the intermediate angular vanishing orders:
first-ideal terms have induced power \(2k-m\), while second-ideal
terms are directly absorbed.

Rank-drop strata remain outside this theorem.  No SOS identity,
normal floor, or ideal membership has yet been certified for the
actual SM--Einstein shell.
\end{assessment}

\begin{decision}[V10: companion audit and singular monomial strata]
\label{dec:v17v9-next}
Perform the mandatory YM/NS audit.  Then treat a rank-drop monomial
sum-of-squares model
\(H_m=\sum_a x^{2\alpha_a}\), compute its distance exponent from the
smallest active monomial order on each stratum, and continue in V11
with a blow-up chart at a stratum intersection.
\end{decision}

\begin{firewall}[Claim boundary after V9]
\label{fire:v17v9}
V9 handles only regular sum-of-squares zero strata.  Singular strata,
the physical shell algebra, and all continuum gates remain open.
\end{firewall}

\section*{Version V9 append-only record}
\addcontentsline{toc}{section}{Version V9 append-only record}
The complete V8 source was retained byte for byte before this
continuation.  It had \(86826\) bytes, \(2782\) lines, and SHA--256
\[
 \texttt{48413D5EEB56B09D53DDCF200E415CA3FAE046D6FC4D1E4CDAC48EE307BE5BD7}.
\]
V9 proved regular sum-of-squares distance and ideal-power
certificates for the V8 angular data.


\section{V10: companion audit and singular monomial strata}
\label{sec:v17v10}

\subsection{Mandatory companion audit}

The checkpoint files are
\[
\begin{array}{lll}
\text{YM V70:}&16097\text{ lines},&
\texttt{44889770DBCFB035B3586E67EAF0B7397447B8E082E33C5BE93439296558B7A3},\\
\text{NS V26:}&5319\text{ lines},&
\texttt{82C887F8C2C69E4F28FAD7C3DC8DE5B9099CB5A4BF431EBA1FFF3E91935978AD}.
\end{array}                                                    \tag{10.1}
\]
YM V69--V70 extends the fifth orientation through a complete second
quarter-face; inversion gives the paired direct-minus face.  The
third coordinate and the missing sign on each paired cell remain
open.  No continuum, entropy, or reflection-positive result follows.
NS V26 is unchanged.

\begin{theorem}[V10 audit decision]
\label{thm:v17v10-audit}
No companion theorem supplies the singular angular distance
exponents required by V9.  The present direction is retained.
\end{theorem}

\begin{proof}
YM V70 is a finite symbol-atlas theorem, and NS V26 is an Oseen
kernel norm theorem.  Neither concerns a monomial ideal or a
Łojasiewicz distance bound.
\end{proof}

\subsection{A monomial ideal near a coordinate stratum}

Let \(y=(y_1,\ldots,y_d)\in[-1,1]^d\), and let
\[
 M_a(y)=c_a y^{\beta_a}
 =
 c_a\prod_{i=1}^dy_i^{\beta_{a,i}},
 \qquad c_a\ne0,                                   \tag{10.2}
\]
for \(1\leq a\leq A\).  Put
\[
 {\cal V}=\{y:M_a(y)=0\text{ for every }a\},\qquad
 H(y)=\sum_{a=1}^AM_a(y)^2.                        \tag{10.3}
\]
Assume \(|\beta_a|>0\) for every \(a\).  Thus every support
\(S_a\) is nonempty, the origin belongs to \({\cal V}\), and the
distance to \({\cal V}\) below is finite.
Let
\[
 S_a=\operatorname{supp}\beta_a
 =\{i:\beta_{a,i}>0\}.                              \tag{10.4}
\]
A set \(S\subset\{1,\ldots,d\}\) is a hitting set when
\[
 S\cap S_a\ne\varnothing
 \quad\text{for every }a.                           \tag{10.5}
\]
Setting all coordinates in a hitting set to zero produces a point
of \({\cal V}\).

\subsection{The worst smallest active order}

For a set \(S\) that is not a hitting set, define
\[
 d(S)=
 \min\{|\beta_a|:S\cap S_a=\varnothing\},           \tag{10.6}
\]
and set
\[
 D_*=
 \max_{S\text{ not a hitting set}}d(S).             \tag{10.7}
\]
This finite combinatorial number is the worst, over possible sets of
small coordinates, of the smallest monomial degree that remains
active.

\begin{theorem}[Monomial distance inequality]
\label{thm:v17v10-distance}
Let
\[
 c_*=\min_a|c_a|>0.
\]
For \(y\in[-1,1]^d\), put
\[
 \Delta(y)=\operatorname{dist}(y,{\cal V}).
\]
Then
\[
 H(y)
 \geq
 c_*^2
 \left(\frac{\Delta(y)}{2\sqrt d}\right)^{2D_*}.
                                                               \tag{10.8}
\]
\end{theorem}

\begin{proof}
If \(\Delta(y)=0\), the claim is trivial.  Otherwise set
\[
 t=\frac{\Delta(y)}{2\sqrt d}
\]
and let
\[
 S_t=\{i:|y_i|<t\}.                                 \tag{10.9}
\]
If \(S_t\) were a hitting set, set \(y_i\) to zero for
\(i\in S_t\) and leave the other coordinates unchanged.  The
resulting point would lie in \({\cal V}\) and be at distance at most
\[
 \sqrt{|S_t|}\,t\leq\sqrt d\,t=\Delta(y)/2,
\]
contradicting the definition of \(\Delta(y)\).  Hence \(S_t\) is not
a hitting set.

By (10.6)--(10.7), there is an index \(a\) with
\[
 S_t\cap S_a=\varnothing,\qquad |\beta_a|\leq D_*.
\]
Every coordinate in \(S_a\) has modulus at least \(t\), so
\[
 |M_a(y)|
 \geq c_*t^{|\beta_a|}
 \geq c_*t^{D_*},                                  \tag{10.10}
\]
because \(0<t\leq1\).  Since \(H\geq M_a^2\), this proves (10.8).
\end{proof}

\begin{corollary}[A certified singular Łojasiewicz exponent]
\label{cor:v17v10-exponent}
The monomial sum of squares (10.3) obeys a distance lower bound with
angular exponent
\[
 R=2D_*.                                            \tag{10.11}
\]
\end{corollary}

\begin{proof}
Equation (10.8) is precisely such a bound.  Smooth chart metrics are
uniformly equivalent to Euclidean distance after shrinking a compact
coordinate chart, changing only the constant.
\end{proof}

The exponent (10.11) is a certified exponent, not necessarily the
smallest possible one.

\subsection{Reduction at a sphere stratum}

Suppose the original homogeneous generators are monomials
\[
 G_a(\omega)=c_a\prod_{i=1}^N\omega_i^{\alpha_{a,i}}
                                                               \tag{10.12}
\]
on the sphere.  Fix a coordinate stratum on which coordinates in
\(I\) vanish and coordinates in \(J\) stay uniformly away from zero.
In a sufficiently small chart, the \(J\)-factors are bounded above
and below by positive constants.  Define the normal multi-index
\[
 \beta_a=(\alpha_{a,i})_{i\in I}.                  \tag{10.13}
\]

\begin{proposition}[Stratum monomial reduction]
\label{prop:v17v10-stratum}
If the chart meets the common zero set, every retained generator has
\(|\beta_a|>0\); a generator with \(\beta_a=0\) is bounded away from
zero there and instead places the chart in the coercive exterior.
On a zero-set chart the leading sum of squares is comparable to the
normal monomial sum
\[
 \sum_a y^{2\beta_a}.                               \tag{10.14}
\]
Consequently Theorem~\ref{thm:v17v10-distance} gives a local angular
distance exponent \(2D_*(I,J)\).
\end{proposition}

\begin{proof}
Each tangent-coordinate factor
\(\prod_{j\in J}\omega_j^{\alpha_{a,j}}\) lies between two positive
constants in absolute value.  Multiplying the normal monomial by
such a factor changes its square only by uniform positive
constants.  Sum over the finite generator set and apply (10.8).
\end{proof}

If a generator has \(\beta_a=0\) and is nonzero on the stratum, the
point is not in the common zero set and belongs to the coercive
exterior region instead.

\subsection{Examples}

\begin{example}
For
\[
 H(x,y)=x^{2p}+y^{2q},                              \tag{10.15}
\]
the zero set is the origin.  The supports are \(\{x\}\) and
\(\{y\}\).  The non-hitting sets
\(\varnothing,\{x\},\{y\}\) give
\[
 D_*=\max\{\min(p,q),q,p\}=\max(p,q).              \tag{10.16}
\]
Thus
\[
 H(x,y)\geq c\,\operatorname{dist}((x,y),0)^{
 2\max(p,q)}.                                      \tag{10.17}
\]
This exponent is sharp along the axis carrying the larger power.
\end{example}

\begin{example}
For
\[
 H(x,y)=x^{2p}y^{2q},                               \tag{10.18}
\]
the zero set is the union of the two axes.  The only monomial has
degree \(p+q\), so
\[
 D_*=p+q,\qquad R=2(p+q).                           \tag{10.19}
\]
Indeed, on the diagonal \(x=y=t\), the distance to the union of axes
is \(|t|\) and \(H=t^{2(p+q)}\).
\end{example}

\begin{example}
For
\[
 H(x,y,z)=x^2y^2+x^2z^2,                            \tag{10.20}
\]
the zero set is
\[
 \{x=0\}\cup\{y=z=0\}.
\]
Both generators have degree two, and \(D_*=2\); hence \(R=4\).
\end{example}

\subsection{Intermediate monomials on a singular zero variety}

Let \({\mathfrak H}\) be the finite family of hitting sets from
(10.5), and put
\[
 Z_S=\{y:y_i=0\text{ for }i\in S\}.
\]
Then
\[
 {\cal V}=\bigcup_{S\in{\mathfrak H}}Z_S,
 \qquad
 \Delta(y)=\min_{S\in{\mathfrak H}}
 \left(\sum_{i\in S}|y_i|^2\right)^{1/2}.          \tag{10.21}
\]
For a multi-index \(\gamma\), define its dual hitting order by
\[
 \nu(\gamma)=
 \min_{S\in{\mathfrak H}}\sum_{i\in S}\gamma_i.  \tag{10.22}
\]
Notice that \(\nu(\gamma)>0\) exactly when \(y^\gamma\) vanishes
identically on every component \(Z_S\), hence on \({\cal V}\).

\begin{lemma}[Dual hitting-set upper bound]
\label{lem:v17v10-dual-upper}
For \(y\in[-1,1]^d\) with \(\Delta(y)\leq1\),
\[
 |y^\gamma|\leq\Delta(y)^{\nu(\gamma)}.            \tag{10.23}
\]
\end{lemma}

\begin{proof}
Choose a hitting set \(S_y\) attaining the finite minimum in
(10.21), and write \(\delta=\Delta(y)\).  Every coordinate indexed
by \(S_y\) has modulus at most \(\delta\), while all remaining
coordinates have modulus at most one.  Therefore
\[
 |y^\gamma|
 \leq
 \delta^{\sum_{i\in S_y}\gamma_i}
 \leq
 \delta^{\nu(\gamma)},                             \tag{10.24}
\]
because \(0\leq\delta\leq1\).
\end{proof}

Let an intermediate homogeneous form have a local monomial
representation
\[
 H_k(y)=\sum_\gamma a_\gamma y^\gamma.              \tag{10.25}
\]
Its certified vanishing order relative to the entire singular zero
variety is
\[
 s_k=\min_{\gamma:a_\gamma\ne0}\nu(\gamma).        \tag{10.26}
\]

\begin{proposition}[Intermediate distance bound]
\label{prop:v17v10-upper}
On the region \(\Delta(y)\leq1\),
\[
 |H_k(y)|
 \leq B_k\Delta(y)^{s_k},
 \qquad B_k=\sum_\gamma|a_\gamma|.                 \tag{10.27}
\]
\end{proposition}

\begin{proof}
Apply Lemma~\ref{lem:v17v10-dual-upper} term by term.  Since
\(\nu(\gamma)\geq s_k\) and \(\Delta(y)\leq1\), each term is at
most \(|a_\gamma|\Delta(y)^{s_k}\).
\end{proof}

Combining (10.11) and (10.27) supplies the V8 data
\[
 R=2D_*,\qquad s=s_k.                               \tag{10.28}
\]
The induced radial power is therefore
\[
 \rho_k=
 \frac{2D_*k-ms_k}{2D_*-s_k}
 \quad(s_k<2D_*).                                  \tag{10.29}
\]
A displayed monomial with \(\nu(\gamma)=0\) does not vanish on at
least one zero component and correctly forces \(s_k=0\).  Algebraic
cancellation among such monomials can yield a better order, but then
that better order requires a separate ideal or divisibility
certificate.
\subsection{Stratified certificate}

\begin{theorem}[Finite monomial-stratum atlas]
\label{thm:v17v10-atlas}
Suppose the compact angular zero set is covered by finitely many
coordinate strata on which the leading generators reduce to
monomials as in Proposition~\ref{prop:v17v10-stratum}.  If on every
stratum:
\begin{enumerate}
\item the intermediate vanishing orders satisfy the V8 Newton
      inequalities using \(R=2D_*\);
\item the critical coefficient margins are positive;
\item the first lower form on the residual zero set is uniformly
      positive;
\end{enumerate}
then the shell has a uniform lower bound on that compact
coarse-field set.
\end{theorem}

\begin{proof}
Theorem~\ref{thm:v17v10-distance} and
Proposition~\ref{prop:v17v10-upper} verify the hypotheses of the V8
stability theorem on each stratum.  V5 combines the finitely many
chartwise lower bounds with the coercive exterior cone.
\end{proof}

\subsection{Limitations of the monomial certificate}

General polynomial generators can cancel after a coordinate change;
replacing them by their separate monomial absolute values may change
the zero set.  Theorem~\ref{thm:v17v10-distance} applies only after an
exact monomial reduction or a two-sided comparability proof.  It also
gives a possibly nonoptimal exponent, which can make the Newton test
appear supercritical even when a sharper local exponent would pass.

\begin{assessment}[What V10 closes]
\label{ass:v17v10-status}
V10 performs the required companion audit and proves a combinatorial
distance exponent for singular monomial sum-of-squares zero sets.
The exponent is twice the worst smallest active monomial degree
\(D_*\), computed from non-hitting coordinate sets.  Intermediate
monomial orders then feed directly into the V8 induced-power formula.

No exact monomial reduction or critical margin has been proved for
the SM--Einstein shell.  Intersections of strata can require
different blow-up charts, treated next.
\end{assessment}

\begin{decision}[V11: blow-up at a crossing stratum]
\label{dec:v17v10-next}
Continue after the audit with
\(H=x^2y^2\) near the crossing of its two zero components.  Use a
real blow-up to separate radial approach to the intersection from
angular approach to each branch, and compare the resulting exponent
with the coarse monomial value \(R=4\).
\end{decision}

\begin{firewall}[Claim boundary after V10]
\label{fire:v17v10}
V10 proves a monomial-ideal distance certificate, not a stability
theorem for the actual coupled action.  The companion programmes do
not close any continuum gate.
\end{firewall}

\section*{Version V10 append-only record}
\addcontentsline{toc}{section}{Version V10 append-only record}
The complete V9 source was retained byte for byte before this
continuation.  It had \(96632\) bytes, \(3109\) lines, and SHA--256
\[
 \texttt{D9AD6B07A829DB4F1AFC6BCFCE3791710AEEF60559CE841D759DCD297C568F74}.
\]
V10 performed the companion audit and proved singular
monomial-stratum distance and vanishing-order certificates.


\section{V11: blow-up analysis of a crossing zero stratum}
\label{sec:v17v11}

\subsection{Exact distance and the sharp uniform exponent}

Consider
\[
 H(x,y)=x^2y^2,\qquad
 Z=\{(x,y):xy=0\},
 \qquad
 \Delta(x,y)=\operatorname{dist}((x,y),Z).          \tag{11.1}
\]

\begin{lemma}[Distance to the crossing]
\label{lem:v17v11-distance}
For every \((x,y)\in\mathbb R^2\),
\[
 \Delta(x,y)=\min\{|x|,|y|\}.                      \tag{11.2}
\]
\end{lemma}

\begin{proof}
The distances to the \(y\)-axis and \(x\)-axis are respectively
\(|x|\) and \(|y|\).  Distance to their union is their minimum.
\end{proof}

\begin{theorem}[Sharp uniform crossing exponent]
\label{thm:v17v11-sharp}
On \([-1,1]^2\),
\[
 H(x,y)\geq\Delta(x,y)^4.                          \tag{11.3}
\]
The exponent \(4\) is the least \(R\) for which
\(H\geq c\Delta^R\), with \(c>0\), can hold near the crossing.
\end{theorem}

\begin{proof}
Both \(|x|\) and \(|y|\) are at least \(\Delta\), which proves
(11.3).  On \(x=y=t\), one has \(H=t^4\) and \(\Delta=|t|\).
For \(R<4\), the ratio \(H/\Delta^R=|t|^{4-R}\) tends to zero.
\end{proof}

Thus the V10 certificate \(R=4\) is optimal for this crossing.

\subsection{Real blow-up}

In polar blow-up coordinates
\[
 x=r\cos\theta,\qquad y=r\sin\theta,\qquad r\geq0, \tag{11.4}
\]
put
\[
 a(\theta)=\min\{|\cos\theta|,|\sin\theta|\},
 \qquad
 b(\theta)=\max\{|\cos\theta|,|\sin\theta|\}.
\]
Then \(2^{-1/2}\leq b\leq1\) and
\[
 \Delta=ra,\qquad H=r^4a^2b^2.                    \tag{11.5}
\]

\begin{proposition}[Resolved factors]
\label{prop:v17v11-polar}
The pullback obeys
\[
 \frac12r^2\Delta^2\leq H\leq r^2\Delta^2.         \tag{11.6}
\]
The exceptional divisor has radial order four away from the strict
transforms of the axes.  Each strict transform has transverse order
two away from the exceptional divisor.
\end{proposition}

\begin{proof}
Use (11.5) and \(1/2\leq b^2\leq1\).  If \(a\) is bounded below,
\(H\) is comparable to \(r^4\).  If \(r\) is bounded below and a
strict transform is approached, the nonzero factor \(b^2\) leaves
the square of its angular defining function.
\end{proof}

In the affine \(x\)-chart,
\[
 x=u,\qquad y=uv,\qquad |v|\leq1,
\]
and hence
\[
 \pi_x^*H=u^4v^2,\qquad
 \pi_x^*\Delta=|u||v|,\qquad
 \pi_x^*H=u^2(\pi_x^*\Delta)^2.                   \tag{11.7}
\]
The \(y\)-chart is symmetric.  The intersection \(u=v=0\) explains
why the branch exponent two is not uniform downstairs.

\begin{proposition}[Regular branch exponent]
\label{prop:v17v11-branch}
If \(0<\varepsilon\leq1\),
\[
 |x|\geq\varepsilon,\qquad |y|\leq\varepsilon/2,
\]
then \(\Delta=|y|\) and
\[
 \varepsilon^2\Delta^2\leq H\leq\Delta^2.          \tag{11.8}
\]
The symmetric assertion holds near the other branch.
\end{proposition}

\begin{proof}
Here \(|y|<|x|\), and \(H=x^2\Delta^2\).
\end{proof}

\subsection{Sharp orders of intermediate monomials}

For \(a,b\in\mathbb N\), let \(K_{a,b}=x^ay^b\).

\begin{theorem}[Crossing order]
\label{thm:v17v11-monomial-order}
On \([-1,1]^2\),
\[
 |K_{a,b}(x,y)|
 \leq\Delta(x,y)^{\min(a,b)}.                     \tag{11.9}
\]
The exponent \(\min(a,b)\) is the largest possible uniform exponent
in every neighborhood of the crossing.
\end{theorem}

\begin{proof}
If \(|x|\leq|y|\), then
\[
 |x|^a|y|^b\leq|x|^a
 \leq\Delta^{\min(a,b)}.
\]
The other region is symmetric.  If \(a\leq b\), fix an arbitrarily
small \(y_0\ne0\) and let \(x=t\to0\).  For \(q>a\),
\[
 \frac{|K_{a,b}(t,y_0)|}{\Delta(t,y_0)^q}
 =|y_0|^b|t|^{a-q}\longrightarrow\infty.
\]
If \(b\leq a\), interchange the variables.
\end{proof}

The two affine pullbacks are
\[
 \pi_x^*K_{a,b}=u^{a+b}v^b,\qquad
 \pi_y^*K_{a,b}=u^{a+b}v^a.                       \tag{11.10}
\]
They record exceptional order \(a+b\) and strict-transform orders
\(b\) and \(a\).  Their worst strict-transform order is precisely
\(\min(a,b)\).

For \(K=\sum c_{a,b}x^ay^b\), the termwise certificate is
\[
 s_K=\min_{c_{a,b}\ne0}\min(a,b),\qquad
 |K|\leq
 \left(\sum|c_{a,b}|\right)\Delta^{s_K}.           \tag{11.11}
\]
Cancellation can improve this order only with a separate algebraic
certificate.

\subsection{The crossing Newton test}

Let \(m\) be the radial degree of a leading shell form with local
angular model (11.1).  If its degree-\(k\) intermediate form has
crossing order \(s_k\), V8 gives
\[
 \rho_k=\frac{4k-ms_k}{4-s_k}
 \quad(0\leq s_k<4).                               \tag{11.12}
\]
Terms with \(s_k\geq4\) are absorbed by the leader.  On the regular
\(y=0\) branch the data instead are
\[
 R=2,\qquad
 s_k^{(y)}=\min_{c_{a,b}\ne0}b,                    \tag{11.13}
\]
and the other branch uses the minimum \(a\).  A global certificate
must pass the two regular tests and the crossing test separately.

\subsection{Quadratic lower jets}

Set
\[
 Q(x,y)=A x^2+2Bxy+C y^2,\qquad
 F(x,y)=x^2y^2+Q(x,y).                             \tag{11.14}
\]

\begin{theorem}[Quadratic-jet stability and residual order]
\label{thm:v17v11-quadratic}
The following are equivalent:
\begin{enumerate}
\item \(F\geq0\) on some neighborhood of the origin;
\item \(Q\) is positive semidefinite;
\item
\[
 A\geq0,\qquad C\geq0,\qquad AC-B^2\geq0.         \tag{11.15}
\]
\end{enumerate}
The residual alternatives are:
\begin{enumerate}
\item if \(Q\) is positive definite, then
      \(F\geq\lambda(x^2+y^2)\) for some \(\lambda>0\);
\item if \(Q\) has rank one and its kernel is not a coordinate
      axis, then \(F\geq c(x^2+y^2)^2\) locally, and exponent four
      is sharp;
\item if \(Q\) has rank one and its kernel is a coordinate axis,
      the zero set is that axis and \(F\) is comparable to its
      squared distance;
\item if \(Q=0\), Theorem~\ref{thm:v17v11-sharp} applies.
\end{enumerate}
\end{theorem}

\begin{proof}
If \(F\geq0\) near zero, then for every \(v\in\mathbb R^2\),
\[
 Q(v)=\lim_{t\to0}\frac{F(tv)}{t^2}\geq0.          \tag{11.16}
\]
Conversely, \(Q\geq0\) implies \(F\geq0\).  The equivalence with
(11.15) is the principal-minor test for a symmetric \(2\times2\)
matrix.

Positive definiteness gives the first lower bound by compactness of
the unit circle.  Suppose next that \(Q\) has rank one.  Away from
its kernel directions, \(Q\geq c_1r^2\geq c_1r^4\) for \(r\leq1\).
If the kernel is not an axis, \(x^2y^2\geq c_2r^4\) in a small
conic neighborhood of it.  These finitely many cones cover the
circle and prove the uniform fourth-order bound.  On the kernel ray
the quadratic term vanishes and the quartic is a nonzero multiple
of \(r^4\), proving sharpness.

If the kernel is the \(x\)-axis, then \(Q=C y^2\), \(C>0\), and
\[
 F=y^2(x^2+C),                                     \tag{11.17}
\]
which is comparable to \(y^2\) on a small box.  The other axis is
symmetric.  The case \(Q=0\) was already proved.
\end{proof}

Thus a negative quadratic direction is an immediate obstruction.
A positive definite jet lifts the crossing to order two.  A
singular nonnegative jet passes its kernel to the next restriction
stage with the residual order proved above.

\subsection{Certificate ledger}

At a candidate physical crossing, one must record
\[
\begin{array}{c|c}
\text{leading pullback}&u^4v^2\text{ up to a positive unit}\\
\text{regular branch exponents}&R=2\text{ on both branches}\\
\text{intersection exponent}&R=4\text{ with sharpness}\\
\text{intermediate pullbacks}&u^{a+b}v^b,\ u^{a+b}v^a\\
\text{quadratic jet}&A,C,AC-B^2\text{ margins}\\
\text{residual kernel}&\text{next restriction stratum, if any}.
\end{array}                                                       \tag{11.18}
\]

\begin{assessment}[What V11 closes]
\label{ass:v17v11-status}
V11 resolves the model crossing \(x^2y^2\).  Its least uniform
distance exponent is four, each regular branch has exponent two,
and \(x^ay^b\) has sharp uniform order \(\min(a,b)\).  The blow-up
retains the exceptional and strict-transform orders separately.
The quadratic-jet theorem gives a necessary and sufficient local
nonnegativity test and classifies the residual exponent.

The actual SM--Einstein angular polynomial and its jets have not
yet been extracted.  No physical stability or continuum gate is
claimed.
\end{assessment}

\begin{decision}[V12: extract the physical leading shell]
\label{dec:v17v11-next}
Return upstream to the grand-tensor and SM--Einstein manuscripts.
Locate the finite-dimensional coarse-field action, write its
highest homogeneous part and first lower jets in a fixed field
basis, and decide which regular, monomial, or crossing certificate
is applicable.  If coefficients are absent, isolate the exact
missing datum instead of adding another abstract model.
\end{decision}

\begin{firewall}[Claim boundary after V11]
\label{fire:v17v11}
V11 is complete for the normal-crossing model.  It does not identify
that model inside the physical shell, construct a Gibbs measure, or
prove a continuum limit.
\end{firewall}

\section*{Version V11 append-only record}
\addcontentsline{toc}{section}{Version V11 append-only record}
The complete V10 source was retained byte for byte before this
continuation.  It had \(109012\) bytes, \(3506\) lines, and SHA--256
\[
 \texttt{032ECF46333BEF1D9494C9C0D1EF60F0E7A55CBE82792B2457724CA92374F027}.
\]
V11 proved the sharp crossing exponent, its blow-up factorization,
the exact intermediate monomial orders, and the quadratic-jet
stability criterion.


\section{V12: extraction of the physical Higgs leader and the total-action gap}
\label{sec:v17v12}

\subsection{Source audit}

The upstream sources used here are
\[
\begin{array}{lll}
\text{Grand Tensor V067:}&68938\text{ lines},&
\texttt{FACD058EC8BB210AD8B296C9D7F4D78A229D08D7A2AE378808700655CB8BC0FE},\\
\text{SM--Einstein V35:}&36310\text{ lines},&
\texttt{AE1FBC0A8125C37465E38E7649E2992725F4579A86F0BB17483B9AE829CFB17C}.
\end{array}                                                       \tag{12.1}
\]
The originally supplied V30 path is absent; V35 is its current later
version in the same source series.

Four distinct source objects must not be conflated:
\begin{enumerate}
\item the classical active-carrier action \textup{(CA.12)} has the
      explicit potential
      \(\lambda_H\sum_xm_x(\|H_x\|^2-v_H^2)^2\);
\item the literal finite shell \textup{(FS.8)} has
      \[
      \kappa_H\sum_{e:x\to y}
      \|H_y-\rho_H(U_e)H_x\|^2
      +\sum_x\bigl(m_H^2\|H_x\|^2+\lambda_H\|H_x\|^4\bigr),
                                                                  \tag{12.2}
      \]
      but restricts \(H\) to a fixed compact ball;
\item the unbounded RPESM cylinder \textup{(RP.5)} specifies only
      that \(S_{H,n}\) contains
      \(\lambda_4\sum_x\|H_x\|^4\), that \(\lambda_4>0\), and that
      all remaining negative parts are dominated by that quartic;
\item the exact checkerboard marginal \textup{(WB.4)} contains the
      nonpolynomial term
      \(-\sum_y\psi_8(\sum_{x\sim y}\phi_x)\), with a proved
      nonzero degree-six-and-higher Taylor tail.
\end{enumerate}
Consequently an exact finite-degree analysis applies directly to
the Higgs part of \textup{(CA.12)} and \textup{(FS.8)}, but not to
the whole Wilsonian or full Einstein--SM action as presently written.

\subsection{The extracted homogeneous forms}

Let \(N=|V_n|\), choose real coordinates
\[
 H_x=(h_{x,1},\ldots,h_{x,4})\in\mathbb R^4\simeq\mathbb C^2,
 \qquad
 r^2=\sum_{x,\alpha}h_{x,\alpha}^2.                \tag{12.3}
\]
For the explicit finite-shell Higgs action, the highest and next
homogeneous parts are
\[
\boxed{
\begin{aligned}
 {\cal H}_{4,n}(H)
 &=\lambda_H\sum_{x\in V_n}
   \left(\sum_{\alpha=1}^4h_{x,\alpha}^2\right)^2,\\
 {\cal H}_{2,n}(U,H)
 &=\kappa_H\|D_UH\|_{\ell^2(E_n)}^2+m_H^2r^2 .
\end{aligned}}                                                     \tag{12.4}
\]
There is no cubic part.  Gauge-link variables are compact, and the
representation \(\rho_H(U_e)\) is unitary.

For the weighted classical potential in \textup{(CA.12)}, replace
the first row by
\[
 {\cal H}_{4,X}(H)
 =\lambda_H\sum_{x\in V_X}m_{x,X}\|H_x\|^4,        \tag{12.5}
\]
and its quadratic potential row is
\(-2\lambda_Hv_H^2\sum_xm_{x,X}\|H_x\|^2\).

\begin{theorem}[Exact angular floor of the physical Higgs leader]
\label{thm:v17v12-angular-floor}
Let \(w_x>0\) and
\[
 P_4(H)=\lambda\sum_{x=1}^Nw_x\|H_x\|^4,
 \qquad \lambda>0.
\]
Then
\[
 P_4(H)\geq
 \frac{\lambda}{\sum_xw_x^{-1}}\,r^4,             \tag{12.6}
\]
and the constant is optimal.  In the unweighted finite shell,
\[
 \min_{r=1}{\cal H}_{4,n}=\frac{\lambda_H}{N}.
                                                                  \tag{12.7}
\]
\end{theorem}

\begin{proof}
Put \(q_x=\|H_x\|^2\).  Weighted Cauchy--Schwarz gives
\[
 \left(\sum_xq_x\right)^2
 \leq
 \left(\sum_xw_xq_x^2\right)
 \left(\sum_xw_x^{-1}\right).
\]
This is (12.6).  Equality holds when \(q_x=cw_x^{-1}\).
Taking \(w_x=1\) proves (12.7).
\end{proof}

\begin{corollary}[No singular angular Higgs stratum]
\label{cor:v17v12-no-stratum}
At every fixed cutoff the leading Higgs form is strictly positive
on its unit sphere.  Its angular zero set is empty.  The V9--V11
regular, monomial, and crossing analyses are therefore unnecessary
for the explicitly displayed pure Higgs leader.
\end{corollary}

\begin{proof}
The lower bound (12.6) is positive when \(r=1\).
\end{proof}

The exact floor \(\lambda_H/N\) tends to zero when the number of
sites grows.  Thus fixed-cutoff positivity does not supply the
cutoff-uniform coercivity required by the continuum compiler.

\subsection{Coercivity with a signed mass}

Write \(a=m_H^2\in\mathbb R\) and
\(a_-=\max\{0,-a\}\).  Assume \(\kappa_H\geq0\).

\begin{theorem}[Finite-cutoff Higgs coercivity]
\label{thm:v17v12-coercivity}
The action in (12.2), now on the unbounded space
\((\mathbb R^4)^N\), obeys
\[
 S_{H,n}(U,H)
 \geq \frac{\lambda_H}{N}r^4-a_-r^2
 \geq-\frac{Na_-^2}{4\lambda_H}.                  \tag{12.8}
\]
More usefully for tails,
\[
 S_{H,n}(U,H)
 \geq
 \frac{\lambda_H}{2N}r^4
 -\frac{Na_-^2}{2\lambda_H}.                      \tag{12.9}
\]
Both bounds are uniform in the compact gauge field \(U\).
\end{theorem}

\begin{proof}
The covariant edge square is nonnegative, and (12.7) gives the
first inequality.  Minimize
\((\lambda_H/N)s^2-a_-s\) over \(s=r^2\) for the second inequality
in (12.8).  To prove (12.9), retain half of the quartic and minimize
\((\lambda_H/(2N))s^2-a_-s\).
\end{proof}

\subsection{The V1 exponential-gradient condition at fixed cutoff}

Let \(d_{\max}\) be the maximum vertex degree.  Since
\[
 \|D_UH\|_{\ell^2(E_n)}^2\leq2d_{\max}r^2,
\]
the positive operator \(D_U^*D_U\) has norm at most
\(2d_{\max}\).

\begin{lemma}[Cubic gradient bound]
\label{lem:v17v12-gradient}
For every \(U\),
\[
 \|\nabla_HS_{H,n}(U,H)\|
 \leq
 4\lambda_Hr^3+
 (4\kappa_Hd_{\max}+2|a|)r.                       \tag{12.10}
\]
\end{lemma}

\begin{proof}
The quartic gradient has squared norm
\[
 16\lambda_H^2\sum_x\|H_x\|^6
 \leq16\lambda_H^2r^6.
\]
The edge gradient is \(2\kappa_HD_U^*D_UH\), and the mass gradient
is \(2aH\).  Apply the operator bound and the triangle inequality.
\end{proof}

\begin{theorem}[Fixed-cutoff exponential-gradient moments]
\label{thm:v17v12-exp-gradient}
Let
\[
 \dd\mu_{n,U}(H)=Z_{n,U}^{-1}e^{-S_{H,n}(U,H)}\,\dd H.
\]
For every fixed \(n\), every \(\alpha>0\), and every polynomial
\(p\),
\[
 \int |p(H)|
 \exp\{\alpha\|\nabla_HS_{H,n}(U,H)\|\}
 \,\dd\mu_{n,U}(H)<\infty,                         \tag{12.11}
\]
uniformly over \(U\).
\end{theorem}

\begin{proof}
By (12.9)--(12.10), the unnormalized integrand is bounded by a
constant, uniform in \(U\), times
\[
 |p(H)|\exp\{-c_nr^4+C_{\alpha,n}(r^3+r)\},
 \qquad c_n=\lambda_H/(2N)>0.
\]
The negative quartic dominates the cubic and linear terms as
\(r\to\infty\), so this function is integrable on
\(\mathbb R^{4N}\), with a majorant independent of \(U\).  On the
unit ball, unitarity of the link representation and finiteness of the
graph give \(S_{H,n}(U,H)\leq C_n\) uniformly in \(U\).  Hence
\(Z_{n,U}\) is bounded below by the ball volume times \(e^{-C_n}\),
and (12.9) gives a uniform finite upper bound.  Dividing the
uniform numerator bound by this positive lower bound proves the claim.
\end{proof}

Finite determinant or Pfaffian modulus factors of polynomial growth
can be included in (12.11).  Thus the V1 sufficient moment
criterion is verified for the isolated fixed-cutoff Higgs shell.
The proof gives constants depending on \(N\), so it does not close a
uniform shell limit.

\subsection{Why the total leading polynomial is not extractable}

\begin{proposition}[Underdetermination of the RPESM leader]
\label{prop:v17v12-underdetermined}
The declarations in \textup{(RP.5)} do not determine a unique
highest homogeneous form for the full bosonic action.
\end{proposition}

\begin{proof}
Starting from any action satisfying the stated RPESM assumptions,
add, for \(\epsilon>0\),
\[
 \epsilon\sum_{\langle xy\rangle}
 \|H_x\|^2\|H_y\|^2.                               \tag{12.12}
\]
This term is finite range, gauge invariant, reflection even, and
nonnegative.  It preserves the positive on-site quartic and every
declared domination of negative terms, but changes the quartic
coefficient tensor.  Hence the written premises admit distinct
leading forms.  The separate symbols
\(S_{g,n},S_{{\rm gf},n},S_{{\rm med},n},S_{{\rm ct},n}\) are also
not expanded in a common unbounded field basis.
\end{proof}

The exact Wilsonian marginal presents a second obstruction to an
exact homogeneous decomposition: its proved nonzero
\(R_{\geq6}\) tail is part of the action and cannot be discarded
without a quantitative remainder theorem.

\begin{definition}[Missing physical coefficient packet]
\label{def:v17v12-missing-packet}
At cutoff \(n\), the missing packet consists of:
\begin{enumerate}
\item a real basis \(z_n\) for every unbounded bosonic mode,
      including the retained gravitational/conformal modes;
\item the complete coefficient tensors
      \(P_{m,n},P_{m-1,n},\ldots\) in
      \(S_{B,n}(z_n)=\sum_{k\leq m}P_{k,n}(z_n)\), or an exact
      nonpolynomial replacement with derivative bounds;
\item numerical domination margins for every signed sector and
      its first three interpolation derivatives;
\item cutoff scaling sufficient to make the resulting exponential
      moments and shell errors summable.
\end{enumerate}
\end{definition}

\begin{assessment}[What V12 closes]
\label{ass:v17v12-status}
V12 extracts the physical Higgs leader and its first lower form.
The leader is a strictly positive radial quartic with exact angular
floor \(\lambda_H/N\), so singular angular geometry is not the
finite-cutoff Higgs obstruction.  The isolated Higgs law satisfies
all exponential-gradient moments at each fixed cutoff, even with a
signed mass term.

The full RPESM source does not specify a unique total quartic tensor
or a common unbounded Einstein--matter field basis.  Its domination
and continuum claims therefore remain conditional at exactly the
packet in Definition~\ref{def:v17v12-missing-packet}.
\end{assessment}

\begin{decision}[V13: replace the global \(1/N\) floor]
\label{dec:v17v12-next}
The fixed-cutoff global angular estimate loses a factor \(N\).
Derive volume-uniform conditional exponential moments from the
on-site quartic and bounded-degree local interactions.  Express the
result as a Dobrushin-style row-sum condition that can be checked
without knowing the missing gravitational coefficient packet.
\end{decision}

\begin{firewall}[Claim boundary after V12]
\label{fire:v17v12}
V12 proves fixed-cutoff Higgs coercivity and identifies the exact
source underdetermination.  It does not promote the RPESM
domination assumption to a theorem or establish cutoff-uniform
Einstein--SM tightness.
\end{firewall}

\section*{Version V12 append-only record}
\addcontentsline{toc}{section}{Version V12 append-only record}
The complete V11 source was retained byte for byte before this
continuation.  It had \(118327\) bytes, \(3802\) lines, and SHA--256
\[
 \texttt{22E0C661CF8A2391D2BD2865F3B5D1EF25604ADED6EDF0C8E2A56A04A729376C}.
\]
V12 extracted the positive Higgs leader, proved its exact angular
floor and fixed-cutoff exponential-gradient moments, and isolated
the missing total-action coefficient packet.


\section{V13: volume-uniform conditional quartic moments}
\label{sec:v17v13}

\subsection{A bounded-degree Higgs graph}

Let \(G=(V,E)\) be any finite graph of maximum incident degree at
most \(d_*\).  Put \(H_x\in\mathbb R^q\), let every oriented edge
carry an orthogonal map \(U_{xy}\), and define
\[
 S_G(H;U)
 =\sum_{x\in V}\bigl(\lambda\|H_x\|^4+a\|H_x\|^2\bigr)
 +\kappa\sum_{\{x,y\}\in E}
 \|H_y-U_{xy}H_x\|^2,                             \tag{13.1}
\]
where
\[
 \lambda>0,\qquad \kappa\geq0,\qquad a\in\mathbb R. \tag{13.2}
\]
Let \(\mu_{G,U}\) be the normalized Lebesgue-density Gibbs law
proportional to \(e^{-S_G}\).  V12 proves its finite-volume
normalizability.  The purpose here is to remove all dependence on
\(|V|\) from local moments.

\subsection{The one-site conditional estimate}

Condition on every field except \(H_x\), and write the conditional
energy as
\[
 E_x(h\mid\eta)
 =\lambda\|h\|^4+a\|h\|^2
 +\kappa\sum_{y\sim x}\|\eta_y-U_{xy}h\|^2.        \tag{13.3}
\]
Terms independent of \(h\) may be retained because they cancel
between numerator and denominator.

\begin{lemma}[Conditional quartic recursion]
\label{lem:v17v13-conditional}
Fix
\[
 0<\theta<\lambda/2,\qquad \delta>0.
\]
There is a finite
\(C=C(q,\lambda,\kappa,a,d_*,\theta,\delta)\), independent of
\(G,U,x\), and \(\eta\), such that
\[
 \mathbb E_{\mu_{G,U}}
 \left[e^{\theta\|H_x\|^4}\mid H_{V\setminus\{x\}}=\eta\right]
 \leq
 C\exp\left\{\delta\sum_{y\sim x}\|\eta_y\|^4\right\}.             \tag{13.4}
\]
\end{lemma}

\begin{proof}
Let \(a_-=\max\{0,-a\}\).  Since the edge squares are nonnegative,
\[
 E_x(h\mid\eta)
 \geq\frac{\lambda}{2}\|h\|^4-\frac{a_-^2}{2\lambda}.             \tag{13.5}
\]
Therefore the unnormalized numerator in (13.4) is at most
\[
 e^{a_-^2/(2\lambda)}
 \int_{\mathbb R^q}
 e^{-(\lambda/2-\theta)\|h\|^4}\,\dd h
 =:C_1<\infty.                                      \tag{13.6}
\]
On the unit ball,
\[
 E_x(h\mid\eta)
 \leq
 \lambda+|a|+2\kappa d_*
 +2\kappa\sum_{y\sim x}\|\eta_y\|^2.                \tag{13.7}
\]
If \(v_q\) is the unit-ball volume, the conditional partition
function is consequently at least
\[
 v_qe^{-\lambda-|a|-2\kappa d_*}
 \exp\left\{-2\kappa\sum_{y\sim x}\|\eta_y\|^2\right\}.             \tag{13.8}
\]
For every \(t\geq0\),
\[
 2\kappa t^2\leq\delta t^4+\frac{\kappa^2}{\delta}.                \tag{13.9}
\]
Divide (13.6) by (13.8), apply (13.9) to each neighbor,
and absorb at most \(d_*\kappa^2/\delta\) into \(C\).
\end{proof}

\subsection{Closing the row-sum recursion}

\begin{theorem}[Volume-uniform one-site exponential moment]
\label{thm:v17v13-uniform}
For every \(0<\theta<\lambda/2\) there is
\[
 M_\theta<\infty                                      \tag{13.10}
\]
depending only on
\((q,\lambda,\kappa,a,d_*,\theta)\), such that for every finite
\(G\), every orthogonal link field \(U\), and every \(x\in V\),
\[
 \mathbb E_{\mu_{G,U}}e^{\theta\|H_x\|^4}
 \leq M_\theta.                                      \tag{13.11}
\]
\end{theorem}

\begin{proof}
Choose \(0<\delta<\theta/d_*\), with the assertion immediate when
\(d_*=0\), and put
\[
 c_{xy}=
 \begin{cases}
 \delta/\theta,&y\sim x,\\
 0,&y\not\sim x.
 \end{cases}
 \qquad
 \gamma:=\sup_x\sum_yc_{xy}<1.                     \tag{13.12}
\]
Let
\[
 M_x=\mathbb E e^{\theta\|H_x\|^4}.
\]
Finite-volume coercivity makes every \(M_x\) finite.  Taking
expectations in (13.4) and applying generalized H\"older, with a
dummy factor \(1\) of exponent weight \(1-\sum_yc_{xy}\), gives
\[
 M_x
 \leq C\,
 \mathbb E\prod_y
 \left(e^{\theta\|H_y\|^4}\right)^{c_{xy}}
 \leq C\prod_yM_y^{c_{xy}}.                       \tag{13.13}
\]
If \(M=\max_xM_x\), then \(M\leq CM^\gamma\), and hence
\[
 M\leq C^{1/(1-\gamma)}.                           \tag{13.14}
\]
This bound contains no volume parameter.
\end{proof}

\begin{remark}[What the row sum means]
\label{rem:v17v13-not-uniqueness}
The number \(\gamma\) in (13.12) is a moment-propagation row sum.
It proves integrability, not Dobrushin uniqueness or decay of
correlations.  The freedom to choose a small \(\delta\) enlarges
the constant \(C\) but leaves it volume independent.
\end{remark}

\subsection{Fixed-support joint moments}

\begin{corollary}[Uniform finite-block quartic moments]
\label{cor:v17v13-block}
Let \(B\subset V\) have at most \(K\) vertices.  For every
\(0<\vartheta<\lambda/(2K)\),
\[
 \sup_{G,U,B}
 \mathbb E_{\mu_{G,U}}
 \exp\left\{\vartheta\sum_{x\in B}\|H_x\|^4\right\}
 <\infty,                                           \tag{13.15}
\]
where the supremum is over graphs of degree at most \(d_*\).
\end{corollary}

\begin{proof}
Apply H\"older with \(K\) factors and use
Theorem~\ref{thm:v17v13-uniform} at parameter \(K\vartheta\).
\end{proof}

\subsection{Uniform local gradient moments}

Let \(g_x=\nabla_{H_x}S_G\).  Direct differentiation gives
\[
 \|g_x\|
 \leq
 4\lambda\|H_x\|^3
 +(2|a|+2\kappa d_*)\|H_x\|
 +2\kappa\sum_{y\sim x}\|H_y\|.                    \tag{13.16}
\]

\begin{theorem}[Local exponential-gradient bound]
\label{thm:v17v13-local-gradient}
For every \(\alpha>0\) and every \(K<\infty\), there is
\(C_{\alpha,K}<\infty\), independent of \(G,U\), such that
\[
 \mathbb E_{\mu_{G,U}}
 \exp\left\{\alpha\sum_{x\in B}\|g_x\|\right\}
 \leq C_{\alpha,K}                                  \tag{13.17}
\]
whenever \(|B|\leq K\).
\end{theorem}

\begin{proof}
The right side of (13.16), summed over \(B\), depends only on the
set \(B\) and its neighbors, which has at most \(K(1+d_*)\)
vertices.  For every \(\epsilon>0\), Young's inequality bounds each
linear or cubic field power by
\[
 c_3t^3+c_1t\leq\epsilon t^4+C_{\epsilon,c_1,c_3}.                 \tag{13.18}
\]
Choose \(\epsilon\) so that the total quartic coefficient over this
fixed block is below the threshold in
Corollary~\ref{cor:v17v13-block}.  Apply that corollary and absorb
the finitely many constants.
\end{proof}

\subsection{Import into the shell cumulant criterion}

\begin{corollary}[Uniform local tilted third cumulants]
\label{cor:v17v13-cumulant}
Let \(X_{n}\) be any shell interpolation derivative supported on at
most \(K\) Higgs sites, polynomial of degree at most three in those
sites and their first covariant differences, with coefficients
bounded independently of \(n\).  On graphs with a common degree
bound and parameters in a compact subset of
\[
 \lambda>0,\qquad\kappa\geq0,\qquad a\in\mathbb R,
\]
there exist \(t_0>0\) and \(C_3<\infty\) such that
\[
 \sup_n\sup_{|t|\leq t_0}
 \left|\operatorname{cum}_{3,t}(X_n)\right|
 \leq C_3.                                         \tag{13.19}
\]
\end{corollary}

\begin{proof}
The polynomial and support hypotheses give
\(|X_n|\leq C(1+\sum_{x\in B_n}\|H_x\|^3)\) on a uniformly bounded
neighborhood block.  Theorem~\ref{thm:v17v13-local-gradient} and
Young's inequality give a uniform exponential moment of
\(|X_n|\).  The V1 tilted-moment argument then bounds the third
centered moment, hence the third cumulant, uniformly for
\(|t|\leq t_0\).
\end{proof}

Thus a local shell remainder of size \(\varepsilon_n\) has a
uniform \(O(\varepsilon_n^3)\) third-cumulant debit.  Any additional
counting factor and the summability condition remain those isolated
in V3.

\subsection{Extension gate for the coupled action}

The preceding proof survives extra local sectors if their
one-site conditional law obeys, for a common \(\theta>0\),
\[
 \mathbb E\left[e^{\theta\|H_x\|^4}\mid{\cal F}_{x^c}\right]
 \leq C
 \exp\left\{\theta\sum_yc_{xy}\|H_y\|^4\right\},
 \qquad
 \sup_x\sum_yc_{xy}<1.                             \tag{13.20}
\]
Equation (13.20), with its explicit constants and cutoff scaling,
is a sufficient row of the missing packet in V12.  The current
RPESM quartic-domination sentence does not specify (13.20), because
it gives neither conditional coefficients nor a row-sum margin.

\begin{assessment}[What V13 closes]
\label{ass:v17v13-status}
V13 replaces the deteriorating global angular floor
\(\lambda/N\) by volume-uniform one-site and finite-block
exponential quartic moments for the explicit Higgs
quartic-plus-hopping action.  It consequently verifies the V1
uniform tilted third-cumulant hypothesis for bounded-support
degree-three Higgs shell derivatives.

The theorem is an integrability result.  It does not prove
uniqueness, a mass gap, tightness at a selected continuum scaling,
or the coupled gravitational conditional inequality (13.20).
\end{assessment}

\begin{decision}[V14: conditional shell transport]
\label{dec:v17v13-next}
Use (13.19) in an exact adjacent-shell log-Laplace calculation.
Separate the number of local shell cells, the interpolation size,
and the conditional row-sum margin, and derive a regulator-uniform
summability criterion for local Higgs counterterms.  Keep the
gravitational coefficient packet as a named external gate.
\end{decision}

\begin{firewall}[Claim boundary after V13]
\label{fire:v17v13}
V13 closes a pure-Higgs local moment bottleneck.  It does not verify
the full RPESM action, fermion-line singularities, or any continuum
nontriviality condition.
\end{firewall}

\section*{Version V13 append-only record}
\addcontentsline{toc}{section}{Version V13 append-only record}
The complete V12 source was retained byte for byte before this
continuation.  It had \(129178\) bytes, \(4115\) lines, and SHA--256
\[
 \texttt{85B2A798CE7986FD8E9D74C6D64F6304A1C60ECC7312809B857D362346FC1C8A}.
\]
V13 proved volume-uniform conditional quartic, local-gradient, and
tilted third-cumulant bounds for the explicit bounded-degree Higgs
action.


\section{V14: conditional log-Laplace shell transport}
\label{sec:v17v14}

\subsection{One conditional normalizer}

Let \({\cal G}\) be a conditioning sigma algebra and let \(X\) be a
real random variable.  For \(s\) in an interval about zero, define
\[
 c(s)=\log\mathbb E[e^{-sX}\mid{\cal G}],           \tag{14.1}
\]
and let \(\mathbb E_s[\cdot\mid{\cal G}]\) denote the conditional
law tilted by \(e^{-sX-c(s)}\).

\begin{lemma}[Exact third-order conditional expansion]
\label{lem:v17v14-loglaplace}
Suppose the conditional exponential moments needed below are
finite.  Then
\[
\begin{aligned}
 c'(s)&=-\mathbb E_s[X\mid{\cal G}],\\
 c''(s)&=\operatorname{Var}_s(X\mid{\cal G}),\\
 c'''(s)&=-\operatorname{cum}_{3,s}(X\mid{\cal G}).
\end{aligned}                                                       \tag{14.2}
\]
Consequently
\[
 c(\varepsilon)
 =-\varepsilon\mathbb E[X\mid{\cal G}]
 +\frac{\varepsilon^2}{2}
   \operatorname{Var}(X\mid{\cal G})
 +R(\varepsilon),                                  \tag{14.3}
\]
where
\[
 R(\varepsilon)
 =-\frac12\int_0^\varepsilon
 (\varepsilon-s)^2
 \operatorname{cum}_{3,s}(X\mid{\cal G})\,\dd s.   \tag{14.4}
\]
For negative \(\varepsilon\), the integral is oriented in the usual
way.
\end{lemma}

\begin{proof}
Differentiate the conditional moment generating function and
divide by it.  The first derivative is the tilted mean of \(-X\);
successive derivatives are its second and third cumulants.  Taylor's
formula with integral remainder gives (14.3)--(14.4).
\end{proof}

\begin{corollary}[Conditional remainder debit]
\label{cor:v17v14-one-debit}
If
\[
 \sup_{|s|\leq|\varepsilon|}
 \mathbb E\left[
 |\operatorname{cum}_{3,s}(X\mid{\cal G})|
 \right]\leq C_3,                                  \tag{14.5}
\]
then
\[
 \mathbb E|R(\varepsilon)|
 \leq\frac{C_3}{6}|\varepsilon|^3.                 \tag{14.6}
\]
\end{corollary}

\begin{proof}
Take absolute values in (14.4), use Tonelli, and evaluate
\(\frac12\int_0^{|\varepsilon|}(|\varepsilon|-s)^2\dd s\).
\end{proof}

\subsection{A sequential local shell}

At adjacent cutoff \(n\to n+1\), order the \(M_n\) new local cells.
Let
\[
 q_{n,j}(\dd y_j\mid{\cal G}_{n,j-1}),
 \qquad 1\leq j\leq M_n,                           \tag{14.7}
\]
be reference conditional kernels, where \({\cal G}_{n,j-1}\)
contains the retained field and the first \(j-1\) new cells.  Let
\(X_{n,j}\) be a local shell writer and let
\[
\begin{aligned}
 c_{n,j}
 &=\log\int e^{-\varepsilon_nX_{n,j}}\,
 q_{n,j}(\dd y_j\mid{\cal G}_{n,j-1}),\\
 k_{n,j}(\dd y_j\mid{\cal G}_{n,j-1})
 &=e^{-\varepsilon_nX_{n,j}-c_{n,j}}\,
 q_{n,j}(\dd y_j\mid{\cal G}_{n,j-1}).
\end{aligned}                                                       \tag{14.8}
\]

\begin{theorem}[Exact normalized shell composition]
\label{thm:v17v14-exact-shell}
Every \(k_{n,j}\) is a probability kernel, and
\[
 K_{n+1/n}=k_{n,1}k_{n,2}\cdots k_{n,M_n}           \tag{14.9}
\]
is a probability kernel from the retained field to the complete
new shell.  Relative to the reference product of conditional
kernels, its exact logarithmic density is
\[
 -\varepsilon_n\sum_{j=1}^{M_n}X_{n,j}
 -C_n,
 \qquad
 C_n=\sum_{j=1}^{M_n}c_{n,j}.                      \tag{14.10}
\]
\end{theorem}

\begin{proof}
The definition of \(c_{n,j}\) makes the integral of \(k_{n,j}\)
equal to one for every conditioning history.  Iterated integration
therefore normalizes (14.9).  Multiplying the densities and taking
their logarithm gives (14.10).
\end{proof}

Thus exact projective transport requires the whole conditional
normalizer \(C_n\).  It is not legitimate to replace it by a fitted
unconditional constant.

\subsection{The finite-head approximation}

Define the second-order local head
\[
 h_{n,j}
 =-\varepsilon_n
   \mathbb E[X_{n,j}\mid{\cal G}_{n,j-1}]
 +\frac{\varepsilon_n^2}{2}
   \operatorname{Var}(X_{n,j}\mid{\cal G}_{n,j-1}),                \tag{14.11}
\]
and set
\[
 H_n=\sum_{j=1}^{M_n}h_{n,j},
 \qquad
 {\cal R}_n=C_n-H_n.                               \tag{14.12}
\]

\begin{theorem}[Shell counting bound]
\label{thm:v17v14-counting}
Assume that for every \(n,j\) and
\(|s|\leq|\varepsilon_n|\),
\[
 \mathbb E
 \left[
 |\operatorname{cum}_{3,s}
 (X_{n,j}\mid{\cal G}_{n,j-1})|
 \right]\leq C_{3,n}.                              \tag{14.13}
\]
Then
\[
 \mathbb E|{\cal R}_n|
 \leq\frac{M_nC_{3,n}}{6}|\varepsilon_n|^3.        \tag{14.14}
\]
\end{theorem}

\begin{proof}
Apply Corollary~\ref{cor:v17v14-one-debit} at each conditional
stage and sum by the triangle inequality.
\end{proof}

No independence between distinct shell cells is used in (14.14).
Dependence is carried by the conditioning filtration.

\subsection{Summability and scaling thresholds}

\begin{theorem}[Convergent counterterm tail]
\label{thm:v17v14-summability}
Suppose all shells are realized on one projective probability space
and
\[
 \sum_{n=1}^{\infty}
 M_nC_{3,n}|\varepsilon_n|^3<\infty.               \tag{14.15}
\]
Then
\[
 \sum_n\mathbb E|{\cal R}_n|<\infty,\qquad
 \sum_n|{\cal R}_n|<\infty\quad\text{almost surely}.               \tag{14.16}
\]
In particular, the accumulated difference between the exact
conditional counterterm and its second-order head converges in
\(L^1\) and almost surely.
\end{theorem}

\begin{proof}
The first assertion follows by summing (14.14).  Tonelli gives
\(\mathbb E\sum_n|{\cal R}_n|<\infty\), so the nonnegative random
series is finite almost surely.  Its tails tend to zero in \(L^1\)
by the same summable majorant.
\end{proof}

\begin{corollary}[Polynomial and dyadic schedules]
\label{cor:v17v14-rates}
Assume \(\sup_nC_{3,n}<\infty\).
\begin{enumerate}
\item If \(M_n\leq Cn^\beta\) and
      \(|\varepsilon_n|\leq Cn^{-\gamma}\), then (14.15) holds when
      \[
      3\gamma-\beta>1.                             \tag{14.17}
      \]
\item If \(M_n\leq C2^{dn}\) and
      \(|\varepsilon_n|\leq C2^{-\alpha n}\), then (14.15) holds
      when
      \[
      3\alpha>d.                                   \tag{14.18}
      \]
\end{enumerate}
At equality these estimates give a nonsummable harmonic or
constant-order debit; cancellation would require a new theorem.
\end{corollary}

\begin{proof}
The summands in (14.15) are bounded respectively by constants times
\(n^{\beta-3\gamma}\) and \(2^{(d-3\alpha)n}\).
\end{proof}

\subsection{Insertion of the conditional row-sum margin}

For shell \(n\), let
\[
 \eta_n
 =1-\sup_x\sum_yc_{xy}^{(n)}>0                     \tag{14.19}
\]
be the moment row-sum margin from V13.  Denote by
\({\mathfrak C}_3(\eta_n)\) the third-cumulant bound obtained by
the V13 conditional proof for the allowed compact parameter set and
support radius.

\begin{corollary}[Higgs shell criterion]
\label{cor:v17v14-higgs}
Suppose:
\begin{enumerate}
\item every \(X_{n,j}\) is supported on at most \(K\) Higgs sites,
      has degree at most three, and has uniformly bounded
      coefficients;
\item graph degrees and the parameters
      \((\lambda,\kappa,a)\) remain in the V13 compact range;
\item \(\inf_n\eta_n>0\);
\item
      \[
      \sum_nM_n|\varepsilon_n|^3<\infty.           \tag{14.20}
      \]
\end{enumerate}
Then the exact Higgs shell counterterm has the convergent
second-order decomposition (14.16).
\end{corollary}

\begin{proof}
V13 supplies a common tilted third-cumulant bound because the
support radius, coefficient range, and row-sum margin are uniform.
Thus \(C_{3,n}\leq C_3\), and
Theorem~\ref{thm:v17v14-summability} applies.
\end{proof}

If the margins are not uniform, the sufficient condition is the
more precise
\[
 \sum_n
 M_n{\mathfrak C}_3(\eta_n)|\varepsilon_n|^3<\infty.               \tag{14.21}
\]
This formula prevents a closing row-sum margin from being hidden
inside big-\(O\) notation.

\subsection{What this transport does and does not identify}

The first term in (14.11) is the conditional mean action writer.
The second is its conditional variance, hence the local
second-order counterterm.  Higher conditional connected responses
are collected in \({\cal R}_n\).  Equation (14.16) proves that their
total absolute contribution is finite under (14.15).

The theorem does not identify the limiting finite head with the
physical Standard-Model coefficient ray.  Such identification
requires matching the conditional means and variances to the
declared gauge, Higgs, gravity, and counterterm basis.  It also does
not prove total-variation equivalence between a measure formed with
the truncated head and the exact projective law; that stronger
claim requires uniform exponential control of the accumulated
random remainder.

\begin{assessment}[What V14 closes]
\label{ass:v17v14-status}
V14 turns the V13 local cumulant bound into an exact adjacent-shell
calculation.  Sequential conditional normalization preserves
projectivity without independence, and the second-order truncation
has total \(L^1\) debit bounded by
\(M_nC_{3,n}|\varepsilon_n|^3/6\).  The polynomial threshold
\(3\gamma-\beta>1\) and dyadic threshold \(3\alpha>d\) are now
proved with the shell-cell count exposed.

The calculation applies to the explicit local Higgs shell.  The
full Einstein--SM action still lacks the conditional coefficient
packet and row-sum margin required to invoke (14.21).
\end{assessment}

\begin{decision}[V15: companion audit and coupled-row gate]
\label{dec:v17v14-next}
Perform the mandatory YM/NS audit.  Then formulate the minimal
block conditional inequality for Higgs, gauge, gravitational, and
mediator variables which would extend (13.20) and (14.21) to the
full RPESM action.  Continue in V16 with the first block that can be
certified from the source formulas.
\end{decision}

\begin{firewall}[Claim boundary after V14]
\label{fire:v17v14}
V14 proves summability of the random counterterm remainder in
\(L^1\) and almost surely.  It does not prove exponential
equivalence of exact and truncated laws, coefficient matching, or
the interacting continuum limit.
\end{firewall}

\section*{Version V14 append-only record}
\addcontentsline{toc}{section}{Version V14 append-only record}
The complete V13 source was retained byte for byte before this
continuation.  It had \(138699\) bytes, \(4409\) lines, and SHA--256
\[
 \texttt{7402AFD465CCF30C101815D37973850E8BC40A14FB7DADFAD3DBC8BAF944BF43}.
\]
V14 proved exact sequential shell normalization and the
shell-counted third-order counterterm summability criterion.


\section{V15: companion audit and the coupled block-conditional gate}
\label{sec:v17v15}

\subsection{Mandatory companion audit}

The checkpoint files are
\[
\begin{array}{lll}
\text{YM V73:}&16600\text{ lines},&
\texttt{7370912684CDE733ECDE016D6C9E8DCC155F19E7A275A906F9222541AF99149D},\\
\text{NS V26:}&5319\text{ lines},&
\texttt{82C887F8C2C69E4F28FAD7C3DC8DE5B9099CB5A4BF431EBA1FFF3E91935978AD}.
\end{array}                                                       \tag{15.1}
\]
YM V72 repairs the two V71 third-coordinate obstructions with
rational half-step bridges and proves the fifth direct-plus
two-step slab, with the inverted direct-minus companion.  YM V73
then certifies twelve of sixteen terminal centers and records exact
failure witnesses for the four remaining sufficient splittings
\(223,233,323,333\).  Those witnesses do not prove negativity of
the physical Wilson pencil.  No complete quarter cell or continuum
measure is claimed.  NS V26 is unchanged and concerns pointwise
rank-one norms of the first two Oseen-kernel derivatives.

\begin{theorem}[V15 audit decision]
\label{thm:v17v15-audit}
Neither companion file supplies a conditional exponential-moment
inequality for the coupled Einstein--SM blocks.  The V14 direction
must therefore be continued from the SM source formulas.
\end{theorem}

\begin{proof}
The new YM statements are finite rational positivity charts for a
Wilson pencil.  The NS statement is a kernel tensor-norm
calculation.  Neither gives a Gibbs conditional energy, a
Lyapunov transfer matrix, or a row-sum margin of the form (13.20).
\end{proof}

\subsection{Abstract block system}

At cutoff \(n\), let \(I_n\) be a finite set of local blocks.  A
block may be a Higgs site, a compact gauge link, a gravitational
cell, a mediator cell, or another retained bosonic coordinate.  Let
\[
 Z_i\in{\cal Z}_i,\qquad
 V_i:{\cal Z}_i\longrightarrow[0,\infty)            \tag{15.2}
\]
be its state and Lyapunov function.  Compact blocks may use
\(V_i=0\), with bounded observables absorbed into constants.

\begin{definition}[Coupled conditional moment certificate]
\label{def:v17v15-certificate}
A family of Gibbs laws has a certificate
\((\theta_i,C_i,A_n)\) when \(\theta_i>0\),
\(A_n=(a_{ij})_{i,j\in I_n}\) is nonnegative, and
\[
 \mathbb E\left[
 e^{\theta_iV_i(Z_i)}\mid{\cal F}_{i^c}\right]
 \leq
 C_i\exp\left\{
 \sum_{j\in I_n}a_{ij}\theta_jV_j(Z_j)\right\}       \tag{15.3}
\]
almost surely for every \(i\).
\end{definition}

\begin{theorem}[Block row-sum closure]
\label{thm:v17v15-row}
Suppose
\[
 C_*:=\sup_{n,i}C_i<\infty,\qquad
 \gamma_*:=\sup_n\sup_{i\in I_n}\sum_ja_{ij}<1.     \tag{15.4}
\]
Then
\[
 \sup_n\sup_{i\in I_n}
 \mathbb E e^{\theta_iV_i(Z_i)}
 \leq C_*^{1/(1-\gamma_*)}.                         \tag{15.5}
\]
\end{theorem}

\begin{proof}
Put \(M_i=\mathbb E e^{\theta_iV_i}\).  Take expectations in
(15.3).  Generalized H\"older, with a dummy factor \(1\) of weight
\(1-\sum_ja_{ij}\), gives
\[
 M_i\leq C_i\prod_jM_j^{a_{ij}}.                   \tag{15.6}
\]
For a fixed finite cutoff let \(M=\max_iM_i\).  Finite-volume
coercivity makes \(M<\infty\), and (15.4)--(15.6) yield
\(M\leq C_*M^{\gamma_*}\).  Rearrangement proves (15.5), uniformly
in the cutoff.
\end{proof}

\subsection{Deriving the certificate from conditional energies}

Let \(E_i(z\mid\eta)\) be a version of the conditional energy for
block \(i\), with reference measure \(\dd z\).  Let
\({\cal A}_i\subset{\cal Z}_i\) be an anchor set of measure at least
\(v_i>0\).

\begin{theorem}[Energy-to-moment compiler]
\label{thm:v17v15-energy-compiler}
Assume there are numbers
\[
 \ell_i>\theta_i>0,\quad K_i<\infty,\quad
 d_{ij},e_{ij}\geq0                                \tag{15.7}
\]
such that
\[
 E_i(z\mid\eta)
 \geq
 \ell_iV_i(z)-K_i-\sum_jd_{ij}V_j(\eta_j)           \tag{15.8}
\]
for all \(z\), while on the anchor
\[
 E_i(z\mid\eta)
 \leq K_i+\sum_je_{ij}V_j(\eta_j).                 \tag{15.9}
\]
Suppose also
\[
 J_i:=\int_{{\cal Z}_i}
 e^{-(\ell_i-\theta_i)V_i(z)}\,\dd z<\infty.        \tag{15.10}
\]
Then (15.3) holds with
\[
 C_i=v_i^{-1}J_i e^{2K_i},\qquad
 a_{ij}=\frac{d_{ij}+e_{ij}}{\theta_j}.            \tag{15.11}
\]
\end{theorem}

\begin{proof}
The numerator of the conditional expectation in (15.3) is at most
\[
 e^{K_i+\sum_jd_{ij}V_j(\eta_j)}J_i.
\]
Integration over the anchor and (15.9) bound the conditional
partition function below by
\[
 v_i e^{-K_i-\sum_je_{ij}V_j(\eta_j)}.
\]
Their ratio is (15.3) with (15.11).
\end{proof}

This theorem exposes two different debits.  The coefficients
\(d_{ij}\) measure how much boundary fields weaken the coercive
lower bound.  The \(e_{ij}\) measure the boundary cost of finding
one positive-measure normalization anchor.  Both enter the same
influence matrix.

\subsection{Spectral-radius formulation}

Often an energy calculation first returns a nonnegative transfer
matrix \(B=(b_{ij})\) before the exponential weights are chosen.
Rescaling \(\theta_i=p_i\bar\theta_i\) changes the H\"older matrix to
\[
 A=D_pBD_p^{-1},
 \qquad D_p=\operatorname{diag}(p_i).               \tag{15.12}
\]

\begin{theorem}[Perron rescaling criterion]
\label{thm:v17v15-perron}
For a finite nonnegative matrix \(B\), the following are equivalent:
\begin{enumerate}
\item \(\rho(B)<1\);
\item there are \(p_i>0\) and \(\gamma<1\) such that
      \[
      \sum_j\frac{p_i}{p_j}b_{ij}\leq\gamma
      \quad\text{for every }i;                     \tag{15.13}
      \]
\item \(I-B\) has a nonnegative inverse
      \(\sum_{k\geq0}B^k\).
\end{enumerate}
\end{theorem}

\begin{proof}
If \(\rho(B)<1\), choose
\[
 q=(I-B)^{-1}{\bf1}>{\bf0}.
\]
Then \(Bq=q-{\bf1}\leq\gamma q\), where
\(\gamma=1-1/\max_iq_i<1\).  Put \(p_i=q_i^{-1}\);
then (15.13) is exactly
\((Bq)_i/q_i\leq\gamma\).  Conversely, (15.13) says that the
similar matrix \(D_pBD_p^{-1}\) has infinity norm below one, so
\(\rho(B)<1\).  The Neumann series equivalence is standard and
follows directly from convergence in a matrix norm after such a
similarity.
\end{proof}

\begin{corollary}[Exact two-sector cycle condition]
\label{cor:v17v15-two-sector}
For
\[
 B=\begin{pmatrix}a&b\\c&d\end{pmatrix},
 \qquad a,b,c,d\geq0,
\]
one has \(\rho(B)<1\) exactly when
\[
 a<1,\qquad d<1,\qquad
 bc<(1-a)(1-d).                                    \tag{15.14}
\]
\end{corollary}

\begin{proof}
The eigenvalue \(1\) lies above the Perron root precisely when
\(I-B\) is a nonsingular \(M\)-matrix.  In dimension two this is
equivalent to positivity of its two diagonal entries and determinant:
\[
 1-a>0,\quad1-d>0,\quad
 \det(I-B)=(1-a)(1-d)-bc>0.
\]
\end{proof}

For a Higgs--gravity pair, (15.14) shows why two individually
subcritical diagonal rows do not suffice: the feedback product
\(bc\) can still close the margin.

\subsection{The minimal Einstein--SM block matrix}

After absorbing compact gauge variables into constants, a local
unbounded packet may be organized as
\[
 {\bf V}_x=
 \bigl(V_H(H_x),V_g(g_x),V_m(m_x),V_c(c_x)\bigr),   \tag{15.15}
\]
for Higgs, gravitational/conformal, mediator, and counterterm
coordinates.  The certificate requires a finite-range nonnegative
matrix
\[
 {\mathbb B}_n=
\begin{pmatrix}
B_{HH}&B_{Hg}&B_{Hm}&B_{Hc}\\
B_{gH}&B_{gg}&B_{gm}&B_{gc}\\
B_{mH}&B_{mg}&B_{mm}&B_{mc}\\
B_{cH}&B_{cg}&B_{cm}&B_{cc}
\end{pmatrix}                                      \tag{15.16}
\]
whose entries include spatial neighbor transfers and sector
transfers.  A sufficient uniform gate is
\[
 \sup_n\rho({\mathbb B}_n)\leq1-\eta_*,
 \qquad\eta_*>0,                                   \tag{15.17}
\]
together with uniform anchor measures, local integrals (15.10),
and finite support degree.

The V13 Higgs calculation certifies the isolated \(HH\) row.  The
source manuscripts do not provide the remaining entries in
(15.16), because the RPESM symbols for gravity, gauge fixing,
mediators, and counterterms are not expanded as conditional
energies in a common field basis.

\subsection{Import into shell transport}

\begin{corollary}[Coupled sufficient gate for V14]
\label{cor:v17v15-v14-gate}
Assume (15.7)--(15.17) uniformly in \(n\), and suppose every local
shell writer is polynomially bounded by the Lyapunov functions on a
fixed-radius block.  Then its tilted third cumulants admit a uniform
bound \(C_3(\eta_*)\).  Consequently the V14 counterterm remainder
is summable whenever
\[
 \sum_nM_nC_3(\eta_*)|\varepsilon_n|^3<\infty.      \tag{15.18}
\]
\end{corollary}

\begin{proof}
Theorems~\ref{thm:v17v15-row} and
\ref{thm:v17v15-perron} give uniform block exponential moments.
Young and H\"older inequalities give the corresponding moments for
every fixed-radius polynomial writer.  The V1 tilted-cumulant
argument supplies \(C_3(\eta_*)\), and V14 gives (15.18).
\end{proof}

\begin{assessment}[What V15 closes]
\label{ass:v17v15-status}
V15 performs the required companion audit and converts the RPESM
quartic-domination sentence into a checkable coupled certificate:
conditional lower bounds, positive-measure anchors, an integrable
single-block tail, and a nonnegative influence matrix with uniform
spectral radius below one.  The exact two-sector condition exposes
the Higgs--gravity feedback obstruction.

Only the isolated Higgs row is currently certified.  No numerical
gravity, mediator, or counterterm row has been derived from the
source manuscripts, so the coupled gate is not claimed to pass.
\end{assessment}

\begin{decision}[V16: certify the compact gauge block]
\label{dec:v17v15-next}
Continue after the audit with the gauge block explicitly present in
\textup{(FS.8)} and \textup{(RP.5)}.  Prove that compact gauge links
add a zero Lyapunov row, and that the first three gauge-direction
derivatives of plaquette plus Higgs hopping terms have at most
quadratic Higgs growth.  Import V13 to obtain uniform local
exponential moments for those gauge shell writers.
\end{decision}

\begin{firewall}[Claim boundary after V15]
\label{fire:v17v15}
The block matrix is a sufficient compiler gate, not evidence that
the unspecified physical coefficient packet satisfies it.  The
companion results do not close that gap.
\end{firewall}

\section*{Version V15 append-only record}
\addcontentsline{toc}{section}{Version V15 append-only record}
The complete V14 source was retained byte for byte before this
continuation.  It had \(149046\) bytes, \(4732\) lines, and SHA--256
\[
 \texttt{3C5755994FECB9A46C0EE4A771B72F42970B6B21820B87A405C22898E5244891}.
\]
V15 performed the companion audit and proved the coupled
block-conditional energy, row-sum, and spectral-radius criteria.


\section{V16: certification of the compact gauge block}
\label{sec:v17v16}

\subsection{The source-defined gauge--Higgs block}

The actions \textup{(FS.8)} and \textup{(RP.5)} use the compact
group
\[
 K=G_{\rm SM}
 =\frac{SU(3)\times SU(2)\times U(1)}{\mathbb Z_6},                \tag{16.1}
\]
with link variables \(U_e\in K\).  The explicitly displayed local
terms are the nonnegative plaquette action and
\[
 S_{\rm hop}(U,H)
 =\kappa\sum_{e:x\to y}
 \|H_y-\rho(U_e)H_x\|^2,\qquad\kappa\geq0,          \tag{16.2}
\]
where \(\rho\) is unitary.  Couplings and plaquette incidence are
assumed uniformly bounded along the cutoff family.

\subsection{Compact blocks have a zero Lyapunov row}

\begin{lemma}[Density-independent compact moment bound]
\label{lem:v17v16-compact}
Let \(\pi(\dd u\mid\eta)\) be any probability kernel on a compact
space \(K\).  If \(F:K\to[0,\infty)\) is bounded, then
\[
 \int_K e^{\theta F(u)}\pi(\dd u\mid\eta)
 \leq e^{\theta\|F\|_\infty}                       \tag{16.3}
\]
for every \(\theta\geq0\), uniformly in the conditioning variable
\(\eta\).
\end{lemma}

\begin{proof}
The integrand is pointwise at most
\(e^{\theta\|F\|_\infty}\), and \(\pi\) has mass one.
\end{proof}

\begin{theorem}[Gauge-row spectral neutrality]
\label{thm:v17v16-zero-row}
For the gauge link block, one may take \(V_G=0\) in the V15
certificate.  Its conditional row is identically zero.  Moreover,
the V13 Higgs conditional estimate is uniform in every gauge
configuration, so the gauge column in the displayed
quartic-plus-hopping system is also zero.  Appending the gauge
block changes an unbounded-sector influence matrix \(B_{\rm unb}\)
to
\[
 \widetilde B=
 \begin{pmatrix}
 B_{\rm unb}&0\\
 0&0
 \end{pmatrix},
 \qquad
 \rho(\widetilde B)=\rho(B_{\rm unb}).              \tag{16.4}
\]
\end{theorem}

\begin{proof}
With \(V_G=0\), the left side of (15.3) equals one.  The first
statement follows with \(C_G=1\).  In the Higgs conditional
calculation, unitarity of \(\rho(U_e)\) is the only property of the
link field used in (13.5)--(13.9); all constants are independent of
\(U\).  Hence no gauge Lyapunov weight appears in the Higgs row.
The block-diagonal spectral identity in (16.4) is immediate.
\end{proof}

This argument is stronger than an energy oscillation estimate: it
does not require the conditional gauge density to be bounded above
or below.

\subsection{Gauge-direction derivatives of the hopping action}

Fix a link \(e:x\to y\).  Let \(T\) be a unit Lie-algebra direction
and write
\[
 U_e(t)=\exp(tT)U_e,\qquad
 A=\dd\rho(T),\qquad \|A\|_{\rm op}\leq L_\rho.      \tag{16.5}
\]
Set
\[
 f_e(t)=\kappa\|H_y-\rho(U_e(t))H_x\|^2.            \tag{16.6}
\]

\begin{lemma}[Exact first three hopping derivatives]
\label{lem:v17v16-hopping-derivatives}
For \(r=1,2,3\),
\[
 f_e^{(r)}(0)
 =-2\kappa\operatorname{Re}
 \langle H_y,A^r\rho(U_e)H_x\rangle,               \tag{16.7}
\]
and therefore
\[
 |f_e^{(r)}(0)|
 \leq
 \kappa L_\rho^r
 \bigl(\|H_x\|^2+\|H_y\|^2\bigr).                  \tag{16.8}
\]
\end{lemma}

\begin{proof}
Unitarity gives
\[
 f_e(t)=\kappa\bigl(
 \|H_x\|^2+\|H_y\|^2
 -2\operatorname{Re}
 \langle H_y,e^{tA}\rho(U_e)H_x\rangle\bigr).
\]
Differentiate \(r\) times.  The bound follows from
Cauchy--Schwarz, \(\|A^r\|\leq L_\rho^r\), and
\(2uv\leq u^2+v^2\).
\end{proof}

\subsection{Plaquette derivative bounds}

Let \(p_*\) be a uniform upper bound on the number of plaquettes
containing a link.  For a nonabelian factor of dimension \(d_a\),
the local plaquette term is
\[
 P_{a,p}(U)
 =\beta_a\left(
 1-\frac1{d_a}\operatorname{Re}\operatorname{tr}U_p^{(a)}
 \right),                                         \tag{16.9}
\]
and the abelian term is
\(\beta_1(1-\cos\alpha_p)\).

\begin{lemma}[Bounded plaquette jets]
\label{lem:v17v16-plaquette}
Assume the nonnegative couplings \(\beta_a\) are uniformly bounded.
For \(r=0,1,2,3\), there is \(B_{g,r}<\infty\), independent of the
cutoff, link, and surrounding gauge field, such that the \(r\)-th
unit-direction derivative of the sum of plaquettes incident on
that link has modulus at most \(B_{g,r}\).
\end{lemma}

\begin{proof}
Every relevant derivative is a continuous function on a finite
product of the compact group \(K\), so it has a finite supremum.
The same supremum works at every cell because the representation,
coupling range, and incidence bound \(p_*\) are common.  Explicitly,
normalized traces of products of unitary matrices and at most
three unit generators are bounded by the corresponding generator
operator norms.  The derivatives of cosine are bounded by one.
\end{proof}

\begin{corollary}[Quadratic growth of local gauge jets]
\label{cor:v17v16-quadratic-jets}
Let \(X_{e,r}\) be the \(r\)-th gauge-direction derivative,
\(1\leq r\leq3\), of the plaquette plus hopping action at link
\(e:x\to y\).  Then
\[
 |X_{e,r}|
 \leq B_{g,r}
 +\kappa L_\rho^r
 \bigl(\|H_x\|^2+\|H_y\|^2\bigr).                  \tag{16.10}
\]
\end{corollary}

\subsection{Uniform exponential moments of gauge shell writers}

\begin{theorem}[Gauge-jet exponential integrability]
\label{thm:v17v16-gauge-jet-moments}
Under the V13 bounded-degree and compact-parameter hypotheses, for
every \(\alpha>0\) and \(r=1,2,3\),
\[
 \sup_n\sup_e
 \mathbb E
 \exp\{\alpha|X_{e,r}|\}<\infty.                   \tag{16.11}
\]
The expectation may be taken under the joint bosonic
gauge--Higgs law.
\end{theorem}

\begin{proof}
Conditional on the entire gauge field, the Higgs law is exactly of
the V13 form, and its quartic moment bounds are uniform in that
field.  For every \(c,\epsilon>0\),
\[
 ct^2\leq\epsilon t^4+\frac{c^2}{4\epsilon}.        \tag{16.12}
\]
Apply (16.12) to the two Higgs terms in (16.10), choose
\(\epsilon\) below the two-site threshold in (13.15), and then
integrate over the gauge marginal.
\end{proof}

\begin{corollary}[Uniform tilted gauge cumulants]
\label{cor:v17v16-gauge-cumulants}
For each \(r\leq3\), the first three cumulants of \(X_{e,r}\) under
a sufficiently small exponential tilt are bounded uniformly in
\(n\) and \(e\).  Hence bounded-support gauge-link shell
interpolations satisfy the V14 third-order remainder hypothesis.
\end{corollary}

\begin{proof}
Equation (16.11) supplies a uniform two-sided exponential moment.
The V1 tilted-moment argument bounds all centered moments through
order three on a smaller tilt interval.
\end{proof}

\subsection{Eliminating one compact gauge link}

Let \(\nu_K\) be normalized Haar measure.  Condition on all fields
except \(U_e\), and let \(P_e(U_e)\) be the sum of nonnegative
plaquette terms involving the link.  Put
\[
 \Psi_e(H_x,H_y)
 =-\log\int_K
 \exp\{-P_e(U)-\kappa\|H_y-\rho(U)H_x\|^2\}
 \,\nu_K(\dd U).                                   \tag{16.13}
\]

\begin{theorem}[Compact-link elimination preserves quartic control]
\label{thm:v17v16-link-elimination}
If \(0\leq P_e\leq B_e\), then
\[
 0\leq\Psi_e(H_x,H_y)
 \leq
 B_e+2\kappa\bigl(\|H_x\|^2+\|H_y\|^2\bigr).       \tag{16.14}
\]
Thus exact integration of a compact gauge link creates at most a
quadratic Higgs debit and cannot remove the positive on-site
quartic leader.
\end{theorem}

\begin{proof}
The integrand in (16.13) lies in \((0,1]\), proving the lower
bound.  Unitarity and
\(\|u-v\|^2\leq2(\|u\|^2+\|v\|^2)\) give the pointwise lower bound
\[
 e^{-P_e(U)-\kappa\|H_y-\rho(U)H_x\|^2}
 \geq
 e^{-B_e-2\kappa(\|H_x\|^2+\|H_y\|^2)}.
\]
Integrate and take minus the logarithm.
\end{proof}

The effective interaction \(\Psi_e\) need not be polynomial, but
(16.14) is the exact growth estimate needed for quartic
integrability.

\subsection{The remaining gauge-dependent singularity}

The overlap sign, chiral determinant, and Pfaffian line are not
covered by the compactness argument.  Their gauge derivatives
contain inverse Wilson or chiral operators and require the
protected floors \textup{(FS.9)} and \textup{(FS.12)}, or a
quantitative divisor-tube theorem.  The SM--Einstein V35 source
explicitly records that support of the unbounded full-coupling
cylinder inside those protected charts has not been proved.

\begin{assessment}[What V16 closes]
\label{ass:v17v16-status}
V16 certifies the compact gauge block of the displayed bosonic
action.  It adds no Lyapunov influence row or column, leaves the
unbounded-sector spectral radius unchanged, and its first three
local action derivatives have only quadratic Higgs growth.  V13
therefore supplies their uniform exponential moments and V14
counterterm control.  Exact compact-link elimination also preserves
the quartic Higgs tail.

This does not control overlap-sign or determinant-line derivatives
near their divisors and does not supply the missing gravitational
row.
\end{assessment}

\begin{decision}[V17: classify the gravitational block]
\label{dec:v17v16-next}
Return upstream to the RPESM scheduler and gravity formulas.
Determine whether the gravitational variables integrated by the
finite cylinder live on a compact protected chart, are deterministic
functions of compact scheduler coordinates, or contain an
independent unbounded conformal mode.  Certify a zero row in the
first two cases; in the third, extract the signed leading form and
test it against V4--V8.
\end{decision}

\begin{firewall}[Claim boundary after V16]
\label{fire:v17v16}
V16 treats only the source-displayed compact plaquette and Higgs
hopping action.  It does not prove a Wilson gap, chiral floor,
fermionic phase, gravitational coercivity, or continuum limit.
\end{firewall}

\section*{Version V16 append-only record}
\addcontentsline{toc}{section}{Version V16 append-only record}
The complete V15 source was retained byte for byte before this
continuation.  It had \(159598\) bytes, \(5048\) lines, and SHA--256
\[
 \texttt{EFA65B96EC76FD65982C415ABA6A5D7824BF280564AE4850DECB5690831E4AC8}.
\]
V16 certified the compact gauge block, its local action jets, and
its quartic-preserving exact elimination.


\section{V17: classification of the scheduler-derived gravity block}
\label{sec:v17v17}

\subsection{The exact scheduler-to-ADM map}

The source scheduler uses
\[
\begin{aligned}
 {\cal B}_h
 &=h^2\sum_{a=1}^{6}(k_a^++k_a^-)r_ar_a^{\mathsf T},\\
 \beta_h
 &=h\sum_{a=1}^{6}(k_a^+-k_a^-)r_a,               \tag{17.1}
\end{aligned}
\]
and a positive speed density \(\varrho_h\).  Its derived ADM
variables are
\[
\boxed{
 g_h=\varrho_h^{2/3}(\det{\cal B}_h)^{1/3}{\cal B}_h^{-1},
 \qquad
 N_h=\varrho_h^{1/3}(\det{\cal B}_h)^{1/6},
 \qquad
 \beta_h=\beta_h.}                                  \tag{17.2}
\]
These formulas are analytic on
\[
 {\cal B}_h\succ0,\qquad\varrho_h>0.                \tag{17.3}
\]
The coframe, connection, volume, and curvature entries used by the
RPESM record are declared deterministic functions of the scheduler
coordinates.

\subsection{Deterministic vacuum branches}

\begin{theorem}[Flat scheduler is a zero block]
\label{thm:v17v17-flat}
For the RPESM reference rates
\[
 k_{\alpha,n}=(8h_n^2)^{-1}
\]
and unit speed density, the source identity is
\[
 {\cal B}_n=I_3,\qquad N_n=1,\qquad\beta_n=0.       \tag{17.4}
\]
If these scheduler values are fixed rather than integrated, every
derived gravitational writer is deterministic and contributes a
zero row and zero column to the V15 influence matrix.
\end{theorem}

\begin{proof}
Equation (17.4) is the source formula \textup{(RP.4)}.  A
deterministic coordinate has conditional law a Dirac mass.  Its
bounded exponential moment is its fixed value, and no boundary
Lyapunov function is needed.  It cannot transfer randomness to
another block.
\end{proof}

\begin{theorem}[de Sitter scheduler on a compact time slab]
\label{thm:v17v17-desitter}
On the source's homogeneous branch
\[
 {\cal B}_h(t)=e^{-2Ht}I_3,\qquad
 \varrho_h(t)=e^{3Ht},\qquad
 g_h(t)=e^{2Ht}I_3,\qquad N_h(t)=1,\quad\beta_h=0. \tag{17.5}
\]
For fixed \(H\) and \(|t|\leq T\), every displayed variable in
(17.5) lies in a compact interval independent of \(h\).  If the
branch is fixed as the background regulator, it also contributes a
zero stochastic gravity row.
\end{theorem}

\begin{proof}
Substitution in (17.2) gives (17.5).  On a compact time slab,
\(e^{\pm2Ht}\) and \(e^{3Ht}\) have finite positive upper and lower
bounds.  The last statement follows as in
Theorem~\ref{thm:v17v17-flat}.
\end{proof}

These two theorems certify the source's vacuum-background
constructions.  They do not describe an integrated fluctuating
Einstein field.

\subsection{A compact scheduler window}

For positive constants
\(b_-,b_+,\rho_-,\rho_+,b_\beta\), define
\[
 {\cal K}_{\rm ADM}
 =\left\{({\cal B},\beta,\varrho):
 b_-I\preceq{\cal B}\preceq b_+I,\ 
 |\beta|\leq b_\beta,\ 
 \rho_-\leq\varrho\leq\rho_+\right\}.              \tag{17.6}
\]

\begin{theorem}[Protected compact gravity chart]
\label{thm:v17v17-compact}
The image of \({\cal K}_{\rm ADM}\) under (17.2) is compact.
The metric and lapse obey explicit bounds
\[
\begin{aligned}
 \rho_-^{2/3}\frac{b_-}{b_+}\,I
 &\preceq g
 \preceq
 \rho_+^{2/3}\frac{b_+}{b_-}\,I,\\
 \rho_-^{1/3}b_-^{1/2}
 &\leq N\leq
 \rho_+^{1/3}b_+^{1/2}.                            \tag{17.7}
\end{aligned}
\]
Every continuous finite-depth gravitational writer, and every
derivative through fixed order of an analytic writer, is uniformly
bounded on a slightly larger such chart.  Hence an RPESM law
supported there has a zero gravitational Lyapunov row.
\end{theorem}

\begin{proof}
The domain (17.6) is compact and lies a positive distance from the
singular boundary of (17.3).  Continuity of (17.2) gives
compactness of its image.  For eigenvalues
\(\lambda_i({\cal B})\in[b_-,b_+]\), the metric eigenvalues are
\[
 \varrho^{2/3}
 \frac{(\lambda_1\lambda_2\lambda_3)^{1/3}}
 {\lambda_i}.
\]
The geometric mean lies in \([b_-,b_+]\), which proves the first
line of (17.7).  Since
\((\det{\cal B})^{1/6}\in[b_-^{1/2},b_+^{1/2}]\), the second line
follows.  Analytic derivatives are continuous on the closure of a
slightly enlarged compact chart.  Lemma~\ref{lem:v17v16-compact}
then gives the zero-row assertion.
\end{proof}

\begin{remark}[Curvature needs a finite-depth hypothesis]
\label{rem:v17v17-curvature}
Pointwise compactness of \(({\cal B},\beta,\varrho)\) does not alone
bound discrete curvature, which also contains neighboring
differences.  The theorem covers curvature writers only when the
finite neighborhood tuple lies in a common compact chart and the
mesh-scaled difference coefficients are uniformly bounded.
\end{remark}

\subsection{The open positive chart is noncompact}

\begin{theorem}[Exact speed-scaling escape]
\label{thm:v17v17-escape}
The full positive scheduler domain (17.3) is noncompact, and the
derived ADM map is not bounded there.  Indeed, fix
\({\cal B}=I_3\), \(\beta=0\), and set
\[
 \varrho=e^{3s}.
\]
Then
\[
 g=e^{2s}I_3,\qquad N=e^s.                         \tag{17.8}
\]
As \(s\to+\infty\), both variables diverge; as
\(s\to-\infty\), the metric degenerates.
\end{theorem}

\begin{proof}
Equation (17.8) is immediate from (17.2).  It exhibits an escaping
sequence in each direction.
\end{proof}

\begin{corollary}[An open tangent chart is not a Gibbs certificate]
\label{cor:v17v17-open-not-certificate}
The source statement that an open positive rate/speed chart
supplies the ten-dimensional ADM tangent does not imply
normalizability or any V15 block moment estimate for fluctuating
gravity.
\end{corollary}

\begin{proof}
The tangent statement is local differential nondegeneracy.
The escape (17.8) shows that the same domain contains unbounded
conformal/lapse directions.  A probability and moment assertion
requires a base measure and conditional energy controlling those
directions; neither follows from local rank.
\end{proof}

\subsection{Audit of the full RPESM cylinder}

The RPESM definition simultaneously states that:
\begin{enumerate}
\item scheduler and derived geometric variables belong to the
      bosonic record;
\item the action contains a symbol \(S_{g,n}\);
\item the cylinder is measured against \(\dd\lambda_n\);
\item an open positive rate/speed chart supplies the ADM tangent.
\end{enumerate}
It does not define the scheduler component of \(\lambda_n\), the
domain actually integrated, or a formula for \(S_{g,n}\).

\begin{theorem}[Source-level gravity trichotomy]
\label{thm:v17v17-trichotomy}
From the written source, exactly the following conditional
classification is justified:
\begin{enumerate}
\item on the fixed flat or de Sitter background branches, gravity
      is deterministic and has a zero influence row;
\item on a law explicitly supported in a uniform chart such as
      (17.6), gravity is compact and has a zero influence row;
\item on the unrestricted positive scheduler cone, no influence
      row can be computed until its reference measure and
      conditional gravitational action are supplied.
\end{enumerate}
\end{theorem}

\begin{proof}
The first item is
Theorems~\ref{thm:v17v17-flat}--\ref{thm:v17v17-desitter}.
The second is Theorem~\ref{thm:v17v17-compact}.  The third follows
from Theorem~\ref{thm:v17v17-escape} and the missing source data
listed above.  These cases exhaust fixed, compactly supported, and
unrestricted interpretations of the stated scheduler record.
\end{proof}

\subsection{The missing gravitational conditional row}

For an integrated gravitational block, the next required datum is
not merely a signed Hessian.  It is a conditional energy
\[
 E_{g,n}(z\mid H,U,m,c,\text{neighboring }z)        \tag{17.9}
\]
and a proper Lyapunov function \(V_g(z)\) for which one can prove
\[
\begin{aligned}
 E_{g,n}(z\mid\eta)
 &\geq\ell_gV_g(z)-K_g-\sum_jd_{gj}V_j(\eta_j),\\
 E_{g,n}(z\mid\eta)
 &\leq K_g+\sum_je_{gj}V_j(\eta_j)
 \quad\text{on a positive-measure anchor},         \tag{17.10}\\
 \int e^{-(\ell_g-\theta_g)V_g(z)}\,\dd\lambda_{g,n}(z)
 &\leq J_g.
\end{aligned}
\]
Uniform values in (17.10) would determine the \(g\)-row of (15.16).
For the scalar escape \(s\) in (17.8), the source presently gives
neither \(E_{g,n}(s\mid\eta)\) nor
\(\dd\lambda_{g,n}(s)\).

\begin{assessment}[What V17 closes]
\label{ass:v17v17-status}
V17 proves that the source's explicit flat and de Sitter
regulators have deterministic gravitational writers, and that a
uniform protected ADM window would be a compact zero-row block.
It also gives an explicit speed-scaling escape proving that the
open scheduler chart is not itself an integrability certificate.

The full RPESM cylinder does not specify which of these
interpretations its measure uses.  The first unresolved
gravitational object is the conditional packet (17.9)--(17.10), not
an angular-zero calculation.
\end{assessment}

\begin{decision}[V18: signed conformal conditional model]
\label{dec:v17v17-next}
Use the exact escape coordinate \(s=\frac13\log\varrho\) to test
the minimal possible gravitational conditional row.  Prove
necessary and sufficient tail conditions for a one-dimensional
signed conformal energy coupled to the Higgs quartic.  This will
show precisely which growth and sign data an explicit \(S_{g,n}\)
must supply before the Perron cycle test (15.14) can be evaluated.
\end{decision}

\begin{firewall}[Claim boundary after V17]
\label{fire:v17v17}
The deterministic and compact-chart zero rows do not establish a
fluctuating quantum-gravity measure.  The open-chart escape proves
missing control, not divergence of an action that has not been
specified.
\end{firewall}

\section*{Version V17 append-only record}
\addcontentsline{toc}{section}{Version V17 append-only record}
The complete V16 source was retained byte for byte before this
continuation.  It had \(169485\) bytes, \(5346\) lines, and SHA--256
\[
 \texttt{FEFC3D4371968EE812125FF865E9CE6E0745425ACDE47500FB54A439BF6F48A4}.
\]
V17 classified the deterministic, compact, and unrestricted
scheduler-derived gravity branches and isolated the missing
conditional gravity row.


\section{V18: the minimal signed conformal--Higgs conditional model}
\label{sec:v17v18}

\subsection{The conformal-volume coordinate}

Along the V17 speed-scaling escape,
\[
 \varrho=e^{3s},\qquad g=e^{2s}I_3,\qquad N=e^s.    \tag{18.1}
\]
Hence the four-dimensional cell-volume factor is
\[
 N\sqrt{\det g}=e^{4s}.                            \tag{18.2}
\]
If a local \(q\)-component Higgs variable is integrated against
Lebesgue measure and only its radial quartic is retained, the
minimal conditional model is
\[
 E(s,h)=U(s)+\lambda e^{\alpha s}\|h\|^4,
 \qquad
 s\in\mathbb R,\quad h\in\mathbb R^q,              \tag{18.3}
\]
where \(\lambda,\alpha>0\).  The geometric specialization of
(18.2) is \(\alpha=4\).  The function \(U\) is the as-yet
unspecified signed gravitational/conformal energy.

\subsection{Exact integration of the Higgs block}

\begin{theorem}[Exact conformal partition criterion]
\label{thm:v17v18-partition}
Assume \(U:\mathbb R\to\mathbb R\) is measurable and finite almost
everywhere.  Then
\[
\begin{aligned}
 Z
 &:=\int_{\mathbb R}\int_{\mathbb R^q}
 e^{-U(s)-\lambda e^{\alpha s}\|h\|^4}\,\dd h\,\dd s\\
 &=C_{q,\lambda}
 \int_{\mathbb R}
 e^{-U(s)-\alpha q s/4}\,\dd s,                   \tag{18.4}
\end{aligned}
\]
where
\[
 C_{q,\lambda}
 =\frac{|\mathbb S^{q-1}|}{4}
 \lambda^{-q/4}\Gamma(q/4).                       \tag{18.5}
\]
Consequently \(0<Z<\infty\) if and only if
\[
 e^{-U(s)-\alpha q s/4}\in L^1(\mathbb R,\dd s).  \tag{18.6}
\]
\end{theorem}

\begin{proof}
For fixed \(s\), substitute
\[
 y=e^{\alpha s/4}h.
\]
The Jacobian is
\(\dd h=e^{-\alpha q s/4}\dd y\).  Polar integration of
\(e^{-\lambda\|y\|^4}\) gives (18.5), and Tonelli gives (18.4).
The constant is finite and strictly positive, proving the
equivalence.
\end{proof}

The Higgs integral therefore shifts the effective conformal energy
by the exact linear term
\[
 U_{\rm eff}(s)=U(s)+\frac{\alpha q}{4}s.           \tag{18.7}
\]
For one complex Higgs doublet, \(q=4\); with the geometric value
\(\alpha=4\), the shift is \(4s\).

\subsection{Exact polynomial moment criterion}

\begin{theorem}[Physical Higgs moments]
\label{thm:v17v18-polynomial-moments}
For every \(p\geq0\),
\[
\begin{aligned}
 &\int_{\mathbb R}\int_{\mathbb R^q}
 \|h\|^p
 e^{-U(s)-\lambda e^{\alpha s}\|h\|^4}\,\dd h\,\dd s\\
 &\qquad =
 C_{q,p,\lambda}
 \int_{\mathbb R}
 e^{-U(s)-\alpha(q+p)s/4}\,\dd s,                 \tag{18.8}
\end{aligned}
\]
where
\[
 C_{q,p,\lambda}
 =\frac{|\mathbb S^{q-1}|}{4}
 \lambda^{-(q+p)/4}\Gamma((q+p)/4).               \tag{18.9}
\]
Thus the normalized law has a finite \(p\)-th physical Higgs
moment exactly when the integral on the right of (18.8) is finite.
\end{theorem}

\begin{proof}
The same substitution gives
\(\|h\|^p=e^{-\alpha ps/4}\|y\|^p\).  Polar integration yields
(18.9).
\end{proof}

Each higher physical Higgs moment therefore demands a stronger
negative-\(s\) conformal tail.

\subsection{Failure of an unweighted quartic exponential moment}

\begin{theorem}[Conformal-collapse obstruction]
\label{thm:v17v18-collapse}
Suppose the conformal support contains sets of positive Lebesgue
measure arbitrarily far toward \(s=-\infty\), and \(U\) is finite
almost everywhere there.  Then, for every \(\theta>0\),
\[
 \int e^{\theta\|h\|^4}
 e^{-U(s)-\lambda e^{\alpha s}\|h\|^4}
 \,\dd h\,\dd s=\infty.                            \tag{18.10}
\]
Thus no volume-uniform V13 certificate with the unweighted
Lyapunov function \(\|h\|^4\) can hold on the unrestricted
conformal cone.
\end{theorem}

\begin{proof}
Choose \(s_0\) so negative that
\(\lambda e^{\alpha s}<\theta\) for every \(s<s_0\).
For each such \(s\) at which \(U(s)\) is finite, the inner integral
in (18.10) contains
\[
 \int_{\mathbb R^q}
 e^{(\theta-\lambda e^{\alpha s})\|h\|^4}\,\dd h
 =\infty.
\]
Tonelli and the positive-measure support hypothesis prove (18.10).
\end{proof}

This is an obstruction to the choice of Lyapunov coordinate, not
to all conformal--Higgs measures.

\subsection{The correct rescaled Higgs coordinate}

Put
\[
 y=e^{\alpha s/4}h.                                \tag{18.11}
\]

\begin{theorem}[Exact product factorization]
\label{thm:v17v18-factorization}
Whenever (18.6) holds, the normalized \((s,y)\) law is the product
of
\[
 Z_s^{-1}e^{-U(s)-\alpha q s/4}\,\dd s
 \quad\text{and}\quad
 C_{q,\lambda}^{-1}e^{-\lambda\|y\|^4}\,\dd y.     \tag{18.12}
\]
In particular,
\[
 \mathbb E e^{\theta\|y\|^4}<\infty
 \quad\Longleftrightarrow\quad
 0\leq\theta<\lambda.                              \tag{18.13}
\]
\end{theorem}

\begin{proof}
The change of variables used in Theorem
\ref{thm:v17v18-partition} transforms the joint density into
\[
 e^{-U(s)-\alpha q s/4}e^{-\lambda\|y\|^4}\,\dd s\,\dd y.
\]
This factorizes.  Equation (18.13) follows by replacing
\(\lambda\) with \(\lambda-\theta\) in the radial integral.
\end{proof}

Thus the natural coupled Lyapunov function is
\[
 V_H(s,h)=e^{\alpha s}\|h\|^4=\|y\|^4,             \tag{18.14}
\]
not \(\|h\|^4\).  In the minimal model it removes the
Higgs-to-conformal influence row exactly.  Reconstruction of the
physical field \(h=e^{-\alpha s/4}y\) still requires the stronger
tail tests (18.8).

\subsection{Signed quadratic conformal energies}

\begin{theorem}[Complete quadratic classification]
\label{thm:v17v18-quadratic}
Let
\[
 U(s)=\kappa_gs^2+b_gs+c_g.                        \tag{18.15}
\]
The joint partition (18.4) is finite if and only if
\[
 \kappa_g>0.                                       \tag{18.16}
\]
When (18.16) holds, every polynomial physical Higgs moment is
finite, but the unweighted exponential moment in (18.10) remains
infinite.
\end{theorem}

\begin{proof}
The conformal integral in (18.4) is
\[
 e^{-c_g}\int_{\mathbb R}
 \exp\left\{-\kappa_gs^2-
 \left(b_g+\frac{\alpha q}{4}\right)s\right\}\dd s. \tag{18.17}
\]
It is Gaussian-integrable when \(\kappa_g>0\).  If
\(\kappa_g<0\), it diverges in both tails.  If \(\kappa_g=0\), an
affine exponential cannot be integrable at both \(+\infty\) and
\(-\infty\).  Replacing \(q\) by \(q+p\) proves finiteness of every
integral in (18.8) when \(\kappa_g>0\).  The last assertion is
Theorem~\ref{thm:v17v18-collapse}.
\end{proof}

\begin{corollary}[Polynomial signed conformal energy]
\label{cor:v17v18-polynomial}
If \(U\) is a nonconstant real polynomial of degree at least two,
then (18.6) holds exactly when its degree is even and its leading
coefficient is positive.
\end{corollary}

\begin{proof}
An even positive leading term tends to \(+\infty\) faster than the
linear shift in both tails.  An even negative leader tends to
\(-\infty\) in both tails, while an odd-degree leader tends to
\(-\infty\) in one tail.  The corresponding exponential is
respectively integrable or divergent.
\end{proof}

\subsection{A checkable gravitational row}

For the minimal model, the gravitational marginal is exactly
\[
 \dd\mu_g(s)
 =Z_s^{-1}
 e^{-U_{\rm eff}(s)}\,\dd s.                       \tag{18.18}
\]
If there is a proper \(V_g\) and numbers
\(\theta_g>0,C_g<\infty\) with
\[
 U_{\rm eff}(s)\geq
 \theta_gV_g(s)-C_g,                               \tag{18.19}
\]
and
\[
 \int e^{-\epsilon V_g(s)}\,\dd s<\infty
 \quad\text{for every }\epsilon>0,                 \tag{18.20}
\]
then a smaller exponential parameter gives a zero-boundary
gravitational moment row by V15.  For the quadratic case
\(\kappa_g>0\), one may take \(V_g=s^2\).

The source provides no value or sign for \(\kappa_g\), and it does
not state that \(S_{g,n}\) reduces to (18.15).  The theorem therefore
identifies a coefficient test; it does not certify the physical
gravity row.

\begin{assessment}[What V18 closes]
\label{ass:v17v18-status}
V18 solves the minimal conformal--Higgs integrability problem
exactly.  Higgs integration adds the Jacobian drift
\(\alpha q s/4\); a quadratic conformal energy normalizes the joint
law exactly when its quadratic coefficient is positive.  The
unweighted Higgs quartic exponential moment necessarily fails on
an unrestricted collapsing conformal tail, while the rescaled
field \(y=e^{\alpha s/4}h\) restores an exact product quartic law.

The RPESM manuscript does not specify \(U\), its sign, or the
scheduler reference measure.  Hence the physical gravitational
row remains open at the now-explicit test (18.6).
\end{assessment}

\begin{decision}[V19: lower-order conformal couplings]
\label{dec:v17v18-next}
Add the signed Higgs mass and nearest-neighbor hopping terms with
their conformal weights.  After the rescaling (18.11), determine
which exponential powers of \(s\) are induced, and give sharp
Young-absorption conditions under which a positive \(U(s)\)
continues to normalize the coupled block.
\end{decision}

\begin{firewall}[Claim boundary after V18]
\label{fire:v17v18}
The model (18.3) is the minimal conformal-volume test inferred from
the scheduler geometry.  It is not asserted to equal the
unspecified RPESM gravitational action, and its normalizability does
not establish reflection positivity or a continuum limit.
\end{firewall}

\section*{Version V18 append-only record}
\addcontentsline{toc}{section}{Version V18 append-only record}
The complete V17 source was retained byte for byte before this
continuation.  It had \(179394\) bytes, \(5632\) lines, and SHA--256
\[
 \texttt{98BC4C7A0ABFD6AF56CEF1B312C4DADB73D388B38089A45E5ECB360CCCACD693}.
\]
V18 proved the exact conformal--Higgs partition, moment, collapse,
rescaling, and signed-quadratic criteria.


\section{V19: lower-order conformal couplings and the sharp tail threshold}
\label{sec:v17v19}

\subsection{A general exponential Young law}

Let \(m>k>0\), \(A,B>0\), and
\[
 F_s(r)=Ae^{\alpha s}r^m-Be^{\beta s}r^k,
 \qquad r\geq0.                                    \tag{19.1}
\]

\begin{theorem}[Exact induced conformal exponent]
\label{thm:v17v19-young}
The minimum of (19.1) is
\[
 \min_{r\geq0}F_s(r)
 =-C_{m,k}(A,B)e^{\sigma s},                       \tag{19.2}
\]
where
\[
\begin{aligned}
 \sigma&=\frac{m\beta-k\alpha}{m-k},\\
 C_{m,k}(A,B)
 &=\frac{m-k}{m}B
 \left(\frac{kB}{mA}\right)^{k/(m-k)}.             \tag{19.3}
\end{aligned}
\]
Thus a negative degree-\(k\) term induces the exponential
conformal power \(e^{\sigma s}\).
\end{theorem}

\begin{proof}
The nonzero critical point satisfies
\[
 r^{m-k}=\frac{kB}{mA}e^{(\beta-\alpha)s}.
\]
At that point
\(Ae^{\alpha s}r^m=(k/m)Be^{\beta s}r^k\).
Substitution gives (19.2)--(19.3).  The value is negative and is
below the value zero at \(r=0\), so it is the global minimum.
\end{proof}

\begin{corollary}[Conformal Newton trichotomy]
\label{cor:v17v19-trichotomy}
As \(s\to+\infty\), the induced negative term:
\begin{enumerate}
\item decays if \(m\beta<k\alpha\);
\item remains of constant order if \(m\beta=k\alpha\);
\item grows like \(e^{\sigma s}\) if \(m\beta>k\alpha\).
\end{enumerate}
In the last case the gravitational energy must dominate that
exponential power.
\end{corollary}

\subsection{The four-dimensional signed mass}

For the geometric value \(\alpha=4\), consider
\[
 E(s,h)
 =U_0(s)+e^{4s}
 \left(\lambda\|h\|^4+a\|h\|^2+c\right),
 \qquad\lambda>0.                                  \tag{19.4}
\]
Here \(a\) is the signed Higgs mass coefficient and \(ce^{4s}\) is
the leading conformal-volume term supplied by the gravitational or
cosmological sector.  Put
\[
 a_-=\max\{0,-a\},\qquad
 \Delta_c=c-\frac{a_-^2}{4\lambda}.                \tag{19.5}
\]

\begin{lemma}[Exact radial minimum]
\label{lem:v17v19-radial-minimum}
For
\[
 W(h)=\lambda\|h\|^4+a\|h\|^2+c,
\]
one has
\[
 \min_{h\in\mathbb R^q}W(h)=\Delta_c.              \tag{19.6}
\]
If \(a<0\), the minimum occurs on the sphere
\(\|h\|^2=a_-/(2\lambda)\); if \(a\geq0\), it occurs at \(h=0\).
\end{lemma}

\begin{proof}
Set \(u=\|h\|^2\geq0\).  The quadratic
\(\lambda u^2+au+c\) has unconstrained minimizer
\(-a/(2\lambda)\).  This lies in \([0,\infty)\) exactly when
\(a<0\).  Substitution gives (19.5)--(19.6).
\end{proof}

\begin{theorem}[Sharp positive-conformal coefficient]
\label{thm:v17v19-threshold}
Assume
\[
 U_0(s)=o(e^{4s})\qquad(s\to+\infty).              \tag{19.7}
\]
Then:
\begin{enumerate}
\item if \(\Delta_c<0\), the joint partition associated with
      (19.4) diverges at \(+\infty\);
\item if \(\Delta_c>0\), its positive-\(s\) tail is finite;
\item if \(\Delta_c=0\), the positive-\(s\) tail is governed by the
      zero-set alternatives in
      Theorem~\ref{thm:v17v19-critical-tail} below.
\end{enumerate}
Hence the sharp nonnegative-action threshold is
\[
 c\geq\frac{a_-^2}{4\lambda}.                      \tag{19.8}
\]
\end{theorem}

\begin{proof}
If \(\Delta_c<0\), choose a small ball or annulus around a minimizer
of \(W\) on which \(W\leq\Delta_c/2<0\).  Its contribution to the
Higgs integral is at least a positive constant times
\[
 \exp\{|\Delta_c|e^{4s}/2-U_0(s)\},
\]
whose \(s\)-integral diverges by (19.7).

If \(\Delta_c>0\), coercivity and continuity of \(W\) give numbers
\(b>0\) and \(C<\infty\) such that
\[
 W(h)\geq\Delta_c/2+b\|h\|^4.                     \tag{19.9}
\]
Indeed choose a large ball on which the quartic gives the second
term, and take the positive minimum of
\((W-\Delta_c/2)/(1+\|h\|^4)\) on the remaining compact set.
The Higgs integral is then at most
\[
 C e^{-qs}
 \exp\{-\Delta_c e^{4s}/2-U_0(s)\},
\]
which is integrable at \(+\infty\) under (19.7).
The equality case has no positive gap and is treated next.
\end{proof}

\subsection{The critical coefficient}

Let
\[
 I(s)=\int_{\mathbb R^q}e^{-e^{4s}W(h)}\,\dd h.    \tag{19.10}
\]

\begin{theorem}[Critical positive-tail asymptotics]
\label{thm:v17v19-critical-tail}
Assume \(\Delta_c=0\).  As \(s\to+\infty\), with two-sided
comparability constants independent of \(s\),
\[
 I(s)\asymp
 \begin{cases}
 e^{-2s},&a<0,\\
 e^{-qs},&a=0,\\
 e^{-2qs},&a>0.
 \end{cases}                                       \tag{19.11}
\]
Consequently the positive-\(s\) part of the joint partition is
finite exactly when
\[
\begin{cases}
\displaystyle\int^\infty e^{-U_0(s)-2s}\,\dd s<\infty,&a<0,\\
\displaystyle\int^\infty e^{-U_0(s)-qs}\,\dd s<\infty,&a=0,\\
\displaystyle\int^\infty e^{-U_0(s)-2qs}\,\dd s<\infty,&a>0.
\end{cases}                                       \tag{19.12}
\]
\end{theorem}

\begin{proof}
If \(a<0\), then
\[
 W(h)=\lambda(\|h\|^2-r_0^2)^2,
 \qquad r_0^2=a_-/(2\lambda).
\]
In polar coordinates, the only shrinking direction is radial.
On a fixed annulus around \(r_0\),
\((r^2-r_0^2)^2\) is comparable to \((r-r_0)^2\), so its integral is
comparable to \(e^{-2s}\).  Outside that annulus \(W\) has a positive
gap, giving an exponentially smaller contribution.

If \(a=0\), then \(c=0\) and the substitution \(y=e^sh\) gives
exactly \(I(s)=C_{q,\lambda}e^{-qs}\).

If \(a>0\), again \(c=0\).  Substitute \(y=e^{2s}h\):
\[
 I(s)=e^{-2qs}\int_{\mathbb R^q}
 e^{-a\|y\|^2-\lambda e^{-4s}\|y\|^4}\,\dd y.
\]
For large \(s\), the remaining integral is bounded above by the
Gaussian integral and below by its integral over the unit ball.
This proves (19.11), and multiplying by \(e^{-U_0(s)}\) proves
(19.12).
\end{proof}

\subsection{The collapsing tail}

\begin{theorem}[Universal negative-tail criterion]
\label{thm:v17v19-negative-tail}
For every fixed \(a,c\), as \(s\to-\infty\),
\[
 I(s)\asymp e^{-qs}.                               \tag{19.13}
\]
Therefore the negative-\(s\) part of the joint partition is finite
exactly when
\[
 \int_{-\infty}
 e^{-U_0(s)-qs}\,\dd s<\infty.                    \tag{19.14}
\]
\end{theorem}

\begin{proof}
Use \(y=e^sh\) in (19.10):
\[
 I(s)=e^{-qs}e^{-ce^{4s}}
 \int_{\mathbb R^q}
 e^{-\lambda\|y\|^4-ae^{2s}\|y\|^2}\,\dd y.
\]
For sufficiently negative \(s\), the last integral is bounded
above and below by positive constants.  An upper bound follows by
absorbing \(|a|e^{2s}\|y\|^2\) into half the quartic; a lower bound
follows by integrating over the unit ball.  Also
\(e^{-ce^{4s}}\to1\).
\end{proof}

\subsection{Mass absorption in the rescaled field}

With \(y=e^sh\), (19.4) becomes
\[
 E(s,y)
 =U_0(s)+qs+\lambda\|y\|^4
 +ae^{2s}\|y\|^2+ce^{4s},                          \tag{19.15}
\]
where \(qs\) is the Jacobian contribution to the effective action.
For \(a<0\), exact Young minimization gives
\[
 \lambda\|y\|^4-a_-e^{2s}\|y\|^2
 \geq-\frac{a_-^2}{4\lambda}e^{4s}.                \tag{19.16}
\]
This is the specialization of
Theorem~\ref{thm:v17v19-young} with
\[
 m=4,\quad k=2,\quad\alpha=0,\quad\beta=2,
 \quad\sigma=4.                                    \tag{19.17}
\]
Equation (19.16) explains the same threshold (19.8) directly in
the rescaled coordinates.

\subsection{The conformally weighted hopping term}

On one common conformal cell, the scalar kinetic scaling is
\[
 e^{2s}\kappa\|h_y-\rho(U)h_x\|^2.                 \tag{19.18}
\]
If
\[
 y_x=e^sh_x,\qquad y_y=e^sh_y,
\]
then
\[
 e^{2s}\kappa\|h_y-\rho(U)h_x\|^2
 =\kappa\|y_y-\rho(U)y_x\|^2.                      \tag{19.19}
\]
Thus the common-cell hopping term induces no conformal power and
falls exactly within the V13 quartic-plus-hopping class in the
rescaled field.

If neighboring cells have different conformal factors, the source
must specify the edge weight and parallel transport between their
rescaled frames.  Formula (19.19) alone does not determine that
two-cell conditional row.

\subsection{The coupled gate produced by V19}

For the minimal local block, the required data are now explicit:
\begin{enumerate}
\item the negative-tail confinement (19.14);
\item the positive coefficient margin
      \[
      \Delta_c
      =c-a_-^2/(4\lambda)\geq0;                    \tag{19.20}
      \]
\item at equality, the appropriate critical integral in (19.12);
\item the common-cell rescaled hopping identity (19.19);
\item for variable conformal factors, an edge formula sufficient
      to compute the off-diagonal V15 influence coefficients.
\end{enumerate}

\begin{assessment}[What V19 closes]
\label{ass:v17v19-status}
V19 derives the exact conformal exponent induced by every signed
lower-degree radial term.  For the four-dimensional Higgs mass, the
sharp leading coefficient is
\(c-a_-^2/(4\lambda)\).  The positive, negative, and critical cases
are classified, including exact critical power laws and the
universal collapsing-tail criterion.  The correctly weighted
common-cell hopping term becomes the standard V13 hopping action
after rescaling.

The RPESM source gives none of \(U_0\), \(c\), or the variable-cell
conformal edge weight.  Therefore (19.20) is a precise physical
coefficient gate, not a certified physical margin.
\end{assessment}

\begin{decision}[V20: companion audit and edge conformal mismatch]
\label{dec:v17v19-next}
Perform the mandatory YM/NS audit.  Then treat two neighboring
conformal factors \(s_x,s_y\).  Derive the exact rescaled hopping
form for symmetric endpoint and geometric-mean edge weights, and
identify which choice yields a block influence cycle satisfying
the V15 spectral-radius test.  Continue in V21 after the audit.
\end{decision}

\begin{firewall}[Claim boundary after V19]
\label{fire:v17v19}
V19 solves a local conformal--Higgs model.  It does not identify the
model with the unspecified RPESM gravitational action, prove a
positive physical cosmological margin, or construct the continuum
measure.
\end{firewall}

\section*{Version V19 append-only record}
\addcontentsline{toc}{section}{Version V19 append-only record}
The complete V18 source was retained byte for byte before this
continuation.  It had \(188844\) bytes, \(5934\) lines, and SHA--256
\[
 \texttt{BCC4D87304579BB978F584F6263B5BB2EF1C6F3C862C408D83B885101E7623D4}.
\]
V19 proved the induced conformal exponent, sharp signed-mass
threshold, critical tail laws, and rescaled hopping identity.


\section{V20: companion audit and two-cell conformal hopping}
\label{sec:v17v20}

\subsection{Mandatory companion audit}

The checkpoint files are
\[
\begin{array}{lll}
\text{YM V74:}&16824\text{ lines},&
\texttt{33A7CCA4A8487371B5A4753736441CFECEEF5F34887F2C5EDE0CD5ABD41F760F},\\
\text{NS V26:}&5319\text{ lines},&
\texttt{82C887F8C2C69E4F28FAD7C3DC8DE5B9099CB5A4BF431EBA1FFF3E91935978AD}.
\end{array}                                                       \tag{20.1}
\]
YM V74 repairs the four terminal V73 columns with eight rational
bridge charts and proves the complete fifth direct-plus quarter
cell.  Momentum inversion supplies direct minus on the opposite
sixth cell.  Direct minus on the fifth and direct plus on the sixth
remain open, as do full-torus, cofinal, measure, and mass-gap
statements.  NS V26 is unchanged.

\begin{theorem}[V20 audit decision]
\label{thm:v17v20-audit}
The new YM cell theorem does not provide the RPESM overlap Wilson
floor, a conditional conformal edge action, or a coupled moment
margin.  Neither companion result changes the V19 coefficient
gate.
\end{theorem}

\begin{proof}
YM V74 is a finite Loewner estimate for one explicitly constructed
Wilson response pencil and two response-sign/orientation pairings
remain absent.  It is not a cutoff-uniform lower singular-value
bound for the chiral overlap coefficient in the RPESM cylinder.
The NS tensor norms concern neither object.
\end{proof}

\subsection{General two-endpoint rescaling}

Let \(s_x,s_y\in\mathbb R\), put
\[
 d=s_y-s_x,\qquad
 y_x=e^{s_x}h_x,\qquad y_y=e^{s_y}h_y,             \tag{20.2}
\]
and let \(U\) be unitary.  For a positive symmetric edge weight
\(w(s_x,s_y)\), the physical hopping term becomes
\[
\begin{aligned}
 &\kappa w(s_x,s_y)\|h_y-Uh_x\|^2\\
 &\quad=
 \kappa w(s_x,s_y)
 \|e^{-s_y}y_y-e^{-s_x}Uy_x\|^2.                  \tag{20.3}
\end{aligned}
\]
This identity shows that conformal cancellation on a constant
cell does not determine the variable-cell edge.

\subsection{Geometric-mean edge weight}

Take
\[
 w_G(s_x,s_y)=e^{s_x+s_y}.                         \tag{20.4}
\]

\begin{theorem}[Geometric edge normal form]
\label{thm:v17v20-geometric}
The rescaled hopping action is
\[
\boxed{
 E_G
 =\kappa
 \|e^{-d/2}y_y-e^{d/2}Uy_x\|^2.}                  \tag{20.5}
\]
Equivalently,
\[
 E_G=\kappa\left(
 e^{-d}\|y_y\|^2+e^d\|y_x\|^2
 -2\operatorname{Re}\langle y_y,Uy_x\rangle
 \right).                                         \tag{20.6}
\]
\end{theorem}

\begin{proof}
Insert (20.4) in (20.3) and collect the two diagonal coefficients
and the cross coefficient.
\end{proof}

\subsection{Arithmetic endpoint edge weight}

Take
\[
 w_A(s_x,s_y)
 =\frac12(e^{2s_x}+e^{2s_y})
 =e^{s_x+s_y}\cosh d.                              \tag{20.7}
\]

\begin{theorem}[Arithmetic edge normal form]
\label{thm:v17v20-arithmetic}
The rescaled hopping action is
\[
\boxed{
 E_A
 =\kappa\cosh d\,
 \|e^{-d/2}y_y-e^{d/2}Uy_x\|^2
 =\cosh d\,E_G.}                                   \tag{20.8}
\]
\end{theorem}

\begin{proof}
Equation (20.7) differs from (20.4) by the factor \(\cosh d\).
\end{proof}

Both actions are nonnegative, symmetric under endpoint exchange,
and reduce to the V19 common-cell hopping when \(d=0\).  Their
large mismatch costs are different.

\subsection{Anchor costs and induced mismatch powers}

When conditioning \(y_x\), use the unit ball
\(\|y_x\|\leq1\) as normalization anchor.  From (20.6),
\[
 E_G
 \leq\kappa\left(
 e^d+e^{-d}\|y_y\|^2+2\|y_y\|
 \right)                                           \tag{20.9}
\]
on that anchor.  For every \(\delta>0\),
\[
 \kappa e^{-d}\|y_y\|^2
 \leq
 \delta\|y_y\|^4+
 \frac{\kappa^2}{4\delta}e^{-2d}.                 \tag{20.10}
\]
The linear cross term is absorbed into
\(\delta\|y_y\|^4+C_{\delta,\kappa}\).
Also \(e^d\leq1+e^{2d}\).  Thus the geometric edge requires at most
the mismatch Lyapunov function
\[
 V_{g,2}(d)=e^{2d}+e^{-2d}=2\cosh(2d).             \tag{20.11}
\]

For the arithmetic edge, the anchor coefficients are
\[
\begin{aligned}
 \cosh d\,e^d&=\tfrac12(1+e^{2d}),\\
 \cosh d\,e^{-d}&=\tfrac12(1+e^{-2d}).
\end{aligned}                                                       \tag{20.12}
\]
Absorbing the second coefficient times \(\|y_y\|^2\) into its
quartic squares that coefficient and produces \(e^{-4d}\).
Endpoint exchange produces \(e^{4d}\).  Hence the natural
arithmetic-edge mismatch function is
\[
 V_{g,4}(d)=e^{4d}+e^{-4d}=2\cosh(4d).             \tag{20.13}
\]

\begin{theorem}[Mismatch-exponent comparison]
\label{thm:v17v20-mismatch}
The V15 anchor compiler applied to:
\begin{enumerate}
\item the geometric edge closes with \(V_{g,2}\);
\item the arithmetic endpoint edge closes with \(V_{g,4}\).
\end{enumerate}
These exponents are sharp for this anchor scheme: along
\(d\to-\infty\) with \(\|y_y\|^2\) optimized against a positive
quartic, the induced debits grow respectively like
\(e^{-2d}\) and \(e^{-4d}\).
\end{theorem}

\begin{proof}
The upper bounds were derived in (20.9)--(20.13).  For sharpness,
minimize
\[
 \lambda r^4-br^2
\]
with \(b=\kappa e^{-d}\) in the geometric case and
\(b=\frac\kappa2(1+e^{-2d})\) in the arithmetic case.  The minimum
is \(-b^2/(4\lambda)\), which has the stated asymptotic orders.
\end{proof}

The geometric-mean edge is therefore strictly less demanding on
the conformal conditional tail.  This comparison does not select
the physical discretization; that choice must come from the common
action and continuum consistency.

\subsection{Block-matrix consequence}

Let \(\theta_H\|y\|^4\) be the Higgs Lyapunov weight and let
\(\theta_gV_{g,2}(d)\) be the geometric-edge mismatch weight.
Equations (20.9)--(20.10) give an \(H\)-row with schematic entries
\[
 B_{HH}(\delta)
 =\frac{C_H\delta}{\theta_H},\qquad
 B_{Hg}(\delta)
 =\frac{C_g}{\theta_g}
 \left(\kappa+\frac{\kappa^2}{\delta}\right).       \tag{20.14}
\]
The constants count only the uniformly bounded incident edges.
Making the neighbor-Higgs entry small enlarges the conformal entry.
The gravity-to-Higgs entry \(B_{gH}\) must therefore be computed
before the two-sector condition
\[
 B_{Hg}B_{gH}
 <(1-B_{HH})(1-B_{gg})                             \tag{20.15}
\]
can be tested.

\begin{assessment}[What V20 closes]
\label{ass:v17v20-status}
V20 performs the required companion audit and derives the exact
two-cell rescaled hopping forms.  The geometric-mean edge induces
the mismatch exponent \(2\), whereas the arithmetic endpoint edge
induces exponent \(4\).  Both preserve positivity, but their
normalization anchors generate different Higgs-to-gravity
influence costs.

No source formula chooses between these weights or supplies the
reverse gravitational row.  The full spectral-radius gate remains
open.
\end{assessment}

\begin{decision}[V21: optimize the two-sector influence cycle]
\label{dec:v17v20-next}
Continue after the audit with the geometric edge.  Derive explicit
conditional Higgs constants, retain the free Young parameter
\(\delta\), and solve the exact minimization of the two-sector
Perron condition over \(\delta\).  State the resulting necessary
coefficient target for the missing gravitational row.
\end{decision}

\begin{firewall}[Claim boundary after V20]
\label{fire:v17v20}
V20 compares two mathematically admissible edge weights.  It does
not identify the RPESM discretization, certify a gravitational
conditional law, or import the finite YM cell theorem as a
continuum Wilson window.
\end{firewall}

\section*{Version V20 append-only record}
\addcontentsline{toc}{section}{Version V20 append-only record}
The complete V19 source was retained byte for byte before this
continuation.  It had \(199062\) bytes, \(6266\) lines, and SHA--256
\[
 \texttt{F5E47C8397D6639404C603EF6EF905FE742C7D46B53F15823CF2F97A2F955007}.
\]
V20 performed the companion audit and proved the exact geometric
and arithmetic conformal-edge mismatch exponents.


\section{V21: exact optimization of the Higgs--gravity influence cycle}
\label{sec:v17v21}

\subsection{Explicit Higgs row for the geometric edge}

Condition on the conformal variables and on every rescaled Higgs
field except \(y_x\).  For one incident edge \(x\sim y\), V20 gives
\[
 E_{G,xy}
 =\kappa\left(
 e^d\|y_x\|^2+e^{-d}\|y_y\|^2
 -2\operatorname{Re}\langle y_y,Uy_x\rangle
 \right),
 \qquad d=s_y-s_x.                                 \tag{21.1}
\]
Use
\[
 V_H(y)=\|y\|^4,\qquad
 V_g(d)=e^{2d}+e^{-2d}.                            \tag{21.2}
\]

\begin{lemma}[Explicit anchor transfer]
\label{lem:v17v21-anchor}
On the anchor \(\|y_x\|\leq1\), for every \(\delta>0\),
\[
 E_{G,xy}
 \leq
 2\delta V_H(y_y)
 +\left(\frac\kappa2+\frac{\kappa^2}{4\delta}\right)V_g(d)
 +C_{\kappa,\delta},                               \tag{21.3}
\]
where one may take
\[
 C_{\kappa,\delta}
 =\frac\kappa2+
 \frac{3\kappa}{2}
 \left(\frac{\kappa}{2\delta}\right)^{1/3}.         \tag{21.4}
\]
\end{lemma}

\begin{proof}
From (21.1) and \(\|y_x\|\leq1\),
\[
 E_{G,xy}
 \leq\kappa e^d+\kappa e^{-d}\|y_y\|^2
 +2\kappa\|y_y\|.
\]
First,
\[
 \kappa e^d\leq\frac\kappa2+\frac\kappa2e^{2d}
 \leq\frac\kappa2+\frac\kappa2V_g(d).
\]
Second,
\[
 \kappa e^{-d}\|y_y\|^2
 \leq\delta\|y_y\|^4+
 \frac{\kappa^2}{4\delta}e^{-2d}.
\]
Finally, maximization of
\(2\kappa r-\delta r^4\) over \(r\geq0\) gives
\[
 2\kappa r
 \leq\delta r^4+
 \frac{3\kappa}{2}
 \left(\frac{\kappa}{2\delta}\right)^{1/3}.
\]
Adding the three bounds proves (21.3)--(21.4).
\end{proof}

\begin{theorem}[Aggregated Higgs influence row]
\label{thm:v17v21-higgs-row}
Suppose at most \(D\) geometric-mean edges meet a Higgs site, and
use exponential parameters \(\theta_H,\theta_g>0\).  Apart from
pre-existing transfers \(a_0,b_{\rm ext}\geq0\), the conditional
Higgs row can be chosen as
\[
\begin{aligned}
 B_{HH}(\delta)
 &=a_0+A\delta,
 &A&=\frac{2D}{\theta_H},\\
 B_{Hg}(\delta)
 &=b_0+\frac{B}{\delta},
 &b_0&=b_{\rm ext}+\frac{D\kappa}{2\theta_g},
 &B&=\frac{D\kappa^2}{4\theta_g}.                  \tag{21.5}
\end{aligned}
\]
\end{theorem}

\begin{proof}
Sum (21.3) over at most \(D\) incident edges.  Divide the
neighbor-Higgs coefficients by \(\theta_H\) and the mismatch
coefficients by \(\theta_g\), as in the energy-to-moment compiler
(15.11).  Constants such as (21.4) affect \(C_i\), not the influence
matrix.
\end{proof}

The positive edge square causes no coercive lower-bound debit.
All entries in (21.5) arise from the normalization anchor.

\subsection{The unknown reverse row}

Write the aggregated gravitational conditional row as
\[
 B_{gH}=c,\qquad B_{gg}=d,                          \tag{21.6}
\]
where \(c,d\geq0\).  These are precisely the two numbers that an
explicit gravitational energy and reference measure must supply.
The two-sector influence matrix is
\[
 {\mathbb B}(\delta)
 =
 \begin{pmatrix}
 a_0+A\delta&b_0+B/\delta\\
 c&d
 \end{pmatrix}.                                    \tag{21.7}
\]

\subsection{Exact optimization}

\begin{theorem}[Optimal Young parameter and Perron margin]
\label{thm:v17v21-optimal}
Assume
\[
 A>0,\quad B>0,\quad c>0,\quad
 0\leq a_0<1,\quad0\leq d<1,\quad b_0\geq0.        \tag{21.8}
\]
Define
\[
\begin{aligned}
 P&=A(1-d),&
 Q&=cB,\\
 R&=(1-d)(1-a_0)-cb_0.                             \tag{21.9}
\end{aligned}
\]
There exists \(\delta>0\) with
\(\rho({\mathbb B}(\delta))<1\) if and only if
\[
\boxed{R>2\sqrt{PQ}.}                              \tag{21.10}
\]
The determinant margin
\(\det(I-{\mathbb B}(\delta))\) is maximized at
\[
\boxed{\delta_*=\sqrt{\frac{Q}{P}}
=\sqrt{\frac{cB}{A(1-d)}}}                         \tag{21.11}
\]
and its maximum is
\[
\boxed{R-2\sqrt{PQ}.}                              \tag{21.12}
\]
\end{theorem}

\begin{proof}
The two-sector criterion (15.14) requires
\[
 a_0+A\delta<1,\qquad d<1,
\]
and
\[
\begin{aligned}
 0&<
 (1-d)(1-a_0-A\delta)
 -c(b_0+B/\delta)\\
 &=R-P\delta-Q/\delta.                             \tag{21.13}
\end{aligned}
\]
The arithmetic--geometric mean inequality gives
\[
 P\delta+Q/\delta\geq2\sqrt{PQ},
\]
with equality exactly at (21.11).  This proves necessity and the
maximum (21.12).  If (21.10) holds, then (21.13) is positive at
\(\delta_*\).  Moreover,
\[
 A\delta_*
 =\sqrt{\frac{AcB}{1-d}}
 <\frac{1-a_0}{2},
\]
where the last inequality follows from
\(2\sqrt{PQ}<R\leq(1-d)(1-a_0)\).  Thus the first diagonal
condition also holds, proving sufficiency.
\end{proof}

\begin{corollary}[Complete admissible interval]
\label{cor:v17v21-interval}
Under (21.8)--(21.10), every
\[
 \delta\in(\delta_-,\delta_+)                      \tag{21.14}
\]
passes the Perron test, where
\[
 \delta_\pm
 =\frac{R\pm\sqrt{R^2-4PQ}}{2P}.                  \tag{21.15}
\]
No other positive \(\delta\) passes.
\end{corollary}

\begin{proof}
Multiplying (21.13) by \(\delta>0\) gives
\[
 P\delta^2-R\delta+Q<0.
\]
Its two positive roots are (21.15), and an upward-opening
quadratic is negative exactly between them.
\end{proof}

\subsection{Degenerate cycle cases}

\begin{proposition}[One-way and zero-anchor limits]
\label{prop:v17v21-degenerate}
If \(cB=0\), the feedback singularity \(Q/\delta\) is absent.
Assuming \(d<1\), a passing \(\delta\) exists exactly when
\[
 (1-d)(1-a_0)-cb_0>0,                              \tag{21.16}
\]
with additionally \(\delta<(1-a_0)/A\) when \(A>0\).
If \(A=0\) but \(cB>0\), arbitrarily large \(\delta\) removes the
\(B/\delta\) debit, and the condition is again (21.16).
\end{proposition}

\begin{proof}
Set \(Q=0\) or \(P=0\) in (21.13) and solve the resulting monotone
inequality.  The explicit diagonal restriction is retained when
the Higgs diagonal grows with \(\delta\).
\end{proof}

\subsection{An explicit target for the missing gravity-to-Higgs row}

Let
\[
 S=\sqrt{A(1-d)B},\qquad
 T=(1-d)(1-a_0).                                   \tag{21.17}
\]
Condition (21.10) is
\[
 b_0c+2S\sqrt c<T.                                 \tag{21.18}
\]

\begin{corollary}[Maximum admissible reverse coupling]
\label{cor:v17v21-cmax}
If \(b_0>0\), (21.18) holds exactly when
\[
\boxed{
 0\leq c<
 \left(
 \frac{\sqrt{S^2+b_0T}-S}{b_0}
 \right)^2.}                                       \tag{21.19}
\]
If \(b_0=0\), it holds exactly when
\[
\boxed{
 0\leq c<
 \frac{(1-d)(1-a_0)^2}{4AB}.}                      \tag{21.20}
\]
\end{corollary}

\begin{proof}
Put \(x=\sqrt c\).  For \(b_0>0\), inequality (21.18) becomes
\[
 b_0x^2+2Sx-T<0.
\]
Its unique positive root is
\((\sqrt{S^2+b_0T}-S)/b_0\).  For \(b_0=0\), solve
\(2Sx<T\) and square.
\end{proof}

Equations (21.19)--(21.20) are the quantitative acquisition target
for the missing gravitational conditional calculation.  Merely
showing \(c<1\) is insufficient.

\subsection{Insertion of the geometric-edge constants}

With no external Higgs transfers,
\[
 a_0=b_{\rm ext}=0,
\]
the coefficients from (21.5) are
\[
 A=\frac{2D}{\theta_H},\qquad
 b_0=\frac{D\kappa}{2\theta_g},\qquad
 B=\frac{D\kappa^2}{4\theta_g}.                    \tag{21.21}
\]
The optimized sufficient and necessary condition within this
anchor scheme is therefore
\[
\boxed{
 (1-d)-\frac{D\kappa}{2\theta_g}c
 >
 2\sqrt{
 \frac{D^2\kappa^2}{2\theta_H\theta_g}
 (1-d)c}.}                                         \tag{21.22}
\]
Every quantity in (21.22) is known from the Higgs edge except the
gravitational entries \(c,d\) and the admissible conformal
exponential parameter \(\theta_g\).

\subsection{Sharp failure witness}

\begin{countertheorem}[No Young choice beyond the cycle boundary]
\label{ctr:v17v21-no-delta}
If
\[
 R\leq2\sqrt{PQ},                                  \tag{21.23}
\]
then no redistribution between neighbor-Higgs absorption and
conformal-mismatch absorption can make (21.7) pass the V15
spectral-radius gate.
\end{countertheorem}

\begin{proof}
For every \(\delta>0\),
\[
 \det(I-{\mathbb B}(\delta))
 =R-P\delta-Q/\delta
 \leq R-2\sqrt{PQ}\leq0.
\]
By the exact two-sector criterion, the Perron root cannot be below
one.
\end{proof}

This countertheorem concerns the stated anchor/Young certificate.
It does not prove that the physical measure diverges; another
Lyapunov function or a cancellation could yield a different
matrix.

\begin{assessment}[What V21 closes]
\label{ass:v17v21-status}
V21 derives the geometric-edge Higgs row with explicit constants
and solves its two-sector Perron optimization exactly.  The optimal
Young parameter, complete admissible interval, cycle margin, and
maximum permitted gravity-to-Higgs influence are now closed
finite-dimensional formulas.

The source does not supply \(c\), \(d\), or \(\theta_g\), so
(21.22) cannot yet be evaluated on the full RPESM cylinder.  The
result replaces an unspecified small-coupling demand by a precise
coefficient acquisition target.
\end{assessment}

\begin{decision}[V22: derive a gravitational row from a model energy]
\label{dec:v17v21-next}
Take the first conformal energy that passes V19,
\(U(s)=\kappa_gs^2+b_gs+c_g\) with \(\kappa_g>0\), and couple
neighboring mismatches through a convex
\(\cosh(2(s_y-s_x))\) term.  Derive its conditional moment row and
test the explicit coefficients against (21.22).  Keep this as a
model certificate until the RPESM source supplies its \(S_{g,n}\).
\end{decision}

\begin{firewall}[Claim boundary after V21]
\label{fire:v17v21}
V21 optimizes a sufficient conditional certificate.  It does not
manufacture the missing gravitational action, select the physical
edge discretization, or prove the Einstein--SM continuum limit.
\end{firewall}

\section*{Version V21 append-only record}
\addcontentsline{toc}{section}{Version V21 append-only record}
The complete V20 source was retained byte for byte before this
continuation.  It had \(206946\) bytes, \(6515\) lines, and SHA--256
\[
 \texttt{60B3511D89B68C22688EDA6AFAD3F08DA9C9F64E86FF11D05536FF493E4AD1E2}.
\]
V21 proved the exact optimized two-sector influence-cycle
criterion for the geometric conformal edge.


\section{V22: a two-cell conformal model and its exact reverse influence}
\label{sec:v17v22}

\subsection{Status of the model}

The RPESM source still leaves the gravitational cylinder action
\(S_{g,n}\) unexpanded.  This section therefore studies the smallest
explicit conformal action that meets the V19 coercivity requirement and
contains the mismatch function forced by the geometric Higgs edge.  The
calculation is a model certificate and a diagnostic for the eventual
physical action; it is not an identification of that action.

Let \(s_1,s_2\in\mathbb R\), put
\[
 t=\frac{s_1+s_2}{2},\qquad d=s_2-s_1,              \tag{22.1}
\]
and consider
\[
 S_g(s_1,s_2)
 =\kappa_g(s_1^2+s_2^2)+b_g(s_1+s_2)
 +J\bigl(\cosh(2d)-1\bigr).                        \tag{22.2}
\]

\begin{theorem}[Exact two-cell normalizability and factorization]
\label{thm:v17v22-factor}
For real \(\kappa_g,b_g,J\), the density
\(\exp(-S_g)\,ds_1ds_2\) has finite, nonzero mass if and only if
\[
 \kappa_g>0\qquad\hbox{and}\qquad J\geq0.           \tag{22.3}
\]
Under (22.3), \(t\) and \(d\) are independent, with unnormalized
densities
\[
 \exp(-2\kappa_gt^2-2b_gt)\,dt                     \tag{22.4}
\]
and
\[
 \exp\left[-\frac{\kappa_g}{2}d^2
 -J\bigl(\cosh(2d)-1\bigr)\right],dd.             \tag{22.5}
\]
\end{theorem}

\begin{proof}
The inverse transformation is
\(s_1=t-d/2\), \(s_2=t+d/2\), and its Jacobian has absolute value one.
Moreover,
\[
 s_1^2+s_2^2=2t^2+\frac{d^2}{2},
 \qquad s_1+s_2=2t.                                \tag{22.6}
\]
This proves the factorization.  If \(\kappa_g>0\) and \(J\geq0\),
both factors are integrable and strictly positive.  If
\(\kappa_g=0\), the \(t\)-factor is either constant or grows on one
half-line; if \(\kappa_g<0\), it grows quadratically.  Thus
\(\kappa_g>0\) is necessary.  With \(\kappa_g>0\) and \(J<0\), the
\(d\)-integrand contains
\(\exp(|J|\cosh(2d)-\kappa_gd^2/2)\), which is not integrable.
Hence \(J\geq0\) is also necessary.
\end{proof}

\subsection{The correct exponential variable is an edge mismatch}

Recall the V21 gravity variable
\[
 V_g(d)=e^{2d}+e^{-2d}=2\cosh(2d).                 \tag{22.7}
\]

\begin{theorem}[Sharp mismatch-moment threshold]
\label{thm:v17v22-mismatch}
Assume (22.3).  If \(J>0\), then for \(\theta\geq0\),
\[
 \mathbb E_g\exp\bigl(\theta V_g(d)\bigr)<\infty
 \quad\Longleftrightarrow\quad 2\theta\leq J.     \tag{22.8}
\]
If \(J=0\), the same moment is finite only for \(\theta=0\).
Furthermore,
\[
 \mathbb E_g e^{\eta d^2}<\infty                  \tag{22.9}
\]
for every \(\eta\geq0\) when \(J>0\), whereas for \(J=0\) it is
finite exactly when \(0\leq\eta<\kappa_g/2\).
\end{theorem}

\begin{proof}
After multiplying (22.5) by \(e^{2\theta\cosh(2d)}\), the only
non-Gaussian tail factor is
\[
 \exp\bigl(-(J-2\theta)\cosh(2d)\bigr).            \tag{22.10}
\]
It gives super-exponential decay when \(2\theta<J\).  At equality it
disappears and the remaining Gaussian is integrable.  If
\(2\theta>J\), (22.10) grows faster than the Gaussian decays.  This
proves (22.8), including the assertion for \(J=0\).  The same tail
comparison proves (22.9); when \(J=0\), it reduces to the elementary
Gaussian threshold.
\end{proof}

\begin{countertheorem}[Absolute conformal cosh is inadmissible]
\label{ctr:v17v22-common-mode}
Under (22.3), for every \(\theta>0\) and either endpoint \(i=1,2\),
\[
 \mathbb E_g\exp\bigl(\theta\cosh(2s_i)\bigr)
 =\infty.                                          \tag{22.11}
\]
Thus the successful mismatch moment (22.8) cannot be replaced by an
absolute site-cosh moment.
\end{countertheorem}

\begin{proof}
Fix \(d\).  For example, \(s_1=t-d/2\), and the integral in the common
mode contains
\[
 \int_{\mathbb R}
 \exp\bigl(-2\kappa_gt^2-2b_gt
 +\theta\cosh(2t-d)\bigr),dt.                     \tag{22.12}
\]
For sufficiently large positive \(t\),
\(\cosh(2t-d)\geq e^{2t-d}/2\), so the positive exponential term
dominates every quadratic polynomial in \(t\).  The integral (22.12)
is infinite for every fixed \(d\).  Tonelli's theorem applied to the
nonnegative joint integrand proves (22.11).  The proof for \(s_2\) is
identical.
\end{proof}

The obstruction is structural.  The mismatch interaction controls the
difference mode, while the common rescaling mode remains Gaussian.  A
site variable growing like \(\cosh(2s_i)\) therefore asks a Gaussian
tail to support a double exponential.

\subsection{Coupling the model to the rescaled Higgs edge}

Let \(y_1,y_2\) be the rescaled Higgs variables of V20 and let
\(\chi>0\) be the geometric-edge hopping coefficient.  Conditional on
the Higgs endpoints, the part of the action depending on \(d\) is
\[
 \Phi_{a,b}(d)
 =\frac{\kappa_g}{2}d^2+J\bigl(\cosh(2d)-1\bigr)
 +\chi\bigl(ae^d+be^{-d}\bigr),                   \tag{22.13}
\]
where
\[
 a=\|y_1\|^2,qquad b=\|y_2\|^2.                  \tag{22.14}
\]
The cross term
\(-2\chi\operatorname{Re}\langle y_2,Uy_1\rangle\)
is independent of \(d\), hence cancels exactly from the conditional
normalization ratio.

\begin{lemma}[Uniform conditional mismatch moment]
\label{lem:v17v22-reverse}
Assume \(\kappa_g>0\), \(J>0\), and
\[
 0<2\theta_g\leq J.                                \tag{22.15}
\]
For every anchor width \(\varepsilon>0\) and every \(\zeta>0\),
there is a finite constant
\(C=C(\kappa_g,J,\theta_g,\chi,\varepsilon,\zeta)\), independent of
\(y_1,y_2\), such that
\[
 \log\mathbb E\left[
 e^{\theta_gV_g(d)}\mid y_1,y_2\right]
 \leq C+chi e^{\varepsilon}\zeta
 \bigl(V_H(y_1)+V_H(y_2)\bigr),                   \tag{22.16}
\]
where \(V_H(y)=\|y\|^4\).
\end{lemma}

\begin{proof}
Let
\[
 K=\int_{\mathbb R}
 \exp\left[-\frac{\kappa_g}{2}d^2
 -(J-2\theta_g)\cosh(2d)\right],dd.              \tag{22.17}
\]
Condition (22.15) and \(\kappa_g>0\) imply \(K<\infty\), including
the endpoint \(2\theta_g=J\).  Dropping the nonnegative hopping terms
in the numerator gives the upper bound \(e^JK\).

On the denominator anchor \(|d|\leq\varepsilon\),
\[
 \Phi_{a,b}(d)
 \leq\frac{\kappa_g}{2}\varepsilon^2
 +J\bigl(\cosh(2\varepsilon)-1\bigr)
 +\chi e^{\varepsilon}(a+b).                      \tag{22.18}
\]
The denominator is consequently at least \(2\varepsilon\) times the
exponential of the negative right-hand side of (22.18).  Taking the
ratio and its logarithm yields
\[
 \log\mathbb E\left[e^{\theta_gV_g(d)}\mid y_1,y_2\right]
 \leq C_0+\chi e^{\varepsilon}(a+b),               \tag{22.19}
\]
with finite \(C_0\).  Finally,
\[
 q\leq\zeta q^2+\frac{1}{4\zeta},
 \qquad q\geq0,                                   \tag{22.20}
\]
applied to \(a\) and \(b\), proves (22.16).
\end{proof}

\subsection{An explicit passing two-cell influence matrix}

Use the maximum envelope for the two Higgs endpoint blocks and the one
gravity-mismatch block.  With no external transfers, V21 gives
\[
 A=\frac{2}{\theta_H},\qquad
 b_0=\frac{\chi}{2\theta_g},\qquad
 B=\frac{\chi^2}{4\theta_g}.                       \tag{22.21}
\]
Lemma~\ref{lem:v17v22-reverse} supplies
\[
 c=\frac{2\chi e^{\varepsilon}\zeta}{\theta_H},
 \qquad d_{
m inf}=0.                             \tag{22.22}
\]
The factor two in \(c\) sums the two endpoint coefficients in
(22.16).  The symbol \(d_{\rm inf}\) distinguishes the gravitational
self-influence entry from the mismatch coordinate \(d\).

\begin{theorem}[Two-cell Perron certificate]
\label{thm:v17v22-perron}
Assume (22.15) and let \(\theta_H>0\) be an admissible Higgs quartic
moment parameter.  Define
\[
 x=\frac{\chi^2e^{\varepsilon}\zeta}
 {\theta_H\theta_g},
 \qquad q=\frac{\chi}{\theta_H}.                  \tag{22.23}
\]
The V21 optimized two-sector certificate passes exactly when
\[
 1-x>2\sqrt{qx},                                   \tag{22.24}
\]
or equivalently when
\[
\boxed{
 0<\zeta<
 \frac{\theta_H\theta_g}{\chi^2e^{\varepsilon}}
 \left(\sqrt{1+\frac{\chi}{\theta_H}}
 -\sqrt{\frac{\chi}{\theta_H}}\right)^2.}         \tag{22.25}
\]
For such a \(\zeta\), the optimal forward Young parameter is
\[
 \delta_*
 =\frac12\sqrt{\frac{\chi^3e^{\varepsilon}\zeta}
 {\theta_g}}.                                      \tag{22.26}
\]
In particular, the admissible interval in (22.25) is nonempty for
every fixed positive \(\chi,\theta_H,\theta_g\).
\end{theorem}

\begin{proof}
Insert (22.21)--(22.22) and \(d_{\rm inf}=0\) into (21.22).  Its left
side becomes \(1-x\), and its right side becomes
\(2\sqrt{qx}\), proving (22.24).  Put \(u=\sqrt{x}\).  The inequality
is
\[
 u^2+2\sqrt q,u-1<0,                              \tag{22.27}
\]
which for \(u\geq0\) is equivalent to
\(u<\sqrt{1+q}-\sqrt q\).  Substitution of (22.23) gives (22.25).
Finally, (21.11) with the coefficients (22.21)--(22.22) gives
(22.26).
\end{proof}

The small parameter in (22.25) is the Young transfer \(\zeta\), not a
restriction that the physical hopping \(\chi\) vanish.  Sending
\(\zeta\) downward enlarges the additive constant in (22.16), but the
constant does not enter the Perron matrix.

\begin{assessment}[What V22 establishes]
\label{ass:v17v22-status}
For an isolated conformal edge, the convex mismatch action gives the
exact exponential window \(2\theta_g\leq J\), an explicit
gravity-to-Higgs influence coefficient, and a nonempty parameter range
in which the complete two-sector matrix contracts.  This supplies the
first explicit value of the V21 reverse row within a declared model.

V22 also proves that an absolute site-cosh Lyapunov function is
impossible for the quadratic common conformal mode.  Subsequent gravity
blocks must therefore retain edge mismatches or strengthen the
common-mode potential beyond a quadratic.
\end{assessment}

\begin{decision}[V23: pass from one edge to a finite graph]
\label{dec:v17v22-next}
Extend (22.13)--(22.16) to a bounded-degree connected graph.  Separate
the common mode with the incidence operator, keep mismatch variables on
edges, and quantify the dependence created by overlapping edges and
cycle constraints.  If edge-by-edge conditioning becomes circular,
condition a spanning-tree mismatch block and treat the remaining cycle
edges as nonnegative coercive terms.
\end{decision}

\begin{firewall}[Claim boundary after V22]
\label{fire:v17v22}
The theorem is exact for the stated two-cell model.  It does not prove a
volume-uniform graph estimate, derive the missing RPESM gravitational
action, or establish the Einstein--SM continuum limit.
\end{firewall}

\section*{Version V22 append-only record}
\addcontentsline{toc}{section}{Version V22 append-only record}
The complete V21 source was retained byte for byte before this
continuation.  It had \(216871\) bytes, \(6881\) lines, and SHA--256
\[
 \texttt{F5BAC6A4F5C4DF9A1D961BE9A9FB5C583522273B735FE2DBC459306999C0E303}.
\]
V22 constructed and closed the isolated two-cell conformal mismatch
certificate.


\section{V23: finite-graph mismatch moments and the global-block bottleneck}
\label{sec:v17v23}

\subsection{A finite conformal graph}

Let \(G=(\mathcal V,\mathcal E)\) be a finite connected simple graph,
with \(n=|\mathcal V|\) and \(m=|\mathcal E|\).  Give every edge an
arbitrary orientation and write
\[
 d_e=s_{e_+}-s_{e_-}.                              \tag{23.1}
\]
Consider the homogeneous conformal action
\[
 S_{g,G}(s)
 =\kappa_g\sum_{x\in\mathcal V}s_x^2
 +b_g\sum_{x\in\mathcal V}s_x
 +J\sum_{e\in\mathcal E}\bigl(\cosh(2d_e)-1\bigr).
                                                               \tag{23.2}
\]

\begin{theorem}[Common-mode factorization on a graph]
\label{thm:v17v23-common}
The measure \(e^{-S_{g,G}(s)}\,ds\) is finite and nonzero exactly when
\[
 \kappa_g>0\qquad\hbox{and}\qquad J\geq0.           \tag{23.3}
\]
Under (23.3), put
\[
 t=\frac1n\sum_xs_x,qquad r=s-t\mathbf1,qquad
 \sum_xr_x=0.                                     \tag{23.4}
\]
Then the common mode \(t\) and the mean-zero field \(r\) are
independent.  Their action, up to the constant Jacobian of the
orthogonal change of variables, is
\[
 n\kappa_gt^2+nb_gt
 +\kappa_g\|r\|_2^2
 +J\sum_e\bigl(\cosh(2(B r)_e)-1\bigr),            \tag{23.5}
\]
where \(B\) is the oriented incidence operator.
\end{theorem}

\begin{proof}
Since \(r\perp\mathbf1\),
\[
 \sum_xs_x^2=nt^2+\|r\|_2^2,
 \qquad \sum_xs_x=nt,
 \qquad Bs=Br.                                    \tag{23.6}
\]
This gives (23.5) and independence.  Under (23.3), the quadratic term
alone supplies an integrable upper bound after completing the square.
If \(\kappa_g\leq0\), the one-dimensional common-mode integral is
infinite.  If \(\kappa_g>0\) but \(J<0\), take any nonconstant vector
\(h\) and set \(s=Lh\).  On a fixed-width tube about this ray, the
positive term \(|J|\cosh(2L(Bh)_e)\) for an edge with
\((Bh)_e\ne0\) dominates the negative quadratic term.  The integral is
therefore infinite.
\end{proof}

\begin{countertheorem}[Graph-wide absolute-site obstruction]
\label{ctr:v17v23-site}
Under (23.3), for every vertex \(x\) and every \(\theta>0\),
\[
 \mathbb E_{g,G}\exp\bigl(\theta\cosh(2s_x)\bigr)=\infty.
                                                               \tag{23.7}
\]
\end{countertheorem}

\begin{proof}
Conditional on any fixed mean-zero field \(r\), one has \(s_x=t+r_x\)
and the \(t\)-density is Gaussian by (23.5).  The integral
\[
 \int_{\mathbb R}
 \exp\bigl(-n\kappa_gt^2-nb_gt
 +\theta\cosh(2t+2r_x)\bigr),dt                  \tag{23.8}
\]
diverges because the last term grows exponentially in \(|t|\) inside
the outer exponential.  Tonelli's theorem completes the proof.
\end{proof}

\subsection{The single-edge threshold is graph independent}

For an edge \(e\), retain
\[
 V_{g,e}(s)=e^{2d_e}+e^{-2d_e}=2\cosh(2d_e).       \tag{23.9}
\]

\begin{theorem}[Sharp finite-graph edge threshold]
\label{thm:v17v23-edge-threshold}
Assume \(J>0\) and \(\kappa_g>0\).  For every edge \(e\) and
\(\theta\geq0\),
\[
 \mathbb E_{g,G}e^{\theta V_{g,e}}<\infty
 \quad\Longleftrightarrow\quad 2\theta\leq J.     \tag{23.10}
\]
The threshold is independent of the size, degree, and cycle structure
of the finite simple graph.
\end{theorem}

\begin{proof}
If \(2\theta\leq J\), the moment numerator merely replaces the
coefficient of \(\cosh(2d_e)\) by \(J-2\theta\geq0\).  Dropping all
nonnegative mismatch terms leaves a Gaussian integral, hence the
moment is finite.

For necessity, write \(e=\{u,v\}\).  Because the graph is simple,
choose a vector \(h\) with
\[
 h_u=0,qquad h_v=1,qquad
 \frac13<h_w<\frac23\quad(w\ne u,v),               \tag{23.11}
\]
choosing the remaining values distinct if desired.  There is then an
\(\eta>0\) such that
\[
 |h_v-h_u|=1,qquad |h_{f_+}-h_{f_-}|\leq1-\eta
 \quad(f\ne e).                                   \tag{23.12}
\]
On any fixed-width box about \(Lh\), the logarithm of the moment
integrand contains
\[
 (2\theta-J)\cosh(2L+O(1))                         \tag{23.13}
\]
from the distinguished edge.  Every other mismatch term is bounded in
absolute value by \(C e^{2(1-\eta)L}\), while the quadratic and linear
terms are \(O(L^2)\).  If \(2\theta>J\), (23.13) dominates all of them
and tends to positive infinity like \(e^{2L}\).  Taking disjoint
fixed-width boxes along a sequence \(L\to\infty\) proves divergence.
\end{proof}

The simplicity assumption only excludes parallel copies of the
distinguished edge.  With parallel edges, their combined mismatch
coefficient replaces \(J\) in (23.10).

\subsection{A coupled whole-graph conditional estimate}

Orient each edge \(e=(x,y)\), let \(a_x=\|y_x\|^2\), and give it a
hopping coefficient \(\chi_e\geq0\).  Conditional on all rescaled Higgs
variables, add
\[
 S_{H\!g}(s\mid y)
 =\sum_{e=(x,y)}\chi_e
 \bigl(a_xe^{d_e}+a_ye^{-d_e}\bigr).              \tag{23.14}
\]
The cross terms of the positive Higgs edge squares are independent of
\(s\) and cancel from every conformal conditional ratio.  Define the
incident hopping weight
\[
 h_x=\sum_{e\ni x}\chi_e.                          \tag{23.15}
\]

\begin{theorem}[Explicit global-block reverse bound]
\label{thm:v17v23-global-bound}
Assume \(\kappa_g>0\), \(J>0\), and
\(0<2\theta_g\leq J\).  For every \(\varepsilon,\zeta>0\), every edge
\(e_0\), and every finite graph as above,
\[
 \log\mathbb E\left[e^{\theta_gV_{g,e_0}}
 \mid (y_x)_{x\in\mathcal V}\right]
 \leq C_{G,\varepsilon,\zeta}
 +e^{2\varepsilon}\zeta
 \sum_{x\in\mathcal V}h_xV_H(y_x),                \tag{23.16}
\]
where \(C_{G,\varepsilon,\zeta}<\infty\) is explicit and independent
of the Higgs values.
\end{theorem}

\begin{proof}
In the numerator, the insertion changes the coefficient of the
distinguished cosh from \(J\) to \(J-2\theta_g\geq0\).  Drop all
nonnegative mismatch and hopping terms.  Completing the square gives
the finite upper bound
\[
 e^{Jm}
 \left(\frac{\pi}{\kappa_g}\right)^{n/2}
 \exp\left(\frac{nb_g^2}{4\kappa_g}\right).       \tag{23.17}
\]

For the denominator, use the cube \(|s_x|\leq\varepsilon\) for every
vertex.  On this cube, \(|d_e|\leq2\varepsilon\) and
\[
\begin{aligned}
 \kappa_g\sum_xs_x^2&\leq n\kappa_g\varepsilon^2,\\
 b_g\sum_xs_x&\leq n|b_g|\varepsilon,\\
 J\sum_e(\cosh(2d_e)-1)
 &\leq Jm(\cosh(4\varepsilon)-1),\\
 S_{H\!g}(s\mid y)
 &\leq e^{2\varepsilon}\sum_xh_xa_x.              \tag{23.18}
\end{aligned}
\]
The cube has volume \((2\varepsilon)^n\), so (23.17)--(23.18) give
\[
 \log\mathbb E[e^{\theta_gV_{g,e_0}}\mid y]
 \leq C_{G,\varepsilon}
 +e^{2\varepsilon}\sum_xh_xa_x.                  \tag{23.19}
\]
Apply \(a_x\leq\zeta a_x^2+(4\zeta)^{-1}\) and absorb the constant
sum into \(C_{G,\varepsilon,\zeta}\).
\end{proof}

For reference, one permissible constant is
\[
\begin{aligned}
 C_{G,\varepsilon,\zeta}
={}&Jm+\frac n2\log\frac{\pi}{\kappa_g}
 +\frac{nb_g^2}{4\kappa_g}-n\log(2\varepsilon)
 +n\kappa_g\varepsilon^2+n|b_g|\varepsilon\\
 &+Jm(\cosh(4\varepsilon)-1)
 +\frac{e^{2\varepsilon}}{4\zeta}\sum_xh_x.
                                                               \tag{23.20}
\end{aligned}
\]

\subsection{Why the global block is not a continuum certificate}

If \(\chi_e\leq\chi_*\), then
\[
 \sum_xh_x=2\sum_e\chi_e\leq2m\chi_*.            \tag{23.21}
\]
After compressing all Higgs sites into one maximum envelope, the
reverse influence obtained from (23.16) is
\[
 c_G
 =\frac{e^{2\varepsilon}\zeta}{\theta_H}
 \sum_xh_x
 \leq\frac{2e^{2\varepsilon}\zeta m\chi_*}{\theta_H}.
                                                               \tag{23.22}
\]
Keeping \(c_G\) bounded along an exhaustion forces
\(\zeta=O(m^{-1})\).  The last term of (23.20) then grows as
\(O(m^2)\), while several other terms already grow linearly in volume.

\begin{countertheorem}[Failure of the displayed global-anchor method]
\label{ctr:v17v23-global}
The whole-graph cube anchor (23.18), followed by a single Young
parameter and maximum-envelope compression, cannot provide both a
volume-uniform reverse influence coefficient and a volume-uniform
additive conditional-moment constant on graph exhaustions with
\(m\to\infty\) and hopping weights bounded below on a positive fraction
of edges.
\end{countertheorem}

\begin{proof}
If \(\sum_e\chi_e\geq c_0m\), (23.22) can remain bounded only if
\(\zeta=O(m^{-1})\).  The final term in (23.20) is then bounded below by
a positive constant times \(m^2\).  Conversely, keeping that term
bounded forces \(\zeta\sum_xh_x=O(1)\) with
\(\zeta^{-1}\sum_xh_x=O(1)\), which is impossible when
\(\sum_xh_x\to\infty\).
\end{proof}

This is a failure of the stated global conditioning proof, not of the
finite-volume measure.  The sharp local moment (23.10) remains valid.

\begin{assessment}[What V23 closes]
\label{ass:v17v23-status}
V23 extends the exact mismatch threshold to every finite connected
simple graph and proves that cycles do not reduce it.  It also isolates
the common Gaussian mode and rules out absolute-site cosh Lyapunov
functions on every volume.

The first whole-graph Higgs-conditional bound is explicit, but its
anchor constant and compressed reverse row cannot both remain uniform.
Thus the next step must localize the normalization anchor rather than
enlarge the gravitational block to the whole cylinder.
\end{assessment}

\begin{decision}[V24: a neighbor-relative star anchor]
\label{dec:v17v23-next}
Condition one conformal site at a time.  Anchor it in a fixed interval
about one selected neighbor so that every normalization cost is a
neighbor-to-neighbor mismatch rather than an absolute cosh.  Quantify
the resulting distance-two influence graph and test whether bounded
degree gives a uniform row sum.
\end{decision}

\begin{firewall}[Claim boundary after V23]
\label{fire:v17v23}
No physical RPESM gravity action or infinite-volume Gibbs measure has
been derived.  The obstruction concerns the displayed global cube
anchor; a localized conditional construction may avoid it.
\end{firewall}

\section*{Version V23 append-only record}
\addcontentsline{toc}{section}{Version V23 append-only record}
The complete V22 source was retained byte for byte before this
continuation.  It had \(227611\) bytes, \(7192\) lines, and SHA--256
\[
 \texttt{A1D50BB51D16B6BAD1B3EDAD578D34D47682B6FF17C2AFED62E0B4950A682571}.
\]
V23 proved the finite-graph mismatch threshold and exposed the
nonuniform whole-graph anchor.


\section{V24: the local star anchor and its unavoidable boundary spread}
\label{sec:v17v24}

\subsection{One-site conditioning}

Let \(x\) be a vertex of positive degree \(p_x\), with neighbor set
\(N(x)\).  When the exterior conformal field is fixed, write
\[
 W_x(u)=\sum_{y\sim x}V_g(s_y-u)
 =2\sum_{y\sim x}\cosh\bigl(2(s_y-u)\bigr).        \tag{24.1}
\]
The boundary spread of the star is
\[
 M_x=M_x(s_{N(x)})=inf_{u\in\mathbb R}W_x(u).     \tag{24.2}
\]

\begin{lemma}[Star center]
\label{lem:v17v24-center}
The infimum in (24.2) is attained at a unique point \(u_x^*\), which
lies in the interval
\[
 \min_{y\sim x}s_y\leq u_x^*\leq\max_{y\sim x}s_y. \tag{24.3}
\]
Consequently,
\[
 (u_x^*)^2\leq\max_{y\sim x}s_y^2
 \leq\sum_{y\sim x}s_y^2.                         \tag{24.4}
\]
For every \(|r|\leq\varepsilon\),
\[
 W_x(u_x^*+r)\leq e^{2\varepsilon}M_x.            \tag{24.5}
\]
\end{lemma}

\begin{proof}
The function \(W_x\) is coercive and strictly convex.  Its derivative
is negative to the left of all neighbor values and positive to their
right, proving (24.3); (24.4) follows.  For every real \(c\),
\[
 e^{2(c-r)}+e^{-2(c-r)}
 \leq e^{2|r|}(e^{2c}+e^{-2c}).                    \tag{24.6}
\]
Apply this termwise with \(c=s_y-u_x^*\) to obtain (24.5).
\end{proof}

Let \(a_z=\|y_z\|^2\), let \(\chi_{xy}\geq0\) be the hopping
coefficient on \(\{x,y\}\), and put
\[
 h_x=\sum_{y\sim x}\chi_{xy},
 \qquad \chi_*=max_{y\sim x}\chi_{xy}.           \tag{24.7}
\]
After constants independent of \(u=s_x\) are cancelled, the
one-site conformal conditional action is
\[
\begin{aligned}
 \Phi_x(u)={}&\kappa_gu^2+b_gu+\frac J2W_x(u)\\
 &+\sum_{y\sim x}\chi_{xy}
 \bigl(a_xe^{s_y-u}+a_ye^{u-s_y}\bigr).           \tag{24.8}
\end{aligned}
\]
Indeed, \(J\cosh(2(s_y-u))=(J/2)V_g(s_y-u)\).  The
Higgs cross term on each edge is independent of \(u\) and cancels from
the conditional ratio.

\subsection{A uniform local conditional bound}

Use the combined local gravitational function
\[
 L_x(u)=\eta u^2+\theta W_x(u),                    \tag{24.9}
\]
where
\[
 0<\eta<\kappa_g,
 \qquad 0<\theta\leq\frac J2.                     \tag{24.10}
\]

\begin{theorem}[Neighbor-relative star estimate]
\label{thm:v17v24-star}
For every \(\varepsilon,\alpha,\tau>0\), there is a finite constant
\(C_x\), independent of all exterior conformal and Higgs values, such
that
\[
\begin{aligned}
 \log\mathbb E_x e^{L_x(s_x)}
 \leq{}&C_x+(2\kappa_g+\tau)(u_x^*)^2+K_M M_x\\
 &+e^{\varepsilon}\alpha
 \left(h_xV_H(y_x)
 +\sum_{y\sim x}\chi_{xy}V_H(y_y)\right),         \tag{24.11}
\end{aligned}
\]
where
\[
 K_M=\frac J2e^{2\varepsilon}
 +\frac{\chi_*e^{\varepsilon}}{4\alpha}.          \tag{24.12}
\]
The estimate depends on the star degree only through the displayed
finite sums and its additive constant.
\end{theorem}

\begin{proof}
The numerator of the conditional moment is
\[
 \int_{\mathbb R}
 \exp\left[-(\kappa_g-\eta)u^2-b_gu
 -\left(\frac J2-\theta\right)W_x(u)
 -S_{H\!g,x}(u)\right]du,                          \tag{24.13}
\]
where \(S_{H\!g,x}\geq0\) is the last line of (24.8).  Conditions
(24.10) allow both nonnegative interaction terms to be dropped, leaving
the finite Gaussian integral
\[
 \sqrt{\frac{\pi}{\kappa_g-\eta}}
 \exp\left(\frac{b_g^2}{4(\kappa_g-\eta)}\right). \tag{24.14}
\]

For the denominator, set \(u=u_x^*+r\) and use
\(|r|\leq\varepsilon\).  The on-site terms obey
\[
 \kappa_gu^2+b_gu
 \leq(2\kappa_g+\tau)(u_x^*)^2+C_0,               \tag{24.15}
\]
where one may take
\[
 C_0=2\kappa_g\varepsilon^2+|b_g|\varepsilon
 +\frac{b_g^2}{4\tau}.                            \tag{24.16}
\]
Lemma~\ref{lem:v17v24-center} bounds the mismatch part by
\((J/2)e^{2\varepsilon}M_x\).

Put \(c_y=s_y-u_x^*\).  On the anchor interval,
\[
\begin{aligned}
 a_xe^{s_y-u}+a_ye^{u-s_y}
 &\leq e^{\varepsilon}
 \bigl(a_xe^{c_y}+a_ye^{-c_y}\bigr)\\
 &\leq e^{\varepsilon}\left[
 \alpha(a_x^2+a_y^2)
 +\frac{V_g(c_y)}{4\alpha}
 \right].                                        \tag{24.17}
\end{aligned}
\]
After summing, the last terms in (24.17) are at most
\(\chi_*e^{\varepsilon}M_x/(4\alpha)\).  The denominator is bounded
below by \(2\varepsilon\) times the exponential of the negative of the
resulting upper bound for \(\Phi_x\).  Dividing (24.14) by this lower
bound proves (24.11)--(24.12).
\end{proof}

Using (24.4), the quadratic part of (24.11) transfers only to the
neighbor site variables.  The new object is \(M_x\): it depends solely
on the boundary of the conditional block, but it is not a sum of the
original edge functions after \(s_x\) has been integrated out.

\subsection{The same-sector absorption obstruction}

\begin{lemma}[Exact lower transfer]
\label{lem:v17v24-lower}
For every exterior configuration,
\[
 \log\mathbb E_x e^{L_x(s_x)}\geq\theta M_x.       \tag{24.18}
\]
\end{lemma}

\begin{proof}
By definition, \(W_x(u)\geq M_x\) for every \(u\), and
\(e^{\eta u^2}\geq1\).  Hence
\(e^{L_x(u)}\geq e^{\theta M_x}\) pointwise under the normalized
conditional measure.
\end{proof}

On a degree-two star with neighbors \(y,z\), the boundary functional
has a closed form.

\begin{proposition}[Half-exponent boundary chord]
\label{prop:v17v24-chord}
If \(N(x)=\{y,z\}\), then
\[
 u_x^*=\frac{s_y+s_z}{2},
 \qquad
 M_x=4\cosh(s_y-s_z).                              \tag{24.19}
\]
Equivalently, with
\[
 U_{yz}=e^{s_y-s_z}+e^{s_z-s_y}=2\cosh(s_y-s_z),  \tag{24.20}
\]
one has \(M_x=2U_{yz}\).
\end{proposition}

\begin{proof}
Let \(m=(s_y+s_z)/2\) and \(\Delta=s_y-s_z\).  The addition identity
for the hyperbolic cosine gives
\[
 W_x(u)=4\cosh(\Delta)\cosh(2(u-m)).               \tag{24.21}
\]
Its unique minimum is attained at \(u=m\), proving (24.19)--(24.20).
\end{proof}

The chord grows as \(e^{|s_y-s_z|}\), half the exponent of the
original edge variable \(V_g(s_y-s_x)\).  This loss of exponent is the
precise trace left by minimizing over the central site.

\begin{countertheorem}[No strict contraction by identifying
\(M_x\) with the original star sector]
\label{ctr:v17v24-self}
Suppose one tries to close the one-site conditional estimate using
only the original star variable with exponential parameter \(\theta\),
by replacing the boundary spread with a same-sector influence.  Any
bound whose remaining boundary terms grow sublinearly in \(M_x\) must
assign normalized influence at least one to that transfer.  In
particular, (24.11) assigns
\[
 \frac{K_M}{\theta}
 \geq\frac{(J/2)e^{2\varepsilon}}{\theta}
 \geq e^{2\varepsilon}>1.                         \tag{24.22}
\]
It therefore cannot pass a strict Perron contraction test.
\end{countertheorem}

\begin{proof}
Lemma~\ref{lem:v17v24-lower} forces the coefficient of \(M_x\) in any
such upper estimate to be at least \(\theta\).  To see that polynomial
boundary terms cannot hide this requirement, use the degree-two
configuration \(s_y=R\), \(s_z=-R\), with zero Higgs variables.  Then
\(u_x^*=0\), while (24.19) grows exponentially in \(R\).  Division by
the input parameter \(\theta\) gives influence at least one.  The
explicit estimate (24.22) follows from (24.10) and (24.12).
\end{proof}

This obstruction does not invalidate localized conditioning.  It says
that the half-exponent chord must be retained as a distinct boundary
sector with its own moment parameter; collapsing it immediately back
into the original star sector loses the strict margin.

\begin{assessment}[What V24 closes]
\label{ass:v17v24-status}
V24 replaces the nonuniform whole-volume anchor by a degree-local
conditional estimate.  Its coefficients are explicit and uniform on
bounded-degree graphs.  The calculation also proves that normalization
necessarily creates the boundary spread \(M_x\), and that same-sector
absorption of this spread has coefficient at least one.

For degree two, the new sector is exactly the half-exponent chord
\(U_{yz}\).  This identifies the next acquisition problem without
confusing failure of one Lyapunov closure with divergence of the model.
\end{assessment}

\begin{decision}[V25: audit and close the three-site chord]
\label{dec:v17v24-next}
Perform the mandatory companion-note audit.  Then analyze the
three-site path exactly, determine the admissible exponential parameter
for \(U_{yz}\), and test the two-level star-to-chord influence cycle.
If the cycle again saturates, move upstream to a joint path block rather
than iterating ever longer-distance chord sectors.
\end{decision}

\begin{firewall}[Claim boundary after V24]
\label{fire:v17v24}
The star estimate is a model conditional theorem.  It neither supplies
the physical gravitational action nor closes the graph Gibbs system;
the boundary-chord sector remains to be controlled.
\end{firewall}

\section*{Version V24 append-only record}
\addcontentsline{toc}{section}{Version V24 append-only record}
The complete V23 source was retained byte for byte before this
continuation.  It had \(237940\) bytes, \(7485\) lines, and SHA--256
\[
 \texttt{95F0A2849966CB8BC888785A25565A38F9E5789E49E268F9D9BBB5E379ADDF57}.
\]
V24 derived the local star bound and isolated its half-exponent boundary
sector.


\section{V25: companion audit and exact three-site chord margin}
\label{sec:v17v25}

\subsection{Mandatory companion-note audit}

The newest active Yang--Mills note inspected at this checkpoint was
\[
 \texttt{Renewal\_Yang\_Mills\_Mass\_Gap\_Research\_Notes\_V76.tex}.
                                                               \tag{25.1}
\]
It had \(630335\) bytes, \(17206\) physical lines, and SHA--256
\[
 \texttt{9C2F5579ADA93AA46C30EAFCA0CCE85AFFD1C5BD6EE8205206991A174BCBE9CF}.
                                                               \tag{25.2}
\]
Its new result proves both direct Wilson-response signs on paired
\(2t\times t\times t\) slabs in the sixth atlas.  It expressly does not
yet prove a complete sixth quarter cell, full transverse-torus
positivity, a cofinal lattice estimate, or a Yang--Mills continuum mass
gap.  The local compact-gauge conclusions used in SM V16 therefore
remain unchanged; no unfinished YM atlas statement is imported here.

The newest active Navier--Stokes note remained
\[
 \texttt{Renewal\_NS\_Stability\_Research\_Notes\_V26.tex}.     \tag{25.3}
\]
It had \(278283\) bytes, \(5319\) physical lines, and SHA--256
\[
 \texttt{82C887F8C2C69E4F28FAD7C3DC8DE5B9099CB5A4BF431EBA1FFF3E91935978AD}.
                                                               \tag{25.4}
\]
Its active block computes pointwise rank-one norms for first and second
Oseen-kernel derivatives; the first gives no improvement and the second
gives a modest local improvement.  These kernel estimates do not alter
the conformal Gibbs or Higgs influence rows studied here.

\begin{audit}[V25 cross-program decision]
\label{aud:v17v25-companions}
YM V76 and NS V26 are mathematically consistent with the present claim
boundary but supply no missing gravitational coefficient.  V25
therefore continues the V24 three-site chord calculation without
importing a companion theorem.
\end{audit}

\subsection{Orthogonal coordinates on the three-site path}

Consider the path \(y\sim x\sim z\).  Put
\[
 m=\frac{s_y+s_z}{2},\qquad
 \Delta=s_y-s_z,qquad
 r=s_x-m,qquad
 t=m+\frac r3.                                     \tag{25.5}
\]
Define the half-exponent chord and the central star by
\[
 U=e^{\Delta}+e^{-\Delta}=2\cosh\Delta,            \tag{25.6}
\]
\[
 W=V_g(s_y-s_x)+V_g(s_z-s_x).                      \tag{25.7}
\]

\begin{lemma}[Exact path identities]
\label{lem:v17v25-identities}
The change of variables in (25.5) gives
\[
 s_y^2+s_x^2+s_z^2
 =3t^2+\frac23r^2+\frac12\Delta^2,
 \qquad s_y+s_x+s_z=3t,                            \tag{25.8}
\]
and
\[
 W=2U\cosh(2r).                                    \tag{25.9}
\]
Consequently, the two-edge mismatch action is
\[
 J\sum_{v\in\{y,z\}}
 \bigl(\cosh(2(s_v-s_x))-1\bigr)
 =J\bigl(U\cosh(2r)-2\bigr).                      \tag{25.10}
\]
\end{lemma}

\begin{proof}
Substitute
\(s_y=m+\Delta/2\), \(s_z=m-\Delta/2\), and \(s_x=m+r\), followed
by \(m=t-r/3\), to obtain (25.8).  The addition identity gives
\[
\begin{aligned}
 W
 &=2\cosh(\Delta-2r)+2\cosh(\Delta+2r)\\
 &=4\cosh\Delta\cosh(2r)=2U\cosh(2r),             \tag{25.11}
\end{aligned}
\]
which proves the remaining claims.
\end{proof}

\subsection{Exact moment threshold for the chord}

Equip the path with the pure gravity density whose action is
\[
 \kappa_g(s_y^2+s_x^2+s_z^2)
 +b_g(s_y+s_x+s_z)
 +J\bigl(U\cosh(2r)-2\bigr),                       \tag{25.12}
\]
where \(\kappa_g>0\) and \(J>0\).

\begin{theorem}[Sharp half-exponent chord window]
\label{thm:v17v25-chord-window}
For \(\lambda\geq0\),
\[
 \mathbb E_{g,\mathrm{path}}e^{\lambda U}<\infty
 \quad\Longleftrightarrow\quad \lambda\leq J.    \tag{25.13}
\]
\end{theorem}

\begin{proof}
The common mode \(t\) factors as a Gaussian by (25.8).  In the relative
variables, \(\cosh(2r)\geq1\), so if \(\lambda\leq J\),
\[
 -JU\cosh(2r)+\lambda U
 \leq-(J-\lambda)U\leq0.                           \tag{25.14}
\]
The remaining Gaussian in \((r,\Delta)\) is integrable, including at
\(\lambda=J\).

If \(\lambda>J\), choose \(r_0>0\) so small that
\(J\cosh(2r)<\lambda\) for \(|r|<r_0\).  On this strip the exponent
contains a positive constant times
\(U=2\cosh\Delta\), which dominates the negative quadratic in
\(\Delta\).  The integral over the strip diverges.
\end{proof}

The V24 lower transfer is \(\theta M_x=2\theta U\).  The exact input
window (25.13) therefore leaves a strict normalized margin whenever
\(2\theta<J\).  To use that margin one must preserve the leading
coefficient \(2\theta\) in an upper conditional estimate; the fixed
anchor of V24 used the larger coefficient \(J\).

\subsection{An adaptive midpoint anchor}

Let the two hopping coefficients be \(\chi_y,\chi_z\geq0\), and put
\(\chi_*=\max(\chi_y,\chi_z)\).  Conditional on the endpoints and the
Higgs variables, the central-site hopping terms are
\[
\begin{aligned}
 H_x(r)={}&\chi_y
 \bigl(a_xe^{\Delta/2-r}+a_ye^{-\Delta/2+r}\bigr)\\
 &+\chi_z
 \bigl(a_xe^{-\Delta/2-r}+a_ze^{\Delta/2+r}\bigr),
                                                               \tag{25.15}
\end{aligned}
\]
where \(a_v=\|y_v\|^2\).  The Higgs cross terms again cancel from the
conditional ratio.

\begin{theorem}[Sharp-leading star-to-chord transfer]
\label{thm:v17v25-adaptive}
Assume
\[
 0<2\theta<J.                                      \tag{25.16}
\]
For every \(\alpha,\tau>0\), there is a finite constant \(C\),
independent of all boundary values, such that
\[
\begin{aligned}
 \log\mathbb E_x e^{\theta W}
 \leq{}&C+\frac{2\kappa_g+\tau}{2}(s_y^2+s_z^2)\\
 &+\left(2\theta+\frac{\mathrm e\,\chi_*}{2\alpha}\right)U\\
 &+\mathrm e\,\alpha\left[
 (\chi_y+\chi_z)V_H(y_x)
 +\chi_yV_H(y_y)+\chi_zV_H(y_z)
 \right].                                         \tag{25.17}
\end{aligned}
\]
Here \(\mathrm e\) denotes the base of the natural logarithm.
\end{theorem}

\begin{proof}
Set
\[
 \rho_U=(1+JU)^{-1/2}\leq1.                       \tag{25.18}
\]
The numerator, after insertion of \(e^{\theta W}\), contains the
interaction coefficient \((J-2\theta)U\cosh(2r)\).  Completing the
square gives
\(\kappa_g(m+r)^2+b_g(m+r)\geq-b_g^2/(4\kappa_g)\).
Dropping the nonnegative hopping and using
\(\cosh(2r)\geq1+2r^2\), the numerator is at most a constant times
\[
 e^{-(J-2\theta)U}
 \sqrt{\frac{\pi}{2(J-2\theta)U}}.                \tag{25.19}
\]

For the denominator, use \(|r|\leq\rho_U\).  The power series for the
cosh shows that, for \(|r|\leq1\),
\[
 \cosh(2r)-1\leq(\cosh2-1)r^2.                    \tag{25.20}
\]
Hence
\[
 JU\cosh(2r)\leq JU+(\cosh2-1).                   \tag{25.21}
\]
Also,
\[
 \kappa_g(m+r)^2+b_g(m+r)
 \leq(2\kappa_g+\tau)m^2+C_0.                    \tag{25.22}
\]
On this anchor, (25.15) and Young's inequality give
\[
\begin{aligned}
 H_x(r)\leq{}&\mathrm e\,\alpha
 \left[(\chi_y+\chi_z)a_x^2+\chi_ya_y^2+\chi_za_z^2\right]\\
 &+\frac{\mathrm e\,\chi_*}{2\alpha}U.           \tag{25.23}
\end{aligned}
\]
The anchor length is \(2\rho_U\).  Dividing (25.19) by the resulting
denominator lower bound produces the leading factor
\(e^{2\theta U}\).  The remaining square-root ratio is uniformly
bounded because
\[
 \frac{1+JU}{U}\leq J+\frac12,
 \qquad U\geq2.                                   \tag{25.24}
\]
Finally, \(m^2\leq(s_y^2+s_z^2)/2\).  Taking logarithms proves
(25.17).
\end{proof}

\begin{corollary}[Nonempty strict chord margin]
\label{cor:v17v25-margin}
Let \(\lambda\) be the exponential parameter assigned to the chord
sector \(U\).  There exist \(\alpha>0\) and \(\lambda\leq J\) for
which the normalized star-to-chord influence is strictly below one,
\[
 \frac{2\theta+\mathrm e\chi_* /(2\alpha)}{\lambda}<1,           \tag{25.25}
\]
if and only if \(2\theta<J\).  One may take any
\[
 \alpha>\frac{\mathrm e\chi_*}{2(J-2\theta)}       \tag{25.26}
\]
and then choose
\[
 2\theta+\frac{\mathrm e\chi_*}{2\alpha}
 <\lambda\leq J.                                  \tag{25.27}
\]
\end{corollary}

\begin{proof}
Under (25.16), (25.26) makes the left endpoint in (25.27) smaller than
\(J\), and Theorem~\ref{thm:v17v25-chord-window} permits every
\(\lambda\leq J\).  Conversely, (25.25) and \(\lambda\leq J\) imply
\(2\theta<J\).
\end{proof}

The cost of enlarging \(\alpha\) is visible in the Higgs row of
(25.17).  Thus V25 removes the gravity-only saturation but creates a
quantitative chord--Higgs tradeoff that must be tested jointly rather
than hidden in an additive constant.

\begin{assessment}[What V25 closes]
\label{ass:v17v25-status}
The mandatory YM/NS audit is complete.  On the three-site path, V25
derives the chord's exact exponential threshold and a boundary-adaptive
conditional estimate whose leading transfer is the sharp
\(2\theta U\), rather than the saturated \(JU\) from the fixed anchor.
The strict chord margin is nonempty precisely when \(2\theta<J\).

The full system is not yet closed because the same parameter
\(\alpha\) that improves the chord transfer increases the reverse
Higgs-quartic transfer.
\end{assessment}

\begin{decision}[V26: optimize the chord--Higgs tradeoff]
\label{dec:v17v25-next}
Form the three-sector comparison matrix for endpoint Higgs quartics,
the central star, and the half-exponent chord.  Eliminate \(\alpha\)
exactly where possible and determine whether a nonempty joint Perron
region survives.  If the chord has no local conditional closure on a
graph, replace the star and chord by one joint three-site gravity block.
\end{decision}

\begin{firewall}[Claim boundary after V25]
\label{fire:v17v25}
The exact chord window and conditional bound concern the declared
three-site conformal model.  They do not construct a graph-wide Gibbs
measure or identify the missing physical RPESM gravitational action.
\end{firewall}

\section*{Version V25 append-only record}
\addcontentsline{toc}{section}{Version V25 append-only record}
The complete V24 source was retained byte for byte before this
continuation.  It had \(247014\) bytes, \(7758\) lines, and SHA--256
\[
 \texttt{2BE9442132FBB459CC705E8A4562DED7C26ADF46179FF8367862DCB3EE43031E}.
\]
V25 completed the companion audit and proved the exact three-site chord
margin.


\section{V26: joint three-site gravity block and closed Higgs feedback}
\label{sec:v17v26}

\subsection{Why the block is enlarged}

V25 showed that one-site conditioning generates a half-exponent chord
and a second Young tradeoff.  On one path, both the star and the chord
are functions of the same three conformal variables.  This section
therefore conditions those variables jointly.  The chord becomes
internal to the gravity block, so only the genuine Higgs--gravity cycle
remains in the influence matrix.

Consider the path \(y\sim x\sim z\), with equal hopping coefficient
\(\chi>0\), and retain
\[
 W=V_g(s_y-s_x)+V_g(s_z-s_x).                      \tag{26.1}
\]
The gravity action is
\[
 S_{g,P}(s_y,s_x,s_z)
 =\kappa_g(s_y^2+s_x^2+s_z^2)
 +b_g(s_y+s_x+s_z)+\frac J2W-2J,                  \tag{26.2}
\]
where \(\kappa_g>0\) and \(J>0\).

\begin{theorem}[Exact joint-star moment window]
\label{thm:v17v26-star-window}
For \(\theta_G\geq0\),
\[
 \mathbb E_{g,P}e^{\theta_GW}<\infty
 \quad\Longleftrightarrow\quad
 \theta_G\leq\frac J2.                            \tag{26.3}
\]
\end{theorem}

\begin{proof}
If \(\theta_G\leq J/2\), insertion of the moment leaves the
nonnegative interaction coefficient \(J/2-\theta_G\).  Dropping that
interaction leaves an integrable three-dimensional Gaussian.

If \(\theta_G>J/2\), take a fixed-width tube about
\((s_y,s_x,s_z)=(R,0,0)\).  The term
\((\theta_G-J/2)V_g(s_y-s_x)\) grows like a positive multiple of
\(e^{2R}\); the other edge is bounded on the tube and the on-site action
is only quadratic.  Disjoint tubes with \(R\to\infty\) give an infinite
moment.
\end{proof}

This theorem controls both original edge variables at their natural
exponent.  The half-exponent chord of V25 needs no separate row when the
whole path is one block.

\subsection{Reverse conditional influence of the joint block}

Let \(a_v=\|y_v\|^2\).  Conditional on the three Higgs variables, add
the nonnegative diagonal hopping action
\[
\begin{aligned}
 S_{H\!g,P}(s\mid y)={}&\chi
 \bigl(a_xe^{s_y-s_x}+a_ye^{s_x-s_y}\bigr)\\
 &+\chi\bigl(a_xe^{s_z-s_x}+a_ze^{s_x-s_z}\bigr).
                                                               \tag{26.4}
\end{aligned}
\]
Both Higgs cross terms are independent of the conformal variables and
cancel from the conditional ratio.

\begin{theorem}[Joint-block reverse row]
\label{thm:v17v26-reverse}
Assume
\[
 0<\theta_G\leq\frac J2.                           \tag{26.5}
\]
For every \(\varepsilon,\zeta>0\), there is a finite constant \(C\),
independent of \(y_x,y_y,y_z\), such that
\[
\begin{aligned}
 \log\mathbb E_P\left[e^{\theta_GW}\mid y_x,y_y,y_z\right]
 \leq{}&C+\chi e^{2\varepsilon}\zeta
 \bigl(2V_H(y_x)+V_H(y_y)+V_H(y_z)\bigr).         \tag{26.6}
\end{aligned}
\]
Consequently, after maximum-envelope compression of the three Higgs
blocks, one may take
\[
 B_{GH}=c(\zeta)
 =\frac{4\chi e^{2\varepsilon}\zeta}{\theta_H},
 \qquad B_{GG}=0.                                  \tag{26.7}
\]
\end{theorem}

\begin{proof}
In the numerator, (26.5) leaves a nonnegative multiple of \(W\).
Drop it and (26.4).  Completing three squares gives a finite constant
depending only on \(\kappa_g,b_g,J\).

For the denominator, use the cube
\[
 |s_x|,|s_y|,|s_z|\leq\varepsilon.                 \tag{26.8}
\]
On this cube, the gravity action is bounded above by a parameter-only
constant and
\[
 S_{H\!g,P}(s\mid y)
 \leq\chi e^{2\varepsilon}(2a_x+a_y+a_z).          \tag{26.9}
\]
The cube has volume \((2\varepsilon)^3\), so the conditional log moment
is at most a finite constant plus the last term in (26.9).  Apply
\(q\leq\zeta q^2+(4\zeta)^{-1}\) to each \(a_v\).  The coefficients of
the three quartics sum to \(4\chi e^{2\varepsilon}\zeta\), which after
division by \(\theta_H\) gives (26.7).  Since the whole conformal path
is conditioned at once, there is no exterior gravity variable and the
self-entry is zero.
\end{proof}

\subsection{Forward Higgs row}

Apply the V21 anchor estimate to the two incident edges of the central
Higgs site.  Since the gravity block function \(W\) is their sum, the
maximum-envelope row over the three Higgs sites is
\[
 B_{HH}(\delta)=\frac{4\delta}{\theta_H},           \tag{26.10}
\]
\[
 B_{HG}(\delta)
 =\frac{\chi}{2\theta_G}
 +\frac{\chi^2}{4\theta_G\delta}.                 \tag{26.11}
\]
The endpoint rows have only one incident edge and are dominated by
(26.10)--(26.11).  Thus the comparison matrix is
\[
 \mathbb B_P(\delta,\zeta)
 =\begin{pmatrix}
 4\delta/\theta_H&
 \chi/(2\theta_G)+\chi^2/(4\theta_G\delta)\\
 4\chi e^{2\varepsilon}\zeta/\theta_H&0
 \end{pmatrix}.                                   \tag{26.12}
\]

\begin{theorem}[Exact optimized joint-path certificate]
\label{thm:v17v26-perron}
Let (26.5) hold.  Define
\[
 X=\frac{2\chi^2e^{2\varepsilon}\zeta}
 {\theta_H\theta_G},
 \qquad q=\frac{\chi}{2\theta_H}.                 \tag{26.13}
\]
There exists \(\delta>0\) for which
\(\rho(\mathbb B_P(\delta,\zeta))<1\) if and only if
\[
 1-X>4\sqrt{qX}.                                   \tag{26.14}
\]
Equivalently,
\[
\boxed{
 0<\zeta<
 \frac{\theta_H\theta_G}
 {2\chi^2e^{2\varepsilon}}
 \left(
 \sqrt{1+\frac{2\chi}{\theta_H}}
 -\sqrt{\frac{2\chi}{\theta_H}}
 \right)^2.}                                      \tag{26.15}
\]
The determinant margin is maximized at
\[
 \boxed{
 \delta_*=\frac12
 \sqrt{\frac{\chi^3e^{2\varepsilon}\zeta}{\theta_G}}.}       \tag{26.16}
\]
The interval (26.15) is nonempty for every fixed positive
\(\chi,\theta_H,\theta_G\).
\end{theorem}

\begin{proof}
Use Theorem~\ref{thm:v17v21-optimal} with
\[
 A=\frac4{\theta_H},\quad
 b_0=\frac\chi{2\theta_G},\quad
 B=\frac{\chi^2}{4\theta_G},\quad
 c=\frac{4\chi e^{2\varepsilon}\zeta}{\theta_H},quad
 d=0.                                              \tag{26.17}
\]
Condition (21.10) becomes (26.14).  Put \(u=\sqrt X\).  Then
\[
 u^2+4\sqrt q,u-1<0,                              \tag{26.18}
\]
so
\[
 u<\sqrt{1+4q}-2\sqrt q.                          \tag{26.19}
\]
Substitution of (26.13) gives (26.15).  Formula (21.11) gives
(26.16).
\end{proof}

\subsection{Comparison with the one-site chord scheme}

The joint block has removed the parameter \(\alpha\) from V25.  The
only reverse Young parameter is now \(\zeta\), and making it small
does not consume the gravity window \(\theta_G\leq J/2\).  The price is
geometric: different three-site paths overlap on a general graph, so
the zero self-entry in (26.7) is valid only when the displayed path is
conditioned as one isolated block.

\begin{assessment}[What V26 closes]
\label{ass:v17v26-status}
V26 follows the mandatory V25 audit with a substantive upstream step.
For the isolated three-site path, it proves the exact joint gravity
moment window, derives both directions of the Higgs--gravity influence
cycle, and solves the Perron test in closed form.  The passing interval
(26.15) is nonempty without requiring the physical hopping coefficient
to vanish.

The remaining obstruction is no longer analytic on one path.  It is
the overlap of such blocks on graph exhaustions and the boundary edges
created by any block partition.
\end{assessment}

\begin{decision}[V27: bounded-overlap path blocking]
\label{dec:v17v26-next}
Construct a finite-color family of vertex-disjoint paths on a
bounded-degree graph.  Compute the conditional influence carried by
edges crossing between path blocks and determine whether the color and
boundary-degree factors preserve a nonempty version of (26.15).  Do
not treat overlapping paths as independent blocks.
\end{decision}

\begin{firewall}[Claim boundary after V26]
\label{fire:v17v26}
The joint-path certificate is rigorous for one declared model block.
It does not yet give a volume-uniform block decomposition, the physical
RPESM gravity action, or the Einstein--SM continuum limit.
\end{firewall}

\section*{Version V26 append-only record}
\addcontentsline{toc}{section}{Version V26 append-only record}
The complete V25 source was retained byte for byte before this
continuation.  It had \(257083\) bytes, \(8055\) lines, and SHA--256
\[
 \texttt{7D737224D7B99DB42DD08B8770013BD47A181FF0A5C768D3E2593BE37903F145}.
\]
V26 replaced the saturated star--chord recursion by a closed joint-path
Higgs--gravity certificate.


\section{V27: bounded-overlap blocking and the exact path bridge}
\label{sec:v17v27}

\subsection{The combinatorial overlap is finite}

Let \(G\) be a finite simple graph of maximum degree \(D\geq2\), and
let \(G^4\) join two vertices whenever their graph distance in \(G\) is
at most four.

\begin{lemma}[Finite coloring by disjoint nonadjacent stars]
\label{lem:v17v27-color}
The vertices of \(G\) admit a coloring with at most
\[
 C_D=1+D\sum_{j=0}^{3}(D-1)^j                    \tag{27.1}
\]
colors such that, for centers \(x,x'\) of the same color, the closed
stars \(N[x]\) and \(N[x']\) are disjoint and no edge joins them.
\end{lemma}

\begin{proof}
A vertex has at most
\(D\sum_{j=0}^{3}(D-1)^j\) other vertices within distance four.
Thus the maximum degree of \(G^4\) is no larger than this number, and
the greedy coloring bound gives (27.1).  Same-color centers have
distance at least five in \(G\).  Their radius-one neighborhoods are
therefore disjoint and have mutual distance at least three, so no edge
joins them.
\end{proof}

Thus a bounded number of block-update colors suffices.  This fact does
not eliminate the conformal edges leaving each star; their exact
normalization cost is the analytic issue.

\subsection{A path with fixed boundary endpoints}

Let \(P_\ell=(v_0,\ldots,v_\ell)\) be a path of \(\ell\geq1\) edges.
Write
\[
 d_i=s_{v_i}-s_{v_{i-1}},\qquad
 \Delta=s_{v_\ell}-s_{v_0}=\sum_{i=1}^{\ell}d_i,   \tag{27.2}
\]
and define
\[
 W_\ell=\sum_{i=1}^{\ell}V_g(d_i)
 =2\sum_{i=1}^{\ell}\cosh(2d_i),                  \tag{27.3}
\]
\[
 U_\ell(\Delta)
 =e^{2\Delta/\ell}+e^{-2\Delta/\ell}
 =2\cosh\left(\frac{2\Delta}{\ell}\right).       \tag{27.4}
\]

\begin{theorem}[Exact bridge minimization]
\label{thm:v17v27-bridge}
For fixed endpoints, the internal path mismatch has the exact minimum
\[
 \inf_{\substack{d_1+\cdots+d_\ell=\Delta}}
 W_\ell
 =\ell U_\ell(\Delta).                             \tag{27.5}
\]
The unique minimizing increments are
\[
 d_1=\cdots=d_\ell=\frac\Delta\ell.                \tag{27.6}
\]
\end{theorem}

\begin{proof}
The function \(q\mapsto\cosh(2q)\) is strictly convex.  Jensen's
inequality gives
\[
 \frac1\ell\sum_{i=1}^{\ell}\cosh(2d_i)
 \geq\cosh\left(\frac2\ell\sum_{i=1}^{\ell}d_i\right)
 =\cosh\left(\frac{2\Delta}{\ell}\right).         \tag{27.7}
\]
Multiplication by \(2\ell\) proves (27.5), and strict convexity gives
the equality condition (27.6).
\end{proof}

Conditioning the \(\ell-1\) internal vertices therefore generates the
boundary chord \(U_\ell\).  Its exponential rate is \(2/\ell\), rather
than the rate two of an original edge.  V24--V25 are the case
\(\ell=2\).

\subsection{The induced chord has an exact matching moment window}

On all \(\ell+1\) path variables, consider the pure gravity action
\[
 S_{g,P_\ell}(s)
 =\kappa_g\sum_{i=0}^{\ell}s_{v_i}^2
 +b_g\sum_{i=0}^{\ell}s_{v_i}
 +\frac J2W_\ell-J\ell,                            \tag{27.8}
\]
with \(\kappa_g>0\) and \(J>0\).

\begin{theorem}[Sharp length-\(\ell\) chord threshold]
\label{thm:v17v27-chord-window}
For every \(\lambda\geq0\),
\[
 \mathbb E_{g,P_\ell}e^{\lambda U_\ell}<\infty
 \quad\Longleftrightarrow\quad
 \lambda\leq\frac{J\ell}{2}.                     \tag{27.9}
\]
\end{theorem}

\begin{proof}
Theorem~\ref{thm:v17v27-bridge} gives
\[
 \frac J2W_\ell\geq\frac{J\ell}{2}U_\ell.        \tag{27.10}
\]
If \(\lambda\leq J\ell/2\), insertion of
\(e^{\lambda U_\ell}\) leaves a nonnegative mismatch action.  Dropping
it leaves an integrable \((\ell+1)\)-dimensional Gaussian.

For the converse, use the ray
\[
 s_{v_i}=\frac{iR}{\ell},\qquad0\leq i\leq\ell.    \tag{27.11}
\]
Along this ray, equality holds in (27.10).  If
\(\lambda>J\ell/2\), the logarithm of the moment integrand contains a
positive multiple of \(e^{2R/\ell}\), which dominates the quadratic
on-site cost.  By continuity, the coefficient remains positive on
sufficiently narrow fixed-width tubes about the ray.  A disjoint
sequence of such tubes proves divergence.
\end{proof}

\begin{corollary}[Length-independent normalized acquisition margin]
\label{cor:v17v27-margin}
For the conditional internal-block moment
\(e^{\theta W_\ell}\), the pointwise bound
\[
 W_\ell\geq\ell U_\ell                             \tag{27.12}
\]
forces boundary transfer at least \(\theta\ell U_\ell\).  Relative to
the largest admissible chord parameter in (27.9), the forced normalized
coefficient is
\[
 \frac{\theta\ell}{J\ell/2}=\frac{2\theta}{J}.     \tag{27.13}
\]
It is strictly below one exactly when \(2\theta<J\), independently of
the path length.
\end{corollary}

\begin{proof}
For every internal configuration,
\(e^{\theta W_\ell}\geq e^{\theta\ell U_\ell}\), so any boundary
upper estimate with lower-order remainder must carry at least the
coefficient \(\theta\ell\).  Divide by the exact endpoint
\(J\ell/2\) from Theorem~\ref{thm:v17v27-chord-window}.
\end{proof}

This calculation rules out one feared loss: longer block boundaries do
not by themselves drive the normalized conformal influence to one.
The exponent becomes softer, but its admissible moment parameter grows
by the compensating factor \(\ell\).

\subsection{What finite coloring does and does not solve}

For a path graph, Lemma~\ref{lem:v17v27-color} permits a finite-color
schedule of disjoint three-site blocks.  Each block update generates a
chord spanning its two exterior boundary sites.  Repeating the blocking
creates longer chords of the form (27.4).  Theorem
\ref{thm:v17v27-chord-window} supplies an exact admissible parameter for
each such scale, and Corollary~\ref{cor:v17v27-margin} shows that the
gravity-only strict margin can remain fixed.

On a graph with branching, a star has several boundary endpoints and
the scalar constraint in (27.2) is replaced by a network Dirichlet
problem.  Jensen along a single path still gives valid pairwise lower
bounds, but summing them can overcount internal edges.  A graph-wide
construction therefore needs a congestion-controlled family of paths
or the exact convex network minimum.

\begin{assessment}[What V27 closes]
\label{ass:v17v27-status}
V27 proves that bounded-degree star overlap has a finite coloring and
computes exactly the boundary variable created by every finite path
block.  Its moment window scales in exact proportion to path length,
so the normalized acquisition margin \(2\theta/J\) does not deteriorate
under longer path blocking.

The outstanding graph issue is path congestion at branching vertices,
not the one-dimensional conformal tail.  Higgs transfers across block
boundaries also remain to be included at each scale.
\end{assessment}

\begin{decision}[V28: convex network bridge and congestion]
\label{dec:v17v27-next}
For a finite connected block with prescribed boundary values, study the
strictly convex minimum of the sum of edge coshes.  Derive lower bounds
from edge-disjoint or fractionally weighted paths, with coefficients
that do not overcount internal edges.  Use those coefficients to define
the multiboundary chord sector before reintroducing Higgs hopping.
\end{decision}

\begin{firewall}[Claim boundary after V27]
\label{fire:v17v27}
Finite coloring and the exact path threshold do not yet constitute a
graph Gibbs construction.  No branching-network Perron matrix,
physical gravitational action, or Einstein--SM continuum limit is
claimed.
\end{firewall}

\section*{Version V27 append-only record}
\addcontentsline{toc}{section}{Version V27 append-only record}
The complete V26 source was retained byte for byte before this
continuation.  It had \(265307\) bytes, \(8299\) lines, and SHA--256
\[
 \texttt{34C5DD88DA0430F12EE0AB5BD29B389F35D982765374F919A066450F51BBFF20}.
\]
V27 established finite star coloring and the exact multiscale path
bridge margin.


\section{V28: convex network bridges and congestion-free boundary sectors}
\label{sec:v17v28}

\subsection{The exact network boundary functional}

Let \(K=(\mathcal V_K,\mathcal E_K)\) be a finite connected simple
graph.  Split its vertices into a nonempty boundary \(\partial K\) and
an interior \(K^\circ\).  For a boundary value
\(\sigma\in\mathbb R^{\partial K}\), define
\[
 \mathcal W_K(\sigma,u)
 =\sum_{\{v,w\}\in\mathcal E_K}V_g(s_v-s_w),       \tag{28.1}
\]
where \(s|_{\partial K}=\sigma\) and
\(s|_{K^\circ}=u\), and put
\[
 \mathcal M_K(\sigma)
 =\inf_{u\in\mathbb R^{K^\circ}}
 \mathcal W_K(\sigma,u).                          \tag{28.2}
\]

\begin{theorem}[Exact convex Dirichlet bridge]
\label{thm:v17v28-dirichlet}
If \(K^\circ\ne\varnothing\), the infimum in (28.2) is attained at a
unique field \(u^\sigma\).  At every interior vertex it satisfies
\[
 \sum_{w\sim v}\sinh\bigl(2(u_v^\sigma-s_w^\sigma)\bigr)=0.
                                                               \tag{28.3}
\]
The functional \(\mathcal M_K\) is finite, continuous, convex, and
invariant under a common shift of all boundary values.
\end{theorem}

\begin{proof}
Each summand \(V_g(q)=e^{2q}+e^{-2q}\) is strictly convex.  The Hessian
in the interior variables is the grounded weighted graph Laplacian
\[
 B_\circ^{\mathsf T}
 \operatorname{diag}\bigl(4V_g((Bs)_e)\bigr)B_\circ,             \tag{28.4}
\]
where \(B_\circ\) is the incidence matrix restricted to interior
columns.  Connectedness and the nonempty fixed boundary make
\(B_\circ\) injective, so (28.4) is positive definite.  Coercivity in
the interior variables gives existence, and strict convexity gives
uniqueness.  Differentiation gives (28.3).

Joint convexity of \(\mathcal W_K\) and partial minimization give
convexity of \(\mathcal M_K\).  Continuity follows from finite
dimensional coercive convex minimization; equivalently, it follows
locally from the implicit function theorem and (28.4).  Since every
edge difference is unchanged by a common shift, the final assertion is
immediate.
\end{proof}

For a path with two boundary endpoints, Theorem
\ref{thm:v17v27-bridge} says exactly that
\[
 \mathcal M_{P_\ell}=\ell U_\ell.                 \tag{28.5}
\]
Thus \(\mathcal M_K\) is the branching-network replacement for the
family of induced chords.

\subsection{Fractional paths without edge overcounting}

Choose boundary pairs \((a_j,b_j)\), paths \(P_j\) joining them, path
lengths \(\ell_j\), and weights \(w_j\geq0\).  Their edge congestion is
\[
 \mathfrak C
 =\max_{e\in\mathcal E_K}
 \sum_{j:e\in P_j}w_j.                             \tag{28.6}
\]
For \(\mathfrak C>0\), define the computable boundary functional
\[
 \mathcal F_{\mathcal P}(\sigma)
 =\sum_jw_j\ell_j
 \,2\cosh\left(
 \frac{2(\sigma_{b_j}-\sigma_{a_j})}{\ell_j}
 \right).                                         \tag{28.7}
\]

\begin{theorem}[Congestion-controlled path lower bound]
\label{thm:v17v28-congestion}
For every boundary datum,
\[
 \boxed{
 \mathcal M_K(\sigma)
 \geq\frac{1}{\mathfrak C}
 \mathcal F_{\mathcal P}(\sigma).}                \tag{28.8}
\]
In particular, fractionally edge-disjoint paths
\((\mathfrak C\leq1)\) may be summed without weakening their individual
path bounds.
\end{theorem}

\begin{proof}
For any extension \(s\) of \(\sigma\), Theorem
\ref{thm:v17v27-bridge} applied on each path gives
\[
 w_j\sum_{e\in P_j}V_g((Bs)_e)
 \geq w_j\ell_j
 \,2\cosh\left(
 \frac{2(\sigma_{b_j}-\sigma_{a_j})}{\ell_j}
 \right).                                         \tag{28.9}
\]
After summing in \(j\), the left side is at most
\(\mathfrak C\mathcal W_K(s)\).  Divide by
\(\mathfrak C\) and minimize over the interior extension.
\end{proof}

The weights in (28.6) can be selected by a finite multicommodity-flow
linear program.  The resulting certificate is weaker than the exact
convex minimum but has fully auditable rational path and congestion
data.

\subsection{Exact moment window of the network minimum}

First suppose that \(\partial K=\{a,b\}\) has two vertices.  On all
vertices of \(K\), use the pure gravity action
\[
 S_{g,K}(s)
 =\kappa_g\sum_{v\in\mathcal V_K}s_v^2
 +b_g\sum_{v\in\mathcal V_K}s_v
 +\frac J2\mathcal W_K(s),                         \tag{28.10}
\]
where \(\kappa_g>0\) and \(J>0\).  Additive edge constants have been
suppressed because they do not affect integrability.

\begin{theorem}[Sharp network-boundary moment threshold]
\label{thm:v17v28-network-window}
For the two-boundary connected network and \(\Lambda\geq0\),
\[
 \mathbb E_{g,K}
 e^{\Lambda\mathcal M_K(s_a,s_b)}<\infty
 \quad\Longleftrightarrow\quad
 \Lambda\leq\frac J2.                             \tag{28.11}
\]
For an arbitrary nonempty boundary, the forward implication
\[
 0\leq\Lambda\leq J/2
 \quad\Longrightarrow\quad
 \mathbb E_{g,K}e^{\Lambda\mathcal M_K}<\infty     \tag{28.12}
\]
remains valid.
\end{theorem}

\begin{proof}
For every full field,
\[
 \mathcal M_K(s|_{\partial K})\leq\mathcal W_K(s). \tag{28.13}
\]
Thus, when \(\Lambda\leq J/2\), the insertion leaves a nonnegative
multiple of \(\mathcal W_K\) in the action.  Dropping it leaves a finite
Gaussian integral, proving (28.12) and sufficiency in (28.11).

For necessity with two boundary vertices, set \(s_a=0\), \(s_b=R\)
and put the interior field at its unique minimizer.  Any fixed simple
\(a\)-to-\(b\) path of length \(L\) and (28.8) give
\[
 \mathcal M_K(0,R)
 \geq2L\cosh(2R/L),                                \tag{28.14}
\]
so the minimum grows exponentially in \(R\).  On the minimizing field,
\(\mathcal W_K=\mathcal M_K\).

It remains to turn this ray into positive-volume sets.  Perturb every
vertex by at most \(\delta\).  Each edge difference changes by at most
\(2\delta\), so
\[
 \mathcal W_K(s+z)
 \leq e^{4\delta}\mathcal W_K(s).                 \tag{28.15}
\]
The two-boundary minimum depends only on the boundary difference.
Changing that difference by at most \(2\delta\) changes its value by at
most a multiplicative factor tending to one with \(\delta\); this also
follows by translating a minimizer with any fixed unit boundary
interpolant and applying (28.15).  If \(\Lambda>J/2\), choose
\(\delta>0\) so small that the inserted minimum coefficient still
exceeds the action coefficient throughout these boxes.  Their
log-integrand then grows like a positive multiple of (28.14), while the
on-site contribution is only \(O(R^2)\).  Indeed, (28.3) has the
discrete maximum principle: an interior value strictly above both
boundary values would give a positive sum in (28.3) at an interior
maximum, and an interior value strictly below them gives the analogous
negative contradiction.  Hence the minimizing field lies in
\([0,R]\).  A disjoint sequence of boxes proves divergence.
\end{proof}

\subsection{The structural block margin}

Condition the interior of \(K\) with its boundary fixed.  For every
\(\theta>0\), the pointwise inequality (28.13), now read as
\(\mathcal W_K\geq\mathcal M_K\), gives
\[
 \log\mathbb E_{K^\circ}
 e^{\theta\mathcal W_K}
 \geq\theta\mathcal M_K.                          \tag{28.16}
\]
If the network-minimum sector is assigned parameter \(\Lambda\), the
forced normalized boundary influence is at least
\(\theta/\Lambda\).  Theorem
\ref{thm:v17v28-network-window} permits
\[
 \theta<\Lambda\leq\frac J2                       \tag{28.17}
\]
exactly when
\[
 2\theta<J.                                        \tag{28.18}
\]
The same strict margin found for paths therefore survives arbitrary
finite branching when the boundary sector is the exact convex network
minimum rather than a sum of overcounted pairwise chords.

The lower bound (28.16) identifies the best possible leading
coefficient.  A useful conditional theorem must still prove a matching
upper estimate with coefficient arbitrarily close to \(\theta\), as V25
did on the degree-two star.

\begin{assessment}[What V28 closes]
\label{ass:v17v28-status}
V28 replaces the proliferating chord list by one canonical convex
Dirichlet functional \(\mathcal M_K\).  It proves uniqueness of its
minimizer, supplies congestion-controlled computable lower bounds, and
derives its sharp two-boundary exponential window.  Branching does not
consume the structural strict margin \(2\theta<J\).

What remains is an upper conditional Laplace-ratio estimate around the
network minimizer, uniform over boundary data and bounded-size block
shapes, together with the Higgs hopping debit on that adaptive tube.
\end{assessment}

\begin{decision}[V29: adaptive tube around the network minimizer]
\label{dec:v17v28-next}
Use the positive grounded Hessian (28.4) to choose a boundary-dependent
ellipsoidal normalization tube.  Compare numerator and denominator
determinants so that their boundary-volume factors cancel, retain the
sharp leading transfer \(\theta\mathcal M_K\), and then add the Higgs
hopping terms with explicit block-size and degree constants.
\end{decision}

\begin{firewall}[Claim boundary after V28]
\label{fire:v17v28}
The network functional and its moment window do not yet give the needed
uniform conditional upper bound.  The RPESM gravitational action and
the Einstein--SM continuum construction remain unresolved.
\end{firewall}

\section*{Version V28 append-only record}
\addcontentsline{toc}{section}{Version V28 append-only record}
The complete V27 source was retained byte for byte before this
continuation.  It had \(273056\) bytes, \(8512\) lines, and SHA--256
\[
 \texttt{5AC1CE6FCF09649E3FFBB7371ECA0DD93E6966A19BDC7ED3CEB1CEEC761667DE}.
\]
V28 constructed the convex network boundary sector and removed path
overcounting by explicit congestion.


\section{V29: a sharp-leading conditional bound for finite networks}
\label{sec:v17v29}

\subsection{Retaining the numerator minimum}

The V28 decision proposed an ellipsoidal Hessian comparison.  A simpler
argument gives the needed leading coefficient.  The essential point is
to retain
\(\mathcal W_K\geq\mathcal M_K\) in the numerator rather than dropping
the entire mismatch action.  A fixed cube about the convex minimizer
then suffices.

Let \(K\) be as in V28, put
\[
 k=|K^\circ|,
 \qquad A=\frac J2,                                \tag{29.1}
\]
and let \(s^\sigma\) denote the extension of the boundary datum by the
unique minimizer from Theorem~\ref{thm:v17v28-dirichlet}.  Write
\[
 q_e=(Bs^\sigma)_e,
 \qquad \mathcal M=\mathcal M_K(\sigma).           \tag{29.2}
\]
Let \(\mathcal E_K^\circ\) be the edges with at least one interior
endpoint.  Give them hopping weights \(0\leq\chi_e\leq\chi_*\), and
define
\[
 h_v^\circ=\sum_{\substack{e\in\mathcal E_K^\circ\\e\ni v}}\chi_e.
                                                               \tag{29.3}
\]
The nonconstant diagonal Higgs hopping is
\[
 H_K(s\mid y)
 =\sum_{e=(v,w)\in\mathcal E_K^\circ}\chi_e
 \bigl(a_ve^{s_w-s_v}+a_we^{s_v-s_w}\bigr),       \tag{29.4}
\]
where \(a_v=\|y_v\|^2\).  Its Higgs cross terms are independent of the
conformal variables and cancel from the conditional ratio.

\subsection{Uniform tube estimates}

\begin{lemma}[Multiplicative control about the minimizer]
\label{lem:v17v29-tube}
If \(u=s^\sigma|_{K^\circ}+v\) with
\(\|v\|_\infty\leq r\), then
\[
 \mathcal W_K(\sigma,u)
 \leq e^{4r}\mathcal M.                           \tag{29.5}
\]
Moreover, for every \(\alpha>0\),
\[
 H_K(s\mid y)
 \leq e^{2r}\alpha\sum_{v\in\mathcal V_K}
 h_v^\circ V_H(y_v)
 +\frac{e^{2r}\chi_*}{4\alpha}\mathcal M.         \tag{29.6}
\]
\end{lemma}

\begin{proof}
Across an edge with one or two perturbed endpoints, the difference
changes by at most \(2r\).  Hence
\[
 V_g(q_e+(Bv)_e)\leq e^{4r}V_g(q_e).               \tag{29.7}
\]
Summation gives (29.5).

For an oriented edge, the corresponding summand in (29.4) is at most
\[
 e^{2r}\chi_e
 \bigl(a_ve^{q_e}+a_we^{-q_e}\bigr).               \tag{29.8}
\]
Young's inequality gives
\[
 a_ve^{q_e}+a_we^{-q_e}
 \leq\alpha(a_v^2+a_w^2)+\frac{V_g(q_e)}{4\alpha}.
                                                               \tag{29.9}
\]
Sum (29.9), use \(\chi_e\leq\chi_*\), and note that the sum of
\(V_g(q_e)\) over \(\mathcal E_K^\circ\) is at most \(\mathcal M\).
\end{proof}

\begin{lemma}[Boundary control of the minimizing field]
\label{lem:v17v29-maximum}
Let
\[
 Q_{\partial K}(\sigma)=\sum_{b\in\partial K}\sigma_b^2.       \tag{29.10}
\]
Then
\[
 \sum_{v\in K^\circ}(s_v^\sigma)^2
 \leq kQ_{\partial K}(\sigma).                    \tag{29.11}
\]
\end{lemma}

\begin{proof}
Equation (28.3) has the discrete maximum principle.  At an interior
maximum strictly above all boundary values every summand in (28.3) is
nonnegative and at least one is positive; the minimum case is
analogous.  Thus every interior value lies between the smallest and
largest boundary values.  Its square is at most
\(\max_{b\in\partial K}\sigma_b^2\), which is bounded by (29.10).
Sum over the \(k\) interior vertices.
\end{proof}

\subsection{The conditional Laplace-ratio theorem}

Condition the interior conformal variables with action
\[
 \Phi_K(u\mid\sigma,y)
 =\kappa_g\sum_{v\in K^\circ}u_v^2
 +b_g\sum_{v\in K^\circ}u_v
 +A\mathcal W_K(\sigma,u)+H_K(s\mid y),            \tag{29.12}
\]
where \(\kappa_g>0\), \(J>0\), and constants independent of \(u\)
have been cancelled.

\begin{theorem}[Sharp-leading network conditional estimate]
\label{thm:v17v29-conditional}
Assume
\[
 0<\theta<A=\frac J2.                              \tag{29.13}
\]
For every \(r,\alpha,\tau>0\), there is a finite constant \(C_K\),
independent of \(\sigma\) and \(y\), such that
\[
\begin{aligned}
 \log\mathbb E_{K^\circ}
 e^{\theta\mathcal W_K}
 \leq{}&C_K+(2\kappa_g+\tau)kQ_{\partial K}(\sigma)\\
 &+q_M(r,\alpha)\mathcal M_K(\sigma)\\
 &+e^{2r}\alpha\sum_{v\in\mathcal V_K}
 h_v^\circ V_H(y_v),                              \tag{29.14}
\end{aligned}
\]
where
\[
 \boxed{
 q_M(r,\alpha)
 =\theta+A(e^{4r}-1)
 +\frac{e^{2r}\chi_*}{4\alpha}.}                  \tag{29.15}
\]
\end{theorem}

\begin{proof}
The conditional-moment numerator contains the mismatch coefficient
\(A-\theta>0\).  By the defining minimum,
\(\mathcal W_K(\sigma,u)\geq\mathcal M\), and the hopping is
nonnegative.  Therefore the numerator is at most
\[
 e^{-(A-\theta)\mathcal M}
 \left(\frac{\pi}{\kappa_g}\right)^{k/2}
 \exp\left(\frac{kb_g^2}{4\kappa_g}\right).       \tag{29.16}
\]

For the denominator, use the cube
\(\|u-s^\sigma|_{K^\circ}\|_\infty\leq r\), whose volume is
\((2r)^k\).  For any \(\tau>0\),
\[
\begin{aligned}
 \kappa_g\sum_{v\in K^\circ}u_v^2
 +b_g\sum_{v\in K^\circ}u_v
 \leq{}&(2\kappa_g+\tau)
 \sum_{v\in K^\circ}(s_v^\sigma)^2+C_0,          \tag{29.17}
\end{aligned}
\]
with \(C_0<\infty\) depending only on
\(k,\kappa_g,b_g,r,\tau\).  Lemmas
\ref{lem:v17v29-tube}--\ref{lem:v17v29-maximum} then give a denominator
lower bound whose boundary-dependent exponent is at most
\[
 (2\kappa_g+\tau)kQ_{\partial K}
 +\left(Ae^{4r}+\frac{e^{2r}\chi_*}{4\alpha}\right)\mathcal M
 +e^{2r}\alpha\sum_vh_v^\circ V_H(y_v).           \tag{29.18}
\]
Divide (29.16) by this lower bound.  The coefficient of
\(\mathcal M\) is
\[
 Ae^{4r}-(A-\theta)
 +\frac{e^{2r}\chi_*}{4\alpha},                   \tag{29.19}
\]
which is exactly (29.15).
\end{proof}

\subsection{A strict network-sector transfer}

Assign exponential parameter \(\Lambda\) to the boundary sector
\(\mathcal M_K\).  V28 permits \(\Lambda\leq A\) in the two-boundary
network, and gives the same sufficient window for any boundary.

\begin{corollary}[Nonempty uniform leading margin]
\label{cor:v17v29-margin}
Let
\[
 0<\theta<\Lambda<A.                               \tag{29.20}
\]
There exist \(r,\alpha>0\) such that
\[
 q_M(r,\alpha)<\Lambda.                            \tag{29.21}
\]
More explicitly, choose
\[
 0<r<\frac14\log\left(1+\frac{\Lambda-\theta}{A}\right)       \tag{29.22}
\]
and then
\[
 \alpha>
 \frac{e^{2r}\chi_*}
 {4\left[\Lambda-\theta-A(e^{4r}-1)\right]}.       \tag{29.23}
\]
The normalized network-boundary influence
\(q_M/\Lambda\) is then strictly below one.
\end{corollary}

\begin{proof}
Condition (29.22) makes the square bracket in (29.23) positive.
Substitution in (29.15) proves (29.21).
\end{proof}

For a family of block shapes with
\[
 k\leq k_*,\qquad \deg v\leq D,qquad
 \chi_e\leq\chi_*,                                \tag{29.24}
\]
the same \(r,\alpha\) and the same displayed influence coefficients
work for every block.  Additive constants are uniformly bounded because
there are only finitely many simple graph shapes with at most
\(k_*+|\partial K|\) vertices once the boundary size is also bounded.

\subsection{The remaining two feedback debits}

The theorem isolates two transfers that were absent from the pure
network margin:
\[
 (2\kappa_g+\tau)kQ_{\partial K}                  \tag{29.25}
\]
to the quadratic conformal site sector, and
\[
 e^{2r}\alpha h_v^\circ V_H(y_v)                 \tag{29.26}
\]
to the Higgs sector.  Making \(r\) small improves both the network
coefficient and the exponential prefactor.  Making \(\alpha\) large
improves the network coefficient but worsens (29.26), reproducing the
genuine Higgs tradeoff found in V25.

Unlike V25, (29.14) is valid for every finite connected block and every
boundary datum.  The tradeoff can now be entered into a finite
three-sector Perron matrix with block-size and degree constants.

\begin{assessment}[What V29 closes]
\label{ass:v17v29-status}
V29 proves the missing upper conditional estimate around the exact
network minimizer.  Its leading boundary coefficient can be made
arbitrarily close to the forced value \(\theta\), so the strict window
\(\theta<\Lambda<J/2\) survives branching, block boundaries, and
nonnegative Higgs hopping on every bounded-size block family.

The full comparison still needs the quadratic boundary row and the
opposing \(\alpha\)-dependence of the Higgs feedback.  These are now
explicit finite coefficients rather than unclassified normalization
costs.
\end{assessment}

\begin{decision}[V30: audit and three-sector Perron optimization]
\label{dec:v17v29-next}
At the next mandatory companion checkpoint, inspect the active YM and
NS notes.  Then assemble the Higgs, quadratic conformal, and network
minimum rows from V21 and (29.14), optimize the \(\alpha\) tradeoff, and
state the exact or best explicit spectral feasibility condition.
Continue to V31 after the audit before treating the checkpoint as
complete.
\end{decision}

\begin{firewall}[Claim boundary after V29]
\label{fire:v17v29}
The conditional theorem concerns the declared convex conformal model
on bounded finite blocks.  It does not identify the RPESM gravitational
action, close an infinite-volume block dynamics, or prove the
Einstein--SM continuum limit.
\end{firewall}

\section*{Version V29 append-only record}
\addcontentsline{toc}{section}{Version V29 append-only record}
The complete V28 source was retained byte for byte before this
continuation.  It had \(282669\) bytes, \(8769\) lines, and SHA--256
\[
 \texttt{D2CCB2CBFE592586E8A1A8CA6E79D61229CA4CF0C65DCD6BAB834D905263493D}.
\]
V29 proved the sharp-leading finite-network conditional Laplace bound.


\section{V30: companion audit and exact three-sector Schur optimization}
\label{sec:v17v30}

\subsection{Mandatory companion-note audit}

The active Yang--Mills note at this checkpoint was
\[
 \texttt{Renewal\_Yang\_Mills\_Mass\_Gap\_Research\_Notes\_V79.tex}.
                                                               \tag{30.1}
\]
It had \(643922\) bytes, \(17592\) physical lines, and SHA--256
\[
 \texttt{F68AEEEB02910AD4538E270EF9A742505C340BB5C46A083F4684A5F7E1A99CE6}.
                                                               \tag{30.2}
\]
Since the V25 audit, YM V77--V79 have closed the sixth first and second
quarter faces and, by inversion, proved both direct response signs on
the paired faces.  The sixth third coordinate, the complete
fifth/sixth cell pair, full-torus positivity, cofinal control, the
continuum measure, clustering, and the Hamiltonian mass gap remain
open.  These new local atlas certificates do not change the compact
gauge-jet bounds used in SM V16 and do not supply a gravitational row.

The active Navier--Stokes note remained
\[
 \texttt{Renewal\_NS\_Stability\_Research\_Notes\_V26.tex},     \tag{30.3}
\]
with \(278283\) bytes, \(5319\) physical lines, and SHA--256
\[
 \texttt{82C887F8C2C69E4F28FAD7C3DC8DE5B9099CB5A4BF431EBA1FFF3E91935978AD}.
                                                               \tag{30.4}
\]
Its rank-one Oseen-kernel calculation remains analytically separate
from the present conformal block problem.

\begin{audit}[V30 cross-program decision]
\label{aud:v17v30-companions}
No theorem is imported from YM V79 or NS V26.  The SM programme retains
their accurate claim boundaries and continues with the explicit
three-sector influence optimization forced by V29.
\end{audit}

\subsection{A structured three-sector matrix}

Let the first two sectors be the Higgs quartic sector \(H\) and the
quadratic conformal boundary sector \(Q\).  Let the third sector \(G\)
be the convex network mismatch block.  Consider
\[
 \mathbb B(\alpha)
 =\begin{pmatrix}
 a_{HH}&a_{HQ}&u_H\\
 a_{QH}&a_{QQ}&u_Q\\
 v_H^0+g\alpha&v_Q&m_0+r/\alpha
 \end{pmatrix},
 \qquad \alpha>0,                                  \tag{30.5}
\]
where the blank \(2\times2\) corner is the nonnegative matrix
\[
 \mathbb A=
 \begin{pmatrix}a_{HH}&a_{HQ}\\a_{QH}&a_{QQ}\end{pmatrix},    \tag{30.6}
\]
and every displayed scalar is nonnegative.  In block notation,
\[
 \mathbb B(\alpha)
 =\begin{pmatrix}
 \mathbb A&u\\v(\alpha)^{\mathsf T}&m(\alpha)
 \end{pmatrix},                                   \tag{30.7}
\]
with
\[
 u=\binom{u_H}{u_Q},\quad
 v(\alpha)=\binom{v_H^0+g\alpha}{v_Q},\quad
 m(\alpha)=m_0+\frac r\alpha.                     \tag{30.8}
\]

\begin{theorem}[Exact Schur criterion for three nonnegative sectors]
\label{thm:v17v30-schur}
Assume \(\rho(\mathbb A)<1\) and define
\[
 z=(I-\mathbb A)^{-1}u
 =\binom{z_H}{z_Q}.                                \tag{30.9}
\]
Then
\[
 \boxed{
 \rho(\mathbb B(\alpha))<1
 \quad\Longleftrightarrow\quad
 m_0+(v_H^0+g\alpha)z_H+v_Qz_Q+\frac r\alpha<1.}  \tag{30.10}
\]
\end{theorem}

\begin{proof}
Because \(\rho(\mathbb A)<1\), the inverse
\((I-\mathbb A)^{-1}=\sum_{n\geq0}\mathbb A^n\) exists and is
entrywise nonnegative.  The Schur complement of
\(I-\mathbb A\) in \(I-\mathbb B(\alpha)\) is
\[
 1-m(\alpha)-v(\alpha)^{\mathsf T}
 (I-\mathbb A)^{-1}u.                              \tag{30.11}
\]
A nonnegative matrix has spectral radius below one exactly when
\(I-\mathbb B\) is a nonsingular \(M\)-matrix.  With the leading block
already a nonsingular \(M\)-matrix, the block Schur-complement criterion
is precisely positivity of (30.11).  Substitution of (30.8)--(30.9)
gives (30.10).
\end{proof}

\subsection{Exact elimination of the Young tradeoff}

Put
\[
 E=m_0+v_H^0z_H+v_Qz_Q,
 \qquad C=gz_H,qquad D=r.                         \tag{30.12}
\]

\begin{theorem}[Optimal three-sector Young parameter]
\label{thm:v17v30-optimal}
Suppose \(C,D>0\).  There exists \(\alpha>0\) with
\(\rho(\mathbb B(\alpha))<1\) if and only if
\[
 \boxed{\rho(\mathbb A)<1
 \quad\hbox{and}\quad E+2\sqrt{CD}<1.}            \tag{30.13}
\]
The optimal parameter and the maximum Schur margin are
\[
 \boxed{\alpha_*=\sqrt{\frac DC}},                \tag{30.14}
\]
\[
 \boxed{1-E-2\sqrt{CD}.}                          \tag{30.15}
\]
When (30.13) holds, all passing parameters form the interval
\[
 \alpha\in(\alpha_-,\alpha_+),                    \tag{30.16}
\]
where
\[
 \alpha_\pm
 =\frac{1-E\pm\sqrt{(1-E)^2-4CD}}{2C}.            \tag{30.17}
\]
\end{theorem}

\begin{proof}
By (30.10), the remaining scalar condition is
\[
 E+C\alpha+D/\alpha<1.                             \tag{30.18}
\]
The arithmetic--geometric mean inequality gives
\(C\alpha+D/\alpha\geq2\sqrt{CD}\), with equality exactly at
(30.14).  This proves (30.13)--(30.15).  Multiplication of (30.18) by
\(\alpha>0\) gives an upward-opening quadratic; its two positive roots
are (30.17), proving (30.16).
\end{proof}

If \(C=0\), increasing \(\alpha\) removes \(D/\alpha\), and existence
reduces to \(E<1\).  If \(D=0\), decreasing \(\alpha\) removes
\(C\alpha\), again leaving \(E<1\).  These are one-way influence
limits, not feedback cycles.

\subsection{Insertion of the V29 network coefficients}

Let \(\Theta_H,\Theta_Q,\Lambda>0\) be the input exponential
parameters for \(H,Q,G\), respectively.  Compress the finite block in
(29.14) by maximum envelopes.  Define
\[
 H_K^\circ=\sum_{v\in\mathcal V_K}h_v^\circ
 =2\sum_{e\in\mathcal E_K^\circ}\chi_e.           \tag{30.19}
\]
Then the V29 gravity-output row has the form (30.5) with
\[
\begin{aligned}
 g&=\frac{e^{2r_0}H_K^\circ}{\Theta_H},
 &v_H^0&=0,\\
 v_Q&=\frac{(2\kappa_g+\tau)k}{\Theta_Q},
 &m_0&=\frac{\theta+(J/2)(e^{4r_0}-1)}{\Lambda},\\
 &&r&=\frac{e^{2r_0}\chi_*}{4\Lambda}.            \tag{30.20}
\end{aligned}
\]
Here \(r_0>0\) is the tube radius from V29; it is distinguished from
the reciprocal coefficient \(r\) in (30.5).

\begin{corollary}[Explicit V29 acquisition inequality]
\label{cor:v17v30-acquisition}
Once the upper \((H,Q)\) block \(\mathbb A\) and its incoming gravity
column \(u\) satisfy \(\rho(\mathbb A)<1\), the V29 row can be closed
by some \(\alpha>0\) exactly when
\[
\boxed{
 \frac{\theta+(J/2)(e^{4r_0}-1)}{\Lambda}
 +\frac{(2\kappa_g+\tau)k}{\Theta_Q}z_Q
 +2\sqrt{
 \frac{e^{4r_0}H_K^\circ\chi_*}
 {4\Theta_H\Lambda},z_H}
 <1.}                                              \tag{30.21}
\]
The optimizing value is
\[
 \boxed{
 \alpha_*
 =\sqrt{
 \frac{\chi_*\Theta_H}
 {4\Lambda H_K^\circ z_H}}}                       \tag{30.22}
\]
when \(H_K^\circ z_H>0\).
\end{corollary}

\begin{proof}
Substitute (30.20) into (30.12)--(30.14).  The factor
\(e^{4r_0}\) in (30.21) is the product of the two
\(e^{2r_0}\) factors in \(g\) and \(r\).
\end{proof}

The Higgs row already has
\[
 a_{HH}=a_0+A_H\delta,qquad
 a_{HQ}=0,qquad
 u_H=b_0+B_H/\delta                                \tag{30.23}
\]
from V21, after matching the gravity block function to the sum of its
incident mismatch variables.  The remaining acquisition is the
quadratic conformal row
\[
 (a_{QH},a_{QQ},u_Q).                              \tag{30.24}
\]
Equation (30.21) states exactly how small that row must be; no
unspecified three-sector smallness condition remains.

\begin{assessment}[What V30 closes]
\label{ass:v17v30-status}
The mandatory YM/NS audit is complete.  V30 reduces the structured
three-sector Perron problem to one scalar Schur complement and optimizes
the opposing \(\alpha\) and \(1/\alpha\) transfers exactly.  Formula
(30.21) is the explicit coefficient target for the only missing row,
the quadratic conformal conditional moment.

The full continuum remains open because that row, block-exhaustion
compatibility, and the physical RPESM gravitational action have not yet
been supplied.
\end{assessment}

\begin{decision}[V31: derive the quadratic conformal row]
\label{dec:v17v30-next}
Apply the V29 numerator-minimum method with insertion
\(\eta\sum_{v\in K^\circ}s_v^2\), where
\(0<\eta<\kappa_g\).  Derive explicit
\((a_{QH},a_{QQ},u_Q)\), insert them into (30.21), and continue past the
audit with a substantive feasibility or obstruction theorem.
\end{decision}

\begin{firewall}[Claim boundary after V30]
\label{fire:v17v30}
The Schur criterion is exact for the displayed influence matrix.  It
does not prove that the missing quadratic row or the physical RPESM
gravity coefficients satisfy the criterion.
\end{firewall}

\section*{Version V30 append-only record}
\addcontentsline{toc}{section}{Version V30 append-only record}
The complete V29 source was retained byte for byte before this
continuation.  It had \(292106\) bytes, \(9059\) lines, and SHA--256
\[
 \texttt{8E30ED633A9084645EB0085D39141FFB529D4DADE10AA5F42A53332D1FD3573B}.
\]
V30 completed the audit and solved the structured three-sector Young
optimization.


\section{V31: exact double-Young closure and the quadratic-sector obstruction}
\label{sec:v17v31}

\subsection{The quadratic conditional row}

Retain the finite network notation of V29 and put
\[
 Q_K(u)=\sum_{v\in K^\circ}u_v^2,
 \qquad A_g=\frac J2.                              \tag{31.1}
\]
Insert \(e^{\eta Q_K}\) in the conditional measure (29.12), with
\[
 0<\eta<\kappa_g.                                  \tag{31.2}
\]

\begin{theorem}[Explicit quadratic conformal row]
\label{thm:v17v31-qrow}
For every tube radius \(r_Q>0\) and every
\(\beta,\omega,\tau_Q>0\), there is a finite constant \(C_Q\),
independent of the boundary conformal and Higgs values, such that
\[
\begin{aligned}
 \log\mathbb E_{K^\circ}e^{\eta Q_K}
 \leq{}&C_Q+c_QQ_{\partial K}(\sigma)\\
 &+\left[s_0+\frac{t}{\beta}\right]
 \mathcal M_K(\sigma)
 +p\beta\,\Theta_H\max_vV_H(y_v),                 \tag{31.3}
\end{aligned}
\]
where, after maximum-envelope compression,
\[
\begin{aligned}
 c_Q&=\bigl((1+\omega)\kappa_g+\tau_Q\bigr)k,\\
 p&=\frac{e^{2r_Q}H_K^\circ}{\Theta_H},\\
 s_0&=A_g(e^{4r_Q}-1),\\
 t&=\frac{e^{2r_Q}\chi_*}{4}.                     \tag{31.4}
\end{aligned}
\]
Equivalently, with quadratic input parameter \(\Theta_Q\) and network
input parameter \(\Lambda\), the influence entries are
\[
 a_{QH}=p\beta,qquad
 a_{QQ}=q:=\frac{c_Q}{\Theta_Q},qquad
 u_Q=s_0/\Lambda+\frac{t}{\Lambda\beta}.           \tag{31.5}
\]
\end{theorem}

\begin{proof}
The numerator has on-site coefficient \(\kappa_g-\eta>0\), retains the
full mismatch coefficient \(A_g\), and has nonnegative hopping.  Since
\(\mathcal W_K\geq\mathcal M_K\), it is at most
\[
 e^{-A_g\mathcal M_K}
 \left(\frac{\pi}{\kappa_g-\eta}\right)^{k/2}
 \exp\left(\frac{kb_g^2}{4(\kappa_g-\eta)}\right).               \tag{31.6}
\]

For the denominator, use the cube of radius \(r_Q\) about the network
minimizer.  If \(u=u^\sigma+v\), then
\[
 (u_v^\sigma+v_v)^2
 \leq(1+\omega)(u_v^\sigma)^2
 +(1+\omega^{-1})v_v^2.                            \tag{31.7}
\]
Also
\[
 |b_g||u_v^\sigma|
 \leq\tau_Q(u_v^\sigma)^2+\frac{b_g^2}{4\tau_Q}.  \tag{31.8}
\]
Lemmas~\ref{lem:v17v29-tube}--\ref{lem:v17v29-maximum}, with
\(r=r_Q\) and \(\alpha=\beta\), bound the remaining denominator
costs.  After division by (31.6), the network-minimum coefficient is
\[
 A_ge^{4r_Q}-A_g+\frac{e^{2r_Q}\chi_*}{4\beta},    \tag{31.9}
\]
and the Higgs coefficient sum is
\(e^{2r_Q}H_K^\circ\beta\).  Equations
(31.3)--(31.5) follow.
\end{proof}

\subsection{Exact two-Young three-sector optimization}

The coefficients from V21, V29, and Theorem
\ref{thm:v17v31-qrow} have the structured matrix form
\[
 \mathbb B(\delta,\alpha,\beta)
 =\begin{pmatrix}
 a_0+A_H\delta&0&b_0+B_H/\delta\\
 p\beta&q&s_0+t/\beta\\
 g\alpha&v&m_0+r/\alpha
 \end{pmatrix},                                   \tag{31.10}
\]
where all coefficients are nonnegative and
\(\delta,\alpha,\beta>0\).  The factors \(\Lambda^{-1}\) have been
absorbed into \(s_0,t,m_0,r\) in (31.10).

\begin{theorem}[Complete nested optimization]
\label{thm:v17v31-nested}
Assume
\[
 A_H,B_H,p,t,g,r,v>0,qquad
 0\leq a_0<1,qquad0\leq q<1.                     \tag{31.11}
\]
Put \(c=1-a_0\) and
\[
 h(\delta)=
 \frac{b_0+B_H/\delta}{c-A_H\delta},
 \qquad0<\delta<\frac c{A_H}.                     \tag{31.12}
\]
For fixed \(\delta\), there exist \(\alpha,\beta>0\) with
\(\rho(\mathbb B)<1\) if and only if
\[
 \boxed{
 m_0+2\sqrt{rgh(\delta)}
 +\frac{v}{1-q}
 \left[s_0+2\sqrt{tph(\delta)}\right]<1.}         \tag{31.13}
\]
The optimizing parameters are
\[
 \boxed{
 \alpha_*(\delta)=\sqrt{\frac{r}{g h(\delta)}},
 \qquad
 \beta_*(\delta)=\sqrt{\frac{t}{p h(\delta)}}.}  \tag{31.14}
\]

If \(b_0>0\), the unique minimizer of \(h\) is
\[
 \boxed{
 \delta_h=
 \frac{\sqrt{A_HB_H(A_HB_H+b_0c)}-A_HB_H}
 {A_Hb_0}.}                                       \tag{31.15}
\]
If \(b_0=0\), it is
\[
 \boxed{\delta_h=\frac{c}{2A_H}.}                 \tag{31.16}
\]
Consequently, existence of all three Young parameters is equivalent to
(31.13) evaluated at \(h_*=h(\delta_h)\).
\end{theorem}

\begin{proof}
For fixed \(\delta,\beta\), the leading \((H,Q)\) block is triangular
and has spectral radius below one exactly when
\(a_0+A_H\delta<1\) and \(q<1\).  In the notation of V30,
\[
 z_H=h(\delta),qquad
 z_Q=\frac{s_0+t/\beta+p\beta h(\delta)}{1-q}.     \tag{31.17}
\]
Theorem~\ref{thm:v17v30-optimal} first minimizes
\(g\alpha h+r/\alpha\), giving the first formula in (31.14) and the
term \(2\sqrt{rgh}\).  The remaining \(\beta\)-dependence is
\(t/\beta+p\beta h\), whose exact minimum is
\(2\sqrt{tph}\) at the second value in (31.14).  This proves
(31.13)--(31.14).

The left side of (31.13) is strictly increasing in \(h\), so it remains
to minimize (31.12).  Differentiation shows that the unique critical
point solves
\[
 A_Hb_0\delta^2+2A_HB_H\delta-B_Hc=0.              \tag{31.18}
\]
The positive root is (31.15); it lies inside
\((0,c/A_H)\) and is the minimum because \(h\) diverges at both
endpoints.  When \(b_0=0\), (31.18) gives (31.16).
\end{proof}

The theorem completely solves the finite-dimensional optimization of
the displayed three-sector certificate.  It remains to check whether
the analytic parameter windows permit its hypothesis \(q<1\).

\subsection{Exact global quadratic moment window}

Let \(G\) be any finite connected conformal graph, including the
nonnegative mismatch and rescaled-Higgs diagonal hopping terms of
V23.  All interactions depend only on differences.  Decompose
\[
 s=t\mathbf1+r,qquad r\perp\mathbf1.              \tag{31.19}
\]

\begin{theorem}[Sharp common-mode quadratic threshold]
\label{thm:v17v31-qwindow}
Conditional on arbitrary rescaled Higgs variables, for
\(\Theta_Q\geq0\),
\[
 \mathbb E\exp\left(
 \Theta_Q\sum_{x\in G}s_x^2\right)<\infty
 \quad\Longleftrightarrow\quad
 \Theta_Q<\kappa_g.                               \tag{31.20}
\]
\end{theorem}

\begin{proof}
If \(n=|G|\), then
\[
 \sum_xs_x^2=nt^2+\|r\|_2^2,qquad
 \sum_xs_x=nt.                                    \tag{31.21}
\]
Every mismatch and hopping term is independent of \(t\), so the common
mode has density proportional to
\(e^{-n\kappa_gt^2-nb_gt}\).  Its inserted quadratic moment is finite
exactly when \(\Theta_Q<\kappa_g\), proving necessity.  Under this
strict inequality, dropping all nonnegative difference interactions
leaves a positive-definite Gaussian in \((t,r)\), proving sufficiency.
\end{proof}

\subsection{The quadratic-sector countertheorem}

\begin{countertheorem}[No Perron closure with absolute quadratic
boundary transfer]
\label{ctr:v17v31-qobstruction}
For every nonempty interior block \(k\geq1\), every admissible positive quadratic
parameter \(\Theta_Q<\kappa_g\), and every
\(\omega,\tau_Q>0\), the quadratic diagonal entry in (31.5) satisfies
\[
 \boxed{
 q=\frac{((1+\omega)\kappa_g+\tau_Q)k}{\Theta_Q}>k\geq1.}       \tag{31.22}
\]
Consequently, the matrix (31.10) has spectral radius greater than one
for every \(\delta,\alpha,\beta>0\).  The exact feasibility condition
(31.13) is never reached by this absolute-quadratic sector.
\end{countertheorem}

\begin{proof}
Equation (31.22) follows immediately from
\(\Theta_Q<\kappa_g\), \(k\geq1\), and positive
\(\omega,\tau_Q\).  Since a nonnegative matrix has spectral radius at
least each diagonal entry, \(\rho(\mathbb B)\geq q>1\).
\end{proof}

The obstruction persists as \(\omega,\tau_Q\downarrow0\).  For a
one-site block the diagonal approaches one from above; for larger
blocks it approaches at least \(k\).  Thus it is not caused by the
earlier factor two in (29.25), nor can it be repaired by retuning the
Young parameters.

\begin{assessment}[What V31 closes]
\label{ass:v17v31-status}
V31 derives the previously missing quadratic row and solves the full
three-Young matrix optimization exactly.  It then proves that the row
cannot satisfy the necessary diagonal Perron condition because the
same Gaussian common mode fixes the sharp input-moment threshold.

This is a rigorous failure witness for the absolute-quadratic
Dobrushin sector, not a divergence theorem for the conformal model.
The successful mismatch and network-minimum estimates remain valid.
\end{assessment}

\begin{decision}[V32: integrate the common mode before blocking]
\label{dec:v17v31-next}
Move upstream.  On each connected finite cylinder, integrate the global
common conformal mode exactly using (31.19), since all rescaled hopping
and mismatch terms are translation invariant.  Formulate block
conditioning only on the mean-zero field, where boundary control should
use the inverse grounded Laplacian rather than absolute site squares.
Derive the resulting covariance and determine whether the fatal
diagonal (31.22) disappears.
\end{decision}

\begin{firewall}[Claim boundary after V31]
\label{fire:v17v31}
V31 rules out one comparison architecture.  It does not rule out the
convex conformal Gibbs measure, provide the physical RPESM gravity
action, or establish the Einstein--SM continuum limit.
\end{firewall}

\section*{Version V31 append-only record}
\addcontentsline{toc}{section}{Version V31 append-only record}
The complete V30 source was retained byte for byte before this
continuation.  It had \(300958\) bytes, \(9325\) lines, and SHA--256
\[
 \texttt{1FADFFD044F383D2498F2BD3281BDF774A71C1789D8962E1BB1F4979CC31FEDC}.
\]
V31 closed the three-Young algebra and proved the absolute-quadratic
sector obstruction.


\section{V32: exact common-mode integration and the mean-zero quotient}
\label{sec:v17v32}

\subsection{The difference-only conformal model}

Let \(G=(\mathcal V,\mathcal E)\) be finite and connected, with
\(n=|\mathcal V|\), oriented incidence matrix \(B\), and
\(a_x=\|y_x\|^2\).  Conditional on the rescaled Higgs field, consider
\[
\begin{aligned}
 S_G(s\mid y)={}&\kappa_g\|s\|_2^2
 +b_g\langle\mathbf1,s\rangle\\
 &+\sum_{e=(v,w)}\left[
 \frac J2V_g((Bs)_e)
 +\chi_e\bigl(a_ve^{(Bs)_e}+a_we^{-(Bs)_e}\bigr)
 \right].                                         \tag{32.1}
\end{aligned}
\]
Additive mismatch constants and Higgs cross terms have been suppressed;
they are independent of the conformal integration variable.  Assume
\[
 \kappa_g>0,qquad J\geq0,qquad\chi_e\geq0.        \tag{32.2}
\]

Write
\[
 s=t\mathbf1+r,qquad
 t=\frac1n\langle\mathbf1,s\rangle,qquad
 r\in\mathcal H_0:=\{r:\langle\mathbf1,r\rangle=0\}.           \tag{32.3}
\]

\begin{theorem}[Exact integration of the common conformal mode]
\label{thm:v17v32-factor}
Under (32.2), the conditional law factorizes into an independent
Gaussian common mode and a mean-zero difference field.  Specifically,
\[
 t\sim N\left(-\frac{b_g}{2\kappa_g},
 \frac{1}{2n\kappa_g}\right),                     \tag{32.4}
\]
and the quotient field has density on \(\mathcal H_0\) proportional to
\[
 e^{-\Phi_y(r)},d\mathcal H_0(r),                 \tag{32.5}
\]
where
\[
\begin{aligned}
 \Phi_y(r)={}&\kappa_g\|r\|_2^2\\
 &+\sum_{e=(v,w)}\left[
 \frac J2V_g((Br)_e)
 +\chi_e\bigl(a_ve^{(Br)_e}+a_we^{-(Br)_e}\bigr)
 \right].                                         \tag{32.6}
\end{aligned}
\]
The common-mode integral is independent of \(y\).  Hence integrating
it out creates no Higgs influence coefficient.
\end{theorem}

\begin{proof}
Orthogonality in (32.3) gives
\[
 \|s\|_2^2=nt^2+\|r\|_2^2,qquad
 \langle\mathbf1,s\rangle=nt,qquad Bs=Br.        \tag{32.7}
\]
Every interaction in the second line of (32.1) is therefore independent
of \(t\).  Completing the square,
\[
 n\kappa_gt^2+nb_gt
 =n\kappa_g\left(t+\frac{b_g}{2\kappa_g}\right)^2
 -\frac{nb_g^2}{4\kappa_g},                       \tag{32.8}
\]
which proves (32.4)--(32.6).  The Gaussian normalizing factor in
(32.8) depends on \(n,\kappa_g,b_g\), but not on the Higgs variables.
\end{proof}

\begin{corollary}[Difference observables live entirely on the quotient]
\label{cor:v17v32-difference}
For every nonnegative or integrable measurable function \(F\),
\[
 \mathbb E\bigl[F(Bs,y)\mid y\bigr]
 =\mathbb E_{\mathcal H_0}
 \bigl[F(Br,y)\mid y\bigr].                       \tag{32.9}
\]
In particular, every mismatch moment and every geometric rescaled-Higgs
edge used in V20--V31 is unchanged by removal of the common mode.
\end{corollary}

\begin{proof}
The integrand is independent of \(t\), so the same Gaussian factor
cancels between numerator and denominator.
\end{proof}

This cancellation is stronger than a moment estimate.  The problematic
absolute quadratic variable from V31 is absent from the state space of
all difference observables.

\subsection{Uniform convexity on the quotient}

For an edge \(e=(v,w)\), define
\[
 h_e(r,y)=\chi_e
 \bigl(a_ve^{(Br)_e}+a_we^{-(Br)_e}\bigr)\geq0.    \tag{32.10}
\]

\begin{theorem}[Uniform quotient Hessian]
\label{thm:v17v32-hessian}
The Hessian of (32.6), restricted to \(\mathcal H_0\), is
\[
 \nabla^2\Phi_y(r)
 =2\kappa_gI_{mathcal H_0}
 +B^{\mathsf T}\operatorname{diag}
 \bigl(2JV_g((Br)_e)+h_e(r,y)\bigr)B,              \tag{32.11}
\]
and therefore
\[
 \nabla^2\Phi_y(r)\succeq2\kappa_gI_{mathcal H_0}             \tag{32.12}
\]
uniformly in the volume, graph, conformal field, and Higgs boundary
values.
\end{theorem}

\begin{proof}
The second derivative of \(V_g(q)=e^{2q}+e^{-2q}\) is \(4V_g(q)\).
After multiplication by \(J/2\), its edge Hessian weight is
\(2JV_g(q)\).  Each hopping summand in (32.10) equals its own second
derivative with respect to the edge difference.  Pulling the edge
Hessians back by the incidence matrix gives (32.11).  All diagonal
weights are nonnegative, proving (32.12).
\end{proof}

\begin{corollary}[Volume-uniform Brascamp--Lieb bound]
\label{cor:v17v32-bl}
For every locally Lipschitz \(f:\mathcal H_0\to\mathbb R\) with finite
variance under (32.5),
\[
 \operatorname{Var}_y(f)
 \leq\frac{1}{2\kappa_g}
 \mathbb E_y\|\nabla_{\mathcal H_0}f\|_2^2.        \tag{32.13}
\]
In particular, if \(\ell\perp\mathbf1\), then
\[
 \operatorname{Var}_y\langle\ell,r\rangle
 \leq\frac{\|\ell\|_2^2}{2\kappa_g}.             \tag{32.14}
\]
For every edge \(e=(v,w)\),
\[
 \operatorname{Var}_y(r_w-r_v)\leq\frac1{\kappa_g}.             \tag{32.15}
\]
\end{corollary}

\begin{proof}
The finite-dimensional Brascamp--Lieb variance inequality gives
\[
 \operatorname{Var}_y(f)
 \leq\mathbb E_y\left[
 \nabla f^{\mathsf T}(\nabla^2\Phi_y)^{-1}\nabla f
 \right].                                         \tag{32.16}
\]
Use (32.12) to bound the inverse by
\((2\kappa_g)^{-1}I\).  For a linear functional the gradient is
\(\ell\), giving (32.14).  The edge vector
\(\ell=\mathbf e_w-\mathbf e_v\) has squared norm two, proving
(32.15).
\end{proof}

The variance estimate is not a substitute for the sharp exponential
mismatch threshold \(2\theta\leq J\).  It supplies the uniform quadratic
control needed after the inadmissible absolute sector has been removed.

\subsection{Converting centered size into mismatch energy}

Let \(K=(\mathcal V_K,\mathcal E_K)\) be a connected finite network,
let \(L_K=B_K^{\mathsf T}B_K\) be its unweighted graph Laplacian, and
let \(\lambda_2(K)>0\) be its first nonzero eigenvalue.

\begin{lemma}[Quotient Poincare conversion]
\label{lem:v17v32-poincare}
For every \(r\perp\mathbf1\),
\[
 \|r\|_2^2
 \leq\frac{1}{\lambda_2(K)}\sum_{e\in\mathcal E_K}(B_Kr)_e^2
 \leq\frac{1}{4\lambda_2(K)}
 \sum_{e\in\mathcal E_K}\bigl(V_g((B_Kr)_e)-2\bigr).          \tag{32.17}
\]
\end{lemma}

\begin{proof}
The first inequality is the Rayleigh characterization of
\(\lambda_2(K)\).  Since
\[
 2\cosh(2q)-2\geq4q^2,                             \tag{32.18}
\]
the second follows term by term.
\end{proof}

\begin{corollary}[Centered network minimizer needs no absolute
quadratic sector]
\label{cor:v17v32-minimizer}
Let \(s^\sigma\) be the V28 mismatch minimizer and center it by
\[
 \widetilde s^\sigma
 =s^\sigma-\frac1{|\mathcal V_K|}
 \langle\mathbf1,s^\sigma\rangle\mathbf1.         \tag{32.19}
\]
Then
\[
 \|\widetilde s^\sigma\|_2^2
 \leq\frac{
 \mathcal M_K(\sigma)-2|\mathcal E_K|}
 {4\lambda_2(K)}
 \leq\frac{\mathcal M_K(\sigma)}{4\lambda_2(K)}. \tag{32.20}
\]
\end{corollary}

\begin{proof}
Centering leaves every edge difference unchanged, so the edge sum at
\(\widetilde s^\sigma\) is still \(\mathcal M_K(\sigma)\).  Apply
Lemma~\ref{lem:v17v32-poincare}.
\end{proof}

Equation (32.20) converts the centered on-site cost into the already
admissible network-minimum sector.  It does not introduce a new
absolute \(Q\)-row.  Its coefficient depends on the block spectral gap,
which is uniform for a fixed finite family of bounded-size block
shapes.

\subsection{Removal of the V31 diagonal obstruction}

\begin{theorem}[The obstruction is absent after exact quotienting]
\label{thm:v17v32-removal}
For the difference-observable and rescaled-Higgs conditional problem
(32.1), exact integration of \(t\) removes the quadratic sector and its
diagonal coefficient \(q>1\) from (31.22).  The remaining quotient
measure is normalizable, uniformly strongly convex, and has the
volume-uniform covariance bound (32.13).  Centered finite-block anchor
costs can be assigned to \(\mathcal M_K\) by (32.20).
\end{theorem}

\begin{proof}
Theorem~\ref{thm:v17v32-factor} proves that the common-mode factor is
independent of every difference and Higgs variable.  Hence it can be
integrated before any comparison matrix is formed.  The quotient action
and its control are given by Theorem~\ref{thm:v17v32-hessian},
Corollary~\ref{cor:v17v32-bl}, and (32.20).  Since no absolute
quadratic moment is used, the row whose diagonal was (31.22) is absent.
\end{proof}

This theorem removes a failed proof architecture; it does not yet prove
that a local quotient block schedule respects the one global mean-zero
constraint.  That compatibility must be constructed explicitly.

\begin{assessment}[What V32 closes]
\label{ass:v17v32-status}
V32 integrates the conformal common mode exactly and proves that it
contributes no Higgs feedback.  The mean-zero quotient action is
uniformly strongly convex, its linear covariances are volume-uniform,
and centered block size is controlled by mismatch energy through the
graph spectral gap.  The fatal absolute-quadratic diagonal of V31 is
therefore unnecessary for all difference observables in the model.

The remaining issue is geometric: local block updates must preserve the
global mean-zero constraint or use coordinates that do so automatically.
\end{assessment}

\begin{decision}[V33: local coordinates on the mean-zero quotient]
\label{dec:v17v32-next}
Choose a spanning tree and an orthonormal or uniformly conditioned basis
of mean-zero local differences.  Express block resampling in those
coordinates, quantify the condition number and overlap of the basis,
and determine whether the V29 network-minimum row remains uniformly
contractive without reintroducing a global quadratic variable.
\end{decision}

\begin{firewall}[Claim boundary after V32]
\label{fire:v17v32}
The quotient theorem is exact for the declared homogeneous
difference-only conformal model.  It does not identify the physical
RPESM gravitational action, construct the infinite-volume block
dynamics, or prove the Einstein--SM continuum limit.
\end{firewall}

\section*{Version V32 append-only record}
\addcontentsline{toc}{section}{Version V32 append-only record}
The complete V31 source was retained byte for byte before this
continuation.  It had \(310349\) bytes, \(9601\) lines, and SHA--256
\[
 \texttt{7398BDC074B91FF6372CEFF84EE0822FDA330C9AD8D6C7938E718BB58E90554C}.
\]
V32 removed the common conformal mode exactly and established the
uniformly convex quotient formulation.


\section{V33: mean-zero coordinates, tree obstruction, and Haar locality}
\label{sec:v17v33}

\subsection{Why raw spanning-tree differences are not uniformly conditioned}

Let \(G\) have \(n\) vertices and let \(T\) be a spanning tree.  On
\(\mathcal H_0=\mathbf1^\perp\), the tree incidence map
\[
 D_T:\mathcal H_0\longrightarrow\mathbb R^{n-1},
 \qquad D_Tr=B_Tr                                  \tag{33.1}
\]
is invertible.

\begin{theorem}[Exact conditioning of tree-difference coordinates]
\label{thm:v17v33-tree}
Let \(L_T=B_T^{\mathsf T}B_T\) be the tree Laplacian.  Then
\[
 \lambda_2(T)\|r\|_2^2
 \leq\|D_Tr\|_2^2
 \leq\lambda_{\max}(T)\|r\|_2^2,
 \qquad r\in\mathcal H_0.                          \tag{33.2}
\]
Consequently,
\[
 \|D_T^{-1}\|_{2\to2}=\lambda_2(T)^{-1/2},
 \qquad
 \operatorname{cond}_2(D_T)
 =\sqrt{\frac{\lambda_{\max}(T)}{\lambda_2(T)}}.  \tag{33.3}
\]
For the path tree \(P_n\),
\[
 \lambda_2(P_n)=2\left(1-\cos\frac{\pi}{n}\right),              \tag{33.4}
\]
so \(\|D_T^{-1}\|_{2\to2}\sim n/\pi\).
\end{theorem}

\begin{proof}
The identity
\(\|B_Tr\|_2^2=r^{\mathsf T}L_Tr\) and the spectral theorem on
\(\mathbf1^\perp\) give (33.2)--(33.3).  The path eigenvalues are
\(2-2\cos(k\pi/n)\), \(0\leq k<n\), giving (33.4) and the stated
asymptotic.
\end{proof}

Thus a spanning tree removes the mean-zero constraint algebraically,
but the reconstruction map is not uniformly bounded on path
exhaustions.  Using its raw edge coordinates would reintroduce a
long-range amplification into the block rows.

\subsection{The discrete Haar basis on a dyadic path}

Let \(n=2^m\), with path vertices \(0,\ldots,n-1\).  For every scale
\(1\leq j\leq m\), put \(L=2^j\).  For
\(0\leq k<n/L\), define the dyadic interval
\[
 I_{j,k}=\{kL,\ldots,(k+1)L-1\}                    \tag{33.5}
\]
and the Haar vector
\[
 \psi_{j,k}(x)=
 \begin{cases}
 L^{-1/2},&kL\leq x<kL+L/2,\\
 -L^{-1/2},&kL+L/2\leq x<(k+1)L,\\
 0,&x\notin I_{j,k}.
 \end{cases}                                      \tag{33.6}
\]

\begin{theorem}[Orthonormal quotient coordinates]
\label{thm:v17v33-haar}
The \(n-1\) vectors \(\psi_{j,k}\) form an orthonormal basis of
\(\mathcal H_0\).  Hence every mean-zero field has the unique expansion
\[
 r=\sum_{j=1}^m\sum_{k=0}^{n/2^j-1}c_{j,k}\psi_{j,k},
 \qquad
 \|r\|_2^2=\sum_{j,k}c_{j,k}^2.                   \tag{33.7}
\]
Updating one coefficient preserves the global mean-zero constraint
exactly, and the on-site quotient energy is the diagonal Gaussian
\(\kappa_g\sum c_{j,k}^2\).
\end{theorem}

\begin{proof}
Every vector in (33.6) has zero sum and unit norm.  Two dyadic intervals
are disjoint or nested.  In the nested case, the coarser wavelet is
constant on the support of the finer wavelet, whose sum is zero; hence
the inner product vanishes.  Finally,
\[
 \sum_{j=1}^m\frac{n}{2^j}=n-1,                   \tag{33.8}
\]
so the orthonormal family has the dimension of \(\mathcal H_0\).
Parseval gives (33.7).
\end{proof}

\subsection{Exact edge sparsity}

Orient each path edge from \(x-1\) to \(x\), \(1\leq x<n\), and let
\[
 b_{x;j,k}=(B\psi_{j,k})_x.                        \tag{33.9}
\]

\begin{theorem}[Three-edge support of a Haar coefficient]
\label{thm:v17v33-column}
For a wavelet of support length \(L\), the incidence column
\(B\psi_{j,k}\) is nonzero on at most three edges: the left support
boundary, the midpoint, and the right support boundary.  Its nonzero
magnitudes are contained in
\[
 \left\{L^{-1/2},,2L^{-1/2},,L^{-1/2}\right\}.  \tag{33.10}
\]
Consequently,
\[
 \|B\psi_{j,k}\|_1\leq\frac4{\sqrt L},
 \qquad
 \|B\psi_{j,k}\|_2^2\leq\frac6L.                 \tag{33.11}
\]
\end{theorem}

\begin{proof}
The wavelet is constant on each half of its support and zero outside.
Its difference therefore vanishes except where one of those three
constant regions changes to another.  The jumps are respectively
\(L^{-1/2}\), \(-2L^{-1/2}\), and \(L^{-1/2}\), with an exterior jump
omitted when the support meets a path endpoint.  Equations (33.10)--
(33.11) follow.
\end{proof}

\begin{theorem}[Uniform row overlap and exact square sum]
\label{thm:v17v33-row}
For every path edge \(x\),
\[
 \sum_{j,k}|b_{x;j,k}|
 \leq2(\sqrt2+1),                                  \tag{33.12}
\]
and exactly
\[
 \boxed{\sum_{j,k}|b_{x;j,k}|^2=2.}               \tag{33.13}
\]
Both bounds are independent of \(n\).
\end{theorem}

\begin{proof}
At scale \(L=2^j\), the edge at position \(x\) is a support boundary
only if \(2^j\) divides \(x\); then at most two wavelets contribute,
each with magnitude \(2^{-j/2}\).  It is a midpoint for at most the
single next scale, with magnitude \(2\,2^{-j/2}\).  If
\(\nu_2(x)=v\), the total is at most
\[
 2\sum_{j=1}^{v}2^{-j/2}+2\,2^{-(v+1)/2}
 \leq2\sum_{j=1}^{\infty}2^{-j/2}
 =2(\sqrt2+1),                                     \tag{33.14}
\]
proving (33.12).

For (33.13), the edge difference vector
\(\mathbf e_x-\mathbf e_{x-1}\) belongs to \(\mathcal H_0\) and has
squared norm two.  Its Haar coefficient against \(\psi_{j,k}\) is
exactly \(b_{x;j,k}\).  Parseval's identity in the complete basis of
Theorem~\ref{thm:v17v33-haar} gives (33.13).
\end{proof}

\subsection{Implication for quotient block updates}

Let \(c_\alpha\) denote one Haar coefficient and write each edge
difference as
\[
 (Br)_e=b_{e\alpha}c_\alpha+\xi_{e,\alpha}.        \tag{33.15}
\]
Theorem~\ref{thm:v17v33-column} shows that the conditional density of
\(c_\alpha\) contains at most three mismatch and hopping edge terms,
regardless of scale or volume.  The on-site term is simply
\(\kappa_gc_\alpha^2\), and no global constraint or absolute boundary
quadratic appears.

Conversely, (33.13) shows that an individual physical edge receives a
fixed total squared coefficient from all quotient coordinates.  Thus
any influence estimate quadratic in the incidence amplitudes has an
exact volume-independent overlap constant two.  Estimates linear in
the amplitudes use the uniform bound (33.12).

\begin{countertheorem}[Raw tree coordinates fail the requested
uniformity on paths]
\label{ctr:v17v33-tree}
No argument whose reconstruction coefficient is bounded below by
\(\|D_{P_n}^{-1}\|_{2\to2}\) can yield a volume-uniform block row on
dyadic path exhaustions.  The coefficient diverges linearly in \(n\)
by (33.3)--(33.4).  The Haar synthesis map has norm one and avoids this
specific obstruction.
\end{countertheorem}

\begin{proof}
This is the asymptotic in Theorem~\ref{thm:v17v33-tree}, compared with
the Parseval identity (33.7).
\end{proof}

\begin{assessment}[What V33 closes]
\label{ass:v17v33-status}
V33 gives an exact coordinate realization of the mean-zero quotient on
dyadic paths.  It proves unit conditioning, three-edge column support,
a uniform absolute row overlap, and the exact squared overlap two.  It
also proves that raw spanning-tree differences cannot supply the same
volume-uniform reconstruction bound.

The remaining task is analytic: conditional Haar-coefficient moments
must be compiled into mismatch and Higgs influence rows, with the
amplitude factors in (33.11)--(33.13) retained rather than replaced by
the number of scales.
\end{assessment}

\begin{decision}[V34: conditional Haar-coordinate compiler]
\label{dec:v17v33-next}
Condition one Haar coefficient, keep its Gaussian on-site term exact,
and anchor near the minimizer of its at-most-three-edge convex
potential.  Derive influence coefficients weighted by
\(|b_{e\alpha}|\) or \(|b_{e\alpha}|^2\), then sum them with
(33.12)--(33.13).  Test whether the resulting quotient row has a strict
volume-uniform Perron margin.
\end{decision}

\begin{firewall}[Claim boundary after V33]
\label{fire:v17v33}
The Haar construction is proved only for dyadic path exhaustions.  It
does not yet give a conditional moment theorem, a higher-dimensional
lattice basis, the physical gravitational action, or the continuum
limit.
\end{firewall}

\section*{Version V33 append-only record}
\addcontentsline{toc}{section}{Version V33 append-only record}
The complete V32 source was retained byte for byte before this
continuation.  It had \(320458\) bytes, \(9892\) lines, and SHA--256
\[
 \texttt{5432081AD40F915C2CC1B9EA573191B97AD266F46813D0378852BB99DCE0D632}.
\]
V33 constructed uniformly conditioned local quotient coordinates on
dyadic paths.


\section{V34: scalar Haar conditioning and the cavity-cancellation obstruction}
\label{sec:v17v34}

\subsection{One quotient coefficient}

Let \(\psi_\alpha\) be one Haar wavelet from V33, let
\[
 c=c_\alpha,qquad
 b_e=(B\psi_\alpha)_e,qquad
 E_\alpha=\{e:b_e\ne0\},qquad
 a_\alpha=\max_{e\in E_\alpha}|b_e|.              \tag{34.1}
\]
Then \(|E_\alpha|\leq3\).  With all other coefficients fixed, write
\[
 (Br)_e=b_ec+\xi_e,qquad e\in E_\alpha.           \tag{34.2}
\]
The scalar conditional action on the quotient is
\[
 \Phi_\alpha(c)
 =\kappa_gc^2+A_gW_\alpha(c)+H_\alpha(c),
 \qquad A_g=\frac J2,                              \tag{34.3}
\]
where
\[
 W_\alpha(c)=\sum_{e\in E_\alpha}V_g(\xi_e+b_ec)  \tag{34.4}
\]
and \(H_\alpha\geq0\) is the sum of the diagonal rescaled-Higgs
hopping terms on those edges.  There is no linear common-mode term and
no mean-zero constraint on \(c\).

\subsection{The zero-anchor compiler}

Let
\[
 R_\alpha(\xi)=\sum_{e\in E_\alpha}V_g(\xi_e)      \tag{34.5}
\]
be the cavity edge energy obtained by deleting the coefficient
\(c_\alpha\).

\begin{theorem}[Scalar Haar zero-anchor estimate]
\label{thm:v17v34-zero}
Assume \(0<\eta<\kappa_g\).  For every
\(\varepsilon,\gamma>0\), there is a finite constant \(C_\alpha\),
independent of the cavity and Higgs values, such that
\[
\begin{aligned}
 \log\mathbb E_\alpha e^{\eta c_\alpha^2}
 \leq{}&C_\alpha+left[
 A_ge^{2a_\alpha\varepsilon}
 +\frac{e^{a_\alpha\varepsilon}\chi_*}{4\gamma}
 \right]R_\alpha(\xi)\\
 &+e^{a_\alpha\varepsilon}\gamma
 \sum_{v}h_{v,\alpha}V_H(y_v),                    \tag{34.6}
\end{aligned}
\]
where \(h_{v,\alpha}\) is the total hopping weight of the affected
edges incident to \(v\).
\end{theorem}

\begin{proof}
The moment numerator has Gaussian coefficient
\(\kappa_g-\eta>0\).  Dropping the nonnegative edge and hopping terms
gives the finite upper bound
\(\sqrt{\pi/(\kappa_g-\eta)}\).

For the denominator, use \(|c|\leq\varepsilon\).  On every affected
edge,
\[
 V_g(\xi_e+b_ec)
 \leq e^{2|b_e|\varepsilon}V_g(\xi_e)
 \leq e^{2a_\alpha\varepsilon}V_g(\xi_e).         \tag{34.7}
\]
For an oriented edge \(e=(v,w)\), its hopping term is at most
\[
 e^{a_\alpha\varepsilon}\chi_e
 \bigl(a_ve^{\xi_e}+a_we^{-\xi_e}\bigr)           \tag{34.8}
\]
and Young's inequality bounds the bracket by
\[
 \gamma(a_v^2+a_w^2)+\frac{V_g(\xi_e)}{4\gamma}.  \tag{34.9}
\]
The on-site anchor cost is at most
\(\kappa_g\varepsilon^2\), and the anchor has length
\(2\varepsilon\).  Sum (34.7)--(34.9), divide the numerator by the
denominator lower bound, and take logarithms.
\end{proof}

The coefficient of \(R_\alpha\) approaches \(A_g\) from above as
\(\varepsilon\downarrow0\) and \(\gamma\uparrow\infty\).  Since the
sharp mismatch input window is at most \(A_g\), this anchor has no
strict gravity-only margin.

\subsection{Anchoring at the edge-energy minimizer}

Define
\[
 M_\alpha(\xi)=\inf_{c\in\mathbb R}W_\alpha(c).    \tag{34.10}
\]

\begin{lemma}[Unique scalar cavity center]
\label{lem:v17v34-center}
The minimum in (34.10) is attained at a unique \(c_\alpha^*(\xi)\),
characterized by
\[
 \sum_{e\in E_\alpha}
 b_e\sinh\bigl(2(\xi_e+b_ec_\alpha^*)\bigr)=0.    \tag{34.11}
\]
For \(|c-c_\alpha^*|\leq\varepsilon\),
\[
 W_\alpha(c)
 \leq e^{2a_\alpha\varepsilon}M_\alpha.           \tag{34.12}
\]
\end{lemma}

\begin{proof}
At least one \(b_e\) is nonzero, so (34.4) is coercive and strictly
convex in \(c\).  Differentiation gives (34.11).  The multiplicative
bound used in (34.7), now centered at \(c_\alpha^*\), gives (34.12).
\end{proof}

For an edge-moment insertion \(e^{\theta W_\alpha}\), with
\(0<\theta<A_g\), retaining
\(W_\alpha\geq M_\alpha\) in the numerator produces the desired sharp
leading coefficient
\[
 \theta+A_g(e^{2a_\alpha\varepsilon}-1)            \tag{34.13}
\]
on \(M_\alpha\), before hopping is added.  However, the denominator
anchor also pays
\[
 \kappa_g(c_\alpha^*)^2.                           \tag{34.14}
\]
The next theorem shows that (34.14) cannot be controlled by
\(M_\alpha\) alone.

\subsection{Exact cross-scale cancellation}

Take the coarsest Haar wavelet \(\psi_\alpha\) on the full dyadic path
of length \(n\geq4\).  It affects only the midpoint edge
\(e_*=(n/2-1,n/2)\), with
\[
 |b_{e_*\alpha}|=\frac2{\sqrt n}.                  \tag{34.15}
\]
Let \(\psi_\beta\) be a wavelet of support length \(n/2\) adjacent to
that midpoint.  The same edge is a support boundary for
\(\psi_\beta\), and
\[
 |b_{e_*\beta}|=\sqrt{\frac2n}.                   \tag{34.16}
\]

\begin{countertheorem}[The cavity minimum does not control its center]
\label{ctr:v17v34-cancellation}
Fix every Haar coefficient except \(c_\beta\) to zero in the cavity of
\(c_\alpha\).  Then
\[
 W_\alpha(c)
 =V_g(b_{e_*\alpha}c+b_{e_*\beta}c_\beta),         \tag{34.17}
\]
and
\[
 M_\alpha=2,qquad
 c_\alpha^*
 =-\frac{b_{e_*\beta}}{b_{e_*\alpha}}c_\beta.     \tag{34.18}
\]
Thus \(M_\alpha\) remains constant while
\(|c_\alpha^*|\to\infty\) as \(|c_\beta|\to\infty\).  In
particular, no constants \(C_0,C_1<\infty\) can satisfy
\[
 \kappa_g(c_\alpha^*)^2\leq C_0+C_1M_\alpha       \tag{34.19}
\]
for all admissible Haar cavities.
\end{countertheorem}

\begin{proof}
Equation (34.17) follows from the exact supports in V33.  The unique
minimum of \(V_g(q)\) is \(V_g(0)=2\), attained when the affine
argument vanishes.  This gives (34.18), and (34.19) is contradicted by
letting \(|c_\beta|\to\infty\).
\end{proof}

For this same configuration, the zero anchor pays
\[
 R_\alpha=V_g(b_{e_*\beta}c_\beta),                \tag{34.20}
\]
which is unbounded, whereas the minimizer anchor pays the quadratic
cost
\[
 \kappa_g\frac{b_{e_*\beta}^2}{b_{e_*\alpha}^2}
 c_\beta^2=\frac{\kappa_g}{2}c_\beta^2.            \tag{34.21}
\]
The cancellation is therefore genuine and occurs between two adjacent
Haar scales on one physical edge.

\begin{countertheorem}[Scalar Haar updates do not inherit the V29
network closure]
\label{ctr:v17v34-scalar}
The V29 minimizer-anchor proof cannot be transferred coefficient by
coefficient to the Haar basis using only the scalar cavity minimum
\(M_\alpha\).  The required anchor-center quadratic cost is independent
of that minimum in the admissible cavity family (34.17).  The zero
anchor avoids this cost but saturates the sharp mismatch parameter at
\(A_g\) by (34.6).
\end{countertheorem}

\begin{proof}
The first assertion is Countertheorem
\ref{ctr:v17v34-cancellation}; the second is the limiting coefficient
following Theorem~\ref{thm:v17v34-zero}.  Both conclusions concern the
specified scalar-anchor compilers.
\end{proof}

The countertheorem does not rule out Haar coordinates.  It says that
coefficients sharing a physical edge must be updated in a block, so the
edge direction and its Gaussian-orthogonal cancellation directions are
treated together.

\begin{assessment}[What V34 closes]
\label{ass:v17v34-status}
V34 derives the full scalar Haar conditional bound and exhibits an
exact two-scale cavity that defeats minimizer-only closure.  The failure
is not caused by volume or by loss of mean zero: it occurs on one edge
between two explicit Haar coefficients.

The next viable architecture is an edge-aligned coefficient block.  Its
incidence row has squared norm exactly two by (33.13), while the
orthogonal coefficient directions retain their diagonal Gaussian
coercivity.
\end{assessment}

\begin{decision}[V35: audit and edge-aligned block rotation]
\label{dec:v17v34-next}
Perform the mandatory YM/NS audit.  Then, for one physical edge, rotate
the Haar coefficients that affect it into the normalized incidence
direction \(z_e=(Br)_e/\sqrt2\) and its orthogonal complement.  Compute
the exact Gaussian form, identify the interactions shared with other
edges, and continue to V36 before judging the checkpoint.
\end{decision}

\begin{firewall}[Claim boundary after V34]
\label{fire:v17v34}
V34 rules out two scalar-anchor proofs, not the quotient Gibbs measure
or all Haar block schemes.  It does not provide the physical RPESM
gravity action or the Einstein--SM continuum limit.
\end{firewall}

\section*{Version V34 append-only record}
\addcontentsline{toc}{section}{Version V34 append-only record}
The complete V33 source was retained byte for byte before this
continuation.  It had \(328736\) bytes, \(10129\) lines, and SHA--256
\[
 \texttt{37E44C8193F114BA981FCB7E09C01FF0EF453ED358D0E73484A51C21BAAEE036}.
\]
V34 derived the scalar Haar compiler and proved its exact cross-scale
cancellation obstruction.


\section{V35: companion audit and edge-aligned quotient rotations}
\label{sec:v17v35}

\subsection{Mandatory companion-note audit}

The active Yang--Mills note inspected at this checkpoint was
\[
 \texttt{Renewal\_Yang\_Mills\_Mass\_Gap\_Research\_Notes\_V80.tex},
                                                               \tag{35.1}
\]
with \(650260\) bytes, \(17771\) physical lines, and SHA--256
\[
 \texttt{4A239C512AA80DC81CFB862C452802D1762AF3AED4A00C3EE587EB379FB7F642}.
                                                               \tag{35.2}
\]
YM V80 performs its audit, certifies ten new sixth-atlas charts, and,
together with four inherited charts, covers fourteen of the sixteen
third-coordinate columns.  It records exact negative witnesses for the
proposed fixed splittings at the remaining centers \(222\) and \(232\),
while explicitly allowing smaller charts or rational bridges.  No
third-coordinate slab, complete cell pair, full-torus positivity,
cofinal continuum measure, or mass gap is claimed.

The active Navier--Stokes note remained
\[
 \texttt{Renewal\_NS\_Stability\_Research\_Notes\_V26.tex},     \tag{35.3}
\]
with \(278283\) bytes, \(5319\) physical lines, and SHA--256
\[
 \texttt{82C887F8C2C69E4F28FAD7C3DC8DE5B9099CB5A4BF431EBA1FFF3E91935978AD}.
                                                               \tag{35.4}
\]
Its rank-one Oseen-kernel block remains separate from the quotient
conformal problem.

\begin{audit}[V35 cross-program decision]
\label{aud:v17v35-companions}
No YM or NS theorem is imported.  The exact YM splitting failures
reinforce the present practice of distinguishing failure of a
certificate from failure of the underlying model, but provide no
gravity coefficient.
\end{audit}

\subsection{The normalized incidence direction}

Let \(G\) be a finite connected simple graph.  For an oriented edge
\(e=(v,w)\), put
\[
 g_e=\mathbf e_w-\mathbf e_v\in\mathcal H_0,
 \qquad \|g_e\|_2^2=2.                             \tag{35.5}
\]
Let \(\Psi:\mathbb R^{n-1}\to\mathcal H_0\) be any orthogonal
synthesis map, including the Haar synthesis of V33, and write
\(r=\Psi c\).  The coefficient-space incidence row is
\[
 b_e=\Psi^{\mathsf T}g_e,\qquad \|b_e\|_2^2=2.     \tag{35.6}
\]
Define
\[
 u_e=\frac{b_e}{\sqrt2}.                            \tag{35.7}
\]

\begin{theorem}[The edge rotation is a local physical move]
\label{thm:v17v35-local}
Decompose the quotient coefficients as
\[
 c=z_eu_e+w_e,\qquad w_e\perp u_e.                 \tag{35.8}
\]
Then
\[
 \Psi u_e=\frac{g_e}{\sqrt2},                     \tag{35.9}
\]
and the chosen edge difference is exactly
\[
 (Br)_e=\sqrt2\,z_e.                               \tag{35.10}
\]
Thus changing \(z_e\) changes only the two endpoint field values,
with opposite increments, and preserves the global mean-zero
constraint.
\end{theorem}

\begin{proof}
Since \(g_e\in\mathcal H_0\) and \(\Psi\Psi^{\mathsf T}\) is the
identity on \(\mathcal H_0\),
\[
 \Psi u_e
 =\frac{\Psi\Psi^{\mathsf T}g_e}{\sqrt2}
 =\frac{g_e}{\sqrt2}.                              \tag{35.11}
\]
Moreover,
\[
 (Br)_e=\langle g_e,r\rangle
 =\langle b_e,c\rangle
 =z_e\langle b_e,u_e\rangle
 =\sqrt2\,z_e,                                     \tag{35.12}
\]
because \(w_e\perp b_e\).  Equation (35.9) gives locality and mean-zero
preservation.
\end{proof}

Unlike the scalar Haar coordinate in V34, the own-edge term (35.10) has
no cavity residual.

\subsection{Exact interaction stencil}

For another oriented edge \(f\), define
\[
 \rho_{fe}=\langle g_f,\Psi u_e\rangle
 =\frac{\langle g_f,g_e\rangle}{\sqrt2}.           \tag{35.13}
\]

\begin{theorem}[Line-graph stencil]
\label{thm:v17v35-stencil}
Under a \(z_e\)-update, the edge differences have the form
\[
 (Br)_f=\xi_{f,e}+\rho_{fe}z_e,                    \tag{35.14}
\]
where
\[
 |\rho_{fe}|=
 \begin{cases}
 \sqrt2,&f=e,\\
 1/\sqrt2,&f\ne e\text{ shares one endpoint with }e,\\
 0,&f\text{ is disjoint from }e.
 \end{cases}                                      \tag{35.15}
\]
The sign in the middle case depends only on the chosen orientations.
If the maximum vertex degree is \(D\), at most \(2D-1\) physical edges
occur in the conditional stencil.  On a path, at most three occur.
\end{theorem}

\begin{proof}
Equation (35.14) is the orthogonal decomposition (35.8) pulled back to
edge \(f\).  Two signed incidence vectors have inner product two when
they are the same edge, absolute inner product one when distinct edges
share one endpoint, and inner product zero when disjoint.  This proves
(35.15).  Edge \(e\) has at most \(D-1\) other incident edges at each
endpoint.
\end{proof}

The stencil is independent of which orthonormal quotient basis was used
to construct it.  Haar coordinates supplied a convenient route to the
quotient, but the rotated update is intrinsically an edge move.

\subsection{The own edge controls the rotated minimizer}

Let
\[
 \mathcal S(e)=\{f:\rho_{fe}\ne0\}                \tag{35.16}
\]
and define, for fixed residuals \(\xi_{f,e}\),
\[
 W_e(z)=\sum_{f\in\mathcal S(e)}
 V_g(\xi_{f,e}+\rho_{fe}z).                        \tag{35.17}
\]
By (35.10), \(\xi_{e,e}=0\).

\begin{theorem}[Edge-center coercivity]
\label{thm:v17v35-center}
The function \(W_e\) has a unique minimizer \(z_e^*\).  With
\[
 M_e=\min_zW_e(z),                                 \tag{35.18}
\]
one has the exact deterministic bound
\[
 \boxed{
 (z_e^*)^2
 \leq\frac{V_g(\sqrt2z_e^*)-2}{8}
 \leq\frac{M_e-2}{8}
 \leq\frac{M_e}{8}.}                              \tag{35.19}
\]
\end{theorem}

\begin{proof}
The own-edge summand \(V_g(\sqrt2z)\) is strictly convex and coercive,
so the finite sum (35.17) is strictly convex and coercive.  At the
minimum it is one of the nonnegative summands of \(M_e\).  The elementary
bound
\[
 V_g(q)-2=2\cosh(2q)-2\geq4q^2                    \tag{35.20}
\]
with \(q=\sqrt2z_e^*\) gives (35.19).
\end{proof}

This is the estimate that fails for a scalar Haar coefficient in V34:
there the distinguished edge can be cancelled by the cavity, whereas
the rotated coordinate owns the entire edge difference.

\subsection{A finite-color update schedule}

\begin{theorem}[Disjoint edge directions by coloring]
\label{thm:v17v35-color}
The edges of a graph of maximum degree \(D\) can be greedily colored
with at most \(2D-1\) colors so that each color is a matching.  For one
color, the directions
\[
 \left\{\frac{g_e}{\sqrt2}:e\text{ in the color}\right\}        \tag{35.21}
\]
are orthonormal and their updates have disjoint vertex supports.  Across
all colors, these directions span \(\mathcal H_0\).
\end{theorem}

\begin{proof}
The line graph has maximum degree at most \(2D-2\), so greedy vertex
coloring of the line graph uses at most \(2D-1\) colors.  Edges in one
color share no vertex, hence their incidence vectors are orthogonal and
their supports are disjoint.  Since the original graph is connected,
its incidence vectors span \(\mathbf1^\perp\).
\end{proof}

The collection over all colors is a spanning frame rather than one
simultaneous orthonormal coordinate basis.  A color-by-color Gibbs or
heat-bath schedule is nevertheless well defined because each update is
a one-dimensional orthogonal decomposition of the current quotient
field.

\begin{assessment}[What V35 closes]
\label{ass:v17v35-status}
The mandatory companion audit is complete.  V35 constructs the
edge-aligned quotient move, proves its bounded line-graph stencil, and
derives the coercive bound (35.19) that removes the exact cancellation
found in V34.  A finite matching-color schedule supplies disjoint
parallel updates and spans the full mean-zero field.

The conditional Laplace ratio and its Higgs debit have not yet been
compiled; those determine whether the new coercivity leaves a strict
Perron margin.
\end{assessment}

\begin{decision}[V36: edge-aligned conditional certificate]
\label{dec:v17v35-next}
Apply the V29 numerator-minimum argument to \(W_e(z)\), use (35.19) for
the Gaussian anchor cost, include all at-most-\((2D-1)\) hopping edges,
and derive the exact inequality on \(J,\kappa_g,\theta\) required for a
strict network-minimum transfer.  Continue past the audit before
assessing the block architecture.
\end{decision}

\begin{firewall}[Claim boundary after V35]
\label{fire:v17v35}
The rotation and coloring are finite-graph geometric results.  They do
not yet prove contraction of the color schedule, construct the physical
RPESM gravity measure, or establish the continuum limit.
\end{firewall}

\section*{Version V35 append-only record}
\addcontentsline{toc}{section}{Version V35 append-only record}
The complete V34 source was retained byte for byte before this
continuation.  It had \(337229\) bytes, \(10385\) lines, and SHA--256
\[
 \texttt{5D1C59CD3914465BF21FC3237B70FE168D00995247FB6A170DAC34DF5FF302F0}.
\]
V35 completed the audit and constructed coercive edge-aligned quotient
updates.


\section{V36: edge-aligned conditional closure and the induced-matching schedule}
\label{sec:v17v36}

\subsection{The sharp own-edge constant}

Define
\[
 C_{\rm edge}
 =\sup_{z\in\mathbb R}
 \frac{z^2}{V_g(\sqrt2z)}.                         \tag{36.1}
\]

\begin{lemma}[Exact evaluation of the coercive ratio]
\label{lem:v17v36-constant}
There is a unique \(x_0>0\) satisfying
\[
 x_0\tanh x_0=2,                                   \tag{36.2}
\]
and
\[
 \boxed{
 C_{\rm edge}=\frac{x_0^2}{16\cosh x_0}<\frac18.} \tag{36.3}
\]
Consequently, the minimizer from V35 obeys the sharper bound
\[
 (z_e^*)^2\leq C_{\rm edge}M_e.                   \tag{36.4}
\]
\end{lemma}

\begin{proof}
Put \(x=2\sqrt2|z|\).  Since
\(V_g(\sqrt2z)=2\cosh x\), the ratio in (36.1) is
\[
 f(x)=\frac{x^2}{16\cosh x},\qquad x\geq0.         \tag{36.5}
\]
Away from zero, the sign of \(f'(x)\) is the sign of
\(2-x\tanh x\).  The function \(x\tanh x\) is strictly increasing
from zero to infinity, so (36.2) has one solution and gives the global
maximum.  The strict comparison with \(1/8\) also follows from
\(V_g(\sqrt2z)-2\geq8z^2\).  Finally, the own-edge summand at the
minimizer is no larger than \(M_e\), proving (36.4).
\end{proof}

\subsection{Conditional moment along an edge direction}

Use the edge rotation of V35.  Conditional on its orthogonal
complement, let
\[
 \Phi_e(z)=\kappa_gz^2+A_gW_e(z)+H_e(z),
 \qquad A_g=\frac J2,                              \tag{36.6}
\]
where \(H_e\geq0\) contains the diagonal Higgs hopping on
\(\mathcal S(e)\).  Put
\[
 h_{v,e}=\sum_{\substack{f\in\mathcal S(e)\\f\ni v}}\chi_f,
 \qquad
 H_e^{\Sigma}=\sum_vh_{v,e}=2\sum_{f\in\mathcal S(e)}\chi_f.    \tag{36.7}
\]
If \(\deg G\leq D\) and \(\chi_f\leq\chi_*\), then
\[
 H_e^{\Sigma}\leq2(2D-1)\chi_*.                  \tag{36.8}
\]

\begin{theorem}[Edge-aligned conditional Laplace bound]
\label{thm:v17v36-conditional}
Assume
\[
 0<\theta<A_g.                                     \tag{36.9}
\]
For every \(\varepsilon,\omega,\gamma>0\), there is a finite constant
\(C_e\), independent of the residual edge differences and Higgs values,
such that
\[
\begin{aligned}
 \log\mathbb E_e e^{\theta W_e(z_e)}
 \leq{}&C_e+q_e(\varepsilon,\omega,\gamma)M_e\\
 &+e^{\sqrt2\varepsilon}\gamma
 \sum_vh_{v,e}V_H(y_v),                            \tag{36.10}
\end{aligned}
\]
where
\[
\boxed{
 q_e(\varepsilon,\omega,\gamma)
 =\theta+A_g(e^{2\sqrt2\varepsilon}-1)
 +(1+\omega)\kappa_gC_{\rm edge}
 +\frac{e^{\sqrt2\varepsilon}\chi_*}{4\gamma}.}  \tag{36.11}
\]
All nonconstant coefficients are independent of the volume and depend
on the graph only through \(D\) and the local hopping bounds.
\end{theorem}

\begin{proof}
The numerator after insertion has mismatch coefficient
\(A_g-\theta>0\).  Since \(W_e\geq M_e\) and \(H_e\geq0\), it is at
most
\[
 e^{-(A_g-\theta)M_e}
 \int_{\mathbb R}e^{-\kappa_gz^2}\,dz.             \tag{36.12}
\]

For the denominator, use
\(|z-z_e^*|\leq\varepsilon\).  The largest incidence amplitude in
(35.15) is \(\sqrt2\), so
\[
 W_e(z)\leq e^{2\sqrt2\varepsilon}M_e.             \tag{36.13}
\]
The Gaussian anchor cost satisfies
\[
 \kappa_gz^2
 \leq(1+\omega)\kappa_g(z_e^*)^2+C_0
 \leq(1+\omega)\kappa_gC_{\rm edge}M_e+C_0.       \tag{36.14}
\]
At the minimizing edge values \(q_f^*\), the hopping terms on the
anchor are at most
\[
\begin{aligned}
 e^{\sqrt2\varepsilon}
 \sum_{f=(v,w)\in\mathcal S(e)}\chi_f
 \bigl(a_ve^{q_f^*}+a_we^{-q_f^*}\bigr)
 \leq{}&e^{\sqrt2\varepsilon}\gamma
 \sum_vh_{v,e}a_v^2\\
 &+\frac{e^{\sqrt2\varepsilon}\chi_*}{4\gamma}M_e.
                                                               \tag{36.15}
\end{aligned}
\]
Here the second line uses Young's inequality and
\(\sum_{f\in\mathcal S(e)}V_g(q_f^*)=M_e\).  The anchor has length
\(2\varepsilon\).  Dividing (36.12) by its denominator lower bound
gives (36.10)--(36.11).
\end{proof}

\subsection{Exact parameter criterion for the gravity transfer}

Assign exponential parameter \(\Lambda\) to the cavity-minimum sector.
Since \(M_e\leq W_e\leq\sum_fV_g((Br)_f)\), the global mismatch action
permits every \(\Lambda\leq A_g\) as a finite-volume integrability
parameter.

\begin{theorem}[Nonempty edge-aligned gravity margin]
\label{thm:v17v36-margin}
There exist
\[
 \varepsilon,\omega,\gamma>0,qquad
 q_e(\varepsilon,\omega,\gamma)<\Lambda\leq A_g  \tag{36.16}
\]
if and only if
\[
\boxed{
 \theta+\kappa_gC_{\rm edge}<\frac J2.}            \tag{36.17}
\]
When (36.17) holds, choose any
\[
 \theta+\kappa_gC_{\rm edge}<\Lambda<A_g,         \tag{36.18}
\]
then take \(\varepsilon,\omega>0\) small enough that
\[
 \theta+A_g(e^{2\sqrt2\varepsilon}-1)
 +(1+\omega)\kappa_gC_{\rm edge}<\Lambda,         \tag{36.19}
\]
and finally take
\[
 \gamma>
 \frac{e^{\sqrt2\varepsilon}\chi_*}
 {4\left[
 \Lambda-\theta-A_g(e^{2\sqrt2\varepsilon}-1)
 -(1+\omega)\kappa_gC_{\rm edge}
 \right]}.                                        \tag{36.20}
\]
\end{theorem}

\begin{proof}
The construction (36.18)--(36.20) proves sufficiency.  Conversely,
every positive \(\varepsilon,\omega,\gamma\) makes (36.11) strictly
larger than \(\theta+\kappa_gC_{\rm edge}\).  If it is smaller than
some \(\Lambda\leq A_g\), condition (36.17) follows.
\end{proof}

The strict condition contains a genuine on-site-to-edge ratio.  It is
independent of the volume and maximum degree; degree enters only the
Higgs coefficient sum (36.8).

\subsection{Correction of the simultaneous-update coloring}

V35 proved that an ordinary matching gives orthogonal, vertex-disjoint
directions.  That is sufficient for sequential updates, but not for
factorized simultaneous heat baths: an edge joining the endpoints of
two matching edges would depend on both update variables.  The correct
parallel schedule uses induced matchings.

\begin{theorem}[Finite induced-matching schedule]
\label{thm:v17v36-induced}
Let \(L(G)\) be the line graph.  A greedy coloring of \(L(G)^2\) uses at
most
\[
 C_D^{\rm edge}=(2D-2)^2+1                        \tag{36.21}
\]
colors.  Each color is an induced matching in \(G\).  For edges in one
color:
\begin{enumerate}[label=\textup{(\roman*)},leftmargin=3em]
\item their normalized incidence directions are orthonormal;
\item no physical edge is affected by two update variables;
\item the quotient Gaussian and every edge interaction split into a
product of one-dimensional conditional factors.
\end{enumerate}
Thus all edge-aligned heat baths in one color may be performed in
parallel, and a sweep through the colors preserves the finite-volume
quotient Gibbs measure.
\end{theorem}

\begin{proof}
The line graph has maximum degree at most \(2D-2\), so its square has
maximum degree at most \((2D-2)^2\); greedy coloring gives (36.21).
Same-color vertices in \(L(G)\) have distance at least three.  Their
edges in \(G\) share no endpoint, and no third physical edge shares an
endpoint with both.  This proves (i)--(ii).  Orthogonality removes
Gaussian cross terms, while (ii) assigns each edge action to at most one
update variable, proving (iii).  Each factor is the exact conditional
law along its orthogonal direction, so its heat-bath kernel preserves
the Gibbs measure; products and finite compositions preserve it as
well.
\end{proof}

Because all graph edges occur in the coloring and their incidence
vectors span \(\mathcal H_0\), the schedule reaches every quotient
direction.  Establishing a uniform mixing or contraction rate still
requires combining (36.10) with the Higgs row and with influence between
successive edge colors.

\begin{assessment}[What V36 closes]
\label{ass:v17v36-status}
V36 continues substantively after the V35 audit.  It derives the full
edge-aligned conditional moment estimate, evaluates the optimal
own-edge coercive constant, and proves the exact gravity-only
feasibility condition (36.17).  It also corrects the update coloring
from matchings to induced matchings and constructs a bounded-color
finite-volume heat-bath schedule that preserves the quotient measure.

The remaining analytic cycle is explicit: increasing \(\gamma\) makes
the gravity transfer contract but enlarges the Higgs-quartic row in
(36.10).  Inter-color influences and the physical gravitational action
are also unresolved.
\end{assessment}

\begin{decision}[V37: optimize Higgs feedback and color overlap]
\label{dec:v17v36-next}
Combine the reverse coefficient
\(e^{\sqrt2\varepsilon}\gamma H_e^{\Sigma}/\Theta_H\)
with the V21 forward geometric-edge row.  Add the finite
\(C_D^{\rm edge}\)-color overlap as a block-scan factor and derive the
exact optimized Perron condition in \(\gamma\) and the Higgs Young
parameter.  Separate invariance of the heat bath from any claimed
uniform mixing rate.
\end{decision}

\begin{firewall}[Claim boundary after V36]
\label{fire:v17v36}
The conditional certificate and invariant color schedule concern the
declared quotient conformal model.  They do not prove uniform mixing,
an infinite-volume Gibbs limit, identification with the RPESM gravity
action, or the Einstein--SM continuum limit.
\end{firewall}

\section*{Version V36 append-only record}
\addcontentsline{toc}{section}{Version V36 append-only record}
The complete V35 source was retained byte for byte before this
continuation.  It had \(346111\) bytes, \(10641\) lines, and SHA--256
\[
 \texttt{08BE4E1E8BF5C3C39D7E48C16FC0B67683C2C6C2A064793BA1206BB6213B69E8}.
\]
V36 proved the edge-aligned conditional certificate and corrected
induced-matching schedule.


\section{V37: the typed cavity no-go and uniform quotient Langevin contraction}
\label{sec:v17v37}

\subsection{Why the two gravity parameters cannot be collapsed}

V36 uses two distinct gravity objects:
\begin{enumerate}[label=\textup{(\roman*)},leftmargin=3em]
\item the updated edge-stencil energy \(W_e\), whose conditional moment
is taken at parameter \(\theta\);
\item the cavity minimum \(M_e\), assigned a larger admissible parameter
\(\Lambda\).
\end{enumerate}
The strict margin in (36.16) requires
\[
 \theta<q_e<\Lambda\leq A_g.                       \tag{37.1}
\]
The first inequality follows from (36.11), even before the strict
choices are imposed.

One always has the deterministic domination
\[
 M_e\leq W_e.                                      \tag{37.2}
\]
If this is used as the only return estimate from the cavity type to the
updated-energy type, the normalized two-type comparison matrix is
\[
 \mathbb B_{WM}
 =\begin{pmatrix}
 0&q_e/\Lambda\\
 \Lambda/\theta&0
 \end{pmatrix}.                                   \tag{37.3}
\]

\begin{countertheorem}[Typed cavity cycle cannot contract]
\label{ctr:v17v37-typed}
For every parameter choice in the V36 conditional certificate,
\[
 \boxed{
 \rho(\mathbb B_{WM})
 =\sqrt{\frac{q_e}{\theta}}>1.}                    \tag{37.4}
\]
Adding nonnegative Higgs or inter-color influences cannot reduce this
spectral radius.  Therefore the V36 inequality is not, by itself, a
Dobrushin contraction after the two gravity types are composed through
(37.2).
\end{countertheorem}

\begin{proof}
The eigenvalues of (37.3) are
\(\pm\sqrt{q_e/\theta}\).  Formula (36.11) gives
\[
 q_e>\theta+\kappa_gC_{\rm edge}>\theta.           \tag{37.5}
\]
For nonnegative matrices, entrywise enlargement cannot decrease the
Perron root, proving the final assertion.
\end{proof}

This result does not contradict Theorem~\ref{thm:v17v36-margin}.
That theorem gives a strict one-way conditional transfer into a more
integrable boundary type.  Countertheorem~\ref{ctr:v17v37-typed} shows
that deterministic return through (37.2) spends more than the entire
gain.

\subsection{The color count is not a substitute for contraction}

The induced-matching schedule of V36 proves invariance of the
finite-volume quotient measure.  Its number of colors is a bounded
scheduling constant.  Multiplying or dividing (37.3) by that color
count has no mathematical justification: a composition of invariant
heat-bath kernels remains invariant, but its contraction rate depends
on how the conditional directions interact.  Thus no uniform mixing
claim is made from coloring alone.

The uniformly convex quotient potential of V32 provides a different,
closed mixing mechanism that has no cavity return row.

\subsection{Projected Langevin dynamics}

Fix the rescaled Higgs field \(y\) and let \(\Phi_y\) be the quotient
potential (32.6) on \(\mathcal H_0\).  Let
\(B_t^0\) be standard Brownian motion on this \((n-1)\)-dimensional
Euclidean space.  Consider
\[
 dr_t=-\nabla_{\mathcal H_0}\Phi_y(r_t)\,dt
 +\sqrt2\,dB_t^0.                                  \tag{37.6}
\]

\begin{theorem}[Pathwise volume-uniform contraction]
\label{thm:v17v37-langevin}
Let \(r_t\) and \(\widetilde r_t\) solve (37.6) with the same Brownian
path and arbitrary initial conditions.  Then
\[
 \boxed{
 \|r_t-\widetilde r_t\|_2
 \leq e^{-2\kappa_gt}
 \|r_0-\widetilde r_0\|_2}                        \tag{37.7}
\]
for every \(t\geq0\), almost surely.  The rate is independent of the
volume, graph degree, mismatch coupling, hopping coefficients, and
fixed Higgs values.
\end{theorem}

\begin{proof}
Put \(d_t=r_t-\widetilde r_t\).  The synchronous noises cancel, so
\[
 \frac{d}{dt}\|d_t\|_2^2
 =-2\langle d_t,
 \nabla\Phi_y(r_t)-\nabla\Phi_y(\widetilde r_t)\rangle.          \tag{37.8}
\]
The Hessian bound (32.12) and integration along the segment between the
two points give
\[
 \langle r-\widetilde r,
 \nabla\Phi_y(r)-\nabla\Phi_y(\widetilde r)\rangle
 \geq2\kappa_g\|r-\widetilde r\|_2^2.              \tag{37.9}
\]
Hence
\[
 \frac{d}{dt}\|d_t\|_2^2
 \leq-4\kappa_g\|d_t\|_2^2.                       \tag{37.10}
\]
Gronwall's inequality gives (37.7).
\end{proof}

The smooth convex drift is coercive; standard finite-dimensional
monotone-drift approximation gives a unique nonexplosive strong solution
to (37.6).  The next theorem records the consequences needed here.

\begin{theorem}[Unique invariant quotient law and Wasserstein mixing]
\label{thm:v17v37-invariant}
The probability measure
\[
 \mu_y(dr)=Z_y^{-1}e^{-\Phi_y(r)}\,d\mathcal H_0(r)              \tag{37.11}
\]
is reversible and is the unique invariant probability measure of
(37.6).  If \(P_t^y\) is the Markov semigroup, then for all probability
measures \(\nu_1,\nu_2\) on \(\mathcal H_0\) with finite second moments,
\[
 \boxed{
 W_2(\nu_1P_t^y,\nu_2P_t^y)
 \leq e^{-2\kappa_gt}W_2(\nu_1,\nu_2).}           \tag{37.12}
\]
In particular,
\[
 W_2(\nu P_t^y,\mu_y)
 \leq e^{-2\kappa_gt}W_2(\nu,\mu_y).              \tag{37.13}
\]
\end{theorem}

\begin{proof}
Integration by parts on \(\mathcal H_0\) shows that the generator
\[
 \mathcal L_y=\Delta_{\mathcal H_0}
 -\nabla\Phi_y\cdot\nabla                         \tag{37.14}
\]
is symmetric in \(L^2(\mu_y)\), so \(\mu_y\) is reversible and
invariant.  Couple optimally distributed initial conditions and then
use the same Brownian path.  Taking the second moment in (37.7) and the
infimum over initial couplings gives (37.12).  A second invariant law
with finite second moment would have distance to \(\mu_y\) multiplied
by at most \(e^{-2\kappa_gt}\) for every \(t\), hence distance zero.
Strong convexity gives Gaussian tails and finite second moments for
every invariant law, proving uniqueness and (37.13).
\end{proof}

\begin{corollary}[Spectral and entropy constants]
\label{cor:v17v37-functional}
The quotient law satisfies the Poincare inequality already recorded in
(32.13), with constant \((2\kappa_g)^{-1}\).  By the
Bakry--Emery criterion applied to (32.12), it also satisfies the
logarithmic Sobolev inequality
\[
 \operatorname{Ent}_{\mu_y}(f^2)
 \leq\frac1{\kappa_g}
 \int\|\nabla f\|_2^2\,d\mu_y.                    \tag{37.15}
\]
Both constants are uniform in the finite volume and in the fixed Higgs
configuration.
\end{corollary}

\begin{proof}
The Poincare statement is Corollary~\ref{cor:v17v32-bl}.  In the
normalization
\(\operatorname{Ent}(f^2)\leq2\rho^{-1}\int|\nabla f|^2\), the
Bakry--Emery curvature is \(\rho=2\kappa_g\) by (32.12), giving
(37.15).
\end{proof}

\subsection{What this changes in the SM strategy}

The edge heat baths remain useful finite-volume samplers and local
conditional diagnostics.  They are no longer the proposed proof of
uniform gravity mixing.  The synchronous quotient Langevin dynamics
gives that result directly for every fixed Higgs field.

The next coupling question is therefore analytic rather than
combinatorial: integrate the uniformly mixing quotient gravity field
and control the resulting effective Higgs action.  The covariance term
in the Hessian of that effective action can be bounded by the
Brascamp--Lieb inverse Hessian from (32.16).

\begin{assessment}[What V37 closes]
\label{ass:v17v37-status}
V37 proves that the naive typed return from \(M_e\) to \(W_e\) cannot
contract, preventing an invalid use of the V36 one-way margin.  It then
replaces that failed architecture by a volume-uniform synchronous
Langevin contraction, unique invariant quotient law, Wasserstein
mixing estimate, and logarithmic Sobolev bound for every fixed Higgs
field in the declared model.

This closes conditional gravity mixing for the model, but not the joint
Higgs--gravity measure or the physical Einstein action.
\end{assessment}

\begin{decision}[V38: effective Higgs curvature after gravity integration]
\label{dec:v17v37-next}
Define the quotient gravitational free energy as a function of the
rescaled Higgs variables.  Differentiate it twice, express the negative
term as a gravity covariance, and use (32.16) to bound that covariance
by explicit hopping derivatives.  Compare the resulting debit with the
local Higgs quartic coercivity and record the exact parameter condition
for convexity or the precise obstruction if it fails.
\end{decision}

\begin{firewall}[Claim boundary after V37]
\label{fire:v17v37}
The Langevin result is conditional on the declared difference-only
conformal action and fixed Higgs variables.  It does not construct the
joint infinite-volume SM--gravity measure, identify the RPESM
gravitational action, or prove the Einstein--SM continuum limit.
\end{firewall}

\section*{Version V37 append-only record}
\addcontentsline{toc}{section}{Version V37 append-only record}
The complete V36 source was retained byte for byte before this
continuation.  It had \(355554\) bytes, \(10915\) lines, and SHA--256
\[
 \texttt{727353E99CE812217B281CD2F000359AD7FC3480F1187FDFC302E3EB32E246AC}.
\]
V37 proved the typed cavity no-go and the uniform quotient Langevin
alternative.


\section{V38: effective Higgs curvature after quotient-gravity integration}
\label{sec:v17v38}

\subsection{The gravitational free energy}

Realify the finite-dimensional Higgs fibres so that inner products are
real Euclidean inner products.  For an oriented edge \(e=(v,w)\), put
\[
 d_e=(Br)_e=r_w-r_v                              \tag{38.1}
\]
and
\[
 T_e(r,y)=\chi_e\left(
 e^{d_e}\|y_v\|^2+e^{-d_e}\|y_w\|^2
 \right).                                         \tag{38.2}
\]
Let
\[
 \Phi_y(r)=\Phi_0(r)+\sum_eT_e(r,y)                \tag{38.3}
\]
be the quotient action (32.6), and define the normalized gravity free
energy
\[
 \mathcal F_g(y)
 =-\log\frac{Z_g(y)}{Z_g(0)},
 \qquad
 Z_g(y)=\int_{\mathcal H_0}e^{-\Phi_y(r)}
 \,d\mathcal H_0(r).                              \tag{38.4}
\]

\begin{lemma}[Positivity and smoothness]
\label{lem:v17v38-basic}
The function \(\mathcal F_g\) is smooth and
\[
 \mathcal F_g(y)\geq0,qquad \mathcal F_g(0)=0.   \tag{38.5}
\]
\end{lemma}

\begin{proof}
Every \(T_e\) is nonnegative, so
\(Z_g(y)\leq Z_g(0)\), proving (38.5).  On compact Higgs sets, all
Higgs derivatives of the integrand are polynomials in
\(e^{\pm d_e}\) times a Gaussian and the nonnegative cosh action.  They
are dominated by an integrable function.  Differentiation under the
finite-dimensional integral to every order is therefore valid.
\end{proof}

\subsection{Exact first and second variations}

For a Higgs variation \(h=(h_v)\), define
\[
 G_h(r,y)=2\sum_{e=(v,w)}\chi_e\left[
 e^{d_e}\langle y_v,h_v\rangle
 +e^{-d_e}\langle y_w,h_w\rangle
 \right],                                         \tag{38.6}
\]
\[
 K_h(r)=2\sum_{e=(v,w)}\chi_e\left[
 e^{d_e}\|h_v\|^2+e^{-d_e}\|h_w\|^2
 \right].                                         \tag{38.7}
\]
Let \(\mathbb E_y\) and \(\operatorname{Var}_y\) refer to the
probability measure proportional to \(e^{-\Phi_y}\).

\begin{theorem}[Effective free-energy Hessian identity]
\label{thm:v17v38-hessian}
For every \(y,h\),
\[
 D\mathcal F_g(y)[h]=\mathbb E_yG_h,               \tag{38.8}
\]
and
\[
 \boxed{
 D^2\mathcal F_g(y)[h,h]
 =\mathbb E_yK_h-\operatorname{Var}_y(G_h).}       \tag{38.9}
\]
\end{theorem}

\begin{proof}
The directional derivatives of \(\Phi_y\) are
\(D_y\Phi_y[h]=G_h\) and
\(D_y^2\Phi_y[h,h]=K_h\).  Differentiating
\(-\log Z_g(y)\) once gives the expectation in (38.8).
Differentiating the expectation gives the expected second derivative
minus the covariance of the first derivatives, which is (38.9).
\end{proof}

The covariance is the only possible negative curvature created by
integrating gravity.

\subsection{An edge-weighted Brascamp--Lieb bound}

Define
\[
 p_e=\chi_e\left(
 e^{d_e}\|y_v\|^2+e^{-d_e}\|y_w\|^2
 \right)=T_e(r,y),                                \tag{38.10}
\]
\[
 k_e=2\chi_e\left(
 e^{d_e}\langle y_v,h_v\rangle
 -e^{-d_e}\langle y_w,h_w\rangle
 \right).                                         \tag{38.11}
\]
Then
\[
 \nabla_rG_h=B^{\mathsf T}k.                      \tag{38.12}
\]

\begin{lemma}[Weighted inverse-Hessian contraction]
\label{lem:v17v38-weighted}
Let \(D=\operatorname{diag}(p_e)\), interpreting terms with
\(p_e=0\) by continuity.  Then pointwise
\[
 (B^{\mathsf T}k)^{\mathsf T}
 (\nabla^2\Phi_y)^{-1}(B^{\mathsf T}k)
 \leq\sum_e\frac{k_e^2}{p_e}.                     \tag{38.13}
\]
\end{lemma}

\begin{proof}
The quotient Hessian (32.11) obeys
\(\nabla^2\Phi_y\succeq B^{\mathsf T}DB\).  For positive
\(p_e\), use the variational representation
\[
 x^{\mathsf T}H^{-1}x
 =\sup_z\bigl(2\langle x,z\rangle-z^{\mathsf T}Hz\bigr).       \tag{38.14}
\]
With \(x=B^{\mathsf T}k\) and
\(H=\nabla^2\Phi_y\), the right side is at most
\[
 \sup_z\left(
 2\langle k,Bz\rangle-\langle Bz,D Bz\rangle
 \right)
 \leq\langle k,D^{-1}k\rangle.                   \tag{38.15}
\]
Add a positive diagonal regularization and pass to the limit when some
\(p_e=0\); then (38.11) also gives \(k_e=0\) on those edges.
\end{proof}

\begin{theorem}[Universal relative semiconvexity]
\label{thm:v17v38-semiconvex}
Put
\[
 S_h(r)=\sum_{e=(v,w)}\chi_e\left[
 e^{d_e}\|h_v\|^2+e^{-d_e}\|h_w\|^2
 \right].                                         \tag{38.16}
\]
Then
\[
 \operatorname{Var}_y(G_h)
 \leq4\mathbb E_yS_h,                              \tag{38.17}
\]
and therefore
\[
 \boxed{
 -2\mathbb E_yS_h
 \leq D^2\mathcal F_g(y)[h,h]
 \leq2\mathbb E_yS_h.}                            \tag{38.18}
\]
\end{theorem}

\begin{proof}
Brascamp--Lieb, (38.12), and Lemma
\ref{lem:v17v38-weighted} give
\[
 \operatorname{Var}_y(G_h)
 \leq\mathbb E_y\sum_e\frac{k_e^2}{p_e}.          \tag{38.19}
\]
For each edge, Cauchy--Schwarz in the direct sum of the two endpoint
fibres gives
\[
\begin{aligned}
 &\left(
 e^{d_e}\langle y_v,h_v\rangle
 -e^{-d_e}\langle y_w,h_w\rangle
 \right)^2\\
 &\quad\leq
 \left(e^{d_e}\|y_v\|^2+e^{-d_e}\|y_w\|^2\right)
 \left(e^{d_e}\|h_v\|^2+e^{-d_e}\|h_w\|^2\right).            \tag{38.20}
\end{aligned}
\]
Thus \(k_e^2/p_e\) is at most four times the corresponding summand of
(38.16), proving (38.17).  Since
\(\mathbb E_yK_h=2\mathbb E_yS_h\), identity (38.9) gives both sides of
(38.18).
\end{proof}

The estimate is uniform in volume but does not assert that
\(\mathcal F_g\) is convex for arbitrary Higgs backgrounds.

\subsection{Exact positivity at the symmetric Higgs origin}

At \(y=0\), one has \(G_h=0\), so the covariance in (38.9) vanishes.
The unperturbed quotient law is invariant under \(r\mapsto-r\).  Define
\[
 m_e=\mathbb E_0e^{d_e}
 =\mathbb E_0e^{-d_e}\geq1.                        \tag{38.21}
\]
The inequality follows from Jensen and \(\mathbb E_0d_e=0\).

Restore the conformal-independent cross term of the geometric Higgs
edge,
\[
 C_H(y)=-2\sum_{e=(v,w)}\chi_e
 \operatorname{Re}\langle y_w,U_ey_v\rangle,       \tag{38.22}
\]
where \(U_e\) is unitary.

\begin{theorem}[Positive effective geometric edge at the origin]
\label{thm:v17v38-origin}
For every Higgs variation \(h\),
\[
\begin{aligned}
 D^2(\mathcal F_g+C_H)(0)[h,h]
 =2\sum_{e=(v,w)}\chi_e\bigl[
 &m_e(\|h_v\|^2+\|h_w\|^2)\\
 &-2\operatorname{Re}\langle h_w,U_eh_v\rangle
 \bigr]                                           \tag{38.23}
\end{aligned}
\]
and every edge bracket is nonnegative.  More precisely,
\[
\begin{aligned}
 &m_e(\|h_v\|^2+\|h_w\|^2)
 -2\operatorname{Re}\langle h_w,U_eh_v\rangle\\
 &\quad=(m_e-1)(\|h_v\|^2+\|h_w\|^2)
 +\|h_w-U_eh_v\|^2.                               \tag{38.24}
\end{aligned}
\]
\end{theorem}

\begin{proof}
At the origin, (38.9) and (38.7) give the diagonal part of (38.23),
using (38.21).  Twice differentiating (38.22) gives the cross part.
Unitarity expands the final square in (38.24); its two terms are
nonnegative because \(m_e\geq1\).
\end{proof}

Thus quotient-gravity integration preserves the positive-square nature
of the geometric Higgs edge at the symmetric point and adds the
fluctuation mass \(m_e-1\).

\subsection{Explicit general-background acquisition target}

Define
\[
 \overline\omega_v(y)
 =\mathbb E_y\sum_{w\sim v}\chi_{vw}e^{r_w-r_v}.   \tag{38.25}
\]
Then
\[
 \mathbb E_yS_h
 =\sum_v\overline\omega_v(y)\|h_v\|^2.            \tag{38.26}
\]
If the local Higgs potential is
\[
 U_H(y)=\sum_v\left(m_H^2\|y_v\|^2
 +\lambda_H\|y_v\|^4\right),qquad\lambda_H>0,   \tag{38.27}
\]
then (38.18) gives the sufficient lower bound
\[
\begin{aligned}
 D^2(U_H+\mathcal F_g+C_H)(y)[h,h]
 \geq\sum_v\bigl[
 2m_H^2+4\lambda_H\|y_v\|^2
 -2\overline\omega_v(y)-2h_v^\chi
 \bigr]\|h_v\|^2,                                \tag{38.28}
\end{aligned}
\]
where
\[
 h_v^\chi=\sum_{w\sim v}\chi_{vw}.                \tag{38.29}
\]
The omitted term
\(8\lambda_H\sum_v\langle y_v,h_v\rangle^2\) is nonnegative.

Equation (38.28) is a sufficient convexity test, not a claim that the
physical broken-symmetry Higgs potential is globally convex.  The new
coefficient acquisition problem is to control
\(\overline\omega_v(y)\) sharply enough to combine effective curvature
with quartic coercivity without assuming convexity near every Higgs
background.

\begin{assessment}[What V38 closes]
\label{ass:v17v38-status}
V38 differentiates the quotient gravitational free energy exactly,
isolates its covariance debit, and bounds that debit by the positive
edge-weighted Hessian.  At the symmetric Higgs origin it proves exact
positive semidefiniteness of every integrated geometric edge and
identifies a nonnegative fluctuation mass.

For general Higgs fields the result is a volume-uniform relative
semiconvexity estimate, not global convexity.  The outstanding scalar
quantities are the tilted endpoint weights
\(\overline\omega_v(y)\).
\end{assessment}

\begin{decision}[V39: tilted hopping-weight bounds and radial coercivity]
\label{dec:v17v38-next}
Scale the Higgs field by a common parameter and use the monotonicity of
the tilted partition function to bound
\(\sum_v\|y_v\|^2\overline\omega_v(y)\) by its
unperturbed value.  Combine the quotient logarithmic Sobolev bound with
the symmetry of the unperturbed measure to obtain an explicit
volume-uniform radial estimate and compare it with the Higgs quartic.
\end{decision}

\begin{firewall}[Claim boundary after V38]
\label{fire:v17v38}
The effective-curvature formulas concern the declared finite-volume
quotient model.  They do not prove global convexity of the physical
Higgs sector, an infinite-volume joint measure, identification of the
RPESM gravitational action, or the Einstein--SM continuum limit.
\end{firewall}

\section*{Version V38 append-only record}
\addcontentsline{toc}{section}{Version V38 append-only record}
The complete V37 source was retained byte for byte before this
continuation.  It had \(364539\) bytes, \(11158\) lines, and SHA--256
\[
 \texttt{6ABD29259845A5499FF5F5F294A73322AD6F9BE90649061680E4F15CF76ACD36}.
\]
V38 derived the effective Higgs Hessian and its edge-weighted covariance
bound.


\section{V39: tilted hopping bounds and finite-volume joint coercivity}
\label{sec:v17v39}

\subsection{Sub-Gaussian unperturbed edge differences}

Let \(\mu_0\) be the quotient gravity law at \(y=0\).  It is invariant
under \(r\mapsto-r\), so
\[
 \mathbb E_0d_e=0.                                 \tag{39.1}
\]

\begin{theorem}[Uniform edge moment-generating bound]
\label{thm:v17v39-mgf}
For every edge \(e\) and every \(t\in\mathbb R\),
\[
 \boxed{
 \log\mathbb E_0e^{t d_e}
 \leq\frac{t^2}{2\kappa_g}.}                      \tag{39.2}
\]
In particular,
\[
 m_e=\mathbb E_0e^{d_e}=\mathbb E_0e^{-d_e}
 \leq e^{1/(2\kappa_g)}.                          \tag{39.3}
\]
\end{theorem}

\begin{proof}
The logarithmic Sobolev inequality (37.15) and the Herbst argument give,
for a centered linear functional \(\langle\ell,r\rangle\),
\[
 \log\mathbb E_0e^{t\langle\ell,r\rangle}
 \leq\frac{t^2\|\ell\|_2^2}{4\kappa_g}.          \tag{39.4}
\]
For \(d_e=\langle g_e,r\rangle\), one has
\(\|g_e\|_2^2=2\), giving (39.2).  Symmetry and \(t=\pm1\) give
(39.3).
\end{proof}

This bound uses only the on-site curvature.  The positive cosh
interaction can improve it, but no improvement is needed for the
volume-uniform estimate below.

\subsection{Monotonicity under the Higgs tilt}

For a fixed Higgs field, write
\[
 T_y(r)=\sum_{e=(v,w)}\chi_e\left(
 e^{d_e}\|y_v\|^2+e^{-d_e}\|y_w\|^2
 \right).                                         \tag{39.5}
\]
For \(0\leq s\leq1\), let
\[
 d\mu_s(r)=Z_s^{-1}e^{-\Phi_0(r)-sT_y(r)}
 \,d\mathcal H_0(r).                              \tag{39.6}
\]

\begin{lemma}[Tilt expectation is decreasing]
\label{lem:v17v39-tilt}
The function
\[
 f(s)=\mathbb E_sT_y                               \tag{39.7}
\]
is differentiable and
\[
 f'(s)=-\operatorname{Var}_s(T_y)\leq0.            \tag{39.8}
\]
\end{lemma}

\begin{proof}
Differentiate the normalized finite-dimensional integral.  The
derivative of the numerator contributes
\(-\mathbb E_sT_y^2\), and differentiating the normalization contributes
\((\mathbb E_sT_y)^2\).  Their sum is (39.8).
\end{proof}

Recall
\[
 \overline\omega_v(y)
 =\mathbb E_y\sum_{w\sim v}\chi_{vw}e^{r_w-r_v},
 \qquad
 h_v^\chi=\sum_{w\sim v}\chi_{vw}.                \tag{39.9}
\]

\begin{theorem}[Volume-uniform radial hopping bound]
\label{thm:v17v39-radial}
For every finite connected graph and every Higgs field,
\[
\boxed{
 \sum_v\|y_v\|^2\overline\omega_v(y)
 \leq e^{1/(2\kappa_g)}
 \sum_vh_v^\chi\|y_v\|^2.}                       \tag{39.10}
\]
If \(\deg G\leq D\) and \(\chi_e\leq\chi_*\), then
\[
 \sum_v\|y_v\|^2\overline\omega_v(y)
 \leq D\chi_*e^{1/(2\kappa_g)}
 \sum_v\|y_v\|^2.                                \tag{39.11}
\]
\end{theorem}

\begin{proof}
The left side of (39.10) is exactly \(\mathbb E_1T_y=f(1)\).
Lemma~\ref{lem:v17v39-tilt} gives \(f(1)\leq f(0)\).  By (39.3),
\[
\begin{aligned}
 f(0)
 &=\sum_{e=(v,w)}\chi_e\left(
 m_e\|y_v\|^2+m_e\|y_w\|^2\right)\\
 &\leq e^{1/(2\kappa_g)}
 \sum_vh_v^\chi\|y_v\|^2,                       \tag{39.12}
\end{aligned}
\]
which is (39.10).  The bounded-degree estimate gives (39.11).
\end{proof}

\subsection{Bounds on the gravitational free energy}

\begin{theorem}[Quadratic upper bound and radial derivative]
\label{thm:v17v39-free}
The normalized free energy (38.4) satisfies
\[
\boxed{
 0\leq\mathcal F_g(y)
 \leq e^{1/(2\kappa_g)}
 \sum_vh_v^\chi\|y_v\|^2.}                       \tag{39.13}
\]
Its radial derivative obeys
\[
 0\leq D\mathcal F_g(y)[y]
 =2\sum_v\|y_v\|^2\overline\omega_v(y)
 \leq2e^{1/(2\kappa_g)}
 \sum_vh_v^\chi\|y_v\|^2.                       \tag{39.14}
\]
\end{theorem}

\begin{proof}
The lower bound is (38.5).  Since
\[
 \frac{Z_g(y)}{Z_g(0)}=\mathbb E_0e^{-T_y},        \tag{39.15}
\]
Jensen's inequality gives
\[
 -\log\mathbb E_0e^{-T_y}\leq\mathbb E_0T_y.      \tag{39.16}
\]
Use (39.12) to prove the upper bound.  Formula (38.8) with \(h=y\)
gives the identity in (39.14), and Theorem
\ref{thm:v17v39-radial} gives its upper bound.  Positivity follows
also from the monotonicity of
\(\mathcal F_g(sy)\) for \(s\geq0\).
\end{proof}

Thus integrating quotient gravity creates a nonnegative effective
Higgs term, but that term grows no faster than quadratically along every
radial Higgs direction.  It cannot defeat the local quartic tail.

\subsection{Effective Higgs coercivity}

Restore the conformal-independent cross term
\[
 C_H(y)=-2\sum_{e=(v,w)}\chi_e
 \operatorname{Re}\langle y_w,U_ey_v\rangle       \tag{39.17}
\]
and the local potential
\[
 U_H(y)=\sum_v\left(m_H^2\|y_v\|^2
 +\lambda_H\|y_v\|^4\right),qquad\lambda_H>0.   \tag{39.18}
\]
Define
\[
 S_{H,\rm eff}(y)=U_H(y)+C_H(y)+\mathcal F_g(y).   \tag{39.19}
\]

\begin{theorem}[Uniform quartic lower bound]
\label{thm:v17v39-coercive}
For every finite graph,
\[
 S_{H,\rm eff}(y)
 \geq\sum_v\left[
 \lambda_H\|y_v\|^4+(m_H^2-h_v^\chi)\|y_v\|^2
 \right].                                         \tag{39.20}
\]
If \(\deg G\leq D\), \(\chi_e\leq\chi_*\), and
\[
 c_*=(D\chi_*-m_H^2)_+,                            \tag{39.21}
\]
then
\[
\boxed{
 S_{H,\rm eff}(y)
 \geq\frac{\lambda_H}{2}\sum_v\|y_v\|^4
 -\frac{|\mathcal V|c_*^2}{2\lambda_H}.}          \tag{39.22}
\]
\end{theorem}

\begin{proof}
Unitarity and Cauchy--Schwarz give, edge by edge,
\[
 -2\chi_e\operatorname{Re}\langle y_w,U_ey_v\rangle
 \geq-\chi_e(\|y_v\|^2+\|y_w\|^2).               \tag{39.23}
\]
Sum this inequality and use \(\mathcal F_g\geq0\) to obtain (39.20).
Under the bounded-degree hypotheses,
\(m_H^2-h_v^\chi\geq-c_*\).  For \(a\geq0\),
\[
 \lambda_Ha^2-c_*a
 \geq\frac{\lambda_H}{2}a^2
 -\frac{c_*^2}{2\lambda_H}.                       \tag{39.24}
\]
Apply (39.24) with \(a=\|y_v\|^2\) and sum.
\end{proof}

\begin{corollary}[Finite-volume joint normalization]
\label{cor:v17v39-normalization}
For the declared quotient conformal model, compact gauge links, and
finite-dimensional Higgs fibres, the joint finite-volume partition
function is finite for every bounded-degree graph whenever
\[
 \kappa_g>0,qquad J\geq0,qquad\lambda_H>0.        \tag{39.25}
\]
Moreover, for every \(0\leq\vartheta<\lambda_H/2\), the effective Higgs
law has the finite moment
\[
 \mathbb E\exp\left(
 \vartheta\sum_v\|y_v\|^4\right)<\infty.          \tag{39.26}
\]
\end{corollary}

\begin{proof}
Integrate the common conformal Gaussian by V32 and the quotient gravity
field by definition of \(\mathcal F_g\).  Compact gauge integration has
finite mass.  Bound (39.22) gives an integrable quartic majorant for the
remaining Higgs integral.  After insertion of (39.26), the residual
quartic coefficient is positive when
\(\vartheta<\lambda_H/2\).
\end{proof}

The factor \(1/2\) in (39.22)--(39.26) is a convenient uniform debit,
not the sharp quartic threshold.  It is enough to establish finite
volume normalization without any global convexity assumption on the
broken-symmetry Higgs potential.

\subsection{Radial curvature remains quartically controlled}

Taking \(h=y\) in (38.18) and using (39.10) gives
\[
 D^2\mathcal F_g(y)[y,y]
 \geq-2e^{1/(2\kappa_g)}
 \sum_vh_v^\chi\|y_v\|^2.                         \tag{39.27}
\]
Together with
\[
 D^2U_H(y)[y,y]
 =\sum_v\left(2m_H^2\|y_v\|^2
 +12\lambda_H\|y_v\|^4\right)                   \tag{39.28}
\]
and the quadratic cross-term bound, this proves that every negative
radial curvature contribution is at most quadratic while the local
positive curvature is quartic.

\begin{assessment}[What V39 closes]
\label{ass:v17v39-status}
V39 proves a volume-uniform tilted hopping estimate, a quadratic upper
bound for the quotient gravitational free energy, and a uniform
quartic lower bound for the effective Higgs action.  It consequently
constructs a finite normalized joint measure for every finite graph in
the declared model, without assuming global Higgs convexity.

The result is not yet an infinite-volume or continuum construction.
Uniform local moments under graph exhaustion, consistency of boundary
conditions, gauge-sector continuum control, and identification of the
physical Einstein action remain open.
\end{assessment}

\begin{decision}[V40: audit and volume-uniform local Higgs moments]
\label{dec:v17v39-next}
At the mandatory companion checkpoint, inspect the active YM and NS
notes.  Then replace the extensive constant in (39.22) by a one-site
conditional quartic estimate uniform in volume and boundary data.  Use
the tilted gravity bound only in the radial aggregate where it is
proved, and do not infer coordinatewise bounds from (39.10).  Continue
to V41 after the audit.
\end{decision}

\begin{firewall}[Claim boundary after V39]
\label{fire:v17v39}
Finite-volume normalization and radial coercivity do not prove a DLR
limit, Osterwalder--Schrader axioms, identification with the RPESM
gravitational action, or the Einstein--SM continuum limit.
\end{firewall}

\section*{Version V39 append-only record}
\addcontentsline{toc}{section}{Version V39 append-only record}
The complete V38 source was retained byte for byte before this
continuation.  It had \(374362\) bytes, \(11485\) lines, and SHA--256
\[
 \texttt{C537CD3BCEB739425C7E4C95C8D73CB5F659B178CA2CAB5C11C5273F0F9B8350}.
\]
V39 proved tilted radial control and finite-volume joint coercivity.


\section{V40: companion audit and monotonicity of the local Higgs moment}
\label{sec:v17v40}

\subsection{Mandatory companion-note audit}

The active Yang--Mills note at this checkpoint was
\[
 \texttt{Renewal\_Yang\_Mills\_Mass\_Gap\_Research\_Notes\_V82.tex},
                                                               \tag{40.1}
\]
with \(661404\) bytes, \(18068\) physical lines, and SHA--256
\[
 \texttt{86D373F6311FAD024953C28BF0E67EC8722801688AE41BB0FB38F9C4180F878B}.
                                                               \tag{40.2}
\]
YM V81 repairs the two previously obstructed columns and proves both
response signs on paired third-coordinate two-step slabs.  V82 treats
the terminal layer: four centers are inherited, eight new full-radius
charts are certified, and four proposed fixed splittings have exact
negative witnesses.  Those four terminal columns remain uncovered, so
the complete sixth cell, full-torus positivity, cofinal continuum
control, and the mass gap are still open.

The active Navier--Stokes note remained
\[
 \texttt{Renewal\_NS\_Stability\_Research\_Notes\_V26.tex},     \tag{40.3}
\]
with \(278283\) bytes, \(5319\) physical lines, and SHA--256
\[
 \texttt{82C887F8C2C69E4F28FAD7C3DC8DE5B9099CB5A4BF431EBA1FFF3E91935978AD}.
                                                               \tag{40.4}
\]
Its rank-one Oseen-kernel estimates remain unrelated to the present
local Higgs normalization problem.

\begin{audit}[V40 cross-program decision]
\label{aud:v17v40-companions}
No companion result is imported.  In particular, the incomplete YM
terminal atlas is not used as a continuum gauge premise.  V40 proceeds
only with the declared finite-volume gauge--Higgs--quotient-gravity
model.
\end{audit}

\subsection{Exact local Higgs conditional form}

Fix a vertex \(x\), all conformal differences, all compact gauge links,
and all Higgs variables except \(y_x\).  Realify the Higgs fibre as
\(\mathbb R^q\).  The terms depending on \(h=y_x\) have the form
\[
 \mathcal U_x(h)
 =\lambda_H\|h\|^4+A_x\|h\|^2-\langle L_x,h\rangle,            \tag{40.5}
\]
where
\[
 A_x=m_H^2+
 \sum_{y\sim x}\chi_{xy}e^{s_y-s_x}
 \geq m_H^2,                                       \tag{40.6}
\]
and
\[
 L_x=2\sum_{y\sim x}\chi_{xy}U_{xy}^*y_y.         \tag{40.7}
\]
The orientation of an edge may be reversed; (40.6) always uses the
positive exponential multiplying the \(x\)-endpoint norm.  Terms
depending only on the exterior variables cancel from the conditional
normalization.

\begin{lemma}[Local source bound]
\label{lem:v17v40-source}
If \(\deg x\leq D\) and \(\chi_{xy}\leq\chi_*\), then
\[
 \|L_x\|
 \leq2\chi_*\sum_{y\sim x}\|y_y\|,               \tag{40.8}
\]
and
\[
 \|L_x\|^{4/3}
 \leq(2\chi_*)^{4/3}D^{1/3}
 \sum_{y\sim x}\|y_y\|^{4/3}.                   \tag{40.9}
\]
\end{lemma}

\begin{proof}
Unitarity of every \(U_{xy}\) and the triangle inequality give
(40.8).  The convex power inequality
\((\sum_{j=1}^D a_j)^{4/3}\leq D^{1/3}\sum_j a_j^{4/3}\)
gives (40.9).
\end{proof}

\subsection{The positive conformal quadratic can only help}

For \(a\in\mathbb R\), \(L\in\mathbb R^q\), and \(c>0\), define
\[
 Z_c(a,L)=\int_{\mathbb R^q}
 \exp\left(-c\|h\|^4-a\|h\|^2+\langle L,h\rangle\right)dh.    \tag{40.10}
\]
For \(0<\vartheta<\lambda_H\), the conditional quartic moment is
\[
 M_\vartheta(a,L)
 =\frac{Z_{\lambda_H-\vartheta}(a,L)}
 {Z_{\lambda_H}(a,L)}.                            \tag{40.11}
\]

\begin{theorem}[Monotonicity in the quadratic coefficient]
\label{thm:v17v40-monotone}
For every fixed \(L\) and
\(0<\vartheta<\lambda_H\), the function
\[
 a\longmapsto M_\vartheta(a,L)                    \tag{40.12}
\]
is nonincreasing on \(\mathbb R\).  Consequently, the local conditional
moment in (40.5) satisfies
\[
 \boxed{
 \mathbb E_xe^{\vartheta\|y_x\|^4}
 \leq M_\vartheta(m_H^2,L_x),}                    \tag{40.13}
\]
uniformly in every conformal difference and compact gauge link.
\end{theorem}

\begin{proof}
Let \(\nu_{a,L}\) be the probability law with density proportional to
the denominator integrand in (40.11), and put \(R=\|h\|^2\).  Then
\[
 M_\vartheta(a,L)=\mathbb E_{\nu_{a,L}}e^{\vartheta R^2}.       \tag{40.14}
\]
Differentiation under the integral gives
\[
 \frac{d}{da}M_\vartheta(a,L)
 =-\operatorname{Cov}_{\nu_{a,L}}
 \left(R,e^{\vartheta R^2}\right).                \tag{40.15}
\]
For any real random variable \(R\) and nondecreasing functions
\(f,g\),
\[
 \operatorname{Cov}(f(R),g(R))
 =\frac12\mathbb E\left[
 (f(R)-f(R'))(g(R)-g(R'))
 \right]\geq0,                                    \tag{40.16}
\]
where \(R'\) is an independent copy.  Apply this with
\(f(R)=R\) and \(g(R)=e^{\vartheta R^2}\), noting that
\(R\geq0\).  Hence (40.15) is nonpositive.  Since
\(A_x\geq m_H^2\), (40.13) follows.
\end{proof}

This monotonicity removes the dangerous factor
\(\sum_{y\sim x}\chi_{xy}e^{s_y-s_x}\) from the normalization bound.
A small fixed anchor at the origin would have charged that factor to a
gravity row even though stronger positive quadratic confinement cannot
increase a radial quartic moment.

\subsection{The remaining source problem is subquartic}

After Theorem~\ref{thm:v17v40-monotone}, every exterior dependence is
carried by \(L_x\), which contains neighboring Higgs variables but no
conformal exponent.  Its natural growth is \(\|L_x\|^{4/3}\), dual to
the quartic power.  The elementary exact inequality
\[
 r^{4/3}\leq\zeta r^4+
 \frac{2}{3\sqrt{3\zeta}},
 \qquad r\geq0,quad\zeta>0,                       \tag{40.17}
\]
follows by maximizing \(u-\zeta u^3\) with
\(u=r^{4/3}\).  It permits an arbitrarily small neighbor-quartic
influence once a bound of the form
\[
 \log M_\vartheta(m_H^2,L)
 \leq C+C_L\|L\|^{4/3}                            \tag{40.18}
\]
has been proved.

\begin{assessment}[What V40 closes]
\label{ass:v17v40-status}
The mandatory companion audit is complete.  V40 identifies the exact
local Higgs conditional potential and proves that the full conformal
quadratic coefficient can only decrease its quartic exponential moment.
The remaining boundary source has subquartic order \(4/3\) and no
gravity variable.

The explicit bound (40.18), its Dobrushin row, and the resulting
volume-uniform local moment are completed in V41.
\end{assessment}

\begin{decision}[V41: close the local quartic moment]
\label{dec:v17v40-next}
Bound the numerator of (40.11) by quartic Young inequalities and the
denominator on a fixed unit ball.  Prove (40.18) with explicit constants,
insert (40.9) and (40.17), and solve the scalar neighbor-row contraction.
Continue past the audit before treating the checkpoint as complete.
\end{decision}

\begin{firewall}[Claim boundary after V40]
\label{fire:v17v40}
Monotonicity of the local moment does not itself prove a volume-uniform
DLR estimate or an infinite-volume joint measure.  The physical RPESM
gravitational action and continuum limit remain unresolved.
\end{firewall}

\section*{Version V40 append-only record}
\addcontentsline{toc}{section}{Version V40 append-only record}
The complete V39 source was retained byte for byte before this
continuation.  It had \(383558\) bytes, \(11785\) lines, and SHA--256
\[
 \texttt{46DF83B2015B3ECB923093CEE42FDAEC4362D347F99A553052AEBE2CE7FB4DF2}.
\]
V40 completed the audit and proved local Higgs moment monotonicity.


\section{V41: volume-uniform local Higgs exponential moments}
\label{sec:v17v41}

\subsection{An explicit quartic source integral}

Fix
\[
 0<\vartheta<\lambda_H,
 \qquad c=\lambda_H-\vartheta>0,                   \tag{41.1}
\]
and write
\[
 a_+=\max(m_H^2,0),
 \qquad a_-=\max(-m_H^2,0).                        \tag{41.2}
\]
Let
\[
 I_q(\alpha)=\int_{\mathbb R^q}e^{-\alpha\|h\|^4}\,dh
 =\frac{\pi^{q/2}\Gamma(q/4)}{2\Gamma(q/2)}
 \alpha^{-q/4},                                    \tag{41.3}
\]
and let
\[
 v_q=\frac{\pi^{q/2}}{\Gamma(q/2+1)}              \tag{41.4}
\]
be the volume of the unit ball.

\begin{theorem}[Explicit quartic Laplace-ratio bound]
\label{thm:v17v41-ratio}
For every \(L\in\mathbb R^q\),
\[
 \boxed{
 \log M_\vartheta(m_H^2,L)
 \leq C_0+C_L\|L\|^{4/3},}                        \tag{41.5}
\]
where one may take
\[
 C_L=1+\frac{3}{4c^{1/3}},                         \tag{41.6}
\]
\[
 C_0=\frac{a_-^2}{c}
 +\log\frac{I_q(c/2)}{v_q}
 +\lambda_H+a_++1.                                \tag{41.7}
\]
\end{theorem}

\begin{proof}
The numerator in (40.11), at \(a=m_H^2\), is at most
\[
 \int_{\mathbb R^q}
 \exp\left(-c\|h\|^4+a_-\|h\|^2+\|L\|\|h\|\right)dh.         \tag{41.8}
\]
For \(r\geq0\),
\[
 a_-r^2\leq\frac c4r^4+\frac{a_-^2}{c},           \tag{41.9}
\]
and direct maximization gives
\[
 \|L\|r
 \leq\frac c4r^4+
 \frac{3}{4c^{1/3}}\|L\|^{4/3}.                  \tag{41.10}
\]
Thus (41.8) is bounded above by
\[
 I_q(c/2)\exp\left(
 \frac{a_-^2}{c}
 +\frac{3}{4c^{1/3}}\|L\|^{4/3}
 \right).                                         \tag{41.11}
\]

For the denominator, integrate only over \(\|h\|\leq1\).  On that
ball,
\[
 -\lambda_H\|h\|^4-m_H^2\|h\|^2+\langle L,h\rangle
 \geq-\lambda_H-a_+-\|L\|.                       \tag{41.12}
\]
Hence the denominator is at least
\[
 v_qe^{-\lambda_H-a_+-\|L\|}.                     \tag{41.13}
\]
Divide (41.11) by (41.13), take logarithms, and use
\(u\leq1+u^{4/3}\) for \(u\geq0\).  Equations
(41.5)--(41.7) follow.
\end{proof}

\subsection{The local Higgs row}

Under the degree and hopping bounds of V40, define
\[
 K_H=C_L(2\chi_*)^{4/3}D^{1/3}.                   \tag{41.14}
\]
Theorems~\ref{thm:v17v40-monotone} and
\ref{thm:v17v41-ratio}, together with (40.9), give
\[
 \log\mathbb E_xe^{\vartheta V_H(y_x)}
 \leq C_0+K_H\sum_{y\sim x}\|y_y\|^{4/3}.       \tag{41.15}
\]

\begin{theorem}[Arbitrarily small neighbor-quartic influence]
\label{thm:v17v41-row}
For every \(\zeta>0\),
\[
 \log\mathbb E_xe^{\vartheta V_H(y_x)}
 \leq C_\zeta+K_H\zeta
 \sum_{y\sim x}V_H(y_y),                          \tag{41.16}
\]
where
\[
 C_\zeta=C_0+
 \frac{2K_HD}{3\sqrt{3\zeta}}.                   \tag{41.17}
\]
Thus the normalized Higgs row has entries
\[
 B_{xy}^{H}=\frac{K_H\zeta}{\vartheta}
 \mathbf1_{\{x\sim y\}},                         \tag{41.18}
\]
and row sum at most
\[
 \beta_H=\frac{DK_H\zeta}{\vartheta}.             \tag{41.19}
\]
\end{theorem}

\begin{proof}
Apply the exact inequality (40.17) to every neighbor norm in (41.15).
The constant is charged at most \(D\) times, giving
(41.16)--(41.17).  Division of the input quartic coefficient by the
exponential parameter \(\vartheta\) gives (41.18)--(41.19).
\end{proof}

Unlike the V21 normalization anchor, this row contains no gravity
entry.  The positive conformal quadratic was removed by monotonicity,
and the geometric cross term contributes only the subquartic source.

\subsection{Uniform finite-volume moment}

Let \(\mu_G\) be the full finite-volume joint law of the declared model
with free or periodic Higgs boundary convention, so every Higgs variable
appearing in an interaction is integrated in the finite-volume law.
Conformal and compact-gauge boundary values may be arbitrary.  V39
supplies finiteness of all moments used below.

\begin{theorem}[Volume-uniform one-site quartic moment]
\label{thm:v17v41-uniform}
Choose
\[
 0<\zeta<\frac{\vartheta}{DK_H},                   \tag{41.20}
\]
so that \(\beta_H<1\).  Then
\[
\boxed{
 \sup_G\sup_{x\in G}
 \log\mathbb E_{\mu_G}
 e^{\vartheta\|y_x\|^4}
 \leq\frac{C_\zeta}{1-\beta_H}.}                 \tag{41.21}
\]
The supremum is over all finite graphs with degree at most \(D\),
hopping weights at most \(\chi_*\), and the same local Higgs parameters.
The bound is independent of the conformal and compact-gauge boundary
values.
\end{theorem}

\begin{proof}
For a fixed finite graph, put
\[
 m_x=\mathbb E_{\mu_G}e^{\vartheta V_H(y_x)},
 \qquad M=\max_x\log m_x.                          \tag{41.22}
\]
Take the full expectation of (41.16).  Generalized Holder's inequality,
with an additional constant factor of exponent \(1-\beta_x\), where
\[
 \beta_x=\sum_yB_{xy}^{H}\leq\beta_H,             \tag{41.23}
\]
gives
\[
 m_x
 \leq e^{C_\zeta}
 \prod_{y\sim x}m_y^{B_{xy}^{H}}.                 \tag{41.24}
\]
Therefore
\[
 \log m_x\leq C_\zeta+\beta_HM.                  \tag{41.25}
\]
Maximize over \(x\) and solve
\(M\leq C_\zeta+\beta_HM\), proving (41.21).
All applications of Holder are legitimate because the finite-volume
quartic moments are finite by Corollary~\ref{cor:v17v39-normalization}.
\end{proof}

\begin{corollary}[Uniform local tightness of Higgs marginals]
\label{cor:v17v41-tight}
For every \(R>0\),
\[
 \sup_G\sup_{x\in G}
 \mu_G(\|y_x\|>R)
 \leq\exp\left(
 \frac{C_\zeta}{1-\beta_H}-\vartheta R^4
 \right).                                         \tag{41.26}
\]
Consequently, the one-site Higgs marginals are uniformly tight along
every bounded-degree graph exhaustion.  Every fixed finite collection
of Higgs sites is tight as well.
\end{corollary}

\begin{proof}
Apply Markov's inequality to (41.21).  A union bound gives tightness for
any fixed finite collection.
\end{proof}

\subsection{Claim boundary of the local estimate}

Theorem~\ref{thm:v17v41-uniform} concerns Higgs variables and is uniform
even before gravity is integrated.  It does not give tightness of the
mean-zero conformal field under Higgs tilting.  The conditional quotient
covariance in V32 is centered at its own mean, so a separate estimate
on that mean is required before Prokhorov compactness can be applied to
the joint local marginals.

\begin{assessment}[What V41 closes]
\label{ass:v17v41-status}
V41 continues substantively after the V40 audit.  It proves the explicit
quartic source bound, an arbitrarily small neighbor-Higgs row, a
volume-uniform one-site exponential quartic moment, and uniform local
Higgs tightness on bounded-degree graph exhaustions.  No gravitational
boundary variable enters the final estimate.

The next unresolved local variable is the mean of each conformal edge
difference under the Higgs-tilted quotient law.
\end{assessment}

\begin{decision}[V42: control the tilted conformal mean]
\label{dec:v17v41-next}
Use the stationarity equation
\(\mathbb E_y\nabla\Phi_y=0\), the strong monotonicity of
\(\nabla\Phi_y\), and the Higgs quartic moment from V41 to bound the
quotient mean or its edge differences uniformly.  Combine that mean
bound with (32.15) to obtain local conformal tightness.  If a vertex
mean is graph-size dependent, work directly with edge differences,
which are the physical variables in the model.
\end{decision}

\begin{firewall}[Claim boundary after V41]
\label{fire:v17v41}
Uniform Higgs tightness does not yet construct a joint DLR limit,
establish reflection positivity or continuum scaling, identify the
physical RPESM gravitational action, or prove the Einstein--SM
continuum limit.
\end{firewall}

\section*{Version V41 append-only record}
\addcontentsline{toc}{section}{Version V41 append-only record}
The complete V40 source was retained byte for byte before this
continuation.  It had \(390881\) bytes, \(11994\) lines, and SHA--256
\[
 \texttt{4E69B64DEAEF40A3BE6A305310D458819D7147D0E5B6CA382E5B20241C951DD5}.
\]
V41 proved volume-uniform local Higgs exponential moments.


\section{V42: transport control of the tilted conformal mean and local tightness}
\label{sec:v17v42}

\subsection{Relative entropy of the Higgs tilt}

For fixed Higgs variables, recall
\[
 d\mu_y=Z_g(y)^{-1}e^{-\Phi_0-T_y}\,d\mathcal H_0,
 \qquad
 d\mu_0=Z_g(0)^{-1}e^{-\Phi_0}\,d\mathcal H_0.     \tag{42.1}
\]
The normalized free energy is
\(\mathcal F_g(y)=-\log(Z_g(y)/Z_g(0))\).

\begin{lemma}[Exact entropy identity]
\label{lem:v17v42-entropy}
The relative entropy is
\[
 \operatorname{Ent}(\mu_y\mid\mu_0)
 =\mathcal F_g(y)-\mathbb E_yT_y
 \leq\mathcal F_g(y).                             \tag{42.2}
\]
\end{lemma}

\begin{proof}
The logarithmic density ratio is
\[
 \log\frac{d\mu_y}{d\mu_0}
 =-T_y+\mathcal F_g(y).                            \tag{42.3}
\]
Integrate it under \(\mu_y\).  Since \(T_y\geq0\), the upper bound
follows.
\end{proof}

\subsection{Talagrand displacement bound}

The unperturbed quotient law has Hessian at least
\(2\kappa_gI\) and satisfies the logarithmic Sobolev inequality
(37.15).

\begin{theorem}[Transport cost of the Higgs tilt]
\label{thm:v17v42-t2}
For every fixed Higgs field,
\[
\boxed{
 W_2^2(\mu_y,\mu_0)
 \leq\frac{1}{\kappa_g}\mathcal F_g(y).}          \tag{42.4}
\]
If
\[
 m(y)=\mathbb E_y r,                               \tag{42.5}
\]
then
\[
 \boxed{
 \|m(y)\|_2^2
 \leq\frac{1}{\kappa_g}\mathcal F_g(y).}          \tag{42.6}
\]
\end{theorem}

\begin{proof}
The logarithmic Sobolev inequality with curvature
\(2\kappa_g\) implies the Talagrand inequality
\[
 W_2^2(\nu,\mu_0)
 \leq\frac1{\kappa_g}
 \operatorname{Ent}(\nu\mid\mu_0).                \tag{42.7}
\]
This is the standard finite-dimensional LSI-to-transport implication,
obtained by applying the entropy inequality to the Hamilton--Jacobi
semigroup.  Apply (42.7) with \(\nu=\mu_y\) and use (42.2), proving
(42.4).

Symmetry gives \(\mathbb E_0r=0\).  For every coupling
\((R_y,R_0)\) of \(\mu_y,\mu_0\), Jensen gives
\[
 \|m(y)\|_2^2
 =\|\mathbb E(R_y-R_0)\|_2^2
 \leq\mathbb E\|R_y-R_0\|_2^2.                   \tag{42.8}
\]
Take the infimum over couplings and use (42.4).
\end{proof}

\subsection{Conditional second moments}

Let \(n=|\mathcal V|\), \(m_E=|\mathcal E|\), and suppose the maximum
degree is at most \(D\).

\begin{theorem}[Mean-zero field and edge averages]
\label{thm:v17v42-conditional}
For every fixed Higgs field,
\[
 \mathbb E_y\|r\|_2^2
 \leq\frac{n-1}{2\kappa_g}
 +\frac{\mathcal F_g(y)}{\kappa_g},                \tag{42.9}
\]
and
\[
 \sum_{e\in\mathcal E}\mathbb E_y d_e^2
 \leq\frac{m_E}{\kappa_g}
 +\frac{2D}{\kappa_g}\mathcal F_g(y).             \tag{42.10}
\]
\end{theorem}

\begin{proof}
The covariance matrix form of (32.14) is
\[
 \operatorname{Cov}_y(r)
 \preceq\frac1{2\kappa_g}I_{\mathcal H_0}.         \tag{42.11}
\]
Taking its trace and adding \(\|m(y)\|^2\), then using (42.6), proves
(42.9).

For an edge, (32.15) gives
\(\operatorname{Var}_y(d_e)\leq1/\kappa_g\).  The vector of edge means
is \(Bm(y)\), and
\[
 \|Bm(y)\|_2^2
 \leq\lambda_{\max}(B^{\mathsf T}B)\|m(y)\|_2^2
 \leq2D\|m(y)\|_2^2.                              \tag{42.12}
\]
Sum variance plus squared mean and apply (42.6).
\end{proof}

\subsection{Averaging over the Higgs field}

Fix the V41 parameters and write
\[
 L_H=\frac{C_\zeta}{1-\beta_H}.                    \tag{42.13}
\]
Theorem~\ref{thm:v17v41-uniform} and Jensen imply
\[
 \sup_{G,x}\mathbb E\|y_x\|^4
 \leq\frac{L_H}{\vartheta},
 \qquad
 \sup_{G,x}\mathbb E\|y_x\|^2
 \leq M_{H,2}:=\sqrt{\frac{L_H}{\vartheta}}.      \tag{42.14}
\]
Let \(\nu_G\) be the joint law of all non-gravity variables after
quotient gravity has been integrated.

\begin{theorem}[Uniform averaged quotient size]
\label{thm:v17v42-average}
On every finite graph of maximum degree \(D\) and hopping bound
\(\chi_*\),
\[
 \frac1n\mathbb E_{\nu_G}\mathbb E_y\|r\|_2^2
 \leq\frac1{2\kappa_g}
 +\frac{D\chi_*e^{1/(2\kappa_g)}M_{H,2}}{\kappa_g}.              \tag{42.15}
\]
Also,
\[
\begin{aligned}
 \frac1{m_E}\sum_e
 \mathbb E_{\nu_G}\mathbb E_y d_e^2
 \leq{}&\frac1{\kappa_g}\\
 &+\frac{2D^2\chi_*e^{1/(2\kappa_g)}nM_{H,2}}
 {\kappa_gm_E}.                                   \tag{42.16}
\end{aligned}
\]
\end{theorem}

\begin{proof}
The free-energy estimate (39.13) and (42.14) give
\[
 \mathbb E_{\nu_G}\mathcal F_g(y)
 \leq e^{1/(2\kappa_g)}
 \sum_vh_v^\chi\mathbb E\|y_v\|^2
 \leq D\chi_*e^{1/(2\kappa_g)}nM_{H,2}.           \tag{42.17}
\]
Average (42.9)--(42.10), use (42.17), and divide by \(n\) or
\(m_E\), respectively.
\end{proof}

\subsection{Local tightness on homogeneous periodic exhaustions}

Consider a sequence of periodic, vertex-transitive graphs with
homogeneous local couplings and translation-invariant gauge action and
boundary convention.  Their joint finite-volume laws are
vertex-transitive.  If the graphs are also edge-transitive, the edge
marginals are identical.

\begin{corollary}[Uniform local conformal second moments]
\label{cor:v17v42-local}
On the vertex-transitive sequence,
\[
 \sup_{G,x}\mathbb E r_x^2
 \leq\frac1{2\kappa_g}
 +\frac{D\chi_*e^{1/(2\kappa_g)}M_{H,2}}{\kappa_g}.              \tag{42.18}
\]
If in addition \(m_E/n\geq c_E>0\) and the sequence is edge-transitive,
then
\[
 \sup_{G,e}\mathbb E d_e^2
 \leq\frac1{\kappa_g}
 +\frac{2D^2\chi_*e^{1/(2\kappa_g)}M_{H,2}}
 {\kappa_gc_E}.                                   \tag{42.19}
\]
\end{corollary}

\begin{proof}
Vertex transitivity makes every term in the left side of (42.15)
equal, giving (42.18).  Edge transitivity does the same in (42.16), and
\(n/m_E\leq c_E^{-1}\) gives (42.19).
\end{proof}

Restore the exactly independent common mode \(t\) from V32.  It has
mean \(-b_g/(2\kappa_g)\) and variance
\((2n\kappa_g)^{-1}\), so
\[
 \sup_G\mathbb E t^2
 \leq\frac{b_g^2}{4\kappa_g^2}+\frac1{2\kappa_g}. \tag{42.20}
\]
Since \(s_x=t+r_x\), (42.18)--(42.20) give a uniform second moment for
every absolute conformal site variable on the homogeneous periodic
exhaustion.

\begin{corollary}[Joint one-site tightness]
\label{cor:v17v42-tight}
On the homogeneous periodic exhaustions above, the one-site marginals
of \((s_x,y_x,U_{x,\cdot})\) are uniformly tight.  The same holds for
every fixed finite set of sites and links.
\end{corollary}

\begin{proof}
The conformal variables have the uniform second moments just proved;
the Higgs variables have the exponential quartic bound (41.21); and the
gauge links range over a compact group.  Markov's inequality and finite
unions give the assertion.
\end{proof}

For nontransitive exhaustions, (42.15)--(42.16) prove averaged tightness
only.  They do not justify a supremum over exceptional vertices or
edges.

\begin{assessment}[What V42 closes]
\label{ass:v17v42-status}
V42 controls the Higgs-induced quotient mean by relative entropy and
Talagrand transport, combines it with the uniform conditional covariance,
and proves averaged conformal second-moment bounds on arbitrary
bounded-degree graphs.  On homogeneous periodic transitive exhaustions,
these become uniform local conformal bounds and hence joint local
tightness with the V41 Higgs estimate and compact gauge sector.

This supplies tight finite-dimensional marginals but not yet their
projective consistency, a DLR equation, or a continuum scaling limit.
\end{assessment}

\begin{decision}[V43: diagonal subsequence and lattice DLR limit]
\label{dec:v17v42-next}
On a fixed locally finite transitive lattice, use the local tightness of
V41--V42 and a diagonal Prokhorov argument to extract a subsequential
infinite-volume measure.  Prove the DLR equations for bounded local
observables by controlling the finite-range conditional kernels.  Keep
this lattice thermodynamic limit distinct from the later continuum and
physical-Einstein identifications.
\end{decision}

\begin{firewall}[Claim boundary after V42]
\label{fire:v17v42}
The tightness result assumes the declared homogeneous periodic model.
It does not yet prove a DLR limit, uniqueness of a joint infinite-volume
phase, Osterwalder--Schrader properties, continuum scaling, the physical
RPESM gravitational action, or the Einstein--SM continuum limit.
\end{firewall}

\section*{Version V42 append-only record}
\addcontentsline{toc}{section}{Version V42 append-only record}
The complete V41 source was retained byte for byte before this
continuation.  It had \(398718\) bytes, \(12252\) lines, and SHA--256
\[
 \texttt{C260DE59885CD36FB2C847BB840C6FE1F97E2019B5DF7A99DDE4116D8746E51A}.
\]
V42 proved transport control of the conformal mean and local tightness
on homogeneous periodic exhaustions.


\section{V43: an infinite-volume lattice Gibbs state}
\label{sec:v17v43}

\subsection{Periodic quotient hypothesis}

The compactness statement of V42 concerns the full site field
\[
 X_x=(s_x,y_x)
\]
and the compact gauge variables, after the common conformal mode has
been restored.  The quotient coordinates used to prove the estimate
will not be used to define the infinite-volume specification.

Let \(\mathbb L=(\mathbb V,\mathbb E)\) be a countable, locally finite
transitive lattice of degree at most \(D\).  Choose one orientation of
each edge and write
\[
 \Omega=
 \prod_{x\in\mathbb V}(\mathbb R\times\mathbb R^{q_H})
 \times\prod_{e\in\mathbb E}K,                    \tag{43.1}
\]
where \(K\) is the compact gauge group.  This is a Polish space in its
product topology.

We make the periodic exhaustion used implicitly in V42 precise.  Let
\(p_j:\mathbb L\to G_j\) be finite periodic quotient maps with
injectivity radii \(R_j\to\infty\).  All local couplings on \(G_j\) are
the quotients of the fixed homogeneous couplings on \(\mathbb L\), and
lattice translations descend to the quotients.  Let \(\mu_j\) be the
full finite-volume periodic Gibbs law of the declared
Higgs--gauge--conformal model on \(G_j\).

The interaction has a fixed finite range \(R_0\).  We use only the
following properties, all satisfied by the declared finite-volume
action: the local potentials are continuous; \(\kappa_g>0\) gives the
on-site conformal quadratic; \(\lambda_H>0\) gives the on-site Higgs
quartic; the hopping terms are nonnegative squares before expansion;
and the pure gauge variables lie in \(K\).  These facts make every
finite-region conditional partition function finite and positive.

\subsection{The infinite-volume local kernels}

For finite \(\Lambda\Subset\mathbb V\), include in \(X_\Lambda\) the
site variables in \(\Lambda\) and a fixed ownership choice for the
gauge links whose variables are integrated in that region.  Let
\(\partial_{R_0}\Lambda\) be the finite set of exterior variables met
by interaction cells intersecting \(\Lambda\).  The local Gibbs kernel
is
\[
 \gamma_\Lambda(d\xi_\Lambda\mid\eta)
 =
 \frac{\exp\{-H_\Lambda(\xi_\Lambda\mid\eta)\}}
 {Z_\Lambda(\eta)}
 \,d\xi_\Lambda,                                  \tag{43.2}
\]
with Lebesgue measure on the noncompact fields and Haar measure on the
owned gauge links.  Only
\(\eta|_{\partial_{R_0}\Lambda}\) enters (43.2).

\begin{lemma}[Proper Feller specification]
\label{lem:v17v43-feller}
For every finite \(\Lambda\), (43.2) is a probability kernel.  If \(f\)
is bounded and continuous in the variables of \(\Lambda\), then
\[
 \eta\longmapsto\gamma_\Lambda f(\eta)
 :=\int f(\xi_\Lambda)
 \gamma_\Lambda(d\xi_\Lambda\mid\eta)             \tag{43.3}
\]
is a bounded continuous function of the finite boundary layer
\(\partial_{R_0}\Lambda\).
\end{lemma}

\begin{proof}
Fix a compact set \(C\) of boundary configurations.  Continuity of the
finite collection of local potentials gives pointwise continuity of
the integrand in (43.2).  The nonnegative hopping terms may be dropped
when seeking an upper envelope.  The compact-group gauge action is
bounded below, while the quadratic and quartic on-site terms give,
after absorbing lower powers by Young's inequality,
\[
 e^{-H_\Lambda(\xi_\Lambda\mid\eta)}
 \leq A_C
 \exp\left\{
 -\frac{\kappa_g}{2}\sum_{x\in\Lambda}s_x^2
 -\frac{\lambda_H}{2}\sum_{x\in\Lambda}\|y_x\|^4
 \right\}                                        \tag{43.4}
\]
for \(\eta\in C\).  Haar volume is finite, so the right side is
integrable.  Dominated convergence proves continuity of
\(Z_\Lambda\) and of the numerator in (43.3).  The partition function
is strictly positive.  Division proves (43.3); boundedness follows
from \(|\gamma_\Lambda f|\leq\|f\|_\infty\).
\end{proof}

\subsection{Diagonal construction}

Choose increasing finite coordinate sets
\[
 \mathcal C_1\subset\mathcal C_2\subset\cdots,
 \qquad
 \bigcup_m\mathcal C_m
 =
 \{\hbox{all site and owned-link coordinates of }\mathbb L\}.   \tag{43.5}
\]
For fixed \(m\) and all large \(j\), the quotient map identifies
\(\mathcal C_m\) with coordinates of \(G_j\).

\begin{theorem}[Subsequential thermodynamic limit]
\label{thm:v17v43-limit}
Under the V41 parameter condition (41.20) and the homogeneous periodic
hypotheses above, there are a subsequence \(j_k\) and a
translation-invariant probability measure \(\mu_\infty\) on \(\Omega\)
such that, for every \(m\),
\[
 (p_{j_k}^{-1})_*
 \bigl(\mu_{j_k}|_{p_{j_k}(\mathcal C_m)}\bigr)
 \Longrightarrow
 \mu_\infty|_{\mathcal C_m}.                       \tag{43.6}
\]
The limit inherits the local estimates
\[
 \mathbb E_{\mu_\infty}s_x^2\leq C_s,\qquad
 \mathbb E_{\mu_\infty}
 e^{\vartheta\|y_x\|^4}\leq e^{L_H},               \tag{43.7}
\]
where
\[
 C_s=
 2\left(
 \frac1{2\kappa_g}
 +\frac{D\chi_*e^{1/(2\kappa_g)}M_{H,2}}{\kappa_g}
 \right)
 +2\left(
 \frac{b_g^2}{4\kappa_g^2}+\frac1{2\kappa_g}
 \right),                                         \tag{43.8}
\]
and \(L_H=C_\zeta/(1-\beta_H)\).
\end{theorem}

\begin{proof}
Corollary~\ref{cor:v17v42-tight} makes the laws of every fixed
\(\mathcal C_m\) uniformly tight.  Prokhorov's theorem gives a weakly
convergent subsequence for \(m=1\).  Extract successively for
\(m=2,3,\ldots\), and take the diagonal subsequence.  Denote the
resulting limit on \(\mathcal C_m\) by \(\nu_m\).  Marginal projection
is continuous, hence uniqueness of weak limits gives
\[
 (\pi_{\mathcal C_m})_*\nu_{m+1}=\nu_m.            \tag{43.9}
\]
The Kolmogorov extension theorem on the countable product of standard
Borel spaces therefore produces a unique measure \(\mu_\infty\) with
these marginals, proving (43.6).

Each lattice translation descends to every sufficiently large
periodic quotient and preserves \(\mu_j\).  Apply this identity to a
bounded continuous cylinder function and pass to (43.6).  Such
functions determine probability measures on \(\Omega\), so
\(\mu_\infty\) is translation invariant.

Equation (42.18), (42.20), and
\((a+b)^2\leq2a^2+2b^2\) give the first uniform bound in
(43.7)--(43.8).  Equation (41.21) gives the second.  Both integrands
are nonnegative and lower semicontinuous, so the Portmanteau theorem
passes the bounds to the weak limit.
\end{proof}

\subsection{Passage of the DLR identities}

\begin{theorem}[DLR equation for the declared lattice model]
\label{thm:v17v43-dlr}
The measure \(\mu_\infty\) of
Theorem~\ref{thm:v17v43-limit} is a Gibbs measure for the specification
\(\{\gamma_\Lambda\}_{\Lambda\Subset\mathbb V}\).  Equivalently, for
every finite \(\Lambda\), every bounded continuous
\(f=f(X_\Lambda)\), and every bounded exterior-measurable \(g\),
\[
 \int g(\eta)f(\eta)\,d\mu_\infty(\eta)
 =
 \int g(\eta)\gamma_\Lambda f(\eta)
 \,d\mu_\infty(\eta).                              \tag{43.10}
\]
\end{theorem}

\begin{proof}
First take \(g\) bounded and continuous with finite support outside
\(\Lambda\).  For all sufficiently large \(j\), the injectivity radius
places the supports of \(f,g\), and
\(\partial_{R_0}\Lambda\) in a copy of the same lattice neighborhood
inside \(G_j\).  The finite-volume conditional-distribution identity
then gives
\[
 \int g\bigl(f-\gamma_\Lambda f\bigr)\,d\mu_j=0.   \tag{43.11}
\]
By Lemma~\ref{lem:v17v43-feller}, the integrand in (43.11) is a bounded
continuous cylinder function.  Its support lies in one
\(\mathcal C_m\), so (43.6) permits passage to the limit and proves
(43.10) for continuous cylinder \(g\).

The bounded continuous exterior cylinder functions form a unital
algebra that generates the exterior Borel sigma-field.  The class of
bounded exterior functions for which (43.10) holds is closed under
bounded monotone limits.  The functional monotone-class theorem
therefore extends (43.10) to every bounded exterior-measurable \(g\).
Finally, bounded continuous functions determine probability measures
on the finite-dimensional Polish state space of \(X_\Lambda\).
Consequently (43.10) identifies the regular conditional law in
\(\Lambda\) with \(\gamma_\Lambda\), which is the DLR property.
\end{proof}

\begin{corollary}[Nonempty Gibbs simplex]
\label{cor:v17v43-nonempty}
The set of translation-invariant DLR measures for the declared
infinite-volume lattice model is nonempty.  Every subsequential limit
constructed above obeys the moment bounds (43.7).
\end{corollary}

\begin{proof}
Theorems~\ref{thm:v17v43-limit} and
\ref{thm:v17v43-dlr} give one such measure.  The same proof applies to
any convergent diagonal subsequence.
\end{proof}

\begin{assessment}[What V43 closes]
\label{ass:v17v43-status}
V43 converts the V41--V42 local estimates into an actual
translation-invariant infinite-volume Gibbs state for the declared
finite-range lattice Higgs--gauge--conformal model.  The proof keeps
the full local conformal variable throughout the DLR passage, so no
global mean-zero constraint contaminates the limiting specification.

The argument proves existence, not uniqueness, clustering, a mass gap,
or independence of the periodic subsequence.
\end{assessment}

\begin{decision}[V44: stability under changes of lattice scale]
\label{dec:v17v43-next}
Introduce an explicit lattice spacing \(a\) and identify scale-dependent
couplings for which the local estimates (41.21) and (43.7) remain
uniform as \(a\downarrow0\).  Begin with the canonical engineering
rescaling of the scalar and conformal variables, and isolate the first
coefficient whose bound degenerates.  If the direct scaling is
circular, move upstream to a dimensionless block-field formulation.
\end{decision}

\begin{firewall}[Claim boundary after V43]
\label{fire:v17v43}
This is a thermodynamic limit for the declared lattice model at fixed
lattice spacing.  It does not prove uniqueness, Osterwalder--Schrader
reflection positivity, continuum tightness, a renormalization-group
trajectory, identification with the physical Standard Model or
Einstein action, or the Einstein--SM continuum limit.
\end{firewall}

\section*{Version V43 append-only record}
\addcontentsline{toc}{section}{Version V43 append-only record}
The complete V42 source was retained byte for byte before this
continuation.  It had \(407241\) bytes, \(12529\) lines, and SHA--256
\[
 \texttt{4945C285C01226CCBC2724B276901FE15DD5743400492C645DFA69F9EFA5085E}.
\]
V43 constructed a subsequential translation-invariant DLR state for
the declared infinite-volume lattice model.


\section{V44: cosh-edge moments and the first continuum-scale debit}
\label{sec:v17v44}

\subsection{A graph-uniform comparison with one isolated edge}

The factor \(e^{1/(2\kappa_g)}\) in V39 came from discarding the
nonlinear mismatch action.  That is harmless at fixed cutoff but
catastrophic under the four-dimensional scaling
\(\kappa_g(a)\asymp a^4\).  The mismatch term gives a sharper bound.

Let \(\mu_0\) be the unperturbed quotient law in (38.4).  Thus, up to an
irrelevant constant, its potential is
\[
 \Phi_0(r)=
 \kappa_g\|r\|_2^2+
 J\sum_{f\in\mathcal E}\cosh\bigl(2(Br)_f\bigr),
 \qquad r\in\mathcal H_0.                         \tag{44.1}
\]
For \(\nu\in\mathbb R\) and \(J>0\), write
\[
 K_\nu(J)=\int_0^\infty e^{-J\cosh u}\cosh(\nu u)\,du             \tag{44.2}
\]
for the modified Bessel function in its integral representation.

\begin{theorem}[Cosh-improved edge exponential moment]
\label{thm:v17v44-bessel}
On every finite connected simple graph, for every edge \(e\),
\[
\boxed{
 \mathbb E_0e^{(Br)_e}
 =\mathbb E_0\cosh((Br)_e)
 \leq M_g(J):=\frac{K_{1/2}(J)}{K_0(J)}.}          \tag{44.3}
\]
The bound is independent of the volume and of \(\kappa_g>0\).  Moreover,
for \(0<J\leq1/2\),
\[
 M_g(J)
 \leq
 \frac{e\sqrt{\pi/(2J)}}{\log(1/J)}.               \tag{44.4}
\]
\end{theorem}

\begin{proof}
Put \(d=(Br)_e\), and disintegrate Lebesgue measure on
\(\mathcal H_0\) along the affine fibres with fixed \(d\).  The
marginal density of \(d\) has the form
\[
 p_e(d)=Z_e^{-1}e^{-J\cosh(2d)}A_e(d),             \tag{44.5}
\]
where \(A_e(d)\) is the fibre integral of the exponential of the
negative potential with the distinguished edge term removed.  That
remaining potential is convex.  Pr\'ekopa's theorem therefore makes
\(A_e\) log-concave.  The transformation \(r\mapsto-r\) gives
\(A_e(-d)=A_e(d)\).  An even log-concave function is nonincreasing on
\([0,\infty)\).

Let \(q_J\) be the symmetric probability density proportional to
\(e^{-J\cosh(2d)}\).  If \(H(|d|)\) is nondecreasing, the elementary
oppositely-ordered covariance inequality gives
\[
 \frac{\mathbb E_{q_J}[H(|d|)A_e(d)]}
 {\mathbb E_{q_J}A_e(d)}
 \leq\mathbb E_{q_J}H(|d|).                       \tag{44.6}
\]
Indeed, for independent \(D,D'\) with law \(q_J\), the product
\((H(|D|)-H(|D'|))(A_e(D)-A_e(D'))\) is nonpositive; expanding its
expectation proves (44.6).

The law (44.5) is symmetric, so
\(\mathbb E_0e^d=\mathbb E_0\cosh d\).  Apply (44.6) with
\(H(|d|)=\cosh d\).  The substitution \(u=2d\) and (44.2) give
\[
 \mathbb E_{q_J}\cosh d
 =\frac{K_{1/2}(J)}{K_0(J)},                      \tag{44.7}
\]
which proves (44.3).

The exact half-order formula is
\[
 K_{1/2}(J)=\sqrt{\frac{\pi}{2J}}e^{-J}.           \tag{44.8}
\]
For \(0\leq u\leq\log(1/J)\) and \(J\leq1/2\),
\[
 J\cosh u\leq\frac{1+J^2}{2}<1,
\]
and hence
\[
 K_0(J)\geq e^{-1}\log(1/J).                      \tag{44.9}
\]
Combine (44.8)--(44.9) and discard \(e^{-J}\leq1\), obtaining (44.4).
\end{proof}

\begin{corollary}[Improved gravity free-energy coefficient]
\label{cor:v17v44-free}
Every occurrence of \(e^{1/(2\kappa_g)}\) in
(39.10)--(39.14), and consequently in
(42.15)--(42.19), may be replaced by \(M_g(J)\).
\end{corollary}

\begin{proof}
The only use of \(e^{1/(2\kappa_g)}\) in the proof of
Theorem~\ref{thm:v17v39-radial} was the unperturbed estimate
\(\mathbb E_0e^{\pm d_e}\leq e^{1/(2\kappa_g)}\).
Symmetry and (44.3) supply the stated replacement.  All later bounds
were algebraic consequences of that estimate.
\end{proof}

\subsection{Canonical four-dimensional coefficient ledger}

Let \(G_a=(a\mathbb Z/L\mathbb Z)^4\), where \(a=L/N\), and set
\[
\begin{gathered}
 \kappa_g(a)=\bar\kappa_g a^4,\qquad
 b_g(a)=\bar b_g a^4,\qquad
 J(a)=\bar J a^2,                                  \tag{44.10}\\
 y_x=aH(ax),\qquad
 \chi_e(a)=\bar\chi_e,\qquad
 m_H^2(a)=\bar m_H^2a^2,\qquad
 \lambda_H(a)=\bar\lambda_H.                       \tag{44.11}
\end{gathered}
\]
For smooth \(s,H\), these powers are forced by Riemann-sum scaling:
\[
\begin{aligned}
 \kappa_g(a)\sum_xs_x^2
 &\longrightarrow\bar\kappa_g\int s^2,\\
 b_g(a)\sum_xs_x
 &\longrightarrow\bar b_g\int s,\\
 J(a)\sum_e\{\cosh(2d_e)-1\}
 &\longrightarrow2\bar J\int|\nabla s|^2,          \tag{44.12}\\
 \sum_e\bar\chi_e\|y_{e_+}-U_ey_{e_-}\|^2
 &\longrightarrow\int|D H|^2,\\
 \sum_x\{\bar m_H^2a^2\|y_x\|^2+
 \bar\lambda_H\|y_x\|^4\}
 &\longrightarrow\int\{\bar m_H^2\|H\|^2+
 \bar\lambda_H\|H\|^4\}.
\end{aligned}
\]
The displayed limits are normalization identities for smooth test
fields.  They do not assert that random fields sampled from the lattice
laws are smooth or converge.

\subsection{What is tight and what still escapes}

For a smooth scalar test function \(f\), put
\[
 \ell_x^{(a)}=a^4f(ax),\qquad
 \|f\|_{2,a}^2=a^4\sum_x|f(ax)|^2.                \tag{44.13}
\]
Conditional on the Higgs field, let \(m_a(y)=\mathbb E_y r\).
Assume here that the V41 parameters
\(\vartheta,C_\zeta,\beta_H\) are chosen uniformly along the
scale-dependent family.

\begin{theorem}[Uniform conditional fluctuation smearing]
\label{thm:v17v44-smear}
Under (44.10),
\[
\boxed{
 \operatorname{Var}_y\left(
 a^4\sum_x f(ax)\{r_x-m_{a,x}(y)\}
 \right)
 \leq\frac{\|f\|_{2,a}^2}{2\bar\kappa_g}.}         \tag{44.14}
\]
The independent common mode also obeys
\[
 \operatorname{Var}\left(t\,a^4\sum_xf(ax)\right)
 =
 \frac{(a^4\sum_xf(ax))^2}
 {2\bar\kappa_gL^4}.                              \tag{44.15}
\]
Thus the conditionally centered conformal distribution is tight
against each fixed test function.
\end{theorem}

\begin{proof}
Apply the covariance bound (42.11) to the vector
\(\ell^{(a)}\):
\[
 \operatorname{Var}_y\langle\ell^{(a)},r\rangle
 \leq\frac{\|\ell^{(a)}\|_2^2}{2\kappa_g(a)}
 =\frac{a^8\sum_x|f(ax)|^2}{2\bar\kappa_ga^4},
\]
which is (44.14).  Since \(n=(L/a)^4\), insertion of
\(\kappa_g(a)=\bar\kappa_ga^4\) into (32.4) gives (44.15).
Chebyshev's inequality proves the final statement.
\end{proof}

The random conditional mean is the first uncontrolled term.  The
present transport estimate gives only
\[
\begin{aligned}
 \mathbb E\left|
 a^4\sum_xf(ax)m_{a,x}(y)\right|^2
 &\leq
 \frac{a^8\sum_x|f(ax)|^2}{\kappa_g(a)}
 \,\mathbb E\mathcal F_g(y)\\
 &\leq
 \frac{D\chi_*M_{H,2}L^4}{\bar\kappa_g}\,
 \|f\|_{2,a}^2\,a^{-4}M_g(\bar Ja^2).             \tag{44.16}
\end{aligned}
\]
By (44.4), the coefficient on the right side is at most
\[
 O\left(\frac{a^{-5}}{\log(1/a)}\right).           \tag{44.17}
\]
This is far better than the artificial \(e^{c/a^4}\) loss from V39.
However \(M_g(J)\geq1\), so the displayed coefficient in (44.16) still
grows at least as \(a^{-4}\); the estimate does not prove continuum
tightness.  Likewise, (41.21) controls the dimensionless variable
\(y_x\); it does not by itself make the point variable
\(H_a(x)=a^{-1}y_x\) tight.

\begin{countertheorem}[The current pointwise ledger does not close the
continuum step]
\label{ctr:v17v44-ledger}
The estimates V41--V44 prove tightness of the conditionally centered
conformal field after testing as in (44.13), but their stated upper
bounds do not give a scale-uniform second moment for the random
conditional-mean term in (44.16), nor tightness of the canonically
rescaled point Higgs field.  Therefore those estimates alone cannot
justify a continuum subsequence.
\end{countertheorem}

\begin{proof}
The first assertion is Theorem~\ref{thm:v17v44-smear}.  The available
bound for the missing mean is exactly (44.16), whose explicit
coefficient diverges because \(M_g(J)\geq1\).  The V41 bound permits
\(\mathbb E\|a^{-1}y_x\|^2\leq a^{-2}M_{H,2}\), which also diverges.
An upper bound that diverges is not a divergence theorem for the
fields; it is a proof that the present estimates do not imply the
required tightness.
\end{proof}

\begin{assessment}[What V44 closes]
\label{ass:v17v44-status}
V44 replaces the super-exponential unperturbed edge estimate by a
graph-uniform Bessel ratio, gives the canonical four-dimensional
coefficient ledger, and proves scale-uniform smearing of the
conditionally centered conformal fluctuation.  It isolates the random
Higgs-induced conditional mean as the first precise continuum
tightness bottleneck.
\end{assessment}

\begin{decision}[V45: companion audit and localization of the tilt mean]
\label{dec:v17v44-next}
Perform the mandatory YM/NS companion audit.  Then seek a localized
representation of
\(\langle\ell,m_a(y)\rangle\), using the divergence form of the hopping
force and the inverse Hessian of
\(2\kappa_gI+4JL\).  The target is a negative-Sobolev estimate that
charges local Higgs moments rather than the extensive free energy in
(44.16).
\end{decision}

\begin{firewall}[Claim boundary after V44]
\label{fire:v17v44}
The Bessel comparison and conditional smearing theorem apply to the
declared positive conformal model.  The scaling ledger is a necessary
engineering comparison, not a constructed renormalization-group
trajectory.  No full joint continuum tightness, gauge scaling,
reflection positivity, physical Einstein action, or Einstein--SM
continuum limit has been proved.
\end{firewall}

\section*{Version V44 append-only record}
\addcontentsline{toc}{section}{Version V44 append-only record}
The complete V43 source was retained byte for byte before this
continuation.  It had \(417637\) bytes, \(12793\) lines, and SHA--256
\[
 \texttt{64344FA0D861F4AC5E45E3C4CFF32FD83E3E603554638B51399D305F0008F249}.
\]
V44 improved the edge moment, fixed the four-dimensional scaling
ledger, and isolated the continuum conditional-mean debit.


\section{V45: mandatory YM/NS audit and the continuum-mean route}
\label{sec:v17v45}

\subsection{Read-only companion checkpoints}

The active companion notes in the shared
\texttt{latex\_notes} directory were inspected before any new SM
direction was chosen.  Their exact checkpoints are
\[
\begin{array}{lll}
\text{programme}&\text{active file}&\text{SHA--256}\\
\text{YM}&
\texttt{Renewal\_Yang\_Mills\_Mass\_Gap\_Research\_Notes\_V84.tex}&
\texttt{7F63FC40435030261151BCB1268EF9AD}\\
&&\texttt{BD1AA24BC2598A59DA1029F7FEA77BA2}\\
\text{NS}&
\texttt{Renewal\_NS\_Stability\_Research\_Notes\_V26.tex}&
\texttt{82C887F8C2C69E4F28FAD7C3DC8DE5B}\\
&&\texttt{9099CB5A4BF431EBA1FFF3E91935978AD}.
\end{array}                                        \tag{45.1}
\]
The YM file has \(673625\) bytes and \(18396\) lines.  The NS file has
\(278283\) bytes and \(5319\) lines.

\subsection{Mathematical findings of the audit}

YM V83 repairs four obstructed terminal columns by rational bridge and
small terminal charts.  It thereby certifies both direct response
signs on six of eight transverse sign cells.  YM V84 begins the
seventh orientation: seven Boolean centres are inherited, one new
radius-\(1/8\) chart is certified at the \(111\) centre, and the first
Boolean cube is covered.  The remaining seventh cell, its inversion
partner, full-torus positivity, cofinal control, reflection positivity,
and a mass gap remain open.

NS V26 computes pairing-specific rank-one norms for the first two
Oseen-kernel derivatives.  The first-derivative norm equals the
all-symmetric-traceless norm, while the second-derivative rank-one norm
is strictly smaller on an intermediate radial window.  Its programme
therefore sharpens the particular contraction actually present in the
nonlinearity instead of charging an ambient operator norm.  No
Navier--Stokes regularity conclusion is claimed there.

\begin{theorem}[Companion-audit decision]
\label{thm:v17v45-decision}
Neither companion theorem can be imported as a missing SM continuum
estimate.  The audit nevertheless fixes two restrictions on the next
step:
\begin{enumerate}
\item the gauge sector remains a compact finite-cutoff spectator in
      the present SM probability argument until the YM atlas reaches
      its full-torus and cofinal stages;
\item the conformal conditional mean must be paired with the actual
      continuum test vector and the inverse conformal Hessian, rather
      than bounded by the ambient Euclidean norm used in (44.16).
\end{enumerate}
\end{theorem}

\begin{proof}
The YM results are positivity certificates for specified finite
Wilson-symbol charts; they contain no probability-space estimate for
the Higgs-induced conformal tilt.  The NS result concerns Oseen-kernel
tensor contractions and is not an SM field estimate.  Thus neither
supplies a premise of V44.  On the other hand, YM's incomplete final
orientation forbids promotion of compactness to gauge continuum
positivity, while the strict NS rank-one improvement demonstrates that
an operator norm taken only in the realized pairing can be smaller
than its ambient counterpart.  These are exactly the two restrictions
listed above.
\end{proof}

\subsection{Adjusted SM acquisition}

For the interpolation
\[
 d\mu_\tau^y(r)
 =
 Z_\tau(y)^{-1}
 e^{-\Phi_0(r)-\tau T_y(r)}\,d\mathcal H_0(r),
 \qquad0\leq\tau\leq1,                             \tag{45.2}
\]
write \(m_\tau(y)=\mathbb E_\tau^y r\).  Direct differentiation already
gives the exact directional identity
\[
 \frac{d}{d\tau}\langle\ell,m_\tau(y)\rangle
 =
 -\operatorname{Cov}_\tau^y(\langle\ell,r\rangle,T_y).           \tag{45.3}
\]
The global estimate in V44 applies Cauchy--Schwarz before using the
finite-range structure of \(T_y\).  V46 will retain the pairing in
(45.3), differentiate \(T_y\) edge by edge, and use an
\(\ell\)-dependent inverse-Hessian norm.  The required target is a
bound of the form
\[
 |\langle\ell,m_1(y)\rangle|
 \leq
 \sum_e w_e(\ell;\kappa_g,J)
 \,\mathcal Q_e(y),                                \tag{45.4}
\]
where \(w_e\) decays away from the support of \(\ell\) and
\(\mathcal Q_e\) is controlled by the V41 local Higgs moment.  Such a
bound would replace the extensive free energy in (44.16).

\begin{assessment}[What V45 changes]
\label{ass:v17v45-status}
The scheduled cross-program review is complete.  It supplies no
continuum theorem to import, but it confirms that the next useful SM
estimate is a pairing-specific, localized response bound for the
conditional conformal mean.  The gauge continuum claims remain behind
the unfinished YM atlas.
\end{assessment}

\begin{decision}[V46: inverse-Hessian response identity]
\label{dec:v17v45-next}
Derive a Helffer--Sj\"ostrand or Brascamp--Lieb representation for the
covariance in (45.3).  Preserve the divergence form
\(\nabla T_y=B^{\mathsf T}q_y\), obtain the exact deterministic
edge-response weights generated by
\((2\kappa_gI+4JL)^{-1}\ell\), and determine which moment of \(q_y\)
is still missing under the V41 law.
\end{decision}

\begin{firewall}[Claim boundary after V45]
\label{fire:v17v45}
V45 is an audit and route selection.  It does not transfer the YM or
NS conclusions into the SM model, prove (45.4), complete the gauge
atlas, establish continuum tightness, or construct the
Einstein--Standard-Model continuum.
\end{firewall}

\section*{Version V45 append-only record}
\addcontentsline{toc}{section}{Version V45 append-only record}
The complete V44 source was retained byte for byte before this
continuation.  It had \(427177\) bytes, \(13069\) lines, and SHA--256
\[
 \texttt{CE7AA5216F5B5248F4B0C971C13E5760564B233A87FBD6D409911984E5CE6FF4}.
\]
V45 performed the required YM/NS audit and selected the localized
conditional-mean route.


\section{V46: Witten-resolvent response of the conformal tilt}
\label{sec:v17v46}

\subsection{One-form covariance representation}

Fix a finite connected graph and a Higgs field \(y\).  For the
interpolation (45.2), set
\[
 \Phi_\tau^y=\Phi_0+\tau T_y,\qquad
 d_e=(Br)_e,\qquad a_x=\|y_x\|^2.                  \tag{46.1}
\]
The edge force of the Higgs tilt is
\[
 q_e(r,y)=
 \chi_e\{a_{e_-}e^{d_e}-a_{e_+}e^{-d_e}\},
 \qquad
 \nabla T_y=B^{\mathsf T}q.                       \tag{46.2}
\]
On vector fields over \(\mathcal H_0\), define the Witten one-form
operator
\[
 \mathscr A_\tau^y
 =
 \left(-\Delta+\nabla\Phi_\tau^y\mathbin{\cdot}\nabla\right)
 \otimes I+\nabla^2\Phi_\tau^y.                   \tag{46.3}
\]
Let \(L=B^{\mathsf T}B\) on \(\mathcal H_0\) and
\[
 A_0=2\kappa_gI+4JL.                              \tag{46.4}
\]
Since \(V_g(d)=2\cosh(2d)\), the mismatch Hessian is at least
\(4JL\), and the hopping Hessian is positive.  Hence
\[
 \nabla^2\Phi_\tau^y(r)\succeq A_0.               \tag{46.5}
\]

\begin{theorem}[Exact paired response and its energy budget]
\label{thm:v17v46-witten}
Assume \(J>0\).  For \(\ell\in\mathcal H_0\), put
\[
 v_\tau^{\ell,y}
 =(\mathscr A_\tau^y)^{-1}\ell,\qquad
 W_{\tau,e}^{\ell,y}=(Bv_\tau^{\ell,y})_e.         \tag{46.6}
\]
Then
\[
\boxed{
 \langle\ell,m_1(y)\rangle
 =
 -\int_0^1
 \mathbb E_\tau^y
 \sum_eW_{\tau,e}^{\ell,y}q_e\,d\tau.}            \tag{46.7}
\]
Moreover,
\[
\begin{aligned}
 &\mathbb E_\tau^y\left[
 \|\nabla v_\tau^{\ell,y}\|_{\rm HS}^2
 +2\kappa_g\|v_\tau^{\ell,y}\|_2^2
 +4J\|Bv_\tau^{\ell,y}\|_2^2
 \right]\\
 &\hspace{35mm}\leq
 \ell^{\mathsf T}A_0^{-1}\ell,                    \tag{46.8}
\end{aligned}
\]
and consequently
\[
\boxed{
 |\langle\ell,m_1(y)\rangle|^2
 \leq
 \frac{\ell^{\mathsf T}A_0^{-1}\ell}{4J}
 \int_0^1\mathbb E_\tau^y\|q(r,y)\|_2^2\,d\tau.}  \tag{46.9}
\]
\end{theorem}

\begin{proof}
The finite-dimensional Helffer--Sj\"ostrand identity gives
\[
 \operatorname{Cov}_\tau^y(\langle\ell,r\rangle,T_y)
 =
 \mathbb E_\tau^y
 \left\langle
 (\mathscr A_\tau^y)^{-1}\ell,\nabla T_y
 \right\rangle.                                  \tag{46.10}
\]
All derivatives have exponential growth controlled by the cosh and
hopping potentials.  Thus the identity follows first for smooth
truncations and then by dominated convergence.  Insert (46.2) into
(46.10), use (45.3), \(m_0=0\), and integrate in \(\tau\), proving
(46.7).

Taking the \(L^2(\mu_\tau^y)\) inner product of
\(\mathscr A_\tau^yv_\tau^{\ell,y}=\ell\) with
\(v_\tau^{\ell,y}\), integration by parts gives
\[
 \mathbb E_\tau^y\{
 \|\nabla v_\tau^{\ell,y}\|_{\rm HS}^2+
 (v_\tau^{\ell,y})^{\mathsf T}
 \nabla^2\Phi_\tau^y v_\tau^{\ell,y}\}
 =
 \mathbb E_\tau^y\langle\ell,v_\tau^{\ell,y}\rangle.             \tag{46.11}
\]
As quadratic forms on vector-valued \(L^2(\mu_\tau^y)\),
\(\mathscr A_\tau^y\succeq A_0\).  Inversion reverses this order, and
the constant vector field \(\ell\) therefore satisfies
\[
 \mathbb E_\tau^y\langle\ell,v_\tau^{\ell,y}\rangle
 \leq\ell^{\mathsf T}A_0^{-1}\ell.                \tag{46.12}
\]
Equations (46.5), (46.11), and (46.12) prove (46.8).

For each \(\tau\), Cauchy--Schwarz and (46.8) give
\[
\left|
 \mathbb E_\tau^y\sum_eW_{\tau,e}^{\ell,y}q_e
 \right|
 \leq
 \sqrt{\ell^{\mathsf T}A_0^{-1}\ell}\,
 \sqrt{\frac{\mathbb E_\tau^y\|q\|_2^2}{4J}}.     \tag{46.13}
\]
Apply Cauchy--Schwarz once more to the \(\tau\)-integral in (46.7).
\end{proof}

\subsection{All unperturbed edge exponential moments}

The Bessel comparison of V44 is not restricted to exponent one.

\begin{lemma}[Bessel hierarchy]
\label{lem:v17v46-bessel}
For \(p\geq0\) and every edge,
\[
 \mathbb E_0e^{p d_e}
 =\mathbb E_0\cosh(pd_e)
 \leq
 M_{g,p}(J):=\frac{K_{p/2}(J)}{K_0(J)}.            \tag{46.14}
\]
For \(0<J\leq1/2\),
\[
 M_{g,2}(J)\leq
 \frac{e}{J\log(1/J)}.                            \tag{46.15}
\]
\end{lemma}

\begin{proof}
The proof of Theorem~\ref{thm:v17v44-bessel} applies with the
nondecreasing function \(H(|d|)=\cosh(p|d|)\).  The substitution
\(u=2d\) gives the ratio in (46.14).  Also
\[
 K_1(J)
 =\int_0^\infty e^{-J\cosh u}\cosh u\,du
 \leq\int_0^\infty e^{-J\sinh u}\cosh u\,du
 =\frac1J,                                        \tag{46.16}
\]
because \(\cosh u\geq\sinh u\).  Combine this with the lower bound
(44.9) for \(K_0(J)\).
\end{proof}

\begin{corollary}[The force is controlled at the unperturbed endpoint]
\label{cor:v17v46-force-zero}
For fixed \(y\),
\[
 \mathbb E_0 q_e^2
 \leq
 2\chi_e^2M_{g,2}(J)
 \{\|y_{e_-}\|^4+\|y_{e_+}\|^4\}.                 \tag{46.17}
\]
If \(y\) is subsequently averaged under any law satisfying (42.14),
then
\[
 \mathbb E_y\mathbb E_0q_e^2
 \leq
 \frac{4\chi_*^2L_H}{\vartheta}M_{g,2}(J).         \tag{46.18}
\]
\end{corollary}

\begin{proof}
Square (46.2), use \((u-v)^2\leq2u^2+2v^2\), and apply (46.14) with
\(p=2\).  Equation (42.14) then gives (46.18).
\end{proof}

\subsection{The deterministic factor has the correct continuum scale}

On the four-torus of V44, let
\(\ell_x^{(a)}=a^4f(ax)\), with the constant Fourier mode removed if
necessary.  From (44.10),
\[
 A_{0,a}
 =
 a^4\left(
 2\bar\kappa_gI+
 4\bar J\,\frac{L_a}{a^2}
 \right).                                        \tag{46.19}
\]

\begin{proposition}[Discrete negative-Sobolev limit]
\label{prop:v17v46-hminus}
For smooth mean-zero \(f\) on the continuum torus,
\[
 (\ell^{(a)})^{\mathsf T}A_{0,a}^{-1}\ell^{(a)}
 \longrightarrow
 \left\langle
 f,(2\bar\kappa_g-4\bar J\Delta)^{-1}f
 \right\rangle_{L^2}.                             \tag{46.20}
\]
In particular,
\[
 (\ell^{(a)})^{\mathsf T}A_{0,a}^{-1}\ell^{(a)}
 \leq\frac{\|f\|_{2,a}^2}{2\bar\kappa_g}.         \tag{46.21}
\]
\end{proposition}

\begin{proof}
The inequality follows directly from
\(A_{0,a}\succeq2\bar\kappa_ga^4I\).  For convergence, diagonalize the
periodic graph Laplacian.  On a fixed Fourier mode \(k\),
\[
 a^{-2}\lambda_k(L_a)
 \longrightarrow |k|^2,                          \tag{46.22}
\]
with the torus normalization absorbed into \(k\).  Smoothness makes
the Fourier coefficients rapidly decreasing, while (46.21) supplies
a summable dominating tail after a standard smooth approximation.
Modewise convergence and Parseval's identity give (46.20).
\end{proof}

\subsection{The exact remaining force quantity}

Let \(\nu_a\) denote the final Higgs--gauge marginal at lattice spacing
\(a\), and define
\[
 \mathcal Q_a
 =
 \mathbb E_{\nu_a}
 \int_0^1\mathbb E_\tau^y
 \sum_e q_e(r,y)^2\,d\tau.                        \tag{46.23}
\]
Averaging (46.9) gives the rigorous reduction
\[
 \mathbb E_{\nu_a}
 |\langle\ell^{(a)},m_1(y)\rangle|^2
 \leq
 \frac{
 (\ell^{(a)})^{\mathsf T}A_{0,a}^{-1}\ell^{(a)}
 }{4J(a)}\,\mathcal Q_a.                          \tag{46.24}
\]
The deterministic numerator has a finite continuum limit by (46.20).
The factor \(J(a)^{-1}\asymp a^{-2}\) and the global edge sum in
\(\mathcal Q_a\) show that (46.24), by itself, is still too coarse.

Equations (46.17)--(46.18) control the endpoint \(\tau=0\), but V41
does not control (46.23) uniformly along the entire asymmetric tilted
path.  This is the precise missing moment.  It cannot be replaced by
the symmetric Bessel comparison without proof, because
\(\tau T_y\) destroys the \(d_e\mapsto-d_e\) symmetry used in
(44.5)--(44.7).

\begin{assessment}[What V46 closes]
\label{ass:v17v46-status}
V46 proves the exact Witten-resolvent representation of the
Higgs-induced conformal mean.  It separates a scale-safe
negative-Sobolev test factor from a path-integrated edge-force factor,
and controls that force at the unperturbed endpoint by a complete
Bessel hierarchy.  The unresolved object is now the localized,
tilted force response, rather than the extensive free energy of V44.
\end{assessment}

\begin{decision}[V47: killed-walk localization]
\label{dec:v17v46-next}
Use the massive graph operator
\(A_{0,a}=2\kappa_g(a)I+4J(a)L_a\) to obtain an explicit
killed-random-walk kernel and physical-distance decay for
\(B A_{0,a}^{-1}\ell^{(a)}\).  Then test whether the Witten response
in (46.7) admits the same domination.  If entrywise domination fails,
record the failure and seek a weighted energy inequality in annuli.
\end{decision}

\begin{firewall}[Claim boundary after V46]
\label{fire:v17v46}
The response identity is finite-volume and conditional on the declared
positive conformal action.  The endpoint force estimate does not
control the full tilted interpolation.  No joint continuum tightness,
renormalization trajectory, reflection positivity, physical Einstein
action, or Einstein--SM continuum limit follows yet.
\end{firewall}

\section*{Version V46 append-only record}
\addcontentsline{toc}{section}{Version V46 append-only record}
The complete V45 source was retained byte for byte before this
continuation.  It had \(432917\) bytes, \(13210\) lines, and SHA--256
\[
 \texttt{926CDE968BC25F201FAB15CEEB8E01C5B324F45A75E51078524176BEC5A38D62}.
\]
V46 derived the Witten response formula, its energy budget, the
Bessel moment hierarchy, and the exact remaining tilted-force term.


\section{V47: killed-walk localization and its conductance obstruction}
\label{sec:v17v47}

\subsection{The baseline massive resolvent}

Let \(G\) be a finite \(D\)-regular graph, let \(P\) be the transition
matrix of simple random walk, and write
\[
 L=D(I-P),\qquad A_0=2\kappa_gI+4JL.               \tag{47.1}
\]
The formula is valid on the full vertex space; it preserves the
mean-zero subspace used in V46.

\begin{theorem}[Exact killed-walk representation]
\label{thm:v17v47-walk}
Put
\[
 c_0=2\kappa_g+4JD,\qquad
 \alpha=\frac{4JD}{2\kappa_g+4JD}\in(0,1).        \tag{47.2}
\]
Then
\[
\boxed{
 A_0^{-1}
 =
 \frac1{c_0}\sum_{n=0}^\infty\alpha^nP^n.}        \tag{47.3}
\]
Thus \(A_0^{-1}(x,z)\geq0\), and
\[
 A_0^{-1}(x,z)
 \leq\frac1{2\kappa_g}
 \alpha^{\operatorname{dist}(x,z)}.               \tag{47.4}
\]
\end{theorem}

\begin{proof}
Equation (47.1) gives
\[
 A_0=c_0(I-\alpha P).
\]
Since \(\|P\|_{2\to2}=1\) and \(\alpha<1\), the Neumann series
converges in operator norm and proves (47.3).  Every entry of \(P^n\)
is nonnegative.  Also \(P^n(x,z)=0\) when
\(n<\operatorname{dist}(x,z)\), so
\[
 A_0^{-1}(x,z)
 \leq c_0^{-1}\sum_{n\geq\operatorname{dist}(x,z)}\alpha^n
 =(2\kappa_g)^{-1}\alpha^{\operatorname{dist}(x,z)}.
\]
\end{proof}

The last estimate uses only the minimum number of steps and has decay
rate \(1-\alpha\asymp\kappa_g/J\).  The correct massive rate is its
square root.  A weighted energy argument recovers it without a heat
kernel estimate.

\begin{theorem}[Agmon bound at the physical mass scale]
\label{thm:v17v47-agmon}
Let \(G\) have maximum degree at most \(D\), let
\(\ell\) be supported in \(S\), and put \(u=A_0^{-1}\ell\).  If
\(\rho(x)=\operatorname{dist}(x,S)\) and
\[
 4JD\{\cosh\theta-1\}\leq\kappa_g,                \tag{47.5}
\]
then
\[
\boxed{
 \left\|e^{\theta\rho}u\right\|_2
 \leq\frac{\|\ell\|_2}{\kappa_g}.}                \tag{47.6}
\]
Consequently,
\[
 |u_x|\leq
 \frac{\|\ell\|_2}{\kappa_g}
 e^{-\theta\operatorname{dist}(x,S)},             \tag{47.7}
\]
and for \(e=\{x,z\}\),
\[
 |(Bu)_e|
 \leq
 \frac{2\|\ell\|_2}{\kappa_g}
 e^{-\theta\operatorname{dist}(e,S)}.             \tag{47.8}
\]
\end{theorem}

\begin{proof}
Write \(w_x=e^{\theta\rho(x)}u_x\).  Multiply
\(A_0u=\ell\) by \(e^{2\theta\rho}u\) and sum over vertices.  For one
edge \(\{x,z\}\),
\[
\begin{aligned}
 &(u_x-u_z)
 (e^{2\theta\rho(x)}u_x-e^{2\theta\rho(z)}u_z)\\
 &\quad=
 (w_x-w_z)^2
 -2\{\cosh(\theta(\rho(x)-\rho(z)))-1\}w_xw_z\\
 &\quad\geq
 -\{\cosh\theta-1\}(w_x^2+w_z^2),                 \tag{47.9}
\end{aligned}
\]
because \(|\rho(x)-\rho(z)|\leq1\).
After summing edges and using the degree bound,
\[
 \{2\kappa_g-4JD(\cosh\theta-1)\}\|w\|_2^2
 \leq|\langle\ell,w\rangle|.                      \tag{47.10}
\]
The weight equals one on \(S\).  Condition (47.5) and
Cauchy--Schwarz give
\(\kappa_g\|w\|_2^2\leq\|\ell\|_2\|w\|_2\), proving
(47.6).  Equations (47.7)--(47.8) follow by taking components and
using the smaller endpoint distance for an edge.
\end{proof}

Under the four-dimensional ledger (44.10), the largest admissible
choice in (47.5) is
\[
 \theta_a=
 \operatorname{arcosh}
 \left(1+\frac{\bar\kappa_g}{4D\bar J}a^2\right),
 \qquad
 \frac{\theta_a}{a}
 \longrightarrow
 \sqrt{\frac{\bar\kappa_g}{2D\bar J}}.             \tag{47.11}
\]
Thus (47.7) decays exponentially in physical distance, with a
nonzero scale-independent rate.

\subsection{Why the baseline kernel cannot dominate the Witten kernel}

The Hessian in V46 is a massive graph operator with random edge
conductances bounded below by \(4J\), but not above.  Loewner
monotonicity controls quadratic forms of its inverse; it does not give
componentwise domination by the inverse at the conductance floor.

\begin{countertheorem}[Failure of entrywise lower-floor domination]
\label{ctr:v17v47-entrywise}
There is no general implication
\[
 C_e\geq c_e\ \hbox{for all edges}
 \quad\Longrightarrow\quad
 |(mI+B^{\mathsf T}CB)^{-1}\ell|_x
 \leq
 |(mI+B^{\mathsf T}cB)^{-1}\ell|_x               \tag{47.12}
\]
for mean-zero \(\ell\).
\end{countertheorem}

\begin{proof}
Take the path \(1\)--\(2\)--\(3\), mass \(m>0\), conductance \(c>0\)
on the first edge, conductance \(C>0\) on the second, and
\(\ell=(1,-1,0)^{\mathsf T}\).  If
\[
 u(C)=(mI+cL_{12}+CL_{23})^{-1}\ell,              \tag{47.13}
\]
then summing the three equations gives \(\sum_xu_x(C)=0\), and direct
elimination gives
\[
 u_3(C)=
 -\frac{C}
 {m^2+2mc+C(2m+3c)}.                              \tag{47.14}
\]
Its absolute value is strictly increasing in \(C\), since
\[
 \frac{d}{dC}|u_3(C)|
 =
 \frac{m^2+2mc}
 {\{m^2+2mc+C(2m+3c)\}^2}>0.                     \tag{47.15}
\]
Choosing \(C>c\) contradicts (47.12) at vertex \(3\).
\end{proof}

\begin{corollary}[No automatic killed-walk bound for the Witten
response]
\label{cor:v17v47-no-transfer}
The lower Hessian bound (46.5) alone cannot imply that the random
response \(v_\tau^{\ell,y}\) or its edge gradient is pointwise
dominated by \(A_0^{-1}\ell\) or \(BA_0^{-1}\ell\).
\end{corollary}

\begin{proof}
The edge conductance in the pointwise Hessian is
\[
 c_{\tau,e}(r,y)
 =
 4J\cosh(2d_e)
 +\tau\chi_e\{a_{e_-}e^{d_e}+a_{e_+}e^{-d_e}\}
 \geq4J.                                         \tag{47.16}
\]
It is unbounded above.  Any domination deduced solely from the lower
floor would apply in particular to constant quadratic potentials,
where the Witten solution for a constant forcing vector is the
ordinary matrix inverse.  Countertheorem~\ref{ctr:v17v47-entrywise}
rules that out.
\end{proof}

The counterexample does not destroy the energy estimate (46.8).
It says that localization must retain the actual conductance in the
energy, rather than replace it pointwise by its lower floor.  This
suggests pairing the force \(q_e\) with the same conductance
\(c_{\tau,e}\).

\begin{assessment}[What V47 closes]
\label{ass:v17v47-status}
V47 gives an exact killed-walk formula and a physical-rate Agmon bound
for the deterministic massive conformal resolvent.  It also proves
that the proposed entrywise transfer to the random Witten response is
false under a conductance lower bound alone.  The surviving route is
an actual-conductance weighted Cauchy--Schwarz estimate.
\end{assessment}

\begin{decision}[V48: force divided by curvature]
\label{dec:v17v47-next}
Use (47.16) directly in the Witten energy.  Bound
\(q_e^2/c_{\tau,e}\) pointwise by endpoint Higgs quartics, average with
V41, and determine the remaining volume factor.  Then formulate the
weighted annular estimate required to exploit physical-distance decay
without an invalid entrywise comparison.
\end{decision}

\begin{firewall}[Claim boundary after V47]
\label{fire:v17v47}
The killed-walk and Agmon bounds concern the deterministic comparison
operator.  The countertheorem explicitly prevents their unproved
pointwise transfer to the random tilted Witten operator.  Joint
continuum tightness, gauge continuum positivity, reflection
positivity, the physical Einstein action, and the Einstein--SM
continuum limit remain open.
\end{firewall}

\section*{Version V47 append-only record}
\addcontentsline{toc}{section}{Version V47 append-only record}
The complete V46 source was retained byte for byte before this
continuation.  It had \(441964\) bytes, \(13506\) lines, and SHA--256
\[
 \texttt{32008D448232149864E72315B8B52378ADCD75D25B6B754C25D3F76D47B33AEC}.
\]
V47 proved deterministic physical-scale localization and the exact
obstruction to transferring it entrywise to random conductances.


\section{V48: the exact edge-force/curvature quotient}
\label{sec:v17v48}

\subsection{Use the conductance that the tilt actually creates}

Retain the interpolation and notation of V46.  The spatial part of
the conditional Hessian is
\[
 \nabla^2\Phi_\tau^y
 =
 2\kappa_gI+
 B^{\mathsf T}\operatorname{diag}(c_{\tau,e})B,    \tag{48.1}
\]
where
\[
\begin{aligned}
 c_{\tau,e}(r,y)
 &=
 4J\cosh(2d_e)+\tau T_e(r,y),\\
 T_e(r,y)
 &=
 \chi_e\{a_{e_-}e^{d_e}+a_{e_+}e^{-d_e}\},\\
 q_e(r,y)
 &=
 \chi_e\{a_{e_-}e^{d_e}-a_{e_+}e^{-d_e}\}.        \tag{48.2}
\end{aligned}
\]
Thus the asymmetric tilt increases the same conductance against which
its force should be measured.

\begin{theorem}[Sharp pointwise force/curvature quotient]
\label{thm:v17v48-quotient}
For every \(J>0\), \(0\leq\tau\leq1\), edge field value \(d_e\), and
Higgs endpoints,
\[
\boxed{
 \frac{q_e(r,y)^2}{c_{\tau,e}(r,y)}
 \leq
 \frac{\chi_e^2}{2J}
 \max\{\|y_{e_-}\|^4,\|y_{e_+}\|^4\}.}            \tag{48.3}
\]
The coefficient \(1/(2J)\) is sharp if one endpoint Higgs norm is
allowed to vanish.
\end{theorem}

\begin{proof}
Since \(\tau T_e\geq0\), it suffices to divide by
\(4J\cosh(2d_e)\).  Put
\[
 A=a_{e_-},\qquad C=a_{e_+},\qquad z=e^{2d_e}>0.
\]
Then
\[
\begin{aligned}
 \frac{q_e^2}{4J\cosh(2d_e)}
 &=
 \frac{\chi_e^2}{2J}
 \frac{A^2z^2+C^2-2ACz}{z^2+1}\\
 &\leq
 \frac{\chi_e^2}{2J}
 \frac{A^2z^2+C^2}{z^2+1}\\
 &\leq
 \frac{\chi_e^2}{2J}\max\{A^2,C^2\}.              \tag{48.4}
\end{aligned}
\]
This is (48.3).  If \(C=0\) and \(z\to\infty\), the quotient in
(48.4) tends to \(\chi_e^2A^2/(2J)\), proving sharpness.
\end{proof}

\subsection{The tilted path moment is no longer missing}

Let \(v_\tau^{\ell,y}\) and
\(W_{\tau,e}^{\ell,y}\) be as in (46.6).  The exact energy identity
(46.11), now retaining all edge curvature, gives
\[
 \mathbb E_\tau^y
 \sum_ec_{\tau,e}|W_{\tau,e}^{\ell,y}|^2
 \leq
 \ell^{\mathsf T}A_0^{-1}\ell.                    \tag{48.5}
\]

\begin{theorem}[Uniform-in-tilt conditional-mean estimate]
\label{thm:v17v48-mean}
For every fixed Higgs field,
\[
\boxed{
 |\langle\ell,m_1(y)\rangle|^2
 \leq
 \ell^{\mathsf T}A_0^{-1}\ell
 \sum_e\frac{\chi_e^2}{2J}
 \max\{\|y_{e_-}\|^4,\|y_{e_+}\|^4\}.}            \tag{48.6}
\]
This estimate requires no moment along the asymmetric interpolation.
\end{theorem}

\begin{proof}
For fixed \(\tau\), use (46.7), weighted Cauchy--Schwarz, and (48.5):
\[
\begin{aligned}
 \left|
 \mathbb E_\tau^y\sum_eW_{\tau,e}^{\ell,y}q_e
 \right|^2
 &\leq
 \left(\mathbb E_\tau^y\sum_ec_{\tau,e}|W_{\tau,e}^{\ell,y}|^2\right)
 \left(\mathbb E_\tau^y\sum_e\frac{q_e^2}{c_{\tau,e}}\right)\\
 &\leq
 \ell^{\mathsf T}A_0^{-1}\ell
 \sum_e\frac{\chi_e^2}{2J}
 \max\{\|y_{e_-}\|^4,\|y_{e_+}\|^4\},             \tag{48.7}
\end{aligned}
\]
where Theorem~\ref{thm:v17v48-quotient} removes the \(r,\tau\)
dependence from the second factor.  Integrate (46.7) over an interval
of length one and use Cauchy--Schwarz in \(\tau\).
\end{proof}

\begin{corollary}[Annealed bound from the V41 moment]
\label{cor:v17v48-annealed}
On a graph of maximum degree \(D\), if \(\chi_e\leq\chi_*\), then
\[
\boxed{
 \mathbb E_{\nu_G}
 |\langle\ell,m_1(y)\rangle|^2
 \leq
 \ell^{\mathsf T}A_0^{-1}\ell\,
 \frac{D\chi_*^2L_H}{2J\vartheta}\,n.}            \tag{48.8}
\]
\end{corollary}

\begin{proof}
Bound each maximum in (48.6) by the sum of its two endpoint quartics.
Each vertex is then charged at most \(D\) times.  Apply (42.14) and
average.
\end{proof}

\subsection{Scale audit and the surviving localization problem}

For the four-dimensional torus and test vector of V46,
Proposition~\ref{prop:v17v46-hminus} makes the first factor in (48.8)
uniform.  However
\[
 \frac{n}{J(a)}
 =
 \frac{L^4}{\bar J}a^{-6}.                         \tag{48.9}
\]
Thus the force/curvature quotient closes the previously unknown
tilted path moment, but global Cauchy--Schwarz still loses six powers
of the lattice spacing.

\begin{countertheorem}[Global energy allocation remains insufficient]
\label{ctr:v17v48-global}
The chain (46.7), (48.5), (48.3), followed by a single global
Cauchy--Schwarz inequality and only the one-site V41 moment, does not
produce a scale-uniform bound for the random conditional mean under
(44.10).
\end{countertheorem}

\begin{proof}
Its resulting estimate is exactly (48.8).  The inverse-Hessian test
factor stays finite, while the explicit coefficient supplied by the
method is proportional to \(n/J(a)=L^4\bar J^{-1}a^{-6}\).  This proves
failure of the stated estimate, not divergence of the conditional
mean itself.
\end{proof}

The obstruction is now spatial rather than probabilistic: the local
quartic moment controls every summand in (48.3), but (48.8) charges
all edges equally.  V47 shows that a naive entrywise comparison with
the baseline Green kernel is false.  The next step must expose a
different exact divergence structure before applying localization.

\begin{assessment}[What V48 closes]
\label{ass:v17v48-status}
V48 proves a sharp, pointwise ratio between the Higgs tilt force and
the conformal curvature it generates.  This removes the uncontrolled
interpolation moment in V46 and gives a finite annealed
conditional-mean estimate using only the V41 quartic moment.  Its
remaining \(a^{-6}\) debit comes solely from allocating the response
energy over the whole lattice.
\end{assessment}

\begin{decision}[V49: Schwinger--Dyson divergence equation]
\label{dec:v17v48-next}
Move upstream from the Witten interpolation.  Integrate the
conditional density by parts to derive an exact equation for
\(m_1(y)\) in which the full nonlinear conformal and Higgs forces
appear under \(B^{\mathsf T}\).  Invert the deterministic operator
\(2\kappa_gI+4JL\) only after this divergence has been exposed, so the
killed-walk edge weights of V47 enter exactly rather than by an invalid
random-conductance comparison.
\end{decision}

\begin{firewall}[Claim boundary after V48]
\label{fire:v17v48}
The force/curvature estimate is a finite-volume theorem for the
declared positive conformal model.  Its global version is not a
continuum tightness estimate.  Gauge cofinality, reflection
positivity, a physical Einstein action, and the Einstein--SM continuum
limit remain open.
\end{firewall}

\section*{Version V48 append-only record}
\addcontentsline{toc}{section}{Version V48 append-only record}
The complete V47 source was retained byte for byte before this
continuation.  It had \(449453\) bytes, \(13746\) lines, and SHA--256
\[
 \texttt{922E0973A077BDB71647CC9456D6C3ED6BD497CA0F5485F2F9394290E2B68860}.
\]
V48 proved the sharp force/curvature quotient and eliminated the
unknown tilted-path moment, leaving a purely spatial volume debit.


\section{V49: the Schwinger--Dyson current equation}
\label{sec:v17v49}

\subsection{Expose the divergence before inverting}

At full tilt, omit constants and write the conditional quotient
potential as
\[
 \Phi_y(r)=
 \kappa_g\|r\|_2^2+
 J\sum_e\cosh(2d_e)+
 \sum_eT_e(r,y).                                  \tag{49.1}
\]
Define the nonlinear mismatch remainder and total residual edge
current by
\[
\begin{aligned}
 N_e(d_e)&=2J\{\sinh(2d_e)-2d_e\},\\
 F_e(y)&=\mathbb E_y\{N_e(d_e)+q_e(r,y)\}.         \tag{49.2}
\end{aligned}
\]
The linear part \(4Jd_e\) has been removed from \(N_e\).

\begin{theorem}[Exact conditional-mean current equation]
\label{thm:v17v49-current}
The conditional quotient mean \(m(y)=\mathbb E_y r\) satisfies
\[
\boxed{
 (2\kappa_gI+4JL)m(y)
 =
 -B^{\mathsf T}F(y).}                             \tag{49.3}
\]
Consequently, for every \(\ell\in\mathcal H_0\),
\[
\boxed{
 \langle\ell,m(y)\rangle
 =
 -\left\langle
 BA_0^{-1}\ell,F(y)
 \right\rangle_{\ell^2(\mathcal E)}.}             \tag{49.4}
\]
\end{theorem}

\begin{proof}
The cosh and hopping tails make every component of
\(\nabla\Phi_y\) integrable and make the boundary terms vanish in
finite-dimensional integration by parts.  Hence
\[
 0=\mathbb E_y\nabla\Phi_y
 =
 2\kappa_gm(y)+
 B^{\mathsf T}\mathbb E_y
 \{2J\sinh(2d)+q\}.                               \tag{49.5}
\]
Since \(\mathbb E_y d=Bm(y)\), insert
\[
 2J\sinh(2d)=4Jd+N(d)
\]
into (49.5).  This gives (49.3).  The operator
\(A_0=2\kappa_gI+4JL\) is positive definite on
\(\mathcal H_0\); invert (49.3), pair with \(\ell\), and move
\(B^{\mathsf T}\) to the other side, proving (49.4).
\end{proof}

Formula (49.4) realizes the route that failed for the random Witten
inverse in V47: the inverse is now exactly the deterministic massive
operator, because integration by parts exposed the full force as a
divergence first.

\subsection{Only the cut component of the current is visible}

Let
\[
 P_{\rm cut}=BL^\dagger B^{\mathsf T}              \tag{49.6}
\]
be the orthogonal projection of edge space onto
\(\operatorname{ran}B\).

\begin{corollary}[Cycle currents do not move the conformal mean]
\label{cor:v17v49-cut}
One has
\[
 \langle\ell,m(y)\rangle
 =
 -\left\langle
 BA_0^{-1}\ell,P_{\rm cut}F(y)
 \right\rangle,                                   \tag{49.7}
\]
and
\[
 \|BA_0^{-1}\ell\|_2^2
 \leq
 \frac1{4J}\ell^{\mathsf T}A_0^{-1}\ell.          \tag{49.8}
\]
\end{corollary}

\begin{proof}
The first vector in the edge pairing belongs to
\(\operatorname{ran}B\), proving (49.7).  If
\(u=A_0^{-1}\ell\), then
\[
 4J\|Bu\|_2^2
 \leq
 2\kappa_g\|u\|_2^2+4J\|Bu\|_2^2
 =
 \ell^{\mathsf T}A_0^{-1}\ell,                   \tag{49.9}
\]
which is (49.8).
\end{proof}

The projection in (49.7) is material: any large circulation in
\(\ker B^{\mathsf T}\) is exactly invisible to \(m(y)\).  The global
force norm in V48 charged that irrelevant sector.

\subsection{A scale-exact current criterion}

Return to the four-dimensional periodic family of V44 and let
\(\nu_a\) be its final Higgs--gauge marginal.  Define
\[
 F_a(y)=
 \mathbb E_y\{N(d)+q(r,y)\}.                       \tag{49.10}
\]

\begin{theorem}[Sufficient current scaling for conformal smearing]
\label{thm:v17v49-criterion}
Suppose there is \(C_F<\infty\), independent of \(a\), such that
\[
\boxed{
 \mathbb E_{\nu_a}
 \|P_{\rm cut}F_a(y)\|_2^2
 \leq C_Fa^2.}                                    \tag{49.11}
\]
Then for every fixed smooth continuum test function \(f\), with
\(\ell_x^{(a)}=a^4f(ax)\) and its constant mode removed,
\[
 \sup_a
 \mathbb E_{\nu_a}
 |\langle\ell^{(a)},m_a(y)\rangle|^2
 <\infty.                                         \tag{49.12}
\]
Together with V44, this gives tightness of every finite family of
full conformal smeared observables
\[
 a^4\sum_x f_j(ax)s_x,\qquad j=1,\ldots,k.        \tag{49.13}
\]
\end{theorem}

\begin{proof}
Equations (49.7)--(49.8), deterministic Cauchy--Schwarz, and (49.11)
give
\[
\begin{aligned}
 \mathbb E_{\nu_a}
 |\langle\ell^{(a)},m_a(y)\rangle|^2
 &\leq
 \frac{
 (\ell^{(a)})^{\mathsf T}A_{0,a}^{-1}\ell^{(a)}
 }{4J(a)}
 \mathbb E_{\nu_a}\|P_{\rm cut}F_a\|_2^2\\
 &\leq
 \frac{C_F}{4\bar J}
 (\ell^{(a)})^{\mathsf T}A_{0,a}^{-1}\ell^{(a)}.  \tag{49.14}
\end{aligned}
\]
Proposition~\ref{prop:v17v46-hminus} bounds the last factor uniformly,
proving (49.12).

Decompose each test function into its constant and mean-zero parts.
The common mode is tight by (44.15), the conditionally centered
quotient fluctuation is tight by (44.14), and its random conditional
mean is tight by (49.12).  A finite union and Chebyshev's inequality
give (49.13).
\end{proof}

\begin{corollary}[A local sufficient form]
\label{cor:v17v49-local}
Since the number of edges is \(O(a^{-4})\), the bound
\[
 \sup_{a,e}
 \mathbb E_{\nu_a}
 |(P_{\rm cut}F_a)_e|^2
 \leq C_{\rm loc}a^6                               \tag{49.15}
\]
implies (49.11).
\end{corollary}

\begin{proof}
Sum (49.15) over \(O(a^{-4})\) edges.
\end{proof}

\subsection{Why the exponent \(a^3\) is the correct local target}

For a smooth comparison configuration,
\[
 y_x=aH(ax),\qquad d_e=O(a).
\]
Across one edge, Taylor expansion gives
\[
 q_e
 =
 \chi_e\{a_{e_-}e^{d_e}-a_{e_+}e^{-d_e}\}
 =O(a^3),                                         \tag{49.16}
\]
because \(a_x=\|y_x\|^2=O(a^2)\) and the endpoint difference is
\(O(a^3)\).  Also,
\[
 N_e(d_e)
 =
 \frac{8J(a)}3d_e^3+O(J(a)d_e^5)
 =O(a^5).                                         \tag{49.17}
\]
Thus (49.15) asks the random cut-current for exactly the squared
scaling suggested by the smooth Higgs current.  Equations
(49.16)--(49.17) are only a consistency check: the lattice Gibbs
fields are not known to be smooth, so they do not prove (49.15).

\begin{assessment}[What V49 closes]
\label{ass:v17v49-status}
V49 replaces the random Witten inverse by an exact deterministic
massive Green operator through the conditional Schwinger--Dyson
identity.  It removes all cycle current, proves a scale-exact
sufficient criterion, and identifies the local target
\(\|(P_{\rm cut}F)_e\|_{L^2}=O(a^3)\).  Establishing that target for
the interacting random fields is now the first continuum bottleneck.
\end{assessment}

\begin{decision}[V50: companion audit and renormalized current test]
\label{dec:v17v49-next}
Perform the mandatory YM/NS companion audit.  Then test (49.15) by
expanding the gauge-covariant Higgs current into its lattice divergence
and fluctuation parts.  Determine whether centering or a local Ward
identity removes the order-one rough-field contribution; if it does
not, record the necessary current renormalization rather than assuming
smooth-field scaling.
\end{decision}

\begin{firewall}[Claim boundary after V49]
\label{fire:v17v49}
The current equation and criterion are rigorous for the declared
finite-volume conformal conditional law.  The \(a^3\) current estimate
is a required hypothesis, not a proved random-field bound.  Full
continuum tightness, gauge cofinality, reflection positivity, the
physical Einstein action, and the Einstein--SM continuum limit remain
open.
\end{firewall}

\section*{Version V49 append-only record}
\addcontentsline{toc}{section}{Version V49 append-only record}
The complete V48 source was retained byte for byte before this
continuation.  It had \(456123\) bytes, \(13958\) lines, and SHA--256
\[
 \texttt{7056DFBCDAD08EED4787CE9C0BCF16DEA59964E20A7ED5C395AC35772D687003}.
\]
V49 derived the exact Schwinger--Dyson current equation and the
scale-exact cut-current criterion.


\section{V50: mandatory companion audit at the current criterion}
\label{sec:v17v50}

\subsection{Exact companion checkpoints}

The newest active companion versions were inspected read-only after
V49.  Their checkpoints are
\[
\begin{array}{lll}
\text{programme}&\text{file}&\text{SHA--256}\\
\text{YM}&
\texttt{Renewal\_Yang\_Mills\_Mass\_Gap\_Research\_Notes\_V86.tex}&
\texttt{3AFC41DE67624728FAE1B7A5F90CB6471}\\
&&\texttt{0D1EE6D4B96620CCE8D28AE10FEE8E8}\\
\text{NS}&
\texttt{Renewal\_NS\_Stability\_Research\_Notes\_V26.tex}&
\texttt{82C887F8C2C69E4F28FAD7C3DC8DE5B}\\
&&\texttt{9099CB5A4BF431EBA1FFF3E91935978AD}.
\end{array}                                        \tag{50.1}
\]
At inspection, YM V86 had \(684946\) bytes and \(18701\) lines; NS
V26 had \(278283\) bytes and \(5319\) lines.  The shared folder also
contained YM V85.  It was left untouched; the higher active version
V86 is the checkpoint used here.

\subsection{Changes since the V45 audit}

YM V85 performed its own scheduled audit and extended the seventh
orientation to a first-coordinate two-step slab.  YM V86 added one
new exact cap chart and closed the full first quarter face of the
\((+-+)\) orientation.  Its other two coordinate directions still
have only seed depth.  Thus the seventh cell, its inversion partner,
full-transverse-torus positivity, cofinal control, reflection
positivity, and the Yang--Mills mass gap remain open.

NS remains at V26.  Its exact rank-one Oseen-kernel contractions and
its stated absence of a global regularity conclusion are unchanged
from the V45 audit.

\begin{theorem}[V50 companion-audit decision]
\label{thm:v17v50-decision}
No companion result proves the cut-current estimate (49.11) or its
local sufficient form (49.15).  The SM route is therefore unchanged:
keep the gauge sector at finite-cutoff compactness, and test the
Higgs--conformal current using symmetries and local identities of the
SM measure itself.
\end{theorem}

\begin{proof}
The new YM theorem concerns exact positivity of a finite Wilson-symbol
face, not the second moment of the scalar edge current \(F_a\).  The
NS tensor norms concern a different PDE kernel.  Neither supplies a
hypothesis in Theorem~\ref{thm:v17v49-criterion}.  YM's unfinished
last orientation also leaves the gauge claim boundary from V45 in
force.
\end{proof}

\begin{assessment}[What V50 records]
\label{ass:v17v50-status}
The required cross-program review is complete.  YM has advanced from a
Boolean seed cube to one quarter face of its final orientation, while
NS is unchanged.  Neither changes the immediate SM current-scaling
bottleneck.
\end{assessment}

\begin{decision}[V51: centering and the one-site-moment no-go]
\label{dec:v17v50-next}
Use edge-reversal symmetry to determine the unconditional mean of the
residual current in (49.2).  Then test whether centering plus the V41
one-site exponential quartic moment can imply (49.15).  If it cannot,
give an explicit counterexample and identify the correlation or Ward
estimate that must be added.
\end{decision}

\begin{firewall}[Claim boundary after V50]
\label{fire:v17v50}
V50 is a read-only audit.  It does not import finite Wilson-symbol
positivity as a continuum gauge theorem, prove the SM current
criterion, or establish the Einstein--Standard-Model continuum.
\end{firewall}

\section*{Version V50 append-only record}
\addcontentsline{toc}{section}{Version V50 append-only record}
The complete V49 source was retained byte for byte before this
continuation.  It had \(463633\) bytes, \(14215\) lines, and SHA--256
\[
 \texttt{86C238B40FDDE8BD9EA9ACC79E6ADDE9A19F7931C8AFDB25F6C0F4AF078D51F0}.
\]
V50 completed the mandatory companion audit and retained the
cut-current acquisition order.


\section{V51: reflection centering and the local-moment no-go}
\label{sec:v17v51}

\subsection{The residual current is centered}

Work on the homogeneous hypercubic periodic family of V44.  For each
unoriented nearest-neighbor edge, assume the periodic lattice
reflection that interchanges its endpoints preserves the couplings and
the full finite-volume action.  Gauge links are pulled back with the
usual inversion when their orientation is reversed.

\begin{theorem}[Exact edge-reflection centering]
\label{thm:v17v51-centering}
For every oriented edge \(e\), the residual conditional current from
(49.2) obeys
\[
\boxed{
 \mathbb E_{\nu_a}F_e(y)=0.}                       \tag{51.1}
\]
Consequently,
\[
 \mathbb E_{\nu_a}(P_{\rm cut}F_a)_e=0.           \tag{51.2}
\]
\end{theorem}

\begin{proof}
Let \(\mathcal R_e\) be a lattice reflection interchanging the two
endpoints of \(e\).  Pullback by \(\mathcal R_e\) preserves the full
joint law.  It sends
\[
 d_e\longmapsto-d_e,\qquad
 (a_{e_-},a_{e_+})\longmapsto(a_{e_+},a_{e_-}).    \tag{51.3}
\]
Both summands of the instantaneous residual current are odd under this
operation:
\[
\begin{aligned}
 N_e(d_e)&\longmapsto-N_e(d_e),\\
 q_e&=\chi_e(a_{e_-}e^{d_e}-a_{e_+}e^{-d_e})
 \longmapsto-q_e.                                 \tag{51.4}
\end{aligned}
\]
Conditional expectation in the conformal variables commutes with the
pullback, so
\[
 F_e(\mathcal R_ey)=-F_e(y).                      \tag{51.5}
\]
The marginal \(\nu_a\) is reflection invariant.  Integrating (51.5)
proves (51.1).  The projection \(P_{\rm cut}\) is deterministic and
linear, giving (51.2).
\end{proof}

Thus the left side of the local target (49.15) is a variance, with no
counterterm needed merely to remove its mean.  The required factor
\(a^6\), however, is a fluctuation statement and does not follow from
centering.

\subsection{One-site quartic control cannot imply current scaling}

\begin{countertheorem}[Centered cut-current no-go]
\label{ctr:v17v51-no-go}
Uniform one-site exponential quartic moments and exact edge-reflection
centering do not imply either (49.11) or (49.15), even when the current
already lies entirely in the cut space.
\end{countertheorem}

\begin{proof}
On each four-dimensional periodic torus, choose independent random
variables
\[
 y_x=
 \begin{cases}
 0,&\text{with probability }1/2,\\
 v,&\text{with probability }1/2,
 \end{cases}
 \qquad \|v\|=1.                                  \tag{51.6}
\]
For every \(\vartheta>0\),
\[
 \sup_{a,x}\mathbb E
 e^{\vartheta\|y_x\|^4}
 =\frac{1+e^\vartheta}{2}<\infty.                 \tag{51.7}
\]
Put \(A_x=\|y_x\|^2\), take \(d_e=0\), \(N_e=0\), and orient each
edge from \(e_-\) to \(e_+\).  The model current
\[
 Q_e=\chi(A_{e_-}-A_{e_+})                       \tag{51.8}
\]
is centered and changes sign under edge reflection.  It is also a
discrete gradient:
\[
 Q=-\chi BA,\qquad P_{\rm cut}Q=Q.                \tag{51.9}
\]
Independence gives, for every edge,
\[
 \mathbb E Q_e^2
 =\chi^2\mathbb E(A_{e_-}-A_{e_+})^2
 =\frac{\chi^2}{2}.                               \tag{51.10}
\]
There are \(O(a^{-4})\) edges, so
\[
 \mathbb E\|P_{\rm cut}Q\|_2^2
 =\frac{\chi^2}{2}|\mathcal E_a|
 \asymp a^{-4},                                   \tag{51.11}
\]
whereas (49.11) requires \(O(a^2)\).  This family satisfies precisely
the one-site moment and symmetry information under discussion but
violates both current targets.
\end{proof}

The counterexample is not asserted to be the declared interacting
Gibbs law.  It proves that the V41 marginal moment and Theorem
\ref{thm:v17v51-centering}, without a correlation identity from that
law, are logically insufficient.

\subsection{The renormalized current target}

Define the projected renormalized conditional current by
\[
 \mathcal J_{a,e}^{\rm ren}
 =
 a^{-3}(P_{\rm cut}F_a)_e.                        \tag{51.12}
\]
The local sufficient criterion (49.15) is exactly
\[
 \sup_{a,e}
 \mathbb E_{\nu_a}
 |\mathcal J_{a,e}^{\rm ren}|^2<\infty.           \tag{51.13}
\]
The exponent agrees with the smooth-field calculation (49.16), but
Countertheorem~\ref{ctr:v17v51-no-go} shows that it cannot be obtained
from a sitewise moment estimate.  It must come from at least one of:
\begin{enumerate}
\item a local Higgs Schwinger--Dyson or Ward identity that cancels the
      order-one radial increment;
\item a scale-dependent block-current estimate with cancellations
      before absolute values are taken;
\item an explicit current counterterm and a proof that the
      renormalized remainder has the \(a^3\) scaling.
\end{enumerate}
These are distinct mathematical acquisitions.  Reflection centering
alone is only the zeroth step in each.

\begin{assessment}[What V51 closes]
\label{ass:v17v51-status}
V51 proves exact unconditional centering of the residual conformal
current under edge reflection.  It then gives a cut-space counterexample
showing that this symmetry and the V41 one-site exponential quartic
moment cannot imply the required variance scaling.  The next input
must use a genuine local identity or multiscale correlation estimate
of the interacting measure.
\end{assessment}

\begin{decision}[V52: local Higgs radial identity]
\label{dec:v17v51-next}
Differentiate the rescaled Higgs density in the radial direction at
one site.  Express the diagonal hopping energy
\(\sum_{e\ni x}\chi_e e^{\pm d_e}\|y_x\|^2\) in terms of the fibre
dimension, the quartic and mass terms, and gauge-covariant cross
terms.  Subtract neighboring identities to test whether the current
in (49.2) gains a discrete derivative or whether a composite-current
counterterm is unavoidable.
\end{decision}

\begin{firewall}[Claim boundary after V51]
\label{fire:v17v51}
The no-go concerns inference from marginal moments and centering; it
is not a divergence theorem for the declared Gibbs current.  The
renormalized current bound, continuum tightness, gauge cofinality,
reflection positivity, physical Einstein action, and Einstein--SM
continuum limit remain open.
\end{firewall}

\section*{Version V51 append-only record}
\addcontentsline{toc}{section}{Version V51 append-only record}
The complete V50 source was retained byte for byte before this
continuation.  It had \(467359\) bytes, \(14307\) lines, and SHA--256
\[
 \texttt{0E01172241B4017AD0167AFFD07DE86978BE1DC25AF7233177C9E709956B1333}.
\]
V51 proved reflection centering and the exact one-site-moment
insufficiency for continuum current scaling.


\section{V52: local Higgs radial Ward identities}
\label{sec:v17v52}

\subsection{The sitewise inserted identity}

Let \(q_H\) denote the real dimension of one rescaled Higgs fibre.
For an oriented edge \(e=(v,w)\), define its two diagonal hopping
energies and gauge cross term by
\[
\begin{aligned}
 D_{v,e}&=\chi_e e^{d_e}\|y_v\|^2,\qquad
 D_{w,e}=\chi_e e^{-d_e}\|y_w\|^2,\\
 C_e&=\chi_e\operatorname{Re}
 \langle y_w,U_ey_v\rangle.                       \tag{52.1}
\end{aligned}
\]
Thus \(T_e=D_{v,e}+D_{w,e}\) and \(q_e=D_{v,e}-D_{w,e}\).
For either endpoint \(x\in e\), write \(D_{x,e}\) for the
corresponding term in (52.1).  The geometric hopping action on \(e\)
is
\[
 D_{v,e}+D_{w,e}-2C_e.                            \tag{52.2}
\]

Let \(X_x=y_x\mathbin{\cdot}\nabla_{y_x}\) be the radial vector field.
For the full finite-volume joint action \(S\), put
\[
\begin{aligned}
 \mathscr R_x
 &:=
 X_xS-q_H\\
 &=
 2m_H^2\|y_x\|^2+4\lambda_H\|y_x\|^4
 +2\sum_{e\ni x}(D_{x,e}-C_e)-q_H.                \tag{52.3}
\end{aligned}
\]

\begin{theorem}[Inserted local radial Ward identity]
\label{thm:v17v52-ward}
For every continuously differentiable insertion \(\Psi\) of polynomial
growth,
\[
\boxed{
 \mathbb E_{\mu_G}[\Psi\,\mathscr R_x]
 =
 \mathbb E_{\mu_G}[X_x\Psi].}                     \tag{52.4}
\]
In particular,
\[
 \mathbb E_{\mu_G}\mathscr R_x=0,\qquad
 \mathbb E_{\mu_G}
 \bigl[\|y_x\|^2\mathscr R_x\bigr]
 =
 2\mathbb E_{\mu_G}\|y_x\|^2.                     \tag{52.5}
\]
\end{theorem}

\begin{proof}
In real coordinates on the Higgs fibre,
\[
 \operatorname{div}_{y_x}(y_x\Psi e^{-S})
 =
 \{q_H\Psi+X_x\Psi-\Psi X_xS\}e^{-S}.             \tag{52.6}
\]
The positive quartic term and V39 coercivity make the boundary
integral vanish and dominate every polynomial insertion.  Integrate
(52.6) over \(y_x\), then over all remaining fields.  Rearrangement
gives (52.4).  Taking \(\Psi=1\) and
\(\Psi=\|y_x\|^2\), for which
\(X_x\Psi=2\|y_x\|^2\), gives (52.5).
\end{proof}

\subsection{A two-site identity for the diagonal current}

Fix \(e=(v,w)\), and define the residual star difference with the
distinguished edge omitted:
\[
\begin{aligned}
 K_e={}&
 2m_H^2(\|y_v\|^2-\|y_w\|^2)
 +4\lambda_H(\|y_v\|^4-\|y_w\|^4)\\
 &+2\sum_{\substack{f\ni v\\f\ne e}}(D_{v,f}-C_f)
 -2\sum_{\substack{f\ni w\\f\ne e}}(D_{w,f}-C_f).               \tag{52.7}
\end{aligned}
\]
The cross term \(C_e\) cancels between the two endpoint radial scores,
and therefore
\[
 \mathscr R_v-\mathscr R_w=2q_e+K_e.              \tag{52.8}
\]

\begin{theorem}[Exact current--star Ward identity]
\label{thm:v17v52-edge}
For every edge \(e=(v,w)\),
\[
\boxed{
 \mathbb E_{\mu_G}
 [q_e(\mathscr R_v-\mathscr R_w)]
 =
 2\mathbb E_{\mu_G}T_e,}                          \tag{52.9}
\]
or equivalently,
\[
\boxed{
 \mathbb E_{\mu_G}q_e^2
 =
 \mathbb E_{\mu_G}T_e
 -\frac12\mathbb E_{\mu_G}(q_eK_e).}              \tag{52.10}
\]
\end{theorem}

\begin{proof}
Apply (52.4) with \(x=v,\Psi=q_e\) and with
\(x=w,\Psi=q_e\).  Since only the corresponding endpoint norm is
differentiated,
\[
 X_vq_e=2D_{v,e},\qquad X_wq_e=-2D_{w,e}.         \tag{52.11}
\]
Subtract the two Ward identities to obtain
\[
 \mathbb E[q_e(\mathscr R_v-\mathscr R_w)]
 =
 2\mathbb E(D_{v,e}+D_{w,e}),
\]
which is (52.9).  Insert (52.8) and divide by two, proving (52.10).
\end{proof}

Equation (52.10) is an exact cancellation identity.  To reach the
continuum target
\(\mathbb E q_e^2=O(a^6)\), the two order-larger terms on its right
must cancel to that accuracy:
\[
 \mathbb E(q_eK_e)
 =
 2\mathbb E T_e+O(a^6).                           \tag{52.12}
\]
Neither V41 nor reflection centering estimates this cancellation.  This is the hopping-current part of (49.2); the nonlinear mismatch remainder still requires its own scale estimate.

\subsection{What the site identity can and cannot determine}

With \(\Psi=1\), (52.4) controls the star sum
\[
 \sum_{e\ni x}\mathbb E(D_{x,e}-C_e)
 =
 \frac{q_H}{2}
 -m_H^2\mathbb E\|y_x\|^2
 -2\lambda_H\mathbb E\|y_x\|^4.                  \tag{52.13}
\]
It does not isolate any one endpoint diagonal \(D_{x,e}\).  This loss
is algebraic: on every vertex of degree at least two, half-edge
increments can be increased on one incident edge and decreased on
another while leaving their star sum unchanged.  Small such
perturbations preserve positivity of the \(D_{x,e}\) but change the
edge differences \(q_e\).  Thus the uninserted identities (52.13)
cannot imply an edge-current variance bound.

The inserted identity (52.10) contains the missing edge information,
but transfers it to the mixed star correlation
\(\mathbb E(q_eK_e)\).  This is a genuine nearest-neighbor correlation
quantity.  Estimating its cancellation requires the coupled Gibbs law,
not only one-site tails.

\begin{assessment}[What V52 closes]
\label{ass:v17v52-status}
V52 derives the full inserted radial Ward identity for the declared
Higgs--gauge--conformal action and an exact two-site formula for the
diagonal conformal current variance.  It shows that the desired
\(a^3\) current scaling is equivalent, at this level, to the
quantified cancellation (52.12).  The ordinary sitewise Ward identity
controls only star sums and cannot supply that cancellation.
\end{assessment}

\begin{decision}[V53: centered two-site Ward expansion]
\label{dec:v17v52-next}
Expand the mixed term in (52.10) into on-site, adjacent-edge, and gauge
cross correlations.  Use reflection to remove odd expectations and
V41 to bound the integrable remainders.  Determine whether the leading
piece equals \(2\mathbb ET_e\) by a second local integration by parts;
if a nonzero contact term remains, identify it as the composite-current
counterterm required before continuum scaling.
\end{decision}

\begin{firewall}[Claim boundary after V52]
\label{fire:v17v52}
The Ward identities are exact finite-volume statements for the
declared bosonic model.  Equation (52.12) is a required cancellation,
not an established estimate.  Fermionic contributions, gauge
cofinality, continuum tightness, reflection positivity, the physical
Einstein action, and the Einstein--SM continuum limit remain open.
\end{firewall}

\section*{Version V52 append-only record}
\addcontentsline{toc}{section}{Version V52 append-only record}
The complete V51 source was retained byte for byte before this
continuation.  It had \(473809\) bytes, \(14491\) lines, and SHA--256
\[
 \texttt{6D4CBBC984E8A8A61927534918FC3588B78304F12F8AC9276CCC0B8B1729AC34}.
\]
V52 proved the local radial and two-site current Ward identities and
isolated their exact continuum cancellation target.


\section{V53: the contact-current obstruction and the block criterion}
\label{sec:v17v53}

\subsection{An exact two-site test of the Ward cancellation}

Freeze the conformal difference at \(d=0\), take a one-dimensional real
Higgs fibre and \(U=1\), and consider the two-site probability density
\[
 d\mu_2(y,z)
 =
 Z_2^{-1}
 \exp\{-\lambda(y^4+z^4)-\chi(y-z)^2\}\,dy\,dz,
 \qquad \lambda,\chi>0.                           \tag{53.1}
\]
This is the positive quartic plus geometric hopping block of the
declared model.  Define
\[
\begin{aligned}
 D_y&=\chi y^2,\qquad D_z=\chi z^2,\qquad
 C=\chi yz,\\
 q&=D_y-D_z=\chi(y^2-z^2),\qquad
 T=D_y+D_z.                                      \tag{53.2}
\end{aligned}
\]

\begin{theorem}[A nonzero local contact current]
\label{thm:v17v53-contact}
The two-site law has every quartic exponential moment with parameter
strictly below \(\lambda\), is invariant under endpoint exchange, and
satisfies
\[
 \mathbb E_{\mu_2}q=0,\qquad
 0<\mathbb E_{\mu_2}q^2<\infty.                   \tag{53.3}
\]
Its exact two-site Ward identity is
\[
\boxed{
 \mathbb E_{\mu_2}q^2
 =
 \mathbb E_{\mu_2}T
 -2\lambda\,
 \mathbb E_{\mu_2}
 \{q(y^4-z^4)\}.}                                 \tag{53.4}
\]
All quantities in (53.4) are independent of a lattice spacing.
\end{theorem}

\begin{proof}
Quartic coercivity proves the moment assertion.  Endpoint exchange
sends \(q\) to \(-q\), proving its mean is zero.  The polynomial
\(q^2\) is nonnegative and is strictly positive away from
\(y^2=z^2\).  That exceptional set has two-dimensional Lebesgue
measure zero, while the density in (53.1) is strictly positive.
Therefore \(\mathbb E q^2>0\); finiteness follows from the quartic
tail.

For this one-edge graph, formula (52.7) has no adjacent-edge terms and
the mass is zero, so
\[
 K=4\lambda(y^4-z^4).
\]
Insert this into (52.10), giving (53.4).
\end{proof}

\begin{countertheorem}[A local Ward identity does not force
\(a^3\) scaling]
\label{ctr:v17v53-local}
The radial Ward identities of V52, reflection centering, and uniform
quartic exponential moments do not imply
\(\mathbb E q_e^2=O(a^6)\) for a dimensionless lattice Higgs current
with fixed four-dimensional \(\lambda,\chi\).
\end{countertheorem}

\begin{proof}
Use the same law (53.1) at every \(a\).  It satisfies all three listed
properties, including the exact inserted identity (53.4), while
\(\mathbb E q^2\) is the fixed strictly positive number in (53.3).
\end{proof}

This counterexample again concerns an implication, not the full
hypercubic Gibbs law.  It shows that the pointwise sufficient
condition (49.15) is too strong to serve as the primary acquisition:
order-one microscopic current fluctuations can survive even when all
local Ward identities hold.  A continuum current must be tested after
spatial summation and cancellation.

\subsection{The exact test-dependent block-current criterion}

For a continuum test function \(f\), let
\[
 \ell_x^{(a)}=a^4f(ax),\qquad
 w_a(f)=BA_{0,a}^{-1}\ell^{(a)}.                  \tag{53.5}
\]
By edge reflection, \(\mathbb E_{\nu_a}F_a=0\) componentwise.
Define the block-current variance
\[
 \mathfrak C_a(f)
 =
 \mathbb E_{\nu_a}
 \left|
 \sum_ew_{a,e}(f)F_{a,e}(y)
 \right|^2.                                      \tag{53.6}
\]

\begin{theorem}[Block-current continuum criterion]
\label{thm:v17v53-block}
If, for every member \(f\) of a countable dense family of smooth
torus test functions,
\[
\boxed{
 \sup_a\mathfrak C_a(f)<\infty,}                  \tag{53.7}
\]
then the corresponding family of full conformal smeared observables
\[
 a^4\sum_xf(ax)s_x                               \tag{53.8}
\]
is tight in every finite collection.  Condition (53.7) allows
order-one single-edge current variances.
\end{theorem}

\begin{proof}
The exact current equation (49.4) says
\[
 \langle\ell^{(a)},m_a(y)\rangle
 =
 -\sum_ew_{a,e}(f)F_{a,e}(y).                    \tag{53.9}
\]
Thus (53.7) is precisely the required second-moment bound for the
random conditional mean.  The conditionally centered quotient
fluctuation is tight by (44.14), and the common mode is tight by
(44.15).  Summing these three pieces and applying Chebyshev proves
finite-family tightness.

No estimate of an individual \(F_{a,e}^2\) was used.  Hence
microscopic variances may remain of order one if their covariances
cancel after pairing with \(w_a(f)\).
\end{proof}

Let
\[
 \Sigma_a(e,e')
 =
 \operatorname{Cov}_{\nu_a}(F_{a,e},F_{a,e'}).    \tag{53.10}
\]
Then the criterion is the exact quadratic form
\[
 \mathfrak C_a(f)
 =
 w_a(f)^{\mathsf T}\Sigma_aw_a(f).                \tag{53.11}
\]
The local criterion (49.15) bounded this form by
\(\|w_a(f)\|_2^2\operatorname{tr}\Sigma_a\), losing all signed
off-diagonal covariance.  V53 therefore moves the continuum problem
from a false pointwise smallness target to the low-momentum covariance
of the conserved residual current.

\begin{assessment}[What V53 closes]
\label{ass:v17v53-status}
V53 proves, on an exact positive two-site Higgs block, that local Ward
identities leave a nonzero microscopic contact-current variance.
Therefore the \(a^3\) pointwise target from V49 cannot follow from
those identities.  The correct weaker acquisition is the
test-dependent block-current variance (53.7), which is exactly
equivalent to control of the conditional-mean contribution for each
continuum test.
\end{assessment}

\begin{decision}[V54: Fourier current-covariance threshold]
\label{dec:v17v53-next}
On the translation-invariant periodic family, diagonalize (53.11) in
lattice momentum.  Compute the multiplier of
\(BA_{0,a}^{-1}\ell^{(a)}\), and derive the necessary low-momentum
growth allowed for the longitudinal current structure factor.
Separate the contact term from the Ward-cancelled part before asking
for a scale-uniform bound.
\end{decision}

\begin{firewall}[Claim boundary after V53]
\label{fire:v17v53}
The two-site counterexample rules out a proof architecture; it does
not prove that the full interacting current lacks a continuum
distribution.  The block-current criterion is sufficient but remains
unproved for the lattice Gibbs laws.  Gauge cofinality, reflection
positivity, the physical Einstein action, and the Einstein--SM
continuum limit remain open.
\end{firewall}

\section*{Version V53 append-only record}
\addcontentsline{toc}{section}{Version V53 append-only record}
The complete V52 source was retained byte for byte before this
continuation.  It had \(480394\) bytes, \(14693\) lines, and SHA--256
\[
 \texttt{358C93CFB57554D303D186658EDF3DE6BD7D08B4A65D476EDC7F5A577B90F9D1}.
\]
V53 proved the local contact-current obstruction and replaced the
pointwise target by the exact block-current covariance criterion.


\section{V54: the longitudinal current structure-factor threshold}
\label{sec:v17v54}

\subsection{Unitary Fourier conventions}

Write the four-dimensional periodic lattice as
\((\mathbb Z/N\mathbb Z)^4\), with \(n=N^4\), and orient every edge in
a positive coordinate direction.  For a vertex field \(h\) and an
edge field \(G=(G_\mu)_{\mu=1}^4\), use the unitary transforms
\[
\begin{aligned}
 \widehat h(k)&=n^{-1/2}\sum_xe^{-ik\cdot x}h_x,\\
 \widehat G_\mu(k)&=n^{-1/2}\sum_xe^{-ik\cdot x}G_{x,\mu},
 \qquad
 k\in\frac{2\pi}{N}(\mathbb Z/N\mathbb Z)^4.       \tag{54.1}
\end{aligned}
\]
Put
\[
 b_\mu(k)=e^{ik_\mu}-1,\qquad
 \lambda(k)=\sum_{\mu=1}^4|b_\mu(k)|^2.           \tag{54.2}
\]
Then \(B\) has multiplier \(b(k)\), \(L=B^{\mathsf T}B\) has
multiplier \(\lambda(k)\), and
\[
 A_{0,a}(k)=2\kappa_g(a)+4J(a)\lambda(k).          \tag{54.3}
\]

Let the centered residual current have translation-invariant
covariance and define its \(4\)-by-\(4\) structure matrix by
\[
 S_a(k)_{\mu\nu}
 =
 \mathbb E_{\nu_a}
 [\widehat F_\mu(k)\overline{\widehat F_\nu(k)}]. \tag{54.4}
\]
It is positive semidefinite.  Only its longitudinal scalar
\[
 \mathscr S_a^{\parallel}(k)
 =
 b(k)^*S_a(k)b(k)                                 \tag{54.5}
\]
can enter the conformal conditional mean.

\begin{theorem}[Exact Fourier form of the block-current variance]
\label{thm:v17v54-fourier}
For every mean-zero vertex test vector \(\ell\),
\[
\boxed{
 \mathfrak C_a(\ell)
 =
 \sum_{k\ne0}
 \frac{|\widehat\ell(k)|^2
 \mathscr S_a^{\parallel}(k)}
 {\{2\kappa_g(a)+4J(a)\lambda(k)\}^2}.}           \tag{54.6}
\]
\end{theorem}

\begin{proof}
The Fourier transform of
\(w=BA_{0,a}^{-1}\ell\) is
\[
 \widehat w(k)=
 \frac{b(k)\widehat\ell(k)}
 {2\kappa_g(a)+4J(a)\lambda(k)}.                  \tag{54.7}
\]
Translation invariance diagonalizes the covariance between distinct
momenta.  Insert (54.7) into
\(\mathfrak C_a(\ell)=w^{\mathsf T}\Sigma_aw\), use Parseval, and
obtain (54.6).
\end{proof}

\subsection{The scale-uniform spectral envelope}

\begin{theorem}[Longitudinal structure-factor criterion]
\label{thm:v17v54-envelope}
The operator inequality
\[
\boxed{
 B^{\mathsf T}\Sigma_aB
 \preceq
 C_{\rm cur}\,a^{-4}A_{0,a}^2}                   \tag{54.8}
\]
on mean-zero vertex fields is equivalent to the pointwise momentum
envelope
\[
\boxed{
 \mathscr S_a^{\parallel}(k)
 \leq
 C_{\rm cur}\,a^{-4}
 \{2\kappa_g(a)+4J(a)\lambda(k)\}^2,
 \quad k\ne0.}                                   \tag{54.9}
\]
Either condition implies
\[
 \mathfrak C_a(\ell)
 \leq C_{\rm cur}a^{-4}\|\ell\|_2^2.              \tag{54.10}
\]
For \(\ell_x^{(a)}=a^4f(ax)\),
\[
 \mathfrak C_a(f)
 \leq C_{\rm cur}\|f\|_{2,a}^2,                  \tag{54.11}
\]
so the block-current criterion (53.7) holds.
\end{theorem}

\begin{proof}
Fourier diagonalization identifies the multiplier of
\(B^{\mathsf T}\Sigma_aB\) with (54.5), while that of \(A_{0,a}^2\)
is the squared denominator in (54.9).  This proves the equivalence.
Substitution in (54.6) gives (54.10).  Finally,
\[
 a^{-4}\|\ell^{(a)}\|_2^2
 =
 a^{-4}a^8\sum_x|f(ax)|^2
 =
 \|f\|_{2,a}^2,
\]
which proves (54.11).
\end{proof}

Under the canonical ledger
\(\kappa_g(a)=\bar\kappa_ga^4\),
\(J(a)=\bar Ja^2\), the right side of (54.9) is
\[
 C_{\rm cur}a^4
 \left\{
 2\bar\kappa_g+
 4\bar J\,\frac{\lambda(k)}{a^2}
 \right\}^2.                                     \tag{54.12}
\]
For a physical low mode \(k=ap+o(a)\), it is \(O(a^4)\).  For a mode
with \(\lambda(k)\asymp1\), it is \(O(1)\).  Thus the criterion
requires two additional powers of momentum in the longitudinal
infrared structure factor while allowing order-one ultraviolet
contact fluctuations.

\subsection{The V53 contact example passes after taking its gradient}

Let \(A_x\) be iid centered scalar variables of variance
\(\sigma_A^2\), and take
\[
 F=-\chi BA.                                      \tag{54.13}
\]
This is the centered version of the cut current in
Countertheorem~\ref{ctr:v17v51-no-go}.  Its structure matrix is
\[
 S_a(k)=\chi^2\sigma_A^2b(k)b(k)^*,               \tag{54.14}
\]
and hence
\[
 \mathscr S_a^{\parallel}(k)
 =
 \chi^2\sigma_A^2\lambda(k)^2.                    \tag{54.15}
\]

\begin{corollary}[Microscopic contact noise is compatible with the
block criterion]
\label{cor:v17v54-gradient}
The current (54.13) satisfies (54.9) uniformly in \(a\), with
\[
 C_{\rm cur}
 =
 \frac{\chi^2\sigma_A^2}{16\bar J^2}.             \tag{54.16}
\]
\end{corollary}

\begin{proof}
Since
\[
 a^{-4}A_{0,a}(k)^2
 \geq
 a^{-4}\{4\bar Ja^2\lambda(k)\}^2
 =
 16\bar J^2\lambda(k)^2,                          \tag{54.17}
\]
equations (54.15)--(54.16) give (54.9).
\end{proof}

This resolves the apparent conflict between V53 and continuum
smearing.  The individual edge variance in (54.13) is independent of
\(a\), but the current has an exact discrete derivative, producing the
\(\lambda(k)^2\) longitudinal factor needed in (54.9).

\subsection{The new precise acquisition}

For the actual residual current
\[
 F_e(y)=\mathbb E_y\{N_e(d_e)+q_e(r,y)\},          \tag{54.18}
\]
V49 proves only that its cycle component is irrelevant.  V52 gives
local radial identities, but has not yet produced the two powers of
low momentum in (54.9).  The next continuum problem is therefore the
operator estimate
\[
 B^{\mathsf T}\operatorname{Cov}_{\nu_a}(F_a)B
 \preceq
 C_{\rm cur}a^{-4}
 (2\kappa_g(a)I+4J(a)L_a)^2,                     \tag{54.19}
\]
or a test-family version sufficient for (54.11).

\begin{assessment}[What V54 closes]
\label{ass:v17v54-status}
V54 diagonalizes the exact block-current variance and derives the
scale-uniform longitudinal structure-factor threshold.  The threshold
allows nonvanishing microscopic contact variance and explains exactly
how a discrete-gradient current passes.  The SM continuum bottleneck
is now the concrete covariance inequality (54.19), rather than
pointwise current smallness.
\end{assessment}

\begin{decision}[V55: companion audit and current covariance]
\label{dec:v17v54-next}
Perform the mandatory YM/NS audit.  Then differentiate the inserted
Ward identity with a low-momentum radial insertion and test whether it
implies (54.19) for the hopping-current component.  Treat the nonlinear
mismatch current separately and retain any contact term explicitly.
\end{decision}

\begin{firewall}[Claim boundary after V54]
\label{fire:v17v54}
The Fourier criterion is exact and sufficient, but it has not been
proved for the interacting residual current (54.18).  It does not
establish a continuum measure, gauge cofinality, reflection positivity,
the physical Einstein action, or the Einstein--SM continuum limit.
\end{firewall}

\section*{Version V54 append-only record}
\addcontentsline{toc}{section}{Version V54 append-only record}
The complete V53 source was retained byte for byte before this
continuation.  It had \(487102\) bytes, \(14891\) lines, and SHA--256
\[
 \texttt{6046B2F0404669F371A0F74721F8D313BC60706938852BF548594D371DAF5656}.
\]
V54 derived the longitudinal structure-factor threshold and showed
why gradient contact noise is compatible with continuum smearing.


\section{V55: mandatory audit before the current-covariance attack}
\label{sec:v17v55}

\subsection{Companion checkpoints}

The active companion files were inspected read-only:
\[
\begin{array}{lll}
\text{programme}&\text{file}&\text{SHA--256}\\
\text{YM}&
\texttt{Renewal\_Yang\_Mills\_Mass\_Gap\_Research\_Notes\_V87.tex}&
\texttt{0B37169AEC2053E2246DE408A106D43A}\\
&&\texttt{BB92C3262A9914A02CE8918DBC71FBC6}\\
\text{NS}&
\texttt{Renewal\_NS\_Stability\_Research\_Notes\_V26.tex}&
\texttt{82C887F8C2C69E4F28FAD7C3DC8DE5B}\\
&&\texttt{9099CB5A4BF431EBA1FFF3E91935978AD}.
\end{array}                                        \tag{55.1}
\]
YM V87 has \(689959\) bytes and \(18839\) lines; NS V26 has
\(278283\) bytes and \(5319\) lines.

YM V87 certifies three new exact rational charts and extends the
seventh direct-plus orientation through second-coordinate depth
\(2t\), uniformly over its completed first coordinate.  The
second-coordinate cap, the third-coordinate extension, the inverted
eighth cell, full-transverse-torus positivity, cofinal control,
reflection positivity, and the mass gap remain open.  NS remains at
its exact rank-one Oseen-kernel comparison and has no global
regularity conclusion.

\begin{theorem}[V55 audit decision]
\label{thm:v17v55-decision}
Neither companion checkpoint implies the longitudinal SM current
covariance estimate (54.19).  No mathematical result is imported.
The next SM step remains a low-momentum inserted Ward calculation for
the residual current.
\end{theorem}

\begin{proof}
YM V87 is a finite-dimensional Wilson-symbol positivity and covering
theorem.  It does not estimate the covariance of the composite
Higgs--conformal current.  NS V26 concerns Oseen-kernel tensor norms.
Neither provides a premise of Theorem~\ref{thm:v17v54-envelope}.
\end{proof}

\begin{assessment}[What V55 records]
\label{ass:v17v55-status}
The scheduled cross-program review is complete.  The YM atlas has
advanced locally but has not reached a continuum input, while NS is
unchanged.  The longitudinal structure-factor threshold from V54
remains the active SM bottleneck.
\end{assessment}

\begin{decision}[V56: low-momentum inserted Ward identity]
\label{dec:v17v55-next}
Apply the Higgs radial integration-by-parts formula with a
Fourier-weighted radial insertion.  Derive the exact covariance
identity for the lattice divergence of the diagonal hopping current,
identify its contact term, and test whether the contact-subtracted
longitudinal structure factor has the two powers of momentum required
by (54.9).
\end{decision}

\begin{firewall}[Claim boundary after V55]
\label{fire:v17v55}
V55 is a read-only audit.  It does not promote the YM atlas to a gauge
continuum theorem, prove the SM current covariance estimate, or
construct the Einstein--Standard-Model continuum.
\end{firewall}

\section*{Version V55 append-only record}
\addcontentsline{toc}{section}{Version V55 append-only record}
The complete V54 source was retained byte for byte before this
continuation.  It had \(494219\) bytes, \(15130\) lines, and SHA--256
\[
 \texttt{DB0E9D4841E11B5D923291DD96F1F3AD1386BE148A451E784FDA5B6F98A9DE58}.
\]
V55 completed the mandatory YM/NS audit and retained the
longitudinal-current acquisition.


\section{V56: the low-momentum radial Ward sum rule}
\label{sec:v17v56}

\subsection{Divergence of the hopping current}

For an oriented edge \(e=(v,w)\), retain
\[
 h_e=D_{v,e}+D_{w,e}-2C_e,\qquad
 q_e=D_{v,e}-D_{w,e}.                             \tag{56.1}
\]
Let \(A\) be the unsigned incidence matrix, so
\((A^{\mathsf T}h)_x=\sum_{e\ni x}h_e\), and put
\[
 O_x=2m_H^2\|y_x\|^2+4\lambda_H\|y_x\|^4,\qquad
 G_x=O_x+(A^{\mathsf T}h)_x-q_H.                 \tag{56.2}
\]
The radial score (52.3) has the exact decomposition
\[
\boxed{
 \mathscr R=G-B^{\mathsf T}q,\qquad
 B^{\mathsf T}q=G-\mathscr R.}                   \tag{56.3}
\]

\begin{proof}
At the tail of \(e\),
\[
 2(D_{v,e}-C_e)=h_e+q_e,
\]
and at its head,
\[
 2(D_{w,e}-C_e)=h_e-q_e.
\]
With the convention that \(B^{\mathsf T}q\) is incoming minus
outgoing current, the total radial hopping term is
\(A^{\mathsf T}h-B^{\mathsf T}q\).  Insert this in (52.3).
\end{proof}

\subsection{The contact-subtracted Ward formula}

For a deterministic real vertex test \(\varphi\), define
\[
 \mathscr R(\varphi)=\sum_x\varphi_x\mathscr R_x,\quad
 G(\varphi)=\sum_x\varphi_xG_x,\quad
 X(\varphi)=\sum_x\varphi_xX_x.                  \tag{56.4}
\]

\begin{theorem}[Fourier-weighted inserted radial identity]
\label{thm:v17v56-contact}
For every deterministic \(\varphi\),
\[
\begin{aligned}
 \mathbb E\mathscr R(\varphi)^2
 &=\mathbb E[X(\varphi)^2S],                      \tag{56.5}\\
 \mathbb E[G(\varphi)\mathscr R(\varphi)]
 &=\mathbb E[X(\varphi)G(\varphi)],               \tag{56.6}
\end{aligned}
\]
and consequently
\[
\boxed{
 \mathbb E\{\varphi^{\mathsf T}B^{\mathsf T}q\}^2
 =
 \mathbb E G(\varphi)^2
 -2\mathbb E[X(\varphi)G(\varphi)]
 +\mathbb E[X(\varphi)^2S].}                     \tag{56.7}
\]
\end{theorem}

\begin{proof}
Sum the inserted Ward identity (52.4) against \(\varphi_x\):
\[
 \mathbb E[\Psi\mathscr R(\varphi)]
 =\mathbb E[X(\varphi)\Psi].                     \tag{56.8}
\]
First take \(\Psi=\mathscr R(\varphi)\).  Since
\(\mathscr R_x=X_xS-q_H\), one has
\[
 X(\varphi)\mathscr R(\varphi)=X(\varphi)^2S,
\]
which proves (56.5).  Taking \(\Psi=G(\varphi)\) gives (56.6).
Finally use (56.3), expand
\(\mathbb E\{G(\varphi)-\mathscr R(\varphi)\}^2\),
and substitute (56.5)--(56.6).
\end{proof}

The three terms on the right of (56.7) are individually contact-size.
The identity retains their exact cancellation before any absolute
value is taken.

\subsection{The second-moment infrared condition}

At a fixed lattice spacing, pass to a translation-invariant
infinite-volume Gibbs state supplied by V43 and set
\[
 D_x=(B^{\mathsf T}q)_x,\qquad
 C_a(z)=\operatorname{Cov}(D_0,D_z).              \tag{56.9}
\]
Assume for this subsection that
\[
 \sum_{z\in\mathbb Z^4}(1+|z|^4)|C_a(z)|<\infty. \tag{56.10}
\]
Its structure factor is
\[
 \widehat C_a(k)=\sum_ze^{-ik\cdot z}C_a(z)
 =b(k)^*S_a^q(k)b(k).                             \tag{56.11}
\]

\begin{theorem}[From the automatic zero to the required double zero]
\label{thm:v17v56-sumrule}
Reflection symmetry and the divergence form give
\[
 \widehat C_a(0)=0,\qquad
 \nabla\widehat C_a(0)=0.                         \tag{56.12}
\]
Define the second-moment tensor
\[
 M_{a,ij}^{(2)}
 =
 -\sum_z z_iz_jC_a(z).                            \tag{56.13}
\]
Then
\[
 \widehat C_a(k)
 =
 \frac12\sum_{i,j}M_{a,ij}^{(2)}k_ik_j
 +O\left(
 |k|^4\sum_z|z|^4|C_a(z)|
 \right).                                        \tag{56.14}
\]
In particular, a scale-uniform bound
\[
 \|M_a^{(2)}\|\leq C_2a^2,\qquad
 \sum_z|z|^4|C_a(z)|\leq C_4                      \tag{56.15}
\]
implies, for every fixed physical momentum \(p\),
\[
 \widehat C_a(ap)=O(a^4).                         \tag{56.16}
\]
This is the low-momentum power required by (54.9) for the hopping
current.
\end{theorem}

\begin{proof}
Equation (56.11) contains one factor \(b(k)\) and one factor
\(b(k)^*\), so its value at \(k=0\) is zero.  Equivalently,
\(\sum_zC_a(z)=0\).  Reflection makes \(C_a(z)=C_a(-z)\), hence the
linear moment vanishes.  Taylor's formula for \(e^{-ik\cdot z}\),
paired between \(z\) and \(-z\), gives (56.14), with the fourth
absolute moment controlling the remainder.  Substitute \(k=ap\) and
(56.15).
\end{proof}

The ordinary conservation zero gives only \(O(|k|^2)\).  V54 needs
the coefficient of that term to vanish or shrink like \(a^2\).  The
precise new Ward acquisition is therefore the second-moment estimate
in (56.15), not another one-site moment.

\subsection{Conditional expectation does not increase the quadratic
form}

Let
\[
 F_a^q(y)=\mathbb E[q\mid y]                      \tag{56.17}
\]
be the hopping part of the conditional residual current.

\begin{lemma}[Annealed covariance contraction]
\label{lem:v17v56-conditioning}
For every deterministic edge vector \(w\),
\[
 \operatorname{Var}_{\nu_a}
 \langle w,F_a^q\rangle
 \leq
 \operatorname{Var}_{\mu_a}\langle w,q\rangle.   \tag{56.18}
\]
\end{lemma}

\begin{proof}
This is the conditional-variance decomposition:
\[
 \operatorname{Var}(\langle w,q\rangle)
 =
 \mathbb E\operatorname{Var}(\langle w,q\rangle\mid y)
 +\operatorname{Var}
 \bigl(\mathbb E[\langle w,q\rangle\mid y]\bigr).
\]
Discard the first nonnegative term.
\end{proof}

Thus a uniform version of (56.15), inserted in (56.14) with the
all-momentum envelope of V54, would pass to the conditional hopping
current.  The nonlinear mismatch part
\(\mathbb E[N(d)\mid y]\) and its cross covariance remain separate.

\begin{assessment}[What V56 closes]
\label{ass:v17v56-status}
V56 derives the exact contact-subtracted radial Ward formula for the
divergence of the diagonal hopping current.  Under explicit fourth
covariance-moment summability, it proves that the V54 infrared bound
is reduced to the second-moment sum rule
\(\|M_a^{(2)}\|=O(a^2)\).  Conditioning on the Higgs field cannot
worsen the resulting quadratic form.
\end{assessment}

\begin{decision}[V57: evaluate the second-moment Ward coefficient]
\label{dec:v17v56-next}
Differentiate the contact identity (56.7) with respect to a small
Fourier momentum and express \(M_a^{(2)}\) as a weighted sum of
on-site, hopping, and gauge-cross correlations.  Test whether a second
radial Ward identity cancels its order-one part.  If the coefficient
survives, identify the precise composite-current subtraction needed
to enforce (56.15).
\end{decision}

\begin{firewall}[Claim boundary after V56]
\label{fire:v17v56}
The fourth covariance-moment summability and the \(O(a^2)\)
second-moment sum rule are hypotheses, not proved properties of the
interacting Gibbs state.  The mismatch-current part, complete
structure-factor envelope, continuum measure, gauge cofinality,
reflection positivity, physical Einstein action, and Einstein--SM
continuum limit remain open.
\end{firewall}

\section*{Version V56 append-only record}
\addcontentsline{toc}{section}{Version V56 append-only record}
The complete V55 source was retained byte for byte before this
continuation.  It had \(497456\) bytes, \(15212\) lines, and SHA--256
\[
 \texttt{89A879668F404D42022E5C0B10E72D3134EAF2F8D67BD8C63AB6DDB04B97BE1E}.
\]
V56 derived the contact-subtracted low-momentum Ward identity and its
second-moment sum-rule threshold.


\section{V57: zero-momentum susceptibility and the exact gradient split}
\label{sec:v17v57}

\subsection{The coefficient of the conservation zero}

Let \(S_a^q(k)\) be the edge-current structure matrix of the
instantaneous hopping current \(q\), and assume its covariance kernel
is absolutely summable, so \(S_a^q\) is continuous at \(k=0\).
Equation (56.11) reads
\[
 \widehat C_a(k)=b(k)^*S_a^q(k)b(k).              \tag{57.1}
\]

\begin{theorem}[The second-moment coefficient is the zero-mode
current susceptibility]
\label{thm:v17v57-zero}
For \(k\to0\),
\[
\boxed{
 \widehat C_a(k)
 =
 k^{\mathsf T}S_a^q(0)k+o(|k|^2).}               \tag{57.2}
\]
If reflection symmetry and the fourth covariance moment in (56.10)
hold, then comparison with (56.14) gives
\[
 M_a^{(2)}=2S_a^q(0)                              \tag{57.3}
\]
on the real longitudinal coordinate sector.  Hence the V56
second-moment acquisition is equivalent to
\[
 \|S_a^q(0)\|=O(a^2).                             \tag{57.4}
\]
\end{theorem}

\begin{proof}
The edge multiplier satisfies
\[
 b_\mu(k)=ik_\mu+O(|k|^2).                        \tag{57.5}
\]
Continuity of \(S_a^q\) at zero and (57.1) give (57.2).  Under
reflection symmetry the quadratic coefficient is real symmetric, and
the remainder after the quadratic term is of fourth order under
(56.10).  Comparing the coefficient of
\(\frac12k_ik_j\) in (56.14) with (57.2) proves (57.3), and (57.4)
follows.
\end{proof}

Thus a radial Ward calculation can prove the required double zero only
if it also proves that the integrated longitudinal current
susceptibility shrinks like \(a^2\).  The local contact identity alone
does not do so.

\subsection{Exact gradient and geometric remainder}

Assume homogeneous hopping coefficient \(\chi\), and put
\[
 A_x=\|y_x\|^2.
\]
For \(e=(v,w)\), decompose
\[
\begin{aligned}
 q_e
 &=
 \chi(A_ve^{d_e}-A_we^{-d_e})\\
 &=
 -\chi(BA)_e+g_e,                                 \tag{57.6}\\
 g_e
 &:=
 \chi\{A_v(e^{d_e}-1)-A_w(e^{-d_e}-1)\}.          \tag{57.7}
\end{aligned}
\]
This is an identity, with no small-field expansion.

Let
\[
 H_e(y)
 =
 \mathbb E_y\{N_e(d_e)+g_e(r,y)\}.                \tag{57.8}
\]

\begin{theorem}[Exact conditional residual-current split]
\label{thm:v17v57-split}
The residual current in (49.10) is
\[
\boxed{
 F_a(y)=-\chi BA+H(y).}                           \tag{57.9}
\]
The first term has two exact incidence multipliers.  All effects of
the nonlinear conformal mismatch and of the conformal weighting of the
Higgs norm are confined to \(H\).
\end{theorem}

\begin{proof}
The Higgs norm field \(A\) is fixed under the conditional conformal
expectation.  Insert (57.6) and
\(F=\mathbb E_y(N+q)\) into (49.10).
\end{proof}

\subsection{The scalar-susceptibility criterion for the gradient part}

Center \(A\), and define its scalar structure factor by
\[
 S_a^A(k)
 =
 \mathbb E_{\nu_a}|\widehat A(k)|^2.              \tag{57.10}
\]

\begin{theorem}[Uniform Higgs-norm susceptibility closes the gradient
current]
\label{thm:v17v57-gradient}
If
\[
 \sup_{a,k}S_a^A(k)\leq C_A<\infty,               \tag{57.11}
\]
then the gradient current \(-\chi BA\) satisfies the V54 envelope
(54.9) with
\[
 C_{\rm grad}
 =
 \frac{\chi^2C_A}{16\bar J^2}.                    \tag{57.12}
\]
\end{theorem}

\begin{proof}
The edge structure matrix of \(-\chi BA\) is
\[
 S_a^{\rm grad}(k)
 =
 \chi^2S_a^A(k)b(k)b(k)^*.
\]
Therefore its longitudinal scalar is
\[
 b(k)^*S_a^{\rm grad}(k)b(k)
 =
 \chi^2S_a^A(k)\lambda(k)^2
 \leq\chi^2C_A\lambda(k)^2.                      \tag{57.13}
\]
As in (54.17),
\[
 a^{-4}A_{0,a}(k)^2
 \geq16\bar J^2\lambda(k)^2.
\]
Equations (57.12)--(57.13) prove (54.9).
\end{proof}

The one-site V41 moment bounds \(S_a^A\) only after summing no
off-diagonal covariances.  It does not imply (57.11): an integrated
susceptibility is a correlation estimate.

\subsection{Remainder envelope}

Let \(S_a^H(k)\) be the structure matrix of the centered remainder
\(H\).  Suppose
\[
 b(k)^*S_a^H(k)b(k)
 \leq
 C_H^{\rm cur}a^{-4}A_{0,a}(k)^2.                 \tag{57.14}
\]

\begin{corollary}[Two independent acquisitions suffice]
\label{cor:v17v57-two}
Under (57.11) and (57.14), the full residual current \(F_a\) satisfies
(54.9) with
\[
 C_{\rm cur}
 \leq2(C_{\rm grad}+C_H^{\rm cur}).               \tag{57.15}
\]
Consequently the conformal smeared observables are tight in every
finite test family.
\end{corollary}

\begin{proof}
For any complex vector \(z\),
\[
 |z^*(F^{\rm grad}+H)|^2
 \leq2|z^*F^{\rm grad}|^2+2|z^*H|^2.
\]
Average, take \(z=b(k)\), and use
Theorem~\ref{thm:v17v57-gradient} and (57.14).
The last assertion is Theorems~\ref{thm:v17v54-envelope} and
\ref{thm:v17v53-block}.
\end{proof}

For smooth comparison fields, \(g_e=O(a^3)\) and
\(N_e=O(a^5)\), as in V49.  For the Gibbs fields, (57.14) remains a
composite-operator estimate.  It cannot be replaced by those formal
Taylor orders.

\begin{assessment}[What V57 closes]
\label{ass:v17v57-status}
V57 identifies the second Ward coefficient with the zero-momentum
edge-current susceptibility and splits the residual current exactly
into a Higgs-norm gradient and a geometric/mismatch remainder.  The
gradient part satisfies the complete continuum envelope once the
Higgs-norm susceptibility is uniformly bounded.  The remaining
acquisitions are now (57.11) and (57.14).
\end{assessment}

\begin{decision}[V58: Higgs-norm susceptibility]
\label{dec:v17v57-next}
Test whether the V41 conditional row can be strengthened to an
oscillation or gradient influence estimate for
\(A_x=\|y_x\|^2\).  Under an explicit high-temperature margin, derive
exponential covariance decay and (57.11).  If the broken-symmetry
Higgs potential prevents such a margin, state the exact obstruction
and retain susceptibility as a phase-selection hypothesis.
\end{decision}

\begin{firewall}[Claim boundary after V57]
\label{fire:v17v57}
Neither the uniform Higgs-norm susceptibility nor the remainder
envelope has been proved.  The decomposition does not establish full
continuum tightness, gauge cofinality, reflection positivity, the
physical Einstein action, or the Einstein--SM continuum limit.
\end{firewall}

\section*{Version V57 append-only record}
\addcontentsline{toc}{section}{Version V57 append-only record}
The complete V56 source was retained byte for byte before this
continuation.  It had \(504657\) bytes, \(15447\) lines, and SHA--256
\[
 \texttt{1FFDB2C63F60817B499049DE130F9FDBB185973D1A1F93692FBE7EEF4D937BBC}.
\]
V57 evaluated the infrared Ward coefficient and split the residual
current into its exact gradient and composite-remainder parts.


\section{V58: quenched Higgs-norm susceptibility and the mass-scaling
obstruction}
\label{sec:v17v58}

\subsection{Conditional Higgs convexity}

Fix the conformal and gauge variables.  Before expansion, the
geometric Higgs hopping action is a sum of nonnegative quadratic
squares.  Hence the Hessian of the full Higgs action in the realified
Higgs variables satisfies
\[
 \nabla_y^2S_H(y\mid r,U)
 \succeq2m_H^2I                                  \tag{58.1}
\]
whenever \(m_H^2>0\); the quartic Hessian and the hopping Hessian are
positive semidefinite.

Let
\[
 A_x=\|y_x\|^2,\qquad
 \bar A_x(r,U)=\mathbb E[A_x\mid r,U].             \tag{58.2}
\]

\begin{theorem}[Uniform quenched susceptibility at positive lattice
mass]
\label{thm:v17v58-quenched}
For every deterministic real vertex vector \(t\),
\[
\boxed{
 \mathbb E\operatorname{Var}
 \left(\sum_xt_xA_x\mathrel{\big|}r,U\right)
 \leq
 \frac{2M_{H,2}}{m_H^2}\sum_xt_x^2.}              \tag{58.3}
\]
Equivalently, the annealed expectation of the conditional covariance
operator of \(A\) is bounded by
\((2M_{H,2}/m_H^2)I\).
\end{theorem}

\begin{proof}
The Brascamp--Lieb inequality under the conditional Higgs law and
(58.1) give
\[
\begin{aligned}
 \operatorname{Var}
 \left(\sum_xt_xA_x\mathrel{\big|}r,U\right)
 &\leq
 \mathbb E\left[
 \nabla\left(\sum_xt_xA_x\right)^{\mathsf T}
 (\nabla_y^2S_H)^{-1}
 \nabla\left(\sum_xt_xA_x\right)
 \mathrel{\big|}r,U\right]\\
 &\leq
 \frac{2}{m_H^2}
 \sum_xt_x^2\mathbb E[A_x\mid r,U],               \tag{58.4}
\end{aligned}
\]
because \(\nabla_{y_x}A_x=2y_x\).
Average over \((r,U)\) and use (42.14).
\end{proof}

\subsection{The environmental mean is a separate susceptibility}

The total-variance identity gives, for every \(t\),
\[
\begin{aligned}
 \operatorname{Var}_{\mu_a}\left(\sum_xt_xA_x\right)
 ={}&
 \mathbb E\operatorname{Var}
 \left(\sum_xt_xA_x\mathrel{\big|}r,U\right)\\
 &+
 \operatorname{Var}_{r,U}
 \left(\sum_xt_x\bar A_x(r,U)\right).             \tag{58.5}
\end{aligned}
\]
On a translation-invariant torus, let
\(S_a^{\bar A}(k)\) be the structure factor of the centered field
\(\bar A(r,U)\).

\begin{corollary}[Exact annealed decomposition]
\label{cor:v17v58-decomposition}
For every momentum,
\[
\boxed{
 S_a^A(k)
 \leq
 \frac{2M_{H,2}}{m_H^2}
 +S_a^{\bar A}(k).}                               \tag{58.6}
\]
Thus a uniform bound on the environmental conditional-mean
susceptibility, together with a positive lattice mass bounded away
from zero, proves (57.11).
\end{corollary}

\begin{proof}
Apply (58.5) to the real and imaginary parts of the unitary Fourier
mode and use (58.3).  The second term is exactly the corresponding
quadratic form of \(S_a^{\bar A}\).
\end{proof}

This decomposition prevents a false annealed conclusion from the
quenched gap: conformal or gauge correlations can move the conditional
mean \(\bar A\) coherently even when the Higgs fluctuation at fixed
environment has a uniform Poincar\'e constant.

\subsection{Canonical mass scaling defeats the Hessian floor}

Under the four-dimensional ledger (44.11),
\[
 m_H^2(a)=\bar m_H^2a^2.                           \tag{58.7}
\]
Consequently, the first coefficient on the right of (58.6) becomes
\[
 \frac{2M_{H,2}}{\bar m_H^2}\,a^{-2}.             \tag{58.8}
\]

\begin{countertheorem}[The uniform-convexity susceptibility route
does not reach canonical scaling]
\label{ctr:v17v58-mass}
The Brascamp--Lieb estimate based only on the Hessian floor
\(2m_H^2(a)I\) cannot prove the scale-uniform Higgs-norm
susceptibility (57.11) under (58.7).
\end{countertheorem}

\begin{proof}
Its explicit output is (58.6), whose quenched coefficient is (58.8).
That coefficient diverges as \(a\downarrow0\).  This proves failure of
the stated estimate, not divergence of the actual susceptibility.
\end{proof}

Keeping \(m_H^2\geq m_0^2>0\) would make (58.3) uniform, but in the
canonical relation \(m_{\rm phys}^2=m_H^2/a^2\) it sends the physical
bare mass to infinity.  That branch cannot represent a finite-mass
interacting continuum without an additional critical tuning.

\subsection{The upstream acquisition}

The conditional single-site potential contains
\[
 \lambda_H\|y_x\|^4+c_x(r)\|y_x\|^2
 -2\langle h_x(y_{\ne x},U),y_x\rangle,
 \qquad c_x(r)\geq m_H^2(a).                      \tag{58.9}
\]
Even when \(m_H^2(a)\to0\), the quartic confines the site.  A
scale-safe replacement for (58.1) would be a Poincar\'e inequality
for the entire tilted quartic family in (58.9), uniform in the linear
field \(h_x\) and in \(a_x\geq0\).  Such a one-site theorem, combined
with a genuine intersite contraction estimate, could control the
quenched part without the factor \(a^{-2}\).  The environmental term
in (58.6) would still require conformal/gauge mixing.

\begin{assessment}[What V58 closes]
\label{ass:v17v58-status}
V58 proves a quenched Higgs-norm susceptibility bound and the exact
law-of-total-variance separation of its environmental component.  It
then proves that the direct uniform-Hessian argument degenerates as
\(a^{-2}\) at canonical Higgs mass scaling.  The next route must use
quartic confinement rather than the vanishing quadratic floor.
\end{assessment}

\begin{decision}[V59: tilted quartic one-site spectral gap]
\label{dec:v17v58-next}
Prove a Poincar\'e inequality for
\(\exp\{-\lambda\|y\|^4-a\|y\|^2+2h\cdot y\}\,dy\)
with a constant uniform over \(a\geq0\) and \(h\in\mathbb R^{q_H}\).
Use uniform convexity of the quartic Bregman remainder and a
one-dimensional radial/angular decomposition.  If the constant is not
uniform, identify the exact escaping parameter regime.
\end{decision}

\begin{firewall}[Claim boundary after V58]
\label{fire:v17v58}
Only the quenched positive-mass susceptibility has been proved.
Neither the tilted-quartic uniform gap nor the environmental
susceptibility is established.  The remainder current, continuum
tightness, gauge cofinality, reflection positivity, physical Einstein
action, and Einstein--SM continuum limit remain open.
\end{firewall}

\section*{Version V58 append-only record}
\addcontentsline{toc}{section}{Version V58 append-only record}
The complete V57 source was retained byte for byte before this
continuation.  It had \(511275\) bytes, \(15675\) lines, and SHA--256
\[
 \texttt{8C4995E4E8B52B7DC186666D7B6C1CCE83D5620BF53CB12785A9F067F3C3A6E9}.
\]
V58 proved the quenched susceptibility decomposition and the
canonical mass-floor obstruction.


\section{V59: a uniform spectral gap for every tilted quartic site}
\label{sec:v17v59}

\subsection{The tilted family}

For \(q\geq1\), \(\lambda>0\), \(a\geq0\), and
\(h\in\mathbb R^q\), let
\[
 d\pi_{a,h}(y)
 =
 Z_{a,h}^{-1}
 \exp\{-\lambda\|y\|^4-a\|y\|^2+2h\cdot y\}\,dy. \tag{59.1}
\]
The linear field is unrestricted.  This is the one-site conditional
family generated by the geometric Higgs hopping in the
nonnegative-mass branch.

\begin{lemma}[Uniform quartic monotonicity]
\label{lem:v17v59-monotone}
For
\[
 W_a(y)=\lambda\|y\|^4+a\|y\|^2,
\]
one has, for all \(y,z\in\mathbb R^q\),
\[
\boxed{
 (\nabla W_a(y)-\nabla W_a(z))\cdot(y-z)
 \geq
 \lambda\|y-z\|^4+2a\|y-z\|^2.}                  \tag{59.2}
\]
\end{lemma}

\begin{proof}
The standard \(p\)-monotonicity inequality at \(p=4\) is
\[
 (\|y\|^2y-\|z\|^2z)\cdot(y-z)
 \geq\frac14\|y-z\|^4.                            \tag{59.3}
\]
It follows directly by expanding in the plane spanned by \(y,z\);
equality occurs at \(z=-y\).  Multiply (59.3) by \(4\lambda\) and add
the quadratic contribution \(2a\|y-z\|^2\).
\end{proof}

\subsection{A two-copy covariance bound}

\begin{theorem}[Uniform centered second moment]
\label{thm:v17v59-covariance}
If \(Y\) has law \(\pi_{a,h}\), then
\[
\boxed{
 \operatorname{tr}\operatorname{Cov}(Y)
 \leq
 \sqrt{\frac{q}{2\lambda}}.}                      \tag{59.4}
\]
The bound is uniform in \(a\geq0\) and \(h\in\mathbb R^q\).
\end{theorem}

\begin{proof}
Let \(Y'\) be an independent copy and put
\[
 V(y)=W_a(y)-2h\cdot y.
\]
Integration by parts gives
\[
 \mathbb E[\partial_iV(Y)]=0,\qquad
 \mathbb E[Y_i\partial_iV(Y)]=1.                 \tag{59.5}
\]
Independence therefore yields
\[
\begin{aligned}
 &\mathbb E[
 (\nabla V(Y)-\nabla V(Y'))\cdot(Y-Y')]\\
 &\qquad=2q.                                      \tag{59.6}
\end{aligned}
\]
The linear tilt cancels from the gradient difference.  Apply
Lemma~\ref{lem:v17v59-monotone}:
\[
 \lambda\mathbb E\|Y-Y'\|^4\leq2q.                \tag{59.7}
\]
By Cauchy--Schwarz and the independent-copy identity,
\[
\begin{aligned}
 2\operatorname{tr}\operatorname{Cov}(Y)
 &=\mathbb E\|Y-Y'\|^2\\
 &\leq
 \{\mathbb E\|Y-Y'\|^4\}^{1/2}
 \leq\sqrt{\frac{2q}{\lambda}},
\end{aligned}
\]
which is (59.4).
\end{proof}

\subsection{Uniform Poincar\'e constant}

We use the following standard log-concave Cheeger--Poincar\'e theorem:
there is a universal numerical constant \(C_{\rm LC}\) such that every
log-concave probability law \(\pi\) on Euclidean space satisfies
\[
 C_P(\pi)
 \leq
 C_{\rm LC}\,
 \operatorname{tr}\operatorname{Cov}_\pi(Y).      \tag{59.8}
\]
This is the dimension-explicit localization bound; no KLS conjecture
is being assumed.

\begin{theorem}[Tilt-uniform quartic spectral gap]
\label{thm:v17v59-gap}
Every law (59.1) satisfies
\[
\boxed{
 \operatorname{Var}_{\pi_{a,h}}(f)
 \leq
 C_{q,\lambda}
 \int\|\nabla f\|^2\,d\pi_{a,h},\qquad
 C_{q,\lambda}
 :=
 C_{\rm LC}\sqrt{\frac{q}{2\lambda}}.}            \tag{59.9}
\]
The constant is independent of the quadratic coefficient and of the
linear field.
\end{theorem}

\begin{proof}
The potential in (59.1) is convex for \(a\geq0\), so
\(\pi_{a,h}\) is log-concave.  Apply (59.8) and then (59.4).
\end{proof}

The scaling \(C_{q,\lambda}\asymp\lambda^{-1/2}\) is the correct one:
the substitution \(y=\lambda^{-1/4}z\) multiplies a Poincar\'e
constant by \(\lambda^{-1/2}\).

\subsection{Application to a conditional Higgs site}

For one Higgs site conditioned on all other fields, the local
potential has form (58.9), with
\[
 q=q_H,\qquad\lambda=\lambda_H,\qquad
 a=c_x(r)\geq0,                                   \tag{59.10}
\]
and an arbitrary neighbor-generated \(h\).  Put
\[
 C_H^{\rm site}
 =
 C_{\rm LC}\sqrt{\frac{q_H}{2\lambda_H}}.         \tag{59.11}
\]

\begin{corollary}[Mass-independent averaged one-site variance]
\label{cor:v17v59-site}
In the nonnegative-mass branch,
\[
\boxed{
 \mathbb E\operatorname{Var}
 (A_x\mid\hbox{all fields outside }y_x)
 \leq
 4C_H^{\rm site}M_{H,2}.}                         \tag{59.12}
\]
The bound remains uniform when \(m_H^2(a)\downarrow0\).
\end{corollary}

\begin{proof}
Apply (59.9) to \(f(y_x)=\|y_x\|^2\):
\[
 \operatorname{Var}(A_x\mid\cdots)
 \leq
 4C_H^{\rm site}
 \mathbb E[A_x\mid\cdots].                        \tag{59.13}
\]
Average and use (42.14).
\end{proof}

\subsection{Why the one-site theorem is not yet susceptibility}

The derivative of a one-site conditional expectation with respect to
its linear field is
\[
 D_h\mathbb E_{a,h}f[\delta h]
 =
 2\operatorname{Cov}_{a,h}(f,\delta h\cdot Y).    \tag{59.14}
\]
For \(f=A=\|Y\|^2\), (59.9) gives
\[
 |D_h\mathbb E A[\delta h]|
 \leq
 4C_H^{\rm site}
 \sqrt{\mathbb E_{a,h}A}\,\|\delta h\|.           \tag{59.15}
\]
The raw conditional mean on the right is not uniformly bounded for
arbitrary \(h\): a large linear field moves the quartic well to
distance \(O(\|h\|^{1/3})\).  Thus (59.15) is not a pointwise
Dobrushin coefficient.

After annealing, V41 controls the needed moment, but turning that
average control into the uniform structure factor (57.11) requires an
approximate tensorization or weighted influence theorem.  That is now
the precise intersite problem; the vanishing mass floor itself is no
longer the one-site obstruction.

\begin{assessment}[What V59 closes]
\label{ass:v17v59-status}
V59 proves a spectral gap for the entire linearly tilted quartic
one-site family, uniform in the tilt and in every nonnegative quadratic
coefficient.  It replaces the \(a^{-2}\) Hessian-floor loss of V58 by
a scale-independent conditional-site constant.  The remaining step is
to assemble these site gaps in the presence of unbounded
neighbor-generated fields.
\end{assessment}

\begin{decision}[V60: audit and weighted influence assembly]
\label{dec:v17v59-next}
Perform the mandatory YM/NS audit.  Then combine (59.15) with the V41
exponential quartic moment in a weighted \(L^2\) influence matrix.
Determine an explicit spectral-radius condition under which the
single-site gaps approximately tensorize and imply (57.11).
\end{decision}

\begin{firewall}[Claim boundary after V59]
\label{fire:v17v59}
The theorem is a one-site conditional result in the
nonnegative-mass branch.  Approximate tensorization, the annealed
Higgs-norm susceptibility, the composite remainder envelope,
continuum tightness, gauge cofinality, reflection positivity, the
physical Einstein action, and the Einstein--SM continuum limit remain
open.
\end{firewall}

\section*{Version V59 append-only record}
\addcontentsline{toc}{section}{Version V59 append-only record}
The complete V58 source was retained byte for byte before this
continuation.  It had \(517689\) bytes, \(15865\) lines, and SHA--256
\[
 \texttt{1F5794E12E10D37F1A3C7085203C85C54AC2800EB253B78FC90A8DCF19E87751}.
\]
V59 proved the uniform tilted-quartic one-site spectral gap and
removed the canonical mass-floor loss at one site.


\section{V60: mandatory audit before weighted Higgs assembly}
\label{sec:v17v60}

\subsection{Companion checkpoints}

The newest companion files were inspected read-only:
\[
\begin{array}{lll}
\text{programme}&\text{file}&\text{SHA--256}\\
\text{YM}&
\texttt{Renewal\_Yang\_Mills\_Mass\_Gap\_Research\_Notes\_V88.tex}&
\texttt{A3F9A99B18D5C6AB16BFCF312D7BCB31}\\
&&\texttt{EDB45E128A399E1431B4308EC6EDA8BD}\\
\text{NS}&
\texttt{Renewal\_NS\_Stability\_Research\_Notes\_V26.tex}&
\texttt{82C887F8C2C69E4F28FAD7C3DC8DE5B}\\
&&\texttt{9099CB5A4BF431EBA1FFF3E91935978AD}.
\end{array}                                        \tag{60.1}
\]
YM V88 has \(695024\) bytes and \(18980\) lines; NS V26 has
\(278283\) bytes and \(5319\) lines.

YM V88 certifies three new terminal charts and closes the first two
quarter coordinates of its final direct-plus orientation.  Its third
coordinate, inverted companion, full-torus positivity, cofinal
control, reflection positivity, and mass gap remain open.  NS is
unchanged at its rank-one Oseen-kernel refinement.

\begin{theorem}[V60 audit decision]
\label{thm:v17v60-decision}
Neither companion theorem supplies approximate tensorization for the
Higgs norm field or the residual-current envelope.  No result is
imported.  V61 will assemble the V59 one-site gap under an explicit
weighted influence condition internal to the SM model.
\end{theorem}

\begin{proof}
The YM result is a finite Wilson-symbol chart covering theorem, and
the NS result concerns a fluid kernel.  Neither controls conditional
Higgs kernels, their intersite derivatives, or the scalar structure
factor \(S_a^A\).
\end{proof}

\begin{assessment}[What V60 records]
\label{ass:v17v60-status}
The required audit is complete.  The YM atlas is one coordinate from
its last direct-plus cell, but it still provides no continuum input.
The scale-independent tilted-quartic site gap of V59 remains the
relevant SM datum.
\end{assessment}

\begin{decision}[V61: weighted influence matrix]
\label{dec:v17v60-next}
Use (59.14)--(59.15) to define the neighbor derivative of the Higgs
single-site kernel in an annealed \(L^2\) metric.  Derive an explicit
nonnegative influence matrix and prove approximate tensorization, and
hence (57.11), whenever its spectral radius is below one.  State
separately whether the canonical couplings verify that condition.
\end{decision}

\begin{firewall}[Claim boundary after V60]
\label{fire:v17v60}
V60 is a read-only audit.  It does not complete the YM atlas, prove
Higgs approximate tensorization, establish the residual-current
envelope, or construct the Einstein--Standard-Model continuum.
\end{firewall}

\section*{Version V60 append-only record}
\addcontentsline{toc}{section}{Version V60 append-only record}
The complete V59 source was retained byte for byte before this
continuation.  It had \(524598\) bytes, \(16103\) lines, and SHA--256
\[
 \texttt{080B88CF4879A4B609543E96B6CA6C972E52ACF83AB011388E5A86762829A80D}.
\]
V60 completed the mandatory YM/NS audit and retained the weighted
Higgs-influence route.


\section{V61: weighted \(L^2\) influence assembly for the Higgs field}
\label{sec:v17v61}

\subsection{A neighbor derivative without a raw-field factor}

Fix the conformal and gauge environment \((r,U)\), and let \(P_x\)
denote conditional expectation in the single Higgs variable \(y_x\).
In the nonnegative-mass branch, Theorem~\ref{thm:v17v59-gap} gives the
uniform one-site Poincar\'e constant
\[
 C_0=C_H^{\rm site}
 =
 C_{\rm LC}\sqrt{\frac{q_H}{2\lambda_H}}.         \tag{61.1}
\]

For neighboring sites, the only mixed \(y_x,y_z\) Hessian comes from
the cross term in the geometric hopping square:
\[
 \|\nabla_{y_x}\nabla_{y_z}S_H\|_{\rm op}
 \leq2\chi_{xz}.                                  \tag{61.2}
\]
The conformal diagonal weights contribute no mixed Higgs Hessian, and
unitarity of \(U_{xz}\) makes (61.2) independent of the gauge field.

\begin{lemma}[Single-step sweeping estimate]
\label{lem:v17v61-sweep}
If \(f=f(y_x)\) is smooth and depends only on \(y_x\), then for
\(z\ne x\),
\[
\boxed{
 \|\nabla_{y_z}P_xf\|
 \leq
 2C_0\chi_{xz}
 \left(P_x\|\nabla_{y_x}f\|^2\right)^{1/2}.}      \tag{61.3}
\]
The right side is zero when \(x,z\) are not neighbors.
\end{lemma}

\begin{proof}
Differentiate the normalized conditional integral.  In a unit
direction \(u\) of the \(y_z\) fibre,
\[
 D_{y_z}P_xf[u]
 =
 -\operatorname{Cov}_x
 (f,D_{y_z}S_H[u]).                               \tag{61.4}
\]
Conditional Cauchy--Schwarz and the one-site Poincar\'e inequality
give
\[
\begin{aligned}
 |D_{y_z}P_xf[u]|
 &\leq
 \sqrt{\operatorname{Var}_x(f)}
 \sqrt{\operatorname{Var}_x(D_{y_z}S_H[u])}\\
 &\leq
 C_0
 \left(P_x\|\nabla_{y_x}f\|^2\right)^{1/2}
 \sup_{y_x}
 \|\nabla_{y_x}D_{y_z}S_H[u]\|.                  \tag{61.5}
\end{aligned}
\]
Use (61.2) and take the supremum over unit \(u\).
\end{proof}

Unlike (59.15), this estimate differentiates through the interaction
Hessian before applying Poincar\'e.  It therefore contains no
unbounded factor \(\sqrt{\mathbb E_{a,h}A}\).

\subsection{The influence matrix}

Define the nonnegative matrix
\[
 M_{xz}=
 \begin{cases}
 2C_0\chi_{xz},&x\sim z,\\
 0,&\text{otherwise},
 \end{cases}
 \qquad M_{xx}=0.                                 \tag{61.6}
\]
Let
\[
 r_H=\|M\|_{\ell^2\to\ell^2}.                     \tag{61.7}
\]
For symmetric hoppings on a graph of degree at most \(D\),
\[
 r_H\leq2C_0D\chi_*.                              \tag{61.8}
\]

\begin{theorem}[Quenched approximate tensorization]
\label{thm:v17v61-tensor}
If
\[
\boxed{r_H<1,}                                    \tag{61.9}
\]
then, uniformly in the finite volume and in \((r,U)\),
\[
\boxed{
 \operatorname{Var}(f\mid r,U)
 \leq
 \frac{C_0}{(1-r_H)^2}
 \mathbb E\left[
 \sum_x\|\nabla_{y_x}f\|^2
 \mathrel{\big|}r,U\right].}                     \tag{61.10}
\]
\end{theorem}

\begin{proof}
Order the sites and expose them successively.  The martingale
variance decomposition writes the total variance as the sum of the
squared conditional increments.  Apply the one-site Poincar\'e
inequality to each increment.  When a conditional expectation is
swept past a neighboring site, differentiating its normalized kernel
produces (61.4); Lemma~\ref{lem:v17v61-sweep} bounds the transferred
gradient by the corresponding entry of \(M\).

After all sweeps, the vector of accumulated gradient norms is bounded
by the convergent Neumann series
\[
 (I-M)^{-1}
 =
 \sum_{j=0}^\infty M^j,                            \tag{61.11}
\]
whose \(\ell^2\) norm is at most \((1-r_H)^{-1}\).
Squaring the accumulated gradient bound and summing the martingale
increments gives
\[
 C_0\|(I-M)^{-1}\|_{2\to2}^2
 \mathbb E\sum_x\|\nabla_{y_x}f\|^2,
\]
which is (61.10).  The estimates (61.1)--(61.3) are uniform in the
environment, so the resulting constant is also uniform.
\end{proof}

\subsection{Quenched Higgs-norm structure factor}

Put
\[
 C_{\rm H,glob}
 =
 \frac{C_0}{(1-r_H)^2}.                            \tag{61.12}
\]

\begin{corollary}[Uniform conditional structure factor]
\label{cor:v17v61-structure}
Under (61.9), for every deterministic vertex vector \(t\),
\[
\boxed{
 \mathbb E\operatorname{Var}
 \left(\sum_xt_xA_x\mathrel{\big|}r,U\right)
 \leq
 4C_{\rm H,glob}M_{H,2}\sum_xt_x^2.}              \tag{61.13}
\]
Consequently, the quenched contribution to \(S_a^A(k)\) is uniformly
bounded by
\[
 C_A^{\rm qu}=4C_{\rm H,glob}M_{H,2}.             \tag{61.14}
\]
\end{corollary}

\begin{proof}
Apply (61.10) to \(f=\sum_xt_x\|y_x\|^2\).  Its squared gradient is
\[
 \sum_x4t_x^2\|y_x\|^2.
\]
Average over the environment and use (42.14).  Fourier modes have
unit \(\ell^2\) norm under the convention (54.1), giving (61.14).
\end{proof}

Combining with the exact variance decomposition (58.5),
\[
 S_a^A(k)
 \leq C_A^{\rm qu}+S_a^{\bar A}(k).               \tag{61.15}
\]
Thus V61 closes the entire within-Higgs contribution to (57.11) under
the explicit high-temperature margin (61.9).  The environmental
conditional-mean structure factor remains.

\subsection{Parameter status}

A simple sufficient form of (61.9) is
\[
\boxed{
 2D\chi_*C_{\rm LC}
 \sqrt{\frac{q_H}{2\lambda_H}}<1.}                \tag{61.16}
\]
This condition is scale independent when \(\chi_*,\lambda_H\) have
the four-dimensional powers in (44.11).  The source manuscripts do
not provide numerical bare values or a running trajectory proving
(61.16).  It is therefore a checkable high-temperature branch, not a
certified physical Standard-Model parameter statement.

\begin{assessment}[What V61 closes]
\label{ass:v17v61-status}
V61 converts the tilt-uniform one-site quartic gap into a
volume-uniform quenched Higgs Poincar\'e inequality under the explicit
spectral-radius condition (61.9).  It proves a uniform quenched
Higgs-norm structure factor without a positive lattice-mass floor.
The remaining susceptibility is the coherent response of the
conditional mean \(\bar A(r,U)\) to the conformal and gauge
environment.
\end{assessment}

\begin{decision}[V62: environmental Higgs response]
\label{dec:v17v61-next}
Differentiate \(\bar A_x(r,U)\) with respect to a conformal edge
difference.  Use the global quenched Poincar\'e inequality (61.10) and
the force/curvature quotient of V48 to obtain a weighted derivative
bound.  Test whether conditional gravity mixing then makes
\(S_a^{\bar A}(k)\) uniformly bounded.
\end{decision}

\begin{firewall}[Claim boundary after V61]
\label{fire:v17v61}
The result is conditional on the nonnegative-mass branch and the
high-temperature margin (61.9), which is not verified for physical
running couplings.  The environmental susceptibility, composite
current remainder, continuum measure, gauge cofinality, reflection
positivity, physical Einstein action, and Einstein--SM continuum limit
remain open.
\end{firewall}

\section*{Version V61 append-only record}
\addcontentsline{toc}{section}{Version V61 append-only record}
The complete V60 source was retained byte for byte before this
continuation.  It had \(527648\) bytes, \(16181\) lines, and SHA--256
\[
 \texttt{F28EEA6BE3672A489A391A267DFD3AF93BD82BDA622C6168AAFD3202539705BC}.
\]
V61 assembled the uniform site gaps into a quenched Higgs
susceptibility theorem under an explicit influence margin.


\section{V62: derivative of the environmental Higgs response}
\label{sec:v17v62}

\subsection{Conformal-edge derivative}

Fix an environment \(\eta=(r,U)\), and let \(\pi_\eta\) be the full
conditional Higgs law.  For a deterministic vertex vector \(t\), put
\[
 A(t)=\sum_xt_x\|y_x\|^2,\qquad
 \bar A_\eta(t)=\mathbb E_{\pi_\eta}A(t).          \tag{62.1}
\]
Changing one conformal edge difference \(d_e\) differentiates the
Higgs action by
\[
 \partial_{d_e}S_H=q_e.                           \tag{62.2}
\]

\begin{theorem}[Exact environmental response derivative]
\label{thm:v17v62-derivative}
For every edge \(e\),
\[
\boxed{
 \partial_{d_e}\bar A_\eta(t)
 =
 -\operatorname{Cov}_{\pi_\eta}(A(t),q_e).}       \tag{62.3}
\]
Under the V61 influence margin, this derivative obeys
\[
\boxed{
 |\partial_{d_e}\bar A_\eta(t)|^2
 \leq
 16(C_{\rm H,glob})^2\chi_e^2
 \left(\sum_xt_x^2\bar A_x(\eta)\right)
 \left(
 e^{2d_e}\bar A_{e_-}(\eta)
 +e^{-2d_e}\bar A_{e_+}(\eta)
 \right).}                                       \tag{62.4}
\]
\end{theorem}

\begin{proof}
Differentiate the normalized conditional Higgs integral.  The
derivative of the logarithmic density is
\(-q_e+\mathbb E_{\pi_\eta}q_e\), which gives (62.3).

The global conditional Poincar\'e inequality (61.10) and
Cauchy--Schwarz imply
\[
 |\operatorname{Cov}_{\pi_\eta}(f,g)|^2
 \leq
 (C_{\rm H,glob})^2
 \mathbb E_{\pi_\eta}\|\nabla_yf\|^2
 \mathbb E_{\pi_\eta}\|\nabla_yg\|^2.             \tag{62.5}
\]
Here
\[
\begin{aligned}
 \mathbb E_{\pi_\eta}\|\nabla_yA(t)\|^2
 &=4\sum_xt_x^2\bar A_x(\eta),\\
 \mathbb E_{\pi_\eta}\|\nabla_yq_e\|^2
 &=4\chi_e^2
 \{e^{2d_e}\bar A_{e_-}(\eta)
 +e^{-2d_e}\bar A_{e_+}(\eta)\}.                 \tag{62.6}
\end{aligned}
\]
Insert (62.6) in (62.5).
\end{proof}

\subsection{The exact new mixed moment}

Define
\[
 \mathcal D_{2,a}
 =
 \sup_e
 \mathbb E_{\eta}
 \left[
 e^{2d_e}\bar A_{e_-}(\eta)
 +e^{-2d_e}\bar A_{e_+}(\eta)
 \right].                                        \tag{62.7}
\]
At every fixed finite cutoff this number is finite: quartic Higgs
coercivity controls powers of \(y\), while the cosh mismatch controls
every fixed exponential of \(d_e\).  The earlier estimates do not
make (62.7) uniform as \(a\downarrow0\).

For a unit Fourier vector \(t\), translation invariance and (42.14)
give
\[
 \mathbb E_\eta\sum_xt_x^2\bar A_x(\eta)
 \leq M_{H,2}.                                    \tag{62.8}
\]
Equation (62.4) therefore identifies
\[
 16(C_{\rm H,glob})^2\chi_*^2M_{H,2}\mathcal D_{2,a}               \tag{62.9}
\]
as an annealed squared conformal-edge derivative budget.  This is a
local mixed moment, rather than the extensive free energy that
appeared in V44.

\subsection{Gauge derivative}

Let \(\xi_e\) be a unit Lie-algebra tangent at \(U_e\).  Only the cross
term depends on this gauge link, and
\[
 \partial_{\xi_e}S_H
 =
 -2\chi_e\operatorname{Re}
 \langle y_{e_+},(\xi_eU_e)y_{e_-}\rangle
 =:j_e^\xi.                                      \tag{62.10}
\]
Consequently,
\[
 \boxed{
 \partial_{\xi_e}\bar A_\eta(t)
 =
 -\operatorname{Cov}_{\pi_\eta}(A(t),j_e^\xi).}   \tag{62.11}
\]
The same Poincar\'e argument bounds it by conditional Higgs moments,
but converting those local derivatives into a structure factor for
\(\bar A\) requires a gauge-environment spectral gap or covariance
decay.  Compactness of the gauge group alone does not provide that
estimate.  The YM audit through V60 supplies finite symbol positivity,
not this Gibbs-environment gap.

\subsection{Two-scale variance ledger}

For a Fourier mode \(t_k\), the exact decomposition remains
\[
 S_a^A(k)
 =
 S_{a,\rm qu}^A(k)+S_a^{\bar A}(k),               \tag{62.12}
\]
where V61 proves
\[
 S_{a,\rm qu}^A(k)\leq C_A^{\rm qu}.              \tag{62.13}
\]
To control the second term from (62.4) and (62.11), it is sufficient
to acquire:
\begin{enumerate}
\item a scale-compatible conformal-environment Poincar\'e or
      negative-Sobolev covariance inequality acting on the edge
      derivatives in (62.4), together with a uniform form of
      \(\mathcal D_{2,a}\);
\item a gauge-environment covariance inequality for (62.11), uniform
      along the selected gauge scaling trajectory.
\end{enumerate}
Neither item follows from one-site Higgs confinement.

\begin{assessment}[What V62 closes]
\label{ass:v17v62-status}
V62 differentiates the environmental Higgs norm response exactly and
uses the assembled Higgs gap to bound its conformal derivative by the
single local mixed moment (62.7).  It separately identifies the gauge
current derivative.  The annealed Higgs susceptibility is now reduced
to two explicit environment covariance estimates rather than an
unresolved conditional mean.
\end{assessment}

\begin{decision}[V63: effective environment Hessian]
\label{dec:v17v62-next}
Differentiate the Higgs conditional free energy twice in the
conformal edge variables.  Use V61 to bound the negative current
covariance and derive an explicit lower Hessian for the conformal
environment.  Determine whether the mismatch curvature absorbs that
debit under a scale-independent parameter condition.
\end{decision}

\begin{firewall}[Claim boundary after V62]
\label{fire:v17v62}
The mixed moment \(\mathcal D_{2,a}\), the conformal environment
inequality, and the gauge environment gap are not yet proved uniformly
in lattice spacing.  No full susceptibility, remainder-current
envelope, continuum measure, gauge cofinality, reflection positivity,
physical Einstein action, or Einstein--SM continuum limit follows.
\end{firewall}

\section*{Version V62 append-only record}
\addcontentsline{toc}{section}{Version V62 append-only record}
The complete V61 source was retained byte for byte before this
continuation.  It had \(534837\) bytes, \(16414\) lines, and SHA--256
\[
 \texttt{36FFA29F63AE57DA56BAB182F11D84E73488B12D61C3D9E6B6D479279C1F2895}.
\]
V62 derived the environmental response derivative and its conformal
and gauge covariance ledger.


\section{V63: the effective conformal-environment Hessian}
\label{sec:v17v63}

\subsection{Higgs free energy at fixed gauge field}

For fixed \(U\), integrate the Higgs variables at a fixed conformal
difference field \(d\):
\[
 \mathcal F_H(d,U)=-\log Z_H(d,U),\qquad
 Z_H(d,U)=\int e^{-S_H(y\mid d,U)}\,dy.            \tag{63.1}
\]
For an edge vector \(u\), define
\[
 Q_u=\sum_eu_eq_e,\qquad
 T_{u^2}=\sum_eu_e^2T_e.                          \tag{63.2}
\]

\begin{theorem}[Exact Higgs free-energy Hessian]
\label{thm:v17v63-exact}
For every edge direction \(u\),
\[
\boxed{
 D^2\mathcal F_H(d,U)[u,u]
 =
 \mathbb E_{\pi_\eta}T_{u^2}
 -\operatorname{Var}_{\pi_\eta}(Q_u).}            \tag{63.3}
\]
\end{theorem}

\begin{proof}
The first directional derivative of the Higgs action is \(Q_u\).
Since \(\partial_{d_e}q_e=T_e\), its second directional derivative is
\(T_{u^2}\).  Differentiate the normalized partition function twice;
the derivative of the expectation contributes
\(-\operatorname{Var}(Q_u)\), proving (63.3).
\end{proof}

\subsection{A local covariance debit}

At a vertex \(x\), define its diagonal conformal hopping coefficient
\[
 H_x(d)=
 \sum_{\substack{e\ {\rm outgoing}\\e_-=x}}
 \chi_ee^{d_e}
 +
 \sum_{\substack{e\ {\rm incoming}\\e_+=x}}
 \chi_ee^{-d_e}.                                  \tag{63.4}
\]
For an incident edge, write
\[
 \sigma_{x,e}=
 \begin{cases}
 +1,&x=e_-,\\
 -1,&x=e_+.
 \end{cases}                                      \tag{63.5}
\]

\begin{lemma}[Gradient of the conformal current]
\label{lem:v17v63-gradient}
Under the V61 influence margin,
\[
\operatorname{Var}_{\pi_\eta}(Q_u)
 \leq
 4C_{\rm H,glob}
 \sum_eu_e^2\,
 \mathbb E_{\pi_\eta}
 \left[
 D_{e_-,e}H_{e_-}
 +D_{e_+,e}H_{e_+}
 \right].                                        \tag{63.6}
\]
\end{lemma}

\begin{proof}
At site \(x\),
\[
 \nabla_{y_x}Q_u
 =
 2y_x\sum_{e\ni x}
 \sigma_{x,e}u_e\chi_e e^{\sigma_{x,e}d_e}.       \tag{63.7}
\]
Weighted Cauchy--Schwarz gives
\[
\begin{aligned}
 \|\nabla_{y_x}Q_u\|^2
 &\leq
 4\|y_x\|^2H_x(d)
 \sum_{e\ni x}u_e^2\chi_e e^{\sigma_{x,e}d_e}.
                                                               \tag{63.8}
\end{aligned}
\]
Sum over \(x\), apply the global conditional Poincar\'e inequality
(61.10), and recognize
\(\chi_e e^{\sigma_{x,e}d_e}\|y_x\|^2=D_{x,e}\).
\end{proof}

Define the explicit edge debit
\[
 \Gamma_e(d,U)
 =
 4C_{\rm H,glob}
 \mathbb E_{\pi_\eta}
 [D_{e_-,e}H_{e_-}+D_{e_+,e}H_{e_+}].            \tag{63.9}
\]

\begin{corollary}[Local lower Hessian]
\label{cor:v17v63-lower}
For every \(u\),
\[
\boxed{
 D^2\mathcal F_H(d,U)[u,u]
 \geq
 \sum_eu_e^2
 \{\mathbb E_{\pi_\eta}T_e-\Gamma_e(d,U)\}.}       \tag{63.10}
\]
\end{corollary}

\begin{proof}
Combine (63.3) and (63.6).
\end{proof}

\subsection{Effective conformal curvature}

After Higgs integration at fixed \(U\), the conformal environment
potential is
\[
 \Phi_{\rm env}(r\mid U)
 =
 \kappa_g\|r\|^2+
 J\sum_e\cosh(2d_e)+\mathcal F_H(Br,U).           \tag{63.11}
\]
Set
\[
 c_e^{\rm env}(d,U)
 =
 4J\cosh(2d_e)
 +\mathbb E_{\pi_\eta}T_e-\Gamma_e(d,U).          \tag{63.12}
\]

\begin{theorem}[Checkable environment-convexity criterion]
\label{thm:v17v63-environment}
If
\[
\boxed{
 c_e^{\rm env}(d,U)\geq0
 \quad\text{for every }e,d,U,}                    \tag{63.13}
\]
then
\[
 \nabla_r^2\Phi_{\rm env}(r\mid U)
 \succeq2\kappa_gI.                               \tag{63.14}
\]
More generally, if \(c_e^{\rm env}\geq c_*>0\), then
\[
 \nabla_r^2\Phi_{\rm env}(r\mid U)
 \succeq2\kappa_gI+c_*L.                          \tag{63.15}
\]
\end{theorem}

\begin{proof}
The mismatch Hessian is
\(B^{\mathsf T}\operatorname{diag}(4J\cosh(2d_e))B\).
Apply (63.10) with \(u=Bh\) to a conformal variation \(h\).  The
resulting quadratic form is bounded below by
\[
 2\kappa_g\|h\|^2+
 \sum_ec_e^{\rm env}(Bh)_e^2.
\]
Conditions (63.13) and \(c_e^{\rm env}\geq c_*\) give
(63.14)--(63.15).
\end{proof}

Condition (63.13) is scale local and contains the correct competition:
the negative Higgs covariance is charged against both the nonlinear
mismatch curvature and the positive mean diagonal hopping curvature.
Replacing \(H_x(d)\) by a global supremum would fail because the
conformal differences are unbounded.

\subsection{A sufficient bounded-difference branch}

On the restricted event
\[
 |d_e|\leq R\quad\text{for all edges},             \tag{63.16}
\]
one has \(H_x(d)\leq D\chi_*e^R\), and therefore
\[
 \Gamma_e
 \leq
 4C_{\rm H,glob}D\chi_*e^R
 \mathbb E_{\pi_\eta}T_e.                         \tag{63.17}
\]
Hence the scale-independent margin
\[
\boxed{
 4C_{\rm H,glob}D\chi_*e^R\leq1}                 \tag{63.18}
\]
implies (63.13) on that event.

The full conformal law is not supported in a fixed difference box, so
(63.18) is a localized branch, not a global theorem.  The mismatch
term in (63.12) grows as \(e^{2|d_e|}\), faster than the single
coefficient \(H_x\), suggesting that the complement of (63.16) should
be treated by direct quartic/cosh absorption rather than a supremum.

\begin{assessment}[What V63 closes]
\label{ass:v17v63-status}
V63 computes the effective Higgs free-energy Hessian and bounds its
negative covariance by the explicit local debit \(\Gamma_e\).  It
gives a pointwise criterion for convexity of the conformal environment
and proves that criterion on every bounded-difference sector under
(63.18).  Extending it through the unbounded conformal tails is the
remaining conformal-environment problem.
\end{assessment}

\begin{decision}[V64: large-difference absorption]
\label{dec:v17v63-next}
Bound
\(\mathbb E_{\pi_\eta}[D_{x,e}H_x]\) deterministically in terms of the
quartic coefficient and the neighbor linear fields.  Compare its
large-\(|d|\) growth with \(4J\cosh(2d_e)\).  Either prove a global
version of (63.13) under an explicit parameter margin or isolate the
precise multi-edge star term that prevents absorption.
\end{decision}

\begin{firewall}[Claim boundary after V63]
\label{fire:v17v63}
Global environment convexity has not been proved; (63.18) applies only
on a bounded-difference sector.  The gauge-environment gap,
environmental susceptibility, composite-current envelope, continuum
measure, physical Einstein action, and Einstein--SM continuum limit
remain open.
\end{firewall}

\section*{Version V63 append-only record}
\addcontentsline{toc}{section}{Version V63 append-only record}
The complete V62 source was retained byte for byte before this
continuation.  It had \(540781\) bytes, \(16602\) lines, and SHA--256
\[
 \texttt{EAF95E0F67D19F35C8D03BE6E592F696E918A503BCBF6F8949A05AB18F204472}.
\]
V63 derived the effective environment Hessian, its local covariance
debit, and a bounded-difference convexity branch.


\section{V64: large-difference reduction to a neighbor-field star term}
\label{sec:v17v64}

\subsection{A one-site radial estimate with arbitrary linear field}

Fix the environment \((d,U)\) and all Higgs variables except \(y_x\).
The conditional one-site potential is
\[
 V_x(y)
 =
 \lambda_H\|y\|^4+c_x(d)\|y\|^2-2h_x\cdot y,
                                                               \tag{64.1}
\]
where
\[
\begin{aligned}
 c_x(d)&=m_H^2+H_x(d),\\
 h_x&=\sum_{z\sim x}\chi_{xz}U_{xz}^{(x)}y_z.     \tag{64.2}
\end{aligned}
\]
Here \(U_{xz}^{(x)}\) is the unitary transport into the fibre at \(x\).
In the nonnegative-mass branch,
\[
 c_x(d)\geq H_x(d)\geq0.                          \tag{64.3}
\]

\begin{lemma}[Diagonal hopping against the star coefficient]
\label{lem:v17v64-star}
Let
\[
 g_{x,e}(d)=\chi_e e^{\sigma_{x,e}d_e},
 \qquad e\ni x.                                   \tag{64.4}
\]
For \(H_x(d)>0\),
\[
\boxed{
 \mathbb E_{\pi_\eta}
 [D_{x,e}H_x]
 \leq
 q_Hg_{x,e}(d)+
 \mathbb E_{\pi_\eta}\|h_x\|^2.}                 \tag{64.5}
\]
If \(H_x=0\), the left side is zero.
\end{lemma}

\begin{proof}
The radial integration-by-parts identity for (64.1), conditional on
all other Higgs variables, is
\[
 4\lambda_H\mathbb E_x\|y_x\|^4
 +2c_x\mathbb E_x\|y_x\|^2
 -2h_x\cdot\mathbb E_xy_x
 =q_H.                                            \tag{64.6}
\]
Young's inequality gives
\[
 2h_x\cdot y_x
 \leq c_x\|y_x\|^2+\frac{\|h_x\|^2}{c_x}.
\]
Discard the nonnegative quartic term in (64.6) to obtain
\[
 c_x\mathbb E_x\|y_x\|^2
 \leq q_H+\frac{\|h_x\|^2}{c_x}.                 \tag{64.7}
\]
Since \(D_{x,e}=g_{x,e}\|y_x\|^2\) and
\(g_{x,e}\leq H_x\leq c_x\),
\[
\begin{aligned}
 g_{x,e}H_x\mathbb E_x\|y_x\|^2
 &\leq
 q_H\frac{g_{x,e}H_x}{c_x}
 +\frac{g_{x,e}H_x}{c_x^2}\|h_x\|^2\\
 &\leq q_Hg_{x,e}+\|h_x\|^2.                     \tag{64.8}
\end{aligned}
\]
Average over the remaining Higgs variables.
\end{proof}

\subsection{The effective-Hessian debit after reduction}

Insert (64.5) at both endpoints in (63.9).

\begin{corollary}[Edge debit with no quadratic conformal exponential]
\label{cor:v17v64-debit}
For \(e=(v,w)\),
\[
\boxed{
 \Gamma_e(d,U)
 \leq
 4C_{\rm H,glob}\left[
 q_H\chi_e(e^{d_e}+e^{-d_e})
 +\mathbb E_{\pi_\eta}
 (\|h_v\|^2+\|h_w\|^2)
 \right].}                                       \tag{64.9}
\]
\end{corollary}

\begin{proof}
Apply Lemma~\ref{lem:v17v64-star} to
\((v,e)\) and \((w,e)\) and sum.
\end{proof}

The first term in (64.9) grows like \(e^{|d_e|}\), whereas the pure
mismatch curvature in (63.12) grows like \(e^{2|d_e|}\).  The
potentially quadratic conformal growth from the crude product
\(D_{x,e}H_x\) has therefore been removed exactly.

\begin{theorem}[A global sufficient environment-curvature margin]
\label{thm:v17v64-margin}
Suppose
\[
 2C_{\rm H,glob}q_H\chi_*\leq J                  \tag{64.10}
\]
and, for every edge and environment,
\[
 \mathbb E_{\pi_\eta}T_e
 \geq
 4C_{\rm H,glob}
 \mathbb E_{\pi_\eta}(\|h_{e_-}\|^2+\|h_{e_+}\|^2).              \tag{64.11}
\]
Then the environment coefficient (63.12) is nonnegative, and
\[
 \nabla_r^2\Phi_{\rm env}(r\mid U)
 \succeq2\kappa_gI.                               \tag{64.12}
\]
\end{theorem}

\begin{proof}
For every real \(d\),
\[
 \cosh(2d)\geq\cosh d.
\]
Condition (64.10) gives
\[
 4C_{\rm H,glob}q_H\chi_e(e^d+e^{-d})
 =
 8C_{\rm H,glob}q_H\chi_e\cosh d
 \leq4J\cosh(2d).                                 \tag{64.13}
\]
Condition (64.11) lets the positive
\(\mathbb E_{\pi_\eta}T_e\) absorb the remaining term in (64.9).
Thus \(c_e^{\rm env}\geq0\), and
Theorem~\ref{thm:v17v63-environment} proves (64.12).
\end{proof}

\subsection{The canonical-scaling obstruction to this simple margin}

Under (44.10)--(44.11),
\[
 J(a)=\bar Ja^2,\qquad\chi_*(a)=\bar\chi_*.
\]
If \(C_{\rm H,glob},q_H,\bar\chi_*\) stay positive, (64.10) becomes
\[
 2C_{\rm H,glob}q_H\bar\chi_*
 \leq\bar Ja^2,                                   \tag{64.14}
\]
which fails for all sufficiently small \(a\).

\begin{countertheorem}[Mismatch-only absorption is not scale safe]
\label{ctr:v17v64-scaling}
The sufficient proof in Theorem~\ref{thm:v17v64-margin}, which charges
the entire radial contact term in (64.9) to the pure mismatch
curvature, cannot hold along the canonical four-dimensional scaling
with fixed nonzero hopping.
\end{countertheorem}

\begin{proof}
Equation (64.14) is necessary for that proof and is incompatible with
\(a\downarrow0\).
\end{proof}

This does not prove loss of environment convexity.  It proves that the
positive mean hopping curvature \(\mathbb ET_e\) must participate in
the cancellation of the radial contact term; allocating it only to
the neighbor-field contribution as in (64.11) is too expensive.

\subsection{The precise remaining star obstruction}

Unitarity and Cauchy--Schwarz give
\[
 \|h_x\|^2
 \leq
 D\chi_*^2\sum_{z\sim x}\|y_z\|^2.                \tag{64.15}
\]
After full annealing, V41 controls its expectation.  Condition
(63.13), however, is pointwise in the environment before conformal
mixing is obtained.  The conditional neighbor norms on the right of
(64.15) are not uniformly bounded in \((d,U)\).  This multi-edge star
term is the exact obstruction left by V64.

\begin{assessment}[What V64 closes]
\label{ass:v17v64-status}
V64 uses a local radial identity to reduce the large-difference
Hessian debit from a product of two exponential edge coefficients to
a single \(e^{|d|}\) term plus a neighbor-generated linear-field
square.  It proves a global convexity branch and then shows that its
mismatch-only contact allocation is incompatible with canonical
scaling.  The remaining task is a joint allocation of
\(\mathbb ET_e\) between the contact and star terms.
\end{assessment}

\begin{decision}[V65: audit and optimal contact allocation]
\label{dec:v17v64-next}
Perform the mandatory YM/NS audit.  Then introduce an allocation
parameter in the radial estimate (64.7), retaining a controlled
fraction of the quartic and mean hopping curvature.  Optimize that
parameter edge by edge and determine whether
\(4J\cosh(2d)+\mathbb ET_e\) absorbs the full debit at canonical
scaling, or whether the neighbor-field star gives a genuine
counterexample.
\end{decision}

\begin{firewall}[Claim boundary after V64]
\label{fire:v17v64}
The scale-safe global environment Hessian has not been proved.
The sufficient margin (64.10)--(64.11) is deliberately stronger than
necessary and fails canonical scaling.  Environmental susceptibility,
the remainder-current envelope, continuum measure, gauge cofinality,
reflection positivity, physical Einstein action, and Einstein--SM
continuum limit remain open.
\end{firewall}

\section*{Version V64 append-only record}
\addcontentsline{toc}{section}{Version V64 append-only record}
The complete V63 source was retained byte for byte before this
continuation.  It had \(547489\) bytes, \(16846\) lines, and SHA--256
\[
 \texttt{809447B48A4DB1499FD955D6FEEE0B16384BE3BA035F11FC2ABAF76DCDE30AC7}.
\]
V64 reduced the large-tail Hessian debit and isolated the
scale-critical neighbor-field star term.


\section{V65: compulsory companion audit and allocation decision}
\label{sec:v17v65}

This version performs the five-version audit required by the programme.
The comparison is logical rather than terminological: a result is imported
only when its hypotheses and conclusion apply to the conformal--Higgs
environment problem isolated in V64.

\subsection{Active companion records}

The newest active Yang--Mills note is
\begin{center}\small
\path{Renewal_Yang_Mills_Mass_Gap_Research_Notes_V89.tex}.
\end{center}
It has \(699969\) bytes, \(19116\) lines, and SHA--256
\[
 \texttt{4BE1C250231A8D665973C9BEA7D5263742FFF9A1E5F96E81F20E2DB1AC63382A}.
                                                               \tag{65.1}
\]
Its new result is the exact certification of nine power-two charts for
the direct-plus Wilson pencil in the \((+-+)\) orientation, completing
the third-coordinate slab through depth \(2t\).  The terminal
\(n_4=3\) layer, its cap, the last inversion pair, and full transverse
torus positivity remain open there.

The newest active Navier--Stokes note remains
\begin{center}\small
\path{Renewal_NS_Stability_Research_Notes_V26.tex}.
\end{center}
It has \(278283\) bytes, \(5319\) lines, and SHA--256
\[
 \texttt{82C887F8C2C69E4F28FAD7C3DC8DE5B9099CB5A4BF431EBA1FFF3E91935978AD}.
                                                               \tag{65.2}
\]
Its new block computes pointwise rank-one norms for the first two
derivatives of the Oseen kernel.  The first-derivative refinement gives
no gain, while the second-derivative rank-one norm is strictly smaller
on an intermediate radial interval.  The certified integrated constants
and the resulting floor are deferred to its V27.

\begin{theorem}[Companion-audit non-import theorem]
\label{thm:v17v65-nonimport}
Neither companion result supplies a missing premise of the V64
environment-curvature problem.  In particular, neither proves a bound
of the form
\[
 \mathbb E_{\pi_\eta}\|h_x\|^2
 \leq C\,\mathbb E_{\pi_\eta}T_e                 \tag{65.3}
\]
uniformly in the fixed conformal and gauge environment, nor a spectral
gap for that environment.
\end{theorem}

\begin{proof}
The V89 Yang--Mills statement is finite-dimensional positivity of a
specified Wilson symbol on explicitly covered momentum boxes.  Its
hypotheses contain neither the conditional Higgs measure \(\pi_\eta\)
nor the random linear field \(h_x\), and its conclusion is not a
position-space covariance or environment Poincar\'e inequality.  The V26
Navier--Stokes statement is an algebraic norm computation for derivatives
of the Oseen kernel; it contains none of the lattice variables in (65.3).
Thus there is no implication to import.  The companion results are
consistent with the SM route, but logical consistency is not a missing
hypothesis.
\end{proof}

\begin{assessment}[Direction after the audit]
\label{ass:v17v65-direction}
The YM atlas has advanced toward its last sign cell, but it has not yet
produced a full-torus gauge covariance estimate that could enter V62.
The NS rank-one calculation illustrates a useful methodological point:
restricting a norm to the tensors actually generated by the nonlinear
source can improve a constant, but only after the exact contraction is
computed.  For the SM problem, the analogous exact contraction is the
one-parameter radial allocation in (64.6), which must be optimized before
abandoning pointwise environment convexity.
\end{assessment}

\begin{decision}[V66: optimize the radial allocation]
\label{dec:v17v65-next}
Insert \(2h\cdot y\leq\delta c\|y\|^2+
\|h\|^2/(\delta c)\), \(0<\delta<2\), into the exact radial identity.
Compute the resulting contact and star coefficients, optimize their
tradeoff, and test whether any choice can make mismatch-only absorption
compatible with \(J(a)=\bar J a^2\).  If it cannot, use the exact
inequality to identify which part of \(\mathbb ET_e\) must be retained
in the next upstream argument.
\end{decision}

\begin{firewall}[Claim boundary after V65]
\label{fire:v17v65}
This audit imports no theorem into the SM chain.  The scale-safe global
environment Hessian, environmental susceptibility, remainder-current
envelope, continuum measure, gauge cofinality, physical Einstein action,
and Einstein--SM continuum limit remain open.
\end{firewall}

\section*{Version V65 append-only record}
\addcontentsline{toc}{section}{Version V65 append-only record}
The complete V64 source was retained byte for byte before this
continuation.  It had \(554565\) bytes, \(17079\) lines, and SHA--256
\[
 \texttt{A67807CE8325E575F6905F793E6CFCEBE545A20E31EED17B69E53944FC594B55}.
\]
V65 audited active YM V89 and NS V26 and fixed the V66 acquisition order.


\section{V66: optimal radial allocation and its scale barrier}
\label{sec:v17v66}

V64 used the symmetric Young split in the one-site radial identity.
This version computes the entire one-parameter family and minimizes it
exactly.  The result quantifies how much of the positive mean hopping
curvature must enter any argument based on that family.

\subsection{The parameterized one-site estimate}

Retain the notation of (64.1)--(64.4).  Conditional on all Higgs
variables other than \(y_x\), put
\[
 A_{x,e}:=q_Hg_{x,e}(d),
 \qquad B_x:=\|h_x\|^2.                            \tag{66.1}
\]

\begin{lemma}[Parameterized radial allocation]
\label{lem:v17v66-delta}
For every \(0<\delta<2\),
\[
 \mathbb E_x[D_{x,e}H_x]
 \leq
 \frac{A_{x,e}}{2-\delta}
 +\frac{B_x}{\delta(2-\delta)}.                  \tag{66.2}
\]
Here \(\mathbb E_x\) denotes the conditional expectation in \(y_x\).
The statement is trivial when \(H_x=0\).
\end{lemma}

\begin{proof}
For \(H_x>0\), (64.3) gives \(c_x>0\).  Young's inequality with a
free parameter gives
\[
 2h_x\cdot y_x
 \leq\delta c_x\|y_x\|^2+
       \frac{\|h_x\|^2}{\delta c_x}.              \tag{66.3}
\]
Insert (66.3) in the exact identity (64.6) and discard only the
nonnegative quartic term.  Then
\[
 (2-\delta)c_x\mathbb E_x\|y_x\|^2
 \leq q_H+\frac{B_x}{\delta c_x}.                 \tag{66.4}
\]
Multiplication by \(g_{x,e}H_x/c_x\), followed by
\(g_{x,e}\leq H_x\leq c_x\), yields
\[
\begin{aligned}
 g_{x,e}H_x\mathbb E_x\|y_x\|^2
 &\leq
 \frac{q_Hg_{x,e}H_x}{(2-\delta)c_x}
 +\frac{g_{x,e}H_xB_x}{\delta(2-\delta)c_x^2}\\
 &\leq
 \frac{A_{x,e}}{2-\delta}
 +\frac{B_x}{\delta(2-\delta)}.
\end{aligned}                                      \tag{66.5}
\]
The left side is \(\mathbb E_x[D_{x,e}H_x]\).
\end{proof}

\subsection{Exact minimization}

For \(A,B\geq0\), define
\[
 \mathcal P(A,B)
 :=\inf_{0<\delta<2}
 \left\{\frac{A}{2-\delta}
 +\frac{B}{\delta(2-\delta)}\right\}.             \tag{66.6}
\]

\begin{theorem}[Closed form of the optimal radial envelope]
\label{thm:v17v66-envelope}
For all \(A,B\geq0\),
\[
\boxed{
 \mathcal P(A,B)
 =\frac{A+B+\sqrt{B^2+2AB}}{2}.}                  \tag{66.7}
\]
If \(A,B>0\), the unique minimizer is
\[
 \delta_*(A,B)
 =\frac{\sqrt{B^2+2AB}-B}{A}
 =\frac{2B}{B+\sqrt{B^2+2AB}}\in(0,1).           \tag{66.8}
\]
The boundary cases are obtained by continuity: \(\delta_*\downarrow0\)
when \(B=0\), and \(\delta_*=1\) when \(A=0<B\).
Moreover,
\[
 \frac A2+B\leq\mathcal P(A,B)\leq A+B.          \tag{66.9}
\]
\end{theorem}

\begin{proof}
For positive \(A,B\), differentiation of
\[
 f(\delta)=\frac{A\delta+B}{\delta(2-\delta)}
\]
shows that its only critical point in \((0,2)\) satisfies
\[
 A\delta^2+2B\delta-2B=0.                         \tag{66.10}
\]
This gives (66.8).  Since \(f\) diverges at both endpoints, the
critical point is the unique minimum.  Substitution gives (66.7).
The same formula follows at \(A=0\) or \(B=0\) by taking the indicated
limits.  Finally,
\[
 B\leq\sqrt{B^2+2AB}\leq A+B,                    \tag{66.11}
\]
which proves (66.9).
\end{proof}

Because \(\delta_*\) depends only on the frozen neighboring variables,
it may be selected before integrating in \(y_x\).  Averaging the
result of Lemma~\ref{lem:v17v66-delta} therefore gives the optimized
conditional estimate
\[
\boxed{
 \mathbb E_{\pi_\eta}[D_{x,e}H_x]
 \leq
 \mathbb E_{\pi_\eta}
 \mathcal P(q_Hg_{x,e},\|h_x\|^2).}               \tag{66.12}
\]

\subsection{Optimized edge debit and the unavoidable contact share}

\begin{corollary}[Best debit supplied by the radial Young family]
\label{cor:v17v66-debit}
For \(e=(v,w)\), the V63 debit obeys
\[
\boxed{
 \Gamma_e
 \leq4C_{\rm H,glob}\,
 \mathbb E_{\pi_\eta}\left[
 \mathcal P(q_H\chi_e e^{d_e},\|h_v\|^2)
 +\mathcal P(q_H\chi_e e^{-d_e},\|h_w\|^2)
 \right].}                                       \tag{66.13}
\]
Consequently, the pointwise condition
\[
\begin{aligned}
 4J\cosh(2d_e)+\mathbb E_{\pi_\eta}T_e
 \geq{}&4C_{\rm H,glob}\,
 \mathbb E_{\pi_\eta}\left[
 \mathcal P(q_H\chi_e e^{d_e},\|h_v\|^2)\right.\\
 &\left.\hspace{33mm}
 +\mathcal P(q_H\chi_e e^{-d_e},\|h_w\|^2)
 \right]
                                                        \tag{66.14}
\end{aligned}
\]
implies \(c_e^{\rm env}\geq0\).
\end{corollary}

\begin{proof}
Apply (66.12) at the two endpoints in (63.9).  Inequality (66.14)
then makes the coefficient (63.12) nonnegative.
\end{proof}

The two coefficients before optimization satisfy
\[
 \frac1{2-\delta}\geq\frac12,
 \qquad
 \frac1{\delta(2-\delta)}\geq1.                  \tag{66.15}
\]
Thus this Young family can reduce the contact coefficient by at most a
factor two, and it cannot reduce the coefficient of \(\|h_x\|^2\)
below the V64 value.  If the entire \(q_Hg\) contribution is charged
to mismatch curvature, its best possible uniform comparison is
\[
 4C_{\rm H,glob}q_H\chi_e\cosh d_e
 \leq4J\cosh(2d_e).                               \tag{66.16}
\]
Since \(\cosh(2d)/\cosh d\) has minimum one at \(d=0\), (66.16) for
all \(d\) is equivalent to
\[
 C_{\rm H,glob}q_H\chi_e\leq J.                  \tag{66.17}
\]

\begin{countertheorem}[No Young parameter repairs canonical scaling]
\label{ctr:v17v66-scale}
Assume \(J(a)=\bar Ja^2\), fixed \(\chi_e>0\), and a positive
scale-independent \(C_{\rm H,glob}\).  No choice of
\(0<\delta<2\), including an environment-dependent choice, makes the
mismatch curvature alone absorb the radial contact term uniformly as
\(a\downarrow0\).
\end{countertheorem}

\begin{proof}
The first inequality in (66.15) shows that every member of the family
requires at least (66.17) for mismatch-only allocation.  Its left side
is fixed and positive, whereas its right side tends to zero.
\end{proof}

This is a limitation of the radial-envelope proof, not a counterexample
to convexity of the effective environment.  It proves that a
scale-safe argument must use a positive fraction of
\(\mathbb E_{\pi_\eta}T_e\) to pay even the contact term.  The remaining
fraction must control the neighbor-field part without a pointwise bound
on the surrounding conformal differences.

\begin{assessment}[What V66 closes]
\label{ass:v17v66-status}
V66 exactly optimizes every estimate obtainable from the quadratic
Young split in the radial Ward identity.  The optimized envelope is
(66.7), and the factor-two barrier (66.15) proves that parameter tuning
cannot restore canonical scaling.  The next step must preserve the
weighted edge structure of \(h_x\), which was discarded in (64.15).
\end{assessment}

\begin{decision}[V67: weighted transport of the star field]
\label{dec:v17v66-next}
Use weighted Cauchy--Schwarz before replacing \(H_x/c_x\) by one:
express \(\|h_x\|^2\) through the diagonal hopping terms on the
opposite ends of edges incident to \(x\).  Determine whether the
resulting oriented two-edge kernel is summable into
\(\sum_e\mathbb ET_e\) at the quadratic-form level.  If a local
row-sum obstruction remains, abandon edgewise positivity and test the
full incidence-form Hessian together with the \(2\kappa_gI\) floor.
\end{decision}

\begin{firewall}[Claim boundary after V66]
\label{fire:v17v66}
The optimal radial estimate does not prove the scale-safe global
environment Hessian.  It supplies a sharper sufficient condition and
rules out one proof architecture.  Environmental susceptibility,
remainder-current control, continuum measure, gauge cofinality,
physical Einstein action, and the Einstein--SM continuum limit remain
open.
\end{firewall}

\section*{Version V66 append-only record}
\addcontentsline{toc}{section}{Version V66 append-only record}
The complete V65 source was retained byte for byte before this
continuation.  It had \(559256\) bytes, \(17186\) lines, and SHA--256
\[
 \texttt{59883E0D627280FF9490D0CFD93ED726798F67081A2B150A584C6362D4B5F4FF}.
\]
V66 optimized the radial allocation and proved its canonical-scaling
barrier.


\section{V67: oriented transport of the neighbor-field star}
\label{sec:v17v67}

The unweighted estimate (64.15) discards the conformal weights carried
by the two ends of every hopping edge.  This version restores those
weights and determines exactly why the resulting estimate cannot be
closed edge by edge.

\subsection{Opposite-end diagonal variables}

For an edge \(f\ni x\), let \(\bar x(f)\) denote its other endpoint and
write
\[
 g_{x,f}:=\chi_f e^{\sigma_{x,f}d_f},
 \qquad
 D_{x,f}:=g_{x,f}\|y_x\|^2,                       \tag{67.1}
\]
where \(\sigma_{x,f}=1\) at the positive end and \(-1\) at the
negative end.  Then
\[
 g_{x,f}g_{\bar x(f),f}=\chi_f^2,
 \qquad
 H_x=\sum_{f\ni x}g_{x,f}.                        \tag{67.2}
\]
Define the diagonal star transported into \(x\) from the opposite
ends by
\[
 S_x^{\rm opp}:=
 \sum_{f\ni x}D_{\bar x(f),f}.                   \tag{67.3}
\]

\begin{lemma}[Weighted star factorization]
\label{lem:v17v67-weighted-star}
Pointwise in the Higgs field and the environment,
\[
\boxed{
 \|h_x\|^2\leq H_x S_x^{\rm opp}.}               \tag{67.4}
\]
\end{lemma}

\begin{proof}
By (64.2) and (67.2),
\[
 h_x=
 \sum_{f\ni x}
 \sqrt{g_{x,f}}
 \left(\sqrt{g_{\bar x(f),f}}
 U_f^{(x)}y_{\bar x(f)}\right).                  \tag{67.5}
\]
The transports preserve the fibre norm.  Cauchy--Schwarz for the
finite sum in (67.5) gives (67.4).
\end{proof}

\begin{lemma}[Radial estimate with the oriented star retained]
\label{lem:v17v67-radial}
For \(e\ni x\) and \(0<\delta<2\),
\[
\boxed{
 \mathbb E_x[D_{x,e}H_x]
 \leq
 \frac{q_Hg_{x,e}}{2-\delta}
 +\frac{g_{x,e}S_x^{\rm opp}}
       {\delta(2-\delta)}.}                      \tag{67.6}
\]
\end{lemma}

\begin{proof}
Before the last estimate in (66.5), its neighbor-field term is
\[
 \frac{g_{x,e}H_x\|h_x\|^2}
      {\delta(2-\delta)c_x^2}.                   \tag{67.7}
\]
Use (67.4) and \(H_x^2/c_x^2\leq1\).  The contact term is treated as
in Lemma~\ref{lem:v17v66-delta}.
\end{proof}

\subsection{Uniform conditional site moments already implied by V61}

At fixed \((r,U)\), the Higgs action is invariant under the global
reflection \(y\mapsto-y\).  Hence every real component of every
\(y_x\) has zero \(\pi_\eta\)-mean.

\begin{lemma}[Conditional second-moment ceiling]
\label{lem:v17v67-moment}
Under the V61 influence condition \(r_{\rm H}<1\),
\[
\boxed{
 \mathbb E_{\pi_\eta}\|y_x\|^2
 \leq q_H C_{\rm H,glob}}                        \tag{67.8}
\]
for every site and every fixed environment.
\end{lemma}

\begin{proof}
Apply the global conditional Poincar\'e inequality (61.11) to each
real coordinate function \(y_x^j\).  Its mean is zero and the squared
full Higgs gradient is one, so
\(\mathbb E(y_x^j)^2\leq C_{\rm H,glob}\).  Sum over
\(j=1,\ldots,q_H\).
\end{proof}

It follows that
\[
 \mathbb E_{\pi_\eta}S_x^{\rm opp}
 \leq q_HC_{\rm H,glob}
 \sum_{f\ni x}g_{\bar x(f),f}.                  \tag{67.9}
\]
Combining (67.6) and (67.9) yields a deterministic two-edge kernel:
\[
\begin{aligned}
 \mathbb E_{\pi_\eta}[D_{x,e}H_x]
 \leq{}&\frac{q_Hg_{x,e}}{2-\delta}\\
 &+\frac{q_HC_{\rm H,glob}}
 {\delta(2-\delta)}
 \sum_{f\ni x}g_{x,e}g_{\bar x(f),f}.           \tag{67.10}
\end{aligned}
\]

\subsection{Why coefficientwise absorption fails}

The cross-edge factor in (67.10) is
\[
 g_{x,e}g_{\bar x(f),f}
 =\chi_e\chi_f
 e^{\sigma_{x,e}d_e-\sigma_{x,f}d_f}.            \tag{67.11}
\]
It mixes two independently variable conformal differences.

\begin{countertheorem}[No single-edge mismatch domination]
\label{ctr:v17v67-edgewise}
Suppose the graph has a vertex \(x\) incident to two distinct edges
\(e,f\) with positive hopping coefficients.  There is no finite
constant \(K\) such that
\[
 g_{x,e}g_{\bar x(f),f}
 \leq K\cosh(2d_e)                                \tag{67.12}
\]
for every conformal environment.
\end{countertheorem}

\begin{proof}
Set \(d_e=0\).  Choose the sign of \(d_f\) so that
\(-\sigma_{x,f}d_f\to+\infty\).  The left side of (67.12) diverges,
whereas its right side remains \(K\).
\end{proof}

Even summing mismatch energies does not justify the diagonalized
coefficient estimate at the incidence-form level.  On the two-edge
star, choose a vertex vector \(u\) with
\((Bu)_e\neq0\) but \((Bu)_f=0\).  Then the term
\(g_{x,e}g_{\bar x(f),f}(Bu)_e^2\) diverges along the environment used
above, while the large coefficient \(\cosh(2d_f)\) is multiplied by
\((Bu)_f^2=0\).  The fixed vertex floor \(2\kappa_g\|u\|^2\) cannot
absorb that divergence.

This is not a failure of the exact effective Hessian.  The offending
cross-edge product entered when \(\|h_x\|^2\) was bounded by
\(H_xS_x^{\rm opp}\).  In the original covariance gradient, an edge
with \((Bu)_f=0\) contributes nothing and cannot create the displayed
divergence.

\begin{assessment}[What V67 closes]
\label{ass:v17v67-status}
V67 replaces the unweighted star by the exact oriented factorization
(67.4) and proves the uniform conditional site-moment bound (67.8).
It then gives a concrete two-edge counterexample to coefficientwise
absorption of the resulting kernel.  The obstruction is caused by
diagonalizing the exact vertex sum, so the next argument must retain
that sum.
\end{assessment}

\begin{decision}[V68: return to the exact incidence gradient]
\label{dec:v17v67-next}
Write the Higgs score coupled to an arbitrary conformal direction
\(u\) before applying Cauchy--Schwarz.  Bound its variance using V61
and (67.8), with the exact vertex expression
\(\sum_{e\ni x}\sigma_{x,e}g_{x,e}(Bu)_e\) intact.
Compare the resulting quadratic form directly with
\(\sum_e4J\cosh(2d_e)(Bu)_e^2\).  This removes the false two-edge
row-sum; its remaining scaling margin will decide whether a weighted
inverse-Hessian estimate is necessary.
\end{decision}

\begin{firewall}[Claim boundary after V67]
\label{fire:v17v67}
The oriented-star estimate does not prove global environment convexity.
Its countertheorem concerns a sufficient diagonal domination, not the
effective Hessian itself.  Environmental susceptibility,
remainder-current control, continuum measure, gauge cofinality,
physical Einstein action, and the Einstein--SM continuum limit remain
open.
\end{firewall}

\section*{Version V67 append-only record}
\addcontentsline{toc}{section}{Version V67 append-only record}
The complete V66 source was retained byte for byte before this
continuation.  It had \(567119\) bytes, \(17427\) lines, and SHA--256
\[
 \texttt{F956F6D45C9AAC360339E283B63CDA86655694CCC0CD3DBF24409FC2C44FD023}.
\]
V67 restored the oriented star weights and proved the failure of their
coefficientwise absorption.


\section{V68: the undiagonalized incidence-gradient bound}
\label{sec:v17v68}

V67 showed that the two-edge divergence is created by diagonalizing
the exact vertex gradient.  We now return to (63.7), retain its signed
incidence sum, and compare the resulting quadratic form directly with
the mismatch Hessian.

\subsection{Exact vertex score}

For an edge direction \(u=(u_e)\), set
\[
 \mathcal L_x^{d}u
 :=\sum_{e\ni x}\sigma_{x,e}g_{x,e}u_e.           \tag{68.1}
\]
Equation (63.7) is the exact identity
\[
 \nabla_{y_x}Q_u=2y_x\mathcal L_x^{d}u.           \tag{68.2}
\]
No neighboring diagonal coefficient occurs in (68.2).

\begin{theorem}[Direct conditional-current variance bound]
\label{thm:v17v68-current}
Assume the V61 influence margin and let \(D\) be the maximum vertex
degree.  Then, for every fixed environment and every edge vector
\(u\),
\[
\boxed{
 \operatorname{Var}_{\pi_\eta}(Q_u)
 \leq
 8Dq_HC_{\rm H,glob}^{,2}\chi_*^2
 \sum_e\cosh(2d_e)u_e^2.}                        \tag{68.3}
\]
\end{theorem}

\begin{proof}
The global Higgs Poincar\'e inequality and (68.2) give
\[
 \operatorname{Var}_{\pi_\eta}(Q_u)
 \leq
 4C_{\rm H,glob}
 \sum_x\mathbb E_{\pi_\eta}\|y_x\|^2
       (\mathcal L_x^{d}u)^2.                    \tag{68.4}
\]
Apply the site-moment ceiling (67.8), followed by ordinary
Cauchy--Schwarz over at most \(D\) incident edges:
\[
\begin{aligned}
 \operatorname{Var}_{\pi_\eta}(Q_u)
 &\leq4q_HC_{\rm H,glob}^{,2}
       \sum_x(\mathcal L_x^{d}u)^2\\
 &\leq4Dq_HC_{\rm H,glob}^{,2}
       \sum_x\sum_{e\ni x}g_{x,e}^2u_e^2.       \tag{68.5}
\end{aligned}
\]
For \(e=(v,w)\),
\[
 g_{v,e}^2+g_{w,e}^2
 =\chi_e^2(e^{2d_e}+e^{-2d_e})
 =2\chi_e^2\cosh(2d_e)
 \leq2\chi_*^2\cosh(2d_e).                      \tag{68.6}
\]
Sum (68.6) over edges in (68.5).
\end{proof}

This estimate has no product involving distinct edges.  In particular,
the two-edge configuration used in Countertheorem~\ref{ctr:v17v67-edgewise}
does not create a divergent coefficient when \(u_f=0\), because the
term \(g_{x,f}u_f\) vanishes before any inequality is applied.

\subsection{A global environment-convexity branch}

\begin{theorem}[Incidence-form sufficient margin]
\label{thm:v17v68-margin}
If
\[
\boxed{
 J\geq2Dq_HC_{\rm H,glob}^{,2}\chi_*^2,}         \tag{68.7}
\]
then, for every fixed gauge field,
\[
 \nabla_r^2\Phi_{\rm env}(r\mid U)
 \succeq2\kappa_gI                               \tag{68.8}
\]
throughout the full conformal space.
\end{theorem}

\begin{proof}
Combine the exact free-energy Hessian (63.3) with (68.3):
\[
\begin{aligned}
 D^2\mathcal F_H(d,U)[u,u]
 \geq{}&\sum_e\mathbb E_{\pi_\eta}T_e\,u_e^2\\
 &-8Dq_HC_{\rm H,glob}^{,2}\chi_*^2
   \sum_e\cosh(2d_e)u_e^2.                      \tag{68.9}
\end{aligned}
\]
Add the mismatch Hessian.  Under (68.7),
\[
 4J-8Dq_HC_{\rm H,glob}^{,2}\chi_*^2\geq0,      \tag{68.10}
\]
and the mean hopping term in (68.9) is nonnegative.  Substitute
\(u=Bh\) for a conformal variation \(h\) and retain the bare
\(2\kappa_g\|h\|^2\) term.
\end{proof}

Theorem~\ref{thm:v17v68-margin} is global in the conformal differences
and avoids the false row-sum obstruction of V67.  It is nevertheless
not scale safe on the presently certified parameter branch.  With
\(J(a)=\bar Ja^2\), fixed \(\chi_*\), and the scale-independent V61
bound on \(C_{\rm H,glob}\), condition (68.7) fails for sufficiently
small \(a\).

\begin{countertheorem}[Uniform Poincar\'e loses the required stiffness]
\label{ctr:v17v68-pi}
Within the proof of (68.3), replacing the full Higgs covariance by the
single constant \(C_{\rm H,glob}\) cannot exploit the growth of the
Higgs Hessian with \(g_{x,e}\).  Consequently its mismatch comparison
requires an order-one lower bound on \(J\) whenever
\(C_{\rm H,glob}\chi_*\) stays order one.
\end{countertheorem}

\begin{proof}
The coefficient on the right of (68.3) is independent of \(J\) and is
order one under the stated scaling.  Absorption by the mismatch term
requires (68.7), while \(J(a)\to0\).
\end{proof}

This identifies the next upstream loss.  The Higgs quadratic hopping
that becomes stiff when \(g_{x,e}\) is large is present in the full
Higgs Hessian, but the scalar Poincar\'e constant replaces that matrix
by a uniform multiple of the identity.  A scale-safe estimate must
retain its inverse as a matrix.

\begin{assessment}[What V68 closes]
\label{ass:v17v68-status}
V68 proves the global incidence-form current estimate (68.3) and a
corresponding full-tail convexity branch (68.7).  It removes the
spurious cross-edge divergence of V67 and locates the surviving
canonical-scaling loss at the scalar Poincar\'e substitution.
\end{assessment}

\begin{decision}[V69: matrix Brascamp--Lieb acquisition]
\label{dec:v17v68-next}
Write the full real Higgs Hessian as the sum of the site quartic/mass
block and the positive transported edge-square blocks.  Apply the
matrix Brascamp--Lieb inequality to \(Q_u\), initially with
\(m_H^2>0\), and preserve the inverse Hessian in the quadratic form.
Test whether the edge-square block cancels one power of
\(g_{x,e}\) in (68.2).  Record the remaining null-space obstruction
before taking the canonical \(m_H^2(a)\downarrow0\) limit.
\end{decision}

\begin{firewall}[Claim boundary after V68]
\label{fire:v17v68}
Condition (68.7) is sufficient but incompatible with the canonical
scaling on the present bounds.  A scale-safe environment gap,
environmental susceptibility, remainder-current envelope, continuum
measure, gauge cofinality, physical Einstein action, and the
Einstein--SM continuum limit remain open.
\end{firewall}

\section*{Version V68 append-only record}
\addcontentsline{toc}{section}{Version V68 append-only record}
The complete V67 source was retained byte for byte before this
continuation.  It had \(573706\) bytes, \(17631\) lines, and SHA--256
\[
 \texttt{E8A69340E18CE98C69941CF0D6C399CCA98083BFDF1FEE0CA12D4CBE77EB94F3}.
\]
V68 retained the exact incidence gradient and isolated the scalar-gap
scaling loss.


\section{V69: matrix Brascamp--Lieb and the hopping-range split}
\label{sec:v17v69}

The scalar Poincar\'e estimate of V68 replaces the full Higgs Hessian
by a constant.  Here the matrix Brascamp--Lieb form is retained and
the current gradient is split according to the weighted hopping
operator.  One part is controlled with the correct conformal weights;
the other identifies the exact null-mode problem.

\subsection{Weighted edge operator and the full Hessian}

Work in the realified Higgs fibres.  For
\(e=(v,w)\), put
\[
 a_e=e^{d_e/2},\qquad b_e=e^{-d_e/2},
 \qquad a_eb_e=1,                                 \tag{69.1}
\]
and use the orthogonal realification of \(U_e\).  Define
\[
 (\mathsf A_dy)_e
 =\sqrt{\chi_e}\left(a_ey_v-b_eU_e^*y_w\right).  \tag{69.2}
\]
Then the full geometric hopping action is
\[
 S_{\rm hop}(y\mid d,U)=\|\mathsf A_dy\|^2.       \tag{69.3}
\]
For an edge direction \(u\), let \(\mathsf C_{d,u}\) be its
directional derivative:
\[
 (\mathsf C_{d,u}y)_e
 =\frac{u_e\sqrt{\chi_e}}2
   \left(a_ey_v+b_eU_e^*y_w\right).               \tag{69.4}
\]
Thus
\[
 Q_u=2\langle\mathsf A_dy,\mathsf C_{d,u}y\rangle.
                                                               \tag{69.5}
\]

Let \(\mathsf K(y)=\nabla_y^2S_H(y\mid d,U)\).  Direct
differentiation gives
\[
\boxed{
 \mathsf K(y)
 =2m_H^2I+\mathsf K_4(y)+2\mathsf A_d^*\mathsf A_d,}             \tag{69.6}
\]
where the block at site \(x\) of the quartic Hessian is
\[
 (\mathsf K_4(y))_x
 =4\lambda_H\|y_x\|^2I+8\lambda_Hy_xy_x^{\mathsf T}succeq0.    \tag{69.7}
\]
For now take \(m_H^2>0\), so that \(\mathsf K(y)\) is invertible.

\begin{theorem}[Exact matrix covariance acquisition]
\label{thm:v17v69-bl}
The conditional current variance satisfies
\[
\boxed{
 \operatorname{Var}_{\pi_\eta}(Q_u)
 \leq
 \mathbb E_{\pi_\eta}
 \langle v_R+v_N,\mathsf K(y)^{-1}(v_R+v_N)\rangle,}             \tag{69.8}
\]
where
\[
 v_R=2\mathsf A_d^*\mathsf C_{d,u}y,
 \qquad
 v_N=2\mathsf C_{d,u}^*\mathsf A_dy.             \tag{69.9}
\]
\end{theorem}

\begin{proof}
Differentiate (69.5) in \(y\):
\[
 \nabla_yQ_u
 =2(\mathsf A_d^*\mathsf C_{d,u}
    +\mathsf C_{d,u}^*\mathsf A_d)y=v_R+v_N.     \tag{69.10}
\]
Apply the matrix Brascamp--Lieb inequality to the strictly convex
conditional Higgs potential and use (69.6).
\end{proof}

\subsection{The range part spends at most the mean hopping curvature}

For later use define the direction-weighted hopping action
\[
 S_{{\rm hop},u^2}(y)
 :=\sum_eu_e^2\|(\mathsf A_dy)_e\|^2.             \tag{69.11}
\]

\begin{theorem}[Exact hopping-range cancellation]
\label{thm:v17v69-range}
Pointwise in \(y,d,U\),
\[
\boxed{
 \langle v_R,\mathsf K(y)^{-1}v_R\rangle
 \leq T_{u^2}(y)-\frac12S_{{\rm hop},u^2}(y).}    \tag{69.12}
\]
\end{theorem}

\begin{proof}
By (69.6), \(\mathsf K\succeq2\mathsf A_d^*\mathsf A_d\), and
\(v_R\) lies in the range of \(\mathsf A_d^*\).  The variational
formula for an inverse quadratic form therefore gives
\[
\begin{aligned}
 \langle v_R,\mathsf K^{-1}v_R\rangle
 &\leq
 \langle v_R,(2\mathsf A_d^*\mathsf A_d)^\dagger v_R\rangle\\
 &\leq2\|\mathsf C_{d,u}y\|^2.                  \tag{69.13}
\end{aligned}
\]
Indeed, the last expression is obtained by inserting
\(v_R=2\mathsf A_d^*\mathsf C_{d,u}y\); the intervening operator is
the orthogonal projection onto the range of \(\mathsf A_d\).

For
\(p_e=a_ey_v\) and \(q_e'=b_eU_e^*y_w\), the parallelogram identity
gives
\[
 \|p_e+q_e'\|^2
 =2(\|p_e\|^2+\|q_e'\|^2)-\|p_e-q_e'\|^2.        \tag{69.14}
\]
Multiply by \(u_e^2\chi_e/2\), sum over edges, and use
(69.2), (69.4), and the definition of \(T_{u^2}\).  The result is
\[
 2\|\mathsf C_{d,u}y\|^2
 =T_{u^2}(y)-\frac12S_{{\rm hop},u^2}(y),         \tag{69.15}
\]
which proves (69.12).
\end{proof}

Thus the range component has exactly the scale needed to be paid by
the positive term \(\mathbb ET_{u^2}\) in (63.3), with half the
direction-weighted hopping square left as slack.  This cancellation
was invisible in the scalar bound (68.3).

\subsection{The residual component and the null mode}

The vector \(v_N\) vanishes whenever \(\mathsf A_dy=0\), but it need
not lie in the range of \(\mathsf A_d^*\).  This distinction already
appears on one edge.

\begin{lemma}[One-edge null projection]
\label{lem:v17v69-null}
Consider one edge, suppress the orthogonal transport, and put
\(R=a y_v-b y_w\), \(s=a^2+b^2\).  The kernel of the hopping
operator consists of vectors \((bz,az)\).  The orthogonal projection
of the full score gradient onto this kernel has squared norm
\[
\boxed{
 \|P_{\ker\mathsf A_d}\nabla_yQ_u\|^2
 =\frac{4u^2\chi}{s},S_{\rm hop}(y),}            \tag{69.16}
\]
where \(S_{\rm hop}(y)=\chi\|R\|^2\).
\end{lemma}

\begin{proof}
The score is
\(Q_u=u\chi(a^2\|y_v\|^2-b^2\|y_w\|^2)\), so
\[
 \nabla_yQ_u=2u\chi(a^2y_v,-b^2y_w).             \tag{69.17}
\]
For a fibre vector \(z\), its inner product with the kernel vector
\((bz,az)\) is
\[
 2u\chi ab\langle ay_v-by_w,z\rangle
 =2u\chi\langle R,z\rangle.                     \tag{69.18}
\]
The map \(z\mapsto(bz,az)\) multiplies squared norms by \(s\).
Therefore the squared projection is
\(4u^2\chi^2\|R\|^2/s\), which is (69.16).
\end{proof}

\begin{countertheorem}[Hopping curvature alone cannot control the
full score]
\label{ctr:v17v69-null}
There is no finite quadratic-form estimate of
\(\nabla_yQ_u\) by the inverse of the hopping Hessian alone, uniform
over Higgs configurations.  The obstruction persists for every
nonzero \(u\) and \(\chi\).
\end{countertheorem}

\begin{proof}
Choose a one-edge configuration with \(R\neq0\).  By (69.16), the
score gradient has a nonzero component in
\(\ker\mathsf A_d=\ker(2\mathsf A_d^*\mathsf A_d)\).  The variational
inverse quadratic form of that singular Hessian is infinite on such a
vector.  Hence the hopping Hessian alone cannot furnish the claimed
bound.
\end{proof}

Using only the mass part of (69.6) controls this component by a factor
\((2m_H^2)^{-1}\), reproducing the canonical \(a^{-2}\) loss of V58.
The remaining candidate is the quartic block (69.7).  Its curvature
vanishes at the origin, but (69.16) shows simultaneously that the
null projection is proportional to the hopping mismatch \(R\).  A
successful estimate must use this joint vanishing rather than bound
the inverse Hessian uniformly.

For any \(\varepsilon>0\), (69.8) and the Hilbert norm induced by
\(\mathsf K^{-1}\) give the exact ledger
\[
\begin{aligned}
 \operatorname{Var}(Q_u)
 \leq{}&(1+\varepsilon)
 \mathbb E\left[T_{u^2}-\tfrac12S_{{\rm hop},u^2}\right]\\
 &+(1+\varepsilon^{-1})
 \mathbb E\langle v_N,\mathsf K^{-1}v_N\rangle. \tag{69.19}
\end{aligned}
\]
The second line is now the sole matrix-Hessian acquisition.

\begin{assessment}[What V69 closes]
\label{ass:v17v69-status}
V69 writes the full matrix Brascamp--Lieb problem and proves that the
range part of the score costs no more than the available mean hopping
curvature, with the explicit slack (69.12).  It also computes the
one-edge null projection and proves why hopping curvature by itself is
insufficient.  The unresolved quantity is the residual inverse-Hessian
form in (69.19), where quartic curvature and hopping mismatch vanish
together.
\end{assessment}

\begin{decision}[V70 audit; then quartic-null interpolation]
\label{dec:v17v69-next}
Perform the mandatory YM/NS audit.  Then estimate
\(\langle v_N,\mathsf K^{-1}v_N\rangle\) by a variational interpolation
between \(\mathsf K_4(y)\) and
\(2\mathsf A_d^*\mathsf A_d\).  Begin with the one-edge scalar model,
where (69.16) diagonalizes the only null channel.  Determine whether
the resulting bound is paid by the half-hopping slack in (69.12) plus
quartic Ward moments, uniformly as \(m_H^2\downarrow0\).
\end{decision}

\begin{firewall}[Claim boundary after V69]
\label{fire:v17v69}
The range cancellation is exact, but the residual term (69.19) has
not been bounded at canonical scaling.  Global environment convexity,
environmental susceptibility, remainder-current control, continuum
measure, gauge cofinality, physical Einstein action, and the
Einstein--SM continuum limit remain open.
\end{firewall}

\section*{Version V69 append-only record}
\addcontentsline{toc}{section}{Version V69 append-only record}
The complete V68 source was retained byte for byte before this
continuation.  It had \(579675\) bytes, \(17804\) lines, and SHA--256
\[
 \texttt{88E432260D0B025606754BA827EC9A7B15FA0A71E4231392B705A4B2DAA243D6}.
\]
V69 proved the hopping-range cancellation and isolated the quartic
null-mode inverse-Hessian acquisition.


\section{V70: compulsory companion audit before quartic-null analysis}
\label{sec:v17v70}

\subsection{Active companion checkpoints}

The newest active Yang--Mills note inspected is
\begin{center}\small
\path{Renewal_Yang_Mills_Mass_Gap_Research_Notes_V91.tex}.
\end{center}
It has \(711720\) bytes, \(19441\) lines, and SHA--256
\[
 \texttt{8F08BD71F91404CEA5F37DEB2E23F66783EBD652ABB26520F4E38781E72382C9}.
                                                               \tag{70.1}
\]
YM V90 completed the seventh direct-plus sign cell and its inverted
direct-minus companion.  V91 begins the eighth direct-plus atlas: seven
Boolean seed centres are inherited exactly and the sole new centre is
certified on a full radius-\(1/8\) product box.  The rest of the eighth
cell, full-torus direct-plus positivity, cofinal control, and a mass gap
remain open.

The newest active Navier--Stokes note remains
\begin{center}\small
\path{Renewal_NS_Stability_Research_Notes_V26.tex}.
\end{center}
It has \(278283\) bytes, \(5319\) lines, and SHA--256
\[
 \texttt{82C887F8C2C69E4F28FAD7C3DC8DE5B9099CB5A4BF431EBA1FFF3E91935978AD}.
                                                               \tag{70.2}
\]
Its last result is still the pointwise rank-one Oseen-kernel norm
calculation described in V65.

\begin{theorem}[V70 companion-audit decision]
\label{thm:v17v70-decision}
No theorem from the active companion notes controls the residual
inverse-Hessian form
\[
 \mathbb E_{\pi_\eta}
 \langle v_N,\mathsf K(y)^{-1}v_N\rangle          \tag{70.3}
\]
in (69.19), so no companion statement is imported into the SM proof
chain at V70.
\end{theorem}

\begin{proof}
YM V91 concerns positivity of a fixed finite-dimensional Wilson symbol
on a momentum cube.  It contains neither the conditional Higgs Hessian
\(\mathsf K(y)\) nor its inverse.  Its completion would be relevant to
the later gauge-sector programme, but its present conclusion is not a
gauge-environment covariance inequality.  NS V26 concerns rank-one
contractions of derivatives of the Oseen kernel and has no variable in
(70.3).  Hence neither conclusion discharges a premise of V69.
\end{proof}

\begin{assessment}[Direction after the audit]
\label{ass:v17v70-direction}
The YM programme has reduced its finite-cutoff symbol gap to the last
direct-plus sign cell.  That is genuine progress toward a future gauge
input, but the current SM bottleneck is internal to the conformal--Higgs
Schur complement.  The next calculation remains the one-edge
quartic/null interpolation fixed in Decision~\ref{dec:v17v69-next}.
\end{assessment}

\begin{decision}[V71: scalar one-edge Schur complement]
\label{dec:v17v70-next}
Diagonalize the two-site, one-real-component Hessian in its hopping
range and null directions.  Retain the site quartic curvatures exactly,
invert the resulting \(2\times2\) matrix, and evaluate the score-gradient
quadratic form.  Determine whether the half-hopping slack in (69.12)
pays the null contribution uniformly at zero mass; if it does not,
identify the precise amplitude ratio that escapes.
\end{decision}

\begin{firewall}[Claim boundary after V70]
\label{fire:v17v70}
The audit supplies no new SM estimate.  The quartic/null residual,
global environment convexity at canonical scaling, environmental
susceptibility, remainder-current control, continuum measure, gauge
cofinality, physical Einstein action, and the Einstein--SM continuum
limit remain open.
\end{firewall}

\section*{Version V70 append-only record}
\addcontentsline{toc}{section}{Version V70 append-only record}
The complete V69 source was retained byte for byte before this
continuation.  It had \(588198\) bytes, \(18059\) lines, and SHA--256
\[
 \texttt{6A7A6736BB53EE79A37E3BD6F95DF5AFEB48D24721C9A6D14B2245F6D6367A49}.
\]
V70 audited YM V91 and NS V26 and retained the quartic/null acquisition.


\section{V71: exact one-edge quartic/null Schur complement}
\label{sec:v17v71}

V69 reduced the matrix problem to the component meeting the hopping
null mode.  The smallest model containing that interaction is one
edge and one real Higgs component.  Its Hessian is a \(2\times2\)
matrix and can be inverted without loss.

\subsection{The scalar block}

Suppress the orthogonal gauge transport and write
\[
 x=y_v,\qquad z=y_w,\qquad
 a=e^{d/2},\qquad b=e^{-d/2},\qquad ab=1.          \tag{71.1}
\]
The scalar action and directional conformal score are
\[
\begin{aligned}
 S(x,z)&=m_H^2(x^2+z^2)+\lambda_H(x^4+z^4)
        +\chi(ax-bz)^2,\\
 Q_u(x,z)&=u\chi(a^2x^2-b^2z^2).                 \tag{71.2}
\end{aligned}
\]
Put
\[
 \alpha=2m_H^2+12\lambda_Hx^2,
 \qquad
 \beta=2m_H^2+12\lambda_Hz^2.                   \tag{71.3}
\]
Then
\[
 \mathsf K=
 \begin{pmatrix}
  \alpha+2\chi a^2&-2\chi\\
  -2\chi&\beta+2\chi b^2
 \end{pmatrix},
 \qquad
 v:=\nabla Q_u=2u\chi(a^2x,-b^2z).               \tag{71.4}
\]

\begin{theorem}[Exact scalar inverse-Hessian form]
\label{thm:v17v71-exact}
Whenever \(\mathsf K\) is invertible,
\[
\boxed{
 v^{\mathsf T}\mathsf K^{-1}v
 =4u^2\chi^2
 \frac{
  \beta a^4x^2+\alpha b^4z^2+2\chi(ax-bz)^2}
 {\alpha\beta+2\chi(\alpha b^2+\beta a^2)}.}      \tag{71.5}
\]
\end{theorem}

\begin{proof}
Because \(a^2b^2=1\),
\[
 \det\mathsf K
 =\alpha\beta+2\chi(\alpha b^2+\beta a^2).       \tag{71.6}
\]
Insert the inverse of the symmetric matrix (71.4).  After factoring
\(4u^2\chi^2\), its numerator is
\[
\begin{aligned}
 &a^4x^2(\beta+2\chi b^2)
 -4\chi xz
 +b^4z^2(\alpha+2\chi a^2)\\
 &\qquad=
 \beta a^4x^2+\alpha b^4z^2
 +2\chi(a^2x^2+b^2z^2-2abxz),
\end{aligned}
\]
which is the numerator in (71.5).
\end{proof}

Formula (71.5) is the exact longitudinal scalar block of the
realified multi-component Hessian when all other components vanish.
It therefore tests every proposed pointwise matrix domination for the
full model.

\subsection{Zero-mass small-amplitude limit}

Set \(m_H^2=0\), choose nonzero real \(X,Z\), and let
\[
 x=tX,qquad z=tZ,qquad t\downarrow0.             \tag{71.7}
\]

\begin{theorem}[Quartic resolution of the null singularity]
\label{thm:v17v71-limit}
Along (71.7),
\[
\boxed{
 \lim_{t\downarrow0}v^{\mathsf T}\mathsf K^{-1}v
 =\frac{u^2\chi^2}{3\lambda_H}
 \frac{(aX-bZ)^2}{b^2X^2+a^2Z^2}.}               \tag{71.8}
\]
The limit is positive unless \(aX=bZ\), the hopping-null alignment.
\end{theorem}

\begin{proof}
Now
\(\alpha=12\lambda_Ht^2X^2\) and
\(\beta=12\lambda_Ht^2Z^2\).  In (71.5), the first two numerator
terms are order \(t^4\), while the last is
\(2\chi t^2(aX-bZ)^2\).  The denominator is
\[
 144\lambda_H^2t^4X^2Z^2
 +24\chi\lambda_Ht^2(b^2X^2+a^2Z^2).             \tag{71.9}
\]
Cancel \(t^2\) and take the limit.
\end{proof}

The quartic Hessian makes the inverse form finite: the score's null
component is order \(t\), while the quartic curvature in that direction
is order \(t^2\).  Their quotient has the nonzero limit (71.8).
In contrast, both available pointwise hopping quantities vanish:
\[
 T_{u^2}=u^2\chi t^2(a^2X^2+b^2Z^2),
 \qquad
 S_{{\rm hop},u^2}=u^2\chi t^2(aX-bZ)^2.          \tag{71.10}
\]

\begin{countertheorem}[The half-hopping slack cannot pay the null
term pointwise]
\label{ctr:v17v71-slack}
At zero Higgs mass there is no finite constant \(C\) such that
\[
 v^{\mathsf T}\mathsf K^{-1}v
 \leq T_{u^2}+C S_{{\rm hop},u^2}                 \tag{71.11}
\]
for every nonzero scalar Higgs configuration, even with \(d\),
\(u\), \(\chi\), and \(\lambda_H\) fixed and positive.
\end{countertheorem}

\begin{proof}
Choose \(X,Z\) with \(aX\neq bZ\) and use (71.7).  The left side of
(71.11) tends to the positive number (71.8), whereas its right side
tends to zero by (71.10).
\end{proof}

This disproves pointwise payment by the half-square slack in (69.12).
It does not prove that the expected Brascamp--Lieb form is too large:
the small-amplitude set has shrinking volume and must be integrated
against the conditional Gibbs density before drawing that conclusion.

\subsection{Sharp conformal size of the escaping channel}

The amplitude quotient in (71.8) has an exact supremum.

\begin{lemma}[Worst amplitude ratio]
\label{lem:v17v71-rayleigh}
For fixed \(d\),
\[
 \sup_{(X,Z)\neq(0,0)}
 \frac{(aX-bZ)^2}{b^2X^2+a^2Z^2}
 =\frac{a^2}{b^2}+\frac{b^2}{a^2}
 =2\cosh(2d).                                    \tag{71.12}
\]
One maximizing direction is
\[
 (X,Z)=\left(\frac a{b^2},-\frac b{a^2}\right). \tag{71.13}
\]
\end{lemma}

\begin{proof}
Equation (71.12) is the Rayleigh quotient of the linear functional
\((a,-b)\) in the diagonal metric
\(\operatorname{diag}(b^2,a^2)\).  Its squared dual norm is
\(a^2/b^2+b^2/a^2\), attained by the inverse-metric image (71.13).
Use \(a/b=e^d\).
\end{proof}

Consequently the largest zero-mass small-amplitude limit is
\[
 \frac{2u^2\chi^2}{3\lambda_H}\cosh(2d).         \tag{71.14}
\]
If one insists on dominating the Brascamp--Lieb integrand pointwise by
the positive mean-hopping integrand plus the mismatch curvature, then
(71.10)--(71.14) force at least
\[
 4J\cosh(2d)
 \geq\frac{2\chi^2}{3\lambda_H}\cosh(2d),
 \qquad
 J\geq\frac{\chi^2}{6\lambda_H}.                 \tag{71.15}
\]
This order-one requirement again fails when \(J(a)=\bar Ja^2\) and
\(\chi,\lambda_H\) remain fixed.

\begin{countertheorem}[Pointwise matrix domination is not scale safe]
\label{ctr:v17v71-pointwise}
The exact matrix Brascamp--Lieb integrand cannot be absorbed pointwise
by \(T_{u^2}+4J(a)\cosh(2d)u^2\) uniformly along the canonical
scaling, even in the one-edge scalar block.
\end{countertheorem}

\begin{proof}
Use the maximizing amplitude direction (71.13), take the limit
(71.7), and compare (71.14) with the proposed mismatch term.  The
necessary condition is (71.15), whose right side is fixed and positive
while \(J(a)\to0\).
\end{proof}

The countertheorem concerns pointwise domination of a covariance upper
bound.  The next available upstream move is to average the exact
rational expression (71.5), where the quartic Gibbs density may control
the shrinking small-amplitude sector without an order-one mismatch
coefficient.

\begin{assessment}[What V71 closes]
\label{ass:v17v71-status}
V71 exactly inverts the scalar two-site Hessian, proves that quartic
curvature resolves the hopping-null singularity, and finds its finite
nonzero small-amplitude limit.  That limit rules out payment by the
half-hopping slack and rules out pointwise canonical-scale matrix
domination.  It leaves open an integrated inverse-Hessian estimate.
\end{assessment}

\begin{decision}[V72: integrate before comparison]
\label{dec:v17v71-next}
Average (71.5) under the exact two-site conditional density.  Split
radial amplitude from the ratio \(X/Z\), use the quartic Ward identity
to control the small-amplitude sector, and seek an expectation bound
in terms of \(\mathbb ET_{u^2}\) plus an error that vanishes with the
lattice scaling.  If the one-edge expectation retains an order-one
\(\chi^2/\lambda_H\) term, record it as an actual obstruction to the
global-convexity route and return to a negative-Sobolev environment
estimate rather than pointwise convexity.
\end{decision}

\begin{firewall}[Claim boundary after V71]
\label{fire:v17v71}
The small-amplitude countertheorem rules out a proof architecture, not
environment convexity or the continuum theory.  The averaged residual,
scale-safe environment mixing, environmental susceptibility,
remainder-current control, continuum measure, gauge cofinality,
physical Einstein action, and the Einstein--SM continuum limit remain
open.
\end{firewall}

\section*{Version V71 append-only record}
\addcontentsline{toc}{section}{Version V71 append-only record}
The complete V70 source was retained byte for byte before this
continuation.  It had \(592044\) bytes, \(18152\) lines, and SHA--256
\[
 \texttt{8B36A2769FE99BF6A323FF5BCA1515AFF90E008F089FE2943BD2E18E9E3E7DDA}.
\]
V71 solved the scalar quartic/null Schur complement and proved the
pointwise canonical-scaling barrier.


\section{V72: integrated scalar reduction and the dimensionless target}
\label{sec:v17v72}

Pointwise comparison failed in V71 because the inverse-Hessian form
has a finite directional limit at zero amplitude.  Integration does
remove the apparent singularity.  This version performs the exact
massless scaling reduction and shows what integration must still prove.

\subsection{Dimensionless two-site law}

Retain the scalar one-edge model of V71 with \(m_H^2=0\), and set
\[
 \rho:=\frac{\chi}{\sqrt{\lambda_H}},
 \qquad x=\lambda_H^{-1/4}X,
 \qquad z=\lambda_H^{-1/4}Z.                      \tag{72.1}
\]
Let \(\nu_{\rho,d}\) be the probability law on \(\mathbb R^2\) with
density
\[
 d\nu_{\rho,d}(X,Z)
 =\mathcal Z_{\rho,d}^{-1}
 \exp\{-X^4-Z^4-\rho(aX-bZ)^2\}\,dX\,dZ.         \tag{72.2}
\]
Put
\[
 A=12X^2,qquad B=12Z^2,                          \tag{72.3}
\]
and, away from its measure-zero indeterminacy at the origin, define
\[
 \iota_{\rho,d}(X,Z)
 =4\frac{
  Ba^4X^2+Ab^4Z^2+2\rho(aX-bZ)^2}
 {AB+2\rho(Ab^2+Ba^2)}.                          \tag{72.4}
\]
Define the three dimensionless functions
\[
\begin{aligned}
 \mathcal T(\rho,d)
 &:=\mathbb E_{\nu_{\rho,d}}(a^2X^2+b^2Z^2),\\
 \mathcal I(\rho,d)
 &:=\mathbb E_{\nu_{\rho,d}}\iota_{\rho,d}(X,Z),\\
 \mathcal V(\rho,d)
 &:=\operatorname{Var}_{\nu_{\rho,d}}
       (a^2X^2-b^2Z^2).                           \tag{72.5}
\end{aligned}
\]

\begin{lemma}[Finiteness of the integrated inverse form]
\label{lem:v17v72-finite}
For every \(\rho>0\) and finite \(d\), the three quantities in
(72.5) are finite.  In fact,
\[
 0\leq\iota_{\rho,d}(X,Z)
 \leq\frac43\cosh(2d)                              \tag{72.6}
\]
almost everywhere, uniformly in \(\rho\geq0\).
\end{lemma}

\begin{proof}
The denominator in (72.4) is at least \(AB\).  Its first two numerator
terms therefore contribute at most
\[
 4\frac{Ba^4X^2}{AB}+4\frac{Ab^4Z^2}{AB}
 =\frac{a^4+b^4}{3}=\frac23\cosh(2d).
\]
Separately, the denominator is at least
\(2\rho(Ab^2+Ba^2)\).  The last numerator term contributes at most
\[
 \frac13\frac{(aX-bZ)^2}{b^2X^2+a^2Z^2}
 \leq\frac23\cosh(2d),
\]
by Lemma~\ref{lem:v17v71-rayleigh}.  Summing proves (72.6), with axis
values supplied by limits.  The quartic density supplies the moments
needed for \(\mathcal T\) and \(\mathcal V\).
\end{proof}

\subsection{Exact scaling identities}

\begin{theorem}[No hidden lattice-spacing gain from integration]
\label{thm:v17v72-scaling}
For the massless scalar one-edge law,
\[
\boxed{
\begin{aligned}
 \mathbb E T_{u^2}
 &=u^2\rho\,\mathcal T(\rho,d),\\
 \mathbb E[v^{\mathsf T}\mathsf K^{-1}v]
 &=u^2\rho^2\,\mathcal I(\rho,d),\\
 \operatorname{Var}(Q_u)
 &=u^2\rho^2\,\mathcal V(\rho,d).
\end{aligned}}                                    \tag{72.7}
\]
In particular, if \(\chi\) and \(\lambda_H\) are fixed along the
canonical lattice scaling, all three coefficients in (72.7) are
independent of the lattice spacing.
\end{theorem}

\begin{proof}
The change of variables (72.1) converts the action to (72.2); its
Jacobian cancels from normalized expectations.  The hopping curvature
is
\[
 T_{u^2}=u^2\rho(a^2X^2+b^2Z^2).                 \tag{72.8}
\]
In (71.5), \(\alpha=\sqrt{\lambda_H}A\),
\(\beta=\sqrt{\lambda_H}B\), and
\(\chi=\rho\sqrt{\lambda_H}\).  The common power of
\(\lambda_H\) cancels, leaving
\(u^2\rho^2\iota_{\rho,d}\).  Finally,
\[
 Q_u=u\rho(a^2X^2-b^2Z^2),                       \tag{72.9}
\]
which proves the variance identity.
\end{proof}

The exact Higgs free-energy curvature in this block is therefore
\[
 D_d^2\mathcal F_H[u,u]
 =u^2\{\rho\mathcal T(\rho,d)
       -\rho^2\mathcal V(\rho,d)\}.              \tag{72.10}
\]
Brascamp--Lieb gives the sufficient replacement
\[
 D_d^2\mathcal F_H[u,u]
 \geq u^2\{\rho\mathcal T(\rho,d)
       -\rho^2\mathcal I(\rho,d)\}.              \tag{72.11}
\]

\begin{corollary}[Dimensionless closure criteria]
\label{cor:v17v72-criteria}
At zero mass, the one-edge scalar effective environment has
nonnegative Higgs free-energy curvature exactly when
\[
 \rho\mathcal V(\rho,d)\leq\mathcal T(\rho,d).   \tag{72.12}
\]
The matrix Brascamp--Lieb route proves that conclusion under the
stronger condition
\[
 \rho\mathcal I(\rho,d)\leq\mathcal T(\rho,d).   \tag{72.13}
\]
With mismatch curvature included, the sufficient condition is
\[
 \rho^2\mathcal I(\rho,d)
 \leq\rho\mathcal T(\rho,d)+4J\cosh(2d).         \tag{72.14}
\]
\end{corollary}

\begin{proof}
Equations (72.12)--(72.14) are respectively (72.10), (72.11), and
(72.11) plus the mismatch Hessian, after division by \(u^2\).
\end{proof}

Thus integration makes the inverse form finite but creates no power of
the lattice spacing.  As \(J(a)\downarrow0\), the Brascamp--Lieb route
must establish the genuine dimensionless inequality (72.13), rather
than an estimate with an order-one error.

\subsection{The weak-hopping checkpoint and its nonuniformity in \(d\)}

At \(\rho=0\), the law (72.2) is the product pure-quartic law.  Its
second moment is
\[
 m_2^{(4)}
 :=\frac{\int_{\mathbb R}X^2e^{-X^4}\,dX}
         {\int_{\mathbb R}e^{-X^4}\,dX}
 =\frac{\Gamma(3/4)}{\Gamma(1/4)}.                \tag{72.15}
\]
Direct cancellation in (72.4) gives, almost everywhere,
\[
\begin{aligned}
 \mathcal T(0,d)&=2m_2^{(4)}\cosh d,\\
 \mathcal I(0,d)&=\frac23\cosh(2d).              \tag{72.16}
\end{aligned}
\]

\begin{lemma}[Fixed-d weak-hopping Brascamp--Lieb branch]
\label{lem:v17v72-weak}
For every fixed finite \(d\), there exists
\(\rho_0(d)>0\) such that (72.13) holds for
\(0<\rho\leq\rho_0(d)\).
\end{lemma}

\begin{proof}
The quartic density provides a common integrable dominator for \(\rho\)
in a compact interval, and (72.6) uniformly dominates the inverse
form.  Dominated convergence gives continuity of
\(\mathcal T\) and \(\mathcal I\) at \(\rho=0\).  The first quantity in
(72.16) is positive and finite, while
\(\rho\mathcal I(\rho,d)\to0\).  The inequality follows for small
enough \(\rho\).
\end{proof}

This checkpoint is not uniform in the conformal tail.  Already at
\(\rho=0\), the ratio of the two coefficients is
\[
 \frac{\mathcal T(0,d)}{\mathcal I(0,d)}
 =3m_2^{(4)}\frac{\cosh d}{\cosh(2d)}
 \longrightarrow0
 \quad (|d|\to\infty).                            \tag{72.17}
\]
Hence a perturbative proof based only on continuity at \(\rho=0\)
cannot supply one positive hopping threshold for all \(d\).  The
mismatch term in (72.14) has the correct \(\cosh(2d)\) tail, but its
coefficient tends to zero canonically, so a matched estimate in
\((\rho,d,J)\) is required.

\begin{assessment}[What V72 closes]
\label{ass:v17v72-status}
V72 integrates the exact scalar inverse-Hessian block, proves its
finiteness, and reduces both the actual and Brascamp--Lieb curvatures
to three dimensionless functions.  It proves a fixed-d weak-hopping
branch and shows why that perturbative argument is nonuniform in the
conformal tail.  Integration alone supplies no lattice-spacing gain.
\end{assessment}

\begin{decision}[V73: use the exact scalar covariance]
\label{dec:v17v72-next}
Attack the necessary and sufficient inequality (72.12) directly,
before further estimating by \(\mathcal I\).  Use the reciprocal
rescaling \((x,z)\mapsto(e^{-d/2}X,e^{d/2}Z)\) to move all \(d\)
dependence into the two on-site quartic potentials, and derive the
second derivative of the resulting partition function.  Test whether
one-dimensional integration by parts proves (72.12) for every
\(\rho,d\).  If it fails, retain the exact sign-changing term and
switch the environment target from pointwise convexity to the
negative-Sobolev covariance inequality required in V62.
\end{decision}

\begin{firewall}[Claim boundary after V72]
\label{fire:v17v72}
Only a scalar fixed-d weak-hopping branch is proved.  Uniform
environment convexity, the multi-edge covariance estimate,
environmental susceptibility, remainder-current control, continuum
measure, gauge cofinality, physical Einstein action, and the
Einstein--SM continuum limit remain open.
\end{firewall}

\section*{Version V72 append-only record}
\addcontentsline{toc}{section}{Version V72 append-only record}
The complete V71 source was retained byte for byte before this
continuation.  It had \(600082\) bytes, \(18401\) lines, and SHA--256
\[
 \texttt{EDD9EA62CE44AC1439C0B279503E0B6B15BA4DF18CAD106855CC2EBE45D8E511}.
\]
V72 derived the integrated dimensionless scalar criteria and the
nonuniform weak-hopping branch.


\section{V73: a scaling Ward identity for the exact scalar curvature}
\label{sec:v17v73}

V72 separated the true covariance \(\mathcal V\) from its
Brascamp--Lieb upper bound \(\mathcal I\).  We now attack the true
quantity.  A volume-preserving reciprocal scaling of the two Higgs
variables gives an exact Ward identity and proves strict convexity on
the symmetric conformal slice for every \(\rho>0\).

\subsection{The reciprocal-scaling generator}

Under the law \(\nu_{\rho,d}\) of (72.2), write
\[
\begin{aligned}
 q_d(X,Z)&:=a^2X^2-b^2Z^2,\\
 T_d(X,Z)&:=a^2X^2+b^2Z^2,\\
 W(X,Z)&:=2(X^4-Z^4).                             \tag{73.1}
\end{aligned}
\]
Let
\[
 \mathcal L:=\frac12(X\partial_X-Z\partial_Z).   \tag{73.2}
\]
This vector field has zero Euclidean divergence.  If
\[
 \mathcal S_{\rho,d}=X^4+Z^4+\rho(aX-bZ)^2,       \tag{73.3}
\]
then direct differentiation gives
\[
 \mathcal L\mathcal S_{\rho,d}=W+\rho q_d,
 \qquad
 \mathcal Lq_d=T_d.                              \tag{73.4}
\]

\begin{theorem}[Exact scalar scaling Ward identities]
\label{thm:v17v73-ward}
For every \(\rho>0\) and \(d\in\mathbb R\),
\[
\boxed{
 \mathbb E(W+\rho q_d)=0,}                       \tag{73.5}
\]
and
\[
\boxed{
 \mathbb ET_d
 =\mathbb E[q_d(W+\rho q_d)].}                   \tag{73.6}
\]
Consequently,
\[
\boxed{
 \mathcal T(\rho,d)-\rho\mathcal V(\rho,d)
 =\operatorname{Cov}_{\nu_{\rho,d}}(q_d,W).}     \tag{73.7}
\]
\end{theorem}

\begin{proof}
Quartic decay justifies integration by parts.  Since
\(\operatorname{div}\mathcal L=0\), for every smooth polynomially
growing \(f\),
\[
 0=\int\mathcal L(fe^{-\mathcal S})
 =\int(\mathcal Lf-f\mathcal L\mathcal S)e^{-\mathcal S}.
                                                               \tag{73.8}
\]
Take \(f=1\) and use the first identity in (73.4) to obtain (73.5).
Take \(f=q_d\) and use both identities in (73.4) to obtain (73.6).
Now subtract \(\rho\operatorname{Var}(q_d)\) from (73.6):
\[
\begin{aligned}
 \mathbb ET_d-\rho\operatorname{Var}(q_d)
 &=\mathbb E(q_dW)+\rho(\mathbb Eq_d)^2\\
 &=\mathbb E(q_dW)-(\mathbb Eq_d)(\mathbb EW),
\end{aligned}
\]
where (73.5) was used in the last line.  This is (73.7).
\end{proof}

Combining (72.10) and (73.7), the exact scalar Higgs free-energy
curvature is
\[
\boxed{
 D_d^2\mathcal F_H[u,u]
 =u^2\rho\operatorname{Cov}_{\nu_{\rho,d}}(q_d,W).}              \tag{73.9}
\]
This identity bypasses the pointwise inverse-Hessian obstruction of
V71 entirely.

\subsection{Strict convexity at zero conformal difference}

At \(d=0\), the law is invariant under \((X,Z)\mapsto(Z,X)\).  Hence
\[
 \mathbb Eq_0=\mathbb EW=0.                       \tag{73.10}
\]
Moreover,
\[
 q_0=X^2-Z^2,qquad T_0=X^2+Z^2,qquad
 W=2q_0T_0.                                      \tag{73.11}
\]

\begin{theorem}[All-coupling central convexity]
\label{thm:v17v73-central}
For every \(\rho>0\),
\[
\boxed{
 D_d^2\mathcal F_H(0)[u,u]
 =2u^2\rho\,
 \mathbb E_{\nu_{\rho,0}}[q_0^2T_0]>0}           \tag{73.12}
\]
whenever \(u\neq0\).
\end{theorem}

\begin{proof}
Use (73.10)--(73.11) in (73.9).  The integrand
\(q_0^2T_0\) is nonnegative and is positive off the union of the two
diagonals and the origin.  That exceptional set has zero measure under
the strictly positive density (72.2).  Thus the expectation is
strictly positive.
\end{proof}

\begin{corollary}[A nonperturbative central interval]
\label{cor:v17v73-interval}
For each fixed \(\rho>0\), there exists
\(d_0(\rho)>0\) such that
\[
 D_d^2\mathcal F_H(d)[u,u]>0
 \quad\text{for }|d|<d_0(\rho),\quad u\neq0.      \tag{73.13}
\]
\end{corollary}

\begin{proof}
On every compact \(d\)-interval, all derivatives of the density are
bounded by a polynomial times a quartic integrable function.  Hence
the partition function and its first two derivatives are continuous
in \(d\).  Apply continuity to the strictly positive value (73.12).
\end{proof}

This interval does not require weak hopping and does not use the
mismatch coefficient \(J\).  It is the first zero-mass environment
convexity statement in the present chain that holds for every positive
dimensionless hopping strength.

\subsection{The exact off-center obstruction}

Equation (73.7) reduces the global scalar problem to the sign of one
covariance.  That sign is not pointwise.  For example, if \(d>0\),
choose squared amplitudes satisfying
\[
 e^{-2d}<\frac{X^2}{Z^2}<1.                       \tag{73.14}
\]
Then \(q_d>0\) while \(W<0\), so \(q_dW<0\).  The central proof cannot
be extended merely by declaring the integrand nonnegative.  A tail
argument or a genuine correlation inequality is necessary.

\begin{assessment}[What V73 closes]
\label{ass:v17v73-status}
V73 derives the exact scaling Ward identity (73.7), replaces the
Brascamp--Lieb target by a single covariance sign, and proves strict
scalar free-energy convexity in a nonzero interval around \(d=0\) for
every hopping strength.  Global convexity is reduced to the off-center
sign of \(\operatorname{Cov}(q_d,W)\).
\end{assessment}

\begin{decision}[V74: large-difference scalar asymptotics]
\label{dec:v17v73-next}
Set \(\varepsilon=e^{-|d|}\), rescale the strongly confined endpoint,
and complete the hopping square exactly.  Expand the scalar partition
function through the first nontrivial even power of \(\varepsilon\)
with a dominated remainder.  Determine the sign of
\(\partial_d^2(-\log\mathcal Z_{\rho,d})\) in both conformal tails.
This will leave, at worst, a compact annulus between the central and
tail convexity regions.
\end{decision}

\begin{firewall}[Claim boundary after V73]
\label{fire:v17v73}
The theorem is a one-edge scalar central-region result.  The conformal
tails, the intervening compact region, the multi-edge Hessian,
environmental susceptibility, remainder-current control, continuum
measure, gauge cofinality, physical Einstein action, and the
Einstein--SM continuum limit remain open.
\end{firewall}

\section*{Version V73 append-only record}
\addcontentsline{toc}{section}{Version V73 append-only record}
The complete V72 source was retained byte for byte before this
continuation.  It had \(608481\) bytes, \(18649\) lines, and SHA--256
\[
 \texttt{23719571CA5ADCD913B1E75A2C0A44B160A02DD6DD42D13712E13AE35A965543}.
\]
V73 proved the scalar scaling Ward identity and all-coupling central
convexity.


\section{V74: strict scalar convexity in both conformal tails}
\label{sec:v17v74}

V73 proves exact scalar convexity near \(d=0\).  We now analyze
\(|d|\to\infty\) directly at the partition-function level.  The
strongly weighted endpoint can be rescaled and the hopping square then
completed exactly.

\subsection{Exact tail representation}

By the exchange symmetry \((X,Z,d)\mapsto(Z,X,-d)\), it suffices to
take \(d>0\).  Put
\[
 \varepsilon=e^{-d},
 \qquad a=\varepsilon^{-1/2},
 \qquad b=\varepsilon^{1/2}.                      \tag{74.1}
\]
Starting from (72.2), first set \(X=\sqrt\varepsilon\,\xi\), then
translate \(\eta=\xi-\sqrt\varepsilon Z\).  The hopping square becomes
\[
 \rho(aX-bZ)^2=\rho\eta^2,                        \tag{74.2}
\]
and
\[
 X^4=\varepsilon^2(\eta+\sqrt\varepsilon Z)^4.   \tag{74.3}
\]
Therefore the unnormalized partition function has the exact form
\[
\boxed{
 \mathcal Z_{\rho,d}
 =C_\rho\varepsilon^{1/2}H_\rho(\varepsilon),}    \tag{74.4}
\]
where
\[
\begin{aligned}
 C_\rho&=
 \int_{\mathbb R^2}e^{-\rho\eta^2-Z^4}\,d\eta\,dZ,\\
 H_\rho(\varepsilon)&=
 \mathbb E_{0,\rho}
 \exp\{-\varepsilon^2(\eta+\sqrt\varepsilon Z)^4\}.             \tag{74.5}
\end{aligned}
\]
Here \(\mathbb E_{0,\rho}\) is expectation under the product law
proportional to \(e^{-\rho\eta^2-Z^4}\).

\subsection{Differentiable asymptotics}

\begin{lemma}[Tail expansion through the first curvature term]
\label{lem:v17v74-expansion}
For every fixed \(\rho>0\), as \(\varepsilon\downarrow0\),
\[
\boxed{
 \log H_\rho(\varepsilon)
 =-\frac{3}{4\rho^2}\varepsilon^2+O_\rho(\varepsilon^3).}         \tag{74.6}
\]
The expansion remains valid after applying
\((\varepsilon\partial_\varepsilon)^j\) for \(j=1,2\) to the
remainder.
\end{lemma}

\begin{proof}
Let
\[
 A_\varepsilon
 =\varepsilon^2(\eta+\sqrt\varepsilon Z)^4\geq0. \tag{74.7}
\]
All mixed moments of \((\eta,Z)\) are finite.  Symmetry makes the odd
terms vanish, and hence
\[
\begin{aligned}
 \mathbb E_{0,\rho}A_\varepsilon
 ={}&\varepsilon^2\mathbb E\eta^4
 +6\varepsilon^3\mathbb E\eta^2\mathbb EZ^2
 +\varepsilon^4\mathbb EZ^4.                    \tag{74.8}
\end{aligned}
\]
The Gaussian factor proportional to \(e^{-\rho\eta^2}\) has variance
\((2\rho)^{-1}\), so
\[
 \mathbb E\eta^4=\frac{3}{4\rho^2}.              \tag{74.9}
\]
For \(A\geq0\),
\[
 |e^{-A}-1+A|\leq\frac12A^2.                     \tag{74.10}
\]
The expectation of \(A_\varepsilon^2\) is
\(O_\rho(\varepsilon^4)\), uniformly for small \(\varepsilon\).
Equations (74.8)--(74.10) therefore give
\[
 H_\rho(\varepsilon)
 =1-\frac{3}{4\rho^2}\varepsilon^2
 +O_\rho(\varepsilon^3).                         \tag{74.11}
\]
Taking the logarithm proves (74.6).

For the differentiated assertion, write
\(t=\sqrt\varepsilon\).  The expectation in (74.5) is unchanged by
\(t\mapsto-t\), after \(\eta\mapsto-\eta\), so its expansion contains
only even powers of \(t\).  Differentiate under the product integral.
Each of the first two \(\varepsilon\partial_\varepsilon\) derivatives
is a finite sum of a polynomial in \(\eta,Z,\sqrt\varepsilon\) times
\(e^{-A_\varepsilon}\).  Product Gaussian--quartic moments dominate
these terms uniformly for small \(\varepsilon\).  Taylor's formula
with integral remainder then preserves the
\(O_\rho(\varepsilon^3)\) order under both derivatives.
\end{proof}

\subsection{Positive tail curvature}

Let
\[
 \mathcal F_\rho(d):=-\log\mathcal Z_{\rho,d}.    \tag{74.12}
\]

\begin{theorem}[Two-sided scalar tail convexity]
\label{thm:v17v74-tail}
For every \(\rho>0\),
\[
\boxed{
 \mathcal F_\rho''(d)
 =\frac{3}{\rho^2}e^{-2|d|}
 +O_\rho(e^{-3|d|})
 \quad (|d|\to\infty).}                          \tag{74.13}
\]
Consequently, there exists \(d_\infty(\rho)<\infty\) such that
\[
 \mathcal F_\rho''(d)>0
 \qquad\text{whenever }|d|>d_\infty(\rho).        \tag{74.14}
\]
\end{theorem}

\begin{proof}
For \(d>0\), (74.4) and (74.6) give
\[
 \mathcal F_\rho(d)
 =\frac d2-\log C_\rho
 +\frac{3}{4\rho^2}e^{-2d}+O_\rho(e^{-3d}).       \tag{74.15}
\]
The differentiated remainder statement in
Lemma~\ref{lem:v17v74-expansion} permits two \(d\)-derivatives and
yields (74.13) on the positive tail.  Exchange of \(X\) and \(Z\)
shows \(\mathcal Z_{\rho,-d}=\mathcal Z_{\rho,d}\), proving the
negative-tail statement.  The leading coefficient is positive, so it
dominates the remainder for sufficiently large \(|d|\).
\end{proof}

Combining (73.9) and (74.13) also gives the exact covariance asymptotic
\[
 \operatorname{Cov}_{\nu_{\rho,d}}(q_d,W)
 =\frac{3}{\rho^3}e^{-2|d|}
 +O_\rho(e^{-3|d|}).                              \tag{74.16}
\]
Thus the off-center Ward covariance eventually has the required sign,
although its positive value tends to zero in the tails.

\begin{corollary}[Reduction to a compact scalar interval]
\label{cor:v17v74-compact}
For each fixed \(\rho>0\), possible failure of convexity of the
massless one-edge scalar Higgs free energy is confined to a compact
set
\[
 d_0(\rho)\leq|d|\leq d_\infty(\rho),             \tag{74.17}
\]
where \(d_0(\rho)\) is supplied by
Corollary~\ref{cor:v17v73-interval}.  If the left endpoint exceeds the
right endpoint, scalar convexity is already global.
\end{corollary}

\begin{proof}
V73 gives strict convexity for \(|d|<d_0(\rho)\), and
Theorem~\ref{thm:v17v74-tail} gives it for
\(|d|>d_\infty(\rho)\).  Only their closed complement can remain.
\end{proof}

The tail result uses the exact covariance, not the pointwise matrix
bound that failed in V71.  It also explains why the scalar curvature
seen numerically becomes small at large conformal difference: the
leading decay is \(e^{-2|d|}\), while its sign remains positive.

\begin{assessment}[What V74 closes]
\label{ass:v17v74-status}
V74 proves strict massless scalar free-energy convexity in both
conformal tails for every hopping strength.  Together with V73, it
reduces the one-edge scalar sign problem to a compact interval.  The
compact interval and the passage from one edge to the full graph are
the remaining environment-curvature acquisitions.
\end{assessment}

\begin{decision}[V75 audit; then compact-interval sign]
\label{dec:v17v74-next}
Perform the mandatory YM/NS audit.  Then use the squared-amplitude
representation of \(\nu_{\rho,d}\) and the Ward covariance (73.7) to
seek a total-positivity or monotone-likelihood proof on the remaining
compact interval.  If an analytic sign theorem is unavailable, derive
a rigorous interval-arithmetic certification whose parameter domain
is stated explicitly; do not extrapolate it to all \(\rho\).
\end{decision}

\begin{firewall}[Claim boundary after V74]
\label{fire:v17v74}
The scalar compact interval is not yet closed, and no multi-edge
environment theorem follows.  Environmental susceptibility,
remainder-current control, continuum measure, gauge cofinality,
physical Einstein action, and the Einstein--SM continuum limit remain
open.
\end{firewall}

\section*{Version V74 append-only record}
\addcontentsline{toc}{section}{Version V74 append-only record}
The complete V73 source was retained byte for byte before this
continuation.  It had \(614776\) bytes, \(18842\) lines, and SHA--256
\[
 \texttt{2B8E837E1A3E91C8D6E5AF80AAED2069C1893C26849B460A7C4AB81248A7CE67}.
\]
V74 proved the exact positive scalar tail asymptotic and reduced the
remaining sign question to a compact interval.


\section{V75: compulsory companion audit at the scalar compact gap}
\label{sec:v17v75}

\subsection{Active companion records}

The highest-numbered active Yang--Mills note is
\begin{center}\small
\path{Renewal_Yang_Mills_Mass_Gap_Research_Notes_V93.tex}.
\end{center}
It has \(717341\) bytes, \(19592\) lines, and SHA--256
\[
 \texttt{5088A71889F0895FD10A54ADDF7DA14F71AB03C7042CA5DB48A675172B50EB89}.
                                                               \tag{75.1}
\]
YM V92 certified the eighth first-coordinate two-step slab, and V93
certified its terminal cap, completing the first quarter face of the
last direct-plus sign cell.  The second and third coordinate
extensions, the complete eighth cell, full-torus positivity, and all
continuum consequences remain open.

The active Navier--Stokes note is still
\begin{center}\small
\path{Renewal_NS_Stability_Research_Notes_V26.tex},
\end{center}
with \(278283\) bytes, \(5319\) lines, and SHA--256
\[
 \texttt{82C887F8C2C69E4F28FAD7C3DC8DE5B9099CB5A4BF431EBA1FFF3E91935978AD}.
                                                               \tag{75.2}
\]

\begin{theorem}[V75 companion-audit decision]
\label{thm:v17v75-decision}
Neither companion note determines the sign of
\[
 \operatorname{Cov}_{\nu_{\rho,d}}(q_d,W)         \tag{75.3}
\]
on the compact scalar interval left by V74, and neither supplies a
multi-edge conformal-environment covariance inequality.  No theorem is
imported at V75.
\end{theorem}

\begin{proof}
The YM conclusions are exact momentum-box positivity statements for a
Wilson pencil.  They contain no scalar Higgs partition function and no
correlation inequality for (75.3).  NS V26 is unchanged and concerns
Oseen-kernel tensor norms.  Therefore neither conclusion implies the
missing scalar sign or the full environment estimate.
\end{proof}

\begin{assessment}[Direction after the audit]
\label{ass:v17v75-direction}
The YM programme is closer to a finite-cutoff full-torus symbol theorem,
but it still cannot enter the SM gauge-environment ledger.  The present
SM route remains the exact scalar Ward covariance.  V76 will expose its
squared-amplitude kernel before deciding whether total positivity is
strong enough to close the compact interval.
\end{assessment}

\begin{decision}[V76: squared-amplitude correlation kernel]
\label{dec:v17v75-next}
Push \(\nu_{\rho,d}\) to \((R,S)=(X^2,Z^2)\), integrate the two sign
choices exactly, and compute the mixed logarithmic derivative of the
resulting \(\cosh(2\rho\sqrt{RS})\) kernel.  Decompose (75.3) into its
four marginal and cross covariances.  Prove the sign if total
positivity gives the required diagonal dominance; otherwise record the
precise inequality still missing and leave the central and tail
theorems intact.
\end{decision}

\begin{firewall}[Claim boundary after V75]
\label{fire:v17v75}
This version is an audit only.  The scalar compact interval, full-graph
environment mixing, environmental susceptibility, remainder-current
control, continuum measure, gauge cofinality, physical Einstein action,
and the Einstein--SM continuum limit remain open.
\end{firewall}

\section*{Version V75 append-only record}
\addcontentsline{toc}{section}{Version V75 append-only record}
The complete V74 source was retained byte for byte before this
continuation.  It had \(622104\) bytes, \(19060\) lines, and SHA--256
\[
 \texttt{5A86A2188B61DA8920690E2595060F94AE565E2E04727A32294B2B921348835D}.
\]
V75 audited YM V93 and NS V26 and retained the scalar compact-interval
acquisition.


\section{V76: squared amplitudes and the conditional-imbalance ledger}
\label{sec:v17v76}

The remaining scalar curvature is the sign of the Ward covariance
\(\operatorname{Cov}(q_d,W)\).  This version pushes the law to squared
amplitudes, proves total positivity of its coupling kernel, and then
separates the covariance into a positive conditional variance and one
explicit one-dimensional remainder.

\subsection{The squared-amplitude law}

Set
\[
 R=X^2,qquad S=Z^2.                               \tag{76.1}
\]
Summing the four sign choices of \((X,Z)\) gives the probability
density on \((0,\infty)^2\)
\[
\boxed{
 p_{\rho,d}(R,S)
 =\frac1{\mathcal N_{\rho,d}}
 (RS)^{-1/2}e^{-R^2-S^2-\rho(a^2R+b^2S)}
 \cosh(2\rho\sqrt{RS}).}                         \tag{76.2}
\]

\begin{proof}
On one sign branch,
\(dX\,dZ=(4\sqrt{RS})^{-1}dR\,dS\).  The cross factor is
\(e^{2\rho\sigma_X\sigma_Z\sqrt{RS}}\).  Its sum over the four sign
pairs is \(4\cosh(2\rho\sqrt{RS})\), proving (76.2) after
normalization.
\end{proof}

Let
\[
 K_\rho(R,S)=\cosh(2\rho\sqrt{RS}),
 \qquad t=2\rho\sqrt{RS}.                         \tag{76.3}
\]

\begin{lemma}[Strict total positivity of the coupling kernel]
\label{lem:v17v76-tp2}
For \(R,S>0\),
\[
\boxed{
 \partial_R\partial_S\log K_\rho(R,S)
 =\frac{t\{\tanh t+t\operatorname{sech}^2t\}}
 {4RS}>0.}                                        \tag{76.4}
\]
Consequently the density (76.2) is multivariate totally positive of
order two.  In particular, increasing functions of \(R\) and of
\(S\) have nonnegative covariance.
\end{lemma}

\begin{proof}
Differentiate
\(\partial_R\log K_\rho=t\tanh(t)/(2R)\) once in \(S\), which gives
(76.4).  All remaining factors in (76.2) are products of a function
of \(R\) and a function of \(S\), so they have zero mixed logarithmic
derivative.  Integrating the resulting four-point inequality over
ordered pairs is the two-variable FKG argument and gives association.
\end{proof}

For reference, if \(d>0\), the squared amplitude at the weakly
penalized endpoint is stochastically larger.

\begin{lemma}[Endpoint stochastic order]
\label{lem:v17v76-order}
If \(d>0\), then
\[
 \mathbb E\varphi(S)\geq\mathbb E\varphi(R)       \tag{76.5}
\]
for every increasing integrable \(\varphi\).  The inequality reverses
when \(d<0\).
\end{lemma}

\begin{proof}
The symmetric factors in (76.2) cancel from the swap ratio, leaving
\[
 \frac{p_{\rho,d}(R,S)}{p_{\rho,d}(S,R)}
 =e^{-2\rho\sinh(d)(R-S)}.                        \tag{76.6}
\]
On \(R>S\), this ratio is below one.  Pair the integrals at
\((R,S)\) and \((S,R)\); the contribution to
\(\mathbb E[\varphi(S)-\varphi(R)]\) is nonnegative on every pair.
\end{proof}

\subsection{What total positivity does and does not prove}

In squared amplitudes,
\[
 q_d=a^2R-b^2S,qquad W=2(R^2-S^2).                \tag{76.7}
\]
Therefore
\[
\begin{aligned}
 \operatorname{Cov}(q_d,W)
 =2\{&a^2\operatorname{Cov}(R,R^2)
      +b^2\operatorname{Cov}(S,S^2)\\
     &-a^2\operatorname{Cov}(R,S^2)
      -b^2\operatorname{Cov}(S,R^2)\}.            \tag{76.8}
\end{aligned}
\]
Lemma~\ref{lem:v17v76-tp2} makes all four covariances on the right
nonnegative.  It does not compare the sum of the two marginal terms
with the sum of the two cross terms.  Thus FKG association alone does
not prove the desired sign.

\subsection{Total amplitude and imbalance}

Introduce
\[
 M=R+S>0,
 \qquad \Delta=R-S\in(-M,M),                     \tag{76.9}
\]
and abbreviate
\[
 c=\cosh d,qquad s=\sinh d.                      \tag{76.10}
\]
Then
\[
 q_d=c\Delta+sM,qquad W=2M\Delta.                \tag{76.11}
\]
At fixed \(M\), the conditional density of \(\Delta\) is proportional
to
\[
\begin{aligned}
 &(M^2-\Delta^2)^{-1/2}
 \exp\left\{-\frac{M^2+\Delta^2}{2}-\rho s\Delta\right\}\\
 &\hspace{24mm}\times
 \cosh\left(\rho\sqrt{M^2-\Delta^2}\right),
 \qquad |\Delta|<M.                              \tag{76.12}
\end{aligned}
\]
The omitted factor depends only on \(M\).  Before the exponential tilt
\(e^{-\rho s\Delta}\), the density in (76.12) is even.

Define
\[
 m_d(M):=\mathbb E[\Delta\mid M],
 \qquad
 v_d(M):=\operatorname{Var}(\Delta\mid M).        \tag{76.13}
\]
For \(d>0\), evenness of the base density and exponential tilting give
\[
 -M<m_d(M)<0,qquad v_d(M)>0.                     \tag{76.14}
\]
Also, with the tilt parameter \(h=\rho s\),
\[
 \partial_hm_d(M)=-v_d(M).                       \tag{76.15}
\]

\begin{theorem}[Exact conditional-imbalance decomposition]
\label{thm:v17v76-decomposition}
For every \(\rho>0\) and \(d\in\mathbb R\),
\[
\boxed{
\begin{aligned}
 \operatorname{Cov}(q_d,W)
 ={}&2c\,\mathbb E[Mv_d(M)]\\
 &+2\operatorname{Cov}
 \bigl(cm_d(M)+sM,\,Mm_d(M)\bigr).
\end{aligned}}                                    \tag{76.16}
\]
\end{theorem}

\begin{proof}
Conditioning on \(M\), the only random variable in (76.11) is
\(\Delta\).  Hence
\[
 \operatorname{Cov}(q_d,W\mid M)
 =2cM\operatorname{Var}(\Delta\mid M).            \tag{76.17}
\]
The conditional means are
\[
 \mathbb E(q_d\mid M)=cm_d(M)+sM,
 \qquad
 \mathbb E(W\mid M)=2Mm_d(M).                    \tag{76.18}
\]
Apply the law of total covariance.
\end{proof}

At \(d=0\), one has \(s=0\) and \(m_0(M)=0\), so the second line of
(76.16) vanishes and the first is strictly positive.  This recovers
Theorem~\ref{thm:v17v73-central} in conditional form.

For \(d\neq0\), the complete scalar sign is now equivalent to the
one-dimensional inequality
\[
\boxed{
 \operatorname{Cov}
 \bigl(cm_d(M)+sM,\,Mm_d(M)\bigr)
 \geq-c\,\mathbb E[Mv_d(M)].}                    \tag{76.19}
\]
Neither total positivity (76.4) nor the endpoint order (76.5) implies
(76.19): both control monotone functions of separate squared
amplitudes, whereas the two functions of \(M\) in (76.19) contain the
tilted conditional mean \(m_d(M)\).

\begin{assessment}[What V76 closes]
\label{ass:v17v76-status}
V76 proves strict total positivity and endpoint stochastic ordering
for the squared-amplitude law.  It then shows exactly why these facts
do not by themselves settle curvature, and reduces the remaining sign
to the conditional-imbalance inequality (76.19).  The central theorem
and the two tail theorems remain intact.
\end{assessment}

\begin{decision}[V77: differentiate the conditional imbalance]
\label{dec:v17v76-next}
Differentiate the even-base conditional law (76.12) with respect to
\(M\).  Obtain exact formulas for \(m_d'(M)\) and
\((Mm_d(M))'\), and test whether their monotonicities make the second
covariance in (76.16) nonnegative or bound it by the first term.  If
the endpoint singularity prevents a direct derivative, use the angle
variable \(\Delta=M\cos\theta\) to remove it before differentiating.
\end{decision}

\begin{firewall}[Claim boundary after V76]
\label{fire:v17v76}
Total positivity has not proved the compact-interval scalar sign.
The full-graph environment gap, environmental susceptibility,
remainder-current control, continuum measure, gauge cofinality,
physical Einstein action, and the Einstein--SM continuum limit remain
open.
\end{firewall}

\section*{Version V76 append-only record}
\addcontentsline{toc}{section}{Version V76 append-only record}
The complete V75 source was retained byte for byte before this
continuation.  It had \(625633\) bytes, \(19148\) lines, and SHA--256
\[
 \texttt{D225C8994BC23610B4B2E8CCE9A231AAF4D8EE64729865BC2F07F12D22D622AD}.
\]
V76 proved squared-amplitude total positivity and isolated the exact
conditional-imbalance inequality.


\section{V77: angle differentiation and the failure of monotone closure}
\label{sec:v17v77}

V76 reduced scalar convexity to the conditional imbalance
\(m_d(M)\).  The substitution \(\Delta=M\cos\theta\) removes the
endpoint singularity in (76.12) and permits exact differentiation in
the total amplitude \(M\).

\subsection{The conditional angle law}

For \(M>0\) and \(0<\theta<\pi\), set
\[
 C(\theta)=\cos\theta,
 \qquad S(\theta)=\sin\theta.                     \tag{77.1}
\]
The conditional density of \(\theta\) at fixed \(M\) is proportional
to
\[
 w_M(\theta)
 =\exp\left\{-\frac{M^2}{2}(1+C^2)-\rho sMC\right\}
 \cosh(\rho MS).                                  \tag{77.2}
\]
Indeed,
\((M^2-\Delta^2)^{-1/2}|d\Delta|=d\theta\).
Write \(\langle\cdot\rangle_M\) for expectation under (77.2) and set
\[
 k(M):=\langle C\rangle_M,
 \qquad
 \varsigma^2(M):=\operatorname{Var}_M(C).         \tag{77.3}
\]
Then
\[
 m_d(M)=Mk(M),
 \qquad
 v_d(M)=M^2\varsigma^2(M).                        \tag{77.4}
\]

\begin{lemma}[Exact total-amplitude derivative]
\label{lem:v17v77-derivative}
Define
\[
 G_M(\theta)
 :=-M(1+C^2)-\rho sC
   +\rho S\tanh(\rho MS).                         \tag{77.5}
\]
Then
\[
\boxed{
 k'(M)=\operatorname{Cov}_M(C,G_M),}              \tag{77.6}
\]
and consequently
\[
\boxed{
\begin{aligned}
 m_d'(M)&=k(M)+M\operatorname{Cov}_M(C,G_M),\\
 (Mm_d(M))'&=2Mk(M)+M^2\operatorname{Cov}_M(C,G_M).
\end{aligned}}                                    \tag{77.7}
\]
\end{lemma}

\begin{proof}
The logarithmic derivative of (77.2) is exactly \(G_M\).  For a
normalized parameter-dependent density,
\[
 \frac d{dM}\langle f\rangle_M
 =\langle\partial_Mf\rangle_M
 +\operatorname{Cov}_M(f,\partial_M\log w_M).     \tag{77.8}
\]
Apply (77.8) to \(f=C\), which is independent of \(M\), and then use
the product rule in (77.4).
\end{proof}

The two functions entering the second covariance of (76.16) are
\[
 A_d(M):=cm_d(M)+sM,
 \qquad
 B_d(M):=Mm_d(M).                                 \tag{77.9}
\]
Lemma~\ref{lem:v17v77-derivative} gives their derivatives without an
endpoint term:
\[
\begin{aligned}
 A_d'(M)&=s+c\{k+M\operatorname{Cov}_M(C,G_M)\},\\
 B_d'(M)&=2Mk+M^2\operatorname{Cov}_M(C,G_M).     \tag{77.10}
\end{aligned}
\]

\subsection{Small-total-amplitude expansion}

\begin{theorem}[Opposite local monotonicities]
\label{thm:v17v77-smallM}
Fix \(\rho>0\) and \(d>0\).  As \(M\downarrow0\),
\[
\begin{aligned}
 k(M)&=-\frac{\rho s}{2}M+O_{\rho,d}(M^3),\\
 m_d(M)&=-\frac{\rho s}{2}M^2+O_{\rho,d}(M^4),\\
 v_d(M)&=\frac12M^2+O_{\rho,d}(M^4),              \tag{77.11}\\
 A_d'(M)&=s-c\rho sM+O_{\rho,d}(M^3),\\
 B_d'(M)&=-\frac{3\rho s}{2}M^2+O_{\rho,d}(M^4).
                                                               \tag{77.12}
\end{aligned}
\]
In particular, \(A_d\) is strictly increasing and \(B_d\) is
strictly decreasing on some interval \((0,M_0(\rho,d))\).
\end{theorem}

\begin{proof}
Uniformly in \(\theta\), expansion of (77.2) gives
\[
 w_M(\theta)=1-\rho sMC+O_{\rho,d}(M^2).          \tag{77.13}
\]
The normalizing integral has no linear term.  Since
\[
 \frac1\pi\int_0^\pi C\,d\theta=0,
 \qquad
 \frac1\pi\int_0^\pi C^2\,d\theta=\frac12,      \tag{77.14}
\]
one obtains \(k(M)=-\rho sM/2+O(M^2)\).  The
coefficient of \(M^2\) vanishes: every second-order term of
\(w_M\) is even in \(C\), and its product with \(C\) integrates to
zero.  Hence the remainder improves to \(O(M^3)\).  The same parity
argument gives
\(\langle C^2\rangle_M=1/2+O(M^2)\).  This proves the first three
lines of (77.11).  Substitute them in (77.9) and differentiate to
obtain (77.12).  Because \(s>0\), the leading derivatives have the
asserted signs for small positive \(M\).
\end{proof}

\begin{countertheorem}[Same-direction Chebyshev closure is unavailable]
\label{ctr:v17v77-chebyshev}
For every \(\rho>0\) and \(d>0\), the functions \(A_d(M)\) and
\(B_d(M)\) are not both nondecreasing and are not both nonincreasing
on \((0,\infty)\).  Therefore the second covariance in (76.16) cannot
be declared nonnegative by a same-monotonicity Chebyshev inequality.
\end{countertheorem}

\begin{proof}
Theorem~\ref{thm:v17v77-smallM} gives
\(A_d'>0\) and \(B_d'<0\) on a common interval next to the origin.
This rules out both possible common monotonicity directions.
\end{proof}

The countertheorem explains the negative second term that can occur in
the law-of-total-covariance decomposition.  Scalar convexity, if
global, follows from a quantitative cancellation between that term and
the positive conditional-variance term
\(2c\mathbb E[Mv_d(M)]\), not from positivity of both terms separately.

\subsection{The remaining one-dimensional inequality}

Using (77.4) and (77.9), (76.19) becomes
\[
\boxed{
 \operatorname{Cov}(A_d(M),B_d(M))
 +c\,\mathbb E[M^3\varsigma^2(M)]\geq0.}          \tag{77.15}
\]
All quantities now concern the single marginal variable \(M\) and
the bounded conditional angle \(C\in[-1,1]\).  Formula (77.6)
provides exact derivatives, while (77.11) shows that any proof must
retain the variance term in (77.15).

\begin{assessment}[What V77 closes]
\label{ass:v17v77-status}
V77 removes the conditional endpoint singularity, computes the exact
total-amplitude derivative of the imbalance, and proves that the
tempting same-monotonicity argument fails for every positive \(d\).
It reduces the compact scalar sign to the quantitative one-dimensional
cancellation (77.15).
\end{assessment}

\begin{decision}[V78: marginal integration by parts in \(M\)]
\label{dec:v17v77-next}
Derive the exact marginal density of \(M\) from (77.2), including its
nonzero value at \(M=0\).  Use an anchored integration-by-parts or
Hardy identity that keeps the boundary term explicit.  Apply it to
\(A_d-\mathbb EA_d\) and \(B_d-\mathbb EB_d\), insert (77.10), and
test whether the negative covariance is bounded by
\(c\mathbb E[M^3\varsigma^2(M)]\).  If the boundary term has the wrong
sign, abandon scalar global convexity as the route to V62 and return to
the required negative-Sobolev environment estimate.
\end{decision}

\begin{firewall}[Claim boundary after V77]
\label{fire:v17v77}
The compact scalar sign is not proved; one monotonicity shortcut is
rigorously excluded.  The full-graph environment gap, environmental
susceptibility, remainder-current control, continuum measure, gauge
cofinality, physical Einstein action, and the Einstein--SM continuum
limit remain open.
\end{firewall}

\section*{Version V77 append-only record}
\addcontentsline{toc}{section}{Version V77 append-only record}
The complete V76 source was retained byte for byte before this
continuation.  It had \(633091\) bytes, \(19385\) lines, and SHA--256
\[
 \texttt{85CCE5445060134DCA2425A89AFF663496149E860DB4D79A9E2DDCD19A768813}.
\]
V77 derived the exact angle-response formula and ruled out monotone
closure of the scalar covariance.


\section{V78: conditional-response cancellation and quartic imbalance}
\label{sec:v17v78}

The marginal integration by parts proposed in V77 is not the shortest
route.  The conditional angle law depends on \(d\) through a single
exponential tilt, and differentiating that tilt makes the two terms in
(76.16) cancel exactly.

\subsection{Joint and marginal total-amplitude laws}

The transformations in V76 give the joint density of
\((M,\theta)\in(0,\infty)\times(0,\pi)\), up to normalization,
\[
\boxed{
 \exp\left\{-\frac{M^2}{2}(1+C^2)
 -\rho M(c+sC)\right\}\cosh(\rho MS)\,dM\,d\theta.}              \tag{78.1}
\]
Hence the marginal density of \(M\) is
\[
 p_{\rho,d}^M(M)
 =\frac{e^{-\rho cM}}{\mathcal Z_{\rho,d}}
 \int_0^\pi w_M(\theta)\,d\theta.                \tag{78.2}
\]
It extends to a positive finite value at \(M=0\).  Nevertheless,
\[
 A_d(0)=B_d(0)=0                                  \tag{78.3}
\]
by (77.11), so integration by parts against either function creates
no lower-end boundary contribution.  The nonzero marginal density is
therefore not the scalar obstruction.

\subsection{Conformal derivative of the conditional imbalance}

At fixed \(M\), all \(d\)-dependence in (76.12) occurs through the
tilt parameter
\[
 h(d)=\rho\sinh d=\rho s.                         \tag{78.4}
\]

\begin{lemma}[Exact conditional conformal response]
\label{lem:v17v78-response}
For every \(M>0\),
\[
\boxed{
 \partial_dm_d(M)=-\rho c\,v_d(M),
 \qquad
 \partial_dB_d(M)=-\rho cM v_d(M).}               \tag{78.5}
\]
\end{lemma}

\begin{proof}
Equation (76.15) gives
\(\partial_hm_d=-v_d\), and \(h'(d)=\rho c\).
The second identity follows from \(B_d(M)=Mm_d(M)\).
\end{proof}

The full density (72.2) has parameter score
\[
 \partial_d\log d\nu_{\rho,d}
 =-\rho\{q_d-\mathbb Eq_d\}.                      \tag{78.6}
\]
Therefore, for a differentiable observable \(F_d\),
\[
 \partial_d\mathbb EF_d
 =\mathbb E(\partial_dF_d)-
 \rho\operatorname{Cov}(F_d,q_d).                \tag{78.7}
\]

\begin{theorem}[Conditional-variance cancellation]
\label{thm:v17v78-cancellation}
The two terms in (76.16) satisfy the exact identity
\[
\boxed{
 \operatorname{Cov}(q_d,W)
 =-\frac2\rho\,\partial_d\mathbb E[B_d(M)].}       \tag{78.8}
\]
Equivalently,
\[
\boxed{
 \operatorname{Cov}(q_d,W)
 =-\frac1\rho\,\partial_d\mathbb E[W].}           \tag{78.9}
\]
\end{theorem}

\begin{proof}
Apply (78.7) to \(F_d=B_d(M)\).  Since \(B_d\) is a function of
\(M\), the law of total covariance gives
\[
 \operatorname{Cov}(B_d,q_d)
 =\operatorname{Cov}(B_d,A_d).                   \tag{78.10}
\]
Use (78.5):
\[
 \partial_d\mathbb EB_d
 =-\rho c\mathbb E[Mv_d(M)]
 -\rho\operatorname{Cov}(B_d,A_d).               \tag{78.11}
\]
Multiplication by \(-2/\rho\) yields exactly the right side of
(76.16), proving (78.8).  Finally,
\(W=2M\Delta\) and conditional expectation give
\(\mathbb EW=2\mathbb EB_d\), which proves (78.9).
\end{proof}

This calculation is also consistent with the first Ward identity:
(73.5) gives
\[
 \rho\mathbb Eq_d=-\mathbb EW.                   \tag{78.12}
\]
Thus \(\partial_d\mathcal F_\rho=\rho\mathbb Eq_d=-\mathbb EW\),
and differentiating reproduces (73.9) and (78.9).

\subsection{The exact monotonicity target}

\begin{corollary}[Quartic-imbalance criterion]
\label{cor:v17v78-criterion}
The massless scalar one-edge free energy is convex at \(d\) if and
only if
\[
\boxed{
 \partial_d\mathbb E_{\nu_{\rho,d}}(R^2-S^2)
 =\partial_d\mathbb E B_d(M)\leq0.}               \tag{78.13}
\]
\end{corollary}

\begin{proof}
Combine (73.9) and (78.8), noting that \(u^2\rho>0\).
\end{proof}

For \(d>0\), Lemma~\ref{lem:v17v76-order} with
\(\varphi(t)=t^2\) proves
\[
 \mathbb E(R^2-S^2)<0.                            \tag{78.14}
\]
Exchange symmetry makes this expectation an odd function of \(d\).
V73 proves that its derivative is negative near zero, and V74 proves
the same in both tails.  What remains is monotonicity on the same
compact interval isolated in (74.17).

The identities (78.8)--(78.13) do not, by themselves, prove that
monotonicity: they are exact reformulations of the free-energy
curvature.  Repeatedly differentiating them would return to the
original covariance and would be circular.

\begin{assessment}[What V78 closes]
\label{ass:v17v78-status}
V78 proves that the conditional variance and marginal covariance in
V76 cancel into a single conformal derivative.  It identifies scalar
convexity with monotone decrease of the quartic endpoint imbalance and
shows that the marginal boundary at \(M=0\) is harmless.  Further Ward
rearrangement is circular; only the compact monotonicity statement
remains in the scalar route.
\end{assessment}

\begin{decision}[V79: go upstream to the required covariance norm]
\label{dec:v17v78-next}
Do not spend another version rewriting (78.13).  Return to the V62
target: a negative-Sobolev covariance inequality for the conformal
environment acting on the derivative (62.4).  Derive the exact
Helffer--Sj\"ostrand/Poisson representation at finite volume under the
bare \(2\kappa_gI\) confinement, and isolate the signed Higgs
free-energy Hessian as a perturbation.  Use the proven scalar central
and tail signs only as local diagnostics, not as a global premise.
\end{decision}

\begin{firewall}[Claim boundary after V78]
\label{fire:v17v78}
Global scalar monotonicity (78.13) is not proved, and no full-graph
environment convexity is claimed.  The scale-compatible environment
covariance norm, environmental susceptibility, remainder-current
control, continuum measure, gauge cofinality, physical Einstein action,
and the Einstein--SM continuum limit remain open.
\end{firewall}

\section*{Version V78 append-only record}
\addcontentsline{toc}{section}{Version V78 append-only record}
The complete V77 source was retained byte for byte before this
continuation.  It had \(639930\) bytes, \(19587\) lines, and SHA--256
\[
 \texttt{AE620D19E57185EBCFD0A37C1190EB85FC6C3E6F66ECFA0D84AB298B3724C9EC}.
\]
V78 proved the conditional-response cancellation and moved the route
upstream from scalar convexity to the needed covariance norm.


\section{V79: exact negative-Sobolev representation for the conformal environment}
\label{sec:v17v79}

The scalar convexity route has reached the compact monotonicity problem
(78.13), while the original V62 target is weaker: a covariance norm
for the environmental response.  This version returns to that target
and derives its exact finite-volume Poisson representation without
assuming pointwise convexity of the effective potential.

\subsection{Finite-volume conformal generator}

Fix the gauge field \(U\), and let
\[
 d\nu_U(r)=Z_U^{-1}e^{-\Phi_{\rm env}(r\mid U)}\,dr              \tag{79.1}
\]
with \(\Phi_{\rm env}\) from (63.11).  The positive quadratic
confinement and the nonnegative mismatch action make this a smooth
finite-volume probability law; the Higgs free energy is smooth by
finite-dimensional differentiation under its confining integral.
Define the nonnegative scalar generator
\[
 \mathscr L_U=-\Delta_r+
 \nabla_r\Phi_{\rm env}\cdot\nabla_r.             \tag{79.2}
\]
For smooth functions of sufficient decay,
\[
 \langle f,\mathscr L_Ug\rangle_{L^2(\nu_U)}
 =\mathbb E_{\nu_U}(\nabla f\cdot\nabla g).        \tag{79.3}
\]
At fixed finite volume the confining potential gives a positive
spectral gap, although no volume- or lattice-uniform lower bound is
asserted here.

On one-forms define the Witten operator
\[
 \mathscr H_U^{(1)}
 :=\mathscr L_U\otimes I+
 \nabla_r^2\Phi_{\rm env}.                        \tag{79.4}
\]
It obeys the commutation identity
\[
 \nabla_r\mathscr L_Uf
 =\mathscr H_U^{(1)}\nabla_rf.                   \tag{79.5}
\]
When the Hessian is not pointwise positive, (79.4) need not be
positive on every arbitrary one-form.  Its inverse below is taken on
the closed exact-gradient subspace generated by (79.5), where the
Poisson representation is well defined.

\begin{theorem}[Finite-volume Helffer--Sj\"ostrand identity]
\label{thm:v17v79-hs}
Let \(F,G\in L^2(\nu_U)\) be centered smooth observables.  Then
\[
\boxed{
 \operatorname{Cov}_{\nu_U}(F,G)
 =\left\langle
 \nabla F,(\mathscr H_U^{(1)})^{-1}\nabla G
 \right\rangle_{L^2(\nu_U)}.}                    \tag{79.6}
\]
In particular,
\[
 \operatorname{Var}_{\nu_U}(F)
 =\left\langle
 \nabla F,(\mathscr H_U^{(1)})^{-1}\nabla F
 \right\rangle_{L^2(\nu_U)}.                     \tag{79.7}
\]
\end{theorem}

\begin{proof}
Solve the centered Poisson equation
\(\mathscr L_U\psi_G=G\).  Finite-volume confinement makes the
inverse well defined on the orthogonal complement of constants.
Integration by parts gives
\[
 \operatorname{Cov}(F,G)
 =\langle F,\mathscr L_U\psi_G\rangle
 =\langle\nabla F,\nabla\psi_G\rangle.            \tag{79.8}
\]
Taking a gradient of the Poisson equation and using (79.5) gives
\[
 \mathscr H_U^{(1)}\nabla\psi_G=\nabla G.         \tag{79.9}
\]
Substitute the exact-gradient inverse of (79.9) in (79.8).
\end{proof}

Equation (79.7) is a negative-Sobolev norm identity.  It remains valid
when the pointwise Hessian has a signed region; a uniform estimate now
requires control of the inverse differential operator rather than
coefficientwise positivity.

\subsection{Insertion of the exact Higgs response}

For a deterministic vertex test vector \(t\), put
\[
 A_t(y)=\sum_xt_x\|y_x\|^2,
 \qquad
 \bar A_t(r,U)=\mathbb E_{\pi_\eta}A_t.           \tag{79.10}
\]
Define the edge response
\[
 \zeta_{t,e}(r,U)
 :=\partial_{d_e}\bar A_t(r,U)
 =-\operatorname{Cov}_{\pi_\eta}(A_t,q_e),        \tag{79.11}
\]
using (62.4).  Since \(d=Br\),
\[
 \nabla_r\bar A_t=B^{\mathsf T}\zeta_t.          \tag{79.12}
\]

\begin{corollary}[Exact conformal conditional-mean variance]
\label{cor:v17v79-response}
For every finite-volume test vector \(t\),
\[
\boxed{
 \operatorname{Var}_{\nu_U}(\bar A_t)
 =\left\langle
 B^{\mathsf T}\zeta_t,
 (\mathscr H_U^{(1)})^{-1}B^{\mathsf T}\zeta_t
 \right\rangle_{L^2(\nu_U)}.}                   \tag{79.13}
\]
\end{corollary}

\begin{proof}
Centering does not change the gradient.  Apply (79.7) to
\(F=\bar A_t-\mathbb E_{\nu_U}\bar A_t\) and use (79.12).
\end{proof}

This is the precise conformal part of the environmental susceptibility
left open in (62.12).  It preserves three structures that were lost in
the pointwise-convexity route: the incidence range
\(B^{\mathsf T}\), the spatial differential operator
\(\mathscr L_U\), and the signed Higgs free-energy Hessian inside
\(\mathscr H_U^{(1)}\).

\subsection{Variational acquisition and comparison with convexity}

On exact one-forms the right side of (79.13) has the variational form
\[
\begin{aligned}
 &\left\langle g,(\mathscr H_U^{(1)})^{-1}g\right\rangle\\
 &\quad=
 \sup_{\omega\in\overline{\operatorname{Ran}\nabla}}
 \left\{2\langle g,\omega\rangle
 -\langle\omega,\mathscr H_U^{(1)}\omega\rangle\right\},
 \qquad g=B^{\mathsf T}\zeta_t.                 \tag{79.14}
\end{aligned}
\]
Thus the scale-compatible acquisition is an upper bound on (79.14)
for the response fields (79.11), not positivity of every edge
coefficient in (63.12).

If the stronger pointwise result
\(\nabla^2\Phi_{\rm env}\succeq2\kappa_gI\) were available, then
\(\mathscr L_U\otimes I\succeq0\) and (79.13) would immediately give
\[
 \operatorname{Var}_{\nu_U}(\bar A_t)
 \leq\frac1{2\kappa_g}
 \mathbb E_{\nu_U}\|B^{\mathsf T}\zeta_t\|^2.   \tag{79.15}
\]
V63--V78 sought precisely this sufficient route.  Formula (79.13)
shows that it is stronger than necessary: signed local curvature can
be tolerated if the differential operator in (79.14) controls the
particular incidence-gradient response.

\begin{theorem}[Exact acquisition replacing global scalar convexity]
\label{thm:v17v79-acquisition}
To close the conformal contribution to the conditional-mean
susceptibility, it is sufficient to prove, uniformly in lattice
spacing, volume, and fixed gauge field, an estimate
\[
\boxed{
 \left\langle
 B^{\mathsf T}\zeta_t,
 (\mathscr H_U^{(1)})^{-1}B^{\mathsf T}\zeta_t
 \right\rangle
 \leq C_{\rm env}\|t\|_2^2}                     \tag{79.16}
\]
for the response class (79.11).  Under (79.16),
\[
 \operatorname{Var}_{\nu_U}(\bar A_t)
 \leq C_{\rm env}\|t\|_2^2.                     \tag{79.17}
\]
\end{theorem}

\begin{proof}
Insert (79.16) into the exact identity (79.13).
\end{proof}

The estimate (79.16) is not proved here.  It is a narrower target than
global pointwise convexity, but it still requires a uniform resolvent
bound on a random interacting environment.  The scalar central and
tail signs from V73--V74 can inform localization of the signed Hessian;
they are not assumed in (79.16).

\begin{assessment}[What V79 closes]
\label{ass:v17v79-status}
V79 derives the exact finite-volume negative-Sobolev representation
for the environmental Higgs response and replaces the stalled global
scalar-convexity route by the response-class resolvent estimate
(79.16).  This is directly aligned with the V62 susceptibility target
and retains the incidence and diffusion structures discarded by the
earlier edgewise bounds.
\end{assessment}

\begin{decision}[V80 audit; then resolvent localization]
\label{dec:v17v79-next}
Perform the mandatory YM/NS audit.  Then decompose the one-form
operator in (79.14) into the bare massive Ornstein--Uhlenbeck part,
the positive mismatch conductance, and the signed Higgs Schur
complement.  Use an IMS localization in conformal-difference space,
placing the V73 central and V74 tail regions in positive cells and
retaining only the compact transition as a form-bounded perturbation.
Do not infer a uniform constant until the localization error and the
volume dependence are explicit.
\end{decision}

\begin{firewall}[Claim boundary after V79]
\label{fire:v17v79}
The finite-volume identity is exact, but the uniform resolvent bound
(79.16) is open.  Gauge-environment covariance, full environmental
susceptibility, remainder-current control, continuum measure, gauge
cofinality, physical Einstein action, and the Einstein--SM continuum
limit remain open.
\end{firewall}

\section*{Version V79 append-only record}
\addcontentsline{toc}{section}{Version V79 append-only record}
The complete V78 source was retained byte for byte before this
continuation.  It had \(646030\) bytes, \(19771\) lines, and SHA--256
\[
 \texttt{1D6496AFAF78EC273700CB0666B024ED6572834DF2062E065BA7244FA46AD93F}.
\]
V79 derived the exact response-class negative-Sobolev representation
and its uniform resolvent acquisition.


\section{V80: compulsory companion audit before resolvent localization}
\label{sec:v17v80}

\subsection{Active companion checkpoints}

The newest active Yang--Mills note is
\begin{center}\small
\path{Renewal_Yang_Mills_Mass_Gap_Research_Notes_V94.tex}.
\end{center}
It has \(728997\) bytes, \(19914\) lines, and SHA--256
\[
 \texttt{645ABB0A97779B963A4E82819ADE7FF7DF1ABFC154FDF5703013C9A972D65899}.
                                                               \tag{80.1}
\]
YM V94 certifies three new charts and extends the last direct-plus
atlas to second-coordinate depth \(2t\), uniformly across its completed
first-coordinate quarter.  Its second-coordinate cap, third-coordinate
extension, full torus, and continuum mass-gap steps remain open.

The newest active Navier--Stokes note remains
\begin{center}\small
\path{Renewal_NS_Stability_Research_Notes_V26.tex},
\end{center}
with \(278283\) bytes, \(5319\) lines, and SHA--256
\[
 \texttt{82C887F8C2C69E4F28FAD7C3DC8DE5B9099CB5A4BF431EBA1FFF3E91935978AD}.
                                                               \tag{80.2}
\]

\begin{theorem}[V80 companion-audit decision]
\label{thm:v17v80-decision}
Neither active companion note supplies a form bound for the Witten
operator \(\mathscr H_U^{(1)}\) or the response-class resolvent
(79.16).  No theorem is imported at V80.
\end{theorem}

\begin{proof}
YM V94 is a finite-cutoff Wilson-symbol covering result, not a Gibbs
one-form resolvent estimate.  Its remaining atlas steps also prevent a
full-torus symbol theorem from being used as a later gauge input.  NS
V26 concerns Oseen-kernel norms.  Neither statement contains the
conformal generator, the Higgs Schur complement, or the response field
\(B^{\mathsf T}\zeta_t\).
\end{proof}

\begin{assessment}[Direction after the audit]
\label{ass:v17v80-direction}
The companion programmes do not change the V79 route.  V81 will derive
the exact IMS formula for the one-form operator and quantify the volume
loss of a naive partition that localizes every conformal edge
independently.
\end{assessment}

\begin{decision}[V81: response-adapted IMS localization]
\label{dec:v17v80-next}
Decompose \(\mathscr H_U^{(1)}\) into diffusion, bare mass, mismatch
conductance, and the signed Higgs Schur complement.  Prove the IMS
identity for a smooth partition of conformal configuration space.
Then compute the localization error for an edge-product partition and
determine whether it stays uniform in volume; if it does not, retain
the response-class incidence structure and replace the product cutoff
by a local or multiscale partition.
\end{decision}

\begin{firewall}[Claim boundary after V80]
\label{fire:v17v80}
This audit proves no resolvent estimate.  Uniform conformal and gauge
environment covariance, environmental susceptibility,
remainder-current control, continuum measure, gauge cofinality,
physical Einstein action, and the Einstein--SM continuum limit remain
open.
\end{firewall}

\section*{Version V80 append-only record}
\addcontentsline{toc}{section}{Version V80 append-only record}
The complete V79 source was retained byte for byte before this
continuation.  It had \(654370\) bytes, \(20007\) lines, and SHA--256
\[
 \texttt{66676F3588C5C3D6B07C26019FCF350314BF870CC312D669CF2ADE9ECC17AAC8}.
\]
V80 audited YM V94 and NS V26 and retained the resolvent-localization
acquisition.


\section{V81: IMS localization and the extensive product-cutoff error}
\label{sec:v17v81}

V79 identifies the response-class inverse of the conformal one-form
operator as the required object.  This version decomposes its quadratic
form, proves an exact IMS identity, and tests the naive strategy of
localizing every conformal edge independently.

\subsection{The one-form quadratic form}

Let
\[
 \mathsf H_H(d,U):=\nabla_d^2\mathcal F_H(d,U)     \tag{81.1}
\]
be the full edge-space Higgs free-energy Hessian.  It is generally not
diagonal: by (63.3), its negative part is the covariance matrix of the
edge currents.  Define the positive mismatch matrix
\[
 \mathsf W_J(d)
 :=\operatorname{diag}_e\{4J\cosh(2d_e)\}.        \tag{81.2}
\]
Then
\[
\boxed{
 \nabla_r^2\Phi_{\rm env}
 =2\kappa_gI+B^{\mathsf T}
   \{\mathsf W_J+\mathsf H_H\}B.}                 \tag{81.3}
\]
Consequently,
\[
 \mathscr H_U^{(1)}
 =\mathscr L_U\otimes I+2\kappa_gI
 +B^{\mathsf T}\mathsf W_JB
 +B^{\mathsf T}\mathsf H_HB.                    \tag{81.4}
\]

For a smooth vertex one-form \(\omega(r)\), integration by parts gives
the closed quadratic form
\[
\begin{aligned}
 \mathcal Q_U[\omega]
 =\mathbb E_{\nu_U}\{&\|\nabla_r\omega\|_{\rm HS}^2
 +2\kappa_g\|\omega\|^2\\
 &+(B\omega)^{\mathsf T}\mathsf W_J(B\omega)
 +(B\omega)^{\mathsf T}\mathsf H_H(B\omega)\}.
                                                               \tag{81.5}
\end{aligned}
\]
The last term may be signed; all preceding terms are nonnegative.

\subsection{Exact IMS identity}

Let \(\{\eta_\alpha\}_{\alpha\in\mathcal A}\) be a finite smooth
partition satisfying
\[
 \sum_\alpha\eta_\alpha(r)^2=1.                  \tag{81.6}
\]
Set
\[
 \Lambda_\eta(r):=
 \sum_\alpha\|\nabla_r\eta_\alpha(r)\|^2.        \tag{81.7}
\]

\begin{theorem}[One-form IMS formula]
\label{thm:v17v81-ims}
For every form-domain one-form \(\omega\),
\[
\boxed{
 \mathcal Q_U[\omega]
 =\sum_\alpha\mathcal Q_U[\eta_\alpha\omega]
 -\mathbb E_{\nu_U}[\Lambda_\eta\|\omega\|^2].} \tag{81.8}
\]
\end{theorem}

\begin{proof}
For each component of \(\omega\),
\[
 \sum_\alpha|\nabla(\eta_\alpha\omega)|^2
 =|\nabla\omega|^2+
 \left(\sum_\alpha|\nabla\eta_\alpha|^2\right)|\omega|^2,       \tag{81.9}
\]
because \(\sum_\alpha\eta_\alpha\nabla\eta_\alpha
=\frac12\nabla\sum_\alpha\eta_\alpha^2=0\).  Every multiplication
term in (81.5), including the signed Higgs Hessian, partitions exactly
by (81.6).  Integrate (81.9) and rearrange.
\end{proof}

\begin{corollary}[Local forms imply a global resolvent bound]
\label{cor:v17v81-local}
Suppose
\[
 \mathcal Q_U[\eta_\alpha\omega]
 \geq m_\alpha\|\eta_\alpha\omega\|_{L^2(\nu_U)}^2              \tag{81.10}
\]
for every \(\alpha\), and put
\(m_*:=\inf_\alpha m_\alpha\).  If
\[
 m_* > \|\Lambda_\eta\|_\infty,                 \tag{81.11}
\]
then
\[
 \mathscr H_U^{(1)}
 \succeq(m_*-\|\Lambda_\eta\|_\infty)I          \tag{81.12}
\]
on the exact-gradient subspace, and
\[
 \langle g,(\mathscr H_U^{(1)})^{-1}g\rangle
 \leq
 \frac{\|g\|_2^2}{m_*-\|\Lambda_\eta\|_\infty}.                 \tag{81.13}
\]
\end{corollary}

\begin{proof}
Insert (81.10) into (81.8), use (81.6), and bound the IMS error by its
essential supremum.  The inverse estimate follows from the spectral
theorem.
\end{proof}

This gives a valid route from central, transition, and tail form bounds
to (79.16), provided the IMS error is uniform in volume.

\subsection{The edge-product partition test}

Let \(\chi_0,\chi_1\) be smooth functions on \(\mathbb R\) with
\[
 \chi_0^2+\chi_1^2=1,
 \qquad
 \sup_s\sum_{j=0}^1|\chi_j'(s)|^2\leq\ell^2.      \tag{81.14}
\]
For every binary edge label
\(\sigma\in\{0,1\}^{E}\), define
\[
 \eta_\sigma(r)=\prod_{e\in E}\chi_{\sigma_e}(d_e),
 \qquad d=Br.                                     \tag{81.15}
\]
These functions form a partition of the type (81.6).

\begin{theorem}[Exact product-partition error]
\label{thm:v17v81-product}
For the incidence matrix of a graph without self-loops,
\[
\boxed{
 \Lambda_{\rm prod}(r)
 =2\sum_{e\in E}\sum_{j=0}^1|\chi_j'(d_e)|^2
 \leq2|E|\ell^2.}                                \tag{81.16}
\]
\end{theorem}

\begin{proof}
Since \(d=Br\),
\(\nabla_r\eta_\sigma=B^{\mathsf T}\nabla_d\eta_\sigma\).  Expand
the squared norm and sum over \(\sigma\).  For distinct edges
\(e\neq f\), the cross term contains
\[
 \sum_{j=0}^1\chi_j(d_e)\chi_j'(d_e)
 =\frac12\frac d{dd_e}\sum_j\chi_j(d_e)^2=0.     \tag{81.17}
\]
All other factors sum to one.  The diagonal term for edge \(e\) is
\((BB^{\mathsf T})_{ee}\sum_j|\chi_j'(d_e)|^2\).
Each incidence row has one \(+1\) and one \(-1\), so
\((BB^{\mathsf T})_{ee}=2\).  This proves (81.16).
\end{proof}

\begin{countertheorem}[Fixed-width product IMS is not volume uniform]
\label{ctr:v17v81-volume}
Assume the one-dimensional cutoff is nonconstant and fixed with the
volume.  On any graph sequence containing matchings
\(\mathcal M_V\subset E\) with \(|\mathcal M_V|\to\infty\),
\[
 \|\Lambda_{\rm prod}\|_\infty\longrightarrow\infty.            \tag{81.18}
\]
Consequently, even if every localized form in (81.10) had a common
volume-independent lower bound \(m_*>0\), the product partition could
not satisfy (81.11) uniformly.
\end{countertheorem}

\begin{proof}
Because the cutoff is nonconstant, there is an \(s_*\) for which
\[
 c_*:=\sum_{j=0}^1|\chi_j'(s_*)|^2>0.             \tag{81.19}
\]
The edges of a matching have disjoint endpoints, so one may choose the
vertex values of \(r\) to impose \(d_e=s_*\) simultaneously for every
\(e\in\mathcal M_V\).  Other edge differences are irrelevant because
all summands in (81.16) are nonnegative.  Hence
\[
 \|\Lambda_{\rm prod}\|_\infty
 \geq2c_*|\mathcal M_V|,                          \tag{81.20}
\]
which proves (81.18) and eventually violates (81.11).  This proves the
failure of the stated localization estimate, not failure of the true
resolvent bound.
\end{proof}

Choosing a transition width growing like \(|E|^{1/2}\) would bound the
IMS error, but it would no longer isolate the fixed compact transition
in each edge variable.  The product partition therefore cannot combine
the scalar central and tail information into a thermodynamic estimate.

\begin{assessment}[What V81 closes]
\label{ass:v17v81-status}
V81 derives the exact one-form decomposition and IMS formula and proves
that naive independent localization of every conformal edge incurs an
extensive error.  The loss occurs before estimating the signed Higgs
Schur complement.  A scale- and volume-safe proof must localize only
the response-relevant directions or use a local/multiscale resolvent
method.
\end{assessment}

\begin{decision}[V82: local response blocks instead of global cells]
\label{dec:v17v81-next}
Use the locality of \(\zeta_{t,e}\) from (79.11) and decompose the
resolvent through overlapping bounded-diameter edge blocks.  Prove a
geometric resolvent identity with a bounded-overlap IMS partition so
that the error is controlled per block rather than by \(|E|\).  Track
the commutator across block boundaries and determine which off-block
decay estimate is still required from the environment generator.
\end{decision}

\begin{firewall}[Claim boundary after V81]
\label{fire:v17v81}
The response resolvent bound (79.16) is not proved; V81 rules out one
global localization scheme.  Uniform conformal and gauge environment
covariance, environmental susceptibility, remainder-current control,
continuum measure, gauge cofinality, physical Einstein action, and the
Einstein--SM continuum limit remain open.
\end{firewall}

\section*{Version V81 append-only record}
\addcontentsline{toc}{section}{Version V81 append-only record}
The complete V80 source was retained byte for byte before this
continuation.  It had \(657729\) bytes, \(20091\) lines, and SHA--256
\[
 \texttt{6E4F1DAC88B72BEA7DD6C42188A6C344519DC0DB9DEC80151D68522C77CA0CC3}.
\]
V81 proved the one-form IMS identity and the extensive product-cutoff
obstruction.


\section{V82: exponential locality of the conditional Higgs response}
\label{sec:v17v82}

The product IMS failure in V81 concerns localization of the whole
configuration space.  The source in (79.13) is much narrower.  Under
the V61 influence margin, the Neumann series that propagates Higgs
gradients also yields an exponentially decaying spatial kernel for
this response.

\subsection{The covariance form of the sweeping estimate}

For a smooth Higgs observable \(f\), define its site gradient vector
\[
 a_x(f):=
 \left(\mathbb E_{\pi_\eta}
 \|\nabla_{y_x}f\|^2\right)^{1/2}.                \tag{82.1}
\]
Let \(M\) be the nonnegative symmetric influence matrix (61.6), and
put
\[
 \mathsf R_H:=(I-M)^{-1}=\sum_{j=0}^\infty M^j.   \tag{82.2}
\]

\begin{theorem}[Spatial covariance form of V61 sweeping]
\label{thm:v17v82-covariance}
Under \(r_H=\|M\|_{2\to2}<1\), for smooth Higgs observables \(f,g\),
\[
\boxed{
 |\operatorname{Cov}_{\pi_\eta}(f,g)|
 \leq C_0
 \left\langle\mathsf R_Ha(f),
                    \mathsf R_Ha(g)\right\rangle_{\ell^2}.}     \tag{82.3}
\]
Equivalently, because \(M\) is symmetric,
\[
 |\operatorname{Cov}(f,g)|
 \leq C_0\langle a(f),\mathsf R_H^2a(g)\rangle.   \tag{82.4}
\]
\end{theorem}

\begin{proof}
Use the same ordered-site martingale decomposition as in the proof of
Theorem~\ref{thm:v17v61-tensor}, now bilinearly for \(f\) and \(g\).
Each martingale covariance increment is bounded by the one-site
Poincar\'e constant times the product of its two conditional gradient
norms.  Whenever a conditional expectation is swept across a site,
Lemma~\ref{lem:v17v61-sweep} transfers the gradient by one factor of
\(M\).  Summing every possible transfer path turns the two accumulated
gradient vectors into \(\mathsf R_Ha(f)\) and
\(\mathsf R_Ha(g)\).  Cauchy--Schwarz over martingale increments gives
(82.3).  Symmetry of \(M\) gives (82.4).  Taking \(f=g\) and
\(\|\mathsf R_H\|\leq(1-r_H)^{-1}\) recovers (61.10), so no stronger
single-site input has been introduced.
\end{proof}

The squared resolvent has the path expansion
\[
 \mathsf R_H^2
 =\sum_{n=0}^\infty(n+1)M^n.                     \tag{82.5}
\]
Because \(M\) is nearest-neighbor, \((M^n)_{xz}=0\) when
\(n<\operatorname{dist}(x,z)\).

\begin{lemma}[Off-diagonal Neumann tail]
\label{lem:v17v82-tail}
For \(L\geq0\), define
\[
 \tau_L(r)
 :=\sum_{n=L}^\infty(n+1)r^n
 =\frac{r^L\{(L+1)-Lr\}}{(1-r)^2}.               \tag{82.6}
\]
If vertex sets \(S,T\) have graph distance at least \(L\), then
\[
 \|1_S\mathsf R_H^2 1_T\|_{2\to2}
 \leq\tau_L(r_H).                                 \tag{82.7}
\]
\end{lemma}

\begin{proof}
In (82.5), all terms with \(n<L\) vanish between the two sets.
Bound the remaining powers by \(\|M^n\|\leq r_H^n\) and sum the
geometric derivative series, giving (82.6).
\end{proof}

\subsection{Application to the conformal response field}

For \(A_t\) in (79.10), the conditional site-moment ceiling (67.8)
gives
\[
 a_x(A_t)
 =2|t_x|\{mathbb E\|y_x\|^2\}^{1/2}
 \leq2\sqrt{q_HC_{\rm H,glob}}\,|t_x|.            \tag{82.8}
\]
For the edge current \(q_e\), (63.7) with a single edge gives
\[
 a_x(q_e)
 \leq
 \begin{cases}
 2\sqrt{q_HC_{\rm H,glob}}\,g_{x,e},&x\in e,\\
 0,&x\notin e.
 \end{cases}                                      \tag{82.9}
\]

\begin{theorem}[Exponential response locality]
\label{thm:v17v82-response}
Let \(S=\operatorname{supp}t\), and suppose the endpoints of edge
\(e\) have distance at least \(L\) from \(S\).  Then
\[
\boxed{
 |\zeta_{t,e}|
 \leq4q_HC_0C_{\rm H,glob}\,
 \tau_L(r_H)\,\|t\|_2
 \{g_{e_-,e}^2+g_{e_+,e}^2\}^{1/2}.}             \tag{82.10}
\]
Equivalently,
\[
 |\zeta_{t,e}|
 \leq4\sqrt2q_HC_0C_{\rm H,glob}\chi_e
 \tau_L(r_H)\sqrt{\cosh(2d_e)}\,\|t\|_2.         \tag{82.11}
\]
\end{theorem}

\begin{proof}
By (79.11),
\(|\zeta_{t,e}|=|\operatorname{Cov}(A_t,q_e)|\).
Insert (82.8)--(82.9) in (82.4).  The vector in (82.9) is supported on
the two endpoints of \(e\); Lemma~\ref{lem:v17v82-tail} bounds the
operator between that set and \(S\) by \(\tau_L(r_H)\).  The product
of the two gradient constants is
\(4q_HC_{\rm H,glob}\), and (82.4) contributes \(C_0\), proving
(82.10).  Finally,
\[
 g_{e_-,e}^2+g_{e_+,e}^2
 =2\chi_e^2\cosh(2d_e),                           \tag{82.12}
\]
which gives (82.11).
\end{proof}

For fixed \(r_H<1\),
\(\tau_L(r_H)=O(Lr_H^L)\).  Thus the Higgs response to a local test
has a volume-uniform exponential spatial envelope, with the exact
local conformal weight shown in (82.11).  This is stronger locality
information than the structure-factor bound (61.13).

\subsection{What is still needed for the environment resolvent}

The source in (79.13) is \(B^{\mathsf T}\zeta_t\).  Theorem
\ref{thm:v17v82-response} shows that its spatial tail can be truncated
at radius \(L\) with an error decaying like \(Lr_H^L\), before solving
the conformal Poisson equation.  To pass that decay through
\((\mathscr H_U^{(1)})^{-1}\), however, one still needs an off-block
resolvent estimate for the environment operator.  Locality of the
source alone does not imply locality of the solution.

The conformal factor in (82.11) is matched to the mismatch conductance:
\[
 \frac{|\zeta_{t,e}|^2}{4J\cosh(2d_e)}
 \leq
 \frac{8q_H^2C_0^2C_{\rm H,glob}^2\chi_e^2}{J}
 \tau_L(r_H)^2\|t\|_2^2.                         \tag{82.13}
\]
This cancellation removes the unbounded \(d_e\)-weight, but the
factor \(J^{-1}\) is not canonical-scale safe.  The bare vertex mass
and the diffusion part of the one-form operator must participate in
the resolvent estimate.

\begin{assessment}[What V82 closes]
\label{ass:v17v82-status}
V82 upgrades the V61 sweep to a spatial covariance kernel and proves a
volume-uniform exponential envelope for the conditional Higgs response
away from the support of its test.  It supplies the locality needed for
bounded response blocks and identifies the remaining acquisition as
off-block decay of the conformal one-form resolvent itself.
\end{assessment}

\begin{decision}[V83: geometric resolvent across a block boundary]
\label{dec:v17v82-next}
Let \(\Pi_L\) project vertex one-forms onto an \(L\)-neighborhood of
\(\operatorname{supp}t\).  Derive the exact block Schur-complement
formula for \((\mathscr H_U^{(1)})^{-1}\) relative to
\(\Pi_L\oplus(1-\Pi_L)\).  Bound the coupling block by commutators of
the diffusion and incidence terms, and state the minimal exterior
coercivity needed for exponential resolvent decay.  Keep the signed
Higgs Hessian explicit.
\end{decision}

\begin{firewall}[Claim boundary after V82]
\label{fire:v17v82}
Response locality does not prove environment-resolvent locality or
(79.16).  Uniform conformal and gauge environmental susceptibility,
remainder-current control, continuum measure, gauge cofinality,
physical Einstein action, and the Einstein--SM continuum limit remain
open.
\end{firewall}

\section*{Version V82 append-only record}
\addcontentsline{toc}{section}{Version V82 append-only record}
The complete V81 source was retained byte for byte before this
continuation.  It had \(665651\) bytes, \(20330\) lines, and SHA--256
\[
 \texttt{CFC9034CB7837EC51995832A3D6EDC3F640FB351763C3DD58E530B9B193BE307}.
\]
V82 proved exponential spatial locality of the conditional Higgs
response under the V61 influence margin.


\section{V83: block Schur complement for the conformal one-form resolvent}
\label{sec:v17v83}

V82 proves locality of the source in the Higgs variables.  To pass
that locality through the conformal resolvent, this version derives the
exact spatial block inverse of \(\mathscr H_U^{(1)}\) and identifies
both contributions to its off-block coupling.

\subsection{Exact spatial block inverse}

Let \(P\) be the orthogonal projection of vertex one-forms onto a
chosen spatial block and put \(Q=I-P\).  These projections act only on
the vertex component and are independent of \(r\).  We use the
Friedrichs realization of \(\mathscr H_U^{(1)}\) on the full one-form
space.  Its restriction to exact gradients is the operator entering
(79.6), while the full space is stable under \(P\) and \(Q\).  Write
\[
 \mathscr H_U^{(1)}=
 \begin{pmatrix}
  \mathsf A&\mathsf C^*\\
  \mathsf C&\mathsf D
 \end{pmatrix},
 \quad
 \mathsf A=P\mathscr H_U^{(1)}P,
 \quad
 \mathsf D=Q\mathscr H_U^{(1)}Q,
 \quad
 \mathsf C=Q\mathscr H_U^{(1)}P.                 \tag{83.1}
\]

\begin{theorem}[Schur-complement resolvent identity]
\label{thm:v17v83-schur}
Assume \(\mathsf D\succeq m_{\rm out}Q\) for some
\(m_{\rm out}>0\), and define
\[
 \mathsf S_P
 :=\mathsf A-\mathsf C^*\mathsf D^{-1}\mathsf C. \tag{83.2}
\]
For \(g=g_P+g_Q\) in the exact-gradient subspace,
\[
\boxed{
\begin{aligned}
 \langle g,(\mathscr H_U^{(1)})^{-1}g\rangle
 ={}&\left\langle
 g_P-\mathsf C^*\mathsf D^{-1}g_Q,
 \mathsf S_P^{-1}
 (g_P-\mathsf C^*\mathsf D^{-1}g_Q)
 \right\rangle\\
 &+\langle g_Q,\mathsf D^{-1}g_Q\rangle .
\end{aligned}}                                    \tag{83.3}
\]
If also \(\mathsf S_P\succeq m_{\rm in}P\), then
\[
\boxed{
 \langle g,(\mathscr H_U^{(1)})^{-1}g\rangle
 \leq
 \frac{(\|g_P\|+\|\mathsf C\|m_{\rm out}^{-1}
              \|g_Q\|)^2}{m_{\rm in}}
 +\frac{\|g_Q\|^2}{m_{\rm out}}.}                \tag{83.4}
\]
\end{theorem}

\begin{proof}
Gaussian elimination of the lower block in (83.1) gives
\[
 \begin{pmatrix}I&-\mathsf C^*\mathsf D^{-1}\\0&I\end{pmatrix}
 \begin{pmatrix}\mathsf A&\mathsf C^*\\
                 \mathsf C&\mathsf D\end{pmatrix}
 \begin{pmatrix}I&0\\-\mathsf D^{-1}\mathsf C&I\end{pmatrix}
 =\begin{pmatrix}\mathsf S_P&0\\0&\mathsf D\end{pmatrix}.     \tag{83.5}
\]
Invert the factors and evaluate the resulting quadratic form, which
is (83.3).  The two spectral lower bounds and
\(\|\mathsf D^{-1}\|\leq m_{\rm out}^{-1}\) give (83.4).
\end{proof}

For a source supported in the block, \(g_Q=0\), and (83.3) reduces to
\[
 P(\mathscr H_U^{(1)})^{-1}P=\mathsf S_P^{-1}.    \tag{83.6}
\]
Thus exterior coercivity alone is insufficient: the boundary feedback
\(\mathsf C^*\mathsf D^{-1}\mathsf C\) must leave the inner Schur
complement positive.

\subsection{The exact off-block coupling}

Using (81.4) and the fact that \(P\) is independent of \(r\),
\[
 [P,\mathscr L_U\otimes I]=[P,2\kappa_gI]=0.      \tag{83.7}
\]
Let
\[
 \mathsf C_H(d,U)
 :=\bigl(\operatorname{Cov}_{\pi_\eta}(q_e,q_f)\bigr)_{e,f}.     \tag{83.8}
\]
The exact Hessian identity (63.3) gives
\[
 \mathsf H_H=operatorname{diag}_e(\mathbb ET_e)-\mathsf C_H.   \tag{83.9}
\]
Consequently,
\[
\boxed{
 \mathsf C
 =QB^{\mathsf T}
 \left\{\mathsf W_J+\operatorname{diag}(\mathbb ET_e)
       -\mathsf C_H\right\}BP.}                  \tag{83.10}
\]
The first two edge matrices in braces are diagonal.  They couple the
two vertex blocks only through edges crossing their spatial boundary.
All genuinely long-range block coupling is carried by the conditional
Higgs current covariance \(\mathsf C_H\).

\subsection{Decay of the normalized Higgs-current covariance}

Put
\[
 w_e(d):=g_{e_-,e}^2+g_{e_+,e}^2
 =2\chi_e^2\cosh(2d_e).                           \tag{83.11}
\]
Let \(\operatorname{dist}(e,f)\) be the minimum graph distance
between their endpoint sets.

\begin{theorem}[Off-diagonal current-covariance envelope]
\label{thm:v17v83-current-decay}
Under the V61 influence margin, for all edges \(e,f\),
\[
\boxed{
 |(\mathsf C_H)_{ef}|
 \leq4q_HC_0C_{\rm H,glob}\,
 \tau_{\operatorname{dist}(e,f)}(r_H)
 \sqrt{w_e(d)w_f(d)}.}                            \tag{83.12}
\]
\end{theorem}

\begin{proof}
Apply the covariance form (82.4) to \(q_e,q_f\).  By (82.9), their
site-gradient vectors are supported on their respective endpoint sets
and have \(\ell^2\) norms at most
\(2\sqrt{q_HC_{\rm H,glob}}\sqrt{w_e}\) and
\(2\sqrt{q_HC_{\rm H,glob}}\sqrt{w_f}\).  The off-diagonal Neumann
tail (82.7) supplies \(\tau_{\operatorname{dist}(e,f)}(r_H)\).
Multiplication by the covariance prefactor \(C_0\) proves (83.12).
\end{proof}

For a graph of maximum degree \(D\), the number of edges at edge
distance \(n\) from a fixed edge is at most
\[
 N_n\leq C_D(D-1)^n                              \tag{83.13}
\]
with a degree-dependent constant \(C_D\).  Define the normalized
matrix
\[
 \widehat{\mathsf C}_H
 :=\operatorname{diag}(w)^{-1/2}
   \mathsf C_H\operatorname{diag}(w)^{-1/2}.      \tag{83.14}
\]

\begin{corollary}[A summable normalized off-block tail]
\label{cor:v17v83-summable}
If the strengthened spatial margin
\[
 \alpha_H:=(D-1)r_H<1,                            \tag{83.15}
\]
holds, then the operator norm of the part of
\(\widehat{\mathsf C}_H\) connecting edge sets at distance at least
\(L\) is at most
\[
\boxed{
 \Xi_L
 :=\frac{4q_HC_0C_{\rm H,glob}C_D}{(1-r_H)^2}
 \frac{\alpha_H^L\{(L+1)-L\alpha_H\}}
 {(1-\alpha_H)^2}.}                               \tag{83.16}
\]
In particular, \(\Xi_L=O(L\alpha_H^L)\).
\end{corollary}

\begin{proof}
Use (83.12), the bound
\[
 \tau_n(r_H)
 \leq\frac{(n+1)r_H^n}{(1-r_H)^2},               \tag{83.17}
\]
and (83.13).  The row and column sums beyond distance \(L\) are
bounded by the series in (83.16).  Schur's test gives the operator
norm bound.
\end{proof}

The strengthened condition (83.15) is not part of V61 and has not
been verified for physical running couplings.  It is an explicit
high-temperature branch under which the nonlocal part of (83.10)
decays in the natural conformal weight.

\subsection{The remaining exterior-coercivity acquisition}

Theorems~\ref{thm:v17v83-schur} and
\ref{thm:v17v83-current-decay} separate the V79 resolvent problem into:
\begin{enumerate}
\item a local boundary conductance from
      \(\mathsf W_J+\operatorname{diag}(\mathbb ET_e)\);
\item an exponentially small distant covariance coupling, under
      (83.15);
\item the lower bounds \(\mathsf D\succeq m_{\rm out}\) and
      \(\mathsf S_P\succeq m_{\rm in}\).
\end{enumerate}
The first two items now have exact formulas.  The third is not a
consequence of source locality and remains the true environment
coercivity input.  In particular, the bare term \(2\kappa_gI\) cannot
simply be cited before controlling the signed local part of
\(B^{\mathsf T}\mathsf H_HB\).

\begin{assessment}[What V83 closes]
\label{ass:v17v83-status}
V83 derives the exact spatial Schur-complement resolvent and identifies
the local and nonlocal off-block couplings.  It proves exponential
decay of the normalized conditional current covariance under the
explicit strengthened margin (83.15).  The unresolved step is a
volume-uniform exterior and inner-Schur coercivity bound.
\end{assessment}

\begin{decision}[V84: local Schur coercivity versus the bare mass]
\label{dec:v17v83-next}
Use the variational form (81.5) on a bounded spatial block.  Bound the
negative current-covariance matrix by (83.12), retain the positive
diagonal \(\mathbb ET_e\), and derive an explicit local relative-form
constant with respect to
\(2\kappa_gI+B^{\mathsf T}\mathsf W_JB\).  Determine whether the
constant is below one under (61.9) and (83.15), or isolate the remaining
parameter deficit before attempting a multiscale iteration.
\end{decision}

\begin{firewall}[Claim boundary after V83]
\label{fire:v17v83}
No uniform lower bound for \(\mathsf D\) or \(\mathsf S_P\) has been
proved.  Hence the response resolvent (79.16), environmental
susceptibility, remainder-current control, continuum measure, gauge
cofinality, physical Einstein action, and the Einstein--SM continuum
limit remain open.
\end{firewall}

\section*{Version V83 append-only record}
\addcontentsline{toc}{section}{Version V83 append-only record}
The complete V82 source was retained byte for byte before this
continuation.  It had \(672909\) bytes, \(20536\) lines, and SHA--256
\[
 \texttt{FC54AEC286D724C2AFF84C0C744C0453EB4F50B0421F8C8D569839504A0096A2}.
\]
V83 derived the block resolvent and normalized current-covariance
decay.


\section{V84: local Schur coercivity and the diagonal parameter deficit}
\label{sec:v17v84}

V83 controls the distant part of the Higgs current covariance.  This
version bounds the full normalized covariance operator, derives a
volume-independent coercive central sector, and determines whether
off-block decay alone can make the Schur complement scale safe.

\subsection{A global normalized current-covariance form bound}

Recall
\(w_e=2\chi_e^2\cosh(2d_e)\) from (83.11).  Put
\[
 K_{\rm cur}:=4Dq_HC_{\rm H,glob}^{,2}.           \tag{84.1}
\]

\begin{theorem}[Normalized current covariance]
\label{thm:v17v84-normalized}
Under the V61 influence margin, for every real edge vector \(u\),
\[
\boxed{
 u^{\mathsf T}\mathsf C_Hu
 =\operatorname{Var}_{\pi_\eta}(Q_u)
 \leq K_{\rm cur}\sum_ew_eu_e^2.}                \tag{84.2}
\]
Equivalently,
\[
 \left\|
 \operatorname{diag}(w)^{-1/2}\mathsf C_H
 \operatorname{diag}(w)^{-1/2}
 \right\|_{2\to2}
 \leq K_{\rm cur}.                                \tag{84.3}
\]
\end{theorem}

\begin{proof}
The direct incidence-gradient estimate (68.3) gives
\[
 \operatorname{Var}(Q_u)
 \leq8Dq_HC_{\rm H,glob}^{,2}\chi_*^2
       \sum_e\cosh(2d_e)u_e^2.                   \tag{84.4}
\]
Repeating its last edgewise step with the actual \(\chi_e\), rather
than replacing it by \(\chi_*\), gives
\[
 \operatorname{Var}(Q_u)
 \leq8Dq_HC_{\rm H,glob}^{,2}
       \sum_e\chi_e^2\cosh(2d_e)u_e^2,
\]
which is (84.2) by (83.11).  Conjugation by the positive diagonal
matrix \(\operatorname{diag}(w)^{1/2}\) gives (84.3).
\end{proof}

Combining (83.9) with (84.2) yields the edge-space lower form
\[
 u^{\mathsf T}\mathsf H_Hu
 \geq\sum_e\{\mathbb ET_e-K_{\rm cur}w_e\}u_e^2. \tag{84.5}
\]

\subsection{A volume-independent central sector}

Let
\[
 \Omega_R:=\{r: |d_e|\leq R\ \text{for every edge}\}.            \tag{84.6}
\]
On a graph of maximum degree \(D\), the incidence operator obeys
\[
 \|B\|_{2\to2}^2=\|B^{\mathsf T}B\|_{2\to2}
 \leq2D.                                          \tag{84.7}
\]

\begin{theorem}[Central-sector one-form coercivity]
\label{thm:v17v84-central}
On \(\Omega_R\), define
\[
 m_R:=2\kappa_g-
 16D^2q_HC_{\rm H,glob}^{,2}\chi_*^2\cosh(2R).  \tag{84.8}
\]
If \(m_R>0\), then for every vertex one-form \(\omega\) supported in
that conformal sector,
\[
\boxed{
 \mathcal Q_U[\omega]\geq m_R\|\omega\|_{L^2(\nu_U)}^2.}         \tag{84.9}
\]
The constant is independent of the number of vertices and edges.
\end{theorem}

\begin{proof}
For \(u=B\omega\), (84.2) and
\(w_e\leq2\chi_*^2\cosh(2R)\) give
\[
\begin{aligned}
 u^{\mathsf T}\mathsf C_Hu
 &\leq8Dq_HC_{\rm H,glob}^{,2}\chi_*^2
        \cosh(2R)\|B\omega\|^2\\
 &\leq16D^2q_HC_{\rm H,glob}^{,2}\chi_*^2
        \cosh(2R)\|\omega\|^2.                  \tag{84.10}
\end{aligned}
\]
In (81.5), retain the bare \(2\kappa_g\) term and discard the
nonnegative diffusion, mismatch, and mean hopping terms.  Subtract
(84.10) and integrate over the support of \(\omega\).
\end{proof}

The checkable central margin is
\[
\boxed{
 8D^2q_HC_{\rm H,glob}^{,2}\chi_*^2\cosh(2R)
 <\kappa_g.}                                      \tag{84.11}
\]
Unlike the mismatch-only margins of V64, V66, and V68, (84.11) is
independent of \(J(a)\) and is therefore compatible with canonical
scaling for every fixed \(R\), provided the bare parameters satisfy
it.

\subsection{The tail comparison and its deficit}

On the full conformal space, (84.2) is absorbed by the mismatch
quadratic form whenever
\[
 K_{\rm cur}w_e
 \leq4J\cosh(2d_e)quad\text{for every }e.         \tag{84.12}
\]
Using (83.11), this is guaranteed by
\[
\boxed{
 J\geq2Dq_HC_{\rm H,glob}^{,2}\chi_*^2,}         \tag{84.13}
\]
which recovers (68.7).

\begin{countertheorem}[Spatial decay does not cure the diagonal
canonical deficit]
\label{ctr:v17v84-diagonal}
Neither the off-block decay (83.16) nor increasing the block radius can
make the sufficient tail condition (84.13) compatible with
\(J(a)=\bar Ja^2\) at fixed positive
\(C_{\rm H,glob}\chi_*\).
\end{countertheorem}

\begin{proof}
Off-block decay estimates matrix elements with positive spatial
separation.  The normalized diagonal element at \(e=f\) is controlled
only by the order-one bound (84.3); it is present inside every spatial
block.  Comparing that local term with the mismatch conductance gives
(84.13), independently of block radius.  Its left side is fixed and
positive, while \(J(a)\to0\).
\end{proof}

This identifies the parameter deficit as local rather than spatial.
The positive mean hopping term \(\mathbb ET_e\) in (84.5) has so far
been discarded.  A scale-safe local Schur estimate must cancel the
diagonal variance \(\operatorname{Var}(q_e)\) against that mean before
using spatial decay on the off-diagonal remainder.  The scalar Ward
identity V73 demonstrates this mechanism on one edge, but does not
prove it in the conditional multi-edge law.

\begin{assessment}[What V84 closes]
\label{ass:v17v84-status}
V84 proves a global normalized current-covariance form bound and a
volume-independent central-sector coercivity theorem under (84.11).
It shows that the surviving canonical deficit is the same-edge current
variance, not the distant block coupling.  Off-block decay is ready for
use only after a local mean-curvature cancellation is proved.
\end{assessment}

\begin{decision}[V85 audit; then conditional inserted Ward identity]
\label{dec:v17v84-next}
Perform the mandatory YM/NS audit.  Then insert the local radial
generator into the fixed-environment Higgs expectation of
\(q_e\), keeping all incident edges.  Derive a conditional analogue of
(52.10) that separates
\(\mathbb ET_e-\operatorname{Var}(q_e)\) into a signed local mixed
correlation plus explicit quartic and neighboring-edge terms.  Use it
to decide whether the diagonal part of (84.5) is nonnegative or only
form bounded.
\end{decision}

\begin{firewall}[Claim boundary after V84]
\label{fire:v17v84}
Central-sector coercivity is conditional on (84.11), and the conformal
tail remains open at canonical scaling.  The response resolvent,
environmental susceptibility, remainder-current control, continuum
measure, gauge cofinality, physical Einstein action, and the
Einstein--SM continuum limit remain open.
\end{firewall}

\section*{Version V84 append-only record}
\addcontentsline{toc}{section}{Version V84 append-only record}
The complete V83 source was retained byte for byte before this
continuation.  It had \(681412\) bytes, \(20786\) lines, and SHA--256
\[
 \texttt{C38C217817E6C97AA957A50371A816A33580EDF435087BF938E8A492DF6E65A9}.
\]
V84 proved central-sector coercivity and isolated the local diagonal
parameter deficit.


\section{V85: compulsory companion audit before diagonal Ward insertion}
\label{sec:v17v85}

\subsection{Active companion checkpoints}

The newest active Yang--Mills note is
\begin{center}\small
\path{Renewal_Yang_Mills_Mass_Gap_Research_Notes_V95.tex}.
\end{center}
It has \(736823\) bytes, \(20126\) lines, and SHA--256
\[
 \texttt{6AF46F7C7F763CD984DE71828D4599128E52E96B4407B44E97F43B8CD6060799}.
                                                               \tag{85.1}
\]
YM V95 performs its own companion audit, certifies three new
exponent-three charts, and closes the first two quarter directions of
the eighth direct-plus sign cell.  Its third-coordinate extension,
full-torus positivity, cofinal control, continuum measure, and mass gap
remain open.

The newest active Navier--Stokes note remains
\begin{center}\small
\path{Renewal_NS_Stability_Research_Notes_V26.tex},
\end{center}
with \(278283\) bytes, \(5319\) lines, and SHA--256
\[
 \texttt{82C887F8C2C69E4F28FAD7C3DC8DE5B9099CB5A4BF431EBA1FFF3E91935978AD}.
                                                               \tag{85.2}
\]

\begin{theorem}[V85 companion-audit decision]
\label{thm:v17v85-decision}
Neither companion note proves a conditional identity or estimate for
\[
 \mathbb E_{\pi_\eta}T_e-
 \operatorname{Var}_{\pi_\eta}(q_e),              \tag{85.3}
\]
which is the local diagonal deficit isolated in V84.  No theorem is
imported at V85.
\end{theorem}

\begin{proof}
YM V95 remains a finite momentum-symbol chart theorem and contains no
conditional Higgs radial generator.  NS V26 concerns Oseen-kernel
tensor norms.  Neither statement has the variables or hypotheses of
(85.3), and the YM atlas is not yet a full-torus or Gibbs covariance
theorem.
\end{proof}

\begin{assessment}[Direction after the audit]
\label{ass:v17v85-direction}
The YM programme is one coordinate direction from completing its last
finite-cutoff direct-plus cell, but that progress cannot yet enter the
SM environment resolvent.  The next step remains internal: derive the
fixed-environment inserted radial Ward identity for (85.3), retaining
all incident-edge terms rather than reducing them to the scalar model.
\end{assessment}

\begin{decision}[V86: conditional inserted radial Ward identity]
\label{dec:v17v85-next}
Apply the site dilation generator \(X_x=y_x\cdot\nabla_{y_x}\) under
\(\pi_\eta\) with insertion \(q_e\).  Expand
\(X_xS_H-q_H\) into mass, quartic, endpoint diagonal, and transported
cross terms.  Sum the two endpoints with the correct signs and solve
exactly for \(\mathbb ET_e-\operatorname{Var}(q_e)\).  Keep the
neighboring-edge correlations explicit and test their spatial decay
with V83.
\end{decision}

\begin{firewall}[Claim boundary after V85]
\label{fire:v17v85}
This version is an audit only.  The diagonal cancellation, response
resolvent, environmental susceptibility, remainder-current control,
continuum measure, gauge cofinality, physical Einstein action, and the
Einstein--SM continuum limit remain open.
\end{firewall}

\section*{Version V85 append-only record}
\addcontentsline{toc}{section}{Version V85 append-only record}
The complete V84 source was retained byte for byte before this
continuation.  It had \(688059\) bytes, \(20983\) lines, and SHA--256
\[
 \texttt{422842AF5DBF007046271C4DD3926709BE8072B799E1C31F076CFADA47FB7CDF}.
\]
V85 audited YM V95 and NS V26 and retained the conditional radial Ward
acquisition.


\section{V86: conditional inserted radial Ward identity for the diagonal curvature}
\label{sec:v17v86}

V84 showed that the canonical deficit is the same-edge quantity
\(\mathbb ET_e-\operatorname{Var}(q_e)\).  This version evaluates that
quantity exactly under the fixed-environment Higgs law by inserting
the edge current into the two endpoint radial identities.

\subsection{Site radial scores}

For a site \(x\), let
\[
 \mathcal X_x:=y_x\cdot\nabla_{y_x},
 \qquad
 \mathcal R_x:=\mathcal X_xS_H-q_H.               \tag{86.1}
\]
Using the geometric hopping square (52.2),
\[
\boxed{
 \mathcal R_x
 =2m_H^2A_x+4\lambda_HA_x^2
 +2\sum_{f\ni x}(D_{x,f}-C_f)-q_H,}               \tag{86.2}
\]
where
\[
 A_x=\|y_x\|^2,qquad
 C_f=\chi_f\operatorname{Re}
 \langle y_{f_+},U_fy_{f_-}\rangle.              \tag{86.3}
\]

\begin{lemma}[Conditional radial integration by parts]
\label{lem:v17v86-ibp}
For every smooth polynomially growing Higgs observable \(\Psi\),
\[
\boxed{
 \mathbb E_{\pi_\eta}(\mathcal X_x\Psi)
 =\mathbb E_{\pi_\eta}(\Psi\mathcal R_x).}        \tag{86.4}
\]
In particular,
\[
 \mathbb E_{\pi_\eta}\mathcal R_x=0.             \tag{86.5}
\]
\end{lemma}

\begin{proof}
Integrate the divergence of
\(y_x\Psi e^{-S_H}\) over the realified \(q_H\)-dimensional fibre.
Quartic confinement removes the boundary term and gives
\[
 0=\mathbb E(q_H\Psi+\mathcal X_x\Psi
             -\Psi\mathcal X_xS_H),
\]
which is (86.4).  Set \(\Psi=1\) for (86.5).
\end{proof}

\subsection{The inserted edge identity}

Fix an oriented edge \(e=(v,w)\).  Since
\(q_e=D_{v,e}-D_{w,e}\),
\[
 \mathcal X_vq_e=2D_{v,e},
 \qquad
 \mathcal X_wq_e=-2D_{w,e}.                      \tag{86.6}
\]
Define the endpoint-star remainder
\[
\begin{aligned}
 \mathcal K_e:={}&
 2m_H^2(A_v-A_w)+4\lambda_H(A_v^2-A_w^2)\\
 &+2\sum_{\substack{f\ni v\\f\neq e}}(D_{v,f}-C_f)
 -2\sum_{\substack{f\ni w\\f\neq e}}(D_{w,f}-C_f).            \tag{86.7}
\end{aligned}
\]
The edge contribution in (86.2) then gives the pointwise identity
\[
 \mathcal R_v-\mathcal R_w=2q_e+\mathcal K_e.     \tag{86.8}
\]

\begin{theorem}[Conditional inserted edge Ward identity]
\label{thm:v17v86-inserted}
At every fixed conformal and gauge environment,
\[
\boxed{
 2\mathbb E_{\pi_\eta}T_e
 =\mathbb E_{\pi_\eta}
 [q_e(\mathcal R_v-\mathcal R_w)],}               \tag{86.9}
\]
and hence
\[
\boxed{
 \mathbb ET_e-\mathbb Eq_e^2
 =\frac12\mathbb E(q_e\mathcal K_e).}             \tag{86.10}
\]
All expectations in (86.10) are conditional on the fixed environment.
\end{theorem}

\begin{proof}
Apply (86.4) with \(\Psi=q_e\) first at \(v\) and then at \(w\).
Equation (86.6) gives
\[
 \mathbb E(q_e\mathcal R_v)=2\mathbb ED_{v,e},
 \qquad
 \mathbb E(q_e\mathcal R_w)=-2\mathbb ED_{w,e}.  \tag{86.11}
\]
Subtract and use \(T_e=D_{v,e}+D_{w,e}\), proving (86.9).
Insert (86.8) and solve for the second moment of \(q_e\), proving
(86.10).
\end{proof}

Taking expectations in (86.8) and using (86.5) gives
\[
 \mathbb E\mathcal K_e=-2\mathbb Eq_e.           \tag{86.12}
\]
This converts (86.10) from a raw moment identity to a centered one.

\begin{corollary}[Exact same-edge free-energy curvature]
\label{cor:v17v86-diagonal}
The diagonal entry of the Higgs free-energy Hessian is
\[
\boxed{
 (\mathsf H_H)_{ee}
 =\mathbb ET_e-\operatorname{Var}(q_e)
 =\frac12\operatorname{Cov}(q_e,\mathcal K_e).}   \tag{86.13}
\]
\end{corollary}

\begin{proof}
By (86.10),
\[
 \mathbb ET_e-\operatorname{Var}(q_e)
 =\frac12\mathbb E(q_e\mathcal K_e)+(\mathbb Eq_e)^2.
\]
Use (86.12) to identify the right side with
\(\frac12\operatorname{Cov}(q_e,\mathcal K_e)\).
The first equality is (63.3) on the coordinate direction \(e\).
\end{proof}

For the massless one-edge scalar model, the neighboring-edge sums in
(86.7) vanish and \(\mathcal K_e\) is the quartic endpoint imbalance.
After the scaling of V72, (86.13) is precisely the Ward covariance in
(73.7).  Thus V73--V78 are the exact scalar specialization of the
multi-edge identity, rather than an unrelated analogy.

\subsection{Diagonal/off-diagonal Hessian separation}

Let \(\mathsf C_H^{\rm off}\) be the edge matrix
\[
 (\mathsf C_H^{\rm off})_{ef}
 =\begin{cases}
   \operatorname{Cov}(q_e,q_f),&e\neq f,\\
   0,&e=f.
  \end{cases}                                     \tag{86.14}
\]
and put
\[
 h_e^{\rm Ward}:=\frac12
 \operatorname{Cov}(q_e,\mathcal K_e).            \tag{86.15}
\]

\begin{theorem}[Ward-renormalized Higgs Hessian]
\label{thm:v17v86-hessian}
The full edge Hessian has the exact decomposition
\[
\boxed{
 \mathsf H_H
 =\operatorname{diag}_e(h_e^{\rm Ward})
 -\mathsf C_H^{\rm off}.}                         \tag{86.16}
\]
\end{theorem}

\begin{proof}
For \(e=f\), use Corollary~\ref{cor:v17v86-diagonal}.  For
\(e\neq f\), the second derivative of the Higgs action vanishes and
(63.3) gives
\((\mathsf H_H)_{ef}=-\operatorname{Cov}(q_e,q_f)\).
\end{proof}

The off-diagonal matrix in (86.16) is governed by the spatial envelope
(83.12) and its summable tail (83.16).  The diagonal has been removed
from that estimate exactly.  What remains at zero spatial separation
is the local star covariance (86.15), whose observable
\(\mathcal K_e\) is supported on the endpoints of \(e\) and their
incident edges.

Spatial decay cannot determine the sign of (86.15): its two
observables have overlapping support.  The V61 influence expansion can
make remote paths small, but its zeroth and first local paths remain.
This is the precise local input required before the V83 Schur
complement can be closed.

\begin{assessment}[What V86 closes]
\label{ass:v17v86-status}
V86 proves the conditional inserted radial Ward identity and exactly
renormalizes the same-edge mean-curvature/variance difference into the
local covariance (86.15).  It decomposes the full Higgs Hessian into
that Ward diagonal and an off-diagonal covariance matrix already
controlled at long distance.  The remaining obstruction is a bounded
endpoint-star calculation, not a volume-scale covariance.
\end{assessment}

\begin{decision}[V87: local endpoint-star form]
\label{dec:v17v86-next}
Expand \(\operatorname{Cov}(q_e,\mathcal K_e)\) into its mass,
quartic, and neighboring-edge pieces.  Use conditional reflection and
the one-site tilted-quartic gap to determine the sign of the on-site
piece.  Bound each neighboring-edge term by one explicit power of the
influence matrix \(M\), retaining its orientation.  Derive a local
relative-form constant for the Ward diagonal plus all incident
off-diagonal current covariances.
\end{decision}

\begin{firewall}[Claim boundary after V86]
\label{fire:v17v86}
The Ward identity does not prove that \(h_e^{\rm Ward}\) is
nonnegative.  Local star coercivity, the response resolvent,
environmental susceptibility, remainder-current control, continuum
measure, gauge cofinality, physical Einstein action, and the
Einstein--SM continuum limit remain open.
\end{firewall}

\section*{Version V86 append-only record}
\addcontentsline{toc}{section}{Version V86 append-only record}
The complete V85 source was retained byte for byte before this
continuation.  It had \(691486\) bytes, \(21071\) lines, and SHA--256
\[
 \texttt{241BB0CD18073020818EB54C8FB2492D72606FE0F2B550DD225560F1A8CF5FBF}.
\]
V86 proved the conditional inserted Ward identity and the
Ward-renormalized Higgs Hessian decomposition.


\section{V87: uniform endpoint-star moments and a local Ward form bound}
\label{sec:v17v87}

V86 reduces the diagonal Higgs curvature to the covariance of two
overlapping local observables.  Reflection symmetry does not assign a
sign to that covariance, because both observables are even under the
global Higgs reflection.  We therefore derive scale-independent
conditional moments from the V61 gap and use them to obtain an explicit
local form bound.

\subsection{Moment recursion from the global conditional gap}

Write
\[
 C_G:=C_{\rm H,glob},
 \qquad A_x=\|y_x\|^2.                             \tag{87.1}
\]
The fixed-environment Higgs law is invariant under
\(y\mapsto-y\), so \(\mathbb Ey_x=0\).  V67 therefore gives
\[
 \mathbb EA_x\leq q_HC_G.                         \tag{87.2}
\]

\begin{theorem}[Uniform conditional site moments]
\label{thm:v17v87-moments}
Under the V61 influence margin, for every integer \(p\geq2\),
\[
\boxed{
 \mathbb EA_x^p
 \leq(p^2+q_H)C_G\,\mathbb EA_x^{p-1}.}           \tag{87.3}
\]
In particular,
\[
\begin{aligned}
 \mathbb EA_x^2
 &\leq q_H(q_H+4)C_G^2=:M_2^G,\\
 \mathbb EA_x^3
 &\leq q_H(q_H+4)(q_H+9)C_G^3=:M_3^G.             \tag{87.4}
\end{aligned}
\]
These bounds are uniform in the finite volume and the fixed conformal
and gauge environment.
\end{theorem}

\begin{proof}
Apply the global conditional Poincar\'e inequality to
\(f=A_x^{p/2}=\|y_x\|^p\).  Since
\[
 \|\nabla f\|^2=p^2A_x^{p-1},                    \tag{87.5}
\]
one has
\[
 \mathbb EA_x^p
 \leq p^2C_G\mathbb EA_x^{p-1}
      +(\mathbb EA_x^{p/2})^2.                   \tag{87.6}
\]
Cauchy--Schwarz gives
\[
 (\mathbb EA_x^{p/2})^2
 \leq(\mathbb EA_x^{p-1})(\mathbb EA_x)
 \leq q_HC_G\mathbb EA_x^{p-1}.                 \tag{87.7}
\]
This proves (87.3).  Use (87.2) at \(p=2\), then iterate once more at
\(p=3\), to obtain (87.4).
\end{proof}

\subsection{Gradient budget for the endpoint-star remainder}

For an edge \(e=(v,w)\), define its radius-one conformal star sector
\[
 \Omega_{e,R}^{\rm star}
 :=\{r:|d_f|\leq R\text{ for every }f\text{ incident to }v
                  \text{ or }w\}.                \tag{87.8}
\]
On this sector, \(g_{x,f}\leq\chi_*e^R\) for every half-edge in the
star.  Put
\[
\begin{aligned}
 G_{\rm site}
 &:=32m_H^4q_HC_G+512\lambda_H^2M_3^G,\\
 G_{\rm hop}(R)
 &:=44q_HC_G\chi_*^2e^{2R},\\
 G_{\mathcal K}(R)
 &:=4D\{G_{\rm site}+(D-1)G_{\rm hop}(R)\}.       \tag{87.9}
\end{aligned}
\]
Here \(m_H^4=(m_H^2)^2\).

\begin{lemma}[Endpoint-star gradient bound]
\label{lem:v17v87-gradient}
On \(\Omega_{e,R}^{\rm star}\),
\[
\boxed{
 \mathbb E\|\nabla_y\mathcal K_e\|^2
 \leq G_{\mathcal K}(R).}                        \tag{87.10}
\]
\end{lemma}

\begin{proof}
The gradient of the on-site part
\(2m_H^2A_x+4\lambda_HA_x^2\) is
\[
 4m_H^2y_x+16\lambda_HA_xy_x.
\]
The square of its norm is at most
\[
 32m_H^4A_x+512\lambda_H^2A_x^3,                 \tag{87.11}
\]
whose expectation is bounded by \(G_{\rm site}\).

For a neighboring edge \(f=(x,z)\), the term
\(2(D_{x,f}-C_f)\) has gradients
\[
 4g_{x,f}y_x-2\chi_fU_f^{(x)}y_z,
 \qquad -2\chi_f(U_f^{(x)})^*y_x                \tag{87.12}
\]
in the two endpoint fibres.  Using
\(\|a-b\|^2\leq2\|a\|^2+2\|b\|^2\), its full
squared gradient is at most
\[
 (32g_{x,f}^2+4\chi_f^2)A_x+8\chi_f^2A_z.        \tag{87.13}
\]
Equations (87.2) and (87.8) bound its expectation by
\(G_{\rm hop}(R)\).

At either endpoint, the star observable in (86.7) is a sum of at most
\(D\) terms.  Bound the squared gradient of a sum by the number of
terms times the sum of squared gradients.  Finally use
\(\|a-b\|^2\leq2\|a\|^2+2\|b\|^2\) for the difference of the two
endpoint stars.  This gives the factor \(4D\) in (87.9).
\end{proof}

The current gradient has the exact budget
\[
\begin{aligned}
 \mathbb E\|\nabla_yq_e\|^2
 &=4\{g_{v,e}^2\mathbb EA_v+g_{w,e}^2\mathbb EA_w\}\\
 &\leq4q_HC_Gw_e.                                 \tag{87.14}
\end{aligned}
\]

\begin{theorem}[Local Ward-diagonal form bound]
\label{thm:v17v87-ward-bound}
On \(\Omega_{e,R}^{\rm star}\),
\[
\boxed{
 |h_e^{\rm Ward}|
 \leq C_G\sqrt{q_HC_GG_{\mathcal K}(R)}\,\sqrt{w_e}.}            \tag{87.15}
\]
Consequently,
\[
 h_e^{\rm Ward}\geq-H_{\rm Ward}(R),             \tag{87.16}
\]
where
\[
 H_{\rm Ward}(R)
 :=\sqrt2\,C_G\chi_*\sqrt{q_HC_GG_{\mathcal K}(R)
                         \cosh(2R)}.              \tag{87.17}
\]
\end{theorem}

\begin{proof}
Conditional Cauchy--Schwarz followed by the global Poincar\'e
inequality gives
\[
 |\operatorname{Cov}(q_e,\mathcal K_e)|
 \leq C_G
 \bigl(\mathbb E\|\nabla q_e\|^2\bigr)^{1/2}
 \bigl(\mathbb E\|\nabla\mathcal K_e\|^2\bigr)^{1/2}.    \tag{87.18}
\]
Use (87.10), (87.14), and
\(h_e^{\rm Ward}=\frac12\operatorname{Cov}(q_e,\mathcal K_e)\)
to prove (87.15).  On the star sector,
\(w_e\leq2\chi_*^2\cosh(2R)\), which gives
(87.16)--(87.17).
\end{proof}

This estimate is scale independent on a fixed star sector.  It does
not determine the sign of \(h_e^{\rm Ward}\).  The reflection
\(y\mapsto-y\) leaves both \(q_e\) and \(\mathcal K_e\) unchanged, so
reflection centering cannot remove their covariance.  Likewise, a
Poincar\'e inequality bounds its magnitude but has no sign content.

\subsection{A bounded-sector full local form}

The normalized covariance theorem (84.3) implies
\[
 \left\|
 \operatorname{diag}(w)^{-1/2}\mathsf C_H^{\rm off}
 \operatorname{diag}(w)^{-1/2}
 \right\|
 \leq2K_{\rm cur},                                \tag{87.19}
\]
because removing the diagonal changes the normalized operator by at
most another \(K_{\rm cur}\).

\begin{corollary}[Explicit star-sector one-form floor]
\label{cor:v17v87-floor}
On a conformal sector where every edge satisfies \(|d_e|\leq R\),
the one-form quadratic form obeys
\[
 \mathcal Q_U[\omega]
 \geq m_R^{\rm Ward}\|\omega\|_2^2,              \tag{87.20}
\]
with
\[
\boxed{
 m_R^{\rm Ward}
 :=2\kappa_g-2D\left\{
 H_{\rm Ward}(R)+4K_{\rm cur}\chi_*^2\cosh(2R)
 \right\}.}                                      \tag{87.21}
\]
Thus \(m_R^{\rm Ward}>0\) is a checkable, volume-independent local
coercivity condition.
\end{corollary}

\begin{proof}
For \(u=B\omega\), (87.16) bounds the Ward diagonal below by
\(-H_{\rm Ward}(R)\|u\|^2\).  Equations (87.19) and
\(w_e\leq2\chi_*^2\cosh(2R)\) bound the off-diagonal subtraction below
by
\(-4K_{\rm cur}\chi_*^2\cosh(2R)\|u\|^2\).
Use \(\|B\omega\|^2\leq2D\|\omega\|^2\), retain the bare mass, and
discard the other nonnegative terms in (81.5).
\end{proof}

The new floor is rigorous but not automatically sharper than V84: it
uses the Ward cancellation before estimating, then pays the resulting
local covariance by Cauchy--Schwarz.  Its value is that it isolates the
only missing improvement: a sign or one-sided estimate for the on-site
part of \(\operatorname{Cov}(q_e,\mathcal K_e)\).  Two-sided spectral
gap estimates cannot supply that improvement.

\begin{assessment}[What V87 closes]
\label{ass:v17v87-status}
V87 proves uniform conditional moments through sixth field order,
an explicit endpoint-star gradient budget, and a volume-independent
Ward-renormalized local one-form floor.  It also proves that reflection
and Poincar\'e estimates do not determine the needed covariance sign.
The next acquisition must be one-sided.
\end{assessment}

\begin{decision}[V88: one-sided radial monotonicity]
\label{dec:v17v87-next}
Condition on all Higgs variables outside the two endpoints of \(e\).
Use the two tilted-quartic radial laws to compare the conditional
expectations of \(A_x\) and \(A_x^2\) as their diagonal coefficients
vary.  Test whether the on-site part of
\(\operatorname{Cov}(q_e,\mathcal K_e)\) is nonnegative after this
two-site conditioning.  Keep the neighboring-edge remainder separate
and charge it by one influence step only after a one-sided on-site
floor is obtained.
\end{decision}

\begin{firewall}[Claim boundary after V87]
\label{fire:v17v87}
The Ward diagonal has only a two-sided bound, not a proved positive
sign.  Global resolvent coercivity, environmental susceptibility,
remainder-current control, continuum measure, gauge cofinality,
physical Einstein action, and the Einstein--SM continuum limit remain
open.
\end{firewall}

\section*{Version V87 append-only record}
\addcontentsline{toc}{section}{Version V87 append-only record}
The complete V86 source was retained byte for byte before this
continuation.  It had \(698792\) bytes, \(21302\) lines, and SHA--256
\[
 \texttt{7E37CB2E61F2723400B292E04FC61A09ADC838D6C634CDF4C231365F8B77DD1F}.
\]
V87 proved uniform star moments and an explicit local Ward form bound.


\section{V88: endpoint conditioning and the two-level Ward-sign obstruction}
\label{sec:v17v88}

V87 shows that a two-sided spectral-gap estimate is sufficient for a
bounded-sector Hessian floor but cannot assign a sign to the Ward
diagonal.  We now carry out the endpoint conditioning proposed there.
The calculation separates two logically different sign questions: one
inside the two-endpoint fibre and one in the exterior environment.

\subsection{The exact conditional two-site law}

Fix an oriented edge \(e=(v,w)\), and let \(\mathscr G_e\) be the
sigma-field generated by all Higgs variables except \(y_v,y_w\).  After
realifying the fibres, the conditional potential has the exact form
\[
\begin{aligned}
 V_e(y_v,y_w)={}&
 \lambda_H(A_v^2+A_w^2)+a_vA_v+a_wA_w\\
 &-2\chi_e\operatorname{Re}\langle y_w,U_ey_v\rangle
 -2\langle b_v,y_v\rangle_{\mathbb R}
 -2\langle b_w,y_w\rangle_{\mathbb R},              \tag{88.1}
\end{aligned}
\]
up to a \(\mathscr G_e\)-measurable additive constant, where
\[
 a_x=m_H^2+\sum_{f\ni x}g_{x,f},
 \qquad
 b_x=\sum_{\substack{f\ni x\\f\ne e}}b_{x,f}.       \tag{88.2}
\]
Here \(b_{x,f}\) is the real vector determined by
\(C_f=\langle b_{x,f},y_x\rangle_{\mathbb R}\) when the other endpoint
of \(f\) is frozen.  Formula (88.1) follows by collecting the two
endpoint terms in the geometric hopping squares.

Define
\[
 P(A):=2m_H^2A+4\lambda_HA^2,
 \qquad
 N_x:=2\sum_{\substack{f\ni x\\f\ne e}}
       (g_{x,f}A_x-\langle b_{x,f},y_x\rangle_{\mathbb R}).           \tag{88.3}
\]
Then the two observables in the Ward diagonal are
\[
 q_e=g_{v,e}A_v-g_{w,e}A_w,
 \qquad
 \mathcal K_e=P(A_v)-P(A_w)+N_v-N_w.                \tag{88.4}
\]

\subsection{Total covariance is an unavoidable second level}

Write \(\mathbb E_e(\cdot)=\mathbb E(\cdot\mid\mathscr G_e)\),
\(\bar q_e=\mathbb E_eq_e\), and
\(\overline{\mathcal K}_e=\mathbb E_e\mathcal K_e\).

\begin{theorem}[Two-level Ward covariance]
\label{thm:v17v88-total-cov}
At every fixed conformal and gauge environment,
\[
\boxed{
 h_e^{\rm Ward}
 =\frac12\mathbb E\operatorname{Cov}_e(q_e,\mathcal K_e)
 +\frac12\operatorname{Cov}
       (\bar q_e,\overline{\mathcal K}_e).}          \tag{88.5}
\]
Consequently, positivity of every conditional two-site covariance is
not by itself sufficient for \(h_e^{\rm Ward}\geq0\); the exterior
covariance in (88.5) also needs a one-sided estimate.
\end{theorem}

\begin{proof}
For square-integrable random variables \(F,G\), expand
\(F=(F-\mathbb E_eF)+\mathbb E_eF\) and the analogous expression for
\(G\).  The two cross terms have zero expectation conditionally on
\(\mathscr G_e\), and therefore
\[
 \operatorname{Cov}(F,G)
 =\mathbb E\operatorname{Cov}_e(F,G)
  +\operatorname{Cov}(\mathbb E_eF,\mathbb E_eG).    \tag{88.6}
\]
Apply this identity to \(F=q_e\), \(G=\mathcal K_e\), and use
(86.15).  The last assertion follows directly from (88.5).
\end{proof}

The second term is not a remote tail.  Both conditional expectations
depend on the radius-one exterior star through \(b_v,b_w\).  It must
therefore be treated before the long-distance estimate of V83 can be
invoked.

\subsection{What radial monotonicity does prove}

The elementary identity
\[
 \operatorname{Cov}(X,F(X))
 =\frac12\mathbb E[(X-X')(F(X)-F(X'))]               \tag{88.7}
\]
holds for an independent copy \(X'\).  Hence it is nonnegative whenever
\(F\) is nondecreasing.

\begin{lemma}[Positive self terms]
\label{lem:v17v88-self}
Assume \(m_H^2\geq0\) and \(\lambda_H>0\).  Under each conditional law
in (88.1),
\[
 \operatorname{Cov}_e(A_x,P(A_x))\geq0,
 \qquad x\in\{v,w\}.                                \tag{88.8}
\]
Moreover the radial part of the conditional covariance expands as
\[
\begin{aligned}
 &\operatorname{Cov}_e
 \bigl(g_{v,e}A_v-g_{w,e}A_w,,P(A_v)-P(A_w)\bigr)\\
 ={}&g_{v,e}\operatorname{Cov}_e(A_v,P(A_v))
    +g_{w,e}\operatorname{Cov}_e(A_w,P(A_w))\\
 &-g_{v,e}\operatorname{Cov}_e(A_v,P(A_w))
  -g_{w,e}\operatorname{Cov}_e(A_w,P(A_v)).         \tag{88.9}
\end{aligned}
\]
Thus scalar monotonicity controls the first line on the right but does
not determine the two mixed terms.
\end{lemma}

\begin{proof}
Since \(P'(A)=2m_H^2+8\lambda_HA\geq0\) for \(A\geq0\), (88.7)
proves (88.8) for the marginal conditional distribution of \(A_x\).
Equation (88.9) is bilinearity of covariance.
\end{proof}

When \(\chi_e=0\), the density (88.1) factorizes between \(v\) and
\(w\), so both mixed covariances in (88.9) vanish.  For
\(\chi_e\ne0\), the transported bilinear term couples the radii and the
mixed terms are genuine.  In addition, the neighboring contribution
in (88.3) contains angular linear observables.  At a fixed value of
\(A_x\), the quantity
\(-\langle b_x,y_x\rangle_{\mathbb R}\) ranges between
\(-\|b_x\|\sqrt{A_x}\) and \(\|b_x\|\sqrt{A_x}\); it is not a function
of \(A_x\).  Therefore a one-dimensional Chebyshev argument cannot be
applied to \(N_x\).

\subsection{Response reformulation of the missing sign}

For \(t\) in any compact interval on which the finite-volume integral
is taken, introduce the anti-diagonal perturbation
\[
 d\pi_t=Z_t^{-1}e^{-S_H-tq_e}\,dy.                  \tag{88.10}
\]
Quartic confinement permits differentiation under the integral.

\begin{proposition}[Ward sign as a monotone response]
\label{prop:v17v88-response}
With \(\mathcal K_e\) fixed as the observable (86.7),
\[
 \frac{d}{dt}\mathbb E_t\mathcal K_e
 =-\operatorname{Cov}_t(\mathcal K_e,q_e).           \tag{88.11}
\]
In particular, at \(t=0\),
\[
\boxed{
 h_e^{\rm Ward}
 =-\frac12\left.\frac{d}{dt}
          \mathbb E_t\mathcal K_e\right|_{t=0}.}    \tag{88.12}
\]
Thus the desired sign is exactly the assertion that the expected
endpoint-star force difference decreases under an anti-diagonal
increase of the endpoint quadratic coefficients.
\end{proposition}

\begin{proof}
Differentiate the numerator and denominator defining
\(\mathbb E_t\mathcal K_e\).  Their quotient derivative is
\(-\mathbb E_t(\mathcal K_eq_e)+
\mathbb E_t\mathcal K_e\,\mathbb E_tq_e\), proving (88.11).
Equation (88.12) is (86.15).
\end{proof}

This reformulation prevents a circular use of the conformal Hessian:
(88.12) asks for monotonicity of a local force expectation, whereas the
Hessian positivity being sought is the conclusion.  Any proof of the
response sign must come from the two-site convex structure, an
order-preserving coupling, or a direct interpolation estimate.

\begin{assessment}[What endpoint conditioning resolves]
\label{ass:v17v88-status}
V88 proves the exact conditional two-site potential, the two-level
covariance decomposition, positivity of the radial self terms, and the
response identity (88.12).  It also locates three pieces not settled by
scalar monotonicity: transported mixed-radius covariances, angular
exterior fields, and the exterior covariance in (88.5).
\end{assessment}

\begin{decision}[V89: transported-coupling interpolation]
\label{dec:v17v88-next}
Introduce a parameter \(s\in[0,1]\) in front of the transported
endpoint term in (88.1).  At \(s=0\) the endpoint law factorizes and
the mixed radial covariances vanish.  Differentiate them in \(s\), use
the conditional gap and the V87 moments, and obtain an explicit
first-order error bound.  Apply the same interpolation to the angular
neighbor term.  Compare the resulting error with the positive self
part without assuming a uniform lower variance that has not been
proved.
\end{decision}

\begin{firewall}[Claim boundary after V88]
\label{fire:v17v88}
The self-covariance sign does not establish the full Ward sign.  A
quantitative lower bound for the positive self part or a direct
monotone-response theorem is still missing.  Global resolvent
coercivity, environmental susceptibility, remainder-current control,
continuum measure, gauge cofinality, physical Einstein action, and the
Einstein--SM continuum limit remain open.
\end{firewall}

\section*{Version V88 append-only record}
\addcontentsline{toc}{section}{Version V88 append-only record}
The complete V87 source was retained byte for byte before this
continuation.  It had \(707359\) bytes, \(21575\) lines, and SHA--256
\[
 \texttt{D8CFC3E955C979747431C65591A3D1DCE7BB7FCE0556C1F486C0D1AF0515A39C}.
\]
V88 proved the exact two-level endpoint-conditioning decomposition and
recast the missing Ward sign as a monotone local response.


\section{V89: failure of pointwise exterior-uniform interpolation and the upstream repair}
\label{sec:v17v89}

V88 suggests interpolating the transported endpoint coupling while
conditioning on the exterior.  Before estimating the interpolation we
must check that its constants are uniform in the conditioning variable.
They are not: a bounded conformal star bounds the hopping coefficients,
but it does not bound the exterior Higgs fields \(b_v,b_w\) in (88.1).

\subsection{A tilted quartic conditional moment obstruction}

Consider first the one-dimensional probability law
\[
 d\nu_B(t)=Z_B^{-1}
 \exp\{-\lambda t^4-at^2+2Bt\}\,dt,
 \qquad \lambda>0,                                  \tag{89.1}
\]
where \(a\) is fixed and \(B>0\).

\begin{lemma}[Large-field scaling]
\label{lem:v17v89-large-field}
For the law (89.1),
\[
 \boxed{
 \lim_{B\to\infty}B^{-2/3}\mathbb E_{\nu_B}t^2
 =(2\lambda)^{-2/3}.}                               \tag{89.2}
\]
In particular, no bound on the conditional second moment can be
uniform in the linear tilt \(B\).
\end{lemma}

\begin{proof}
Put \(t=B^{1/3}z\).  Apart from the Jacobian, the density becomes
\[
 \exp\{-B^{4/3}\phi_B(z)\}\,dz,
 \qquad
 \phi_B(z)=\lambda z^4-2z+aB^{-2/3}z^2.             \tag{89.3}
\]
The limiting function \(\phi(z)=\lambda z^4-2z\) is coercive and has
the unique minimizer
\(z_*=(2\lambda)^{-1/3}\).  Fix \(\varepsilon>0\).  Coercivity and
uniqueness give numbers \(R,c_\varepsilon>0\) such that, for all large
\(B\),
\[
 \inf_{\substack{|z-z_*|\geq\varepsilon\\|z|\leq R}}
 [\phi_B(z)-\phi_B(z_*)]\geq c_\varepsilon,
 \qquad
 \phi_B(z)-\phi_B(z_*)\geq\tfrac{\lambda}{4}z^4
 \quad(|z|\geq R).                                  \tag{89.4}
\]
Integrating (89.4), and using a fixed small interval about \(z_*\) for
a lower bound on the denominator, shows that the normalized
\(z\)-law assigns probability tending to one to
\(|z-z_*|<\varepsilon\).  The quartic tail in (89.4) also makes
\(z^2\) uniformly integrable.  Hence
\(\mathbb E z^2\to z_*^2\).  Since \(t^2=B^{2/3}z^2\), this is (89.2).
\end{proof}

The same conclusion holds in a real fibre of any fixed dimension by
restricting to the direction of the tilt and applying the identical
Laplace argument; the transverse Gaussian-scale fluctuations do not
change the minimizing radius.

\begin{theorem}[No exterior-uniform pointwise moment budget]
\label{thm:v17v89-no-uniform}
Suppose an endpoint \(x\) has a neighboring edge \(f\ne e\) with
\(\chi_f>0\).  Even on \(\Omega_{e,R}^{\rm star}\), there is no finite
constant depending only on the fixed model parameters and \(R\) that
bounds
\[
 \mathbb E_e A_x                                      \tag{89.5}
\]
for every exterior configuration.
\end{theorem}

\begin{proof}
Freeze all exterior variables except the other endpoint \(y_z\) of
\(f\), and choose \(y_z\) along a transported unit fibre vector with
norm tending to infinity.  Then the vector \(b_{x,f}\) in (88.2) has
norm \(B=\chi_f\|y_z\|\to\infty\), whereas the bounded conformal-star
condition keeps the quadratic coefficients finite.  Restricting the
conditional potential to the tilt direction gives (89.1), with
additional fixed-dimensional transverse variables and the endpoint
coupling to \(y_{e\setminus x}\).  The joint coercive Laplace principle
has a minimizing \(x\)-radius of order \(B^{1/3}\); equivalently, the
directional argument in Lemma~\ref{lem:v17v89-large-field}, after
integrating the remaining variables, gives
\(\mathbb E_eA_x\geq cB^{2/3}\) for all sufficiently large \(B\).
Thus (89.5) is unbounded over exterior configurations.
\end{proof}

This does not contradict V87.  The moment estimates (87.2)--(87.4)
are expectations under the full fixed-environment Higgs Gibbs law;
they average over the neighboring Higgs variables.  They were never
pointwise claims after conditioning on an arbitrary Higgs exterior.

\subsection{Why the naive transported interpolation is not closed}

Insert \(s\in[0,1]\) in front of the transported term in (88.1), and
write \(\operatorname{Cov}_{e,s}\) for the resulting conditional
covariance.  For fixed exterior data and integrable observables
\(F,G\), differentiation gives
\[
 \frac d{ds}\operatorname{Cov}_{e,s}(F,G)
 =2\,\kappa_{e,s}(F,G,C_e),                          \tag{89.6}
\]
where \(\kappa\) is the centered third joint cumulant.  Indeed, the
derivative of any conditional expectation is
\(2\operatorname{Cov}_{e,s}(\,cdot\,,C_e)\), and
the product rule yields (89.6).

At \(s=0\) the two endpoint variables factorize conditionally, so the
mixed covariances in (88.9) vanish.  Formally,
\[
 \operatorname{Cov}_{e,1}(A_v,P(A_w))
 =2\int_0^1\kappa_{e,s}(A_v,P(A_w),C_e)\,ds.         \tag{89.7}
\]
A H\"older estimate of the integrand necessarily contains conditional
moments of \(A_v,A_w\).  Theorem~\ref{thm:v17v89-no-uniform} proves that
the V87 constants cannot be substituted into (89.7) pointwise in the
exterior.  Doing so would be an invalid interchange of an unconditional
bound and a conditional supremum.

\subsection{The admissible averaged interpolation}

There is nevertheless an exact route that retains the exterior
average.  Let \(\rho(d\zeta)\) be the original marginal law of the
exterior configuration, let \(\pi_{e,s}^{\zeta}\) be the normalized
two-site conditional law with transported parameter \(s\), and define
the hybrid probability measure
\[
 d\mu_s(\zeta,y_v,y_w)
 :=\rho(d\zeta)\,d\pi_{e,s}^{\zeta}(y_v,y_w).        \tag{89.8}
\]
Then for every integrable \(F\),
\[
 \frac d{ds}\mathbb E_{\mu_s}F
 =2\mathbb E_\rho
   \operatorname{Cov}_{e,s}^{\zeta}(F,C_e).         \tag{89.9}
\]
No derivative of the exterior marginal appears because \(\rho\) is
held fixed.  At \(s=1\), \(\mu_1\) is the original Gibbs law.

\begin{proposition}[Exact averaged mixed-term representation]
\label{prop:v17v89-averaged}
For any endpoint observables \(F_v(y_v)\), \(G_w(y_w)\) for which the
following quantities are integrable,
\[
\begin{aligned}
 &\mathbb E_\rho\operatorname{Cov}_{e,1}(F_v,G_w)\\
 &\quad=2\int_0^1\mathbb E_\rho
       \kappa_{e,s}(F_v,G_w,C_e)\,ds.                \tag{89.10}
\end{aligned}
\]
\end{proposition}

\begin{proof}
For every fixed \(\zeta\), conditional factorization at \(s=0\)
gives \(\operatorname{Cov}_{e,0}^{\zeta}(F_v,G_w)=0\).  Integrate
(89.6) from zero to one and then integrate against \(\rho\).  Fubini's
theorem applies under the stated integrability hypothesis.
\end{proof}

The remaining analytic acquisition is now stated at the correct
level: prove moments for the hybrid family \(\mu_s\), averaged over
\(\rho\), rather than conditional suprema over \(\zeta\).  Since
\(\rho\) is the interacting exterior marginal, those moments require a
comparison between \(\mu_s\) and \(\mu_1\); V87 alone only treats the
endpoint \(s=1\).

\begin{assessment}[What V89 changes]
\label{ass:v17v89-status}
V89 proves that the originally proposed pointwise conditional
interpolation cannot use uniform V87 moments.  It replaces that step by
the exact averaged hybrid identity (89.10), which preserves the
exterior law and exposes the precise missing estimate: uniform averaged
moments along \(s\in[0,1]\).
\end{assessment}

\begin{decision}[V90: cross-programme audit before hybrid moments]
\label{dec:v17v89-next}
Perform the required YM/NS audit.  Check in particular whether the new
YM cell bounds provide an interpolation comparison or finite-range
conditional normalization estimate applicable to (89.8).  If they do
not, return to the SM action and compare the hybrid normalizers by
Young domination of \(C_e\) against the retained endpoint diagonal
terms.
\end{decision}

\begin{firewall}[Claim boundary after V89]
\label{fire:v17v89}
No uniform hybrid moment estimate or Ward sign is yet proved.  Global
resolvent coercivity, environmental susceptibility, remainder-current
control, continuum measure, gauge cofinality, physical Einstein
action, and the Einstein--SM continuum limit remain open.
\end{firewall}

\section*{Version V89 append-only record}
\addcontentsline{toc}{section}{Version V89 append-only record}
The complete V88 source was retained byte for byte before this
continuation.  It had \(715679\) bytes, \(21797\) lines, and SHA--256
\[
 \texttt{591FEA9ADA2D5CEF1F5C68E784BCE5038AF664CC86DDBDFD3BADA431693F5676}.
\]
V89 rejected an exterior-nonuniform interpolation and replaced it by
an exact averaged hybrid representation.


\section{V90: mandatory YM/NS audit before the hybrid-moment estimate}
\label{sec:v17v90}

This version performs the five-version companion review required by
the programme protocol.  The live SM obligation from V89 is an
averaged moment estimate for the endpoint hybrid family (89.8).

\begin{audit}[Newest active companion notes]
\label{aud:v17v90-companions}
The newest active Yang--Mills note inspected was
\begin{center}\footnotesize
\path{Renewal_Yang_Mills_Mass_Gap_Research_Notes_V96.tex},
\end{center}
with \(742670\) bytes, \(20292\) physical lines, and SHA--256
\[
 \texttt{673501E061B3E6CEEDDDC5BAF4517CFC68165A9F2AB9A1FDB41D53122E8684E4}.
\]
Its newest theorem certifies nine new radius-\(1/8\) product charts
and proves positivity of the eighth direct-plus Wilson pencil through
third-coordinate depth \(2t\), uniformly on the already completed
first and second quarter directions.  Its terminal third layer,
complete eighth cell, full transverse torus, and continuum obligations
remain open.

The newest active Navier--Stokes note inspected was
\begin{center}\footnotesize
\path{Renewal_NS_Stability_Research_Notes_V26.tex},
\end{center}
with \(278283\) bytes, \(5319\) physical lines, and SHA--256
\[
 \texttt{82C887F8C2C69E4F28FAD7C3DC8DE5B9099CB5A4BF431EBA1FFF3E91935978AD}.
\]
Its latest results remain the exact pointwise rank-one norms for the
first two Oseen-kernel derivatives.  Its stability and regularity
theorem remains open.
\end{audit}

\begin{theorem}[V90 no-import decision]
\label{thm:v17v90-no-import}
Neither companion supplies an estimate that closes the V89 hybrid
moment obligation.  No theorem is imported at V90.
\end{theorem}

\begin{proof}
YM V96 is a coefficient-exact finite-dimensional momentum-symbol
cover.  Its probability-free chart certificates contain neither a
conditional Higgs normalizer nor an averaged Gibbs interpolation of
the form (89.8).  NS V26 estimates deterministic Oseen-kernel tensors
and has neither the Higgs variables nor the quartic conditional law.
Thus neither conclusion matches the antecedent of a moment or
normalizer estimate for \(\mu_s\).
\end{proof}

The audit does reinforce the claim boundary: finite-symbol positivity
in the gauge programme is approaching completion for one cell, but it
cannot substitute for the matter-sector probability estimate needed
by the SM environment resolvent.

\begin{decision}[V91: averaged radial integration by parts]
\label{dec:v17v90-next}
Retain the original exterior marginal \(\rho\) in (89.8).  Apply radial
integration by parts inside each normalized two-site conditional law.
Use positivity of the geometric edge square, Young's inequality for
the exterior linear fields, and the unconditional V87 moments of
\(\rho\).  First prove a uniform second moment for every
\(s\in[0,1]\), then bootstrap all integer endpoint moments needed in
the cumulant (89.10).
\end{decision}

\begin{firewall}[Claim boundary after V90]
\label{fire:v17v90}
This audit proves no new SM estimate.  Hybrid moments, the Ward sign,
global resolvent coercivity, environmental susceptibility,
remainder-current control, continuum measure, gauge cofinality,
physical Einstein action, and the Einstein--SM continuum limit remain
open.
\end{firewall}

\section*{Version V90 append-only record}
\addcontentsline{toc}{section}{Version V90 append-only record}
The complete V89 source was retained byte for byte before this
continuation.  It had \(724076\) bytes, \(22005\) lines, and SHA--256
\[
 \texttt{4D2CEE5FEE7557C04FFF3EE87861A30643D9F69B5B8E629D1539602016B558D4}.
\]
V90 audited YM V96 and NS V26 and imported no theorem.


\section{V91: averaged hybrid moments and a finite transported-radius error}
\label{sec:v17v91}

V89 identified the correct hybrid law \(\mu_s\), and V90 found no
companion theorem that estimates it.  We now obtain the estimate
directly.  Throughout this version we remain in the nonnegative-mass
branch used in V59--V61:
\[
 m_H^2\geq0,
 \qquad \lambda_H>0.                                \tag{91.1}
\]
All conformal and gauge variables remain fixed.  The exterior marginal
\(\rho\) in (89.8) is the marginal of the original \(s=1\) Higgs law.

\subsection{Exterior-field moments under the retained marginal}

For integers \(k\geq0\), define
\[
 M_0^G:=1,
 \qquad M_1^G:=q_HC_G,
 \qquad
 M_k^G:=(k^2+q_H)C_GM_{k-1}^G\quad(k\geq2).          \tag{91.2}
\]
Theorem~\ref{thm:v17v87-moments} gives
\(\mathbb E_\rho A_z^k\leq M_k^G\) for every exterior site \(z\).

Let \(B_x=\|b_x\|\), with \(b_x\) as in (88.2), and let
\(d_x^{\rm ext}\leq D-1\) be the number of exterior neighbors of
\(x\in\{v,w\}\).  Since each transported neighbor vector has norm at
most \(\chi_*\sqrt{A_z}\), for every \(\alpha\geq1\),
\[
 B_x^\alpha
 \leq(d_x^{\rm ext})^{\alpha-1}\chi_*^\alpha
      \sum_{z\sim x,\,z\ne e} A_z^{\alpha/2}.        \tag{91.3}
\]
Consequently, for every integer \(p\geq2\),
\[
\boxed{
 \mathbb E_\rho(B_v^{2p/3}+B_w^{2p/3})
 \leq\mathcal B_p,}                                 \tag{91.4}
\]
where, with \(k_p=\lceil p/3\rceil\),
\[
 \mathcal B_p
 :=2(D-1)^{2p/3}\chi_*^{2p/3}
       (M_{k_p}^G)^{p/(3k_p)}.                       \tag{91.5}
\]
Indeed, use (91.3) with \(\alpha=2p/3\) and Lyapunov's inequality
\(\mathbb EA^{p/3}\leq(\mathbb EA^{k_p})^{p/(3k_p)}\).

\subsection{Radial integration by parts along the hybrid}

Set
\[
 H_p(s):=\mathbb E_{\mu_s}(A_v^p+A_w^p).             \tag{91.6}
\]
The elementary sharp Young remainder needed below is
\[
 2Br^{2p-3}\leq\varepsilon r^{2p}
 +Y_p(\varepsilon)B^{2p/3},                          \tag{91.7}
\]
where
\[
 Y_p(\varepsilon)
 :=\frac3p\left(\frac{2p-3}{p\varepsilon}
          \right)^{(2p-3)/3}.                       \tag{91.8}
\]
To verify it, scale \(r=B^{1/3}t\) and maximize
\(2t^{2p-3}-\varepsilon t^{2p}\); the maximizer satisfies
\(t^3=(2p-3)/(p\varepsilon)\).

\begin{theorem}[Uniform averaged hybrid moments]
\label{thm:v17v91-hybrid-moments}
For every \(s\in[0,1]\),
\[
 H_2(s)\leq\mathcal H_2,                             \tag{91.9}
\]
where
\[
\boxed{
 \mathcal H_2
 :=\frac{2q_H}{3\lambda_H}
 +\frac{\mathcal B_2}
        {2^{4/3}\lambda_H^{4/3}}.}                  \tag{91.10}
\]
Put \(\mathcal H_0=2\),
\(\mathcal H_1=(2\mathcal H_2)^{1/2}\), and for
\(p\geq3\) define recursively
\[
\boxed{
 \mathcal H_p
 :=\frac{(q_H+2p-4)\mathcal H_{p-2}
          +2\chi_e\mathcal H_{p-1}
          +Y_p(\lambda_H)\mathcal B_p}
         {3\lambda_H}.}                             \tag{91.11}
\]
Then
\[
 \sup_{0\leq s\leq1}H_p(s)\leq\mathcal H_p
 \qquad(p=0,1,2,\ldots).                            \tag{91.12}
\]
Every constant is finite and independent of the finite volume.
\end{theorem}

\begin{proof}
Inside the conditional law \(\pi_{e,s}^{\zeta}\), radial integration
by parts at \(x\) gives
\[
 q_H=\mathbb E_{e,s}
 \{4\lambda_HA_x^2+2a_xA_x
   -2\langle b_x,y_x\rangle_{\mathbb R}-2sC_e\}.    \tag{91.13}
\]
Sum (91.13) over \(v,w\).  The edge-diagonal and transported terms
combine as
\[
 2\{D_{v,e}+D_{w,e}-2sC_e\}\geq0,                  \tag{91.14}
\]
because \(2|C_e|\leq D_{v,e}+D_{w,e}\).  All mass and
remaining diagonal contributions are nonnegative by (91.1).  Hence
\[
 4\lambda_HH_2(s)
 \leq2q_H+2\mathbb E_{\mu_s}
       (B_v\sqrt{A_v}+B_w\sqrt{A_w}).                \tag{91.15}
\]
Use (91.7) with \(p=2\), \(\varepsilon=\lambda_H\), then (91.4).
Since \(Y_2(\lambda_H)=3/(2^{4/3}\lambda_H^{1/3})\), absorbing
\(\lambda_HH_2(s)\) gives (91.9)--(91.10).

Cauchy--Schwarz yields
\(H_1(s)\leq(2H_2(s))^{1/2}\).  For \(p\geq3\), apply radial
integration by parts to the vector field
\(y_xA_x^{p-2}\).  Its divergence is
\((q_H+2p-4)A_x^{p-2}\).  After summing the two endpoints and
discarding nonnegative diagonal terms, one obtains
\[
\begin{aligned}
 4\lambda_HH_p(s)
 \leq{}&(q_H+2p-4)H_{p-2}(s)\\
 &+2\chi_eH_{p-1}(s)
 +2\mathbb E_{\mu_s}
   [B_vA_v^{p-3/2}+B_wA_w^{p-3/2}].                 \tag{91.16}
\end{aligned}
\]
For the transported term we used
\[
 A_v^{p-3/2}A_w^{1/2}+A_w^{p-3/2}A_v^{1/2}
 \leq A_v^{p-1}+A_w^{p-1},                          \tag{91.17}
\]
which is weighted arithmetic--geometric mean applied to each
summand.  Apply (91.7) with \(\varepsilon=\lambda_H\) separately at
the two endpoints, average the resulting exterior powers with
(91.4), and absorb \(\lambda_HH_p(s)\).  Induction gives
(91.11)--(91.12).
\end{proof}

The theorem repairs exactly the invalid conditional-supremum step
identified in V89: the unbounded \(B_x\) is never maximized.  Its
moments are instead taken under \(\rho\), where the V87 estimates are
available.

\subsection{The transported mixed-radius term}

Let
\[
 \mathcal P_3
 :=32(m_H^2)^3\mathcal H_3
   +256\lambda_H^3\mathcal H_6.                     \tag{91.18}
\]
Since \((u+v)^3\leq4(u^3+v^3)\) for \(u,v\geq0\),
\[
 \sup_s\mathbb E_{\mu_s}|P(A_x)|^3\leq\mathcal P_3.\tag{91.19}
\]

\begin{theorem}[Finite averaged transported-radius error]
\label{thm:v17v91-mixed-error}
The mixed conditional covariance in (88.9) satisfies
\[
\boxed{
 \left|\mathbb E_\rho
  \operatorname{Cov}_{e,1}(A_v,P(A_w))\right|
 \leq16\chi_e\mathcal H_3^{2/3}\mathcal P_3^{1/3}.} \tag{91.20}
\]
The same bound holds after interchanging \(v,w\).
\end{theorem}

\begin{proof}
For any conditional law and three observables,
\[
 |\kappa(F,G,H)|
 \leq8\|F\|_{L^3}\|G\|_{L^3}\|H\|_{L^3},          \tag{91.21}
\]
because centering increases each \(L^3\)-norm by at most a factor two.
Apply (91.21) conditionally, then H\"older's inequality to the exterior
average.  Theorem~\ref{thm:v17v91-hybrid-moments} and (91.19) give
\[
 \mathbb E_\rho|\kappa_{e,s}(A_v,P(A_w),C_e)|
 \leq8\chi_e\mathcal H_3^{2/3}\mathcal P_3^{1/3},  \tag{91.22}
\]
where
\(|C_e|^3\leq\chi_e^3A_v^{3/2}A_w^{3/2}\) and
Cauchy--Schwarz were used.  Integrate (91.22) in (89.10); its factor
two gives (91.20).
\end{proof}

Thus the two transported mixed terms in (88.9) now have an explicit,
volume-independent averaged error.  The estimate is intentionally
one-sided only after comparison with the positive self terms: V91
does not invent a lower bound for those terms.  The angular neighbor
piece and the exterior covariance in (88.5) also remain to be charged.

\begin{assessment}[What V91 closes]
\label{ass:v17v91-status}
V91 proves uniform averaged moments of every integer order for the
entire endpoint hybrid family and uses the sixth radial moment to
bound the transported mixed-radius covariance.  This validates the
averaged interpolation architecture and removes the exterior-field
nonuniformity found in V89.
\end{assessment}

\begin{decision}[V92: angular and exterior response ledger]
\label{dec:v17v91-next}
Differentiate the conditional endpoint expectations with respect to
the exterior vectors \(b_v,b_w\).  Use the tilt-uniform one-site gap
of V59 together with the hybrid moments of V91 to bound the angular
neighbor contribution after exterior averaging.  Then apply the law
of total covariance to express the second term in (88.5) through the
V61 influence kernel, keeping its radius-one part separate from the
summable V83 tail.
\end{decision}

\begin{firewall}[Claim boundary after V91]
\label{fire:v17v91}
The transported error bound has not been shown smaller than the
positive radial self contribution.  The angular and exterior Ward
terms, the Ward sign, global resolvent coercivity, environmental
susceptibility, remainder-current control, continuum measure, gauge
cofinality, physical Einstein action, and the Einstein--SM continuum
limit remain open.
\end{firewall}

\section*{Version V91 append-only record}
\addcontentsline{toc}{section}{Version V91 append-only record}
The complete V90 source was retained byte for byte before this
continuation.  It had \(727695\) bytes, \(22093\) lines, and SHA--256
\[
 \texttt{200D85BD6E934068FE2F86CB3E4450D95C89FFFD77E2248234C0527F9BE8AFC1}.
\]
V91 proved uniform averaged hybrid moments and an explicit
transported mixed-radius covariance bound.


\section{V92: angular endpoint control and the radius-one exterior remainder}
\label{sec:v17v92}

V91 bounds the transported mixed-radius terms after exterior
averaging.  We now incorporate the non-edge diagonal hopping into the
positive radial force and control the remaining angular linear term by
a two-site version of the V61 gap.  The law-of-total-covariance
remainder is then located precisely in space.

\subsection{A uniform gap for the conditioned endpoint pair}

Put
\[
 m_e:=2C_0\chi_e,
 \qquad
 C_{2,e}:=\frac{C_0}{(1-m_e)^2}.                    \tag{92.1}
\]
Since the full influence matrix is symmetric and contains \(m_e\) as
an off-diagonal entry, the V61 hypothesis \(r_H<1\) implies
\(m_e\leq r_H<1\).

\begin{lemma}[Tilt-uniform two-endpoint Poincar\'e inequality]
\label{lem:v17v92-two-site-gap}
For every exterior configuration and every \(s\in[0,1]\), the
conditional law \(\pi_{e,s}^{\zeta}\) from (89.8) satisfies
\[
 \operatorname{Var}_{e,s}^{\zeta}(F)
 \leq C_{2,e}\,
 \mathbb E_{e,s}^{\zeta}
 \bigl(\|\nabla_{y_v}F\|^2+\|\nabla_{y_w}F\|^2\bigr).              \tag{92.2}
\]
The constant is independent of the exterior linear fields
\(b_v,b_w\).
\end{lemma}

\begin{proof}
Conditioning in either endpoint produces the tilted quartic one-site
law of V59, with Poincar\'e constant \(C_0\).  The only mixed Hessian
between the two endpoints is the transported hopping term and has
operator norm at most \(2s\chi_e\).  Thus the two-by-two influence
matrix has zero diagonal and off-diagonal entries at most
\(2sC_0\chi_e\leq m_e\).  Apply the two-site specialization of the
martingale sweeping proof of Theorem~\ref{thm:v17v61-tensor}; the
Neumann resolvent has norm at most \((1-m_e)^{-1}\), giving (92.2).
Tilt independence follows from the one-site theorem.
\end{proof}

\subsection{Radial repacking and its one-sided estimate}

Define the exterior diagonal coefficient and repacked radial force
\[
 a_x^\star:=m_H^2+
 \sum_{\substack{f\ni x\\f\ne e}}g_{x,f},
 \qquad
 P_x^\star(A):=2a_x^\star A+4\lambda_HA^2.          \tag{92.3}
\]
Also put
\[
 \mathcal L_e
 :=-2\langle b_v,y_v\rangle_{\mathbb R}
    +2\langle b_w,y_w\rangle_{\mathbb R}.           \tag{92.4}
\]
Then (86.7) is exactly
\[
 \mathcal K_e=P_v^\star(A_v)-P_w^\star(A_w)
                    +\mathcal L_e.                  \tag{92.5}
\]
On \(\Omega_{e,R}^{\rm star}\), let
\[
 a_R^\star:=m_H^2+(D-1)\chi_*e^R,                  \tag{92.6}
\]
and define
\[
 \mathcal P_{3,R}^\star
 :=32(a_R^\star)^3\mathcal H_3
   +256\lambda_H^3\mathcal H_6,                    \tag{92.7}
\]
\[
 \mathcal E_{\rm mix}(R)
 :=16\chi_e\mathcal H_3^{2/3}
              (\mathcal P_{3,R}^\star)^{1/3}.       \tag{92.8}
\]

\begin{lemma}[Repacked positive self terms]
\label{lem:v17v92-repacked-self}
For every fixed exterior configuration,
\[
 \operatorname{Cov}_{e,1}(A_x,P_x^\star(A_x))\geq0,
 \qquad x\in\{v,w\}.                               \tag{92.9}
\]
Moreover, on the bounded star sector,
\[
 \left|\mathbb E_\rho
  \operatorname{Cov}_{e,1}(A_v,P_w^\star(A_w))\right|
 \leq\mathcal E_{\rm mix}(R),                       \tag{92.10}
\]
and the same estimate holds with \(v,w\) interchanged.
\end{lemma}

\begin{proof}
The derivative
\(2a_x^\star+8\lambda_HA\) is nonnegative, so the independent-copy
identity (88.7) proves (92.9).  For (92.10), repeat the transported
interpolation (89.10) with \(P_w^\star\) in place of \(P\).  On the
star sector \(0\leq a_w^\star\leq a_R^\star\), and
\[
 \sup_s\mathbb E_{\mu_s}|P_w^\star(A_w)|^3
 \leq\mathcal P_{3,R}^\star.                       \tag{92.11}
\]
The cumulant estimate (91.21), the hybrid moments (91.12), and the
factor two in (89.10) give exactly (92.8)--(92.10).
\end{proof}

Let
\[
 \mathcal K_e^{\rm rad}
 :=P_v^\star(A_v)-P_w^\star(A_w).                   \tag{92.12}
\]
Bilinearity gives the exact expansion
\[
\begin{aligned}
 \operatorname{Cov}_{e,1}(q_e,\mathcal K_e^{\rm rad})
 ={}&g_{v,e}\operatorname{Cov}_{e,1}(A_v,P_v^\star(A_v))
    +g_{w,e}\operatorname{Cov}_{e,1}(A_w,P_w^\star(A_w))\\
 &-g_{v,e}\operatorname{Cov}_{e,1}(A_v,P_w^\star(A_w))
  -g_{w,e}\operatorname{Cov}_{e,1}(A_w,P_v^\star(A_v)).            \tag{92.13}
\end{aligned}
\]
Therefore Lemma~\ref{lem:v17v92-repacked-self} implies
\[
 \mathbb E_\rho
 \operatorname{Cov}_{e,1}(q_e,\mathcal K_e^{\rm rad})
 \geq-(g_{v,e}+g_{w,e})\mathcal E_{\rm mix}(R).     \tag{92.14}
\]

\subsection{The angular neighbor term}

The endpoint gradients of (92.4) obey
\[
 \|\nabla_{v,w}\mathcal L_e\|^2
 =4(B_v^2+B_w^2),                                   \tag{92.15}
\]
whereas
\[
 \mathbb E_{e,1}\|\nabla_{v,w}q_e\|^2
 =4\{g_{v,e}^2\mathbb E_{e,1}A_v
       +g_{w,e}^2\mathbb E_{e,1}A_w\}.              \tag{92.16}
\]

\begin{theorem}[Averaged angular Ward bound]
\label{thm:v17v92-angular}
Under the V61 influence margin,
\[
\boxed{
 \left|\mathbb E_\rho
   \operatorname{Cov}_{e,1}(q_e,\mathcal L_e)\right|
 \leq4C_{2,e}\sqrt{q_HC_Gw_e\mathcal B_3}.}         \tag{92.17}
\]
\end{theorem}

\begin{proof}
The covariance form of (92.2), followed by (92.15)--(92.16), gives
\[
 |\operatorname{Cov}_{e,1}(q_e,\mathcal L_e)|
 \leq2C_{2,e}
 \bigl(\mathbb E_{e,1}\|\nabla q_e\|^2\bigr)^{1/2}
 (B_v^2+B_w^2)^{1/2}.                               \tag{92.18}
\]
Average and apply Cauchy--Schwarz.  Equations (87.14) and (91.4) with
\(p=3\) bound the two resulting factors by
\(4q_HC_Gw_e\) and \(\mathcal B_3\), respectively, proving (92.17).
\end{proof}

Combining (92.5), (92.14), and (92.17) yields the first genuinely
one-sided estimate for the endpoint-conditioned Ward contribution.

\begin{corollary}[First-level Ward floor]
\label{cor:v17v92-first-level}
On \(\Omega_{e,R}^{\rm star}\),
\[
\boxed{
 \frac12\mathbb E_\rho
  \operatorname{Cov}_{e,1}(q_e,\mathcal K_e)
 \geq-\mathcal W_e^{(1)}(R),}                       \tag{92.19}
\]
where
\[
 \mathcal W_e^{(1)}(R)
 :=\frac12(g_{v,e}+g_{w,e})\mathcal E_{\rm mix}(R)
   +2C_{2,e}\sqrt{q_HC_Gw_e\mathcal B_3}.           \tag{92.20}
\]
\end{corollary}

\begin{proof}
Add the lower radial estimate (92.14) and the negative of the absolute
angular estimate (92.17), then divide by two.
\end{proof}

\subsection{The exterior covariance is local, not a V83 tail}

Let
\[
 S_e:=\{z:z\sim v\text{ or }z\sim w\}\setminus\{v,w\}.             \tag{92.21}
\]
By the nearest-neighbor Markov property of the Higgs action, both
\(\bar q_e=\mathbb E_e q_e\) and
\(\overline{\mathcal K}_e=\mathbb E_e\mathcal K_e\) are functions
only of the exterior variables indexed by \(S_e\).  In particular,
their supports coincide rather than being spatially separated.

\begin{proposition}[Radius-one status of the second Ward level]
\label{prop:v17v92-radius-one}
The second term in (88.5),
\[
 \frac12\operatorname{Cov}_\rho
  (\bar q_e,\overline{\mathcal K}_e),                \tag{92.22}
\]
is a radius-one exterior-star covariance.  Neither the positive
distance hypothesis of Lemma~\ref{lem:v17v82-tail} nor the distant-edge
hypothesis behind (83.16) applies to it.  Thus it cannot be charged to
\(\Xi_L\) with \(L\geq1\) at the endpoint-conditioning level.
\end{proposition}

\begin{proof}
In the conditional potential (88.1), exterior variables enter only
through \(b_v,b_w\), each formed from the neighbors in \(S_e\).
The observable \(q_e\) has no other exterior dependence, while
\(\mathcal K_e\) has the same neighbor dependence explicitly through
(92.3)--(92.5).  Hence both conditional expectations are
\(\sigma(y_z:z\in S_e)\)-measurable.  Their two support sets have
distance zero.  The tail bounds (82.7) and (83.16) become exponentially
small only between support sets at positive and increasing distance,
so they give no tail factor here.
\end{proof}

This corrects the tentative V91 plan: invoking V83 directly on
(92.22) would confuse a same-star covariance with a distant current
covariance.  To create a genuine tail one must condition on the
exterior of a larger ball and propagate the dependence from the
central star to that ball's boundary.

\begin{assessment}[What V92 closes]
\label{ass:v17v92-status}
V92 proves a tilt-uniform two-endpoint gap, repacks every non-edge
diagonal hopping term into a positive radial self covariance, bounds
the transported and angular errors after exterior averaging, and
obtains the explicit first-level Ward floor (92.19).  It also proves
that the remaining second-level covariance is local at radius one and
is not covered by the V83 distant tail.
\end{assessment}

\begin{decision}[V93: expanding-block martingale ledger]
\label{dec:v17v92-next}
Let \(\Lambda_L\) be the radius-\(L\) neighborhood of the endpoint
star and condition on its exterior.  Derive the exact total-covariance
identity for \(q_e,\mathcal K_e\) at every \(L\), then telescope the
change from \(L\) to \(L+1\) by boundary martingale increments.  Use
the V61 sweeping matrix to assign at least \(L\) influence steps to
the terminal exterior remainder.  Keep the finite inner martingale
increments explicit; do not label them as a tail without a sign or
absolute summability proof.
\end{decision}

\begin{firewall}[Claim boundary after V92]
\label{fire:v17v92}
The first-level Ward floor does not control the radius-one exterior
covariance.  The Ward sign, global resolvent coercivity,
environmental susceptibility, remainder-current control, continuum
measure, gauge cofinality, physical Einstein action, and the
Einstein--SM continuum limit remain open.
\end{firewall}

\section*{Version V92 append-only record}
\addcontentsline{toc}{section}{Version V92 append-only record}
The complete V91 source was retained byte for byte before this
continuation.  It had \(735997\) bytes, \(22344\) lines, and SHA--256
\[
 \texttt{862F591B9DD14AB61EB48C31F212EB5818E2633EC5157CDF847E8FF32B001791}.
\]
V92 proved the averaged angular bound and isolated the remaining Ward
covariance as a radius-one exterior term.


\section{V93: expanding-block martingales and closure of the exterior Ward bound}
\label{sec:v17v93}

V92 leaves the exterior covariance in (88.5) on the radius-one star.
This version expands the conditioned block shell by shell.  The first
shell is estimated by differentiating the two-endpoint conditional
means; all later shells carry explicit powers of the V61 influence
matrix.

\subsection{Reverse-martingale covariance ledger}

Let \(C_e=\{v,w\}\) denote the endpoint set in this subsection only,
and define
\[
 \Lambda_j:=\{x:\operatorname{dist}(x,C_e)\leq j+1\},
 \qquad j\geq-1.                                    \tag{93.1}
\]
Thus \(\Lambda_{-1}=C_e\) and \(\Lambda_0\) contains the support of
\(\mathcal K_e\).  Put
\[
 \mathscr G_j:=\sigma(y_x:x\notin\Lambda_j),
 \quad
 Q_j:=\mathbb E(q_e\mid\mathscr G_j),
 \quad
 K_j:=\mathbb E(\mathcal K_e\mid\mathscr G_j),      \tag{93.2}
\]
and
\[
 \Delta_j^Q:=Q_j-Q_{j+1},
 \qquad
 \Delta_j^K:=K_j-K_{j+1}.                           \tag{93.3}
\]

\begin{lemma}[Exact shell telescoping]
\label{lem:v17v93-telescope}
For every \(L\geq0\),
\[
\boxed{
 \operatorname{Cov}(Q_{-1},K_{-1})
 =\sum_{j=-1}^{L-1}
    \mathbb E(\Delta_j^Q\Delta_j^K)
  +\operatorname{Cov}(Q_L,K_L).}                   \tag{93.4}
\]
In a finite graph the final covariance vanishes once
\(\Lambda_L\) is the whole graph.
\end{lemma}

\begin{proof}
The sigma-fields decrease:
\(\mathscr G_{j+1}\subset\mathscr G_j\).  The tower property gives
\(\mathbb E(\Delta_j^Q\mid\mathscr G_{j+1})=0\), and similarly for
\(K\).  Therefore the two cross terms vanish in
\[
 \operatorname{Cov}(Q_j,K_j)
 =\operatorname{Cov}(Q_{j+1}+\Delta_j^Q,
                     K_{j+1}+\Delta_j^K),           \tag{93.5}
\]
and
\[
 \operatorname{Cov}(Q_j,K_j)-
 \operatorname{Cov}(Q_{j+1},K_{j+1})
 =\mathbb E(\Delta_j^Q\Delta_j^K).                  \tag{93.6}
\]
Sum (93.6).  Conditional expectation on the trivial sigma-field is a
constant, proving the finite-graph terminal statement.
\end{proof}

The \(j=-1\) increment integrates the radius-one shell singled out in
V92.  It will be kept separate.  Every \(j\geq0\) increment is a true
distant-shell term.

\subsection{Swept conditional means at a distant boundary}

For a finite vertex set \(\Lambda\), write
\(P_\Lambda F=\mathbb E(F\mid y_{\Lambda^c})\).  Recall the gradient
vector \(a(F)\) from (82.1).

\begin{lemma}[Exterior gradient after a block sweep]
\label{lem:v17v93-block-sweep}
Suppose \(F\) is supported in \(S\subset\Lambda\) and
\(\operatorname{dist}(S,\Lambda^c)\geq\ell\).  Under \(r_H<1\),
\[
\boxed{
 \left(
  \mathbb E\sum_{z\notin\Lambda}
  \|\nabla_{y_z}P_\Lambda F\|^2
 \right)^{1/2}
 \leq\frac{r_H^\ell}{1-r_H}\,|a(F)\|_{\ell^2}.}   \tag{93.7}
\]
Consequently,
\[
 \operatorname{Var}(P_\Lambda F)
 \leq C_G\frac{r_H^{2\ell}}{(1-r_H)^2}
          \mathbb E\|\nabla_yF\|^2.                 \tag{93.8}
\]
\end{lemma}

\begin{proof}
Expose the variables of \(\Lambda\) in the martingale-sweeping proof
of V61.  A gradient transferred from a site to a neighboring
unexposed site pays the corresponding entry of \(M\), by (61.3).
After every site of \(\Lambda\) is integrated, the exterior gradient
vector is bounded in \(\ell^2\) by
\[
 1_{\Lambda^c}(I-M)^{-1}1_Sa(F).                    \tag{93.9}
\]
Because \(M\) is nearest-neighbor, all powers below \(\ell\) vanish
between \(S\) and \(\Lambda^c\).  Hence
\[
 \|1_{\Lambda^c}(I-M)^{-1}1_S\|
 \leq\sum_{n=\ell}^\infty r_H^n
 =\frac{r_H^\ell}{1-r_H}.                           \tag{93.10}
\]
This proves (93.7).  The function \(P_\Lambda F\) depends only on
exterior sites.  Apply the global Poincar\'e inequality (61.10) to it
and then (93.7), proving (93.8).
\end{proof}

For \(j\geq0\), the support of \(q_e\) is at least \(j+2\) steps from
\(\Lambda_j^c\), and the support of \(\mathcal K_e\) is at least
\(j+1\) steps away.  We use the common lower distance \(j+1\).
Equations (87.10) and (87.14) then give
\[
\begin{aligned}
 \operatorname{Var}(Q_j)
 &\leq4q_HC_G^2w_e
       \frac{r_H^{2j+2}}{(1-r_H)^2},\\
 \operatorname{Var}(K_j)
 &\leq C_GG_{\mathcal K}(R)
       \frac{r_H^{2j+2}}{(1-r_H)^2}.                \tag{93.11}
\end{aligned}
\]
Also
\[
 \mathbb E(\Delta_j^Q)^2
 =\operatorname{Var}(Q_j)-\operatorname{Var}(Q_{j+1})
 \leq\operatorname{Var}(Q_j),                      \tag{93.12}
\]
and analogously for \(K\).

\begin{corollary}[Summable distant-shell ledger]
\label{cor:v17v93-distant-shells}
On \(\Omega_{e,R}^{\rm star}\),
\[
\boxed{
 \sum_{j=0}^{\infty}
 \left|\mathbb E(\Delta_j^Q\Delta_j^K)\right|
 \leq
 \frac{2C_Gr_H^2}{(1-r_H)^2(1-r_H^2)}
 \sqrt{q_HC_Gw_eG_{\mathcal K}(R)}.}                \tag{93.13}
\]
The same expression bounds the right side uniformly over all finite
volumes, with the sum stopped when the graph is exhausted.
\end{corollary}

\begin{proof}
Cauchy--Schwarz, (93.11), and (93.12) bound the \(j\)-th summand by
\[
 \frac{2C_Gr_H^{2j+2}}{(1-r_H)^2}
 \sqrt{q_HC_Gw_eG_{\mathcal K}(R)}.                 \tag{93.14}
\]
Sum the geometric series.
\end{proof}

\subsection{The first exterior shell}

Let
\[
 \Sigma_e^\chi
 :=\sum_{x\in\{v,w\}}
   \sum_{\substack{f\ni x\\f\ne e}}\chi_f^2
 \leq2(D-1)\chi_*^2.                               \tag{93.15}
\]
The endpoint-conditioned means \(Q_{-1},K_{-1}\) are functions of the
radius-one shell \(S_e\) from (92.21).

\begin{lemma}[First-shell gradient budgets]
\label{lem:v17v93-first-gradient}
On \(\Omega_{e,R}^{\rm star}\),
\[
 \mathbb E\sum_{z\in S_e}\|\nabla_{y_z}Q_{-1}\|^2
 \leq\Gamma_Q,                                      \tag{93.16}
\]
\[
 \mathbb E\sum_{z\in S_e}\|\nabla_{y_z}K_{-1}\|^2
 \leq\Gamma_K,                                      \tag{93.17}
\]
where
\[
 \Gamma_Q
 :=32C_{2,e}^2\Sigma_e^\chi q_HC_Gw_e,             \tag{93.18}
\]
\[
 \Gamma_K
 :=16\Sigma_e^\chi
       \{q_HC_G+C_{2,e}^2G_{\mathcal K}(R)\}.       \tag{93.19}
\]
\end{lemma}

\begin{proof}
Fix an incidence from \(x\in\{v,w\}\) to a shell site \(z\), with
hopping coefficient \(\chi_f\).  Differentiating a normalized
two-endpoint conditional expectation in a unit \(y_z\)-direction
gives
\[
 D_{y_z}\mathbb E_eF
 =\mathbb E_eD_{y_z}F
  -\operatorname{Cov}_e(F,D_{y_z}V_e).              \tag{93.20}
\]
The endpoint gradient of the nonconstant part of
\(D_{y_z}V_e\) has norm at most \(2\chi_f\).  By the two-site
Poincar\'e inequality (92.2),
\[
 \|D_{y_z}Q_{-1}\|
 \leq2C_{2,e}\chi_f
       (\mathbb E_e\|\nabla_{v,w}q_e\|^2)^{1/2}.    \tag{93.21}
\]
For \(F=\mathcal K_e\), its direct derivative has conditional
squared norm at most \(4\chi_f^2\mathbb E_eA_x\), while the covariance
term is at most
\(2C_{2,e}\chi_f
(\mathbb E_e\|\nabla_{v,w}\mathcal K_e\|^2)^{1/2}\).

A shell site can be incident to both endpoints.  Bounding the square
of the resulting sum by twice the sum of incidence squares, summing
all incidences, and averaging gives the conservative factors in
(93.18)--(93.19).  Finally use (87.2), (87.10), and (87.14).
\end{proof}

Conditioned on \(\mathscr G_0\), the shell variables again form a
subsystem of the V61 Gibbs specification.  Its influence submatrix has
norm at most \(r_H\), so its conditional Poincar\'e constant is at most
\(C_G\).  Therefore
\[
 \mathbb E(\Delta_{-1}^Q)^2\leq C_G\Gamma_Q,
 \qquad
 \mathbb E(\Delta_{-1}^K)^2\leq C_G\Gamma_K.        \tag{93.22}
\]

\begin{theorem}[Closed exterior Ward bound]
\label{thm:v17v93-exterior}
On \(\Omega_{e,R}^{\rm star}\),
\[
\boxed{
 |\operatorname{Cov}(Q_{-1},K_{-1})|
 \leq C_G\sqrt{\Gamma_Q\Gamma_K}
 +\frac{2C_Gr_H^2}{(1-r_H)^2(1-r_H^2)}
  \sqrt{q_HC_Gw_eG_{\mathcal K}(R)}.}               \tag{93.23}
\]
This estimate is uniform in the finite volume.
\end{theorem}

\begin{proof}
Use Lemma~\ref{lem:v17v93-telescope} until the terminal covariance
vanishes.  Bound the first increment by Cauchy--Schwarz and (93.22),
and the remaining increments by
Corollary~\ref{cor:v17v93-distant-shells}.
\end{proof}

Combining (88.5), Corollary~\ref{cor:v17v92-first-level}, and
Theorem~\ref{thm:v17v93-exterior} removes the unnamed second-level
term.

\begin{corollary}[Explicit block-expanded Ward floor]
\label{cor:v17v93-ward-floor}
On \(\Omega_{e,R}^{\rm star}\),
\[
\boxed{
 h_e^{\rm Ward}\geq-\mathcal W_e^{\rm block}(R),}   \tag{93.24}
\]
where
\[
\begin{aligned}
 \mathcal W_e^{\rm block}(R):={}&
 \mathcal W_e^{(1)}(R)
 +\frac{C_G}{2}\sqrt{\Gamma_Q\Gamma_K}\\
 &+\frac{C_Gr_H^2}{(1-r_H)^2(1-r_H^2)}
  \sqrt{q_HC_Gw_eG_{\mathcal K}(R)}.                \tag{93.25}
\end{aligned}
\]
\end{corollary}

\begin{proof}
The first term of (88.5) is bounded below by (92.19).  Bound half the
absolute value of its second term by half of (93.23).
\end{proof}

The new floor is not asserted to be positive or smaller than the V87
two-sided floor.  Its purpose is structural: every negative term now
has a local first-shell constant or a summable influence-tail factor.
There is no remaining unspecified exterior covariance.

\begin{assessment}[What V93 closes]
\label{ass:v17v93-status}
V93 proves the exact expanding-block martingale ledger, exponential
decay of the terminal conditional means, absolute summability of all
distant shells, a direct first-shell response bound, and the explicit
Ward floor (93.24).  The earlier radius-one exterior placeholder has
therefore been converted into checkable constants.
\end{assessment}

\begin{decision}[V94: insert the block Ward floor into the one-form operator]
\label{dec:v17v93-next}
Insert (93.24) into the Ward-renormalized Hessian (86.16), retain the
V83 off-diagonal covariance tail by distance, and derive inner and
outer block coercivity constants for the Schur complement (83.2).
Compare the result with the simpler V87 floor and keep the better
constant edge by edge.  Determine whether the bare
\(2\kappa_g\) term closes a nonempty scale-independent parameter
region without using the canonical mismatch coefficient \(J(a)\).
\end{decision}

\begin{firewall}[Claim boundary after V93]
\label{fire:v17v93}
The explicit Ward floor is conditional on the high-temperature and
bounded-star hypotheses and has not been shown positive.  Global
resolvent coercivity, environmental susceptibility, remainder-current
control, continuum measure, gauge cofinality, physical Einstein
action, and the Einstein--SM continuum limit remain open.
\end{firewall}

\section*{Version V93 append-only record}
\addcontentsline{toc}{section}{Version V93 append-only record}
The complete V92 source was retained byte for byte before this
continuation.  It had \(745846\) bytes, \(22627\) lines, and SHA--256
\[
 \texttt{411465B1CD322856CA78DFA908A1FC9C0CE24B1F1918D4B35B2E4B54CB80ACCC}.
\]
V93 closed the exterior Ward estimate by an expanding-block
martingale and influence-tail ledger.


\section{V94: Ward-improved one-form coercivity and Schur-complement closure}
\label{sec:v17v94}

V93 converts the entire Ward diagonal into explicit local and
influence-tail constants.  This version inserts that result into the
one-form Hessian, combines it with the V83 spatial split of the
off-diagonal current covariance, and obtains checkable inner and outer
Schur floors on a bounded conformal sector.

\subsection{The better Ward diagonal constant}

For an edge \(e\), define
\[
 \mathcal W_e^{\rm best}(R)
 :=\min\{H_{\rm Ward}(R),\mathcal W_e^{\rm block}(R)\},             \tag{94.1}
\]
and
\[
 \mathcal W_*^{\rm best}(R)
 :=\sup_e\mathcal W_e^{\rm best}(R).               \tag{94.2}
\]
Both entries in (94.1) are rigorous lower-form constants: the first is
(87.16), and the second is (93.24).  Hence
\[
 h_e^{\rm Ward}\geq-\mathcal W_e^{\rm best}(R)      \tag{94.3}
\]
on \(\Omega_{e,R}^{\rm star}\).  Taking the minimum does not assume
that either derivation is sharper in every parameter regime.

\subsection{Near/far splitting of the off-diagonal covariance}

Let
\[
 A_{\rm cur}:=4q_HC_0C_G,                           \tag{94.4}
\]
and, for an integer \(L\geq1\), set
\[
 \mathcal N_L
 :=\frac{A_{\rm cur}C_D}{(1-r_H)^2}
   \sum_{n=0}^{L-1}(n+1)\{(D-1)r_H\}^n.            \tag{94.5}
\]
Define
\[
 K_{<L}:=\mathcal N_L.                              \tag{94.6}
\]

\begin{lemma}[Normalized near-current matrix]
\label{lem:v17v94-near}
Let \(\mathsf C_H^{\rm off,<L}\) retain the entries of
\(\mathsf C_H^{\rm off}\) with edge distance below \(L\).  Then
\[
 \left\|
 \operatorname{diag}(w)^{-1/2}
 \mathsf C_H^{\rm off,<L}
 \operatorname{diag}(w)^{-1/2}
 \right\|\leq K_{<L}.                               \tag{94.7}
\]
\end{lemma}

\begin{proof}
Normalize (83.12), use \(N_n\leq C_D(D-1)^n\), and then use
(83.17).  The absolute row and column sums below distance \(L\) are
at most \(\mathcal N_L\).  Schur's test proves (94.7).
\end{proof}

Under the strengthened spatial margin
\(\alpha_H=(D-1)r_H<1\), the complementary matrix at edge distance at
least \(L\) has normalized norm at most \(\Xi_L\) by (83.16).  Put
\[
 K_{\rm off}^{\rm best}(L)
 :=\min\{2K_{\rm cur},K_{<L}+\Xi_L\}.                \tag{94.8}
\]

\begin{corollary}[Best available off-diagonal form]
\label{cor:v17v94-off}
For every real edge vector \(u\),
\[
\boxed{
 |u^{\mathsf T}\mathsf C_H^{\rm off}u|
 \leq K_{\rm off}^{\rm best}(L)
        \sum_ew_eu_e^2.}                            \tag{94.9}
\]
\end{corollary}

\begin{proof}
Split the matrix at distance \(L\), apply (94.7) and (83.16), and use
the triangle inequality.  Independently, (87.19) bounds the full
off-diagonal matrix by \(2K_{\rm cur}\).  Take the smaller bound.
\end{proof}

The minimum in (94.8) is important.  Spatial splitting is retained
only when its finite near ledger plus its tail is better than the
global normalized estimate.

\subsection{A scale-independent one-form margin}

On \(\Omega_R\), put
\[
 w_R:=2\chi_*^2\cosh(2R),                            \tag{94.10}
\]
and
\[
\boxed{
 \mu_{R,L}^{\rm Ward}
 :=2\kappa_g-2D\left\{
 \mathcal W_*^{\rm best}(R)
 +w_RK_{\rm off}^{\rm best}(L)
 \right\}.}                                        \tag{94.11}
\]

\begin{theorem}[Ward-improved global one-form coercivity]
\label{thm:v17v94-coercivity}
Assume the V61 margin, the strengthened V83 spatial margin, and a
fixed environment in \(\Omega_R\).  For every vertex one-form
\(\omega\),
\[
\boxed{
 \mathcal Q_U[\omega]
 \geq\mu_{R,L}^{\rm Ward}\|\omega\|_2^2.}           \tag{94.12}
\]
In particular, if
\[
\boxed{
 \kappa_g>D\left\{
 \mathcal W_*^{\rm best}(R)
 +2\chi_*^2\cosh(2R)K_{\rm off}^{\rm best}(L)
 \right\},}                                        \tag{94.13}
\]
then the one-form Hessian is uniformly positive with floor
\(\mu_{R,L}^{\rm Ward}>0\).
\end{theorem}

\begin{proof}
Set \(u=B\omega\).  The Ward decomposition (86.16), (94.3), and
(94.9) give
\[
 u^{\mathsf T}\mathsf H_Hu
 \geq-\left\{
 \mathcal W_*^{\rm best}(R)
 +w_RK_{\rm off}^{\rm best}(L)
 \right\}\|u\|^2.                                  \tag{94.14}
\]
Use \(\|B\omega\|^2\leq2D\|\omega\|^2\).  In (81.5), retain the
bare term \(2\kappa_g\|\omega\|^2\) and discard the nonnegative
diffusion, mismatch, and other mean hopping terms.  This proves
(94.12), and (94.13) is exactly positivity of (94.11).
\end{proof}

Every constant on the right of (94.13) depends on the fixed local
parameters, \(R,L,D\), and the high-temperature margins, but not on
the number of lattice sites and not on the mismatch coefficient
\(J(a)\).  Therefore (94.13) defines a nonempty, scale-independent
mathematical parameter region: after fixing all other admissible
constants, it holds for sufficiently large \(\kappa_g\).  The source
manuscripts do not establish that the physical running couplings lie
in this region.

\subsection{Inner and outer Schur floors}

Let \(P,Q\) be any spatial block projections used in (83.1).  Write
\(\mathsf A,\mathsf C,\mathsf D\) and \(\mathsf S_P\) as in
(83.1)--(83.2).

\begin{theorem}[Uniform block and Schur coercivity]
\label{thm:v17v94-schur}
If \(\mu:=\mu_{R,L}^{\rm Ward}>0\), then
\[
 \mathsf A\succeq\mu P,
 \qquad
 \mathsf D\succeq\mu Q,
 \qquad
 \mathsf S_P\succeq\mu P.                          \tag{94.15}
\]
Consequently Theorem~\ref{thm:v17v83-schur} applies with
\[
 m_{\rm in}=m_{\rm out}=\mu.                        \tag{94.16}
\]
For an exact-gradient source \(g=g_P+g_Q\),
\[
 \langle g,(\mathscr H_U^{(1)})^{-1}g\rangle
 \leq
 \frac{(\|g_P\|+\|\mathsf C\|\mu^{-1}\|g_Q\|)^2
       +\|g_Q\|^2}{\mu}.                            \tag{94.17}
\]
\end{theorem}

\begin{proof}
Compression of (94.12) gives the first two inequalities in (94.15).
Since the full operator is at least \(\mu I\), its inverse is at most
\(\mu^{-1}I\).  The block inverse identity (83.6) gives
\[
 \mathsf S_P^{-1}
 =P(\mathscr H_U^{(1)})^{-1}P
 \preceq\mu^{-1}P,                                  \tag{94.18}
\]
and hence \(\mathsf S_P\succeq\mu P\).  Substitute these floors in
(83.4) to obtain (94.17).
\end{proof}

Theorem~\ref{thm:v17v94-schur} closes the V83 inner/outer coercivity
obligation on the declared central conformal sector and parameter
branch.  It does not control configurations outside \(\Omega_R\),
and (94.17) still contains the off-block norm when a source has an
exterior component.

\begin{assessment}[What V94 closes]
\label{ass:v17v94-status}
V94 combines the sharper of two rigorous Ward floors with a
near/far current-covariance split, proves a volume-independent
one-form floor, identifies a nonempty scale-independent sufficient
parameter region, and supplies equal positive inner, outer, and Schur
floors for every spatial block.  This is the first completed
conditional conformal-resolvent coercivity theorem in the current
lineage.
\end{assessment}

\begin{decision}[V95: mandatory audit and central-sector response]
\label{dec:v17v94-next}
First perform the required YM/NS audit.  Then insert the Schur floor
(94.15) into the response identity (79.6) for a local Higgs-norm test.
Use the V82 source tail and optimize the block radius, retaining the
central-sector restriction explicitly.  Do not extend the estimate to
the full conformal measure until an outside-sector probability or
coercivity theorem has been proved.
\end{decision}

\begin{firewall}[Claim boundary after V94]
\label{fire:v17v94}
The coercivity and Schur bounds hold only under the stated
high-temperature, spatial-margin, bounded-conformal-sector, and
large-enough-\(\kappa_g\) conditions.  Full environmental
susceptibility, outside-sector control, remainder-current control,
continuum measure, gauge cofinality, physical Einstein action, and the
Einstein--SM continuum limit remain open.
\end{firewall}

\section*{Version V94 append-only record}
\addcontentsline{toc}{section}{Version V94 append-only record}
The complete V93 source was retained byte for byte before this
continuation.  It had \(756544\) bytes, \(22959\) lines, and SHA--256
\[
 \texttt{594C5DABB78C9CCA53CABAC1A3D4209876043EA109325356EC258B05B4077FC3}.
\]
V94 proved Ward-improved one-form and Schur-complement coercivity on a
bounded conformal sector.


\section{V95: mandatory audit after the central-sector Schur theorem}
\label{sec:v17v95}

This version performs the required five-version audit before the
V94 Schur floor is inserted into the Higgs environmental response.

\begin{audit}[Newest active YM and NS checkpoints]
\label{aud:v17v95-companions}
The newest Yang--Mills note inspected was
\begin{center}\footnotesize
\path{Renewal_Yang_Mills_Mass_Gap_Research_Notes_V99.tex},
\end{center}
with \(764947\) bytes, \(20862\) physical lines, and SHA--256
\[
 \texttt{AE0023D070E3CA6DE8C886B9C3BA47EADF569B87986174124F22509FEF39D88B}.
\]
After completing its finite full-torus Wilson-symbol atlas in V97 and
proving in V98 that a fixed symbol floor alone does not determine a
physical mass gap, YM V99 begins a common-regulator compatibility
ledger.  It identifies the concrete fixed-tree weak-norm inequality
for the Wilson remainder and marked vertices as the earliest unpaid
row in its current good/bad branch.  Its running marginal path,
volume-uniform contraction, continuum state, and mass gap remain open.

The newest Navier--Stokes note remains
\begin{center}\footnotesize
\path{Renewal_NS_Stability_Research_Notes_V26.tex},
\end{center}
with \(278283\) bytes, \(5319\) physical lines, and SHA--256
\[
 \texttt{82C887F8C2C69E4F28FAD7C3DC8DE5B9099CB5A4BF431EBA1FFF3E91935978AD}.
\]
Its latest theorem remains the closed-form rank-one Oseen-kernel
derivative calculation, with global stability and regularity open.
\end{audit}

\begin{theorem}[V95 no-import decision]
\label{thm:v17v95-no-import}
Neither companion checkpoint closes a hypothesis of the V94
central-sector conformal response.  No theorem is imported at V95.
\end{theorem}

\begin{proof}
YM V99 is a compatibility audit and a statement of an unproved
fixed-tree inequality.  It contains no Higgs conditional covariance,
conformal one-form resolvent, or outside-sector probability estimate.
NS V26 concerns deterministic fluid kernels and contains none of
these objects.  Their conclusions therefore do not match any open
antecedent in V94.
\end{proof}

The YM audit nevertheless changes the bookkeeping direction.  Any SM
response theorem derived next will display its dependence on the
lattice spacing, volume, running couplings, sector cutoff, and
central-sector radius.  A finite-cutoff floor will not be described as
a cofinal continuum estimate unless those dependencies are uniform on
one regulator family.

\begin{decision}[V96: central-sector environmental response]
\label{dec:v17v95-next}
Insert \(\mu_{R,L}^{\rm Ward}\) into the exact response identity
(79.6).  For a Higgs-norm test supported in a spatial block, use V82 to
split its conformal source into a local part and an exponentially
small tail.  Apply the V94 Schur inverse on the local part and optimize
the truncation radius.  Record every regulator dependence and retain
the indicator of \(\Omega_R\); no outside-sector inference is allowed.
\end{decision}

\begin{firewall}[Claim boundary after V95]
\label{fire:v17v95}
V95 is an audit and bookkeeping correction.  It does not extend V94
beyond \(\Omega_R\), prove full environmental susceptibility, control
the remainder current, construct a continuum measure, establish gauge
cofinality or a physical Einstein action, or prove the Einstein--SM
continuum limit.
\end{firewall}

\section*{Version V95 append-only record}
\addcontentsline{toc}{section}{Version V95 append-only record}
The complete V94 source was retained byte for byte before this
continuation.  It had \(764724\) bytes, \(23206\) lines, and SHA--256
\[
 \texttt{9712832B0F3CDCF19228F7A5F9BD20EF5296705EB730A488EF2F5D98FF71FA3C}.
\]
V95 audited YM V99 and NS V26 and retained the central-sector response
acquisition with an explicit regulator ledger.


\section{V96: correction of the V94 scope and a valid central Dirichlet response}
\label{sec:v17v96}

The V95 acquisition proposed inserting the V94 floor directly into the
full Helffer--Sj\"ostrand inverse (79.13).  That step is not valid as
stated.  The operator \(\mathscr H_U^{(1)}\) acts on functions of the
entire conformal configuration \(r\), whereas the pointwise
multiplication estimate used in V94 holds only where
\(r\in\Omega_R\).  This version corrects the scope before using the
estimate.

\subsection{Erratum to the global wording in V94}

Let
\[
 V_R(r):=2\kappa_gI+B^{\mathsf T}
 \{\mathsf W_J(r)+\mathsf H_H(r)\}B                 \tag{96.1}
\]
be the multiplication part of the one-form operator after the
nonnegative scalar diffusion term is removed.

\begin{proposition}[Correct scope of the V94 form estimate]
\label{prop:v17v96-erratum}
Under the hypotheses of Theorem~\ref{thm:v17v94-coercivity},
\[
 V_R(r)\succeq\mu_{R,L}^{\rm Ward}I
 \qquad\text{for every }r\in\Omega_R.               \tag{96.2}
\]
Consequently,
\[
 \mathcal Q_U[\omega]
 \geq\mu_{R,L}^{\rm Ward}\|\omega\|^2              \tag{96.3}
\]
for one-forms \(\omega\) whose configuration-space support is
contained in \(\Omega_R\).

The unrestricted assertion (94.12), and hence the unrestricted Schur
conclusion (94.15), does not follow unless the conformal law is
supported in \(\Omega_R\) or a separate lower form is proved on its
complement.  Equations (94.12)--(94.18) are therefore superseded by
their Dirichlet interpretation below.
\end{proposition}

\begin{proof}
The proof of V94 is pointwise up to the final integration: (94.3) and
(94.9) require \(|d_e(r)|\leq R\), and (94.14) is exactly (96.2).
Integrating it together with the nonnegative diffusion energy proves
(96.3) for supported forms.

For a general form, its values at \(r\notin\Omega_R\) enter both the
potential and diffusion terms.  A spatial projection \(P\) from
(83.1) acts on vertex-one-form indices and does not restrict the
configuration variable \(r\).  It therefore cannot remove the
uncontrolled complement.  This proves the limitation.
\end{proof}

This correction preserves the Ward diagonal, near/far covariance
split, and local constant (94.11).  It withdraws only the inference
from a central-sector coefficient bound to the inverse of the full
configuration-space operator.

\subsection{The valid Dirichlet realization}

Let \(\mathscr H_{U,R,D}^{(1)}\) be the Friedrichs operator associated
with \(\mathcal Q_U\) on the closure of smooth one-forms compactly
supported in \(\Omega_R\).  This is the Dirichlet realization in
conformal configuration space.

\begin{theorem}[Central-sector Dirichlet coercivity]
\label{thm:v17v96-dirichlet}
If
\[
 \mu:=\mu_{R,L}^{\rm Ward}>0,                       \tag{96.4}
\]
then
\[
\boxed{
 \mathscr H_{U,R,D}^{(1)}\succeq\mu I,
 \qquad
 \| (\mathscr H_{U,R,D}^{(1)})^{-1}\|
 \leq\mu^{-1}.}                                    \tag{96.5}
\]
For every spatial block projection \(P\), the corresponding
Dirichlet inner, outer, and Schur blocks have the same floor \(\mu\).
\end{theorem}

\begin{proof}
Equation (96.3) is the closed quadratic-form inequality defining the
Friedrichs realization, so the spectral theorem gives (96.5).
Compression gives the inner and outer block floors.  The block inverse
argument (94.18), now applied to
\(\mathscr H_{U,R,D}^{(1)}\), gives the Dirichlet Schur floor.
\end{proof}

\subsection{A volume-uniform norm for the response source}

The response vector \(\zeta_t\) in (79.11) also admits a global
operator bound that does not sum its pointwise spatial envelope.

\begin{lemma}[Normalized response-source bound]
\label{lem:v17v96-source}
Under the V61 influence margin, for every deterministic vertex vector
\(t\),
\[
\boxed{
 \sum_e\frac{|\zeta_{t,e}|^2}{w_e}
 \leq16q_H^2C_G^4D\,\|t\|_2^2.}                    \tag{96.6}
\]
Consequently, on \(\Omega_R\),
\[
 \|B^{\mathsf T}\zeta_t\|_2^2
 \leq32q_H^2C_G^4D^2w_R\,\|t\|_2^2.               \tag{96.7}
\]
Both constants are independent of the number of sites.
\end{lemma}

\begin{proof}
For a real edge vector \(u\), put \(Q_u=\sum_eu_eq_e\).  By
(79.11),
\[
 \sum_eu_e\zeta_{t,e}=-\operatorname{Cov}(A_t,Q_u).\tag{96.8}
\]
The global conditional Poincar\'e inequality and Cauchy--Schwarz give
\[
 |\operatorname{Cov}(A_t,Q_u)|
 \leq C_G
  (\mathbb E\|\nabla A_t\|^2)^{1/2}
  (\mathbb E\|\nabla Q_u\|^2)^{1/2}.                \tag{96.9}
\]
Equation (87.2) implies
\[
 \mathbb E\|\nabla A_t\|^2
 \leq4q_HC_G\|t\|_2^2.                              \tag{96.10}
\]
At a site \(x\), the gradient of \(Q_u\) is twice \(y_x\) times a
signed sum of at most \(D\) incident coefficients \(g_{x,e}u_e\).
Cauchy--Schwarz over the incident edges and (87.2) give
\[
 \mathbb E\|\nabla Q_u\|^2
 \leq4Dq_HC_G\sum_ew_eu_e^2.                       \tag{96.11}
\]
Insert (96.10)--(96.11) in (96.9), then take the dual norm with
respect to \(\sum_ew_eu_e^2\).  This proves (96.6).

On \(\Omega_R\), \(w_e\leq w_R\), so (96.6) gives
\(\|\zeta_t\|_2^2\leq16q_H^2C_G^4Dw_R\|t\|^2\).
Finally use \(\|B\|^2\leq2D\), proving (96.7).
\end{proof}

\subsection{Central Dirichlet negative-Sobolev response}

Let
\[
 g_{t,R}:=1_{\Omega_R}B^{\mathsf T}\zeta_t,          \tag{96.12}
\]
viewed as an \(L^2(\nu_U)\) one-form on \(\Omega_R\).

\begin{theorem}[Valid central-sector response bound]
\label{thm:v17v96-central-response}
Under (96.4),
\[
\boxed{
 \left\langle g_{t,R},
 (\mathscr H_{U,R,D}^{(1)})^{-1}g_{t,R}
 \right\rangle
 \leq C_{\rm cent}(R,L)\|t\|_2^2,}                 \tag{96.13}
\]
where
\[
 C_{\rm cent}(R,L)
 :=\frac{32q_H^2C_G^4D^2w_R}
          {\mu_{R,L}^{\rm Ward}}.                  \tag{96.14}
\]
The bound is uniform in finite volume whenever the displayed model
constants and margins are uniform.
\end{theorem}

\begin{proof}
By (96.5), the left side of (96.13) is at most
\(\mu^{-1}\|g_{t,R}\|_2^2\).  Integrate the pointwise bound (96.7)
over \(\Omega_R\); its probability is at most one.  This gives
(96.13)--(96.14).
\end{proof}

Theorem~\ref{thm:v17v96-central-response} is an auxiliary Dirichlet
negative-Sobolev estimate.  It is not equal to
\(\operatorname{Var}_{\nu_U}(\bar A_t)\) in (79.13).  Cutting the
source and imposing a Dirichlet boundary changes the inverse, and no
operator comparison with the full inverse is available while the
outside-sector form is uncontrolled.

\subsection{Regulator ledger}

For a regulator tuple \((a,N,U)\), the proven constant is
\[
 C_{\rm cent}(a,N,U;R,L)
 =\frac{64q_H^2D^2C_G(a)^4\chi_*(a)^2\cosh(2R)}
        {\mu_{R,L}^{\rm Ward}(a,U)}.                \tag{96.15}
\]
It is independent of \(N\) if \(D,q_H,R,L,C_G,\chi_*\), and the
positive floor \(\mu_{R,L}^{\rm Ward}\) are uniform along the chosen
family.  V61 and V94 make those requirements explicit but do not prove
them for the physical running Standard-Model couplings.  Dependence on
the gauge background enters through the requirement that the declared
uniform conditional estimates hold; no gauge averaging has occurred.

\begin{assessment}[What V96 corrects and proves]
\label{ass:v17v96-status}
V96 corrects the configuration-space scope of V94, proves the valid
central Dirichlet Schur theorem, derives a volume-uniform normalized
response-source norm, and obtains the rigorous Dirichlet response
constant (96.14).  It does not identify that auxiliary form with the
full environmental variance.
\end{assessment}

\begin{decision}[V97: outside-sector form and IMS interface]
\label{dec:v17v96-next}
Return to the exact IMS identity (81.8).  Construct a single radial
configuration-space cutoff between \(\Omega_R\) and
\(\Omega_{R+\delta}\), compute its dimension and incidence dependence,
and pair the central Dirichlet floor with the best available lower form
on the complement.  Test whether mismatch growth, bare confinement,
and the Ward/current bounds yield a positive outside floor.  If not,
state the precise outside coefficient deficit rather than comparing
the Dirichlet inverse to the full inverse.
\end{decision}

\begin{firewall}[Claim boundary after V96]
\label{fire:v17v96}
No lower form is proved on \(\Omega_R^c\), so the full resolvent
estimate (79.16), environmental susceptibility, gauge-environment
covariance, remainder-current control, continuum measure, gauge
cofinality, physical Einstein action, and the Einstein--SM continuum
limit remain open.
\end{firewall}

\section*{Version V96 append-only record}
\addcontentsline{toc}{section}{Version V96 append-only record}
The complete V95 source was retained byte for byte before this
continuation.  It had \(768489\) bytes, \(23294\) lines, and SHA--256
\[
 \texttt{4DE19C8F85D905DB5BF94EC7C41637D1B5202281D354818F5C3BF773B53468E9}.
\]
V96 corrected the V94 scope and proved a central Dirichlet response
bound with an explicit regulator ledger.


\section{V97: a volume-safe soft-maximum IMS interface and the strong-mismatch response branch}
\label{sec:v17v97}

V96 leaves the coupling between the central Dirichlet region and its
complement open.  V81 already rules out a product cutoff over all
edges.  We replace that product by one smooth soft maximum.  Its
gradient has a volume-independent norm even though it detects the
largest edge difference.

\subsection{A smooth maximum with bounded configuration gradient}

For \(\beta>0\), define
\[
 \mathfrak m_\beta(d)
 :=\frac1\beta\log
   \left(\sum_{e\in E}2\cosh(\beta d_e)\right).      \tag{97.1}
\]

\begin{lemma}[Soft-maximum comparison and gradient]
\label{lem:v17v97-softmax}
For every finite edge set,
\[
 \max_e|d_e|
 \leq\mathfrak m_\beta(d)
 \leq\max_e|d_e|+\frac{\log(2|E|)}\beta.            \tag{97.2}
\]
Moreover, with \(d=Br\),
\[
\boxed{
 \|\nabla_r\mathfrak m_\beta(Br)\|^2\leq2D.}       \tag{97.3}
\]
The right side is independent of \(|E|\) and \(\beta\).
\end{lemma}

\begin{proof}
The sum in (97.1) is at least
\(e^{\beta\max_e|d_e|}\) and at most
\(2|E|e^{\beta\max_e|d_e|}\), proving (97.2).  Its edge gradient is
\[
 p_e(d):=\partial_{d_e}\mathfrak m_\beta
 =\frac{2\sinh(\beta d_e)}
        {\sum_f2\cosh(\beta d_f)}.                  \tag{97.4}
\]
Thus
\[
 \|p(d)\|_1
 \leq\frac{\sum_e2|\sinh(\beta d_e)|}
              {\sum_e2\cosh(\beta d_e)}\leq1,       \tag{97.5}
\]
so \(\|p\|_2\leq1\).  Since
\(\nabla_r\mathfrak m_\beta=B^{\mathsf T}p\) and
\(\|B\|^2\leq2D\), (97.3) follows.
\end{proof}

Choose smooth \(c_0,c_1:\mathbb R\to[0,1]\) satisfying
\[
 c_0^2+c_1^2=1,
 \quad c_0(s)=1\ (s\leq0),
 \quad c_0(s)=0\ (s\geq1),                          \tag{97.6}
\]
and put
\[
 C_c:=\sup_s\{|c_0'(s)|^2+|c_1'(s)|^2\}<\infty.    \tag{97.7}
\]
For \(R,\delta>0\), define the two-cell partition
\[
 \eta_j(r):=c_j\left(
 \frac{\mathfrak m_\beta(Br)-R}{\delta}
 \right),\qquad j=0,1.                              \tag{97.8}
\]

\begin{theorem}[Volume-safe maximum IMS partition]
\label{thm:v17v97-ims}
The partition (97.8) satisfies
\[
\boxed{
 \Lambda_{\rm max}(r)
 :=\sum_{j=0}^1\|\nabla_r\eta_j(r)\|^2
 \leq\frac{2DC_c}{\delta^2}.}                      \tag{97.9}
\]
The central cutoff is supported in \(\Omega_{R+\delta}\).  If
\(\varepsilon>0\) and
\[
 \beta\geq\frac{\log(2|E|)}{\varepsilon},           \tag{97.10}
\]
then \(\eta_0=1\) throughout \(\Omega_{R-\varepsilon}\).
\end{theorem}

\begin{proof}
The chain rule, (97.3), and (97.7) give (97.9).  If \(\eta_0\ne0\),
then \(\mathfrak m_\beta<R+\delta\), and the first inequality in
(97.2) gives \(\max_e|d_e|<R+\delta\).  If
\(\max_e|d_e|\leq R-\varepsilon\), the second inequality in (97.2)
and (97.10) give \(\mathfrak m_\beta\leq R\), so (97.6) gives
\(\eta_0=1\).
\end{proof}

Unlike the V81 product partition, (97.9) stays fixed in the
thermodynamic limit.  The smoothing parameter needs to grow only
logarithmically with \(|E|\) to retain a fixed inner core, and that
growth does not enter the IMS error.

\subsection{Exact central/outside interface}

Let \(m_{\rm out}(R,\delta,\beta)\) be any verified constant such that
\[
 \mathcal Q_U[\eta_1\omega]
 \geq m_{\rm out}\|\eta_1\omega\|^2                \tag{97.11}
\]
for every form-domain one-form.  The central estimate (96.3), at
radius \(R+\delta\), gives
\[
 \mathcal Q_U[\eta_0\omega]
 \geq\mu_{R+\delta,L}^{\rm Ward}
       \|\eta_0\omega\|^2.                          \tag{97.12}
\]

\begin{corollary}[Two-cell global form criterion]
\label{cor:v17v97-two-cell}
If
\[
 m_*:=\min\{\mu_{R+\delta,L}^{\rm Ward},m_{\rm out}\}
 >\frac{2DC_c}{\delta^2},                            \tag{97.13}
\]
then
\[
 \mathscr H_U^{(1)}
 \succeq\left(m_*-\frac{2DC_c}{\delta^2}\right)I   \tag{97.14}
\]
on the exact-gradient subspace.
\end{corollary}

\begin{proof}
Insert (97.11)--(97.12) and (97.9) in the exact IMS identity (81.8),
then use \(\eta_0^2+\eta_1^2=1\).
\end{proof}

Thus the localization architecture itself is no longer extensive.
The only missing input in (97.13) is an outside form.

\subsection{The strongest currently proved outside branch}

The normalized current-covariance estimate (84.2), together with the
nonnegative diagonal \(\operatorname{diag}(\mathbb ET_e)\) in
(83.9), gives for \(u=B\omega\)
\[
 u^{\mathsf T}(\mathsf W_J+\mathsf H_H)u
 \geq\sum_e
 \{4J-2K_{\rm cur}\chi_e^2\}
 \cosh(2d_e)u_e^2.                                  \tag{97.15}
\]
Define the strong-mismatch surplus
\[
 c_J:=4J-2K_{\rm cur}\chi_*^2.                     \tag{97.16}
\]

\begin{theorem}[Global strong-mismatch one-form floor]
\label{thm:v17v97-strong-floor}
If \(c_J\geq0\), then on the full conformal configuration space
\[
 \mathcal Q_U[\omega]
 \geq2\kappa_g\|\omega\|^2
 +c_J\mathbb E_{\nu_U}
   \sum_e\cosh(2d_e)|(B\omega)_e|^2.                \tag{97.17}
\]
In particular, one may take \(m_{\rm out}=2\kappa_g\) in (97.11),
and the full operator has floor \(2\kappa_g\) even without IMS.
\end{theorem}

\begin{proof}
Use (97.15) in (81.5) and retain the nonnegative diffusion term.
When \(c_J\geq0\), the weighted incidence term is nonnegative and may
be discarded for the last assertion.
\end{proof}

The nonnegative condition in this theorem is the mismatch-only margin
already isolated in V84.  The new point is that its strictly positive
version controls the response source directly in its natural
conformal weight.

\begin{theorem}[Full conformal response on the strong-mismatch branch]
\label{thm:v17v97-full-response}
If \(c_J>0\), then for every deterministic vertex test \(t\),
\[
\boxed{
 \operatorname{Var}_{\nu_U}(\bar A_t)
 \leq C_{\rm env}^{J}\|t\|_2^2,}                   \tag{97.18}
\]
where
\[
\boxed{
 C_{\rm env}^{J}
 :=\frac{32q_H^2C_G^4D\chi_*^2}{c_J}.}              \tag{97.19}
\]
The bound is uniform in finite volume and requires no restriction to
\(\Omega_R\).
\end{theorem}

\begin{proof}
Use the variational identity (79.14), write
\(g=B^{\mathsf T}\zeta_t\), and retain only the second term in
(97.17).  Edgewise completion of the square gives
\[
 \left\langle g,(\mathscr H_U^{(1)})^{-1}g\right\rangle
 \leq\frac1{c_J}\mathbb E_{\nu_U}
       \sum_e\frac{|\zeta_{t,e}|^2}{\cosh(2d_e)}.   \tag{97.20}
\]
Since \(w_e=2\chi_e^2\cosh(2d_e)\),
\[
 \sum_e\frac{|\zeta_{t,e}|^2}{\cosh(2d_e)}
 \leq2\chi_*^2\sum_e\frac{|\zeta_{t,e}|^2}{w_e}.  \tag{97.21}
\]
Apply (96.6) and then the exact identity (79.13), proving
(97.18)--(97.19).
\end{proof}

This is a genuine completion of the conformal conditional-mean
susceptibility on a nonempty mathematical branch.  It is not a
physical cofinal result.  Since
\(K_{\rm cur}=4Dq_HC_G^2\), positivity of (97.16) requires
\[
 J>2Dq_HC_G^2\chi_*^2.                              \tag{97.22}
\]
For the canonical scaling \(J(a)=\bar J a^2\) with a nonvanishing
\(\chi_*\) and a scale-independent \(C_G\), (97.22) eventually fails.
This reproduces the V84 coefficient deficit at the exact response
level.  The soft-maximum IMS theorem remains available, but it needs a
different outside floor to cover the canonical branch.

\begin{assessment}[What V97 closes]
\label{ass:v17v97-status}
V97 proves a two-cell maximum cutoff with a volume-independent IMS
error, thereby repairing the geometric defect of the V81 product
partition.  It also proves the full conformal response estimate
(97.18) on the strong-mismatch branch.  The canonical branch remains
open because its mismatch surplus tends negative.
\end{assessment}

\begin{decision}[V98: canonical outside floor from response-adapted incidence]
\label{dec:v17v97-next}
Keep the soft-maximum partition, but do not ask the vanishing mismatch
coefficient to dominate the entire normalized current covariance.
For the particular source \(B^{\mathsf T}\zeta_t\), decompose edge
currents into the incidence range and cycle space.  Test whether the
bare vertex mass and diffusion control the incidence component while
the mismatch pays only the cycle leakage.  Derive the exact weighted
projection and its smallest singular value before making any
canonical-scale claim.
\end{decision}

\begin{firewall}[Claim boundary after V97]
\label{fire:v17v97}
The full response theorem is conditional on the noncanonical
strong-mismatch surplus \(c_J>0\).  No canonical outside floor,
gauge-environment covariance, remainder-current control, continuum
measure, gauge cofinality, physical Einstein action, or Einstein--SM
continuum limit has been proved.
\end{firewall}

\section*{Version V97 append-only record}
\addcontentsline{toc}{section}{Version V97 append-only record}
The complete V96 source was retained byte for byte before this
continuation.  It had \(777314\) bytes, \(23541\) lines, and SHA--256
\[
 \texttt{91A272E2A240AE6B5D339A08017418FCFE5A705F47AC74D3444270C4FBDE4669}.
\]
V97 proved a volume-safe soft-maximum IMS partition and closed the full
conformal response on the strong-mismatch branch.


\section{V98: weighted incidence compression of the current covariance}
\label{sec:v17v98}

V97 closes the full conformal response when the mismatch dominates the
entire normalized edge-current covariance.  That is stronger than the
response problem requires.  The conformal source and Hessian contain
the incidence map \(B\), so cycle-space covariance directions are
annihilated exactly.  This version constructs the corresponding
weighted projection and replaces \(K_{\rm cur}\) by the optimal
incidence-compressed constant.

\subsection{The weighted incidence/cycle decomposition}

Let
\[
 W(d):=\operatorname{diag}_e(w_e(d)),
 \qquad
 T_d:=W(d)^{1/2}B.                                  \tag{98.1}
\]
On the orthogonal complement of the vertex constants, define
\[
 \Pi_d:=T_d(T_d^*T_d)^\dagger T_d^*,                \tag{98.2}
\]
where \(\dagger\) is the Moore--Penrose inverse.  Then \(\Pi_d\) is
the orthogonal projection in edge space onto
\(\operatorname{Ran}T_d\).  Its complement is the weighted cycle
space.

Normalize the current covariance by
\[
 \widehat{\mathsf C}_H(d)
 :=W(d)^{-1/2}\mathsf C_H(d)W(d)^{-1/2}.             \tag{98.3}
\]
For zero hopping edges, all formulas are read after removing the
identically zero current coordinate.

\begin{definition}[Incidence-compressed covariance constant]
\label{def:v17v98-kinc}
Put
\[
 K_{\rm inc}(d)
 :=\|\Pi_d\widehat{\mathsf C}_H(d)\Pi_d\|_{2\to2}.  \tag{98.4}
\]
Equivalently,
\[
 K_{\rm inc}(d)
 =\sup_{\substack{\omega\perp\ker B\\B\omega\ne0}}
 \frac{(B\omega)^{\mathsf T}\mathsf C_H(B\omega)}
      {(B\omega)^{\mathsf T}W(B\omega)}.            \tag{98.5}
\]
\end{definition}

\begin{theorem}[Exact removal of cycle covariance]
\label{thm:v17v98-cycle-removal}
For every vertex one-form \(\omega\),
\[
 (B\omega)^{\mathsf T}\mathsf C_H(B\omega)
 \leq K_{\rm inc}(d)
       (B\omega)^{\mathsf T}W(B\omega).             \tag{98.6}
\]
Moreover, if \(s_t=W^{-1/2}\zeta_t\), then
\[
 \zeta_t^{\mathsf T}B\omega
 =(\Pi_ds_t)^{\mathsf T}W^{1/2}B\omega.             \tag{98.7}
\]
Thus neither the quadratic perturbation nor the response pairing sees
the weighted cycle component \((I-\Pi_d)s_t\).
\end{theorem}

\begin{proof}
Set \(z=W^{1/2}B\omega=T_d\omega\), so
\(z\in\operatorname{Ran}\Pi_d\).  Then
\[
 (B\omega)^{\mathsf T}\mathsf C_H(B\omega)
 =z^{\mathsf T}\widehat{\mathsf C}_Hz
 =z^{\mathsf T}\Pi_d\widehat{\mathsf C}_H\Pi_dz,  \tag{98.8}
\]
which proves (98.6) and (98.5).  Also
\[
 \zeta_t^{\mathsf T}B\omega=s_t^{\mathsf T}z
 =(\Pi_ds_t)^{\mathsf T}z,                      \tag{98.9}
\]
because \((I-\Pi_d)z=0\).  This proves (98.7).
\end{proof}

The decision after V97 anticipated a possible cycle leakage.  Theorem
\ref{thm:v17v98-cycle-removal} proves that there is none at this
linear-response level; the only relevant covariance is the compression
(98.4).

\subsection{An explicit analytic ceiling for the compression}

For an edge vector \(u\), define the weighted endpoint divergence
\[
 (G_du)_x
 :=\sum_{e\ni x}\sigma_{x,e}g_{x,e}u_e,            \tag{98.10}
\]
where the sign is chosen so that
\[
 Q_u=\sum_eu_eq_e=\sum_x(G_du)_xA_x.                \tag{98.11}
\]
Set
\[
 \sigma_{\rm inc}(d)^2
 :=\sup_{\substack{u\in\operatorname{Ran}B\\u\ne0}}
 \frac{\|G_du\|_2^2}{u^{\mathsf T}W(d)u}.           \tag{98.12}
\]

\begin{proposition}[Incidence-symbol covariance ceiling]
\label{prop:v17v98-ceiling}
Under the V61 influence margin,
\[
\boxed{
 K_{\rm inc}(d)
 \leq4q_HC_G^2\sigma_{\rm inc}(d)^2
 \leq4Dq_HC_G^2=K_{\rm cur}.}                      \tag{98.13}
\]
The first inequality can be strict even when the global edge bound is
not.
\end{proposition}

\begin{proof}
From (98.11),
\(\nabla_{y_x}Q_u=2(G_du)_xy_x\).  The global Poincar\'e inequality
and (87.2) give
\[
 \operatorname{Var}(Q_u)
 \leq4q_HC_G^2\|G_du\|_2^2.                         \tag{98.14}
\]
Restrict to \(u\in\operatorname{Ran}B\), divide by \(u^TWu\), and
use (98.5), proving the first inequality.  At each site,
\[
 |(G_du)_x|^2
 \leq D\sum_{e\ni x}g_{x,e}^2u_e^2.                \tag{98.15}
\]
Summing sites gives
\(\|G_du\|^2\leq D\sum_ew_eu_e^2\), so
\(\sigma_{\rm inc}^2\leq D\) and the second inequality follows.
\end{proof}

Formula (98.12) is a weighted finite-dimensional incidence symbol.  It
can be evaluated or bounded without controlling the cycle sector of
\(\widehat{\mathsf C}_H\).

\subsection{The optimal mismatch surplus seen by conformal response}

Let
\[
 \Theta_{\rm inc}
 :=\inf_d\left\{
 \min_e\frac{2J}{\chi_e^2}-K_{\rm inc}(d)
 \right\}.                                         \tag{98.16}
\]
The infimum is over the conformal configurations and fixed gauge
backgrounds in the declared regulator.  Indeed,
\[
 \sum_e4J\cosh(2d_e)u_e^2
 \geq\left(\min_e\frac{2J}{\chi_e^2}\right)u^TWu.  \tag{98.17}
\]

\begin{theorem}[Incidence-compressed full response]
\label{thm:v17v98-response}
If
\[
\boxed{\Theta_{\rm inc}>0,}                         \tag{98.18}
\]
then the full conformal conditional-mean response satisfies
\[
\boxed{
 \operatorname{Var}_{\nu_U}(\bar A_t)
 \leq\frac{16q_H^2C_G^4D}{\Theta_{\rm inc}}
       \|t\|_2^2.}                                  \tag{98.19}
\]
The bound is uniform in volume whenever \(\Theta_{\rm inc}\) and the
V61 constants are uniform.
\end{theorem}

\begin{proof}
Use the nonnegative diagonal
\(\operatorname{diag}(\mathbb ET_e)\) in (83.9), then combine
(98.6), (98.16), and (98.17).  The one-form form obeys
\[
 \mathcal Q_U[\omega]
 \geq\Theta_{\rm inc}\,
 \mathbb E\|W^{1/2}B\omega\|^2.                    \tag{98.20}
\]
In the variational identity (79.14), use (98.7) and complete the square
on \(\operatorname{Ran}\Pi_d\):
\[
 \left\langle B^T\zeta_t,
  (\mathscr H_U^{(1)})^{-1}B^T\zeta_t\right\rangle
 \leq\Theta_{\rm inc}^{-1}
       \mathbb E\|\Pi_ds_t\|^2.                    \tag{98.21}
\]
Since orthogonal projection is contractive, (96.6) bounds the last
expectation by
\(16q_H^2C_G^4D\|t\|^2\).  Apply (79.13).
\end{proof}

The V97 strong-mismatch theorem is recovered by replacing
\(K_{\rm inc}\) with \(K_{\rm cur}\) and
\(\min_e2J/\chi_e^2\) with \(2J/\chi_*^2\).  Hence V98 is never weaker
and can be strictly stronger.

\subsection{Exact canonical scale demand}

For \(J(a)=\bar Ja^2\), condition (98.18) requires
\[
 \sup_dK_{\rm inc}(a,d)
 <\min_e\frac{2\bar Ja^2}{\chi_e(a)^2}.             \tag{98.22}
\]
If \(\chi_e(a)\to\chi_e^0>0\), a sufficient cofinal acquisition is
therefore
\[
\boxed{
 \sup_dK_{\rm inc}(a,d)=O(a^2)
 \quad\text{with leading coefficient below}
 \quad\min_e\frac{2\bar J}{(\chi_e^0)^2}.}          \tag{98.23}
\]
No estimate in V61--V98 proves this decay: (98.13) is scale
independent.  Conversely, cycle-space improvements alone cannot help,
because the cycle component has already disappeared exactly in
(98.6)--(98.7).  The canonical bottleneck is decay of the covariance
compressed to the physical incidence range.

\begin{assessment}[What V98 closes]
\label{ass:v17v98-status}
V98 proves exact weighted cycle removal, constructs the optimal
incidence-compressed current covariance, derives its analytic ceiling,
and proves the full response theorem (98.19) under the strictly weaker
surplus (98.18).  It identifies the exact canonical scale demand
(98.23), which remains unproved.
\end{assessment}

\begin{decision}[V99: Ward cancellation inside the incidence compression]
\label{dec:v17v98-next}
Use the exact Ward-renormalized decomposition (86.16) before taking the
incidence compression.  Express
\(B^T\mathsf H_HB\) as a vertex divergence form whose same-edge part is
\(h_e^{\rm Ward}\), and test whether the shell bounds of V92--V93
produce an \(O(a^2)\) compressed negative part even though the full
edge covariance is \(O(1)\).  If the local Ward constants remain
\(O(1)\), record the precise term preventing (98.23) and move upstream
to the running Higgs normalization.
\end{decision}

\begin{firewall}[Claim boundary after V98]
\label{fire:v17v98}
No \(O(a^2)\) incidence-compressed covariance estimate has been
proved.  Thus the canonical full response, gauge-environment
covariance, remainder-current control, continuum measure, gauge
cofinality, physical Einstein action, and the Einstein--SM continuum
limit remain open.
\end{firewall}

\section*{Version V98 append-only record}
\addcontentsline{toc}{section}{Version V98 append-only record}
The complete V97 source was retained byte for byte before this
continuation.  It had \(786178\) bytes, \(23811\) lines, and SHA--256
\[
 \texttt{268D737E653603B0A582FE3591EB805538DD7B5205A8C53BC5F8805E3434E647}.
\]
V98 removed the cycle sector exactly and replaced the full current
constant by its optimal weighted incidence compression.


\section{V99: the optimal Ward-incidence negative constant and its canonical scale demand}
\label{sec:v17v99}

V98 compresses the full current covariance to the incidence range but
still discards the positive mean matrix
\(\operatorname{diag}(\mathbb ET_e)\).  The Ward identity shows that
this discard loses the same-edge cancellation.  We now retain that
cancellation before taking the incidence compression.

\subsection{The exact negative Higgs operator on the physical range}

Recall
\[
 \mathsf H_H
 =\operatorname{diag}(h_e^{\rm Ward})
  -\mathsf C_H^{\rm off}                             \tag{99.1}
\]
from (86.16).  Define the normalized negative Higgs matrix
\[
 \widehat{\mathsf N}_H(d)
 :=W(d)^{-1/2}\{-\mathsf H_H(d)\}W(d)^{-1/2}         \tag{99.2}
\]
and its incidence compression
\[
\boxed{
 K_{\rm WI}(d)
 :=\max\left\{0,
 \lambda_{\max}
 [\Pi_d\widehat{\mathsf N}_H(d)\Pi_d]
 \right\}.}                                        \tag{99.3}
\]
The subscript WI denotes Ward--incidence.

\begin{theorem}[Optimal negative Higgs constant seen by conformal modes]
\label{thm:v17v99-optimal}
For every vertex one-form \(\omega\),
\[
\boxed{
 (B\omega)^{\mathsf T}\mathsf H_H(B\omega)
 \geq-K_{\rm WI}(d)
       (B\omega)^{\mathsf T}W(B\omega).}            \tag{99.4}
\]
The constant \(K_{\rm WI}(d)\) is the smallest nonnegative number for
which (99.4) holds.  Equivalently,
\[
 K_{\rm WI}(d)
 =\max\left\{0,
 \sup_{\substack{\omega\perp\ker B\\B\omega\ne0}}
 \frac{-(B\omega)^T\mathsf H_H(B\omega)}
      {(B\omega)^TW(B\omega)}\right\}.              \tag{99.5}
\]
\end{theorem}

\begin{proof}
Set \(z=W^{1/2}B\omega\in\operatorname{Ran}\Pi_d\).  Then
\[
 -(B\omega)^T\mathsf H_H(B\omega)
 =z^T\Pi_d\widehat{\mathsf N}_H\Pi_dz.             \tag{99.6}
\]
The Rayleigh--Ritz principle gives (99.4)--(99.5), including
optimality.
\end{proof}

Because
\[
 -\mathsf H_H
 =\mathsf C_H-\operatorname{diag}(\mathbb ET_e)     \tag{99.7}
\]
and the subtracted matrix is positive semidefinite,
\[
\boxed{K_{\rm WI}(d)\leq K_{\rm inc}(d).}           \tag{99.8}
\]
Thus the Ward-incidence constant is never worse than the V98 current
constant and is strictly smaller whenever the mean diagonal improves
the top compressed Rayleigh quotient.

\subsection{A rigorous bounded-sector ledger}

Let
\[
 \chi_{\min}:=\inf_e\chi_e>0.                       \tag{99.9}
\]
On \(\Omega_R\), one has \(w_e\geq2\chi_{\min}^2\).
Define the incidence-compressed off-diagonal norm
\[
 K_{\rm off,inc}(d)
 :=\|\Pi_dW^{-1/2}\mathsf C_H^{\rm off}
             W^{-1/2}\Pi_d\|.                      \tag{99.10}
\]

\begin{proposition}[Ward-incidence upper ledger]
\label{prop:v17v99-ledger}
On \(\Omega_R\),
\[
\boxed{
 K_{\rm WI}(d)
 \leq K_{\rm off,inc}(d)
 +\frac{\mathcal W_*^{\rm best}(R)}
        {2\chi_{\min}^2}.}                          \tag{99.11}
\]
Furthermore,
\[
 K_{\rm off,inc}(d)
 \leq K_{\rm off}^{\rm best}(L),                   \tag{99.12}
\]
so
\[
 K_{\rm WI}(d)
 \leq K_{\rm off}^{\rm best}(L)
 +\frac{\mathcal W_*^{\rm best}(R)}
        {2\chi_{\min}^2}.                           \tag{99.13}
\]
\end{proposition}

\begin{proof}
By (99.1)--(99.2), the negative normalized matrix is the normalized
off-diagonal covariance minus
\(W^{-1/2}\operatorname{diag}(h^{\rm Ward})W^{-1/2}\).
Equation (94.3) bounds the positive part of the latter negative term by
\(\mathcal W_*^{\rm best}(R)/(2\chi_{\min}^2)\).
Compression is contractive in operator norm, proving (99.11).
Equation (94.9) bounds the uncompressed normalized off-diagonal norm by
\(K_{\rm off}^{\rm best}(L)\), proving (99.12)--(99.13).
\end{proof}

The two terms in (99.13) have distinct meanings.  The first is a
transported covariance between different edges; the second is the
same-edge Ward deficit after the exact cancellation of V86.  Any
canonical-scale proof based on this ledger must improve both, rather
than only the full covariance norm.

\subsection{The best current full-response criterion}

Define
\[
 \Theta_{\rm WI}
 :=\inf_d\left\{
 \min_e\frac{2J}{\chi_e^2}-K_{\rm WI}(d)
 \right\}.                                         \tag{99.14}
\]

\begin{theorem}[Ward-incidence full conformal response]
\label{thm:v17v99-response}
If
\[
\boxed{\Theta_{\rm WI}>0,}                          \tag{99.15}
\]
then
\[
\boxed{
 \operatorname{Var}_{\nu_U}(\bar A_t)
 \leq\frac{16q_H^2C_G^4D}{\Theta_{\rm WI}}
       \|t\|_2^2.}                                  \tag{99.16}
\]
This criterion is no stronger than (98.18), because
\(\Theta_{\rm WI}\geq\Theta_{\rm inc}\).
\end{theorem}

\begin{proof}
Equations (99.4) and (98.17) give
\[
 \mathcal Q_U[\omega]
 \geq\Theta_{\rm WI}
 \mathbb E\|W^{1/2}B\omega\|^2.                    \tag{99.17}
\]
The source pairing still satisfies the exact cycle removal (98.7).
Complete the square as in (98.21), apply (96.6), and then use (79.13).
The comparison with \(\Theta_{\rm inc}\) follows from (99.8).
\end{proof}

Theorem~\ref{thm:v17v99-response} is optimal within a proof that uses
only a uniform weighted lower bound for the Higgs multiplication
operator on \(\operatorname{Ran}B\): any smaller negative constant
would contradict the Rayleigh quotient (99.5).

\subsection{Canonical scaling test}

For \(J(a)=\bar Ja^2\) and
\(\chi_e(a)\to\chi_e^0>0\), condition (99.15) requires
\[
\boxed{
 \sup_dK_{\rm WI}(a,d)=O(a^2),}                     \tag{99.18}
\]
with a leading coefficient below
\(\min_e2\bar J/(\chi_e^0)^2\).  The current analytic ledger does not
prove (99.18).  Under scale-independent V61 constants and fixed
\(R,L\), the right side of (99.13) is only \(O(1)\):
\[
 K_{\rm off}^{\rm best}(L)=O(1),
 \qquad
 \frac{\mathcal W_*^{\rm best}(R)}
      {2\chi_{\min}^2}=O(1).                        \tag{99.19}
\]
Equation (99.19) is a limitation of the proved upper bound, not a lower
bound on the exact constant.  It therefore does not disprove a hidden
Ward cancellation of order \(a^2\); it proves that V82--V94 do not
establish one.

The unresolved scale is already visible in the one-site law.  The
constants \(C_0,C_G\) were made independent of the vanishing lattice
mass, while the geometric hopping \(\chi_e\) was retained at order
one.  Those choices are appropriate for a high-temperature gap but
contain no small parameter that could force either term in (99.19) to
be \(O(a^2)\).  The next step must therefore inspect the running
normalization of the Higgs hopping and of the conformal mismatch on one
common regulator family.

\begin{assessment}[What V99 closes]
\label{ass:v17v99-status}
V99 constructs the optimal negative Higgs constant on the physical
weighted incidence range, proves it improves the full-current
compression, derives the strongest current response criterion, and
shows precisely why the existing Ward and off-diagonal ledgers do not
verify its canonical \(O(a^2)\) requirement.
\end{assessment}

\begin{decision}[V100: common-regulator Higgs normalization ledger]
\label{dec:v17v99-next}
At V100 perform the mandatory companion audit and then return to the
source actions and scaling tables.  Put \(J(a),\chi_e(a),\lambda_H(a),
C_0(a),C_G(a)\), and the field normalization on one regulator tuple.
Determine whether the order-one \(\chi_e\) used in V59--V99 is a
dimensionless hopping after field rescaling or a physical coefficient
carrying powers of \(a\).  Derive the induced scaling of
\(K_{\rm WI}\) without assuming the answer.  Then freeze this lineage
at \(\texttt{V100\_f17}\) and begin the next compact V1 with a proof
status summary.
\end{decision}

\begin{firewall}[Claim boundary after V99]
\label{fire:v17v99}
The canonical \(O(a^2)\) Ward-incidence estimate is unproved.  Hence
the canonical full response, gauge-environment covariance,
remainder-current control, continuum measure, gauge cofinality,
physical Einstein action, and the Einstein--SM continuum limit remain
open.
\end{firewall}

\section*{Version V99 append-only record}
\addcontentsline{toc}{section}{Version V99 append-only record}
The complete V98 source was retained byte for byte before this
continuation.  It had \(794841\) bytes, \(24077\) lines, and SHA--256
\[
 \texttt{C6299AEDC35803EA9DEA062D0F8C2EED574C0981914BC87A431AC3DAEA0631F9}.
\]
V99 constructed the optimal Ward-incidence negative constant and its
canonical scale demand.


\section{V100: common-regulator Higgs scaling and the massless incidence gate}
\label{sec:v17v100}

This terminal version of the seventeenth lineage performs the required
companion audit, then places every coefficient entering the V99
Ward--incidence criterion on the canonical four-dimensional regulator
family.  The resulting scale reduction is exact at fixed finite
volume: the required \(O(a^2)\) estimate is equivalent to positivity of
the massless incidence Hessian plus a uniform mass-response bound.

\subsection{Mandatory V100 companion audit}

\begin{audit}[Newest active YM and NS notes]
\label{aud:v17v100-companions}
The newest Yang--Mills note inspected was
\begin{center}\footnotesize
\path{Renewal_Yang_Mills_Mass_Gap_Research_Notes_V99.tex},
\end{center}
with \(772429\) bytes, \(21070\) physical lines, and SHA--256
\[
 \texttt{AAD95349F5132425E2A06A358DA6DEE405BBA0884ADC02986DF58A904261A457}.
\]
YM V99 identifies its concrete fixed-tree strong-to-weak inequality as
the first unpaid post-atlas row.  Conditional on that inequality, it
proves that a tree constant growing like \(\beta^p\) closes the local
activity branch for \(p<1/5\), or at \(p=1/5\) under an explicit
critical constant.  It does not construct the inequality, the running
path \(\beta(a)\), or a continuum state.

The newest Navier--Stokes note remains
\begin{center}\footnotesize
\path{Renewal_NS_Stability_Research_Notes_V26.tex},
\end{center}
with \(278283\) bytes, \(5319\) physical lines, and SHA--256
\[
 \texttt{82C887F8C2C69E4F28FAD7C3DC8DE5B9099CB5A4BF431EBA1FFF3E91935978AD}.
\]
Its latest result remains the rank-one Oseen-kernel derivative norm
calculation, with global regularity open.
\end{audit}

\begin{theorem}[V100 no-import decision]
\label{thm:v17v100-no-import}
Neither companion theorem determines the massless Higgs incidence
Hessian or its mass derivative.  No companion result is imported.
\end{theorem}

\begin{proof}
The YM theorem is a conditional asymptotic estimate for a gauge-polymer
tree constant.  It contains no rescaled Higgs Gibbs covariance or
conformal Hessian.  The NS theorem concerns a deterministic fluid
kernel.  Neither conclusion supplies an antecedent of (99.18).
\end{proof}

\subsection{One canonical regulator tuple}

Let
\[
 \mathfrak R_a=(G_a,U_a,d_a,y_a;
 \bar\kappa_g,\bar J,\bar\chi,\bar m_H^2,
 \bar\lambda_H),
 \qquad G_a=(a\mathbb Z/L\mathbb Z)^4.              \tag{100.1}
\]
The canonical normalization proved in V44 is
\[
\begin{gathered}
 y_x=aH(ax),
 \qquad \chi_e(a)=\bar\chi_e,
 \qquad \lambda_H(a)=\bar\lambda_H,                \tag{100.2}\\
 m_H^2(a)=\bar m_H^2a^2,
 \qquad J(a)=\bar Ja^2,
 \qquad \kappa_g(a)=\bar\kappa_ga^4.               \tag{100.3}
\end{gathered}
\]
Thus \(\chi_e\) is an order-one dimensionless hopping coefficient
after the field rescaling \(y=aH\); it is not a physical coefficient
whose missing power of \(a\) may be inserted later.  Equations
(20.5)--(20.6) independently show that the geometric-mean physical
edge weight becomes exactly this dimensionless hopping form in the
rescaled variables.

The one-site and assembled gap constants are
\[
 C_0(a)=C_{\rm LC}
 \sqrt{\frac{q_H}{2\bar\lambda_H}},
 \qquad
 C_G(a)=\frac{C_0(a)}{(1-r_H(a))^2},                \tag{100.4}
\]
with
\[
 r_H(a)=\|M(a)\|,
 \qquad M_{xz}(a)=2C_0(a)\bar\chi_{xz}.             \tag{100.5}
\]
Hence these constants are order one, provided the same uniform
high-temperature margin \(\sup_ar_H(a)<1\) holds.  The normalized edge
weight
\[
 w_e(d)=2\bar\chi_e^2\cosh(2d_e)                   \tag{100.6}
\]
is also independent of \(a\).  No factor in V61, V83, or V99 therefore
forces the normalized Higgs Hessian to vanish with \(a\).

\begin{record}[Canonical scaling status]
\label{rec:v17v100-ledger}
On the tuple (100.1), the proved orders are:
\begin{enumerate}[label=\textup{(\roman*)},leftmargin=3.4em]
\item \(J=O(a^2)\), \(m_H^2=O(a^2)\), and
      \(\kappa_g=O(a^4)\);
\item \(\chi_e,\lambda_H,C_0,C_G,w_e=O(1)\) under the V61 margin;
\item the existing upper bounds for \(K_{\rm WI}\) are \(O(1)\);
\item the V99 response surplus needs
      \(K_{\rm WI}=O(a^2)\) with a sufficiently small leading
      coefficient.
\end{enumerate}
Items (iii)--(iv) do not contradict one another: item (iii) describes
the strength of the proved ceiling, while item (iv) is the unproved
sharper cancellation needed by the canonical branch.
\end{record}

\subsection{Analytic massless limit at fixed regulator volume}

Fix a finite graph, \(d\), and \(U\), and write
\[
 \mu=m_H^2,
 \qquad
 \mathcal A(y):=\sum_x\|y_x\|^2.                   \tag{100.7}
\]
The conditional Higgs law has density
\[
 d\pi_\mu=Z_\mu^{-1}
 e^{-S_0(y)-\mu\mathcal A(y)}\,dy,                  \tag{100.8}
\]
where \(S_0\) contains the positive quartic and geometric hopping
terms but no mass.  Quartic confinement makes every polynomial moment
finite uniformly for \(\mu\) in a compact interval about zero.

\begin{lemma}[Mass differentiation]
\label{lem:v17v100-mass-derivative}
For every polynomial observable \(F\),
\[
\boxed{
 \frac d{d\mu}\mathbb E_{\pi_\mu}F
 =-\operatorname{Cov}_{\pi_\mu}(F,\mathcal A).}     \tag{100.9}
\]
The expectation is real analytic in \(\mu\) near zero.
\end{lemma}

\begin{proof}
Differentiate the numerator and denominator in (100.8), giving
(100.9).  On every compact \(\mu\)-interval, the signed quadratic term
is absorbed into half the positive quartic by Young's inequality.
The remaining quartic tail dominates every derivative, which is a
polynomial times the same density.  Dominated differentiation to all
orders proves analyticity.
\end{proof}

Let \(\mathsf H_H(\mu;d,U)\) be the Higgs free-energy Hessian and let
\(K_{\rm WI}(\mu;d,U)\) be (99.3).  The matrices \(W(d)\) and
\(\Pi_d\) do not depend on \(\mu\).

\begin{theorem}[Finite-volume massless Ward-incidence expansion]
\label{thm:v17v100-massless-expansion}
For every fixed finite graph and fixed \((d,U)\), there is a finite
constant \(L_{d,U}\) and a neighborhood of \(\mu=0\) such that
\[
 \|\mathsf H_H(\mu;d,U)-\mathsf H_H(0;d,U)\|
 \leq L_{d,U}|\mu|,                                 \tag{100.10}
\]
and
\[
\boxed{
 |K_{\rm WI}(\mu;d,U)-K_{\rm WI}(0;d,U)|
 \leq\frac{L_{d,U}}{\min_ew_e(d)}|\mu|.}           \tag{100.11}
\]
Consequently, under \(\mu(a)=\bar m_H^2a^2\),
\[
 K_{\rm WI}(a;d,U)
 =K_{\rm WI}(0;d,U)+O_{d,U}(a^2).                  \tag{100.12}
\]
\end{theorem}

\begin{proof}
Every entry of \(\mathsf H_H\) is a finite sum of expectations and
covariances of the polynomial observables \(q_e,T_e\).  Lemma
\ref{lem:v17v100-mass-derivative} makes each entry analytic in \(\mu\),
so its derivative is bounded on a sufficiently small compact interval.
Finite dimensionality gives (100.10).

Conjugation by \(W^{-1/2}\) changes the operator-norm difference by at
most the factor \((\min_ew_e)^{-1}\); compression by \(\Pi_d\) is
contractive.  Weyl's eigenvalue inequality and the one-Lipschitz map
\(x\mapsto\max\{0,x\}\) give (100.11).  Insert
\(\mu=\bar m_H^2a^2\) to obtain (100.12).
\end{proof}

\begin{corollary}[Exact fixed-volume massless gate]
\label{cor:v17v100-massless-gate}
At fixed finite graph and fixed \((d,U)\),
\[
 K_{\rm WI}(a;d,U)=O(a^2)                           \tag{100.13}
\]
if and only if
\[
\boxed{
 K_{\rm WI}(0;d,U)=0,
 \quad\text{equivalently}\quad
 B^T\mathsf H_H(0;d,U)B\succeq0.}                  \tag{100.14}
\]
\end{corollary}

\begin{proof}
If (100.14) holds, (100.11) and nonnegativity of
\(K_{\rm WI}\) give (100.13).  Conversely, (100.13) and continuity in
(100.12) force \(K_{\rm WI}(0;d,U)=0\).  By the Rayleigh quotient
(99.5), this is equivalent to nonnegativity of the massless Higgs
Hessian on \(\operatorname{Ran}B\).
\end{proof}

The equivalence is local in regulator volume.  The cofinal version
needs two additional uniform statements:
\[
 B^T\mathsf H_H(0;d,U)B\succeq0
 \quad\text{for every regulator and every }(d,U),   \tag{100.15}
\]
and
\[
 \sup_{a,d,U}
 \frac{\|\mathsf H_H(\bar m_H^2a^2;d,U)
             -\mathsf H_H(0;d,U)\|}{a^2}<\infty.   \tag{100.16}
\]
Neither statement follows from finite-dimensional analyticity because
its constant \(L_{d,U}\) may grow with volume or the conformal field.

\subsection{The leading-coefficient test}

If (100.15)--(100.16) hold with a uniform normalized Lipschitz constant
\(L_{\rm WI}\), then
\[
 \sup_{d,U}K_{\rm WI}(a;d,U)
 \leq L_{\rm WI}|\bar m_H^2|a^2.                   \tag{100.17}
\]
The V99 surplus is positive for all small \(a\) whenever
\[
\boxed{
 L_{\rm WI}|\bar m_H^2|
 <\min_e\frac{2\bar J}{\bar\chi_e^2}.}              \tag{100.18}
\]
Under this condition, (99.16) has a response constant of order
\(a^{-2}\) in lattice normalization.  Whether that order is sufficient
for a continuum-smeared observable depends on the exact test-vector
normalization and is a separate calculation; it is not called a
continuum susceptibility here.

The massless positivity in (100.15) is the graph-level version of the
scalar monotonicity problem reached in V73--V78.  V86--V99 supplied
Ward identities, local error ledgers, incidence compression, and
conditional branches, but did not prove this global massless sign.
The common-regulator audit shows that the sign cannot be replaced by a
missing factor of \(a\) in \(\chi_e\).

\begin{assessment}[What V100 closes]
\label{ass:v17v100-status}
V100 fixes the canonical normalization on one regulator tuple, proves
the exact finite-volume expansion
\(K_{\rm WI}(a)=K_{\rm WI}(0)+O(a^2)\), and reduces the canonical
response problem to uniform massless incidence positivity plus the
uniform derivative and leading-coefficient tests (100.16)--(100.18).
It also proves that the order-one hopping used throughout the current
lineage is the correct dimensionless coefficient after \(y=aH\).
\end{assessment}

\begin{decision}[Next compact lineage: massless incidence positivity]
\label{dec:v17v100-next}
Freeze this file as \(\texttt{V100\_f17}\).  Restart at contextual V1
with a concise proof-status map.  The first new acquisition is a
massless graph theorem for (100.15): derive the Hessian under an
arbitrary incidence perturbation \(d\mapsto d+tB\omega\), use the
geometric hopping-square structure and the radial Ward identity, and
test convexity of the massless Higgs free energy on trees before adding
cycle edges by interpolation.  If tree convexity fails, retain an exact
finite counterexample and move upstream to the required counterterm or
running normalization.
\end{decision}

\begin{firewall}[Claim boundary after V100]
\label{fire:v17v100}
Uniform massless incidence positivity and the uniform mass derivative
are unproved.  Thus the canonical full response, gauge-environment
covariance, remainder-current control, continuum measure, gauge
cofinality, physical Einstein action, and the Einstein--SM continuum
limit remain open.
\end{firewall}

\section*{Version V100 append-only record}
\addcontentsline{toc}{section}{Version V100 append-only record}
The complete V99 source was retained byte for byte before this
continuation.  It had \(803115\) bytes, \(24326\) lines, and SHA--256
\[
 \texttt{021CF4027DEF997576945BFC5B2B5C831E0B12B6D272800FF556BB9E646C0739}.
\]
V100 proved the canonical massless Ward-incidence reduction and closes
the seventeenth active lineage.

\end{document}
